\documentclass[a4paper,11pt]{article}

\usepackage{amsmath}
\usepackage{amssymb}
\usepackage{amsthm}
\usepackage{hyperref}
\usepackage[hmargin=2cm]{geometry}
\usepackage{fvextra}
\DefineVerbatimEnvironment{Verbatim}{Verbatim}{breaklines,fontsize=\small}
\fvset{breaklines}

\usepackage{microtype}
\emergencystretch=3em

\newtheorem{theorem}{Theorem}
\newtheorem{lemma}[theorem]{Lemma}
\newtheorem{definition}[theorem]{Definition}
\newtheorem{proposition}[theorem]{Proposition}
\newtheorem{remark}[theorem]{Remark}

\title{lattice-system Formalization: \\
       A Mathematical Guide to the Lean 4 Code}
\author{lattice-system project}
\date{2026-04-15}

\begin{document}
\maketitle

\begin{abstract}
This document is a guide to reading the Lean~4 source code of the
\texttt{lattice-system} project. The project generalizes the earlier
\texttt{ising-model} project toward a broad class of many-body
systems on graphs (classical spin systems, quantum spin systems,
Hubbard, BCS, CAR-algebraic, and eventually lattice QCD). Despite
its name, the underlying combinatorial datum is a graph
$(\Lambda, E_\Lambda)$ rather than a specific lattice; concrete
lattices appear as graph instances. The current scope is finite
graphs and finite-volume quantum spin systems, but a major
long-term goal is the \emph{infinite-volume / thermodynamic limit},
which intrinsically uses infinite graphs. This guide explains which
file contains what, which textbook material each formalization
corresponds to, and what is and is not yet proved.
\end{abstract}

\tableofcontents

%======================================================================
\section{Graph-centric framework}
%======================================================================

Despite the name, the primary combinatorial datum underlying
Heisenberg, Ising, and Hubbard in this library is a
\emph{graph} $(\Lambda, E_\Lambda)$, not a specific lattice.
Concrete lattices (1D chain, 2D / 3D grids, torus, infinite
$\mathbb{Z}$ chain, $\mathbb{Z}^d$ grid, $\ldots$) appear as
instances. The framework supports both finite-volume work
(with $[\mathrm{Fintype}~\Lambda]$ for traces and finite sums)
and the infinite-volume / thermodynamic limit work that is a
major long-term goal.

\subsection{The \texttt{couplingOf} bridge}

The canonical pairwise-coupling function associated with a
\texttt{SimpleGraph G} and uniform edge weight $J \in \mathbb{C}$
lives in \texttt{LatticeSystem.Lattice.Graph}:

\begin{align*}
  \texttt{couplingOf}~G~J~x~y \;=\; \begin{cases}
    J & \text{if } G.\mathrm{Adj}~x~y, \\
    0 & \text{otherwise.}
  \end{cases}
\end{align*}

This function is symmetric (because graph adjacency is),
vanishes on the diagonal (because \texttt{SimpleGraph} has no
self-loops), and is real whenever $J$ is real
(\texttt{couplingOf\_real}).

\subsection{Named Hamiltonian wrappers on graphs}

Three named wrappers package the graph-form Hamiltonians:

\begin{itemize}
  \item \texttt{isingHamiltonianOnGraph G J h} $=$
    $\sum_{x,y} (\texttt{couplingOf}~G~J~x~y) \cdot \hat{\sigma}^z_x \hat{\sigma}^z_y
     - h \sum_x \hat{\sigma}^x_x$.
  \item \texttt{heisenbergHamiltonianOnGraph G J} $=$
    $\sum_{x,y} (\texttt{couplingOf}~G~J~x~y) \cdot \hat{S}_x \cdot \hat{S}_y$.
  \item \texttt{hubbardHamiltonianOnGraph N G J U} $=$
    $\sum_\sigma \sum_{x,y} (\texttt{couplingOf}~G~J~x~y) \cdot c_{x,\sigma}^\dagger c_{y,\sigma}
     + U \sum_i n_{i\uparrow} n_{i\downarrow}$.
\end{itemize}

Each carries Hermiticity, charge/spin conservation, and Gibbs
state theorems in the \emph{generic} form; concrete lattice
instances (chain, torus, cubic) are corollaries obtained by
plugging in \texttt{pathGraph}, \texttt{cycleGraph}, and
box products thereof.

\subsection{Chain $\leftrightarrow$ graph bridges}

\texttt{openChainCoupling}, \texttt{periodicChainCoupling},
\texttt{squareLatticeCoupling}, \texttt{squareTorusCoupling}, and
\texttt{cubicLatticeCoupling} are primary \texttt{couplingOf}
instances; their bridge theorems hold by \texttt{rfl}.

For the chain Hamiltonians:

\begin{itemize}
  \item \texttt{quantumIsingHamiltonian\_eq\_isingHamiltonianGeneric}:
    the 1D Ising chain equals
    \texttt{isingHamiltonianGeneric (couplingOf (pathGraph (N+1)) (-J/2)) h}.
    The factor $1/2$ absorbs the symmetric double-sum convention.
  \item \texttt{openChainHeisenbergGibbsState\_eq\_onGraph},
    \texttt{periodicChainHeisenbergGibbsState\_eq\_onGraph},
    \texttt{quantumIsingGibbsState\_eq\_isingGibbsStateOnGraph},
    \texttt{hubbardChainGibbsState\_eq\_onGraph}:
    chain Gibbs states are the graph form applied to the corresponding
    \texttt{pathGraph} / \texttt{cycleGraph}.
\end{itemize}

\subsection{Edge-sum decomposition helper}

The technical lemma
\texttt{LatticeSystem.Lattice.sum\_pathGraph\_adj} converts a
double sum over ordered adjacent pairs of \texttt{pathGraph (N+1)}
into a single sum over $\mathrm{Fin}~N$:

\begin{align*}
  \sum_{x,y \in \mathrm{Fin}(N+1)} [G.\mathrm{Adj}~x~y]\, f(x,y)
  = \sum_{i \in \mathrm{Fin}~N} \bigl(f(i.\mathrm{cs}, i.\mathrm{s})
    + f(i.\mathrm{s}, i.\mathrm{cs})\bigr).
\end{align*}

This is the missing piece used to prove the generic-$N$ Ising
chain bridge.

\subsection{Future directions}

\begin{itemize}
  \item Jordan--Wigner abstraction to any
    $[\mathrm{LinearOrder}~\Lambda]~[\mathrm{Fintype}~\Lambda]$
    (currently \texttt{Fin (N+1)}-specific).
  \item Infinite-volume many-body operators on graphs like
    $\mathbb{Z}$ or $\mathbb{Z}^d$ (already present at the graph
    level; requires new many-body operator infrastructure
    beyond the current finite-matrix representation).
  \item Bipartite-graph predicates for Lieb / MLM-style theorems.
\end{itemize}

%======================================================================
\section{How to read this codebase}
%======================================================================

\subsection{File dependency graph}

Read the files in this order. Each file imports only the files above it.

\begin{Verbatim}
LatticeSystem/
  Basic.lean              -- placeholder for shared infrastructure
  Quantum/
    Pauli.lean            -- single-site Pauli matrices, basic relations
    SpinHalf.lean         -- spin-1/2 operators S^(alpha) = sigma^(alpha)/2
                             and algebra (Tasaki section 2.1)
    SpinHalfBasis.lean    -- basis states |up>, |down>, raising/lowering S^+/-
                             and their actions (Tasaki section 2.1)
    SpinHalfDecomp.lean   -- Pauli basis of M_2(C) (Tasaki Problem 2.1.a, S=1/2)
    SpinHalfRotation.lean -- Closed-form S=1/2 spin rotation operators
                             (Tasaki section 2.1, eq 2.1.26)
    TotalSpin.lean        -- Total spin operator on an abstract lattice
                             (Tasaki section 2.2, eq 2.2.7)
    SpinOne.lean          -- S = 1 matrix representations (3x3)
                             (Tasaki section 2.1, eq 2.1.9)
    ManyBody.lean         -- multi-site operator space, onSite embedding,
                             commutativity at distinct sites
    IsingChain.lean       -- 1D open-chain quantum Ising Hamiltonian
    GibbsState.lean       -- Gibbs state rho = exp(-beta H)/Z,
                             trace normalization, thermal expectation
                             (Tasaki section 3.3)
    HeisenbergChain.lean  -- open and periodic Heisenberg chain Hamiltonians,
                             Hermiticity (Tasaki section 3.5, p. 89)
\end{Verbatim}

\subsection{Textbook correspondence}

Each module lists its primary textbook reference in its header and in
the per-module sections below. The references assembled for the project
overall are listed in Section~\ref{sec:refs}.

%======================================================================
\section{Single-site Pauli operators}
\label{sec:pauli}
%======================================================================

\emph{Source file:} \texttt{LatticeSystem/Quantum/Pauli.lean}.

\subsection{Definitions}

\begin{definition}[Pauli matrices]
The Pauli matrices are defined as elements of
$\operatorname{Matrix}(\mathrm{Fin}\,2, \mathrm{Fin}\,2, \mathbb{C})$ by
\[
  \sigma^x = \begin{pmatrix} 0 & 1 \\ 1 & 0 \end{pmatrix},\quad
  \sigma^y = \begin{pmatrix} 0 & -i \\ i & 0 \end{pmatrix},\quad
  \sigma^z = \begin{pmatrix} 1 & 0 \\ 0 & -1 \end{pmatrix}.
\]
Lean names: \texttt{pauliX}, \texttt{pauliY}, \texttt{pauliZ}.
\end{definition}

\subsection{Formalized facts}

\begin{lemma}[\texttt{pauliX\_isHermitian}, \texttt{pauliY\_isHermitian},
\texttt{pauliZ\_isHermitian}]
Each Pauli matrix is Hermitian: $(\sigma^\alpha)^\dagger = \sigma^\alpha$
for $\alpha \in \{x, y, z\}$.
\end{lemma}

\begin{lemma}[\texttt{pauliX\_mul\_self}, \texttt{pauliY\_mul\_self},
\texttt{pauliZ\_mul\_self}]
Each Pauli matrix squares to the identity: $(\sigma^\alpha)^2 = I$.
\end{lemma}

\begin{lemma}[\texttt{pauliX\_mul\_pauliY}, \texttt{pauliY\_mul\_pauliZ},
\texttt{pauliZ\_mul\_pauliX}]
The cyclic Pauli products hold:
\[
  \sigma^x \sigma^y = i \sigma^z,\quad
  \sigma^y \sigma^z = i \sigma^x,\quad
  \sigma^z \sigma^x = i \sigma^y.
\]
\end{lemma}

\begin{lemma}[\texttt{pauliX\_anticomm\_pauliY},
\texttt{pauliY\_anticomm\_pauliZ}, \texttt{pauliZ\_anticomm\_pauliX}]
Distinct Pauli matrices anticommute:
$\sigma^\alpha \sigma^\beta + \sigma^\beta \sigma^\alpha = 0$
for $\alpha \neq \beta$.
\end{lemma}

\paragraph{Proof technique.}
All four families are proved by \texttt{ext} on the matrix entries
followed by \texttt{fin\_cases} on each row/column index and
\texttt{simp} with the definitions.

\paragraph{References.}
Tasaki~\cite{Tasaki:QMB}, §2.1, eq.~(2.1.8), p.~15, defines the Pauli
matrices and remarks on $(\sigma^\alpha)^2 = I$ and the anticommutation
relations for spin-$\tfrac{1}{2}$ just before the definition (Tasaki
uses $\hat S^{(\alpha)}$ with
$\hat\sigma^{(\alpha)} = 2\hat S^{(\alpha)}$).
Independently, Nielsen--Chuang~\cite{NielsenChuang} gives the same
matrices in §2.1.3, Figure~2.2, pp.~65--66, and the four families of
relations we formalized are stated as exercises:
Ex.~2.19 (``Pauli matrices: Hermitian and unitary'') on p.~70 for
Hermiticity; Ex.~2.41 (``Anti-commutation relations'') on p.~78 for
$(\sigma^\alpha)^2 = I$ and $\{\sigma^\alpha, \sigma^\beta\} = 0$ for
$\alpha\neq\beta$; Ex.~2.40 (``Commutation relations'') on p.~77 for the
commutator $[\sigma^\alpha, \sigma^\beta] = 2i\varepsilon_{\alpha\beta\gamma}
\sigma^\gamma$ (combining this with the anticommutator gives the cyclic
products $\sigma^x\sigma^y = i\sigma^z$, etc.).

%======================================================================
\section{Basis states and raising/lowering operators (Tasaki §2.1)}
\label{sec:spinhalfbasis}
%======================================================================

\emph{Source file:} \texttt{LatticeSystem/Quantum/SpinHalfBasis.lean}.

\subsection{Definitions}

\begin{definition}[Tasaki eq.~(2.1.6), p.~14]
Identify the $S = 1/2$ computational basis states with the column vectors
\[
  |\psi^\uparrow\rangle
    = \begin{pmatrix} 1 \\ 0 \end{pmatrix},\quad
  |\psi^\downarrow\rangle
    = \begin{pmatrix} 0 \\ 1 \end{pmatrix}.
\]
Lean names: \texttt{spinHalfUp}, \texttt{spinHalfDown}, both of type
$\mathrm{Fin}\,2 \to \mathbb{C}$.
\end{definition}

\begin{definition}
The raising and lowering operators are
\[
  \hat S^+ = \begin{pmatrix} 0 & 1 \\ 0 & 0 \end{pmatrix},\quad
  \hat S^- = \begin{pmatrix} 0 & 0 \\ 1 & 0 \end{pmatrix}.
\]
Lean names: \texttt{spinHalfOpPlus}, \texttt{spinHalfOpMinus}.
\end{definition}

\subsection{Formalized facts}

\begin{lemma}[Tasaki eq.~(2.1.4), p.~14]
$\hat S^{(3)}|\psi^\uparrow\rangle = \tfrac{1}{2}|\psi^\uparrow\rangle$
and
$\hat S^{(3)}|\psi^\downarrow\rangle = -\tfrac{1}{2}|\psi^\downarrow\rangle$.
Lean names: \texttt{spinHalfOp3\_mulVec\_spinHalfUp/Down}.
\end{lemma}

\begin{lemma}[Consistency with Tasaki's standard definition]
$\hat S^+ = \hat S^{(1)} + i\,\hat S^{(2)}$ and
$\hat S^- = \hat S^{(1)} - i\,\hat S^{(2)}$.
Lean names: \texttt{spinHalfOpPlus\_eq\_add}, \texttt{spinHalfOpMinus\_eq\_sub}.
\end{lemma}

\begin{lemma}[Tasaki eq.~(2.1.5), p.~14]
\[
  \hat S^+|\psi^\uparrow\rangle = 0,\quad
  \hat S^-|\psi^\uparrow\rangle = |\psi^\downarrow\rangle,\quad
  \hat S^+|\psi^\downarrow\rangle = |\psi^\uparrow\rangle,\quad
  \hat S^-|\psi^\downarrow\rangle = 0.
\]
Lean names: \texttt{spinHalfOpPlus/Minus\_mulVec\_spinHalfUp/Down}.
\end{lemma}

\paragraph{Proof technique.}
Matrix-vector multiplication is \texttt{Matrix.mulVec}, which unfolds to
a sum computable via \texttt{Fin.sum\_univ\_two}. Each equation reduces
entry-wise by \texttt{ext i; fin\_cases i; simp} with the relevant
definitions. The equivalence
\texttt{spinHalfOpPlus = spinHalfOp1 + I \textbullet\ spinHalfOp2} needs
the single nontrivial complex identity $i^2 = -1$ (\texttt{Complex.I\_sq});
it is dispatched by \texttt{ring\_nf; rw [Complex.I\_sq]; ring}.

\paragraph{References.}
Tasaki~\cite{Tasaki:QMB}:
\begin{itemize}
  \item eq.~(2.1.4), p.~14: $\hat S^{(3)}$ eigenvalue equation
    for $S = 1/2$.
  \item eq.~(2.1.5), p.~14: action of $\hat S^\pm$.
  \item eq.~(2.1.6), p.~14: identification of basis states with column
    vectors.
  \item The raising/lowering definition $\hat S^\pm := \hat S^{(1)}
    \pm i \hat S^{(2)}$ appears in the paragraph above eq.~(2.1.3) on
    p.~14.
\end{itemize}

%======================================================================
\section{Spin-1/2 rotation operators (Tasaki §2.1 eq 2.1.26)}
\label{sec:spinhalfrot}
%======================================================================

\emph{Source file:} \texttt{LatticeSystem/Quantum/SpinHalfRotation.lean}.

\subsection{Definitions}

\begin{definition}[Tasaki eq.~(2.1.26), p.~17]
For $\alpha \in \{1, 2, 3\}$ and $\theta \in \mathbb{R}$, the
$S = 1/2$ spin rotation operator about the $\alpha$-axis by angle
$\theta$ is
\[
  \hat U^{(\alpha)}_\theta
    := \cos(\theta/2)\,\hat 1 - 2i\,\sin(\theta/2)\,\hat S^{(\alpha)}
    \in M_2(\mathbb{C}).
\]
Lean names: \texttt{spinHalfRot1}, \texttt{spinHalfRot2}, \texttt{spinHalfRot3}.
\end{definition}

For $S = 1/2$ this closed form coincides with the matrix exponential
$\exp(-i\theta\hat S^{(\alpha)})$ because $(\hat S^{(\alpha)})^2 = I/4$;
establishing that equivalence with Mathlib's \texttt{Matrix.exp} is
future work.

\subsection{Formalized facts}

\begin{lemma}[identity at $\theta = 0$]
$\hat U^{(\alpha)}_0 = \hat 1$. Lean name:
\texttt{spinHalfRot1/2/3\_zero}.
\end{lemma}

\begin{lemma}[adjoint is the reverse rotation]
$(\hat U^{(\alpha)}_\theta)^\dagger = \hat U^{(\alpha)}_{-\theta}$.
Lean name: \texttt{spinHalfRot1/2/3\_adjoint}.
\end{lemma}

\begin{theorem}[$2\pi$ rotation, Tasaki eq.~(2.1.23), p.~16, $S = 1/2$]
$\hat U^{(\alpha)}_{2\pi} = -\hat 1$. Lean name:
\texttt{spinHalfRot1/2/3\_two\_pi}.
\end{theorem}

\begin{theorem}[group law (Tasaki p.~15)]
$\hat U^{(\alpha)}_\theta \hat U^{(\alpha)}_\varphi
  = \hat U^{(\alpha)}_{\theta + \varphi}$.
Lean name: \texttt{spinHalfRot1/2/3\_mul}.
\end{theorem}

\begin{theorem}[unitarity]
$\hat U^{(\alpha)}_\theta (\hat U^{(\alpha)}_\theta)^\dagger = \hat 1$.
Lean name: \texttt{spinHalfRot1/2/3\_unitary}.
\end{theorem}

\subsection{The $\pi$-rotations}

\begin{lemma}[$\hat U^{(\alpha)}_\pi$ in closed form]
$\hat U^{(\alpha)}_\pi = -2i\,\hat S^{(\alpha)}$. Proof: instantiate the
closed form at $\theta = \pi$ using $\cos(\pi/2) = 0$ and
$\sin(\pi/2) = 1$. Lean name: \texttt{spinHalfRot1/2/3\_pi}.
\end{lemma}

\begin{lemma}
$(\hat U^{(\alpha)}_\pi)^2 = -\hat 1$. Proof: group law
$\hat U^{(\alpha)}_\pi \hat U^{(\alpha)}_\pi = \hat U^{(\alpha)}_{2\pi}
= -\hat 1$. Lean name: \texttt{spinHalfRot1/2/3\_pi\_sq}.
\end{lemma}

\begin{theorem}[Tasaki eq.~(2.1.25), $S = 1/2$ case]
For $\alpha \neq \beta$,
$\hat U^{(\alpha)}_\pi \hat U^{(\beta)}_\pi
  + \hat U^{(\beta)}_\pi \hat U^{(\alpha)}_\pi = 0$.
Lean names: \texttt{spinHalfRot1\_pi\_anticomm\_spinHalfRot2\_pi}
and the two cyclic variants.
\end{theorem}

\begin{proof}
Use $\hat U^{(\alpha)}_\pi = -2i\,\hat S^{(\alpha)}$ to rewrite
\[
  \hat U^{(\alpha)}_\pi \hat U^{(\beta)}_\pi
    + \hat U^{(\beta)}_\pi \hat U^{(\alpha)}_\pi
    = (-2i)^2 (\hat S^{(\alpha)}\hat S^{(\beta)}
      + \hat S^{(\beta)}\hat S^{(\alpha)})
    = (-2i)^2 \cdot 0 = 0
\]
by the $\hat S^{(\alpha)}$ anticommutation at distinct axes.
\end{proof}

\subsection{$\pi$-rotation products and coordinate transformations}

\begin{theorem}[Tasaki eq.~(2.1.29), $S = 1/2$ case]
$\hat U^{(1)}_\pi \hat U^{(2)}_\pi = \hat U^{(3)}_\pi$ and the two
cyclic permutations. Lean names:
\texttt{spinHalfRot1\_pi\_mul\_spinHalfRot2\_pi} etc.
\end{theorem}

\begin{proof}
$\hat U^{(1)}_\pi \hat U^{(2)}_\pi = (-2i)^2 \hat S^{(1)}\hat S^{(2)}
  = -4 \cdot \tfrac{i}{2}\hat S^{(3)} = -2i\,\hat S^{(3)} = \hat U^{(3)}_\pi$,
using the cyclic product $\hat S^{(1)}\hat S^{(2)} = \tfrac{i}{2}\hat S^{(3)}$
formalized as \texttt{spinHalfOp1\_mul\_spinHalfOp2} in
\texttt{Quantum/SpinHalf.lean}.
\end{proof}

\begin{theorem}[conjugation by $\hat U^{(\alpha)}_\pi$]
For any $\alpha, \beta \in \{1, 2, 3\}$,
\[
  (\hat U^{(\alpha)}_\pi)^\dagger \hat S^{(\beta)} \hat U^{(\alpha)}_\pi
    = \begin{cases}
        +\hat S^{(\alpha)} & \text{if } \alpha = \beta, \\
        -\hat S^{(\beta)} & \text{if } \alpha \neq \beta.
      \end{cases}
\]
Lean names: \texttt{spinHalfRot1/2/3\_pi\_conj\_spinHalfOp1/2/3} for
the nine distinct pairings. This is the Tasaki eq.~(2.1.15)
($\alpha = \beta$) and eq.~(2.1.21) ($\alpha \neq \beta$) both
specialized at $\theta = \pi$.
\end{theorem}

\begin{proof}
Using $(\hat U^{(\alpha)}_\pi)^\dagger = 2i\,\hat S^{(\alpha)}$ and
$\hat U^{(\alpha)}_\pi = -2i\,\hat S^{(\alpha)}$, the conjugation
reduces to
\[
  (\hat U^{(\alpha)}_\pi)^\dagger \hat S^{(\beta)} \hat U^{(\alpha)}_\pi
    = (2i)(-2i)\,\hat S^{(\alpha)} \hat S^{(\beta)} \hat S^{(\alpha)}
    = -4 i^2\,\hat S^{(\alpha)} \hat S^{(\beta)} \hat S^{(\alpha)}
    = 4\,\hat S^{(\alpha)} \hat S^{(\beta)} \hat S^{(\alpha)}.
\]
For $\alpha = \beta$: $\hat S^{(\alpha)} \hat S^{(\alpha)} \hat S^{(\alpha)}
  = (\hat 1/4) \hat S^{(\alpha)} = \hat S^{(\alpha)}/4$, giving
$\hat S^{(\alpha)}$.
For $\alpha \neq \beta$: anticommutation
$\hat S^{(\alpha)} \hat S^{(\beta)} = -\hat S^{(\beta)} \hat S^{(\alpha)}$
plus $\hat S^{(\alpha)} \hat S^{(\alpha)} = \hat 1/4$ give
$\hat S^{(\alpha)} \hat S^{(\beta)} \hat S^{(\alpha)}
  = -\hat S^{(\beta)}/4$, giving $-\hat S^{(\beta)}$.
Both cases require $i^2 = -1$ (\texttt{Complex.I\_sq}).
\end{proof}

\paragraph{Proof technique.}
All statements are parametrised through a private helper
\texttt{rotOf S theta := cos(theta/2) * 1 - (2*I*sin(theta/2)) * S} and then
instantiated at \texttt{S = spinHalfOp\{1,2,3\}}. The
\texttt{adjoint} lemma needs \texttt{Complex.conj\_I} (for
$\overline{i} = -i$), \texttt{Complex.conj\_ofReal} (for real sin),
and \texttt{Real.cos\_neg} / \texttt{Real.sin\_neg}.
The $2\pi$ lemma follows from \texttt{Real.cos\_pi} and
\texttt{Real.sin\_pi}.

\paragraph{Proof technique for the group law and unitarity.}
The key step is an algebraic expansion lemma
\texttt{rot\_mul\_helper}: for any $S$ and $k$ with $S^2 = k I$, and
any $a, b, c, d \in \mathbb{C}$,
\[
  (aI - bS)(cI - dS) = (ac + bdk)\,I - (ad + bc)\,S.
\]
The proof expands the product via \texttt{sub\_mul}, \texttt{mul\_sub},
\texttt{Matrix.smul\_mul}, \texttt{Matrix.mul\_smul},
\texttt{Matrix.one\_mul}, \texttt{Matrix.mul\_one} and
\texttt{smul\_smul}, substitutes $S^2 = kI$, and closes with the
Mathlib \texttt{module} tactic that collects the $I$ and $S$ components.
For the group law, the remaining ring identity in $\cos$ and $\sin$
uses the addition formulas \texttt{Real.cos\_add} / \texttt{Real.sin\_add}
and a \texttt{linear\_combination} with \texttt{Complex.I\_sq} to
eliminate $i^2 = -1$. Unitarity follows by combining the adjoint
identity with the group law:
$\hat U^{(\alpha)}_\theta (\hat U^{(\alpha)}_\theta)^\dagger
  = \hat U^{(\alpha)}_\theta \hat U^{(\alpha)}_{-\theta}
  = \hat U^{(\alpha)}_{\theta + (-\theta)}
  = \hat U^{(\alpha)}_0 = \hat 1$.

\paragraph{Not yet formalized.}
\begin{itemize}
  \item The exponential equivalence
    $\hat U^{(\alpha)}_\theta = \exp(-i\theta\,\hat S^{(\alpha)})$
    via Mathlib's \texttt{Matrix.exp}.
\end{itemize}

\paragraph{References.}
Tasaki~\cite{Tasaki:QMB}, §2.1, eq.~(2.1.23) (p.~16) and eq.~(2.1.26)
(p.~17).

%======================================================================
\section{Pauli-basis decomposition for $M_2(\mathbb{C})$ (Tasaki §2.1 Problem 2.1.a, S = 1/2)}
\label{sec:spinhalfdecomp}
%======================================================================

\emph{Source file:} \texttt{LatticeSystem/Quantum/SpinHalfDecomp.lean}.

\subsection{Definitions and main theorem}

Tasaki's Problem~2.1.a (p.~15) asks to show that every operator on
$\mathcal{H}_0 \simeq \mathbb{C}^{2S+1}$ is a polynomial of $\hat 1,
\hat S^{(1)}, \hat S^{(2)}, \hat S^{(3)}$. For $S = 1/2$ the operator
space has dimension $4 = (2S+1)^2$, so a \emph{linear} combination
suffices.

\begin{definition}
Given $A \in M_2(\mathbb{C})$, define the four scalar coefficients
\begin{align*}
  c_0(A) &= \tfrac{A_{00} + A_{11}}{2}, &
  c_1(A) &= \tfrac{A_{01} + A_{10}}{2}, \\
  c_2(A) &= \tfrac{i (A_{01} - A_{10})}{2}, &
  c_3(A) &= \tfrac{A_{00} - A_{11}}{2}.
\end{align*}
Lean names: \texttt{pauliCoeff0}, \texttt{pauliCoeff1},
\texttt{pauliCoeff2}, \texttt{pauliCoeff3}.
\end{definition}

\begin{theorem}[\texttt{pauli\_decomposition}]
For every $A \in M_2(\mathbb{C})$,
\[
  A = c_0(A)\,I + c_1(A)\,\sigma^x + c_2(A)\,\sigma^y
      + c_3(A)\,\sigma^z.
\]
\end{theorem}

\begin{theorem}[\texttt{spinHalf\_decomposition}]
The same decomposition expressed via $\hat S^{(\alpha)} =
\sigma^{(\alpha)}/2$:
\[
  A = c_0(A)\,I + 2 c_1(A)\,\hat S^{(1)}
      + 2 c_2(A)\,\hat S^{(2)} + 2 c_3(A)\,\hat S^{(3)}.
\]
\end{theorem}

\begin{theorem}[\texttt{pauli\_linearIndep}]
If $a_0 I + a_1 \sigma^x + a_2 \sigma^y + a_3 \sigma^z = 0$ for some
$a_0, a_1, a_2, a_3 \in \mathbb{C}$, then all $a_i = 0$. Combined with
\texttt{pauli\_decomposition}, this says $\{I, \sigma^x, \sigma^y,
\sigma^z\}$ is a $\mathbb{C}$-basis of $M_2(\mathbb{C})$.
\end{theorem}

\paragraph{Proof technique.}
The decomposition reduces to an entry-wise calculation by
\texttt{ext; fin\_cases; simp; ring\_nf; simp only [Complex.I\_sq]; ring}.
Linear independence extracts four scalar equations by evaluating both
sides at the four matrix entries $(0,0), (0,1), (1,0), (1,1)$; the
resulting $4 \times 4$ linear system is solved by
\texttt{linear\_combination} (the off-diagonal case requires
\texttt{Complex.I\_ne\_zero} to conclude $a_2 = 0$ from
$a_2 \cdot i = 0$).

\paragraph{References.}
Tasaki~\cite{Tasaki:QMB}, §2.1, Problem~2.1.a, p.~15.
The author's solution on p.~493 proves the general-$S$ case; we cover
only $S = 1/2$ here.

%======================================================================
\section{S = 1 matrix representations (Tasaki §2.1 eq 2.1.9)}
\label{sec:spinone}
%======================================================================

\emph{Source file:} \texttt{LatticeSystem/Quantum/SpinOne.lean}.

\subsection{Definitions}

\begin{definition}[Tasaki eq.~(2.1.9), p.~15]
The $S = 1$ spin operators are the $3 \times 3$ matrices
\[
  \hat S^{(1)} = \tfrac{1}{\sqrt{2}}
    \begin{pmatrix} 0 & 1 & 0 \\ 1 & 0 & 1 \\ 0 & 1 & 0 \end{pmatrix},\quad
  \hat S^{(2)} = \tfrac{1}{\sqrt{2}}
    \begin{pmatrix} 0 & -i & 0 \\ i & 0 & -i \\ 0 & i & 0 \end{pmatrix},\quad
  \hat S^{(3)} =
    \begin{pmatrix} 1 & 0 & 0 \\ 0 & 0 & 0 \\ 0 & 0 & -1 \end{pmatrix}.
\]
Lean names: \texttt{spinOneOp1}, \texttt{spinOneOp2}, \texttt{spinOneOp3}.
\end{definition}

\subsection{Formalized facts}

\begin{lemma}[Hermiticity]
Each $\hat S^{(\alpha)}$ (for $\alpha \in \{1, 2, 3\}$) is Hermitian
(\texttt{spinOneOp1/2/3\_isHermitian}).
\end{lemma}

\begin{theorem}[Commutation relations, Tasaki eq.~(2.1.1), $S = 1$ case]
\[
  [\hat S^{(1)}, \hat S^{(2)}] = i\,\hat S^{(3)},\quad
  [\hat S^{(2)}, \hat S^{(3)}] = i\,\hat S^{(1)},\quad
  [\hat S^{(3)}, \hat S^{(1)}] = i\,\hat S^{(2)}.
\]
Lean names: \texttt{spinOneOp1\_commutator\_spinOneOp2} and cyclic.
\end{theorem}

\begin{theorem}[Total spin, $S = 1$ case of $\hat S^2 = S(S+1) I$]
\[
  (\hat S^{(1)})^2 + (\hat S^{(2)})^2 + (\hat S^{(3)})^2 = 2\,I.
\]
Lean name: \texttt{spinOne\_total\_spin\_squared}.
\end{theorem}

\paragraph{Proof technique.}
Every statement is verified by entry-wise computation:
\texttt{ext; fin\_cases i; fin\_cases j; simp} expands the $3 \times 3$
matrix products via the \texttt{!!} notation's built-in simp rules.
The auxiliary \texttt{invSqrt2\_sq : invSqrt2\textasciicircum 2 = 1/2}
(proved from \texttt{Real.mul\_self\_sqrt}) handles the squared
$\tfrac{1}{\sqrt{2}}$ coefficient, and \texttt{Complex.I\_sq : I\textasciicircum 2 = -1}
handles the squared imaginary unit after \texttt{ring\_nf} normalization.
The \texttt{!!} notation's simp set already rewrites bare $I \cdot I$ to
$-1$, so \texttt{Complex.I\_sq} is only required where
\texttt{ring\_nf} has introduced \texttt{I\textasciicircum 2}.

\paragraph{References.}
Tasaki~\cite{Tasaki:QMB}, §2.1, eq.~(2.1.9), p.~15.
The commutation and total-spin identities are the $S = 1$ instances of
the abstract relations eq.~(2.1.1) on p.~13 and the convention
$\hat S^2 = S(S+1) I$ stated just below it (p.~14).

%======================================================================
\section{Abstract lattice, site operators, and total spin (Tasaki §2.2)}
\label{sec:sec22}
%======================================================================

\emph{Source files:} \texttt{LatticeSystem/Quantum/ManyBody.lean},
\texttt{Quantum/TotalSpin.lean}.

\subsection{Abstract lattice}

Tasaki~§2.2, p.~21 formulates quantum spin systems on any finite set
$\Lambda$. Our \texttt{ManyBodyOp} now takes $\Lambda : \text{Type}^*$
equipped with \texttt{Fintype} and \texttt{DecidableEq} instances, with
the Hilbert space model
\[
  \texttt{ManyBodyOp}\,\Lambda
    := \operatorname{Matrix}\bigl((\Lambda \to \mathrm{Fin}\,2),
        (\Lambda \to \mathrm{Fin}\,2), \mathbb{C}\bigr).
\]
All earlier theorems about \texttt{onSite}, \texttt{onSite\_isHermitian},
and \texttt{onSite\_mul\_onSite\_of\_ne} carry over verbatim with the
abstract $\Lambda$.

\subsection{Distinct-site commutation (Tasaki eq.~(2.2.6), S = 1/2)}

\begin{theorem}
For $x \neq y$ in $\Lambda$ and any single-site operators $S_\alpha$,
$S_\beta$,
$\hat S_x^{(\alpha)} \hat S_y^{(\beta)} = \hat S_y^{(\beta)} \hat S_x^{(\alpha)}$.
\end{theorem}

This is the $S = 1/2$ case of Tasaki eq.~(2.2.6) on p.~22. Lean name:
\texttt{onSite\_mul\_onSite\_of\_ne} (or its S = 1/2 named wrapper
\texttt{spinHalfOp\_onSite\_comm\_of\_ne}).

\subsection{Total spin (Tasaki eq.~(2.2.7))}

\begin{definition}
$\hat S_{\mathrm{tot}}^{(\alpha)}
  := \sum_{x \in \Lambda} \hat S_x^{(\alpha)}
  = \sum_{x \in \Lambda} \texttt{onSite}\,x\,\hat S^{(\alpha)}$
for $\alpha \in \{1, 2, 3\}$. Lean names:
\texttt{totalSpinHalfOp1}, \texttt{totalSpinHalfOp2},
\texttt{totalSpinHalfOp3}.
\end{definition}

\begin{lemma}
$\hat S_{\mathrm{tot}}^{(\alpha)}$ is Hermitian, for each $\alpha$.
Lean names: \texttt{totalSpinHalfOp1/2/3\_isHermitian}.
\end{lemma}

\begin{proof}
A sum of Hermitian matrices is Hermitian, and each
$\texttt{onSite}\,x\,\hat S^{(\alpha)}$ is Hermitian by the lemma
\texttt{onSite\_isHermitian}.
\end{proof}

\subsection{Same-site commutation (Tasaki eq.~(2.2.6), $x = y$ case)}

\begin{lemma}
For any $i \in \Lambda$ and any $A, B \in M_2(\mathbb C)$,
\[
  \texttt{onSite}\,i\,A \cdot \texttt{onSite}\,i\,B
    = \texttt{onSite}\,i\,(A \cdot B).
\]
Lean name: \texttt{onSite\_mul\_onSite\_same}. As a direct consequence,
\[
  [\texttt{onSite}\,i\,A,\,\texttt{onSite}\,i\,B]
    = \texttt{onSite}\,i\,[A, B]
\]
(\texttt{onSite\_commutator\_same}). Specializing $A, B$ to
$\hat S^{(\alpha)}, \hat S^{(\beta)}$ recovers the $x = y$ case of
Tasaki eq.~(2.2.6), namely
$[\hat S_x^{(\alpha)}, \hat S_x^{(\beta)}]
  = i \sum_\gamma \varepsilon^{\alpha\beta\gamma} \hat S_x^{(\gamma)}$
for $\alpha, \beta \in \{1, 2, 3\}$. Lean names:
\texttt{spinHalfOp1\_onSite\_commutator\_spinHalfOp2\_onSite},
\texttt{spinHalfOp2\_onSite\_commutator\_spinHalfOp3\_onSite},
\texttt{spinHalfOp3\_onSite\_commutator\_spinHalfOp1\_onSite}.
\end{lemma}

\begin{proof}
Expand the matrix product via \texttt{Matrix.mul\_apply} and observe
that $\texttt{onSite}\,i\,A$ is block-diagonal along the fibers
$\{\tau : \tau \upharpoonright_{\Lambda\setminus\{i\}} = \sigma
 \upharpoonright_{\Lambda\setminus\{i\}}\}$. Within each fiber, the
product is the $2\times 2$ matrix multiplication $A \cdot B$.
Off-fiber contributions vanish because the constraint
``$\tau$ agrees with $\sigma'$ off $i$'' and
``$\tau$ agrees with $\sigma$ off $i$'' together force $\sigma'$ to
agree with $\sigma$ off $i$. See
\texttt{LatticeSystem/Quantum/ManyBody.lean}.
\end{proof}

\subsection{Total raising/lowering operators (Tasaki eq.~(2.2.8))}

\begin{definition}
The total raising and lowering operators are
\[
  \hat S^{\pm}_{\mathrm{tot}}
    := \sum_{x \in \Lambda} \hat S_x^{\pm}
    = \sum_{x \in \Lambda} \texttt{onSite}\,x\,\hat S^{\pm}.
\]
Lean names: \texttt{totalSpinHalfOpPlus}, \texttt{totalSpinHalfOpMinus}.
\end{definition}

\begin{lemma}[Tasaki eq.~(2.2.8)]
$\hat S^{\pm}_{\mathrm{tot}}
  = \hat S^{(1)}_{\mathrm{tot}} \pm i \hat S^{(2)}_{\mathrm{tot}}$.
Lean names: \texttt{totalSpinHalfOpPlus\_eq\_add},
\texttt{totalSpinHalfOpMinus\_eq\_sub}.
\end{lemma}

\begin{proof}
At each site $x$ this is
$\hat S_x^{+} = \hat S_x^{(1)} + i\hat S_x^{(2)}$
(\texttt{spinHalfOpPlus\_eq\_add}, applied through
\texttt{onSite\_add} and \texttt{onSite\_smul}), which sums over
$x \in \Lambda$ by linearity.
\end{proof}

\subsection{Total magnetization (Tasaki eq.~(2.2.2))}

\begin{definition}
For $\sigma : \Lambda \to \mathrm{Fin}\,2$, let
\[
  \operatorname{spinSign}(s) := \begin{cases} +1 & s = 0, \\ -1 & s = 1, \end{cases}
  \qquad
  |\sigma| := \sum_{x \in \Lambda} \operatorname{spinSign}(\sigma_x).
\]
Lean names: \texttt{spinSign}, \texttt{magnetization}.
\end{definition}

\subsection{Two-site spin inner product (Tasaki eq.~(2.2.16))}

\begin{definition}
\[
  \hat S_x \cdot \hat S_y
    := \sum_{\alpha = 1}^{3}
       \texttt{onSite}\,x\,\hat S^{(\alpha)}
       \cdot
       \texttt{onSite}\,y\,\hat S^{(\alpha)}.
\]
Lean name: \texttt{spinHalfDot}.
\end{definition}

\begin{lemma}[Tasaki eq.~(2.2.16)]
For $S = 1/2$,
\[
  \hat S_x \cdot \hat S_y
    = \tfrac12 \bigl( \hat S_x^{+}\hat S_y^{-}
                    + \hat S_x^{-}\hat S_y^{+}\bigr)
      + \hat S_x^{(3)} \hat S_y^{(3)}.
\]
Lean name: \texttt{spinHalfDot\_eq\_plus\_minus}.
\end{lemma}

\begin{proof}
Expand the right-hand side via
$\hat S_x^{\pm} = \hat S_x^{(1)} \pm i\hat S_x^{(2)}$ and the companion
identity at site $y$, then use $i^2 = -1$. Writing $a_k, b_k$ for
$\texttt{onSite}\,x\,\hat S^{(k)}$, $\texttt{onSite}\,y\,\hat S^{(k)}$
respectively, the expansion gives
\[
  (a_1 + i a_2)(b_1 - i b_2) + (a_1 - i a_2)(b_1 + i b_2)
    = 2(a_1 b_1 + a_2 b_2),
\]
so the right-hand side of the identity collapses to
$a_1 b_1 + a_2 b_2 + a_3 b_3$, which is the definition of
$\hat S_x \cdot \hat S_y$.
\end{proof}

\paragraph{Not yet formalized.}
\begin{itemize}
  \item Global rotation operator
    $\hat U^{(\alpha)}_\theta = \exp(-i\theta \hat S_{\mathrm{tot}}^{(\alpha)})$
    (eq.~(2.2.11)).
  \item SU(2) / U(1) invariance of operators commuting with
    $\hat S_{\mathrm{tot}}^{(\alpha)}$ (eqs.~(2.2.12)-(2.2.13)).
  \item SU(2) invariance of $\hat S_x \cdot \hat S_y$ and its
    eigenvalues (eqs.~(2.2.17)-(2.2.19)).
  \item Magnetization eigenspace decomposition
    $\mathcal H = \bigoplus_M \mathcal H_M$ (eqs.~(2.2.9)-(2.2.10)).
  \item Problem 2.2.a-c.
\end{itemize}

%======================================================================
\section{S = 1 basis states and raising/lowering
         (Tasaki \S2.1 eqs.~(2.1.2)-(2.1.6), S = 1)}
\label{sec:spinonebasis}
%======================================================================

\emph{Source file:} \texttt{LatticeSystem/Quantum/SpinOneBasis.lean}.

\begin{definition}
Column vectors in $\mathrm{Fin}\,3 \to \mathbb C$:
\[
  |\psi^{+1}\rangle := (1, 0, 0)^T, \quad
  |\psi^{0}\rangle  := (0, 1, 0)^T, \quad
  |\psi^{-1}\rangle := (0, 0, 1)^T.
\]
Lean names: \texttt{spinOnePlus}, \texttt{spinOneZero},
\texttt{spinOneMinus}.
\end{definition}

\begin{lemma}[Tasaki eq.~(2.1.2), S = 1]
\[
  \hat S^{(3)} |\psi^{+1}\rangle = |\psi^{+1}\rangle, \quad
  \hat S^{(3)} |\psi^{0}\rangle = 0, \quad
  \hat S^{(3)} |\psi^{-1}\rangle = -|\psi^{-1}\rangle.
\]
Lean names:
\texttt{spinOneOp3\_mulVec\_spinOnePlus/Zero/Minus}.
\end{lemma}

\begin{definition}
The raising/lowering matrices (cf.~(2.1.3), $S = 1$):
\[
  \hat S^{+} := \begin{pmatrix} 0 & \sqrt 2 & 0 \\ 0 & 0 & \sqrt 2 \\ 0 & 0 & 0 \end{pmatrix}, \qquad
  \hat S^{-} := \begin{pmatrix} 0 & 0 & 0 \\ \sqrt 2 & 0 & 0 \\ 0 & \sqrt 2 & 0 \end{pmatrix}.
\]
Lean names: \texttt{spinOneOpPlus}, \texttt{spinOneOpMinus}.
\end{definition}

\begin{lemma}[Tasaki eq.~(2.1.3), S = 1]
\[
  \hat S^{+} |\psi^{+1}\rangle = 0, \quad
  \hat S^{+} |\psi^{0}\rangle = \sqrt 2 \,|\psi^{+1}\rangle, \quad
  \hat S^{+} |\psi^{-1}\rangle = \sqrt 2 \,|\psi^{0}\rangle,
\]
and the corresponding actions of $\hat S^{-}$. Lean names:
\texttt{spinOneOpPlus/Minus\_mulVec\_spinOnePlus/Zero/Minus}.
\end{lemma}

\paragraph{Not yet formalized.}
\begin{itemize}
  \item Identification
    $\hat S^{\pm} = \hat S^{(1)} \pm i\hat S^{(2)}$ for $S = 1$
    (follows from the definitions~(2.1.9)).
  \item Problem 2.1.c: $S = 1$ analogue of the closed form~(2.1.26).
  \item Problem 2.1.g: derivation of the explicit $\pi$-rotation
    matrices~(2.1.33)/(2.1.34) for $S = 1$.
\end{itemize}

%======================================================================
\section{Multi-body operator space and site embedding}
\label{sec:manybody}
%======================================================================

\emph{Source file:} \texttt{LatticeSystem/Quantum/ManyBody.lean}.

\subsection{Definitions}

\begin{definition}[Many-body operator space]
For an $N$-site lattice, the many-body operator space is
\[
  \texttt{ManyBodyOp}\,N
    := \operatorname{Matrix}\bigl(
      (\mathrm{Fin}\,N \to \mathrm{Fin}\,2),\;
      (\mathrm{Fin}\,N \to \mathrm{Fin}\,2),\;
      \mathbb{C}\bigr).
\]
Equivalently, this is the space of operators on $\mathbb{C}^{2^N}$
indexed by the computational-basis configurations
$\sigma : \mathrm{Fin}\,N \to \mathrm{Fin}\,2$.
\end{definition}

\begin{definition}[Site-embedded operator]
For $i \in \mathrm{Fin}\,N$ and
$A \in \operatorname{Matrix}(\mathrm{Fin}\,2, \mathrm{Fin}\,2, \mathbb{C})$,
define $\texttt{onSite}\,i\,A \in \texttt{ManyBodyOp}\,N$ by
\[
  (\texttt{onSite}\,i\,A)_{\sigma',\sigma}
    = \begin{cases}
        A_{\sigma'(i),\sigma(i)} & \text{if } \sigma'(k) = \sigma(k)
          \text{ for every } k \neq i, \\
        0 & \text{otherwise.}
      \end{cases}
\]
\end{definition}

\subsection{Formalized facts}

\begin{lemma}[\texttt{onSite\_isHermitian}]
If $A$ is Hermitian, so is $\texttt{onSite}\,i\,A$.
\end{lemma}

\begin{lemma}[\texttt{onSite\_mul\_onSite\_of\_ne}]
For $i \neq j$ and any single-site $A$, $B$,
\[
  (\texttt{onSite}\,i\,A)(\texttt{onSite}\,j\,B)
    = (\texttt{onSite}\,j\,B)(\texttt{onSite}\,i\,A).
\]
\end{lemma}

\begin{lemma}[\texttt{Matrix.IsHermitian.mul\_of\_commute}]
If $A$, $B$ are Hermitian matrices with $AB = BA$, then $AB$ is
Hermitian.
\end{lemma}

\paragraph{Proof technique.}
The site-commutativity proof is the technical core. Expanding the
matrix product gives
\[
  \bigl((\texttt{onSite}\,i\,A)(\texttt{onSite}\,j\,B)\bigr)_{\sigma',\sigma}
    = \sum_\tau
      (\texttt{onSite}\,i\,A)_{\sigma',\tau}\,
      (\texttt{onSite}\,j\,B)_{\tau,\sigma}.
\]
Since $i \neq j$, at most one $\tau$ makes both factors nonzero:
$\tau = \mathrm{Function.update}\,\sigma\,j\,(\sigma'(j))$, i.e.,
$\tau(j) = \sigma'(j)$ and $\tau(k) = \sigma(k)$ otherwise. The proof
isolates this pivot by \texttt{Fintype.sum\_eq\_single} and shows all
other $\tau$ kill at least one factor. A matching computation for the
swapped ordering, combined with the commutativity of~$\mathbb{C}$,
yields the equality.

\paragraph{References.}
The matrix-level formulation we use (configuration-indexed matrices on
$\mathbb{C}^{\mathrm{Fin}\,N \to \mathrm{Fin}\,2}$) is a standard
computational-basis setup; see Tasaki~\cite{Tasaki:QMB}, §2.2,
pp.~21--26, where the tensor-product structure of the Hilbert space
and the general form of site-local operators are developed (our
\texttt{onSite}$\,i\,A$ is the matrix realization of the operator
$\bigotimes_{k\neq i} \hat I_k \otimes A_i$ in the computational basis).

%======================================================================
\section{Quantum Ising chain on an open boundary}
\label{sec:ising}
%======================================================================

\emph{Source file:} \texttt{LatticeSystem/Quantum/IsingChain.lean}.

\subsection{Definition}

\begin{definition}[Open-chain quantum Ising Hamiltonian]
Let $N \in \mathbb{N}$, $J \in \mathbb{R}$, and $h \in \mathbb{R}$.
On the $(N+1)$-site lattice, define
\[
  H_{J,h}^{N+1}
    = -J \sum_{i=0}^{N-1} \sigma^z_i \sigma^z_{i+1}
      - h \sum_{i=0}^{N} \sigma^x_i,
\]
where $\sigma^\alpha_i := \texttt{onSite}\,i\,\sigma^\alpha$. The
coupling sum is indexed by $i \in \mathrm{Fin}\,N$ with
$i \mapsto (\texttt{castSucc}\,i, \texttt{succ}\,i)$, so it is empty
when $N = 0$ (one-site system) and yields the usual $N$ coupling terms
on an $(N+1)$-site open chain otherwise. Lean name:
\texttt{quantumIsingHamiltonian}.
\end{definition}

\subsection{Formalized fact}

\begin{theorem}[\texttt{quantumIsingHamiltonian\_isHermitian}]
For any $N \in \mathbb{N}$ and real $J$, $h$, the Hamiltonian
$H_{J,h}^{N+1}$ is Hermitian.
\end{theorem}

\paragraph{Proof technique.}
The bond term $\sigma^z_i \sigma^z_{i+1}$ is a product of two Hermitian
operators at distinct sites (via
\texttt{castSucc\_ne\_succ}), which commute by
\texttt{onSite\_mul\_onSite\_of\_ne} and whose product is then
Hermitian by \texttt{Matrix.IsHermitian.mul\_of\_commute}.
The field term $\sigma^x_i$ is Hermitian by
\texttt{onSite\_isHermitian}. Hermiticity is closed under finite sums
(local helper \texttt{isHermitian\_sum}) and under real scalar
multiplication (local helper \texttt{isHermitian\_smul\_real} using
$\overline{c} = c$ for $c \in \mathbb{R}$).

\paragraph{References.}
Tasaki~\cite{Tasaki:QMB}, §3.3, eq.~(3.3.1) on p.~55 states the
one-dimensional transverse-field quantum Ising Hamiltonian on an open
chain:
\[
  \hat H
    = -\sum_{x=1}^{L-1}\hat S^{(3)}_x \hat S^{(3)}_{x+1}
      -\lambda \sum_{x=1}^{L} \hat S^{(1)}_x,
\]
explicitly noting the choice of open boundary conditions. Our
formalization uses the Pauli convention
$\hat\sigma^{(\alpha)} = 2\hat S^{(\alpha)}$ and introduces an
explicit coupling $J$ in front of the bond term, so our $H_{J,h}^{N+1}$
matches Tasaki's (3.3.1) up to the constants $J$ and the Pauli rescaling.

\subsection{Gibbs state instantiation}
\label{subsec:ising-gibbs}

\emph{Reference:} Tasaki~\cite{Tasaki:QMB}, §3.3, p.~80.

\begin{definition}[\texttt{quantumIsingGibbsState}]
The Gibbs state of the open-boundary 1D quantum Ising chain on
$\mathrm{Fin}(N+1)$ sites with real $J$, $h$ and inverse temperature
$\beta$ is
\[
  \rho_\beta^{\mathrm{Ising}}
  \;:=\;
  \frac{e^{-\beta H_{J,h}^{N+1}}}{\operatorname{Tr}\,e^{-\beta H_{J,h}^{N+1}}},
\]
i.e.\ \texttt{gibbsState $\beta$ (quantumIsingHamiltonian $N$ $J$ $h$)}.
\end{definition}

\begin{theorem}[Hermiticity]
$\rho_\beta^{\mathrm{Ising}}$ is Hermitian. Lean name:
\texttt{quantumIsingGibbsState\_isHermitian}.
\end{theorem}

\begin{proof}[Proof sketch]
Apply the abstract \texttt{gibbsState\_isHermitian} with the Hermitian
$H_{J,h}^{N+1}$ from \texttt{quantumIsingHamiltonian\_isHermitian}.
\end{proof}

\begin{theorem}[Commutativity with the Hamiltonian]
$[\rho_\beta^{\mathrm{Ising}}, H_{J,h}^{N+1}] = 0$. Lean name:
\texttt{quantumIsingGibbsState\_commute\_hamiltonian}.
\end{theorem}

\begin{proof}[Proof sketch]
Specialise the abstract \texttt{gibbsState\_commute\_hamiltonian} to
$H_{J,h}^{N+1}$.
\end{proof}

\begin{theorem}[Infinite-temperature closed form]
At $\beta = 0$, for any observable $A$,
\[
  \langle A \rangle_0^{\mathrm{Ising}}
  \;=\; \dim^{-1}\,\operatorname{Tr}(A),
\]
independent of the parameters $J, h$. Lean name:
\texttt{quantumIsingGibbsExpectation\_zero}.
\end{theorem}

\begin{proof}[Proof sketch]
Specialise \texttt{gibbsExpectation\_zero} to $H_{J,h}^{N+1}$; at
$\beta = 0$ the Hamiltonian drops out because $e^{-0\cdot H} = I$.
\end{proof}

\begin{theorem}[Ising expectation realness, conservation, energy realness]
For the open Ising Hamiltonian $H_{J,h}^{N+1}$:
\begin{enumerate}
  \item For any Hermitian observable $O$,
    $\Im\langle O \rangle_\beta = 0$.
  \item For any observable $A$, $\langle [H, A] \rangle_\beta = 0$.
  \item $\Im\langle H \rangle_\beta = 0$.
  \item For any Hermitian observable $O$,
    $\Im\langle H \cdot O \rangle_\beta = 0$.
  \item $\Im\langle H^2 \rangle_\beta = 0$.
  \item $\Im\langle H^n \rangle_\beta = 0$ for any $n \in \mathbb{N}$.
  \item For any Hermitian $A, B$,
    $\Im\langle A B + B A \rangle_\beta = 0$.
  \item For any Hermitian $A, B$,
    $\Re\langle A B - B A \rangle_\beta = 0$.
  \item $\Im\operatorname{Var}_\beta(H) = 0$ (energy variance real).
\end{enumerate}
Lean names:
\texttt{quantumIsingGibbsExpectation\_im\_of\_isHermitian},
\texttt{quantumIsingGibbsExpectation\_commutator\_hamiltonian},
\texttt{quantumIsingGibbsExpectation\_hamiltonian\_im},
\texttt{quantumIsingGibbsExpectation\_mul\_hamiltonian\_im},
\texttt{quantumIsingGibbsExpectation\_hamiltonian\_sq\_im},
\texttt{quantumIsingGibbsExpectation\_hamiltonian\_pow\_im},
\texttt{quantumIsingGibbsExpectation\_anticommutator\_im},
\texttt{quantumIsingGibbsExpectation\_commutator\_re},
\texttt{quantumIsingGibbsHamiltonianVariance\_im}.
\end{theorem}

\begin{proof}[Proof sketch]
Direct specialisations of
\texttt{gibbsExpectation\_im\_of\_isHermitian} and
\texttt{gibbsExpectation\_commutator\_hamiltonian} at $H_{J,h}^{N+1}$,
using \texttt{quantumIsingHamiltonian\_isHermitian}.
\end{proof}

\begin{theorem}[\texttt{quantumIsingGibbsExpectation\_self\_eq}]
The Ising-chain energy expectation splits into bond and
transverse-field contributions:
\[
  \langle H_{J,h}^{N+1} \rangle_\beta
  \;=\;
  -J \cdot \sum_{i=0}^{N-1}
      \langle \sigma^z_i \sigma^z_{i+1} \rangle_\beta
  \;+\;
  -h \cdot \sum_{i=0}^{N}
      \langle \sigma^x_i \rangle_\beta.
\]
\end{theorem}

\begin{proof}[Proof sketch]
Unfold $H_{J,h}^{N+1} = -J \cdot \sum \sigma^z\sigma^z + (-h) \cdot \sum \sigma^x$
and apply
\texttt{gibbsExpectation\_add},
\texttt{gibbsExpectation\_smul} (twice, for $-J$ and $-h$), and
\texttt{gibbsExpectation\_sum} (twice, on the two index sets).
\end{proof}

%======================================================================
%======================================================================
\section{Tasaki §2.1/§2.2 supplementary content}
\label{sec:supplementary}
%======================================================================

The following additional theorems (added in PR \#16) extend the
formalization of Tasaki \S2.1 and \S2.2.

\subsection{Tasaki \S2.1 supplementary}

\begin{itemize}
\item \texttt{spinHalfRot\{1,2,3\}\_conj\_spinHalfOp\{2,3,1\}}: general-$\theta$
  cyclic conjugation $(\hat U^{(\alpha)}_\theta)^\dagger \hat S^{(\beta)}
  \hat U^{(\alpha)}_\theta = \cos\theta\,\hat S^{(\beta)} - \sin\theta\,\hat S^{(\gamma)}$
  (Tasaki eq.~(2.1.16)), 6 cyclic + 6 anti-cyclic + 3 same-axis (Tasaki
  eq.~(2.1.13)-(2.1.15)).
\item \texttt{spinHalfRot\{1,2,3\}\_half\_pi\_conj\_spinHalfOp\_*}:
  $\pi/2$-rotation conjugation (Tasaki eq.~(2.1.22)), 6 cases.
\item \texttt{spinHalfRot\{1,2,3\}\_eq\_exp}: $\hat U^{(\alpha)}_\theta =
  \exp(-i\theta\hat S^{(\alpha)})$ for all three axes (Tasaki
  Problem 2.1.b). Axis 3 uses \texttt{Matrix.exp\_diagonal} +
  Euler; axes 1, 2 are derived from axis 3 via \texttt{Matrix.exp\_conj}
  with Hadamard / y-axis basis-change unitaries.
\item \texttt{hadamard}, \texttt{yDiag}: basis-change unitaries with
  \texttt{hadamard\_mul\_self}, \texttt{yDiag\_mul\_yDiagAdj}, etc.
  Used to lift the axis-3 exponential proof to axes 1, 2.
\item \texttt{spinHalfRot3\_mul\_spinHalfRot2\_mulVec\_spinHalfUp}:
  coherent state $\hat U^{(3)}_\varphi \hat U^{(2)}_\theta |\psi^\uparrow\rangle$
  (Tasaki Problem 2.1.d).
\item \texttt{spinHalfRot3\_half\_pi\_mul\_spinHalfRot2\_half\_pi\_mulVec\_spinHalfUp}:
  the $\theta=\varphi=\pi/2$ specialization (Tasaki Problem 2.1.e).
\item \texttt{spinHalfDotVec}, \texttt{spinHalfDotVec\_isHermitian},
  \texttt{spinHalfRot3\_commute\_spinHalfOp3\_smul}: vector inner
  product $\hat S \cdot v$ and same-axis commutation (Tasaki
  eq.~(2.1.19)-(2.1.20), partial).
\item \texttt{rot3D\{1,2,3\}}, \texttt{rot3D\{1,2,3\}\_zero},
  \texttt{rot3D\{1,2,3\}\_pi}: explicit 3D rotation matrices
  $R^{(\alpha)}_\theta$ (Tasaki eq.~(2.1.11)).
\item \texttt{rot3D\{1,2,3\}Pi}, \texttt{rot3D\{1,2,3\}Pi\_sq},
  \texttt{rot3D\{1,2,3\}Pi\_mul\_rot3D\{2,3,1\}Pi},
  \texttt{rot3D\{1,2,3\}Pi\_comm\_*}: $\pi$-rotation matrices
  $R^{(\alpha)}_\pi$ and their group structure (Tasaki eq.~(2.1.28),
  Problem 2.1.f).
\item \texttt{spinOnePiRot\{1,2,3\}}, \texttt{spinOnePiRot3\_eq},
  \texttt{spinOnePiRot\{1,2,3\}\_sq}, \texttt{spinOnePiRot\{1,2,3\}\_comm\_*}:
  $S = 1$ $\pi$-rotation matrices and their integer-spin Z$_2 \times$ Z$_2$
  relations (Tasaki eq.~(2.1.31)/(2.1.33)/(2.1.34)).
\item \texttt{spinOneOpPlus}, \texttt{spinOneOpMinus},
  \texttt{spinOneOp3\_mulVec\_*}, \texttt{spinOneOpPlus/Minus\_mulVec\_*},
  \texttt{spinOneOpPlus/Minus\_conjTranspose}: $S = 1$ basis,
  raising/lowering, adjoints (Tasaki eq.~(2.1.2)-(2.1.6) for $S=1$).
\item \texttt{spinOneRot\{1,2,3\}},
  \texttt{spinOneRot\{1,2,3\}\_zero},
  \texttt{spinOneRot\{1,2,3\}\_pi},
  \texttt{spinOnePiRot\{1,2\}\_eq},
  \texttt{spinOneOp\{1,2\}\_mul\_self}:
  $S = 1$ closed-form rotation $\hat U^{(\alpha)}_\theta = \hat 1 -
  i\sin\theta\,\hat S^{(\alpha)} - (1-\cos\theta)(\hat S^{(\alpha)})^2$
  for all 3 axes (Tasaki Problem 2.1.c, eq.~(2.1.30)) and the
  squared-operator helpers used in the boundary checks $\hat U^{(\alpha)}_\pi
  = \hat u_\alpha$.
\item \texttt{spinOnePiRot\{1,2,3\}\_mulVec\_spinOne\{Plus,Zero,Minus\}}:
  action of the $S = 1$ $\pi$-rotation matrices on the
  $|\psi^{+1,0,-1}\rangle$ basis (Tasaki eq.~(2.1.34) / Problem 2.1.g
  for $S = 1$).
\item \texttt{spinOneProj\{Plus,Zero,Minus\}},
  \texttt{spinOneProj\{Plus,Zero,Minus\}\_eq\_polynomial},
  \texttt{spinOneUnit\{01,02,10,12,20,21\}},
  \texttt{spinOneUnit*\_eq\_polynomial},
  \texttt{spinOne\_decomposition}: Tasaki Problem 2.1.a for $S = 1$.
  Three diagonal projectors $|\psi^\sigma\rangle\langle\psi^\sigma|$
  expressed as polynomials in $\hat S^{(3)}$ via Lagrange
  interpolation, six off-diagonal matrix units expressed as
  $\hat S^{\pm}$ acting on diagonal projectors, and the spanning
  theorem that every $3 \times 3$ complex matrix is a linear
  combination of the nine matrix units (entry-wise).
  See \texttt{Quantum/SpinOneDecomp.lean}.
\item \texttt{Quantum/Z2Z2.lean}: documentation module unifying the
  Z$_2 \times$ Z$_2$ structure (Tasaki eq.~(2.1.27)-(2.1.34)) for
  both $S=1/2$ (projective) and $S=1$ (genuine) representations.
\end{itemize}

\subsection{Tasaki \S2.2 supplementary}

\begin{itemize}
\item \texttt{onSite\_add}, \texttt{onSite\_sub}, \texttt{onSite\_zero},
  \texttt{onSite\_smul}, \texttt{onSite\_one}: linearity and identity
  preservation of \texttt{onSite}.
\item \texttt{onSite\_mul\_onSite\_same},
  \texttt{onSite\_commutator\_same}: same-site product / commutator
  (Tasaki eq.~(2.2.6) diagonal).
\item \texttt{spinHalfOp\{1,2,3\}\_onSite\_commutator\_spinHalfOp\{2,3,1\}\_onSite}:
  same-site spin commutators.
\item \texttt{onSite\_commutator\_totalOnSite}: generic
  $[\texttt{onSite}\,x\,A, \sum_z \texttt{onSite}\,z\,B]
   = \texttt{onSite}\,x\,[A, B]$.
\item \texttt{basisVec}, \texttt{onSite\_mulVec\_basisVec}: basis
  states and the explicit tensor-product action (Tasaki
  eq.~(2.2.1)/(2.2.4)).
\item \texttt{totalSpinHalfOpPlus}, \texttt{totalSpinHalfOpMinus},
  \texttt{totalSpinHalfOpPlus\_eq\_add},
  \texttt{totalSpinHalfOpMinus\_eq\_sub},
  \texttt{totalSpinHalfOpPlus/Minus\_conjTranspose}:
  $\hat S^{\pm}_{\rm tot}$ and adjoints (Tasaki eq.~(2.2.8)).
\item \texttt{totalSpinHalfOp\{1,2,3\}\_commutator\_totalSpinHalfOp\{2,3,1\}}:
  total-spin commutators.
\item \texttt{totalSpinHalfOp3\_commutator\_totalSpinHalfOpPlus/Minus}:
  Cartan ladder $[\hat S_{\rm tot}^{(3)}, \hat S^{\pm}_{\rm tot}]
   = \pm\hat S^{\pm}_{\rm tot}$.
\item \texttt{totalSpinHalfOpPlus\_commutator\_totalSpinHalfOpMinus}:
  $[\hat S^+_{\rm tot}, \hat S^-_{\rm tot}] = 2\hat S^{(3)}_{\rm tot}$.
\item \texttt{totalSpinHalfSquared}, \texttt{totalSpinHalfSquared\_isHermitian},
  \texttt{totalSpinHalfSquared\_commutator\_*}:
  Casimir operator $(\hat S_{\rm tot})^2$ and its $\mathrm{SU}(2)$
  invariance (Tasaki eq.~(2.2.12)).
\item \texttt{spinSign}, \texttt{magnetization}: total magnetization
  $|\sigma| := \sum_x \mathrm{sign}(\sigma_x)$ (Tasaki eq.~(2.2.2)).
\item \texttt{spinHalfSign}, \texttt{onSite\_spinHalfOp3\_mulVec\_basisVec},
  \texttt{totalSpinHalfOp3\_mulVec\_basisVec},
  \texttt{totalSpinHalfOp3\_mulVec\_basisVec\_eq\_magnetization}:
  $\hat S^{(3)}_{\rm tot}$ eigenvalues on basis states (Tasaki
  eq.~(2.2.10)).
\item \texttt{onSite\_spinHalfOpPlus/Minus\_mulVec\_basisVec},
  \texttt{totalSpinHalfOpPlus/Minus\_mulVec\_basisVec}: raising and
  lowering action on basis states.
\item \texttt{spinHalfDot}, \texttt{spinHalfDot\_eq\_plus\_minus}:
  two-site inner product $\hat S_x \cdot \hat S_y$ and its
  raising/lowering decomposition (Tasaki eq.~(2.2.16)).
\item \texttt{spinHalfDot\_comm}, \texttt{spinHalfDot\_self},
  \texttt{spinHalfDot\_isHermitian},
  \texttt{spinHalfPairSpinSq}, \texttt{spinHalfPairSpinSq\_eq}:
  symmetry, $\hat S_x \cdot \hat S_x = (3/4)\cdot 1$, Hermiticity, and
  the squared-sum identity (Tasaki eq.~(2.2.18)).
\item \texttt{totalSpinHalfSquared\_eq\_sum\_dot}:
  $(\hat S_{\rm tot})^2 = \sum_{x,y} \hat S_x \cdot \hat S_y$.
\item \texttt{spinHalfDot\_commutator\_totalSpinHalfOp\{1,2,3\}},
  \texttt{spinHalfDot\_commutator\_totalSpinHalfOpPlus/Minus}:
  $\mathrm{SU}(2)$ invariance of $\hat S_x \cdot \hat S_y$
  (Tasaki eq.~(2.2.17)).
\item \texttt{IsInMagnetizationSubspace},
  \texttt{magnetizationSubspace}, \texttt{basisVec\_mem\_magnetizationSubspace},
  \texttt{magnetizationSubspace\_disjoint},
  \texttt{iSup\_magnetizationSubspace\_eq\_top},
  \texttt{magnetizationSubspace\_eq\_eigenspace},
  \texttt{magnetizationSubspace\_iSupIndep},
  \texttt{magnetizationSubspace\_isInternal}:
  magnetization subspace as a \texttt{Submodule}, basis-state
  membership, pairwise disjointness, supremum-equals-top, identification
  with \texttt{Module.End.eigenspace} (used to inherit
  \texttt{eigenspaces\_iSupIndep}), and the full internal direct-sum
  decomposition $\mathcal H = \bigoplus_M \mathcal H_M$
  (Tasaki eq.~(2.2.9)/(2.2.10)).
\item \texttt{totalSpinHalfRot\{1,2,3\}Pi},
  \texttt{totalSpinHalfRot\{1,2,3\}\_theta},
  \texttt{totalSpinHalfRot\{1,2,3\}\_zero},
  \texttt{totalSpinHalfRot\{1,2,3\}Pi\_eq}: global rotation
  $\hat U^{(\alpha)}_{\rm tot,\theta}$ via \texttt{Finset.noncommProd}
  (Tasaki eq.~(2.2.11)).
\item \texttt{totalSpinHalfRot\{1,2,3\}Pi\_mul\_totalSpinHalfRot\{2,3,1\}Pi}:
  $\hat U^{(\alpha)}_{\pi,{\rm tot}} \cdot
   \hat U^{(\beta)}_{\pi,{\rm tot}}
   = \hat U^{(\gamma)}_{\pi,{\rm tot}}$ (Tasaki Problem 2.2.a).
\item \texttt{onSiteRingHom}, \texttt{onSiteLinearMap},
  \texttt{continuous\_onSite}, \texttt{onSite\_pow}: site embedding
  packaged as a ring/linear/continuous map and lifted to powers.
\item \texttt{totalSpinHalfRot\{1,2,3\}Pi\_two\_site},
  \texttt{totalSpinHalfRot\{1,2,3\}\_two\_site}: for $\Lambda = {\rm Fin}\,2$,
  the global rotation factors as $\texttt{onSite}\,0\,(\hat U^{(\alpha)}_\theta)
  \cdot \texttt{onSite}\,1\,(\hat U^{(\alpha)}_\theta)$ (Tasaki Problem
  2.2.b at $\theta = \pi$ and general $\theta$).
\item \texttt{spinHalfDot\_mulVec\_basisVec\_parallel},
  \texttt{spinHalfDot\_mulVec\_basisVec\_both\_up/down}: for $x \neq y$
  with $\sigma_x = \sigma_y$, $\hat S_x \cdot \hat S_y |\sigma\rangle =
  (1/4)|\sigma\rangle$ (Tasaki eq.~(2.2.19), parallel case; gives the
  triplet $m = \pm 1$ eigenvalue).
\item \texttt{basisSwap}, \texttt{basisSwap\_involutive},
  \texttt{basisSwap\_antiparallel}: site-swap of $\sigma$ at $x, y$,
  involution, and preservation of anti-parallelism.
\item \texttt{spinHalfDot\_mulVec\_basisVec\_antiparallel}: for
  $\sigma_x \neq \sigma_y$, $\hat S_x \cdot \hat S_y |\sigma\rangle =
  (1/2)|\mathrm{swap}\,\sigma\rangle - (1/4)|\sigma\rangle$.
\item \texttt{spinHalfDot\_mulVec\_singlet}: singlet eigenvalue
  $\hat S_x \cdot \hat S_y (|\sigma\rangle - |\mathrm{swap}\,\sigma\rangle)
  = -(3/4)(|\sigma\rangle - |\mathrm{swap}\,\sigma\rangle)$ (Tasaki
  eq.~(2.2.19), $S = 0$ singlet).
\item \texttt{spinHalfDot\_mulVec\_triplet\_anti}: triplet $m = 0$
  eigenvalue $\hat S_x \cdot \hat S_y (|\sigma\rangle + |\mathrm{swap}\,\sigma\rangle)
  = (1/4)(|\sigma\rangle + |\mathrm{swap}\,\sigma\rangle)$ (Tasaki
  eq.~(2.2.19), $S = 1$, $m = 0$).
\item \texttt{heisenbergHamiltonian},
  \texttt{heisenbergHamiltonian\_commutator\_totalSpinHalfOp\{1,2,3\}},
  \texttt{heisenbergHamiltonian\_commutator\_totalSpinHalfOp\{Plus,Minus\}}:
  the general Heisenberg-type Hamiltonian
  $H = \sum_{x,y} J(x,y)\,\hat S_x \cdot \hat S_y$ commutes with
  $\hat S^{(\alpha)}_{\rm tot}$ and $\hat S^{\pm}_{\rm tot}$
  (Tasaki eq.~(2.2.13), SU(2) invariance).
\item \texttt{totalSpinHalfSquared\_mulVec\_basisVec\_const},
  \texttt{totalSpinHalfSquared\_mulVec\_basisVec\_all\_up/down}:
  $\hat S_{\rm tot}^2 |s\,s\,\dots\,s\rangle = (N(N+2)/4)|s\,s\,\dots\,s\rangle$
  for any constant $s : {\rm Fin}\,2$ — the Casimir eigenvalue at
  the maximum total spin $S = N/2$.
\item \texttt{onSite\_exp\_eq\_exp\_onSite},
  \texttt{totalSpinHalfRot\{1,2,3\}\_eq\_exp}: Tasaki eq.~(2.2.11),
  $\hat U^{(\alpha)}_{\theta,{\rm tot}}
   = \exp[-i\theta \hat S^{(\alpha)}_{\rm tot}]$. The proof composes
  the per-site exp form, the bridge
  $\texttt{onSite}\,x\,(\exp A) = \exp(\texttt{onSite}\,x\,A)$
  (proved by installing the Frobenius normed structure locally and
  invoking \texttt{NormedSpace.map\_exp} via
  \texttt{LinearMap.continuous\_of\_finiteDimensional}), and
  \texttt{Matrix.exp\_sum\_of\_commute} for the commuting site
  embeddings.
\item \texttt{totalSpinHalfRot\{1,2,3\}\_commute\_of\_commute},
  \texttt{totalSpinHalfOp\{Plus,Minus\}\_exp\_commute\_of\_commute}:
  Tasaki eq.~(2.2.12)$\to$ commutativity with $\hat U$.
  $\mathrm{Commute}\,A\,\hat S^{(\alpha)}_{\rm tot}
    \Rightarrow \mathrm{Commute}\,A\,\hat U^{(\alpha)}_{\theta,{\rm tot}}$
  (and its $\hat S^{\pm}_{\rm tot}$ ladder analogue). Proof: rewrite
  $\hat U$ via \texttt{\_eq\_exp}, then invoke
  \texttt{Commute.smul\_right} and \texttt{Commute.exp\_right}.
\item \texttt{totalSpinHalfRot\{1,2,3\}\_conjTranspose\_mul\_self}:
  unitarity $(\hat U^{(\alpha)}_{\theta,{\rm tot}})^\dagger
  \hat U^{(\alpha)}_{\theta,{\rm tot}} = \hat 1$. Derived from
  \texttt{NormedSpace.exp\_mem\_unitary\_of\_mem\_skewAdjoint} after
  recognising $-i\theta\,\hat S^{(\alpha)}_{\rm tot}$ as skew-adjoint
  (the helper \texttt{neg\_I\_mul\_real\_skewAdjoint\_smul\_isHermitian}).
  Same Frobenius topology bridge with explicit \texttt{@}-notation as
  \texttt{onSite\_exp\_eq\_exp\_onSite}.
\item \texttt{totalSpinHalfRot\{1,2,3\}\_conj\_eq\_self\_of\_commute}:
  Tasaki eq.~(2.2.13) finite form,
  $\mathrm{Commute}\,A\,\hat S^{(\alpha)}_{\rm tot}
    \Rightarrow (\hat U^{(\alpha)}_{\theta,{\rm tot}})^\dagger\,A\,
    \hat U^{(\alpha)}_{\theta,{\rm tot}} = A$. Combines
  \texttt{\_commute\_of\_commute} and \texttt{\_conjTranspose\_mul\_self}
  via standard $\hat U^\dagger A \hat U = \hat U^\dagger \hat U A
  = A$.
\item \textbf{Problem 2.2.c} (Tasaki §2.2, eq.~(2.2.15), p.~24):
  The SU(2)-averaged two-site state is the spin singlet.

  \textit{Statement.}
  \[
    \frac{1}{4\pi}
    \int_0^{2\pi}\!\!d\varphi\int_0^\pi\!\!d\theta\,\sin\theta\;
    \hat U^{(3)}_{\varphi,{\rm tot}}\,\hat U^{(2)}_{\theta,{\rm tot}}\,
    |\!\uparrow\downarrow\rangle
    \;=\;
    \tfrac{1}{2}\bigl(|\!\uparrow\downarrow\rangle
                     - |\!\downarrow\uparrow\rangle\bigr),
  \]
  where $|\!\uparrow\downarrow\rangle := |\sigma_0=0,\sigma_1=1\rangle
  \in \mathcal{H}_{\rm Fin\,2}$.

  \textit{Proof outline.}
  \begin{enumerate}
  \item \textbf{Factorization of the global rotation.}
    Using \texttt{noncommProd\_mul\_distrib} and the \texttt{onSite}
    commutativity at distinct sites, the product
    $\hat U^{(3)}_{\varphi,{\rm tot}}\,\hat U^{(2)}_{\theta,{\rm tot}}$
    factors as
    $(\hat U^{(3)}_\varphi)_0\,(\hat U^{(2)}_\theta)_0\;
     \cdot\;
     (\hat U^{(3)}_\varphi)_1\,(\hat U^{(2)}_\theta)_1$,
    where the subscript denotes the site embedding via \texttt{onSite}.

  \item \textbf{Tensor-product decomposition.}
    The two-site action on $|\sigma_0,\sigma_1\rangle$ produces
    a sum weighted by $M(\tau_0,\sigma_0)\cdot M(\tau_1,\sigma_1)$,
    where $M = \hat U^{(3)}_\varphi\,\hat U^{(2)}_\theta \in M_2(\mathbb{C})$
    (the single-site combined rotation matrix).

  \item \textbf{Explicit single-site action (Problem 2.1.d).}
    From the coherent-state results (\texttt{spinHalfRot3} then
    \texttt{spinHalfRot2}):
    \begin{align*}
      M\,|\!\uparrow\rangle &=
        e^{-i\varphi/2}\cos\tfrac{\theta}{2}|\!\uparrow\rangle
        - e^{-i\varphi/2}\sin\tfrac{\theta}{2}|\!\downarrow\rangle,\\
      M\,|\!\downarrow\rangle &=
        e^{+i\varphi/2}\sin\tfrac{\theta}{2}|\!\uparrow\rangle
        + e^{+i\varphi/2}\cos\tfrac{\theta}{2}|\!\downarrow\rangle.
    \end{align*}

  \item \textbf{Integration over four configuration cases.}
    Expanding the two-site product $M|\sigma_0\rangle\otimes M|\sigma_1\rangle$
    for the input $(\sigma_0,\sigma_1)=(\uparrow,\downarrow)$ yields
    four output components $|\tau_0,\tau_1\rangle$:
    \begin{itemize}
      \item $|\!\uparrow\downarrow\rangle$: the phase $e^{-i\varphi/2}
        \cdot e^{+i\varphi/2} = 1$ cancels; the angular integral gives
        $\int_0^\pi \sin\theta\,\cos^2\!\tfrac{\theta}{2}\,d\theta = 1$,
        contributing $+\tfrac{1}{2}$.
      \item $|\!\downarrow\uparrow\rangle$: the phase again cancels;
        the angular integral gives
        $-\int_0^\pi \sin\theta\,\sin^2\!\tfrac{\theta}{2}\,d\theta = -1$,
        contributing $-\tfrac{1}{2}$.
      \item $|\!\uparrow\uparrow\rangle$: contains the factor
        $e^{-i\varphi/2}\cdot e^{+i\varphi/2}\cdot e^{+i\varphi}
         = e^{+i\varphi}$; the azimuthal integral
        $\int_0^{2\pi}e^{\pm i\varphi}\,d\varphi = 0$ forces this to zero.
      \item $|\!\downarrow\downarrow\rangle$: contains $e^{-i\varphi}$;
        again vanishes by the azimuthal integral.
    \end{itemize}

  \item \textbf{Conclusion.}
    The normalisation factor $\tfrac{1}{4\pi}\cdot 2\pi = \tfrac{1}{2}$
    combines with the angular integrals to give
    $\tfrac{1}{2}|\!\uparrow\downarrow\rangle - \tfrac{1}{2}|\!\downarrow\uparrow\rangle$,
    which is the (unnormalised) spin singlet up to the factor
    $1/\sqrt{2}$ (the statement records the un-normalised form as
    in Tasaki eq.~(2.2.15)).
  \end{enumerate}

  \textit{Key helper lemmas.}
  \begin{itemize}
    \item \texttt{integral\_cexp\_I\_mul\_zero\_two\_pi}:
      $\int_0^{2\pi}e^{i\varphi}\,d\varphi = 0$.
    \item \texttt{integral\_cexp\_neg\_I\_mul\_zero\_two\_pi}:
      $\int_0^{2\pi}e^{-i\varphi}\,d\varphi = 0$.
    \item \texttt{integral\_sin\_mul\_cos\_sq\_half\_zero\_pi}:
      $\int_0^\pi\sin\theta\,\cos^2\!\tfrac{\theta}{2}\,d\theta = 1$.
    \item \texttt{integral\_sin\_mul\_sin\_sq\_half\_zero\_pi}:
      $\int_0^\pi\sin\theta\,\sin^2\!\tfrac{\theta}{2}\,d\theta = 1$.
  \end{itemize}

  Reference: Tasaki~\cite{Tasaki:QMB}, §2.2, Problem~2.2.c,
  eq.~(2.2.15), p.~24.
\end{itemize}

%======================================================================
\section{Gibbs state (Tasaki §3.3, pp.\,75--78)}
\label{sec:gibbs}
%======================================================================

\emph{Source file:} \texttt{LatticeSystem/Quantum/GibbsState.lean}.

\subsection{Definitions}

Let $n \in \mathbb{N}$ and let $H \in M_n(\mathbb{C})$ be a Hermitian matrix
(the Hamiltonian). Fix an inverse temperature $\beta \in \mathbb{R}$.

\begin{definition}[\texttt{gibbsExp}]
The \emph{Gibbs exponential} is
\[
  e^{-\beta H} \;:=\; \exp(-\beta H)
  \;\in\; M_n(\mathbb{C}),
\]
computed via the matrix exponential \texttt{Matrix.exp} applied to
$-\beta \cdot H$.
\end{definition}

\begin{definition}[\texttt{partitionFn}]
The \emph{partition function} is
\[
  Z \;:=\; \operatorname{Tr}(e^{-\beta H}).
\]
Lean name: \texttt{partitionFn~$\beta$~$H$}.
\end{definition}

\begin{definition}[\texttt{gibbsState}]
The \emph{Gibbs state} (density matrix) is
\[
  \rho \;:=\; \frac{1}{Z}\,e^{-\beta H}.
\]
Lean name: \texttt{gibbsState~$\beta$~$H$}.
\end{definition}

\begin{definition}[\texttt{gibbsExpectation}]
The \emph{thermal expectation value} of an observable $O \in M_n(\mathbb{C})$
is
\[
  \langle O \rangle_\beta \;:=\; \operatorname{Tr}(\rho\, O).
\]
Lean name: \texttt{gibbsExpectation~$\beta$~$H$~$O$}.
\end{definition}

\subsection{Formalized facts}

\begin{lemma}[\texttt{gibbsExp\_isHermitian}]
If $H$ is Hermitian then $e^{-\beta H}$ is Hermitian.
\end{lemma}

\begin{lemma}[\texttt{gibbsState\_trace}]
$\operatorname{Tr}(\rho) = 1$.
\end{lemma}

\begin{proof}[Proof sketch]
By linearity of the trace,
$\operatorname{Tr}(\rho)
 = \operatorname{Tr}\!\bigl(\tfrac{1}{Z} e^{-\beta H}\bigr)
 = \tfrac{1}{Z}\operatorname{Tr}(e^{-\beta H})
 = \tfrac{Z}{Z} = 1$,
provided $Z \neq 0$.
\end{proof}

\begin{lemma}[\texttt{gibbsState\_isHermitian}]
$\rho$ is Hermitian.
\end{lemma}

\begin{lemma}[\texttt{gibbsExpectation\_one}]
$\langle \mathbf{1} \rangle_\beta = 1$, i.e.\
$\operatorname{Tr}(\rho \cdot \mathbf{1}) = 1$.
\end{lemma}

\begin{lemma}[\texttt{gibbsExp\_zero}]
At inverse temperature $\beta = 0$,
\[
  e^{-0 \cdot H} \;=\; \mathbf{1}.
\]
\end{lemma}

\begin{proof}[Proof sketch]
$-0 \cdot H = 0$, so $\exp(-0 \cdot H) = \exp(0) = \mathbf{1}$ by
\texttt{Matrix.exp\_zero}.
\end{proof}

\subsubsection*{One-parameter group in $\beta$}

\begin{lemma}[\texttt{gibbsExp\_add}]
For all $\beta_1, \beta_2 \in \mathbb{R}$ and Hamiltonian $H$,
$e^{-(\beta_1 + \beta_2) H} = e^{-\beta_1 H} \cdot e^{-\beta_2 H}$.
\end{lemma}

\begin{proof}[Proof sketch]
The arguments $-\beta_1 H$ and $-\beta_2 H$ commute (both are scalar
multiples of the same $H$). Apply
\texttt{Matrix.exp\_add\_of\_commute} after rewriting
$-(\beta_1+\beta_2) H = -\beta_1 H + -\beta_2 H$.
\end{proof}

\begin{lemma}[\texttt{gibbsExp\_add\_of\_commute\_hamiltonians}]
For commuting Hamiltonians $H_1, H_2$,
$e^{-\beta (H_1 + H_2)} = e^{-\beta H_1} \cdot e^{-\beta H_2}$.
\end{lemma}

\begin{proof}[Proof sketch]
The scalar multiples $-\beta H_1$ and $-\beta H_2$ commute (since
$[H_1, H_2] = 0$).  Apply \texttt{Matrix.exp\_add\_of\_commute}
after $-\beta(H_1 + H_2) = -\beta H_1 + -\beta H_2$ via
\texttt{smul\_add}.

The basic identity behind exact decoupling, mean-field block
decompositions, and the simplest cases of the Suzuki–Trotter formula
(where the commutator-correction terms vanish for commuting
decompositions).
\end{proof}

\begin{lemma}[\texttt{gibbsExp\_neg\_mul\_self} / \texttt{gibbsExp\_self\_mul\_neg}]
$e^{\beta H} \cdot e^{-\beta H} = e^{-\beta H} \cdot e^{\beta H} = I$.
\end{lemma}

\begin{proof}[Proof sketch]
Specialise \texttt{gibbsExp\_add} with $\beta_2 = -\beta_1$ (or
$\beta_1 = -\beta_2$) and use \texttt{gibbsExp\_zero} to reduce
$e^{0} = I$.
\end{proof}

\begin{lemma}[\texttt{gibbsExp\_isUnit}]
$e^{-\beta H}$ is a unit (invertible element) for every $\beta$ and
every Hamiltonian $H$.
\end{lemma}

\begin{proof}[Proof sketch]
Witness the inverse via $e^{-\beta H} \cdot e^{\beta H} = I$ from the
previous lemma; \texttt{IsUnit.of\_mul\_eq\_one} produces the
\texttt{IsUnit} certificate.
\end{proof}

\begin{lemma}[\texttt{gibbsExp\_ne\_zero}]
$e^{-\beta H} \neq 0$.
\end{lemma}

\begin{proof}[Proof sketch]
Direct corollary of \texttt{gibbsExp\_isUnit} via
\texttt{IsUnit.ne\_zero}: a unit in a nontrivial ring is nonzero.
\end{proof}

\begin{lemma}[\texttt{gibbsState\_ne\_zero}]
If $Z(\beta, H) \neq 0$, then $\rho_\beta \neq 0$.
\end{lemma}

\begin{proof}[Proof sketch]
$\rho_\beta = (1/Z) \cdot e^{-\beta H}$.  The scalar $1/Z$ is nonzero
by hypothesis (\texttt{one\_div\_ne\_zero}), and the operator factor
is nonzero by \texttt{gibbsExp\_ne\_zero}; the scalar-multiple of
two nonzeros is nonzero (\texttt{smul\_ne\_zero}).
\end{proof}

\begin{lemma}[\texttt{gibbsState\_inv}]
If $Z(\beta, H) \neq 0$, then
$\rho_\beta^{-1} = Z(\beta, H) \cdot e^{\beta H}$
(in our convention, $e^{\beta H} = \texttt{gibbsExp}\,(-\beta)\,H$).
\end{lemma}

\begin{proof}[Proof sketch]
Apply \texttt{Matrix.inv\_eq\_left\_inv} to the witness
\[
  (Z \cdot e^{\beta H})\,\bigl((1/Z) \cdot e^{-\beta H}\bigr)
  \;=\;
  (Z \cdot 1/Z) \cdot (e^{\beta H} \cdot e^{-\beta H})
  \;=\; 1 \cdot 1 \;=\; 1,
\]
discharged by \texttt{smul\_mul\_smul\_comm},
\texttt{gibbsExp\_neg\_mul\_self} (which gives
$e^{\beta H}\cdot e^{-\beta H} = 1$), \texttt{mul\_one\_div\_cancel}
applied with the partition-function nonzero hypothesis, and
\texttt{one\_smul}.

This generalises \texttt{gibbsState\_zero\_inv}: at $\beta = 0$,
$Z = \dim$ and $e^{0 \cdot H} = I$, recovering $\rho_0^{-1} = \dim \cdot I$.
\end{proof}

\begin{lemma}[\texttt{partitionFn\_smul\_gibbsState\_eq\_gibbsExp}]
If $Z(\beta, H) \neq 0$, then
$Z(\beta, H) \cdot \rho_\beta = e^{-\beta H}$.
\end{lemma}

\begin{proof}[Proof sketch]
Unfold $\rho_\beta = (1/Z) \cdot e^{-\beta H}$ in the LHS:
$Z \cdot ((1/Z) \cdot e^{-\beta H})
 = (Z \cdot 1/Z) \cdot e^{-\beta H}$
via \texttt{smul\_smul}.  Then \texttt{mul\_one\_div\_cancel} (with
the partition-function nonzero hypothesis) reduces the scalar to $1$,
and \texttt{one\_smul} closes the goal.

Inverts the defining relation $\rho_\beta = (1/Z) \cdot e^{-\beta H}$
into the canonical "rescaled" form, useful when factoring out the
partition function in expectation computations.
\end{proof}

\begin{lemma}[\texttt{partitionFn\_mul\_gibbsExpectation\_eq}]
If $Z(\beta, H) \neq 0$, then for any observable $A$,
$Z(\beta, H) \cdot \langle A \rangle_\beta
 = \operatorname{Tr}(e^{-\beta H} \cdot A)$.
\end{lemma}

\begin{proof}[Proof sketch]
Push the scalar $Z$ inside the trace via
\texttt{Matrix.trace\_smul} (rewriting
$Z \cdot \operatorname{Tr}(\rho_\beta A)
 = \operatorname{Tr}(Z \cdot (\rho_\beta A))$),
then \texttt{Matrix.smul\_mul} repackages the scalar onto the first
factor: $\operatorname{Tr}((Z \cdot \rho_\beta) A)$.
The bracket is exactly $e^{-\beta H}$ by
\texttt{partitionFn\_smul\_gibbsState\_eq\_gibbsExp}, giving
$\operatorname{Tr}(e^{-\beta H} A)$.

The "unnormalised" expectation $\operatorname{Tr}(e^{-\beta H} A)$ is
the form that appears throughout canonical-ensemble derivations
(free-energy expansions, Trotter / path-integral arguments).
\end{proof}

\begin{lemma}[\texttt{gibbsExp\_natCast\_mul}]
For every $n \in \mathbb{N}$,
$e^{-(n\beta) H} = (e^{-\beta H})^n$.
\end{lemma}

\begin{proof}[Proof sketch]
Induction on $n$. Base case $n = 0$: $e^{0} = I = M^0$. Inductive
step: $((n+1)\beta) = n\beta + \beta$, hence by \texttt{gibbsExp\_add}
and the IH,
$e^{-((n+1)\beta) H} = e^{-(n\beta) H} \cdot e^{-\beta H}
                    = (e^{-\beta H})^n \cdot e^{-\beta H}
                    = (e^{-\beta H})^{n+1}$
via \texttt{pow\_succ}.
\end{proof}

\begin{lemma}[\texttt{gibbsExp\_two\_mul}]
$e^{-(2\beta) H} = e^{-\beta H} \cdot e^{-\beta H}$.
\end{lemma}

\begin{proof}[Proof sketch]
Specialise the previous lemma at $n = 2$ and unfold $M^2 = M \cdot M$
via \texttt{pow\_two}.
\end{proof}

\begin{lemma}[\texttt{gibbsExp\_inv}]
$(e^{-\beta H})^{-1} = e^{\beta H}$, i.e.\ the matrix inverse of the
Gibbs exponential at $\beta$ is the Gibbs exponential at $-\beta$.
\end{lemma}

\begin{proof}[Proof sketch]
By \texttt{gibbsExp\_neg\_mul\_self},
$e^{\beta H} \cdot e^{-\beta H} = I$.  Apply
\texttt{Matrix.inv\_eq\_left\_inv}, which says $A^{-1} = B$ whenever
$B A = I$.
\end{proof}

\begin{lemma}[\texttt{gibbsExp\_intCast\_mul}]
For every $n \in \mathbb{Z}$,
$e^{-(n\beta) H} = (e^{-\beta H})^n$.
\end{lemma}

\begin{proof}[Proof sketch]
Case split on $n = m$ or $n = -m$ for $m \in \mathbb{N}$ via
\texttt{Int.eq\_nat\_or\_neg}. The non-negative branch reduces to
\texttt{gibbsExp\_natCast\_mul} after \texttt{zpow\_natCast}. The
negative branch uses \texttt{Matrix.zpow\_neg} (which requires
$e^{-\beta H}$ to be a unit, witnessed via
\texttt{gibbsExp\_isUnit} composed with
\texttt{Matrix.isUnit\_iff\_isUnit\_det}), and then
\texttt{gibbsExp\_inv} to flip the sign of $\beta$ inside the
$m$-th power.
\end{proof}

\begin{lemma}[\texttt{partitionFn\_zero}]
At $\beta = 0$,
\[
  Z(0, H) \;=\; \operatorname{Tr}(\mathbf{1}) \;=\; \dim(\mathcal{H}),
\]
i.e.\ \texttt{partitionFn 0 H = Fintype.card ($\Lambda$ → Fin 2)},
where $\Lambda \to \mathrm{Fin}\,2$ is the configuration space.
\end{lemma}

\begin{proof}[Proof sketch]
Combines \texttt{gibbsExp\_zero} with \texttt{Matrix.trace\_one}, which
evaluates $\operatorname{Tr}(\mathbf{1})$ as the cardinality of the
finite index type.
\end{proof}

\begin{lemma}[\texttt{partitionFn\_zero\_ne\_zero}]
$Z(0, H) \neq 0$.
\end{lemma}

\begin{proof}[Proof sketch]
By \texttt{partitionFn\_zero}, $Z(0,H) = \dim(\mathcal{H})$.  Since
$\mathcal{H}$ is a non-empty finite type, its cardinality is a positive
natural number, which is nonzero as a complex number.  The cast
$\mathbb{N} \to \mathbb{C}$ preserves the non-zero property via
\texttt{Nat.cast\_ne\_zero}.
\end{proof}

\begin{remark}[Why $\beta = 0$ suffices]
The general case $Z(\beta, H) > 0$ for all $\beta$ follows from
spectral theory ($Z = \sum_i e^{-\beta\lambda_i}$ with real eigenvalues),
but the CFC instance chain required for that argument causes a
deterministic timeout in Lean's typeclass resolution.  The $\beta = 0$
case proved here is a concrete, sorry-free witness that the partition
function can be shown nonzero by elementary means, avoiding the CFC
timeout entirely.
\end{remark}

\begin{lemma}[\texttt{Matrix.IsHermitian.trace\_im}]
For any Hermitian matrix $A \in M_n(\mathbb{C})$,
$\Im\operatorname{Tr}(A) = 0$.
\end{lemma}

\begin{proof}[Proof sketch]
$\operatorname{Tr}(A^\dagger) = \overline{\operatorname{Tr}(A)}$
(\texttt{Matrix.trace\_conjTranspose}), and $A^\dagger = A$ by
Hermiticity, so $\operatorname{Tr}(A) = \overline{\operatorname{Tr}(A)}$.
Then \texttt{Complex.conj\_eq\_iff\_im} gives the imaginary-part
vanishing.
\end{proof}

\begin{lemma}[\texttt{partitionFn\_im\_of\_isHermitian}]
For a Hermitian Hamiltonian $H$,
$\Im Z(\beta, H) = 0$, i.e.\ the partition function is real.
\end{lemma}

\begin{proof}[Proof sketch]
$Z(\beta, H) = \operatorname{Tr}(e^{-\beta H})$.  By
\texttt{gibbsExp\_isHermitian}, $e^{-\beta H}$ is Hermitian; the
previous lemma then gives $\Im Z = 0$.
\end{proof}

\begin{lemma}[\texttt{gibbsState\_pow\_trace} (Rényi-$n$ precursor)]
For any $n \in \mathbb{N}$,
$\operatorname{Tr}(\rho_\beta^n)
 = Z(n\beta, H) / Z(\beta, H)^n$.
\end{lemma}

\begin{proof}[Proof sketch]
Unfold $\rho_\beta = (1/Z) \cdot e^{-\beta H}$ and apply
\texttt{smul\_pow} to get
$\rho_\beta^n = (1/Z)^n \cdot (e^{-\beta H})^n$.  Replace
$(e^{-\beta H})^n$ by $e^{-(n\beta) H}$ using
\texttt{gibbsExp\_natCast\_mul}.  Pull the scalar out of the trace
via \texttt{Matrix.trace\_smul} and rewrite
$(1/Z)^n = (Z^n)^{-1}$ via \texttt{one\_div} / \texttt{inv\_pow},
then \texttt{div\_eq\_inv\_mul} matches the RHS without needing a
$Z \neq 0$ hypothesis (Lean's $1/0 = 0$ convention handles the
edge case).

The Rényi-$n$ entropy precursor; the n-th moment of the density
matrix expressed entirely in terms of partition functions evaluated
at temperature $n \beta$, the canonical form behind the replica trick
and entanglement-entropy computations.
\end{proof}

\begin{lemma}[\texttt{gibbsState\_mul\_self\_trace} (purity)]
$\operatorname{Tr}(\rho_\beta^2)
 = Z(2\beta, H) / Z(\beta, H)^2$. (Special case $n = 2$ of the
preceding lemma, with a direct \texttt{gibbsExp\_two\_mul} proof.)
\end{lemma}

\begin{proof}[Proof sketch]
Unfold $\rho_\beta = (1/Z) \cdot e^{-\beta H}$:
$\rho_\beta \cdot \rho_\beta
 = (1/Z) \cdot \bigl(e^{-\beta H} \cdot ((1/Z) \cdot e^{-\beta H})\bigr)
 = (1/Z)^2 \cdot (e^{-\beta H} \cdot e^{-\beta H})
 = (1/Z)^2 \cdot e^{-(2\beta) H}$,
the last step by \texttt{gibbsExp\_two\_mul}.  Taking the trace and
using \texttt{Matrix.trace\_smul} gives $(1/Z)^2 \cdot Z(2\beta, H)$,
which \texttt{field\_simp} converts to $Z(2\beta, H) / Z^2$.
The $Z = 0$ edge case is handled by Lean's $1/0 = 0$ convention:
both sides collapse to zero, so no nonvanishing-$Z$ hypothesis is
needed.

The "purity" measure of the canonical density matrix; equal to $1$
for pure states, less than $1$ for mixed states.  Expressed entirely
in terms of partition functions evaluated at two temperatures, the
canonical form used in Rényi-2 entropy and replica-trick computations.
\end{proof}

\begin{lemma}[\texttt{gibbsState\_zero}]
At infinite temperature ($\beta = 0$),
$\rho_0 = \dim(\mathcal{H})^{-1} \cdot I$, the maximally mixed state.
\end{lemma}

\begin{proof}[Proof sketch]
By \texttt{gibbsExp\_zero}, $e^{-0\cdot H} = I$.  By
\texttt{partitionFn\_zero}, $Z(0, H) = \dim(\mathcal{H})$.  Substituting
into the definition $\rho_\beta = (1/Z)\,e^{-\beta H}$ gives
$\rho_0 = \dim(\mathcal{H})^{-1} \cdot I$.
\end{proof}

\begin{lemma}[\texttt{gibbsState\_zero\_inv}]
At $\beta = 0$, the matrix inverse of the Gibbs state is
$\rho_0^{-1} = \dim(\mathcal{H}) \cdot I$.
\end{lemma}

\begin{proof}[Proof sketch]
Substitute $\rho_0 = \dim^{-1} \cdot I$ via \texttt{gibbsState\_zero}.
By \texttt{Matrix.inv\_eq\_left\_inv}, it suffices to verify the
left-inverse identity
$(\dim \cdot I)\,(\dim^{-1} \cdot I) = I$, which
\texttt{smul\_mul\_smul\_comm} reduces to
$(\dim \cdot \dim^{-1}) \cdot (I \cdot I) = 1 \cdot I = I$, then
\texttt{inv\_mul\_cancel\_zero} closes the goal using
$\dim \neq 0$ (as a positive natural number cast to $\mathbb{C}$).
\end{proof}

\begin{lemma}[\texttt{gibbsExpectation\_zero}]
At $\beta = 0$, the expectation of any observable $A$ reduces to the
normalised trace
$\langle A \rangle_0 = \dim(\mathcal{H})^{-1} \cdot \operatorname{Tr}(A)$.
\end{lemma}

\begin{proof}[Proof sketch]
$\langle A \rangle_0 = \operatorname{Tr}(\rho_0 \cdot A)
= \operatorname{Tr}(\dim^{-1}\,I \cdot A)
= \dim^{-1} \cdot \operatorname{Tr}(A)$,
using \texttt{gibbsState\_zero}, \texttt{Matrix.smul\_mul},
\texttt{Matrix.one\_mul}, and \texttt{Matrix.trace\_smul}.
\end{proof}

\begin{lemma}[\texttt{gibbsExpectation\_add}]
For observables $O_1, O_2 \in M_n(\mathbb{C})$,
\[
  \langle O_1 + O_2 \rangle_\beta
  \;=\;
  \langle O_1 \rangle_\beta + \langle O_2 \rangle_\beta.
\]
\end{lemma}

\begin{proof}[Proof sketch]
Follows from linearity of $\operatorname{Tr}$ and distributivity of
matrix multiplication: $\rho(O_1 + O_2) = \rho O_1 + \rho O_2$, so
$\operatorname{Tr}(\rho(O_1+O_2))
 = \operatorname{Tr}(\rho O_1) + \operatorname{Tr}(\rho O_2)$.
\end{proof}

\begin{lemma}[\texttt{gibbsExpectation\_smul}]
For $c \in \mathbb{C}$ and $O \in M_n(\mathbb{C})$,
\[
  \langle c \cdot O \rangle_\beta
  \;=\;
  c \cdot \langle O \rangle_\beta.
\]
\end{lemma}

\begin{proof}[Proof sketch]
$\rho \cdot (c \cdot O) = c \cdot (\rho \cdot O)$ by scalar associativity,
so $\operatorname{Tr}(\rho \cdot (c \cdot O))
   = c \cdot \operatorname{Tr}(\rho \cdot O)$
by linearity of the trace.
\end{proof}

\begin{lemma}[\texttt{gibbsExpectation\_neg} / \texttt{gibbsExpectation\_sub}]
$\langle -O \rangle_\beta = -\langle O \rangle_\beta$
and $\langle A - B \rangle_\beta
     = \langle A \rangle_\beta - \langle B \rangle_\beta$.
\end{lemma}

\begin{proof}[Proof sketch]
Negation: $\rho \cdot (-O) = -(\rho \cdot O)$ via
\texttt{Matrix.mul\_neg}, then \texttt{Matrix.trace\_neg} commutes the
sign out of the trace.  Subtraction: rewrite $A - B$ as $A + (-B)$ and
combine \texttt{gibbsExpectation\_add} with \texttt{gibbsExpectation\_neg}.
\end{proof}

\begin{lemma}[\texttt{gibbsExpectation\_sum}]
For any \texttt{Finset s} and family $f : \iota \to M_n(\mathbb{C})$,
\[
  \Bigl\langle \sum_{i \in s} f(i) \Bigr\rangle_\beta
  \;=\; \sum_{i \in s} \langle f(i) \rangle_\beta.
\]
\end{lemma}

\begin{proof}[Proof sketch]
Induction on the finset.  Base case $s = \emptyset$:
$\sum_{i \in \emptyset} f(i) = 0$, and
$\langle 0 \rangle_\beta
 = \operatorname{Tr}(\rho \cdot 0) = \operatorname{Tr}(0) = 0
 = \sum_{i \in \emptyset} \langle f(i) \rangle_\beta$.
Inductive step: \texttt{Finset.sum\_insert} splits off the new element,
\texttt{gibbsExpectation\_add} pushes the expectation through the sum,
and the IH closes the rest.
\end{proof}

Reference: Tasaki~\cite{Tasaki:QMB}, §3.3, pp.\,75--78.

\subsection{Commutativity with the Hamiltonian and reality of expectation}
\label{subsec:gibbs-commute-real}

\emph{Reference:} Tasaki~\cite{Tasaki:QMB}, §3.3, p.~80.

\begin{lemma}[\texttt{gibbsExp\_commute\_hamiltonian}]
$[e^{-\beta H}, H] = 0$.
\end{lemma}

\begin{proof}[Proof sketch]
$H$ commutes with itself, hence the scalar multiple $-\beta\,H$ commutes
with $H$, and therefore so does the matrix exponential $e^{-\beta H}$
(\texttt{Commute.exp\_left} in Mathlib's
\texttt{NormedSpace.exp} formalism).
\end{proof}

\begin{lemma}[\texttt{gibbsState\_commute\_hamiltonian}]
$[\rho_\beta, H] = 0$, i.e.\ the Gibbs state is stationary under the
unitary dynamics generated by $H$.
\end{lemma}

\begin{proof}[Proof sketch]
$\rho_\beta = (1/Z) \cdot e^{-\beta H}$ is a scalar multiple of
$e^{-\beta H}$, so commutativity is inherited from
\texttt{gibbsExp\_commute\_hamiltonian} via \texttt{Commute.smul\_left}.
\end{proof}

\begin{lemma}[\texttt{Matrix.trace\_mul\_star\_of\_isHermitian}]
For Hermitian matrices $A, B \in M_n(\mathbb{C})$,
$\overline{\operatorname{Tr}(A B)} = \operatorname{Tr}(A B)$.
This is the Gibbs-independent algebraic core used by the next two
lemmas.
\end{lemma}

\begin{proof}[Proof sketch]
By \texttt{Matrix.trace\_conjTranspose} and
\texttt{Matrix.conjTranspose\_mul}, the conjugate of the trace equals
$\operatorname{Tr}(B^\dagger A^\dagger)$. Substituting Hermiticity
$A^\dagger = A$, $B^\dagger = B$ gives $\operatorname{Tr}(B A)$, which
equals $\operatorname{Tr}(A B)$ by the cyclic property
\texttt{Matrix.trace\_mul\_comm}.
\end{proof}

\begin{lemma}[\texttt{gibbsExpectation\_star\_of\_isHermitian}]
For Hermitian $H$ and Hermitian observable $O \in M_n(\mathbb{C})$,
$\overline{\langle O \rangle_\beta} = \langle O \rangle_\beta$.
\end{lemma}

\begin{proof}[Proof sketch]
$\rho_\beta$ is Hermitian by \texttt{gibbsState\_isHermitian}, so this
is an immediate corollary of
\texttt{Matrix.trace\_mul\_star\_of\_isHermitian} applied to
$A = \rho_\beta$ and $B = O$.
\end{proof}

\begin{lemma}[\texttt{gibbsExpectation\_im\_of\_isHermitian}]
For Hermitian $H$ and Hermitian observable $O$,
$\Im\langle O \rangle_\beta = 0$.
\end{lemma}

\begin{proof}[Proof sketch]
Immediate corollary of
\texttt{gibbsExpectation\_star\_of\_isHermitian} via
\texttt{Complex.conj\_eq\_iff\_im}: a complex number is fixed by
conjugation iff its imaginary part vanishes.
\end{proof}

\begin{lemma}[\texttt{gibbsExpectation\_ofReal\_re\_eq\_of\_isHermitian}]
For Hermitian $H$ and Hermitian observable $O$,
$((\langle O \rangle_\beta).\Re : \mathbb{C}) = \langle O \rangle_\beta$.
\end{lemma}

\begin{proof}[Proof sketch]
Constructive complement to the previous lemma.  By
\texttt{Complex.ext}, equality of two complex numbers reduces to equality
of real and imaginary parts.  The real part is trivially equal; the
imaginary parts agree because $(\Re z, 0)$ has zero imaginary part and
$\Im\langle O \rangle_\beta = 0$ by
\texttt{gibbsExpectation\_im\_of\_isHermitian}.
\end{proof}

\subsection{Conservation laws}
\label{subsec:gibbs-conservation}

\emph{Reference:} Tasaki~\cite{Tasaki:QMB}, §3.3, p.~80.

\begin{lemma}[\texttt{gibbsExpectation\_mul\_hamiltonian\_comm}]
For any observable $A \in M_n(\mathbb{C})$,
$\langle H A \rangle_\beta = \langle A H \rangle_\beta$.
\end{lemma}

\begin{proof}[Proof sketch]
Inside the trace,
\[
  \operatorname{Tr}(\rho_\beta H A)
  = \operatorname{Tr}(H \rho_\beta A)
  = \operatorname{Tr}(\rho_\beta A H),
\]
the first equality by
\texttt{gibbsState\_commute\_hamiltonian} (i.e.\ $\rho_\beta H = H \rho_\beta$)
and the second by the cyclic property
\texttt{Matrix.trace\_mul\_comm}, applied as
$\operatorname{Tr}(H (\rho_\beta A)) = \operatorname{Tr}((\rho_\beta A) H)$.
\end{proof}

\begin{lemma}[\texttt{gibbsExpectation\_mul\_comm\_of\_commute\_hamiltonian}]
For any conserved quantity $A$ (i.e.\ $[A, H] = 0$) and any
observable $O$,
$\langle A O \rangle_\beta = \langle O A \rangle_\beta$.
\end{lemma}

\begin{proof}[Proof sketch]
Generalisation of the previous lemma from $A = H$ to any conserved
quantity.  The key observation is that $A$ commutes with the Gibbs
state: $\rho_\beta = (1/Z) \cdot e^{-\beta H}$ is a power-series
function of $H$, so
\texttt{Commute.exp\_right} (with the scalar
$-\beta$ inserted via \texttt{Commute.smul\_right}) gives
$[A, e^{-\beta H}] = 0$, and \texttt{Commute.smul\_right}
extends to $[A, \rho_\beta] = 0$.  Then
$\operatorname{Tr}(\rho A O) = \operatorname{Tr}(A \rho O) =
 \operatorname{Tr}(\rho O A)$,
the latter by \texttt{Matrix.trace\_mul\_comm}.

This is the algebraic statement that conserved quantities can be
"factored through" the Gibbs expectation, the basis for selection
rules and conservation-induced simplifications.
\end{proof}

\begin{lemma}[\texttt{gibbsExpectation\_commutator\_eq\_zero\_of\_commute\_hamiltonian}]
For any conserved quantity $A$ and any observable $O$,
$\langle A O - O A \rangle_\beta = 0$.
\end{lemma}

\begin{proof}[Proof sketch]
Linearity (\texttt{gibbsExpectation\_sub}) plus the previous lemma:
$\langle A O - O A \rangle_\beta = \langle A O \rangle_\beta - \langle O A \rangle_\beta = 0$.

The standard "selection rule" form of Noether's theorem at finite
volume: when $A$ is a conserved quantity, the Heisenberg equation
$\dot A = i[H, A] = 0$ implies that the Gibbs ensemble cannot
distinguish $A O$ from $O A$.
\end{proof}

\begin{lemma}[\texttt{gibbsExpectation\_commutator\_hamiltonian}]
For any observable $A$,
$\langle [H, A] \rangle_\beta = 0$.
\end{lemma}

\begin{proof}[Proof sketch]
Expanding $[H, A] = H A - A H$ and using
\texttt{gibbsExpectation\_add}, \texttt{gibbsExpectation\_smul} for
linearity, the result follows immediately from
\texttt{gibbsExpectation\_mul\_hamiltonian\_comm}.
\end{proof}

\begin{lemma}[\texttt{gibbsExpectation\_hamiltonian\_im}]
For a Hermitian Hamiltonian $H$,
$\Im\langle H \rangle_\beta = 0$, i.e.\ the energy expectation
$\langle H \rangle_\beta$ is real.
\end{lemma}

\begin{proof}[Proof sketch]
Specialise \texttt{gibbsExpectation\_im\_of\_isHermitian} to $O = H$.
\end{proof}

\begin{lemma}[\texttt{gibbsExpectation\_sq\_im\_of\_isHermitian}]
For Hermitian $H$ and Hermitian observable $O$,
$\Im\langle O \cdot O \rangle_\beta = 0$.
\end{lemma}

\begin{proof}[Proof sketch]
$O \cdot O$ is Hermitian: it is the product of two commuting
Hermitian operators (\texttt{Matrix.IsHermitian.mul\_of\_commute}
applied with \texttt{Commute.refl O}). The result then specialises
\texttt{gibbsExpectation\_im\_of\_isHermitian} at observable
$O \cdot O$.

This is the variance precursor: the second moment
$\langle O^2 \rangle_\beta$ being real lets the variance
$\operatorname{Var}_\beta(O) = \langle O^2 \rangle_\beta - \langle O \rangle_\beta^2$
inherit its standard interpretation as a non-negative real number
(positivity itself requires the spectral / CFC infrastructure not
yet developed).
\end{proof}

\begin{lemma}[\texttt{gibbsExpectation\_pow\_im\_of\_isHermitian}]
For Hermitian $H$, Hermitian observable $O$, and any $n \in \mathbb{N}$,
$\Im\langle O^n \rangle_\beta = 0$.
\end{lemma}

\begin{proof}[Proof sketch]
$O^n$ is Hermitian by \texttt{Matrix.IsHermitian.pow}, the
specialisation of \texttt{IsSelfAdjoint.pow}. The result then
specialises \texttt{gibbsExpectation\_im\_of\_isHermitian}
at observable $O^n$. This generalises the second-moment realness to
all natural-power moments, the algebraic foundation of the canonical
moment-generating function and cumulant expansion.
\end{proof}

\subsection{Variance}
\label{subsec:gibbs-variance}

\emph{Reference:} Tasaki~\cite{Tasaki:QMB}, §3.3 (canonical-ensemble
fluctuations).

\begin{definition}[\texttt{gibbsVariance}]
The canonical-ensemble \emph{variance} of an observable $O$ at
inverse temperature $\beta$ under Hamiltonian $H$ is
\[
  \operatorname{Var}_\beta(O)
  \;:=\;
  \langle O \cdot O \rangle_\beta - \langle O \rangle_\beta^2.
\]
\end{definition}

\begin{lemma}[\texttt{gibbsVariance\_im\_of\_isHermitian}]
For Hermitian $H$ and Hermitian observable $O$,
$\Im\operatorname{Var}_\beta(O) = 0$.
\end{lemma}

\begin{proof}[Proof sketch]
Both $\langle O \cdot O \rangle_\beta$ (by
\texttt{gibbsExpectation\_sq\_im\_of\_isHermitian}) and
$\langle O \rangle_\beta$ (by
\texttt{gibbsExpectation\_im\_of\_isHermitian}) are real. The square
of a real complex number $z$ has imaginary part $2 \cdot \Re z \cdot \Im z = 0$
(here computed by \texttt{Complex.mul\_im} after the \texttt{pow\_two}
unfold). The difference of two reals is real
(\texttt{Complex.sub\_im}).

Non-negativity $\operatorname{Var}_\beta(O) \geq 0$ requires the
spectral / CFC infrastructure (positivity of $\rho_\beta$) and is
deferred.
\end{proof}

\begin{lemma}[\texttt{gibbsVariance\_zero}]
At infinite temperature ($\beta = 0$),
\[
  \operatorname{Var}_0(O)
  \;=\;
  \dim^{-1} \operatorname{Tr}(O^2)
  - \bigl(\dim^{-1} \operatorname{Tr}(O)\bigr)^2.
\]
\end{lemma}

\begin{proof}[Proof sketch]
Substitute the $\beta = 0$ closed form
$\langle A \rangle_0 = \dim^{-1} \cdot \operatorname{Tr}(A)$
(\texttt{gibbsExpectation\_zero}) into the variance definition for both
the second-moment and first-moment occurrences.
\end{proof}

\begin{lemma}[\texttt{gibbsVariance\_eq\_centered\_sq}]
For $Z(\beta, H) \neq 0$,
\[
  \operatorname{Var}_\beta(O)
  \;=\;
  \langle (O - \langle O \rangle_\beta \cdot I) \cdot (O - \langle O \rangle_\beta \cdot I) \rangle_\beta.
\]
\end{lemma}

\begin{proof}[Proof sketch]
Expand the matrix square:
$(O - c I)(O - c I) = O^2 - c \cdot O - c \cdot O + c^2 \cdot I$
where $c = \langle O \rangle_\beta$.  Apply linearity of the
expectation (\texttt{gibbsExpectation\_add}, \texttt{gibbsExpectation\_sub},
\texttt{gibbsExpectation\_smul}) and use
$\langle I \rangle_\beta = 1$ (\texttt{gibbsExpectation\_one},
requiring $Z \neq 0$) to obtain
$\langle O^2 \rangle_\beta - 2 c^2 + c^2 = \langle O^2 \rangle_\beta - c^2
 = \operatorname{Var}_\beta(O)$.

Identifies the canonical-ensemble variance with the second moment of
the centered observable $\hat O = O - \langle O \rangle_\beta \cdot I$,
the standard form for concentration / cumulant arguments.
\end{proof}

\subsection{Covariance}
\label{subsec:gibbs-covariance}

\begin{definition}[\texttt{gibbsCovariance}]
The canonical-ensemble (complex) covariance of two observables
$A, B$ at inverse temperature $\beta$ under Hamiltonian $H$ is
\[
  \operatorname{Cov}_\beta(A, B)
  \;:=\;
  \langle A \cdot B \rangle_\beta
  - \langle A \rangle_\beta \cdot \langle B \rangle_\beta.
\]
\end{definition}

\begin{lemma}[\texttt{gibbsCovariance\_self\_eq\_variance}]
$\operatorname{Cov}_\beta(O, O) = \operatorname{Var}_\beta(O)$.
\end{lemma}

\begin{proof}[Proof sketch]
Unfold both sides; the variance uses $(\langle O \rangle_\beta)^2$,
which equals $\langle O \rangle_\beta \cdot \langle O \rangle_\beta$
by \texttt{pow\_two}.
\end{proof}

\begin{remark}
For non-commuting Hermitian $A, B$, $\operatorname{Cov}_\beta(A, B)$
is generally complex: its imaginary part captures
$(1/2) \Im\langle [A, B] \rangle_\beta$.  The symmetric / real-valued
covariance is treated below.
\end{remark}

\begin{lemma}[\texttt{gibbsCovariance\_sub\_swap\_eq\_commutator}]
$\operatorname{Cov}_\beta(A, B) - \operatorname{Cov}_\beta(B, A)
 = \langle A B - B A \rangle_\beta$.
\end{lemma}

\begin{proof}[Proof sketch]
The product terms $\langle A \rangle_\beta \cdot \langle B \rangle_\beta$
cancel under subtraction.  The remaining
$\langle A B \rangle_\beta - \langle B A \rangle_\beta$ is exactly
$\langle A B - B A \rangle_\beta$ by linearity
(\texttt{gibbsExpectation\_sub}).
\end{proof}

\begin{lemma}[Bilinearity of $\operatorname{Cov}_\beta$]
The complex covariance is bilinear in its two observable arguments:
\begin{align*}
\operatorname{Cov}_\beta(A_1 + A_2, B) &= \operatorname{Cov}_\beta(A_1, B) + \operatorname{Cov}_\beta(A_2, B), \\
\operatorname{Cov}_\beta(A, B_1 + B_2) &= \operatorname{Cov}_\beta(A, B_1) + \operatorname{Cov}_\beta(A, B_2), \\
\operatorname{Cov}_\beta(c \cdot A, B) &= c \cdot \operatorname{Cov}_\beta(A, B), \\
\operatorname{Cov}_\beta(A, c \cdot B) &= c \cdot \operatorname{Cov}_\beta(A, B).
\end{align*}
Lean names: \texttt{gibbsCovariance\_add\_left},
\texttt{gibbsCovariance\_add\_right},
\texttt{gibbsCovariance\_smul\_left},
\texttt{gibbsCovariance\_smul\_right}.
\end{lemma}

\begin{proof}[Proof sketch]
Each follows from matrix-level distributivity (\texttt{add\_mul},
\texttt{mul\_add}, \texttt{Matrix.smul\_mul}, \texttt{Matrix.mul\_smul})
on the inner $A \cdot B$ factor, plus linearity of
\texttt{gibbsExpectation} (additivity \texttt{gibbsExpectation\_add},
scalar pullout \texttt{gibbsExpectation\_smul}), and a
\texttt{ring} closing of the resulting complex-arithmetic equality.
\end{proof}

\subsection{Symmetric covariance}
\label{subsec:gibbs-covariance-symm}

\begin{definition}[\texttt{gibbsCovarianceSymm}]
The symmetric (real-valued for Hermitian observables)
canonical-ensemble covariance is
\[
  \operatorname{Cov}^{\mathrm{s}}_\beta(A, B)
  \;:=\;
  \frac{1}{2} \langle A B + B A \rangle_\beta
  - \langle A \rangle_\beta \cdot \langle B \rangle_\beta.
\]
\end{definition}

\begin{lemma}[\texttt{gibbsCovarianceSymm\_self\_eq\_variance}]
$\operatorname{Cov}^{\mathrm{s}}_\beta(O, O) = \operatorname{Var}_\beta(O)$.
\end{lemma}

\begin{proof}[Proof sketch]
The anticommutator $O \cdot O + O \cdot O = 2 (O \cdot O)$ collapses
the $1/2$ prefactor: $\tfrac{1}{2} \langle 2 (O \cdot O) \rangle_\beta
= \langle O \cdot O \rangle_\beta$.  Combined with
$(\langle O \rangle_\beta)^2 = \langle O \rangle_\beta \cdot \langle O \rangle_\beta$,
the LHS equals the variance.
\end{proof}

\begin{lemma}[\texttt{gibbsCovarianceSymm\_im\_of\_isHermitian}]
For Hermitian $H$ and Hermitian $A, B$,
$\Im\operatorname{Cov}^{\mathrm{s}}_\beta(A, B) = 0$.
\end{lemma}

\begin{proof}[Proof sketch]
The anticommutator expectation
$\langle A B + B A \rangle_\beta$ is real by
\texttt{gibbsExpectation\_anticommutator\_im}, and the $1/2$ prefactor
preserves realness.  Each first moment $\langle A \rangle_\beta$,
$\langle B \rangle_\beta$ is real by
\texttt{gibbsExpectation\_im\_of\_isHermitian}, so the product is
real.  The difference of two reals is real.
\end{proof}

\begin{lemma}[\texttt{gibbsCovarianceSymm\_comm}]
$\operatorname{Cov}^{\mathrm{s}}_\beta(A, B) = \operatorname{Cov}^{\mathrm{s}}_\beta(B, A)$.
\end{lemma}

\begin{proof}[Proof sketch]
The anticommutator $A B + B A = B A + A B$ is unchanged under swap
(matrix-level \texttt{add\_comm}), and the product
$\langle A \rangle_\beta \cdot \langle B \rangle_\beta
 = \langle B \rangle_\beta \cdot \langle A \rangle_\beta$
is commutative in $\mathbb{C}$.
\end{proof}

\begin{lemma}[Bilinearity of $\operatorname{Cov}^{\mathrm{s}}_\beta$]
The symmetric covariance is bilinear in its observable arguments
(\texttt{add\_left}, \texttt{add\_right}, \texttt{smul\_left},
\texttt{smul\_right}).  Combined with
\texttt{gibbsCovarianceSymm\_comm} and the realness for Hermitian
observables, this makes $\operatorname{Cov}^{\mathrm{s}}_\beta$ a
genuine real-valued symmetric bilinear pairing on Hermitian
observables.
\end{lemma}

\begin{proof}[Proof sketch]
The left-additivity is direct: split
$(A_1 + A_2) B + B (A_1 + A_2) = (A_1 B + B A_1) + (A_2 B + B A_2)$
via \texttt{add\_mul} / \texttt{mul\_add} / \texttt{abel}, then apply
\texttt{gibbsExpectation\_add} on the resulting sums and close with
\texttt{ring}.  Right-additivity follows from left-additivity and
\texttt{gibbsCovarianceSymm\_comm}.  Scalar pull-out uses
\texttt{Matrix.smul\_mul} / \texttt{Matrix.mul\_smul} on the
anticommutator and \texttt{gibbsExpectation\_smul}.
\end{proof}

\begin{lemma}[Vanishing on constant-scalar arguments]
For $c \in \mathbb{C}$ and $Z(\beta, H) \neq 0$,
\[
  \operatorname{Cov}_\beta(c \cdot I, B)
  = \operatorname{Cov}_\beta(A, c \cdot I)
  = 0,
\]
and analogously for $\operatorname{Cov}^{\mathrm{s}}_\beta$.
Lean names: \texttt{gibbsCovariance\_const\_smul\_one\_left\_eq\_zero},
\texttt{gibbsCovariance\_const\_smul\_one\_right\_eq\_zero},
\texttt{gibbsCovarianceSymm\_const\_smul\_one\_left\_eq\_zero},
\texttt{gibbsCovarianceSymm\_const\_smul\_one\_right\_eq\_zero}.
\end{lemma}

\begin{proof}[Proof sketch]
Constant scalars are deterministic observables.  Concretely
$\langle (c \cdot I) B \rangle_\beta = c \langle B \rangle_\beta$
(via \texttt{Matrix.smul\_mul} / \texttt{Matrix.one\_mul} and
\texttt{gibbsExpectation\_smul}), and
$\langle c \cdot I \rangle_\beta = c \cdot \langle I \rangle_\beta = c$
(using \texttt{gibbsExpectation\_one} which requires $Z \neq 0$).
The two terms cancel.  The $\operatorname{Cov}^{\mathrm{s}}$ versions
follow from the symmetrisation identity
\texttt{gibbsCovarianceSymm\_eq\_half\_add\_swap}.
\end{proof}

\begin{lemma}[\texttt{gibbsVariance\_add}]
$\operatorname{Var}_\beta(A + B)
 = \operatorname{Var}_\beta(A) + \operatorname{Var}_\beta(B)
   + 2 \operatorname{Cov}^{\mathrm{s}}_\beta(A, B)$.
\end{lemma}

\begin{proof}[Proof sketch]
Expand
$(A + B)(A + B) = A^2 + A B + B A + B^2$
via \texttt{add\_mul} / \texttt{mul\_add} / \texttt{abel}, push the
expectation through with \texttt{simp [gibbsExpectation\_add]}, and
close with \texttt{ring}.  The cross term collects into
$\langle A B + B A \rangle_\beta$, which equals
$2 \operatorname{Cov}^{\mathrm{s}}_\beta(A, B) + 2 \langle A \rangle_\beta \langle B \rangle_\beta$.
\end{proof}

\begin{lemma}[Variance under transformations]
\begin{itemize}
  \item $\operatorname{Var}_\beta(c \cdot A) = c^2 \operatorname{Var}_\beta(A)$
    (\texttt{gibbsVariance\_smul}).
  \item $\operatorname{Var}_\beta(-A) = \operatorname{Var}_\beta(A)$
    (\texttt{gibbsVariance\_neg}).
  \item For $Z \neq 0$: $\operatorname{Var}_\beta(I) = 0$
    (\texttt{gibbsVariance\_one}),
    $\operatorname{Var}_\beta(c \cdot I) = 0$
    (\texttt{gibbsVariance\_smul\_one}), and
    $\operatorname{Var}_\beta(A + c \cdot I) = \operatorname{Var}_\beta(A)$
    (\texttt{gibbsVariance\_add\_const\_smul\_one}).
\end{itemize}
\end{lemma}

\begin{proof}[Proof sketch]
Variance scales as $c^2$ via \texttt{Matrix.smul\_mul} /
\texttt{Matrix.mul\_smul} / \texttt{smul\_smul} on the inner product
and \texttt{gibbsExpectation\_smul} on each occurrence; the
$\operatorname{Var}(-A)$ case is the $c = -1$ specialization.
$\operatorname{Var}(I) = \langle I \rangle - \langle I \rangle^2 = 1 - 1 = 0$
when $\langle I \rangle = 1$, requiring $Z \neq 0$.  The constant-shift
invariance follows by direct expansion of $(A + c \cdot I)(A + c \cdot I)$
together with $\langle I \rangle = 1$; the cross terms cancel.
\end{proof}

\begin{lemma}[Symmetric / antisymmetric decomposition]
$\operatorname{Cov}_\beta(A, B)
 = \operatorname{Cov}^{\mathrm{s}}_\beta(A, B)
 + \tfrac{1}{2} \langle A B - B A \rangle_\beta$,
and equivalently
$\operatorname{Cov}^{\mathrm{s}}_\beta(A, B)
 = \tfrac{1}{2}\,\bigl(\operatorname{Cov}_\beta(A, B)
                       + \operatorname{Cov}_\beta(B, A)\bigr)$.
Lean names: \texttt{gibbsCovariance\_eq\_symm\_add\_half\_commutator},
\texttt{gibbsCovarianceSymm\_eq\_half\_add\_swap}.
\end{lemma}

\begin{proof}[Proof sketch]
Both follow from unfolding both covariances and applying
\texttt{gibbsExpectation\_add} (on the anticommutator
$A B + B A$) or \texttt{gibbsExpectation\_sub} (on the commutator
$A B - B A$); a single \texttt{ring} step then closes the
complex-arithmetic equality.

This is the canonical decomposition into symmetric (anticommutator-based,
real for Hermitian observables) and antisymmetric (commutator-based,
purely imaginary for Hermitian observables) parts, the algebraic
backbone of Kubo / fluctuation-dissipation arguments.
\end{proof}

\begin{lemma}[\texttt{gibbsCovariance\_eq\_symm\_of\_commute}]
For commuting observables $A, B$ (i.e.\ $[A, B] = 0$),
$\operatorname{Cov}_\beta(A, B) = \operatorname{Cov}^{\mathrm{s}}_\beta(A, B)$.
\end{lemma}

\begin{proof}[Proof sketch]
Apply the symmetric / antisymmetric decomposition.  The
commutator-based correction
$\tfrac{1}{2} \langle A B - B A \rangle_\beta$ vanishes when
$A B - B A = 0$ (\texttt{Commute.eq}), and
$\langle 0 \rangle_\beta = 0$ via
\texttt{Matrix.mul\_zero} / \texttt{Matrix.trace\_zero}.
\end{proof}

\subsection{Anticommutator real, commutator purely imaginary}
\label{subsec:gibbs-anti-comm}

\emph{Reference:} Tasaki~\cite{Tasaki:QMB}, §3.3, p.~80.

\begin{lemma}[\texttt{Matrix.trace\_mul\_conjTranspose\_swap\_of\_isHermitian}]
For Hermitian $\rho$ and arbitrary $X \in M_n(\mathbb{C})$,
$\overline{\operatorname{Tr}(\rho X)} = \operatorname{Tr}(\rho X^\dagger)$.
\end{lemma}

\begin{proof}[Proof sketch]
$\overline{\operatorname{Tr}(\rho X)}
 = \operatorname{Tr}((\rho X)^\dagger)
 = \operatorname{Tr}(X^\dagger \rho^\dagger)
 = \operatorname{Tr}(X^\dagger \rho)
 = \operatorname{Tr}(\rho X^\dagger)$,
using \texttt{Matrix.trace\_conjTranspose},
\texttt{Matrix.conjTranspose\_mul}, $\rho^\dagger = \rho$, and the
cyclic property \texttt{Matrix.trace\_mul\_comm}.
\end{proof}

\begin{lemma}[\texttt{gibbsExpectation\_star\_swap\_of\_isHermitian}]
For Hermitian $H$ and Hermitian $A, B$,
$\overline{\langle A B \rangle_\beta} = \langle B A \rangle_\beta$.
\end{lemma}

\begin{proof}[Proof sketch]
Apply the previous lemma at $\rho = \rho_\beta$, $X = A B$, then use
$(A B)^\dagger = B^\dagger A^\dagger = B A$ by Hermiticity.
\end{proof}

\begin{lemma}[\texttt{gibbsExpectation\_anticommutator\_im}]
For Hermitian $H$ and Hermitian $A, B$,
$\Im\langle A B + B A \rangle_\beta = 0$.
\end{lemma}

\begin{proof}[Proof sketch]
By the swap lemma, $\langle A B + B A \rangle_\beta = z + \overline{z}$
where $z = \langle A B \rangle_\beta$. The imaginary part of $z + \bar z$
vanishes ($\Im z + (-\Im z) = 0$).
\end{proof}

\begin{lemma}[\texttt{gibbsExpectation\_commutator\_re}]
For Hermitian $H$ and Hermitian $A, B$,
$\Re\langle A B - B A \rangle_\beta = 0$.
\end{lemma}

\begin{proof}[Proof sketch]
By the swap lemma, $\langle A B - B A \rangle_\beta = z - \overline{z}
= 2 i\,\Im z$, which is purely imaginary. Equivalently
$\Re(z + (-1)\bar z) = \Re z + (-1)\,\Re\bar z = \Re z - \Re z = 0$.
\end{proof}

\begin{lemma}[\texttt{gibbsExpectation\_mul\_hamiltonian\_im}]
For Hermitian $H$ and Hermitian observable $O$,
$\Im\langle H O \rangle_\beta = 0$.
\end{lemma}

\begin{proof}[Proof sketch]
By the conservation law $\langle H O \rangle_\beta = \langle O H \rangle_\beta$
combined with the swap lemma at $A = H, B = O$
($\overline{\langle H O \rangle_\beta} = \langle O H \rangle_\beta = \langle H O \rangle_\beta$),
the expectation is fixed by complex conjugation, hence real.
\end{proof}

%======================================================================
\section{Heisenberg chain}
\label{sec:heisenberg-chain}
%======================================================================

\emph{Source file:} \texttt{LatticeSystem/Quantum/HeisenbergChain.lean}.

\emph{Primary reference:} Tasaki~\cite{Tasaki:QMB}, §3.5, p.~89.

\subsection{Coupling functions}

\begin{definition}[\texttt{openChainCoupling}]
For an $(N{+}1)$-site chain, the \emph{open-chain coupling} is the
function $J_{\mathrm{open}} : \mathrm{Fin}\,(N{+}1) \to
\mathrm{Fin}\,(N{+}1) \to \mathbb{C}$ defined by
\[
  J_{\mathrm{open}}(i,j) \;:=\;
  \begin{cases}
    -J & \text{if } i_{\mathrm{val}} + 1 = j_{\mathrm{val}}
        \text{ or } j_{\mathrm{val}} + 1 = i_{\mathrm{val}},\\
    0 & \text{otherwise,}
  \end{cases}
\]
where $J \in \mathbb{R}$ is a real coupling parameter.  The negative
sign encodes Tasaki's ferromagnetic convention $H = -J\sum_{\langle
x,y\rangle} \hat S_x\!\cdot\!\hat S_y$ (Tasaki eq.~(2.4.1)) so that
$J > 0$ is ferromagnetic.  The function is symmetric in $(i,j)$.
Lean name: \texttt{openChainCoupling~$N$~$J$}.
\end{definition}

\begin{definition}[\texttt{periodicChainCoupling}]
The \emph{periodic-chain coupling} on $\mathrm{Fin}\,(N{+}2)$ (at
least 2 sites) imposes nearest-neighbour bonds with periodic boundary:
\[
  J_{\mathrm{per}}(i,j) \;:=\;
  \begin{cases}
    -J & \text{if } i + 1 = j \text{ or } j + 1 = i \text{ in } \mathbb{Z}/(N{+}2),\\
    0 & \text{otherwise.}
  \end{cases}
\]
Same ferromagnetic sign convention.
Lean name: \texttt{periodicChainCoupling~$N$~$J$}.
\end{definition}

\subsection{Formalized facts}

\begin{theorem}[\texttt{heisenbergHamiltonian\_isHermitian\_of\_real\_symm}]
Let $\Lambda$ be a finite lattice and let
$J : \Lambda \to \Lambda \to \mathbb{R}$ be a real-valued symmetric
coupling ($J(x,y) = J(y,x)$ for all $x,y$).  Then the Heisenberg
Hamiltonian
\[
  H \;:=\; \sum_{x,y \in \Lambda} J(x,y)\,\hat S_x \cdot \hat S_y
\]
is Hermitian.
\end{theorem}

\begin{proof}[Proof sketch]
Each two-site term $\hat S_x \cdot \hat S_y$ is Hermitian
(proved in \texttt{spinHalfDot\_isHermitian}).  Scaling by the real
coefficient $J(x,y)$ preserves Hermiticity, and a finite sum of
Hermitian matrices is Hermitian.  The symmetry $J(x,y) = J(y,x)$
together with the commutativity $\hat S_x \cdot \hat S_y =
\hat S_y \cdot \hat S_x$ ensures that the double sum produces each bond
only once up to the real scalar.
\end{proof}

\begin{theorem}[\texttt{openChainHeisenberg\_isHermitian}]
The open-chain Heisenberg Hamiltonian
$H_{\mathrm{open}} := \sum_{i=0}^{N-2} \hat S_i \cdot \hat S_{i+1}$
is Hermitian.
\end{theorem}

\begin{proof}[Proof sketch]
Instantiate \texttt{heisenbergHamiltonian\_isHermitian\_of\_real\_symm}
with $J = J_{\mathrm{open}}$, verifying that $J_{\mathrm{open}}$ is
real and symmetric by construction.
\end{proof}

\begin{theorem}[\texttt{periodicChainHeisenberg\_isHermitian}]
The periodic-chain Heisenberg Hamiltonian
$H_{\mathrm{per}} := H_{\mathrm{open}} + \hat S_0 \cdot \hat S_{N-1}$
is Hermitian.
\end{theorem}

\begin{proof}[Proof sketch]
Instantiate \texttt{heisenbergHamiltonian\_isHermitian\_of\_real\_symm}
with $J = J_{\mathrm{per}}$.  Symmetry of $J_{\mathrm{per}}$ follows
from the explicit definition, which treats the boundary bond
$(0, N-1)$ and $(N-1, 0)$ symmetrically.
\end{proof}

\subsection{Gibbs state instantiation}
\label{subsec:heisenberg-gibbs}

\emph{Reference:} Tasaki~\cite{Tasaki:QMB}, §3.3, p.~80.

\begin{definition}[\texttt{openChainHeisenbergGibbsState}]
The Gibbs state of the open-boundary 1D Heisenberg chain on
$\mathrm{Fin}(N+1)$ sites with real coupling $J$ and inverse temperature
$\beta$ is
\[
  \rho_\beta^{\mathrm{open}} \;:=\;
  \frac{e^{-\beta H_{\mathrm{open}}(J,N)}}{\operatorname{Tr}\,e^{-\beta H_{\mathrm{open}}(J,N)}},
\]
i.e.\ the abstract \texttt{gibbsState} applied to
$H_{\mathrm{open}}(J,N) := \texttt{heisenbergHamiltonian (openChainCoupling N J)}$.
\end{definition}

\begin{definition}[\texttt{periodicChainHeisenbergGibbsState}]
Analogously for the periodic chain on $\mathrm{Fin}(N+2)$ sites:
$\rho_\beta^{\mathrm{per}} := \texttt{gibbsState}\,\beta\,
H_{\mathrm{per}}(J,N)$.
\end{definition}

\begin{theorem}[Hermiticity of the Heisenberg Gibbs states]
$\rho_\beta^{\mathrm{open}}$ and $\rho_\beta^{\mathrm{per}}$ are Hermitian.
Lean names: \texttt{openChainHeisenbergGibbsState\_isHermitian},
\texttt{periodicChainHeisenbergGibbsState\_isHermitian}.
\end{theorem}

\begin{proof}[Proof sketch]
Apply the abstract \texttt{gibbsState\_isHermitian} with the Hermitian
Hamiltonians $H_{\mathrm{open}}(J,N)$ /
$H_{\mathrm{per}}(J,N)$ from the previous subsection.
\end{proof}

\begin{theorem}[Commutativity with the Hamiltonian]
$[\rho_\beta^{\mathrm{open}}, H_{\mathrm{open}}] = 0$ and
$[\rho_\beta^{\mathrm{per}}, H_{\mathrm{per}}] = 0$, i.e.\ both Gibbs
states are stationary under the dynamics generated by their respective
Hamiltonians.
Lean names:
\texttt{openChainHeisenbergGibbsState\_commute\_hamiltonian},
\texttt{periodicChainHeisenbergGibbsState\_commute\_hamiltonian}.
\end{theorem}

\begin{proof}[Proof sketch]
Specialise the abstract
\texttt{gibbsState\_commute\_hamiltonian} to the corresponding
Heisenberg Hamiltonians.
\end{proof}

\begin{theorem}[High-temperature closed form]
At $\beta = 0$, for any observable $A$,
\[
  \langle A \rangle_0^{\mathrm{open}}
  \;=\; \dim^{-1}\,\operatorname{Tr}(A),
\]
and analogously for the periodic chain.
Lean names:
\texttt{openChainHeisenbergGibbsExpectation\_zero},
\texttt{periodicChainHeisenbergGibbsExpectation\_zero}.
\end{theorem}

\begin{proof}[Proof sketch]
Specialise \texttt{gibbsExpectation\_zero} to the Heisenberg
Hamiltonians; the Hamiltonian itself is irrelevant at $\beta = 0$ because
$e^{-0 \cdot H} = I$.
\end{proof}

\begin{theorem}[Expectation realness, conservation, energy realness]
For both the open and periodic Heisenberg Hamiltonians:
\begin{enumerate}
  \item For any Hermitian observable $O$,
    $\Im\langle O \rangle_\beta = 0$.
  \item For any observable $A$, $\langle [H, A] \rangle_\beta = 0$.
  \item $\Im\langle H \rangle_\beta = 0$ (energy expectation is real).
  \item For any Hermitian observable $O$,
    $\Im\langle H \cdot O \rangle_\beta = 0$.
  \item $\Im\langle H^2 \rangle_\beta = 0$ (energy-squared expectation is real).
  \item $\Im\langle H^n \rangle_\beta = 0$ for any $n \in \mathbb{N}$
    (all natural-power energy moments are real).
  \item For any Hermitian $A, B$,
    $\Im\langle A B + B A \rangle_\beta = 0$ (anticommutator real).
  \item For any Hermitian $A, B$,
    $\Re\langle A B - B A \rangle_\beta = 0$ (commutator purely imaginary).
  \item $\Im\operatorname{Var}_\beta(H) = 0$ (energy variance real).
\end{enumerate}
Lean names (open chain):
\texttt{openChainHeisenbergGibbsExpectation\_im\_of\_isHermitian},
\texttt{openChainHeisenbergGibbsExpectation\_commutator\_hamiltonian},
\texttt{openChainHeisenbergGibbsExpectation\_hamiltonian\_im},
\texttt{openChainHeisenbergGibbsExpectation\_mul\_hamiltonian\_im},
\texttt{openChainHeisenbergGibbsExpectation\_hamiltonian\_sq\_im},
\texttt{openChainHeisenbergGibbsExpectation\_hamiltonian\_pow\_im},
\texttt{openChainHeisenbergGibbsExpectation\_anticommutator\_im},
\texttt{openChainHeisenbergGibbsExpectation\_commutator\_re},
\texttt{openChainHeisenbergGibbsHamiltonianVariance\_im};
analogous for the periodic chain.
\end{theorem}

\begin{proof}[Proof sketch]
Direct specialisations of
\texttt{gibbsExpectation\_im\_of\_isHermitian} and
\texttt{gibbsExpectation\_commutator\_hamiltonian} at the Heisenberg
Hamiltonian, using the Hermiticity lemmas
\texttt{openChainHeisenberg\_isHermitian} /
\texttt{periodicChainHeisenberg\_isHermitian}.
\end{proof}

\subsection{Bond-sum decomposition of the energy expectation}
\label{subsec:heisenberg-bondsum}

\begin{theorem}[\texttt{heisenbergHamiltonian\_gibbsExpectation\_eq}]
For any Gibbs Hamiltonian $H$ and any coupling
$J : \Lambda \to \Lambda \to \mathbb{C}$,
\[
  \langle \texttt{heisenbergHamiltonian}\,J \rangle_\beta
  \;=\;
  \sum_{x \in \Lambda} \sum_{y \in \Lambda}
    J(x,y)\,\langle \hat S_x \cdot \hat S_y \rangle_\beta.
\]
\end{theorem}

\begin{proof}[Proof sketch]
Apply \texttt{gibbsExpectation\_sum} to the outer and inner sums of
$\sum_{x,y} J(x,y) \cdot \hat S_x \cdot \hat S_y$, then
\texttt{gibbsExpectation\_smul} pulls each scalar $J(x,y)$ out of the
expectation.
\end{proof}

\begin{remark}[Open- and periodic-chain energy bond-sums]
Specialising at
$H = \texttt{heisenbergHamiltonian (openChainCoupling N J)}$
(resp.\ \texttt{periodicChainCoupling}) yields
\[
  \langle H_{\mathrm{open}} \rangle_\beta
  \;=\;
  \sum_{x,y \in \mathrm{Fin}\,(N+1)}
    J_{\mathrm{open}}(x,y)\,\langle \hat S_x \cdot \hat S_y \rangle_\beta,
\]
and analogously for the periodic chain on $\mathrm{Fin}(N+2)$.
Lean names: \texttt{openChainHeisenbergGibbsExpectation\_self\_eq},
\texttt{periodicChainHeisenbergGibbsExpectation\_self\_eq}.
\end{remark}

\subsection{2-site explicit form and spectrum}
\label{subsec:two-site-spectrum}

The simplest concrete instance is $\Lambda = \mathrm{Fin}\,2$
($N = 1$ in \texttt{openChainCoupling 1 J}).  Direct evaluation gives
\[
  \texttt{heisenbergHamiltonian}\,(\texttt{openChainCoupling 1 J})
  \;=\; -2J\,\hat S_{0}\!\cdot\!\hat S_{1},
\]
a single bond term.  The four computational basis decompositions
$|{\uparrow}{\uparrow}\rangle$, $|{\downarrow}{\downarrow}\rangle$,
the singlet $|{\uparrow}{\downarrow}\rangle - |{\downarrow}{\uparrow}\rangle$,
and the triplet $|{\uparrow}{\downarrow}\rangle + |{\downarrow}{\uparrow}\rangle$
are exact eigenstates with eigenvalues $-J/2$, $-J/2$, $+3J/2$, and
$-J/2$ respectively (Lean theorems
\texttt{openChainHeisenbergHamiltonian\_two\_site\_mulVec\_basisVec\_all\_up},
\texttt{\_all\_down},
\texttt{\_singlet}, and \texttt{\_triplet\_zero}).
The four eigenvalues form the explicit Clebsch--Gordan decomposition
$\tfrac{1}{2}\otimes\tfrac{1}{2} = 0\oplus 1$: the singlet
($S=0$, eigenvalue $+3J/2$) and the triplet ($S=1$, three-fold
degenerate at $-J/2$).  Under our convention $H = -J\sum_{\langle
x,y\rangle}\hat S_x\!\cdot\!\hat S_y$ (Tasaki eq.~(2.4.1)), the
ferromagnetic regime is $J > 0$ (triplet ground state, energy
$-J/2$) and the antiferromagnetic regime is $J < 0$ (singlet ground
state, energy $+3J/2$).

\subsection{3-site explicit form}
\label{subsec:three-site-explicit}

For $N = 2$ (the 3-site open chain on $\mathrm{Fin}\,3$ with 2 bonds),
the Heisenberg Hamiltonian collapses to two bond terms:
\[
  \texttt{heisenbergHamiltonian}\,(\texttt{openChainCoupling 2 J})
  \;=\; -2J\,\bigl(\hat S_{0}\!\cdot\!\hat S_{1}
                    + \hat S_{1}\!\cdot\!\hat S_{2}\bigr).
\]
Acting on the all-up state gives
$\hat H \cdot |{\uparrow}{\uparrow}{\uparrow}\rangle
   = -J \cdot |{\uparrow}{\uparrow}{\uparrow}\rangle$, confirming the
linear scaling $E(|{\uparrow}\dots{\uparrow}\rangle) = -N\cdot J/2$
($N = 2$ bonds here).  Lean theorems
\texttt{openChainHeisenbergHamiltonian\_three\_site\_eq} and
\texttt{openChainHeisenbergHamiltonian\_three\_site\_mulVec\_basisVec\_all\_up}.

\subsection{General open-chain ferromagnetic ground-state energy
(Tasaki §2.4 (2.4.5)/(2.4.1))}
\label{subsec:openchain-allup-general}

Generalising the 2-site / 3-site cases, for any chain length
$N + 1$ the all-up state is an exact eigenvector of the open-chain
Heisenberg Hamiltonian with eigenvalue $-N\cdot J/2$, matching
Tasaki's ferromagnetic ground-state energy
$E_{\mathrm{GS}} = -|B|\cdot S^{2}$ for $S = 1/2$ and $|B| = N$
unordered bonds.

\begin{lemma}[\texttt{openChainCoupling\_sum\_eq}]
For any $N \in \mathbb{N}$ and $J \in \mathbb{R}$,
\[
  \sum_{x \in \mathrm{Fin}(N+1)} \sum_{y \in \mathrm{Fin}(N+1)}
    J_{\mathrm{open}}(x,y) \;=\; -2N\cdot J.
\]
\end{lemma}

\begin{proof}[Proof sketch]
Split the disjunction $i+1 = j \lor j+1 = i$ into two disjoint
subsums (the conjunction is impossible).  For each fixed $x$, the
right-neighbour count is $1$ if $x_{\mathrm{val}} < N$ else $0$; the
left-neighbour count is $1$ if $x_{\mathrm{val}} > 0$ else $0$
(\texttt{openChain\_row\_sum\_succ} and \texttt{openChain\_row\_sum\_pred}).
Each row sum, when summed over $x$, contributes $N$ via Finset
filter cardinalities (\texttt{sum\_lt\_eq} and \texttt{sum\_pos\_eq}).
Combining gives $-J\cdot N + (-J)\cdot N = -2N\cdot J$.
\end{proof}

\begin{theorem}[\texttt{openChainHeisenbergHamiltonian\_mulVec\_basisVec\_all\_up}]
For any $N \in \mathbb{N}$ and $J \in \mathbb{R}$,
\[
  \texttt{heisenbergHamiltonian (openChainCoupling }N\, J)
    \cdot |{\uparrow}\dots{\uparrow}\rangle
  \;=\; -\frac{N\cdot J}{2}\cdot
    |{\uparrow}\dots{\uparrow}\rangle.
\]
This is the Tasaki ferromagnetic ground-state energy
$E_{\mathrm{GS}} = -|B|\cdot S^{2}$ for $S = 1/2$, $|B| = N$ bonds.
\end{theorem}

\begin{proof}[Proof sketch]
By the constant-config eigenvalue
(\texttt{heisenbergHamiltonian\_mulVec\_basisVec\_const}), the
eigenvalue is $\sum_{x,y} J_{\mathrm{open}}(x,y)\cdot \chi_{x,y}$.
The diagonal terms vanish since $J_{\mathrm{open}}(x,x) = 0$, so the
sum equals $(1/4)\cdot \sum_{x,y} J_{\mathrm{open}}(x,y)$. Apply the
counting lemma above to obtain $(1/4)\cdot (-2N\cdot J) = -N\cdot J/2$.
\end{proof}

\begin{remark}[Symmetric specialisations]
The proof above generalises to any constant configuration
$|s\,s\,\dots\,s\rangle$ ($s \in \{0,1\}$):
$\hat H_{\mathrm{open}}\cdot|s\,s\,\dots\,s\rangle = -\tfrac{N\cdot J}{2}
\cdot|s\,s\,\dots\,s\rangle$, since the constant-config eigenvalue
formula does not depend on $s$ (Lean
\texttt{openChainHeisenbergHamiltonian\_mulVec\_basisVec\_const}).
The all-down case
\texttt{openChainHeisenbergHamiltonian\_mulVec\_basisVec\_all\_down}
follows by SU(2) symmetry of the Heisenberg Hamiltonian.
\end{remark}

\begin{theorem}[\texttt{openChainHeisenbergHamiltonian\_mulVec\_totalSpinHalfOpMinus\_pow\_basisVec\_all\_up}]
For any $N \in \mathbb{N}$, $J \in \mathbb{R}$, and $k \in \mathbb{N}$,
\[
  \hat H_{\mathrm{open}} \cdot
    \bigl((\hat S_{\mathrm{tot}}^{-})^{k} \cdot
            |{\uparrow}\dots{\uparrow}\rangle\bigr)
  \;=\;
  -\frac{N\cdot J}{2}\;
    \bigl((\hat S_{\mathrm{tot}}^{-})^{k} \cdot
            |{\uparrow}\dots{\uparrow}\rangle\bigr).
\]
The dual ladder from $|{\downarrow}\dots{\downarrow}\rangle$
($\hat S_{\mathrm{tot}}^{+}$) shares the same eigenvalue
(\texttt{openChainHeisenbergHamiltonian\_mulVec\_totalSpinHalfOpPlus\_pow\_basisVec\_all\_down}).
This makes Tasaki eq.~(2.4.9) explicit for the open chain.
\end{theorem}

\begin{proof}[Proof sketch]
Apply
\texttt{heisenbergHamiltonian\_mulVec\_totalSpinHalfOpMinus\_pow\_basisVec\_const}
(\ref{subsec:heisenberg-ladder}, iterated lowering preserves the
H-eigenvalue) to reduce the goal to showing $c_{J} = -N\cdot J/2$ for
the chain coupling.  Then proceed exactly as in the const-config
proof above: drop the diagonal contribution and apply the bond
counting lemma.
\end{proof}

\subsection{Action on the fully-polarised state (Tasaki §2.4 (2.4.5))}
\label{subsec:heisenberg-allup}

Tasaki \emph{Physics and Mathematics of Quantum Many-Body Systems},
§2.4 ``The Ferromagnetic Heisenberg Model'', eq.~(2.4.5) (p.~32),
asserts that the fully-polarised state
$|\Phi^{\uparrow}\rangle = \bigotimes_{x} |\psi_{x}^{S}\rangle$
satisfies
\[
  -\hat S_x \cdot \hat S_y \, |\Phi^{\uparrow}\rangle
  \;=\; -\hat S_x^{(3)} \hat S_y^{(3)} \, |\Phi^{\uparrow}\rangle
  \;=\; -S^{2}\, |\Phi^{\uparrow}\rangle,
\]
so each bond term contributes the minimum eigenvalue $-S^{2}$ and
the ferromagnetic ground-state energy is $E_{\mathrm{GS}}=-|B|S^{2}$.

For $S=1/2$, every Heisenberg-type Hamiltonian acts as a scalar on
every uniformly-aligned basis state.

\begin{theorem}[\texttt{heisenbergHamiltonian\_mulVec\_basisVec\_const}]
For every $J : \Lambda \to \Lambda \to \mathbb{C}$ and every
$s \in \{0,1\} = \{\uparrow,\downarrow\}$,
\[
  \texttt{heisenbergHamiltonian}\,J \cdot \,
    |\,\underbrace{s\,s\,\dots\,s}_{|\Lambda|}\,\rangle
  \;=\;
  \Bigl(\sum_{x,y\in\Lambda} J(x,y)\,
    \chi_{x,y}\Bigr) \cdot
    |\,s\,s\,\dots\,s\,\rangle,
\]
where $\chi_{x,y} := \tfrac{3}{4}$ if $x=y$ (the diagonal Casimir
contribution $S(S+1)=3/4$) and $\chi_{x,y}:=\tfrac{1}{4}$ otherwise
(the maximum eigenvalue $S^{2}=1/4$ of $\hat S_x\cdot\hat S_y$ on
parallel spins).
\end{theorem}

\begin{proof}[Proof sketch]
Unfold \texttt{heisenbergHamiltonian} to
$\sum_{x,y} J(x,y)\,\hat S_x\cdot\hat S_y$. By
\texttt{Matrix.sum\_mulVec} each summand acts on the constant basis
state independently. The internal lemma
\texttt{spinHalfDot\_mulVec\_const} provides the eigenvalue
$\tfrac{3}{4}$ for $x=y$ (using \texttt{spinHalfDot\_self}, which
gives $\hat S_x\cdot\hat S_x = (3/4)\,I$) and $\tfrac{1}{4}$ for
$x\ne y$ (using \texttt{spinHalfDot\_mulVec\_basisVec\_parallel};
both sites carry the same value~$s$). Distributivity of scalar
multiplication and \texttt{Finset.sum\_smul} consolidate the bilinear
sum into a single eigenvalue.
\end{proof}

\begin{remark}[Specialisations]
The all-up case $s=0$ (\texttt{heisenbergHamiltonian\_mulVec\_basisVec\_all\_up})
and the all-down case $s=1$
(\texttt{heisenbergHamiltonian\_mulVec\_basisVec\_all\_down}) are
direct corollaries.

To recover Tasaki's ferromagnetic ground-state energy
$E_{\mathrm{GS}}=-|B|\,S^{2}=-|B|/4$ from
$H = -\sum_{\{x,y\}\in B}\hat S_x\cdot\hat S_y$ (sum over unordered
bonds, no self-bonds) in our ordered-pair convention, take
$J(x,y)=-\tfrac{1}{2}$ for both orientations of every bond
$\{x,y\}\in B$ and $J(x,y)=0$ otherwise.  Then the bilinear sum
contains $2|B|$ nonzero terms each contributing
$(-\tfrac{1}{2})(\tfrac{1}{4})=-\tfrac{1}{8}$, so
$\sum_{x,y}J(x,y)\,\chi_{x,y}=-\tfrac{|B|}{4}=-|B|\,S^{2}$.

The present theorem proves only the \emph{eigenvector} property — that
the constant configuration is an eigenvector of every Heisenberg-type
Hamiltonian.  Whether the eigenvalue equals the ground-state energy
depends on the sign and structure of $J$ and is not asserted here.
Tasaki's argument that $|\Phi^{\uparrow}\rangle$ \emph{minimises} the
ferromagnetic energy uses Lemma~A.9 (frustration-free Hamiltonians);
that step is not yet formalised.
\end{remark}

\begin{remark}[Relation to Tasaki (2.4.5)]
Tasaki's eq.~(2.4.5) is the bond statement
$-\hat S_x\!\cdot\!\hat S_y\,|\Phi^{\uparrow}\rangle
   = -S^{2}\,|\Phi^{\uparrow}\rangle$ for $x\ne y$.  The present theorem
is the bilinear-sum lift of that statement (and of the diagonal
identity $\hat S_x\!\cdot\!\hat S_x = S(S+1)\,I$ from
\texttt{spinHalfDot\_self}); it subsumes (2.4.5) for the off-diagonal
indices and additionally records the diagonal Casimir contribution.
\end{remark}

\subsection{Ladder step on the constant configuration (Tasaki §2.4 (2.4.7), (2.4.9))}
\label{subsec:heisenberg-ladder}

Tasaki's enumeration of ferromagnetic ground states proceeds by acting
the global lowering operator on the all-up state:
\[
  |\Phi_{M}\rangle
  \;:=\; \frac{(\hat S_{\mathrm{tot}}^{-})^{|\Lambda|S - M}\,
                  |\Phi^{\uparrow}\rangle}
                {\bigl\|(\hat S_{\mathrm{tot}}^{-})^{|\Lambda|S - M}\,
                          |\Phi^{\uparrow}\rangle\bigr\|},
  \qquad
  \hat H\,|\Phi_{M}\rangle = E_{\mathrm{GS}}\,|\Phi_{M}\rangle.
\]
The eigenvalue is preserved at each ladder step because $\hat
S_{\mathrm{tot}}^{\pm}$ commutes with the SU(2)-invariant
Hamiltonian.  Our spin-1/2 formalisation captures both the single-step
form (below) and the iterated form (further below) for arbitrary
constant configurations $|s\,s\,\dots\,s\rangle$.

\begin{lemma}[\texttt{heisenbergHamiltonian\_mulVec\_totalSpinHalfOpMinus\_basisVec\_const}]
For every $J : \Lambda \to \Lambda \to \mathbb{C}$ and every
$s \in \{0,1\}$,
\[
  \texttt{heisenbergHamiltonian}\,J \cdot
    \bigl(\hat S_{\mathrm{tot}}^{-} \cdot |s\,s\,\dots\,s\rangle\bigr)
  \;=\;
  c_{J}\;
    \bigl(\hat S_{\mathrm{tot}}^{-} \cdot |s\,s\,\dots\,s\rangle\bigr),
\]
where $c_{J} = \sum_{x,y}J(x,y)\,\chi_{x,y}$ is the same eigenvalue as
in Theorem~\ref{subsec:heisenberg-allup}.  An identical statement holds
for $\hat S_{\mathrm{tot}}^{+}$
(\texttt{heisenbergHamiltonian\_mulVec\_totalSpinHalfOpPlus\_basisVec\_const}).
\end{lemma}

\begin{proof}[Proof sketch]
Combine
\texttt{heisenbergHamiltonian\_commutator\_totalSpinHalfOpMinus} (which
gives $\hat H \cdot \hat S_{\mathrm{tot}}^{-} = \hat S_{\mathrm{tot}}^{-}
\cdot \hat H$ via \texttt{sub\_eq\_zero}) with
\texttt{Matrix.mulVec\_mulVec} to swap $\hat H$ past
$\hat S_{\mathrm{tot}}^{-}$, apply the constant-configuration eigenvalue
\texttt{heisenbergHamiltonian\_mulVec\_basisVec\_const} on the inner
$\hat H \cdot |s\,s\,\dots\,s\rangle$, and pull the resulting scalar
out via \texttt{Matrix.mulVec\_smul}.
\end{proof}

\begin{remark}[Toward the ferromagnetic ground-state lattice]
Iterating the lemma gives
$\hat H \cdot (\hat S_{\mathrm{tot}}^{-})^{k}\,|\Phi^{\uparrow}\rangle
   = c_{J}\,(\hat S_{\mathrm{tot}}^{-})^{k}\,|\Phi^{\uparrow}\rangle$
for every $k \geq 0$, which is the unnormalised version of Tasaki
eq.~(2.4.9): all $|\Phi_{M}\rangle$ share the eigenvalue of
$|\Phi^{\uparrow}\rangle$ under any Heisenberg Hamiltonian.  The
normalisation factor of (2.4.9) and the corresponding
\emph{completeness} of $\{|\Phi_{M}\rangle\}_{M}$ as the ferromagnetic
ground-state space (Tasaki Theorem~2.1) are not yet formalised.
\end{remark}

\begin{theorem}[\texttt{heisenbergHamiltonian\_mulVec\_totalSpinHalfOpMinus\_pow\_basisVec\_const}]
For every $J : \Lambda \to \Lambda \to \mathbb{C}$, every
$s \in \{0,1\}$, and every $k \in \mathbb{N}$,
\[
  \texttt{heisenbergHamiltonian}\,J \cdot
    \bigl((\hat S_{\mathrm{tot}}^{-})^{k} \cdot
            |s\,s\,\dots\,s\rangle\bigr)
  \;=\;
  c_{J}\;
    \bigl((\hat S_{\mathrm{tot}}^{-})^{k} \cdot
            |s\,s\,\dots\,s\rangle\bigr).
\]
A companion statement holds for $\hat S_{\mathrm{tot}}^{+}$
(\texttt{heisenbergHamiltonian\_mulVec\_totalSpinHalfOpPlus\_pow\_basisVec\_const}).
\end{theorem}

\begin{proof}[Proof sketch]
Induction on $k$. The base case $k=0$ uses
\texttt{pow\_zero} and \texttt{Matrix.one\_mulVec} to reduce to
Theorem~\ref{subsec:heisenberg-allup}. The inductive step uses
\texttt{pow\_succ'}
($M^{k+1} = M\cdot M^{k}$) to peel one factor of
$\hat S_{\mathrm{tot}}^{-}$, swaps it past $\hat H$ via
\texttt{heisenbergHamiltonian\_commutator\_totalSpinHalfOpMinus} (turned
into the commutation $\hat H\cdot\hat S_{\mathrm{tot}}^{-}
=\hat S_{\mathrm{tot}}^{-}\cdot\hat H$ by
\texttt{sub\_eq\_zero.mp}), applies the inductive hypothesis to the
inner $\hat H \cdot \bigl((\hat S_{\mathrm{tot}}^{-})^{k}\cdot
|s\,s\,\dots\,s\rangle\bigr)$, and pulls the resulting scalar out of
the leftmost $\hat S_{\mathrm{tot}}^{-}$ via
\texttt{Matrix.mulVec\_smul}.
\end{proof}

\subsection{Commutativity with global rotation (Tasaki §2.4 (2.4.7))}
\label{subsec:heisenberg-rot-commute}

Tasaki's argument that the spin-coherent state
$|\Xi_{\theta,\phi}\rangle = \hat U^{(3)}_{\phi}\hat U^{(2)}_{\theta}
|\Phi^{\uparrow}\rangle$ of eq.~(2.4.6) is a ground state of $\hat H$
runs through the operator-level identity
$[\hat H, \hat U^{(\alpha)}_{\theta}] = 0$.  In the formalisation we
already have $[\hat H, \hat S_{\mathrm{tot}}^{(\alpha)}] = 0$
(see~\ref{subsec:heisenberg-bondsum}), and
\texttt{totalSpinHalfRot\{1,2,3\}\_commute\_of\_commute} from
\texttt{TotalSpin.lean} lifts this to the exponential.

\begin{lemma}[\texttt{heisenbergHamiltonian\_commute\_totalSpinHalfRot\{1,2,3\}}]
For every $J : \Lambda \to \Lambda \to \mathbb{C}$, every $\theta \in
\mathbb{R}$, and every axis $\alpha \in \{1,2,3\}$,
\[
  \texttt{heisenbergHamiltonian}\,J \cdot
    \hat U^{(\alpha)}_{\theta}
  \;=\;
  \hat U^{(\alpha)}_{\theta} \cdot
    \texttt{heisenbergHamiltonian}\,J.
\]
\end{lemma}

\begin{proof}[Proof sketch]
Convert
\texttt{heisenbergHamiltonian\_commutator\_totalSpinHalfOp\{$\alpha$\}}
($H\cdot S - S\cdot H = 0$) into a \texttt{Commute} witness via
\texttt{sub\_eq\_zero.mp}, then apply
\texttt{totalSpinHalfRot\{$\alpha$\}\_commute\_of\_commute}, which
internally lifts \texttt{Commute} past
\texttt{NormedSpace.exp} via \texttt{Commute.exp\_right}.
\end{proof}

\begin{theorem}[\texttt{heisenbergHamiltonian\_mulVec\_totalSpinHalfRot\{1,2,3\}\_basisVec\_const}]
For every $J$, every $\theta$, every axis $\alpha$, and every
$s \in \{0,1\}$,
\[
  \texttt{heisenbergHamiltonian}\,J \cdot
    \bigl(\hat U^{(\alpha)}_{\theta} \cdot
            |s\,s\,\dots\,s\rangle\bigr)
  \;=\;
  c_{J}\;
    \bigl(\hat U^{(\alpha)}_{\theta} \cdot
            |s\,s\,\dots\,s\rangle\bigr).
\]
Composing two axes (axis~2 then axis~3, as in Tasaki eq.~(2.4.6)) gives
$\hat H \cdot |\Xi_{\theta,\phi}\rangle
   = c_{J} \cdot |\Xi_{\theta,\phi}\rangle$
of Tasaki eq.~(2.4.7).  This composition is formalised below as
\texttt{heisenbergHamiltonian\_mulVec\_totalSpinHalfRot32\_basisVec\_const}.
\end{theorem}

\begin{proof}[Proof sketch]
Use \texttt{Matrix.mulVec\_mulVec} to merge the two outer matrices,
the lemma above to swap $\hat H$ past $\hat U^{(\alpha)}_{\theta}$,
\texttt{Matrix.mulVec\_mulVec} again to expand back, the
constant-config eigenvalue
\texttt{heisenbergHamiltonian\_mulVec\_basisVec\_const} on the inner
product, and \texttt{Matrix.mulVec\_smul} to pull the scalar out.
\end{proof}

\begin{theorem}[\texttt{heisenbergHamiltonian\_mulVec\_totalSpinHalfRot32\_basisVec\_const}]
For every $J : \Lambda \to \Lambda \to \mathbb{C}$, every
$\theta,\phi \in \mathbb{R}$, and every $s \in \{0,1\}$,
\[
  \texttt{heisenbergHamiltonian}\,J \cdot
    \bigl(\hat U^{(3)}_{\phi} \cdot \hat U^{(2)}_{\theta}
            \cdot |s\,s\,\dots\,s\rangle\bigr)
  \;=\;
  c_{J}\;
    \bigl(\hat U^{(3)}_{\phi} \cdot \hat U^{(2)}_{\theta}
            \cdot |s\,s\,\dots\,s\rangle\bigr).
\]
Specialised to $s = 0$, this is Tasaki eq.~(2.4.7) for the
spin-coherent state $|\Xi_{\theta,\phi}\rangle = \hat U^{(3)}_{\phi}
\hat U^{(2)}_{\theta} |\Phi^{\uparrow}\rangle$ of eq.~(2.4.6).
\end{theorem}

\begin{proof}[Proof sketch]
Apply the single-axis rotated-eigenvector lemma (above) twice.
Using \texttt{Matrix.mulVec\_mulVec} merge $\hat H \cdot \hat
U^{(3)}_{\phi}$, swap them via
\texttt{heisenbergHamiltonian\_commute\_totalSpinHalfRot3}, expand
back to $\hat U^{(3)}_{\phi} \cdot (\hat H \cdot \hat U^{(2)}_{\theta}
\cdot |s\,s\,\dots\,s\rangle)$, then invoke
\texttt{heisenbergHamiltonian\_mulVec\_totalSpinHalfRot2\_basisVec\_const}
on the inner application and pull the resulting scalar out via
\texttt{Matrix.mulVec\_smul}.
\end{proof}

\subsection{Magnetic-quantum-number ladder on the all-up state (Tasaki §2.4 (2.4.9))}
\label{subsec:magnetic-ladder}

The companion to the \emph{$H$-eigenvalue} ladder is the
\emph{$\hat S_{\mathrm{tot}}^{(3)}$-eigenvalue} ladder.  The
all-up state has $\hat S_{\mathrm{tot}}^{(3)}$ eigenvalue
$S_{\max} = |\Lambda|/2$, and each application of $\hat
S_{\mathrm{tot}}^{-}$ lowers the eigenvalue by one (the Cartan
ladder relation $[\hat S_{\mathrm{tot}}^{(3)},\hat
S_{\mathrm{tot}}^{-}] = -\hat S_{\mathrm{tot}}^{-}$).

\begin{theorem}[\texttt{totalSpinHalfOp3\_mulVec\_totalSpinHalfOpMinus\_pow\_basisVec\_all\_up}]
For every $k \in \mathbb{N}$,
\[
  \hat S_{\mathrm{tot}}^{(3)} \cdot
    \bigl((\hat S_{\mathrm{tot}}^{-})^{k} \cdot
            |{\uparrow}\dots{\uparrow}\rangle\bigr)
  \;=\;
  \Bigl(\tfrac{|\Lambda|}{2} - k\Bigr) \cdot
    \bigl((\hat S_{\mathrm{tot}}^{-})^{k} \cdot
            |{\uparrow}\dots{\uparrow}\rangle\bigr).
\]
This is the magnetic-quantum-number $M = S_{\max} - k$ labelling of
the iterated lowering ladder
$|\Phi_{M}\rangle \propto (\hat S_{\mathrm{tot}}^{-})^{S_{\max} - M}
|\Phi^{\uparrow}\rangle$ in Tasaki eq.~(2.4.9), p.~33.  The
normalisation factor and the restriction $M \in \{-S_{\max},\ldots,
+S_{\max}\}$ are not asserted.
\end{theorem}

\begin{proof}[Proof sketch]
Induction on $k$.  The base case $k=0$ uses
\texttt{totalSpinHalfOp3\_mulVec\_basisVec} and the identity
$\sum_{x \in \Lambda} \tfrac{1}{2} = |\Lambda|/2$ on the constant
all-up configuration.  The inductive step rewrites $\hat
S_{\mathrm{tot}}^{(3)} \cdot \hat S_{\mathrm{tot}}^{-}$ as $\hat
S_{\mathrm{tot}}^{-} \cdot \hat S_{\mathrm{tot}}^{(3)} - \hat
S_{\mathrm{tot}}^{-}$ via
\texttt{totalSpinHalfOp3\_commutator\_totalSpinHalfOpMinus} (the
Cartan relation), then uses \texttt{Matrix.sub\_mulVec},
\texttt{Matrix.mulVec\_mulVec}, and \texttt{pow\_succ'} to apply the
inductive hypothesis to the inner $\hat S_{\mathrm{tot}}^{(3)}\cdot
(\hat S_{\mathrm{tot}}^{-})^{k}\cdot|{\uparrow}\dots{\uparrow}\rangle$,
and finishes with $c\,v - v = (c-1)\,v$ (with $c = |\Lambda|/2 - k$
the previous-step eigenvalue, so $c-1 = |\Lambda|/2 - (k+1)$) via the
\texttt{module} tactic on the resulting scalar smul.
\end{proof}

\begin{theorem}[\texttt{totalSpinHalfOp3\_mulVec\_totalSpinHalfOpPlus\_pow\_basisVec\_all\_down}]
For every $k \in \mathbb{N}$,
\[
  \hat S_{\mathrm{tot}}^{(3)} \cdot
    \bigl((\hat S_{\mathrm{tot}}^{+})^{k} \cdot
            |{\downarrow}\dots{\downarrow}\rangle\bigr)
  \;=\;
  \Bigl(-\tfrac{|\Lambda|}{2} + k\Bigr) \cdot
    \bigl((\hat S_{\mathrm{tot}}^{+})^{k} \cdot
            |{\downarrow}\dots{\downarrow}\rangle\bigr).
\]
This is the dual ladder of the previous theorem, parameterising the
same ferromagnetic-ground-state lattice from the lowest weight
$M = -S_{\max}$ upward.
\end{theorem}

\begin{proof}[Proof sketch]
Identical structure to the lowering version, replacing $\hat
S_{\mathrm{tot}}^{-} \mapsto \hat S_{\mathrm{tot}}^{+}$, the Cartan
sign $-\hat S_{\mathrm{tot}}^{-} \mapsto +\hat S_{\mathrm{tot}}^{+}$
(\texttt{totalSpinHalfOp3\_commutator\_totalSpinHalfOpPlus}), and
$\sum_{x} \tfrac{1}{2} \mapsto \sum_{x} (-\tfrac{1}{2}) =
-|\Lambda|/2$ in the base case.  The inductive step closes via $c\,v
+ v = (c+1)\,v$ in place of subtraction.
\end{proof}

\begin{theorem}[\texttt{totalSpinHalfSquared\_mulVec\_totalSpinHalfOpMinus\_pow\_basisVec\_const}]
For every $s \in \{0,1\}$ and every $k \in \mathbb{N}$,
\[
  (\hat S_{\mathrm{tot}})^{2}
    \cdot \bigl((\hat S_{\mathrm{tot}}^{-})^{k} \cdot
                 |s\,s\,\dots\,s\rangle\bigr)
  \;=\; S_{\max}(S_{\max}+1)\;
    \bigl((\hat S_{\mathrm{tot}}^{-})^{k} \cdot
            |s\,s\,\dots\,s\rangle\bigr),
\]
with $S_{\max} = |\Lambda|/2$.  The companion theorem holds with
$\hat S_{\mathrm{tot}}^{+}$ replacing $\hat S_{\mathrm{tot}}^{-}$.
\end{theorem}

\begin{proof}[Proof sketch]
Induction on $k$, identical template to the $\hat H$ case
(\ref{subsec:heisenberg-ladder}): the base case is the constant-config
Casimir eigenvalue $\hat S_{\mathrm{tot}}^{2}\,
|s\,s\,\dots\,s\rangle = S_{\max}(S_{\max}+1)\,|s\,s\,\dots\,s\rangle$
(\texttt{totalSpinHalfSquared\_mulVec\_basisVec\_const}); the step
swaps $\hat S_{\mathrm{tot}}^{2}$ past $\hat S_{\mathrm{tot}}^{\mp}$
via \texttt{totalSpinHalfSquared\_commutator\_totalSpinHalfOp\{Plus,Minus\}}
(both zero, since $\hat S_{\mathrm{tot}}^{2}$ is the SU(2) Casimir),
applies the inductive hypothesis, and pulls the scalar out via
\texttt{Matrix.mulVec\_smul}.
\end{proof}

\begin{remark}[Simultaneous-eigenvector packaging]
For the matched highest- and lowest-weight ladders, combining this
Casimir invariance with the $\hat H$ ladder
(\ref{subsec:heisenberg-ladder}) and the $\hat S_{\mathrm{tot}}^{(3)}$
ladder (\ref{subsec:magnetic-ladder}) shows that
$(\hat S_{\mathrm{tot}}^{-})^{k}\,|{\uparrow}\dots{\uparrow}\rangle$
and $(\hat S_{\mathrm{tot}}^{+})^{k}\,|{\downarrow}\dots{\downarrow}\rangle$
are simultaneous eigenvectors of $(\hat H,\,\hat S_{\mathrm{tot}}^{2},\,
\hat S_{\mathrm{tot}}^{(3)})$ with eigenvalues
$(c_{J},\,S_{\max}(S_{\max}+1),\,S_{\max}\mp k)$ — a complete SU(2)
highest- (or lowest-) weight labelling.  The Casimir invariance holds
for any constant configuration, but the matching $\hat
S_{\mathrm{tot}}^{(3)}$ eigenvalue is only formalised for these
matched ladders; for mismatched cases (e.g.\ $\hat
S_{\mathrm{tot}}^{-}|{\downarrow}\dots{\downarrow}\rangle$) the
iterates already vanish at $k = 1$, so the eigenvalue equation holds
trivially.
\end{remark}

\begin{remark}[Boundary annihilation]
The ladder eventually terminates at the opposite-weight constant
configuration:
\begin{align*}
  \hat S_{\mathrm{tot}}^{-} \cdot
    |{\downarrow}\dots{\downarrow}\rangle &= 0,\\
  \hat S_{\mathrm{tot}}^{+} \cdot
    |{\uparrow}\dots{\uparrow}\rangle &= 0.
\end{align*}
Each on-site lowering $\hat S_{x}^{-}\,|{\downarrow}\rangle = 0$
already at the minimum, summed over $x \in \Lambda$ gives the global
identity (and dually for raising at the maximum).  Lean names:
\texttt{totalSpinHalfOpMinus\_mulVec\_basisVec\_all\_down} and
\texttt{totalSpinHalfOpPlus\_mulVec\_basisVec\_all\_up}.  These are
the boundary conditions ensuring $\bigl((\hat
S_{\mathrm{tot}}^{-})^{k}\,|{\uparrow}\dots{\uparrow}\rangle\bigr)_{k}$
of Tasaki eq.~(2.4.9) traverses the magnetisation sectors
$\mathcal H_{S_{\max}},\mathcal H_{S_{\max}-1},\dots,
\mathcal H_{-S_{\max}}$ before becoming zero on further iteration.
\end{remark}

\begin{remark}[Magnetisation-subspace placement]
Combining the two ladders with \texttt{mem\_magnetizationSubspace\_iff}
yields
\begin{align*}
  (\hat S_{\mathrm{tot}}^{-})^{k}\,|{\uparrow}\dots{\uparrow}\rangle
    &\in \mathcal H_{S_{\max} - k},\\
  (\hat S_{\mathrm{tot}}^{+})^{k}\,|{\downarrow}\dots{\downarrow}\rangle
    &\in \mathcal H_{-S_{\max} + k},
\end{align*}
where $\mathcal H_{M}$ is the magnetisation-$M$ subspace of Tasaki
eq.~(2.2.10) (formalised as \texttt{magnetizationSubspace} in
\texttt{Quantum/MagnetizationSubspace.lean}).  Lean names:
\texttt{totalSpinHalfOpMinus\_pow\_basisVec\_all\_up\_mem\_magnetizationSubspace}
and
\texttt{totalSpinHalfOpPlus\_pow\_basisVec\_all\_down\_mem\_magnetizationSubspace}.
\end{remark}

\subsection{Two-site Néel and singlet/triplet placement (Tasaki §2.5/§A.3)}
\label{subsec:two-site-singlet-mag}

For $\Lambda = \mathrm{Fin}\,2$, the antiparallel basis states
$|{\uparrow}{\downarrow}\rangle$ and $|{\downarrow}{\uparrow}\rangle$
both have total magnetisation zero, so they lie in $\mathcal H_{0}$.
By Submodule closure, so do their sum $|{\uparrow}{\downarrow}\rangle
+ |{\downarrow}{\uparrow}\rangle$ (the triplet $m=0$ state) and
difference $|{\uparrow}{\downarrow}\rangle -
|{\downarrow}{\uparrow}\rangle$ (the singlet).

\begin{theorem}[\texttt{basisVec\_upDown\_mem\_magnetizationSubspace\_zero},
\texttt{basisVec\_basisSwap\_upDown\_mem\_magnetizationSubspace\_zero}]
The two-site Néel state $|{\uparrow}{\downarrow}\rangle$ and the
swapped state $|{\downarrow}{\uparrow}\rangle$ both lie in
$\mathcal H_{0}$ on $\Lambda = \mathrm{Fin}\,2$.
\end{theorem}

\begin{proof}[Proof sketch]
Apply \texttt{basisVec\_mem\_magnetizationSubspace} to the
configurations $\sigma = \texttt{upDown}$ and
$\sigma = \texttt{basisSwap}\,\texttt{upDown}\,0\,1$, then compute
the magnetisation $\sum_{x}\texttt{spinSign}(\sigma\,x) = 0$ via
\texttt{Fin.sum\_univ\_two}.
\end{proof}

\begin{proposition}[\texttt{singlet\_mem\_magnetizationSubspace\_zero},
\texttt{triplet\_zero\_mem\_magnetizationSubspace\_zero}]
The spin-singlet $|{\uparrow}{\downarrow}\rangle -
|{\downarrow}{\uparrow}\rangle$ and the triplet $m=0$ state
$|{\uparrow}{\downarrow}\rangle + |{\downarrow}{\uparrow}\rangle$
both lie in $\mathcal H_{0}$, by \texttt{Submodule.\{add,sub\}\_mem}.
\end{proposition}

\begin{theorem}[\texttt{mulVec\_mem\_magnetizationSubspace\_of\_commute\_of\_mem}]
For any operator $A$ with $[A,\hat S_{\mathrm{tot}}^{(3)}] = 0$ and
any $M \in \mathbb{C}$,
$v \in \mathcal H_{M}$ implies $A \cdot v \in \mathcal H_{M}$.

A direct corollary
(\texttt{totalSpinHalfSquared\_mulVec\_mem\_magnetizationSubspace\_of\_mem})
is that the Casimir $\hat S_{\mathrm{tot}}^{2}$ preserves every
magnetisation sector, since
$[\hat S_{\mathrm{tot}}^{2}, \hat S_{\mathrm{tot}}^{(3)}] = 0$
(\texttt{totalSpinHalfSquared\_commutator\_totalSpinHalfOp3}).
\end{theorem}

\begin{proof}[Proof sketch]
\texttt{mem\_magnetizationSubspace\_iff} reduces membership to the
eigenvalue equation $\hat S_{\mathrm{tot}}^{(3)} \cdot v = M \cdot v$.
Using \texttt{Matrix.mulVec\_mulVec} merge $\hat S_{\mathrm{tot}}^{(3)}
\cdot A$, swap them via \texttt{h.symm} (the supplied $\texttt{Commute}\,
A\,\hat S_{\mathrm{tot}}^{(3)}$ witness reversed), expand back via
\texttt{← Matrix.mulVec\_mulVec}, apply the inner eigenvalue equation,
and pull the scalar out via \texttt{Matrix.mulVec\_smul}.
\end{proof}

\begin{theorem}[\texttt{heisenbergHamiltonian\_mulVec\_mem\_magnetizationSubspace\_of\_mem}]
For any $J : \Lambda \to \Lambda \to \mathbb{C}$ and any $M \in
\mathbb{C}$, if $v \in \mathcal H_{M}$ then
$\texttt{heisenbergHamiltonian}\,J \cdot v \in \mathcal H_{M}$.

Equivalently, the Heisenberg Hamiltonian block-diagonalises against
the magnetisation decomposition of Tasaki eq.~(2.2.10).
\end{theorem}

\begin{proof}[Proof sketch]
Unfold $v \in \mathcal H_{M}$ to the eigenvalue equation
$\hat S_{\mathrm{tot}}^{(3)} \cdot v = M \cdot v$.  Convert the SU(2)
invariance
\texttt{heisenbergHamiltonian\_commutator\_totalSpinHalfOp3} into the
commutation $\hat S_{\mathrm{tot}}^{(3)} \cdot \hat H = \hat H \cdot
\hat S_{\mathrm{tot}}^{(3)}$ via \texttt{abel}.  Then
\texttt{Matrix.mulVec\_mulVec} merges the two outer matrices, the
commutation swaps them, \texttt{Matrix.mulVec\_mulVec} again expands
back, and the inner $\hat S_{\mathrm{tot}}^{(3)} \cdot v = M \cdot v$
applies, leaving $\hat S_{\mathrm{tot}}^{(3)} \cdot (\hat H \cdot v)
= \hat H \cdot (M \cdot v) = M \cdot (\hat H \cdot v)$ via
\texttt{Matrix.mulVec\_smul}.
\end{proof}

\begin{theorem}[\texttt{totalSpinHalfOpMinus\_mulVec\_mem\_magnetizationSubspace\_of\_mem}]
For any $M \in \mathbb{C}$, if $v \in \mathcal H_{M}$ then
$\hat S_{\mathrm{tot}}^{-} \cdot v \in \mathcal H_{M-1}$.  Dually,
$\hat S_{\mathrm{tot}}^{+} \cdot v \in \mathcal H_{M+1}$
(\texttt{totalSpinHalfOpPlus\_mulVec\_mem\_magnetizationSubspace\_of\_mem}).
This is the SU(2) ladder action on the magnetisation eigenspaces of
Tasaki eq.~(2.2.10).
\end{theorem}

\begin{proof}[Proof sketch]
Apply the Cartan relation $\hat S_{\mathrm{tot}}^{(3)} \cdot \hat
S_{\mathrm{tot}}^{-} = \hat S_{\mathrm{tot}}^{-} \cdot \hat
S_{\mathrm{tot}}^{(3)} - \hat S_{\mathrm{tot}}^{-}$ (rearranged from
\texttt{totalSpinHalfOp3\_commutator\_totalSpinHalfOpMinus} via
\texttt{abel}).  Then $\hat S_{\mathrm{tot}}^{(3)} \cdot
(\hat S_{\mathrm{tot}}^{-} \cdot v)
= \hat S_{\mathrm{tot}}^{-} \cdot \hat S_{\mathrm{tot}}^{(3)} \cdot v
- \hat S_{\mathrm{tot}}^{-} \cdot v
= \hat S_{\mathrm{tot}}^{-} \cdot M \cdot v
- \hat S_{\mathrm{tot}}^{-} \cdot v
= (M-1) \cdot (\hat S_{\mathrm{tot}}^{-} \cdot v)$ via
\texttt{Matrix.mulVec\_mulVec}, \texttt{Matrix.sub\_mulVec},
\texttt{Matrix.mulVec\_smul}, and \texttt{sub\_smul, one\_smul}.
The dual case for $\hat S_{\mathrm{tot}}^{+}$ replaces subtraction
with addition throughout.
\end{proof}

\subsection{Generic graph-centric Néel state and per-bond
expectations (Tasaki §2.5 (2.5.2)--(2.5.4), Phase 3 generic API)}
\label{subsec:neel-state-of}

The chain, 2D, and 3D parity-coloured Néel states share a common
structure: an alternating spin assignment from a Boolean
sublattice indicator, and a per-bond $\hat S_x \cdot \hat S_y$
action / expectation governed by whether the bond crosses the
sublattice partition.  Tasaki §2.5 abstracts this to any bipartite
graph; the Lean formalisation captures the abstraction
(graph-centric design, \texttt{CLAUDE.local.md}) via a single
generic $\texttt{neelStateOf}$ from a Boolean indicator
$A : V \to \texttt{Bool}$.

\begin{definition}[\texttt{neelConfigOf} / \texttt{neelStateOf}]
For a finite vertex type $V$ and sublattice indicator
$A : V \to \texttt{Bool}$,
\[
  \texttt{neelConfigOf}\,A\,x := \begin{cases}
    \uparrow\,(=0:\mathrm{Fin}\,2) & \text{if } A\,x \\
    \downarrow\,(=1:\mathrm{Fin}\,2) & \text{otherwise,}
  \end{cases}
\]
and $\texttt{neelStateOf}\,A := \texttt{basisVec}\,(\texttt{neelConfigOf}\,A)$.
The chain / 2D / 3D \texttt{neelXyzConfig} / \texttt{neelXyzState}
definitions are bridged via \texttt{neelXyzConfig\_eq\_neelConfigOf}
and \texttt{neelXyzState\_eq\_neelStateOf}.  Tasaki §2.5
eq.~(2.5.2) graph-centric form.
\end{definition}

\begin{theorem}[\texttt{spinHalfDot\_mulVec\_neelStateOf\_antiparallel},
Tasaki §2.5 eq.~(2.5.3)]
For any sublattice indicator $A$ and bond $(x, y)$ with $x \ne y$
and $A\,x \ne A\,y$,
\[
  \hat S_x \cdot \hat S_y \cdot \Phi_{\mathrm{N\acute eel}}(A)
    = \frac{1}{2}\,|\mathrm{swap}_{x,y}\,\Phi_{\mathrm{N\acute eel}}(A)\rangle
      - \frac{1}{4}\,\Phi_{\mathrm{N\acute eel}}(A).
\]
The 11 chain / 2D / 3D adjacent and wrap-around bond actions are
1-line corollaries via the \texttt{\_eq\_neelStateOf} bridges.
\end{theorem}

\begin{theorem}[\texttt{inner\_neelStateOf\_spinHalfDot\_neelStateOf\_antiparallel},
Tasaki §2.5 (2.5.4) ingredient]
Under the same hypotheses,
\[
  \langle \Phi_{\mathrm{N\acute eel}}(A),\,
      \hat S_x \cdot \hat S_y \cdot \Phi_{\mathrm{N\acute eel}}(A) \rangle
    = -\frac{1}{4}.
\]
The 11 chain / 2D / 3D \texttt{\_eq\_neg\_one\_quarter} per-bond
expectation lemmas reduce to this via the bridges.
\end{theorem}

\begin{theorem}[\texttt{inner\_neelStateOf\_szsz\_neelStateOf\_antiparallel}]
The diagonal $\hat S^z_x \cdot \hat S^z_y$ correlator on the
generic Néel state at any antiparallel bond:
\[
  \langle \Phi_{\mathrm{N\acute eel}}(A),\,
      \hat S^z_x \cdot \hat S^z_y \cdot \Phi_{\mathrm{N\acute eel}}(A) \rangle
    = -\frac{1}{4}.
\]
Proof: \texttt{inner\_basisVec\_onSite\_spinHalfOp3\_mul\_onSite\_spinHalfOp3\_basisVec}
returns $\texttt{spinHalfSign}\,(\sigma\,x)\cdot\texttt{spinHalfSign}\,(\sigma\,y)$,
which equals $-1/4$ at antiparallel $\sigma$ via
\texttt{spinHalfSign\_mul\_antiparallel}.
\end{theorem}

\begin{remark}[Marshall sign generalisation, PR \#329]
Analogously, the parity-coloured chain / 2D / 3D Marshall signs
are now specialisations of a single generic $\texttt{marshallSignOf}\,A
:= \prod_{x \in A} (-1)^{\sigma\,x}$ for any sublattice indicator
$A : V \to \texttt{Bool}$.  The all-up corollary
$\texttt{marshallSignOf}\,A\,(\mathrm{const}\,0) = 1$
(\texttt{marshallSignOf\_const\_zero}) collapses the chain / 2D / 3D
\texttt{\_const\_zero} variants to 1-line proofs through the
\texttt{\_eq\_marshallSignOf} bridges.  As of 2026-04-22 the chain
/ 2D / 3D \texttt{marshallSignChainConfig} / \texttt{\_square\_} /
\texttt{\_cubic\_} variants are marked
\texttt{@[deprecated]} in favour of the generic form (deprecation
window 6 months from \texttt{since}; see the project page
\href{https://phasetr.github.io/lattice-system/deprecations/}%
{\textsl{Deprecation window}} entry for the live tracking table
and removal-PR checklist).
\end{remark}

\subsection{Marshall-dressed standard basis (Tasaki §2.5 (2.5.8),
Issue \#412 PR $\alpha$-1)}
\label{ssec:marshall-dressed-basis}

\emph{Source file:}
\texttt{LatticeSystem/Quantum/MarshallDressedBasis.lean}.

For a finite vertex type $V$, a sublattice indicator
$A : V \to \texttt{Bool}$, and a spin-$1/2$ configuration
$\sigma : V \to \texttt{Fin}\,2$, the
\emph{Marshall-dressed standard basis state} is
\[
  |\widetilde{\Psi}^\sigma\rangle
    \;:=\; \texttt{marshallSignOf}\,A\,\sigma \cdot |\sigma\rangle,
\]
where $|\sigma\rangle = \texttt{basisVec}\,\sigma$ is the standard
computational basis vector and
$\texttt{marshallSignOf}\,A\,\sigma = \prod_{x \in A} (-1)^{\sigma_x}$
(\texttt{marshallDressedBasis}).  This is Tasaki §2.5 eq.~(2.5.8),
p.~41: the change of basis underlying the Marshall--Lieb--Mattis
theorem.

\begin{theorem}[Pointwise rules]
\label{thm:marshall-dressed-pointwise}
\[
  |\widetilde{\Psi}^\sigma\rangle\,\sigma
    = \texttt{marshallSignOf}\,A\,\sigma,
\qquad
  |\widetilde{\Psi}^\sigma\rangle\,\tau = 0 \quad (\tau \neq \sigma).
\]
Lean: \texttt{marshallDressedBasis\_self} and
\texttt{marshallDressedBasis\_of\_ne}.
\end{theorem}

\begin{theorem}[Marshall sign squares to one]
\label{thm:marshall-sign-sq}
$(\texttt{marshallSignOf}\,A\,\sigma)^2 = 1.$ Each factor is
$\pm 1$ since $(-1)^{\sigma_x}$ is $\pm 1$ for $\sigma_x \in
\{0, 1\}$.  Lean: \texttt{marshallSignOf\_sq\_eq\_one}.
\end{theorem}

\begin{theorem}[Orthonormality]
\label{thm:marshall-dressed-inner}
Under the real bilinear pairing,
\[
  \sum_\tau\, |\widetilde{\Psi}^\sigma\rangle\,\tau \cdot
    |\widetilde{\Psi}^\rho\rangle\,\tau
    \;=\; \begin{cases} 1 & \rho = \sigma, \\ 0 & \rho \neq \sigma. \end{cases}
\]
Combines orthonormality of $\texttt{basisVec}$
(\texttt{basisVec\_inner}) with
Theorem~\ref{thm:marshall-sign-sq}; no complex conjugation
appears because the Marshall sign and the standard basis values
are real.  Lean: \texttt{marshallDressedBasis\_inner}.
\end{theorem}

\begin{theorem}[$H_M$-membership]
\label{thm:marshall-dressed-hm}
$|\widetilde{\Psi}^\sigma\rangle$ lies in the same magnetisation
sector $H_M = H_{\bar\sigma/2}$ as $|\sigma\rangle$ (Tasaki
eq.~(2.2.10)), since the dressing is a scalar multiplication and
$H_M$ is a $\mathbb{C}$-submodule.  Lean:
\texttt{marshallDressedBasis\_mem\_magnetizationSubspace}; the
$\bar\sigma = 0$ specialisation places it in $H_0$
(\texttt{\_zero}).
\end{theorem}

This file is the first step of the Marshall--Lieb--Mattis
formalisation track (Issue \#412 = phase B-$\alpha$ for $S = 1/2$).
The next steps establish:
(i) realness of the dressed Heisenberg matrix elements
(see~\ref{ssec:mlm-realness});
(ii) non-positivity of the off-diagonal elements (the
Marshall sign trick proper);
(iii) connectivity in the Perron--Frobenius sense; and finally
the uniqueness and sign structure of the antiferromagnetic
Heisenberg ground state on a connected bipartite graph.

\subsection{Realness of the dressed Heisenberg matrix elements
(Tasaki §2.5 p. 41 Property (i), Issue \#412 PR $\alpha$-2)}
\label{ssec:mlm-realness}

\emph{Source file:}
\texttt{LatticeSystem/Quantum/MarshallLiebMattis/Realness.lean}.

Property (i) in Tasaki's Perron--Frobenius proof of the
Marshall--Lieb--Mattis theorem is
$\langle \widetilde{\Psi}^\sigma | \hat H |
 \widetilde{\Psi}^{\sigma'}\rangle \in \mathbb{R}$
for the antiferromagnetic Heisenberg Hamiltonian
$\hat H = \sum_{x,y} J_{x,y}\,\hat S_x \cdot \hat S_y$ with real
coupling $J$ (formalised as $(J\,x\,y).{\rm im} = 0$).

\begin{theorem}[Two-site spin product matrix entries are real]
\label{thm:spinHalfDot-real}
For every $x, y \in \Lambda$ and every $\sigma, \sigma'$,
$\big((\hat S_x \cdot \hat S_y)_{\sigma,\sigma'}\big).{\rm im} = 0$.
By case analysis on $x = y$ / parallel / antiparallel via the
\texttt{spinHalfDot\_mulVec\_basisVec\_\{parallel,antiparallel\}}
action lemmas, the entry is in $\{3/4, 1/4, -1/4, 1/2, 0\}$.
Lean: \texttt{spinHalfDot\_apply\_im\_eq\_zero}.
\end{theorem}

\begin{theorem}[Heisenberg matrix entries real for real $J$]
\label{thm:heisenberg-real}
If $(J\,x\,y).{\rm im} = 0$ for all $x, y$, then
$\big((H_J)_{\sigma,\sigma'}\big).{\rm im} = 0$.  Proof:
$\mathbb{R}$-linearity of $\mathrm{Complex.im}$ and
Theorem~\ref{thm:spinHalfDot-real}.  Lean:
\texttt{heisenbergHamiltonian\_apply\_im\_eq\_zero}.
\end{theorem}

\begin{theorem}[Marshall sign is real]
\label{thm:marshall-sign-real}
$(\texttt{marshallSignOf}\,A\,\sigma).{\rm im} = 0$.  Each factor
is $\pm 1 \in \mathbb{R}$.  Lean:
\texttt{marshallSignOf\_im\_eq\_zero}.
\end{theorem}

\begin{theorem}[MLM Property (i)]
\label{thm:dressed-bilinear-real}
For real $J$, the dressed bilinear pairing
$\sum_\tau |\widetilde{\Psi}^\sigma\rangle\,\tau \cdot
 (H_J |\widetilde{\Psi}^{\sigma'}\rangle)\,\tau$
has zero imaginary part.  The pairing reduces to
$\texttt{marshallSignOf}\,A\,\sigma \cdot
 \texttt{marshallSignOf}\,A\,\sigma' \cdot (H_J)_{\sigma,\sigma'}$,
a product of three reals.  Lean:
\texttt{dot\_marshallDressed\_heisenbergHamiltonian\_marshallDressed\_im\_eq\_zero}.
\end{theorem}

\subsection{The Marshall sign trick: non-positive off-diagonal
dressed elements (Tasaki §2.5 p. 41 Property (ii), Issue \#412
PR $\alpha$-3)}
\label{ssec:mlm-marshall-sign-trick}

\emph{Source file:}
\texttt{LatticeSystem/Quantum/MarshallLiebMattis/MarshallSignTrick.lean}.

The second of three properties needed in Tasaki's
Perron--Frobenius proof of Theorem 2.2 is

\begin{quote}
  \textbf{(ii)}
  $\langle \widetilde{\Psi}^\sigma | \hat H |
   \widetilde{\Psi}^{\sigma'}\rangle \le 0$
  for $\sigma \neq \sigma'$
\end{quote}

for the spin-$1/2$ antiferromagnetic Heisenberg Hamiltonian on a
bipartite graph $(\Lambda, A)$ with real non-negative coupling
$J$ supported on bipartite bonds (i.e. $J_{x,y} = 0$ whenever
$A\,x = A\,y$).

\begin{theorem}[Pairing factorisation]
\label{thm:dressed-pairing-factor}
For any operator $M$,
$\sum_\tau |\widetilde{\Psi}^\sigma\rangle\,\tau \cdot
 (M |\widetilde{\Psi}^{\sigma'}\rangle)\,\tau =
 \texttt{marshallSignOf}\,A\,\sigma \cdot
 \texttt{marshallSignOf}\,A\,\sigma' \cdot M_{\sigma,\sigma'}$.
Lean: \texttt{dot\_marshallDressed\_mulVec\_marshallDressed\_eq}.
\end{theorem}

\begin{theorem}[Marshall sign relation]
\label{thm:marshall-swap-sign}
For a bond $\{x, y\}$ with $A\,x \neq A\,y$ and a configuration
$\sigma$ antiparallel at $\{x, y\}$ ($\sigma\,x \neq \sigma\,y$),
$\texttt{marshallSignOf}\,A\,\sigma \cdot
 \texttt{marshallSignOf}\,A\,(\texttt{basisSwap}\,\sigma\,x\,y) = -1$.
Outside $\{x, y\}$ each pairwise factor squares to $1$. At the
unique site in $A \cap \{x, y\}$ the pair contributes
$(-1)^{\sigma_x + \sigma_y} = -1$ since $\sigma_x \neq \sigma_y$
in $\texttt{Fin}\,2$. The other site of $\{x, y\}$ lies outside
$A$ and contributes $1$. Lean:
\texttt{marshallSignOf\_mul\_marshallSignOf\_basisSwap\_of\_bipartite\_antiparallel}.
\end{theorem}

\begin{theorem}[Per-bond non-positivity]
\label{thm:bond-dressed-nonpos}
For any bond $(x, y)$, real non-negative $J_{x,y}$ supported on
bipartite bonds, and $\sigma \neq \sigma'$, the contribution
$\texttt{marshallSignOf}\,A\,\sigma \cdot
 \texttt{marshallSignOf}\,A\,\sigma' \cdot J_{x,y} \cdot
 (\hat S_x \cdot \hat S_y)_{\sigma,\sigma'}$
has non-positive real part. By case analysis on
$(\hat S_x \cdot \hat S_y)_{\sigma,\sigma'}$: it is zero except
at $x \neq y$, $\sigma'_x \neq \sigma'_y$, and
$\sigma = \texttt{basisSwap}\,\sigma'\,x\,y$, where it equals
$1/2$. In that case, if $A\,x = A\,y$ the bipartite hypothesis
forces $J_{x,y} = 0$; if $A\,x \neq A\,y$ the Marshall sign
relation (Theorem~\ref{thm:marshall-swap-sign}) gives sign $-1$,
yielding $-(1/2)\,(J_{x,y}).\mathrm{re} \le 0$. Lean:
\texttt{bond\_dressed\_contribution\_re\_nonpos}.
\end{theorem}

\begin{theorem}[MLM Property (ii)]
\label{thm:dressed-offdiagonal-nonpos}
For real non-negative $J$ supported on bipartite bonds and
$\sigma \neq \sigma'$, the dressed off-diagonal Heisenberg
pairing has non-positive real part:
\[
  \mathrm{Re}\,\sum_\tau\,
    |\widetilde{\Psi}^\sigma\rangle\,\tau \cdot
    (H_J\,|\widetilde{\Psi}^{\sigma'}\rangle)\,\tau \;\le\; 0.
\]
Sum the bond contributions using
Theorem~\ref{thm:bond-dressed-nonpos}. Lean:
\texttt{dot\_marshallDressed\_heisenbergHamiltonian\_marshallDressed\_re\_nonpos\_of\_ne}.
\end{theorem}

\subsection{Swap-connectivity (Tasaki §2.5 p. 41--42 Property (iii),
Issue \#412 PR $\alpha$-4)}
\label{ssec:mlm-connectivity}

\emph{Source file:}
\texttt{LatticeSystem/Quantum/MarshallLiebMattis/Connectivity.lean}.

The third property in Tasaki's Perron--Frobenius proof of
Theorem 2.2 is connectivity: any two configurations
$\sigma, \sigma'$ with the same total magnetisation are connected
by a sequence of nonvanishing matrix elements of
$\langle\widetilde{\Psi}^\sigma|\hat H|\widetilde{\Psi}^{\sigma'}\rangle$.
For a connected graph $G$ this reduces to the purely combinatorial
statement that any pair $(x, y)$ with $\sigma_x \neq \sigma_y$
admits a $G$-edge swap chain producing
$\texttt{basisSwap}\,\sigma\,x\,y$ from $\sigma$.

\begin{definition}[\texttt{SwapStep}, \texttt{SwapReachable}]
\label{def:swap-step}
$\texttt{SwapStep}\,G\,\sigma\,\sigma'$ holds iff
$\sigma' = \texttt{basisSwap}\,\sigma\,x\,y$ for some $G$-adjacent
$(x, y)$ with $\sigma_x \neq \sigma_y$.
$\texttt{SwapReachable}\,G$ is the reflexive-transitive closure.
Lean: \texttt{SwapStep}, \texttt{SwapReachable}.
\end{definition}

\begin{theorem}[Walk-induction Lemma, Tasaki p. 41]
\label{thm:swap-walk}
For any $G$-walk from $x$ to $y$ and any configuration
$\sigma$ with $\sigma_x \neq \sigma_y$,
$\texttt{SwapReachable}\,G\,\sigma\,(\texttt{basisSwap}\,\sigma\,x\,y)$.

By induction on the walk. The cons case
$x \to v \to \cdots \to y'$ (after the head edge $(x, v)$) splits
on $\sigma_v$:
\begin{itemize}
\item Case A ($\sigma_v \neq \sigma_x$, so $\sigma_v = \sigma_y$):
  swap $(x, v)$ first via the head edge; the resulting
  configuration $\sigma_1$ has $\sigma_1{}_v = \sigma_x \neq
  \sigma_y = \sigma_1{}_y$, so the inductive hypothesis on the
  remaining walk applies.
\item Case B ($\sigma_v = \sigma_x$): apply the inductive
  hypothesis on the remaining walk first, then swap $(x, v)$
  via the head edge on the resulting configuration.
\end{itemize}
Both cases produce a configuration equal to
$\texttt{basisSwap}\,\sigma\,x\,y$ off the support
$\{x, v, y\}$ and at the three special sites by direct check.
Lean: \texttt{swapReachable\_of\_walk\_of\_ne}.
\end{theorem}

\begin{theorem}[MLM Property (iii) ingredient]
\label{thm:swap-reachable}
For a connected (or merely preconnected) graph $G$ and any
$\sigma$ with $\sigma_x \neq \sigma_y$,
$\texttt{SwapReachable}\,G\,\sigma\,(\texttt{basisSwap}\,\sigma\,x\,y)$.
Combined with iteration over the magnetisation distance
$\sum_x |\sigma_x - \sigma'_x|$, this gives Perron--Frobenius
irreducibility of the dressed Heisenberg matrix on the
magnetisation-$M$ subspace. Lean:
\texttt{swapReachable\_of\_reachable\_of\_ne},
\texttt{swapReachable\_of\_preconnected\_of\_ne}.
\end{theorem}

\subsection{Real-valued dressed Heisenberg matrix on $H_0$ (Tasaki
§2.5 setup, Issue \#412 PR $\alpha$-5a)}
\label{ssec:mlm-h0-matrix}

\emph{Source file:}
\texttt{LatticeSystem/Quantum/MarshallLiebMattis/H0Matrix.lean}.

To apply Perron--Frobenius to the dressed Heisenberg Hamiltonian
on the zero-magnetisation subspace $H_0$, we assemble its
real-valued matrix. The index type is
$\texttt{H}_0\texttt{Index}\,\Lambda :=
 \{\sigma : \Lambda \to \texttt{Fin}\,2 \mid
   \texttt{magnetization}\,\Lambda\,\sigma = 0\}$, and the matrix is
\[
  M_{\sigma, \tau}
   := \mathrm{Re}\big(\texttt{marshallSignOf}\,A\,\sigma \cdot
                       \texttt{marshallSignOf}\,A\,\tau \cdot
                       (H_J)_{\sigma, \tau}\big)
\]
(\texttt{dressedHeisenbergMatrixH0}).

\begin{theorem}[Symmetry]
\label{thm:dressed-matrix-symm}
For real symmetric $J$ (i.e.\ $(J\,x\,y).{\rm im} = 0$ and
$J\,x\,y = J\,y\,x$), the matrix $M$ is symmetric:
$M_{\sigma, \tau} = M_{\tau, \sigma}$.
Hermiticity of the Heisenberg Hamiltonian
(\texttt{heisenbergHamiltonian\_isHermitian\_of\_real\_symm})
combined with the realness of matrix entries
(Theorem~\ref{thm:heisenberg-real}) implies the entries are
real symmetric. Lean:
\texttt{dressedHeisenbergMatrixH0\_isSymm}.
\end{theorem}

\begin{theorem}[Off-diagonal non-positivity]
\label{thm:dressed-matrix-offdiag}
For real non-negative $J$ supported on bipartite bonds and
distinct $\sigma \neq \tau$, $M_{\sigma, \tau} \le 0$.
This packages Theorem~\ref{thm:dressed-offdiagonal-nonpos} via
the factorisation
Theorem~\ref{thm:dressed-pairing-factor}. Lean:
\texttt{dressedHeisenbergMatrixH0\_offdiag\_nonpos}.
\end{theorem}

The remaining work for $H_0$ uniqueness (subsequent PRs) is to
shift $M$ to a non-negative form $cI - M$, prove its
irreducibility from Theorem~\ref{thm:swap-reachable}, and apply
Perron--Frobenius to extract the unique strictly positive
eigenvector that yields Tasaki (2.5.4).

\subsection{Magnetisation preservation under \texttt{basisSwap}
(Tasaki §2.5 p. 42 ingredient, Issue \#412 PR $\alpha$-5b)}
\label{ssec:mlm-magnetization-swap}

\emph{Source file:}
\texttt{LatticeSystem/Quantum/MarshallLiebMattis/EqMagnetization.lean}.

\begin{theorem}[Magnetisation preservation under bond swap]
\label{thm:mag-swap}
$\texttt{magnetization}\,\Lambda\,(\texttt{basisSwap}\,\sigma\,x\,y) =
 \texttt{magnetization}\,\Lambda\,\sigma$.

The bond swap is the composition with the transposition permutation
$\texttt{Equiv.swap}\,x\,y : \Lambda \simeq \Lambda$:
$\texttt{basisSwap}\,\sigma\,x\,y = \sigma \circ \texttt{Equiv.swap}\,x\,y$.
Total magnetisation $\sum_z \texttt{spinSign}\,(\sigma\,z)$ is
invariant under permutation reindexing of the underlying lattice
(\texttt{Equiv.sum\_comp}). Lean: \texttt{magnetization\_basisSwap},
with $H_0$-membership corollary
\texttt{magnetization\_basisSwap\_eq\_zero\_iff}.
\end{theorem}

This is the key ingredient for the Tasaki §2.5 p. 42 Proposition
(any two configurations with equal magnetisation are connected via
a chain of single-edge bond swaps), proved next.

\subsection{Tasaki §2.5 p. 42 Proposition: equal-magnetisation
reachability (Issue \#412 PR $\alpha$-5c)}
\label{ssec:mlm-eq-magnetization-reachable}

\emph{Source file:}
\texttt{LatticeSystem/Quantum/MarshallLiebMattis/EqMagnetizationReachable.lean}.

\begin{theorem}[Tasaki §2.5 p. 42 Proposition]
\label{thm:eq-mag-reachable}
For a connected graph $G$ and any two configurations
$\sigma, \sigma'$ with the same total magnetisation
$\bar\sigma = \bar\sigma'$,
$\texttt{SwapReachable}\,G\,\sigma\,\sigma'$.

Strong induction on the **configuration distance**
$\mathrm{configDist}(\sigma, \sigma') :=
 \#\{x \in \Lambda \mid \sigma_x \neq \sigma'_x\}$:
\begin{itemize}
\item Base $\mathrm{configDist} = 0$: $\sigma = \sigma'$, refl.
\item Inductive step: pigeonhole on the magnetisation equality
gives sites $x$ with $(\sigma_x, \sigma'_x) = (0, 1)$ and $y$ with
$(\sigma_y, \sigma'_y) = (1, 0)$
(\texttt{exists\_swap\_pair\_of\_eq\_magnetization}); apply
Theorem~\ref{thm:swap-reachable} to get $\sigma \to
\texttt{basisSwap}\,\sigma\,x\,y$; the new configuration has
$\mathrm{configDist}$ strictly smaller (sites $x, y$ newly agree
with $\sigma'$, see \texttt{configDist\_basisSwap\_lt}) and the
same magnetisation as $\sigma'$
(Theorem~\ref{thm:mag-swap}); apply IH.
\end{itemize}
Lean: \texttt{swapReachable\_of\_eq\_magnetization}.
\end{theorem}

This completes the combinatorial content needed for
Perron--Frobenius irreducibility of the dressed Heisenberg matrix
on the magnetisation-$M$ subspace. Subsequent PRs will:
\begin{itemize}
\item upgrade $\mathrm{configDist}$ reachability to a positive
matrix-power entry $(B^n)_{\sigma, \sigma'} > 0$ for the shifted
non-negative matrix $B = cI - M$ defined below;
\item apply Perron--Frobenius
(\texttt{LatticeSystem.Math.PerronFrobenius.exists\_pos\_eigenvec\_max}
and \texttt{pos\_eigenvec\_unique}) to extract the unique
strictly-positive eigenvector;
\item translate back to the antiferromagnetic Heisenberg
ground-state expansion
$|\Phi_{\rm GS}\rangle = \sum_\sigma c_\sigma |\widetilde{\Psi}^\sigma\rangle$
with $c_\sigma > 0$ (Tasaki (2.5.4)).
\end{itemize}

\subsection{Shifted dressed Heisenberg matrix $B = cI - M$ on
$H_0$ (Issue \#412 PR $\alpha$-5d)}
\label{ssec:mlm-h0-shifted}

\emph{Source file:}
\texttt{LatticeSystem/Quantum/MarshallLiebMattis/H0Shifted.lean}.

To apply Perron--Frobenius, the dressed Heisenberg matrix $M$ must
be transformed into a non-negative form. The $cI - M$ shift
preserves the eigenvector structure (eigenvalue $\mu$ of $M$ becomes
$c - \mu$ of $B$ with the same eigenvector), and translates the
\emph{minimum} eigenvalue of $M$ (= ground state of the restricted
Heisenberg) into the \emph{maximum} eigenvalue of $B$ (which is
where Perron--Frobenius gives a unique positive eigenvector).

\begin{definition}[Shifted matrix]
\label{def:shifted-matrix}
$B := \texttt{dressedHeisenbergShifted}\,A\,J\,c$, with
$B_{\sigma, \tau} := (c \cdot I)_{\sigma, \tau} - M_{\sigma, \tau}$.
Lean: \texttt{dressedHeisenbergShifted}.
\end{definition}

\begin{theorem}[Symmetry]
\label{thm:shifted-symm}
For real symmetric $J$ and any $c \in \mathbb{R}$, $B$ is symmetric.
Inherits symmetry from
Theorem~\ref{thm:dressed-matrix-symm}; the $cI$ term is trivially
symmetric. Lean: \texttt{dressedHeisenbergShifted\_isSymm}.
\end{theorem}

\begin{theorem}[Non-negativity]
\label{thm:shifted-nonneg}
For real non-negative $J$ on bipartite bonds:
\begin{itemize}
\item Off-diagonal entries of $B$ are non-negative (sign flip of
$M$'s non-positive off-diagonal,
Theorem~\ref{thm:dressed-matrix-offdiag});
\item Diagonal entries are non-negative when $c \ge M_{\sigma, \sigma}$
for every $\sigma$.
\end{itemize}
Lean: \texttt{dressedHeisenbergShifted\_offdiag\_nonneg},
\texttt{dressedHeisenbergShifted\_diag\_nonneg},
\texttt{dressedHeisenbergShifted\_nonneg}.
\end{theorem}

%======================================================================
\section{Perron--Frobenius theorem for symmetric non-negative irreducible matrices}
\label{sec:perron-frobenius}
%======================================================================

\emph{Source files:}
\texttt{LatticeSystem/Math/PerronFrobeniusPrimitive.lean},
\texttt{LatticeSystem/Math/CollatzWielandt.lean},
\texttt{LatticeSystem/Math/PerronFrobeniusMain.lean},
\texttt{LatticeSystem/Math/PerronFrobenius.lean}.

\medskip
\textbf{Setting.}
Let $n$ be a finite type and $A : \mathrm{Matrix}(n,n,\mathbb{R})$ be:
\begin{itemize}
  \item \emph{symmetric} (\texttt{A.IsHermitian}),
  \item \emph{non-negative}: $A_{ij} \geq 0$ for all $i,j$,
  \item \emph{irreducible} (\texttt{A.IsIrreducible}): for every $i,j$
    there exists $k>0$ with $(A^k)_{ij} > 0$.
\end{itemize}

\subsection{Non-negative max eigenvector (\texttt{exists\_nonneg\_eigenvec\_max})}

\textbf{Statement (sorry).}
There exists a maximum eigenvalue $\mu$ and a non-zero non-negative eigenvector
$v : n \to \mathbb{R}$ satisfying $Av = \mu v$, $v \neq 0$,
$v_i \geq 0$ for all $i$, and $\lambda_k \leq \mu$ for every eigenvalue $\lambda_k$.

\textbf{Mathematical proof sketch.}
Let $v$ be any max eigenvector (from the spectral theorem) and set $w_i = |v_i|$.
Since $A_{ij} \geq 0$,
\[
  (Aw)_i = \sum_j A_{ij}|v_j| \geq \left|\sum_j A_{ij}v_j\right| = |\mu v_i| = \mu|v_i|,
\]
so $w \cdot (Aw) \geq \mu\|w\|^2$.
The Rayleigh quotient bound gives $w \cdot (Aw) \leq \mu\|w\|^2$, hence equality.
By the eigenbasis expansion $w \cdot (Aw) = \sum_k \lambda_k c_k^2$ and
$\|w\|^2 = \sum_k c_k^2$ where $c_k = \langle e_k, w\rangle$,
equality forces $c_k = 0$ for $\lambda_k < \mu$, giving $Aw = \mu w$.

\textbf{Lean status.}
The proof is \texttt{sorry}: the Rayleigh equality condition
$\langle w, Aw\rangle = \mu\|w\|^2 \Rightarrow Aw = \mu w$
requires computing $\langle w, Aw\rangle$ via the eigenbasis inner product,
which is blocked by the \texttt{EuclideanSpace}/\texttt{n $\to$ $\mathbb{R}$} API gap in Mathlib.
Tracked in Issue~\#405.

\subsection{Collatz--Wielandt main results (\texttt{PerronFrobeniusMain})}
\label{sec:pf-main}

\emph{Source file:} \texttt{LatticeSystem/Math/PerronFrobeniusMain.lean}.

\medskip
The key results eliminating the sorry from \texttt{exists\_pos\_eigenvec\_max}.

\begin{itemize}
\item \textbf{\texttt{pos\_of\_nonneg\_eigenvec}.}
  If $A$ is irreducible nonneg, $Av = \mu v$, $v \geq 0$, $v \neq 0$, then $v > 0$.
  Proof: $v_i = 0$ forces all $v_j = 0$ via irreducibility.

\item \textbf{\texttt{exists\_positive\_eigenvector\_of\_primitive}.}
  For a primitive nonneg matrix: $\exists\,r>0$, $v>0$, $Av = rv$.
  Proof: (1) CW maximizer $v$ on $\Delta_n$ satisfies $r_0 \cdot v \leq Av$ (fundamental ineq).
  (2) If $z = Av - r_0 v \neq 0$, primitivity gives $A^k z > 0$, so all ratios
  $(A^{k+1}v)_i/(A^k v)_i > r_0$; normalising $A^k v$ contradicts maximality.
  (3) Hence $Av = r_0 v$. Positivity follows from irreducibility; $r_0 > 0$
  from $A^k v = r_0^k v > 0$.
  Ref: Seneta \S1.2 (pp.~27--28); MCMC \texttt{Primitive.lean}.

\item \textbf{\texttt{exists\_positive\_eigenvector\_of\_irreducible}.}
  For an irreducible nonneg matrix: $\exists\,r>0$, $v>0$, $Av = rv$.
  Proof: Set $B = I + A$; it is primitive by \texttt{IsPrimitive.of\_irreducible\_pos\_diagonal}.
  Apply the primitive theorem to $B$: get $r_B > 0$, $v > 0$, $Bv = r_B v$.
  Then $Av = (r_B-1)v$. Since $B \geq I$, $r_B \geq 1$.
  If $r_B = 1$ then $Av = 0$, so $A = 0$, contradicting irreducibility.
  Hence $r_B > 1$.
  Ref: Seneta \S1.2 (pp.~27--28); MCMC \texttt{Irreducible.lean}.
\end{itemize}

\subsection{Strictly positive max eigenvector (\texttt{exists\_pos\_eigenvec\_max})}

\textbf{Statement.}
Under the above hypotheses,
$\exists\,\mu\,\exists\,v,\; Av = \mu v \;\land\; v \neq 0 \;\land\; \forall i,\; v_i > 0$.

\textbf{Proof.}
Call \texttt{exists\_positive\_eigenvector\_of\_irreducible} directly.
The sorry-bearing \texttt{exists\_nonneg\_eigenvec\_max} is bypassed entirely.
(The sorry is retained in that function for documentation of the unresolved
Rayleigh equality approach.)

\subsection{Uniqueness (\texttt{pos\_eigenvec\_unique})}

\textbf{Statement.}
If $Av = \mu v$ and $Aw = \mu w$ with $v_i > 0$ and $w_i > 0$ for all $i$,
then $\exists\, r > 0,\; w = r\cdot v$.

\textbf{Proof.}
Set $r = \sup_i(w_i/v_i)$ (via \texttt{Finset.sup'}).
Then $w \leq rv$ componentwise.
Let $u = rv - w \geq 0$; since $A(rv-w) = \mu(rv-w)$, $u$ is a non-negative
$\mu$-eigenvector.
The maximizer $i_0$ satisfies $r = w_{i_0}/v_{i_0}$, so $u_{i_0} = 0$.
By the propagation argument (same as the strict-positivity lemma), $u = 0$, hence $w = rv$.

\begin{center}
\begin{tabular}{lll}
\hline
Lean name & Status & Notes \\
\hline
\texttt{IsPrimitive.of\_irreducible\_pos\_diagonal} & proved & PR A \\
\texttt{CollatzWielandt.*} & proved & PR B \\
\texttt{pos\_of\_nonneg\_eigenvec} & proved & PR C \\
\texttt{exists\_positive\_eigenvector\_of\_primitive} & proved & PR C; Seneta \S1.2 \\
\texttt{exists\_positive\_eigenvector\_of\_irreducible} & proved & PR C; Seneta \S1.2 \\
\texttt{exists\_nonneg\_eigenvec\_max} & sorry & Rayleigh step; retained for docs \\
\texttt{exists\_pos\_eigenvec\_max} & \textbf{sorry-free} & bypasses sorry via PR C \\
\texttt{pos\_eigenvec\_unique} & proved & $r = \sup_i w_i/v_i$ \\
\hline
\end{tabular}
\end{center}

%======================================================================
\section{Single-mode fermion (Phase 2 skeleton)}
\label{sec:fermion-mode}
%======================================================================

\emph{Source file:} \texttt{LatticeSystem/Fermion/Mode.lean}.

\emph{Convention.}  A single fermion mode acts on the two-dimensional
Hilbert space $\mathbb{C}^2$ with computational basis
$|0\rangle = e_0$ (vacuum) and $|1\rangle = e_1$ (occupied).

\subsection{Definitions}

\begin{definition}[\texttt{fermionAnnihilation}, \texttt{fermionCreation}, \texttt{fermionNumber}]
\[
  c \;=\; |0\rangle\langle 1| \;=\;
  \begin{pmatrix} 0 & 1 \\ 0 & 0 \end{pmatrix},
  \quad
  c^\dagger \;=\; |1\rangle\langle 0| \;=\;
  \begin{pmatrix} 0 & 0 \\ 1 & 0 \end{pmatrix},
  \quad
  n \;=\; c^\dagger c \;=\;
  \begin{pmatrix} 0 & 0 \\ 0 & 1 \end{pmatrix}.
\]
The number-operator decomposition $n = c^\dagger c$ is recorded
explicitly as \texttt{fermionNumber\_eq\_creation\_mul\_annihilation}.
\end{definition}

\subsection{CAR relations}

\begin{lemma}[\texttt{fermionAnnihilation\_sq}, \texttt{fermionCreation\_sq}]
$c^2 = 0$ and $(c^\dagger)^2 = 0$.
\end{lemma}

\begin{lemma}[\texttt{fermionAnticomm\_self}]
The single-mode canonical anticommutation relation:
\[
  \{c, c^\dagger\} \;=\; c\,c^\dagger + c^\dagger\,c \;=\; I.
\]
\end{lemma}

\begin{proof}[Proof sketch]
Direct $2 \times 2$ matrix computation; \texttt{fin\_cases} on row /
column followed by \texttt{simp}.
\end{proof}

\begin{lemma}[\texttt{fermionNumber\_sq}]
The number operator is idempotent: $n^2 = n$ (eigenvalues $0$ and $1$).
\end{lemma}

\subsection{Hermiticity}

\begin{lemma}[\texttt{fermionAnnihilation\_conjTranspose}, \texttt{fermionCreation\_conjTranspose}, \texttt{fermionNumber\_isHermitian}]
$c^\dagger = (c)^\dagger$, $c = (c^\dagger)^\dagger$, and $n$ is
Hermitian.
\end{lemma}

\subsection{Basis vectors and basis action}

\begin{definition}[\texttt{fermionVacuum}, \texttt{fermionOccupied}]
$|0\rangle = (1, 0)$ and $|1\rangle = (0, 1)$.
\end{definition}

\begin{lemma}[Basis action]
\begin{align*}
  c\,|0\rangle &= 0, & c\,|1\rangle &= |0\rangle, \\
  c^\dagger\,|0\rangle &= |1\rangle, & c^\dagger\,|1\rangle &= 0, \\
  n\,|0\rangle &= 0, & n\,|1\rangle &= |1\rangle.
\end{align*}
Lean names: \texttt{fermionAnnihilation\_mulVec\_vacuum} /
\texttt{\_occupied},
\texttt{fermionCreation\_mulVec\_vacuum} / \texttt{\_occupied},
\texttt{fermionNumber\_mulVec\_vacuum} / \texttt{\_occupied}.
\end{lemma}

\subsection{Identification with spin-1/2 raising / lowering operators}

In the computational basis,
\[
  c \;=\; \sigma^+,
  \qquad
  c^\dagger \;=\; \sigma^-.
\]
Lean names: \texttt{fermionAnnihilation\_eq\_spinHalfOpPlus},
\texttt{fermionCreation\_eq\_spinHalfOpMinus}.

%======================================================================
\section{Multi-mode fermion via Jordan–Wigner (Phase 2 backbone)}
\label{sec:fermion-jw}
%======================================================================

\emph{Source file:} \texttt{LatticeSystem/Fermion/JordanWigner.lean}.

\subsection{Definitions}

\begin{definition}[\texttt{jwString}]
The Jordan–Wigner string at site $i \in \mathrm{Fin}\,(N+1)$ is
\[
  J_i \;:=\; \prod_{\substack{j : \mathrm{Fin}\,(N+1) \\ j.\mathrm{val} < i.\mathrm{val}}}
              \sigma^z_j,
\]
implemented as a \texttt{Finset.noncommProd} (because
$\texttt{ManyBodyOp}\,\Lambda$ is a non-commutative ring), with the
pairwise-commutativity certificate supplied by
\texttt{onSite\_mul\_onSite\_of\_ne}.
\end{definition}

\begin{definition}[\texttt{fermionMultiAnnihilation}, \texttt{fermionMultiCreation}]
\[
  c_i \;:=\; J_i \cdot \sigma^+_i,
  \qquad
  c_i^\dagger \;:=\; J_i \cdot \sigma^-_i.
\]
\end{definition}

\subsection{Boundary identifications}

\begin{lemma}[\texttt{jwString\_zero}]
$J_0 = I$ (empty product at the leftmost site).
\end{lemma}

\begin{lemma}[\texttt{fermionMultiAnnihilation\_zero}, \texttt{fermionMultiCreation\_zero}]
$c_0 = \sigma^+_0$ and $c_0^\dagger = \sigma^-_0$ (the JW string is
the identity at the leftmost site).
\end{lemma}

\subsection{On-site CAR vanishing (Pauli exclusion)}

\begin{lemma}[\texttt{jwString\_commute\_onSite}]
For any $A \in M_2(\mathbb{C})$, $[J_i, \texttt{onSite } i\, A] = 0$.
\end{lemma}

\begin{proof}[Proof sketch]
Each factor $\sigma^z_j$ (with $j.\mathrm{val} < i.\mathrm{val}$,
hence $j \neq i$) commutes with the site-$i$ operator via
\texttt{onSite\_mul\_onSite\_of\_ne}, and pairwise commutativity lifts
to the noncomm-product via \texttt{Finset.noncommProd\_commute}.
\end{proof}

\begin{lemma}[\texttt{fermionMultiAnnihilation\_sq}, \texttt{fermionMultiCreation\_sq}]
$c_i^2 = 0$ and $(c_i^\dagger)^2 = 0$.
\end{lemma}

\begin{proof}[Proof sketch]
$(J_i \sigma^+_i)(J_i \sigma^+_i)
 = J_i (\sigma^+_i J_i) \sigma^+_i
 = J_i (J_i \sigma^+_i) \sigma^+_i
 = J_i^2 (\sigma^+_i \sigma^+_i)$,
where the middle commutation uses the previous lemma.  The single-site
square $\sigma^+ \cdot \sigma^+ = 0$ then gives
$J_i^2 \cdot \texttt{onSite } i\, 0 = J_i^2 \cdot 0 = 0$.  Same
argument for $(c_i^\dagger)^2$ with $\sigma^-$ in place of $\sigma^+$.
\end{proof}

\subsection{JW string Hermiticity and adjoint of multi-mode operators}

\begin{lemma}[\texttt{jwString\_isHermitian}]
$J_i$ is Hermitian.
\end{lemma}

\begin{proof}[Proof sketch]
The string is a noncomm-product of pairwise-commuting Hermitian
$\sigma^z_j$ factors.  An induction over the Finset, using
\texttt{Matrix.IsHermitian.mul\_of\_commute} at each step (the new
factor commutes with the rest by
\texttt{Finset.noncommProd\_commute}), establishes the result.
\end{proof}

\begin{lemma}[\texttt{fermionMultiAnnihilation\_conjTranspose},
            \texttt{fermionMultiCreation\_conjTranspose}]
$(c_i)^\dagger = c_i^\dagger$ and $(c_i^\dagger)^\dagger = c_i$.
\end{lemma}

\begin{proof}[Proof sketch]
$(J_i \sigma^+_i)^\dagger = (\sigma^+_i)^\dagger \cdot J_i^\dagger
 = \sigma^-_i \cdot J_i = J_i \cdot \sigma^-_i = c_i^\dagger$,
using \texttt{Matrix.conjTranspose\_mul},
\texttt{jwString\_isHermitian},
\texttt{onSite\_conjTranspose} (a small private lemma that lifts
conjTranspose into \texttt{onSite}),
\texttt{spinHalfOpPlus\_conjTranspose}, and finally
\texttt{jwString\_commute\_onSite} to swap the order back.  The
$(c_i^\dagger)^\dagger = c_i$ direction is symmetric.
\end{proof}

\subsection{Site-occupation number operator}

\begin{lemma}[\texttt{jwString\_sq}]
$J_i^2 = I$.
\end{lemma}

\begin{proof}[Proof sketch]
Each factor $\sigma^z_j$ satisfies $(\sigma^z)^2 = I$ at the
single-site level, hence
$(\texttt{onSite}\,j\,\sigma^z)^2 = \texttt{onSite}\,j\,((\sigma^z)^2) = \texttt{onSite}\,j\,I = I$
via \texttt{onSite\_mul\_onSite\_same} and the auxiliary
\texttt{onSite\_one}.  An induction on the noncomm-product (using
the same commutativity/factor-cancellation pattern as the Hermiticity
proof) extends to the full string.
\end{proof}

\begin{definition}[\texttt{fermionMultiNumber}]
The site-occupation number operator is $n_i := c_i^\dagger \cdot c_i$.
\end{definition}

\begin{lemma}[\texttt{fermionMultiNumber\_eq\_onSite}]
$n_i = \texttt{onSite}\,i\,(\sigma^- \cdot \sigma^+)$.
\end{lemma}

\begin{proof}[Proof sketch]
$(J_i \sigma^-)(J_i \sigma^+)
 = J_i (\sigma^- J_i) \sigma^+
 = J_i (J_i \sigma^-) \sigma^+
 = J_i^2 (\sigma^- \sigma^+)
 = I \cdot \texttt{onSite}\,i\,(\sigma^- \sigma^+)
 = \texttt{onSite}\,i\,(\sigma^- \sigma^+)$,
using \texttt{jwString\_commute\_onSite}, \texttt{jwString\_sq}, and
\texttt{onSite\_mul\_onSite\_same}.
\end{proof}

\begin{lemma}[\texttt{fermionMultiNumber\_isHermitian},
            \texttt{fermionMultiNumber\_sq}]
$n_i$ is Hermitian and idempotent ($n_i^2 = n_i$, eigenvalues $0, 1$).
\end{lemma}

\begin{proof}[Proof sketch]
Reduce both via the previous lemma to the single-site
$\sigma^- \sigma^+ = \texttt{!![0,0;0,1]}$.  This $2 \times 2$
matrix is manifestly Hermitian (real-diagonal) and idempotent
(diagonal of $0, 1$).  Lift through \texttt{onSite\_isHermitian} and
\texttt{onSite\_mul\_onSite\_same}.
\end{proof}

\subsection{Number operator: commutativity and total}

\begin{lemma}[\texttt{fermionMultiNumber\_commute}]
For any $i, j \in \mathrm{Fin}\,(N+1)$, $[n_i, n_j] = 0$.
\end{lemma}

\begin{proof}[Proof sketch]
Both $n_i$ and $n_j$ equal $\texttt{onSite}\,k\,(\sigma^- \sigma^+)$
for $k = i, j$.  If $i = j$, commutativity is reflexive; if
$i \neq j$, different-site embeddings commute via
\texttt{onSite\_mul\_onSite\_of\_ne}.
\end{proof}

\begin{definition}[\texttt{fermionTotalNumber}]
$\hat N := \sum_{i \in \mathrm{Fin}\,(N+1)} n_i$.
\end{definition}

\begin{lemma}[\texttt{fermionTotalNumber\_isHermitian}]
$\hat N$ is Hermitian.
\end{lemma}

\begin{proof}[Proof sketch]
Finset induction: empty sum is the zero matrix (Hermitian); insert
step uses \texttt{Matrix.IsHermitian.add} to combine the new
\texttt{fermionMultiNumber} (Hermitian) with the inductive sum.
\end{proof}

\subsection{Same-site canonical anticommutation}

\begin{lemma}[\texttt{fermionMultiAnticomm\_self}]
$\{c_i, c_i^\dagger\} = c_i \cdot c_i^\dagger + c_i^\dagger \cdot c_i = I$.
\end{lemma}

\begin{proof}[Proof sketch]
Both products rearrange as
$(J \sigma^\pm)(J \sigma^\mp)
 = J (\sigma^\pm J) \sigma^\mp
 = J (J \sigma^\pm) \sigma^\mp
 = J^2 (\sigma^\pm \sigma^\mp)$,
using \texttt{jwString\_commute\_onSite} for the middle commutation.
After \texttt{jwString\_sq} reduces $J^2$ to $I$ and
\texttt{onSite\_mul\_onSite\_same} folds the products into the single
site, the sum becomes
$\texttt{onSite}\,i\,(\sigma^+\sigma^- + \sigma^-\sigma^+)
 = \texttt{onSite}\,i\,I = I$
via \texttt{onSite\_add} and the small auxiliary
\texttt{spinHalfOpPlus\_anticomm\_spinHalfOpMinus} (a $2\times 2$
\texttt{fin\_cases} computation showing the spin raising/lowering
anticommutator equals the identity).
\end{proof}

\subsection*{General cross-site CAR at site zero}

\noindent\textbf{Theorem (general \(\{c_0, c_k\} = 0\)).} For every
\(N : \mathbb{N}\) and every \(k : \texttt{Fin}\,(N + 1)\) with
\(0 < k.\texttt{val}\),
\[
  c_0 \cdot c_k + c_k \cdot c_0 \;=\; 0.
\]
The companion relation \(c_0^\dagger c_k^\dagger + c_k^\dagger c_0^\dagger = 0\)
follows by matrix \texttt{conjTranspose}.

\begin{proof}
The proof reduces to the \emph{operator-level} anticommutator
\(\{\sigma^+_0,\,\texttt{jwString}\,N\,k\} = 0\), which is proved
by induction on the length of the JW string.

\textbf{Base case} (\(k.\texttt{val} = 1\)):
\(\texttt{jwString}\,N\,\langle 1\rangle = \sigma^z_0\), and the
single-site Pauli anticommutator
\(\sigma^+\sigma^z + \sigma^z\sigma^+ = 0\) (immediate from
\texttt{spinHalfOpPlus\_mul\_pauliZ} \(= -\sigma^+\) and
\texttt{pauliZ\_mul\_spinHalfOpPlus} \(= \sigma^+\)) lifts to
\texttt{onSite} via additivity and \texttt{onSite\_zero}.

\textbf{Inductive step}: suppose
\(\{\sigma^+_0, \texttt{jwString}\,N\,\langle m+1\rangle\} = 0\).
Then using \texttt{jwString\_succ\_eq},
\(\texttt{jwString}\,N\,\langle m+2\rangle =
\texttt{jwString}\,N\,\langle m+1\rangle \cdot \sigma^z_{m+1}\), and
\(\sigma^z_{m+1}\) commutes with \(\sigma^+_0\) (different sites
\(m+1 \neq 0\), via \texttt{onSite\_mul\_onSite\_of\_ne}). Then
\begin{align*}
  \sigma^+_0 (\texttt{JW} \cdot \sigma^z_{m+1}) +
  (\texttt{JW} \cdot \sigma^z_{m+1})\sigma^+_0
  &= (\sigma^+_0\,\texttt{JW} + \texttt{JW}\,\sigma^+_0)\,\sigma^z_{m+1}
    \;=\; 0 \cdot \sigma^z_{m+1} \;=\; 0,
\end{align*}
where the first step uses the commutation
\(\sigma^z_{m+1}\sigma^+_0 = \sigma^+_0\sigma^z_{m+1}\) and the
distributivity of matrix multiplication over addition.

The main theorem follows by observing
\(c_k = \texttt{JW}\,\sigma^+_k\), so that
\begin{align*}
  c_0 c_k + c_k c_0
  &= \sigma^+_0 (\texttt{JW}\,\sigma^+_k) + (\texttt{JW}\,\sigma^+_k) \sigma^+_0 \\
  &= \sigma^+_0\,\texttt{JW}\,\sigma^+_k + \texttt{JW}\,\sigma^+_0\,\sigma^+_k \\
  &= (\sigma^+_0\,\texttt{JW} + \texttt{JW}\,\sigma^+_0)\,\sigma^+_k
    \;=\; 0\,\sigma^+_k \;=\; 0,
\end{align*}
using \(\sigma^+_k\sigma^+_0 = \sigma^+_0\sigma^+_k\) (different
sites).
\end{proof}

\subsection*{Fully general cross-site CAR for \(i < j\)}

\noindent\textbf{Theorem (general \(\{c_i, c_j\} = 0\)).} For every
\(N : \mathbb{N}\) and every \(i, j : \texttt{Fin}\,(N+1)\) with
\(i < j\),
\[
  c_i c_j + c_j c_i = 0,\quad
  c_i^\dagger c_j^\dagger + c_j^\dagger c_i^\dagger = 0,\quad
  c_i c_j^\dagger + c_j^\dagger c_i = 0,\quad
  c_i^\dagger c_j + c_j c_i^\dagger = 0.
\]

\begin{proof}
Write \(c_i = \texttt{JW}_i \sigma^+_i\), \(c_j = \texttt{JW}_j \sigma^+_j\).
The two key operator-level facts (tracked in Lean as
\texttt{jwString\_anticomm\_onSite\_pos\_spinHalfOpPlus} and
\texttt{jwString\_commute\_jwString}) are:
\begin{enumerate}
\item \emph{Interior-site anticommutation.}
For \(i < j\), \(\sigma^+_i \cdot \texttt{JW}_j +
\texttt{JW}_j \cdot \sigma^+_i = 0\). Proof: factor
\(\texttt{JW}_j = \sigma^z_i \cdot M\) where \(M =
\texttt{jwStringExceptAt}\,N\,j\,i\) is a product of \(\sigma^z\)
at sites \(\neq i\); then the single-site Pauli identities
\(\sigma^+\sigma^z = -\sigma^+\), \(\sigma^z\sigma^+ = \sigma^+\)
and the commutation \([\sigma^+_i, M] = 0\) close the argument.
\item \emph{JW string commutativity.}
\(\texttt{JW}_i \cdot \texttt{JW}_j = \texttt{JW}_j \cdot \texttt{JW}_i\)
always — both are \texttt{Finset.noncommProd}'s of \(\sigma^z_k\) at
distinct sites, hence pairwise commute.
\end{enumerate}

Using (1) as \(\sigma^+_i \cdot \texttt{JW}_j =
-(\texttt{JW}_j \cdot \sigma^+_i)\), together with the
single-site commutations
\([\sigma^+_i, \texttt{JW}_i] = 0\),
\([\sigma^+_j, \texttt{JW}_i] = 0\), and \([\sigma^+_i, \sigma^+_j] = 0\),
\begin{align*}
  c_i \cdot c_j
  &= \texttt{JW}_i \cdot \sigma^+_i \cdot \texttt{JW}_j \cdot \sigma^+_j
  \;=\; -\,\texttt{JW}_i \cdot \texttt{JW}_j \cdot \sigma^+_i \cdot \sigma^+_j,\\
  c_j \cdot c_i
  &= \texttt{JW}_j \cdot \sigma^+_j \cdot \texttt{JW}_i \cdot \sigma^+_i
  \;=\; \texttt{JW}_j \cdot \texttt{JW}_i \cdot \sigma^+_i \cdot \sigma^+_j.
\end{align*}
Adding and invoking (2),
\[
  c_i c_j + c_j c_i =
    (-\,\texttt{JW}_i \texttt{JW}_j + \texttt{JW}_j \texttt{JW}_i)
      \,\sigma^+_i \sigma^+_j \;=\; 0.
\]

The three remaining forms follow by:
\(\{c_i^\dagger, c_j^\dagger\}\) --- matrix
\texttt{conjTranspose};
\(\{c_i, c_j^\dagger\}\) --- same argument with \(\sigma^-_j\)
instead of \(\sigma^+_j\) at site \(j\), which still commutes with
\(\sigma^+_i\) (different sites);
\(\{c_i^\dagger, c_j\}\) --- \texttt{conjTranspose} again.
\end{proof}

\subsection{Hubbard chain (open and periodic boundary conditions)}
\label{subsec:hubbard-chain}

The Hubbard chain Hamiltonian on \(\mathrm{Fin}\,(N+1)\) sites is
\[
  \hat H_{\rm Hubbard}
    = -t \sum_{\langle i, j\rangle, \alpha}
        \bigl(c_{i,\alpha}^\dagger c_{j,\alpha}
              + c_{j,\alpha}^\dagger c_{i,\alpha}\bigr)
      + U \sum_i n_{i,\uparrow} n_{i,\downarrow},
\]
with the spinful operators
\(c_{i,\alpha}\) for \(\alpha \in \{\uparrow,\downarrow\}\)
constructed via Jordan--Wigner from the spinless multi-mode
operators of the previous section.  The Lean formalisation is
\texttt{LatticeSystem/Fermion/JordanWigner/Hubbard.lean}.

\textbf{Status}: Formalised on
\(\mathrm{Fin}\,(N+1)\) (open BC) and \(\mathrm{Fin}\,(N+2)\)
(periodic BC) with the corresponding Gibbs state +
expectation companion family
(\texttt{hubbardChain\{Hamiltonian, Gibbs\}\{State,Expectation\}\_*}
and the periodic counterparts, Phase 2 work, issue \#241).
The companion family is at parity with the 1D / 2D / 3D
Heisenberg variants (Hermiticity, commute-with-H, $\beta = 0$
closed form, im-of-Hermitian, commutator, mul-Hamiltonian-im,
hamiltonian-pow-im, anticommutator, commutator-re,
partition function im, ofReal-re-eq, pow-trace).

\subsection{Hubbard spin symmetry: U(1)\texorpdfstring{$\times$}{×}U(1) invariance (Tasaki §9.3.3)}
\label{subsec:hubbard-spin-symmetry}

\noindent\textbf{Theorem} (Tasaki §9.3.3, eq.~(9.3.35), p.~333).
The Hubbard Hamiltonian $\hat H = \hat H_{\rm kin} + \hat H_{\rm int}$
satisfies
\[
  [N_\uparrow, \hat H] = [N_\downarrow, \hat H] = [S^z_{\rm tot}, \hat H] = 0,
\]
where $N_\sigma = \sum_x \hat n_{x,\sigma}$ is the total particle count
per spin and $S^z_{\rm tot} = (N_\uparrow - N_\downarrow)/2$.

\noindent\textbf{Key auxiliary commutators}.
\begin{align*}
  [N_\uparrow, \hat c^\dagger_{i,\uparrow}] &= \hat c^\dagger_{i,\uparrow}, &
  [N_\uparrow, \hat c_{i,\uparrow}] &= -\hat c_{i,\uparrow}.
\end{align*}
These are proved by writing $N_\uparrow = \sum_k n_{k,\uparrow}$ and
isolating the $k = i$ term (same-site commutator $[n_i, c^\dagger_i] = c^\dagger_i$)
while all other terms vanish by
$[n_k, c^\dagger_i] = 0$ ($k \neq i$, different JW positions).

\noindent\textbf{Proof sketch}.
For the kinetic term, each bond contributes
$[N_\uparrow, \hat c^\dagger_{i,\uparrow}\hat c_{j,\uparrow}] =
 [N_\uparrow,\hat c^\dagger_{i,\uparrow}]\hat c_{j,\uparrow} +
 \hat c^\dagger_{i,\uparrow}[N_\uparrow,\hat c_{j,\uparrow}] =
 \hat c^\dagger_{i,\uparrow}\hat c_{j,\uparrow} -
 \hat c^\dagger_{i,\uparrow}\hat c_{j,\uparrow} = 0$.
The cross-species kinetic term $\hat c^\dagger_{i,\downarrow}\hat c_{j,\downarrow}$
commutes with $N_\uparrow$ because the JW positions are of opposite parity
(odd vs.\ even); this was already in \texttt{Charges.lean}.
For the interaction, $[N_{\uparrow/\downarrow}, \hat H_{\rm int}] = 0$
was established in \texttt{Charges.lean}.

\noindent\textbf{Lean names} (file \texttt{Fermion/JordanWigner/Hubbard/SpinSymmetry.lean}):
\begin{itemize}
  \item \texttt{fermionTotalUpNumber\_commutator\_fermionUpCreation}:
    $[N_\uparrow, c^\dagger_{i,\uparrow}] = c^\dagger_{i,\uparrow}$.
  \item \texttt{fermionTotalUpNumber\_commute\_upHopping}:
    $[N_\uparrow, c^\dagger_{i,\uparrow}c_{j,\uparrow}] = 0$.
  \item \texttt{fermionTotalUpNumber\_commute\_hubbardHamiltonian}:
    $[N_\uparrow, H] = 0$.
  \item \texttt{fermionTotalDownNumber\_commute\_hubbardHamiltonian}:
    $[N_\downarrow, H] = 0$.
  \item \texttt{fermionTotalSpinZ\_commute\_hubbardHamiltonian}:
    $[S^z_{\rm tot}, H] = 0$.
\end{itemize}

\subsection{Hubbard spin symmetry: full SU(2) invariance (Tasaki §9.3.3)}
\label{subsec:hubbard-su2-symmetry}

\noindent\textbf{Theorem} (Tasaki §9.3.3, eq.~(9.3.35), p.~333).
The Hubbard Hamiltonian also commutes with the SU(2) raising and lowering operators:
\[
  [\hat S^+_{\rm tot}, \hat H] = [\hat S^-_{\rm tot}, \hat H] = 0,
\]
where
\[
  \hat S^+_{\rm tot} = \sum_{x} \hat c^\dagger_{x,\uparrow}\hat c_{x,\downarrow},
  \qquad
  \hat S^-_{\rm tot} = (\hat S^+_{\rm tot})^\dagger
    = \sum_x \hat c^\dagger_{x,\downarrow}\hat c_{x,\uparrow}.
\]

\noindent\textbf{Key auxiliary commutators} (Tasaki §9.3.3, eq.~(9.3.36), p.~333):
\begin{align*}
  [\hat c^\dagger_{j,\uparrow}, \hat S^+_{\rm tot}] &= 0, &
  [\hat c_{j,\uparrow},   \hat S^+_{\rm tot}] &= \hat c_{j,\downarrow}, \\
  [\hat c^\dagger_{j,\downarrow}, \hat S^+_{\rm tot}] &= -\hat c^\dagger_{j,\uparrow}, &
  [\hat c_{j,\downarrow},  \hat S^+_{\rm tot}] &= 0.
\end{align*}
Each is proved by the fermionic Leibniz rule
$[A, BC] = \{A,B\}C - B\{A,C\}$
and the CAR identities (same-site $\{c_{j,\sigma}, c^\dagger_{j,\sigma}\} = 1$,
cross-site $\{c_{j,\sigma}, c^\dagger_{k,\tau}\} = 0$ for $(j,\sigma)\neq(k,\tau)$).

\noindent\textbf{Proof sketch for $[\hat S^+_{\rm tot}, \hat H] = 0$}.
\emph{Kinetic term}: for each bond $(i,j)$,
\[
  [\hat S^+_{\rm tot},\, \hat c^\dagger_{i,\uparrow}\hat c_{j,\uparrow}
     + \hat c^\dagger_{i,\downarrow}\hat c_{j,\downarrow}]
  = \hat c^\dagger_{i,\uparrow}\hat c_{j,\downarrow}
    + (-\hat c^\dagger_{i,\uparrow})\hat c_{j,\downarrow} = 0.
\]
\emph{Interaction term}: at the same site ($k = x$) both products vanish by
Pauli exclusion; at different sites ($k \neq x$) all four constituent operators
commute pairwise (different JW positions), giving zero commutator.
The result $[\hat S^-_{\rm tot}, \hat H] = 0$ follows by taking adjoints.

\noindent\textbf{Lean names} (file \texttt{Fermion/JordanWigner/Hubbard/SpinSymmetry.lean}):
\begin{itemize}
  \item \texttt{fermionTotalSpinPlus} / \texttt{fermionTotalSpinMinus}: definitions.
  \item \texttt{fermionTotalSpinPlus\_conjTranspose}:
    $(\hat S^+_{\rm tot})^\dagger = \hat S^-_{\rm tot}$.
  \item \texttt{fermionUpAnnihilation\_commutator\_fermionTotalSpinPlus}:
    $[\hat c_{j,\uparrow}, \hat S^+_{\rm tot}] = \hat c_{j,\downarrow}$.
  \item \texttt{fermionDownCreation\_commutator\_fermionTotalSpinPlus}:
    $[\hat c^\dagger_{j,\downarrow}, \hat S^+_{\rm tot}] = -\hat c^\dagger_{j,\uparrow}$.
  \item \texttt{fermionTotalSpinPlus\_commute\_hubbardHamiltonian}:
    $[\hat S^+_{\rm tot}, \hat H] = 0$.
  \item \texttt{fermionTotalSpinMinus\_commute\_hubbardHamiltonian}:
    $[\hat S^-_{\rm tot}, \hat H] = 0$ (by adjoint).
\end{itemize}

%======================================================================
\subsection{Hubbard all-up-spin state: no double occupancy (Tasaki §11.1.1)}
\label{subsec:hubbard-allup-state}
%======================================================================

\noindent\textbf{Definition.}
The \emph{all-up-spin state} for the Hubbard model on $N+1$ sites is
\[
  |\!\uparrow\cdots\uparrow\rangle
  = \hat c^\dagger_{N,\uparrow}\cdots\hat c^\dagger_{1,\uparrow}
    \hat c^\dagger_{0,\uparrow}|\Phi_{\rm vac}\rangle.
\]
In the Jordan--Wigner representation on $\mathrm{Fin}(2N+2)$
(even indices $= $ spin-up, odd indices $=$ spin-down), this is
$\mathrm{basisVec}(\lambda k.\; \text{if } k \bmod 2 = 0 \text{ then } 1 \text{ else } 0)$.

\noindent\textbf{Theorem} (Tasaki §11.1.1, p.~373; eq.~(10.1.5), p.~344).
\begin{enumerate}
\item $\hat n_{i,\uparrow}\,|\!\uparrow\cdots\uparrow\rangle
        = |\!\uparrow\cdots\uparrow\rangle$ for every site $i$.
\item $\hat n_{i,\downarrow}\,|\!\uparrow\cdots\uparrow\rangle
        = 0$ for every site $i$ (no down-spin electrons present).
\item $\hat H_{\rm int}\,|\!\uparrow\cdots\uparrow\rangle = 0$
        (no double occupancy $\Rightarrow$ interaction vanishes).
\item $\hat H\,|\!\uparrow\cdots\uparrow\rangle
        = \hat H_{\rm hop}\,|\!\uparrow\cdots\uparrow\rangle$
        (the Hubbard model in the all-up sector reduces to
        a non-interacting hopping problem).
\end{enumerate}

\noindent\textbf{Proof sketch for item 3.}
$\hat H_{\rm int} = U\sum_x \hat n_{x,\uparrow}\hat n_{x,\downarrow}$.
By item 2, $\hat n_{x,\downarrow}|\!\uparrow\cdots\uparrow\rangle = 0$ for
every $x$, so each term vanishes and the sum is zero.

\noindent\textbf{Kinetic eigenvalue} (Tasaki §11.1.1, p.~373).
The hopping Hamiltonian acts as
\[
  \hat H_{\rm hop}\,|\!\uparrow\cdots\uparrow\rangle
  = \Bigl(\sum_i t_{ii}\Bigr)\,|\!\uparrow\cdots\uparrow\rangle.
\]
\begin{itemize}
  \item \emph{Down-spin terms vanish}: $\hat c_{i,\downarrow}|\!\uparrow\cdots\uparrow\rangle = 0$
    because the spin-down JW index $2i{+}1$ is unoccupied, so $\sigma^+$ maps it to~0.
  \item \emph{Up-spin off-diagonal ($i\neq j$) terms vanish}: by the CAR
    $\{\hat c_{j,\uparrow}, \hat c^\dagger_{i,\uparrow}\} = 0$ for $i\neq j$,
    $\hat c^\dagger_{i,\uparrow}\hat c_{j,\uparrow}
      = -\hat c_{j,\uparrow}\hat c^\dagger_{i,\uparrow}$.
    Then $\hat c^\dagger_{i,\uparrow}|\!\uparrow\cdots\uparrow\rangle = 0$
    (site $2i$ already occupied, $\sigma^-$ column 1 is all-zero).
  \item \emph{Diagonal terms}: $\hat c^\dagger_{i,\uparrow}\hat c_{i,\uparrow}
      = \hat n_{i,\uparrow}$ and $\hat n_{i,\uparrow}|\!\uparrow\cdots\uparrow\rangle
      = |\!\uparrow\cdots\uparrow\rangle$.
\end{itemize}

\noindent\textbf{Lean names}
(file \texttt{Fermion/JordanWigner/Hubbard/AllUpState.lean}):
\begin{itemize}
  \item \texttt{hubbardAllUpState N}: the basis vector for the all-up state.
  \item \texttt{fermionUpNumber\_mulVec\_allUpState}:
    $\hat n_{i,\uparrow}\cdot|\!\uparrow\cdots\uparrow\rangle = |\!\uparrow\cdots\uparrow\rangle$.
  \item \texttt{fermionDownNumber\_mulVec\_allUpState}:
    $\hat n_{i,\downarrow}\cdot|\!\uparrow\cdots\uparrow\rangle = 0$.
  \item \texttt{fermionDownAnnihilation\_mulVec\_allUpState}:
    $\hat c_{i,\downarrow}\cdot|\!\uparrow\cdots\uparrow\rangle = 0$.
  \item \texttt{fermionUpCreation\_mulVec\_allUpState}:
    $\hat c^\dagger_{i,\uparrow}\cdot|\!\uparrow\cdots\uparrow\rangle = 0$.
  \item \texttt{hubbardOnSiteInteraction\_mulVec\_allUpState}:
    $\hat H_{\rm int}\cdot|\!\uparrow\cdots\uparrow\rangle = 0$.
  \item \texttt{hubbardHamiltonian\_mulVec\_allUpState}:
    $\hat H\cdot|\!\uparrow\cdots\uparrow\rangle
      = \hat H_{\rm hop}\cdot|\!\uparrow\cdots\uparrow\rangle$.
  \item \texttt{hubbardKinetic\_mulVec\_allUpState}:
    $\hat H_{\rm hop}\cdot|\!\uparrow\cdots\uparrow\rangle
      = \bigl(\sum_i t_{ii}\bigr)\,|\!\uparrow\cdots\uparrow\rangle$.
  \item \texttt{hubbardHamiltonian\_mulVec\_allUpState\_eigenstate}:
    $\hat H\cdot|\!\uparrow\cdots\uparrow\rangle
      = \bigl(\sum_i t_{ii}\bigr)\,|\!\uparrow\cdots\uparrow\rangle$.
\end{itemize}

%======================================================================
\subsection{Hubbard hard-core subspace (Tasaki §11.2)}
\label{subsec:hubbard-hardcore-subspace}
%======================================================================

\noindent\textbf{Definition.}
For the Hubbard model on \(N+1\) sites, the same-site
double-occupancy operator is
\[
  \hat D_i := \hat n_{i,\uparrow}\hat n_{i,\downarrow}.
\]
The hard-core subspace is the linear subspace
\[
  \mathcal H_{\rm hc}
  :=
  \{\psi \mid \hat D_i\psi = 0 \text{ for every site } i\}.
\]
This is the no-double-occupancy sector used in the \(U=\infty\)
and one-hole Nagaoka discussion as unnumbered infrastructure for
Tasaki Theorems 11.5 and 11.7 (1st edition, §11.2, pp.~381--388).

\noindent\textbf{Lean theorem} (unnumbered hard-core infrastructure for
Tasaki Theorems 11.5 and 11.7; 1st edition, §11.2, pp.~381--388).
If \(\psi \in \mathcal H_{\rm hc}\), then the on-site Hubbard
interaction annihilates \(\psi\):
\[
  \hat H_{\rm int}\psi
  =
  U\sum_i \hat n_{i,\uparrow}\hat n_{i,\downarrow}\psi
  =
  0.
\]

\noindent\textbf{Proof.}
Each summand is \(U\,\hat D_i\psi\), and the defining hard-core
condition gives \(\hat D_i\psi=0\) for every \(i\). Therefore the
finite sum is zero.

\noindent\textbf{Lean names}
(file \texttt{Fermion/JordanWigner/Hubbard/HardcoreSubspace.lean}):
\begin{itemize}
  \item \texttt{hubbardDoubleOccupancy N i}: \(\hat D_i\).
  \item \texttt{hubbardHardcoreSubspace N}: the linear subspace
    \(\mathcal H_{\rm hc}\).
  \item \texttt{mem\_hubbardHardcoreSubspace\_iff}: membership is
    equivalent to all double-occupancy actions vanishing.
  \item \texttt{hubbardDoubleOccupancy\_mulVec\_eq\_zero\_of\_mem\_hardcore}:
    each \(\hat D_i\) annihilates a hard-core vector.
  \item \texttt{hubbardOnSiteInteraction\_mulVec\_eq\_zero\_of\_mem\_hardcore}
    and
    \texttt{hubbardOnSiteInteraction\_apply\_eq\_zero\_of\_mem\_hardcore}:
    vector and pointwise forms of \(\hat H_{\rm int}\psi=0\).
\end{itemize}

%======================================================================
\subsection{Hubbard hard-core projection (Tasaki §11.2)}
\label{subsec:hubbard-hardcore-projection}
%======================================================================

\noindent\textbf{Construction.}
The same-site double-occupancy operators
\(\hat D_i = n_{i,\uparrow}n_{i,\downarrow}\) are pairwise-commuting
idempotents: \(\hat D_i^2 = \hat D_i\) and
\([\hat D_i,\hat D_j]=0\) (the four underlying Jordan--Wigner number
operators live at distinct positions). Hence each
\emph{hard-core factor}
\[
  \hat P_i := 1 - \hat D_i
\]
is an idempotent (\(\hat P_i^2 = 1 - 2\hat D_i + \hat D_i^2
= 1 - \hat D_i = \hat P_i\)), the factors pairwise commute, and
\(\hat D_i\hat P_i = \hat D_i - \hat D_i^2 = 0\). The
\emph{hard-core projection} is the non-commutative product
\[
  \hat P_{\rm hc} := \prod_{i} \hat P_i = \prod_{i}\bigl(1 - \hat D_i\bigr),
\]
well-defined because the factors commute.

\medskip
\noindent\textbf{Properties} (unnumbered infrastructure for Tasaki
Theorems 11.5 and 11.7; 1st edition, §11.2, pp.~381--388).
\begin{enumerate}
  \item \emph{Idempotence}: \(\hat P_{\rm hc}^2 = \hat P_{\rm hc}\).
    A non-commutative product of pairwise-commuting idempotents is
    idempotent, proved by induction on the index finset using
    \(\hat P_i^2=\hat P_i\) and commutativity to regroup.
  \item \emph{Annihilation}: for every site \(j\),
    \(\hat D_j\,\hat P_{\rm hc} = 0\). Extract the \(j\)-th factor to the
    front (the factors commute) and use \(\hat D_j\hat P_j = 0\).
  \item \emph{Fixed points}: \(\hat P_{\rm hc}\psi = \psi\) for every
    \(\psi\in\mathcal H_{\rm hc}\), since each factor fixes \(\psi\)
    (\(\hat P_i\psi = \psi - \hat D_i\psi = \psi\)).
  \item \emph{Range}: \(\hat P_{\rm hc}\psi\in\mathcal H_{\rm hc}\) for
    every \(\psi\), since
    \(\hat D_j(\hat P_{\rm hc}\psi) = (\hat D_j\hat P_{\rm hc})\psi = 0\).
  \item \emph{Hermiticity}: \(\hat P_{\rm hc}^\dagger = \hat P_{\rm hc}\).
    Each \(\hat D_i\) is Hermitian (product of commuting Hermitian number
    operators), so each factor \(\hat P_i = 1 - \hat D_i\) is Hermitian, and
    a non-commutative product of pairwise-commuting Hermitian matrices is
    Hermitian. Together with idempotence this makes \(\hat P_{\rm hc}\) an
    orthogonal projection.
\end{enumerate}
Thus \(\hat P_{\rm hc}\) is the projection onto the no-double-occupancy
(hard-core) subspace \(\mathcal H_{\rm hc}\) used in the infinite-\(U\) /
one-hole Nagaoka sector.

\medskip
\noindent\textbf{Lean declarations}
(file \texttt{Fermion/JordanWigner/Hubbard/HardcoreProjection.lean}):
\begin{itemize}
  \item \texttt{hubbardHardcoreFactor N i}: \(\hat P_i = 1 - \hat D_i\).
  \item \texttt{hubbardHardcoreFactor\_mul\_self}: \(\hat P_i^2 = \hat P_i\).
  \item \texttt{hubbardDoubleOccupancy\_mul\_hardcoreFactor}:
    \(\hat D_i\hat P_i = 0\).
  \item \texttt{hubbardHardcoreFactor\_commute}: factors commute.
  \item \texttt{hubbardDoubleOccupancy\_isHermitian} and
    \texttt{hubbardHardcoreFactor\_isHermitian}: \(\hat D_i\) and
    \(\hat P_i\) are Hermitian.
  \item \texttt{hubbardHardcoreFactor\_mulVec\_eq\_self\_of\_mem}:
    each factor fixes a hard-core vector.
  \item \texttt{hubbardHardcoreProjection N}:
    \(\hat P_{\rm hc} = \prod_i \hat P_i\).
  \item \texttt{hubbardHardcoreProjection\_mul\_self}:
    \(\hat P_{\rm hc}^2 = \hat P_{\rm hc}\).
  \item \texttt{hubbardHardcoreProjection\_isHermitian}:
    \(\hat P_{\rm hc}^\dagger = \hat P_{\rm hc}\).
  \item \texttt{hubbardDoubleOccupancy\_mul\_hardcoreProjection}:
    \(\hat D_j\hat P_{\rm hc} = 0\).
  \item \texttt{hubbardHardcoreProjection\_mulVec\_eq\_self\_of\_mem}:
    \(\hat P_{\rm hc}\psi=\psi\) on \(\mathcal H_{\rm hc}\).
  \item \texttt{hubbardHardcoreProjection\_mulVec\_mem}:
    \(\hat P_{\rm hc}\psi\in\mathcal H_{\rm hc}\).
\end{itemize}

%======================================================================
\subsection{One-hole hard-core basis states (Tasaki §11.2)}
\label{subsec:hubbard-hardcore-basis}
%======================================================================

\noindent\textbf{Construction.}
The one-hole hard-core basis states \(|\Phi_{x,\sigma}\rangle\) (Tasaki
eq.~(11.2.3)) label the infinite-\(U\) / one-hole Nagaoka sector: one
spinful site \(x\in\Lambda\) is the hole (both orbitals empty), and every
other site \(j\neq x\) carries exactly one electron of spin
\(\sigma(j)\in\{\uparrow,\downarrow\}\). In the Jordan--Wigner occupation
representation this is the computational basis configuration
\(c_{x,\sigma}\) with
\[
  c_{x,\sigma}(2j)=\begin{cases}0 & j=x\\ 1 & j\neq x,\ \sigma(j)=\uparrow\\
    0 & j\neq x,\ \sigma(j)=\downarrow\end{cases}
  \qquad
  c_{x,\sigma}(2j{+}1)=\begin{cases}0 & j=x\\ 0 & j\neq x,\ \sigma(j)=\uparrow\\
    1 & j\neq x,\ \sigma(j)=\downarrow,\end{cases}
\]
and \(|\Phi_{x,\sigma}\rangle := \mathtt{basisVec}\,(c_{x,\sigma})\). The
states are defined sign-free as occupation configurations; the relation to
Tasaki's ordered creation-operator products (carrying fermion signs,
relevant to the hopping matrix elements (11.2.5)) is deferred to the
effective-Hamiltonian layer.

\medskip
\noindent\textbf{Properties} (unnumbered one-hole hard-core basis
infrastructure for Tasaki Theorems 11.5 and 11.7; 1st edition, §11.2,
pp.~381--388).
\begin{enumerate}
  \item \emph{Hard-core membership}: at every site at least one orbital is
    empty, so \(\hat D_i|\Phi_{x,\sigma}\rangle = 0\) and
    \(|\Phi_{x,\sigma}\rangle\in\mathcal H_{\rm hc}\). Concretely, at the
    hole the up-orbital is empty; off the hole the empty orbital is the
    down-orbital for \(\sigma=\uparrow\) and the up-orbital for
    \(\sigma=\downarrow\), and \(\hat D_i = n_{i,\uparrow}n_{i,\downarrow}\)
    is killed by whichever number operator sees the empty orbital (using
    \([n_{i,\uparrow},n_{i,\downarrow}]=0\)).
  \item \emph{Projection-fixed}: \(\hat P_{\rm hc}|\Phi_{x,\sigma}\rangle =
    |\Phi_{x,\sigma}\rangle\), immediate from membership.
  \item \emph{Orthonormality}: since the states are distinct computational
    basis vectors, \(\langle\Phi_{x,\sigma}|\Phi_{x',\sigma'}\rangle = 1\) if
    \(c_{x,\sigma}=c_{x',\sigma'}\) and \(0\) otherwise; in particular each
    state is normalised.
\end{enumerate}

\medskip
\noindent\textbf{Lean declarations}
(file \texttt{Fermion/JordanWigner/Hubbard/HardcoreBasis.lean}):
\begin{itemize}
  \item \texttt{hubbardOneHoleConfig N x }\(\sigma\): the configuration
    \(c_{x,\sigma}\).
  \item \texttt{hubbardOneHoleConfig\_apply\_up} /
    \texttt{hubbardOneHoleConfig\_apply\_down}: its orbital occupation
    values per site.
  \item \texttt{hubbardHardcoreBasisState N x }\(\sigma\):
    \(|\Phi_{x,\sigma}\rangle\).
  \item \texttt{hubbardHardcoreBasisState\_mem\_hardcoreSubspace}:
    \(|\Phi_{x,\sigma}\rangle\in\mathcal H_{\rm hc}\).
  \item \texttt{hubbardHardcoreProjection\_mulVec\_basisState}:
    \(\hat P_{\rm hc}|\Phi_{x,\sigma}\rangle = |\Phi_{x,\sigma}\rangle\).
  \item \texttt{hubbardHardcoreBasisState\_inner} and
    \texttt{hubbardHardcoreBasisState\_self\_inner}: orthonormality and
    normalisation.
\end{itemize}

%======================================================================
\subsection{Jordan--Wigner string action on basis states (Tasaki §11.2)}
\label{subsec:jwstring-basisvec}
%======================================================================

\noindent\textbf{Statement.}
The Jordan--Wigner string \(\hat J_i := \prod_{j<i}\sigma^z_j\) is diagonal in
the computational basis, since \(\sigma^z = \operatorname{diag}(1,-1)\). On a
basis state \(|c\rangle\) each factor contributes the parity sign
\((-1)^{c_j}\), so
\[
  \hat J_i\,|c\rangle = \Bigl(\prod_{j<i}(-1)^{c_j}\Bigr)\,|c\rangle,
\]
the fermion-parity sign of the modes occupied below \(i\). This is the origin
of the fermion signs in the hopping matrix elements (11.2.4)/(11.2.5), and is
the foundation for computing the action of \(\hat c_i\), \(\hat c_i^\dagger\) on
basis states.

\medskip
\noindent\textbf{Lean declarations}
(file \texttt{Fermion/JordanWigner/StringBasisVecAction.lean}):
\begin{itemize}
  \item \texttt{onSite\_pauliZ\_mulVec\_basisVec}:
    \(\sigma^z_j|c\rangle = (-1)^{c_j}|c\rangle\).
  \item \texttt{jwString\_mulVec\_basisVec}:
    \(\hat J_i|c\rangle = \bigl(\prod_{j<i}(-1)^{c_j}\bigr)|c\rangle\).
\end{itemize}

\medskip
\noindent\textbf{Annihilation and creation on basis states.}
Writing the string sign \(\varepsilon_j(c) := \prod_{k<j}(-1)^{c_k}\) and using
\(\hat c_j = \hat J_j\,\sigma^+_j\), \(\hat c_j^\dagger = \hat J_j\,\sigma^-_j\)
with \(\sigma^+=\bigl(\begin{smallmatrix}0&1\\0&0\end{smallmatrix}\bigr)\),
\(\sigma^-=\bigl(\begin{smallmatrix}0&0\\1&0\end{smallmatrix}\bigr)\) flipping the
site-\(j\) occupation, the single-mode operators act on a basis state by
\[
  \hat c_j|c\rangle =
    \begin{cases}\varepsilon_j(c)\,|c^{j\mapsto0}\rangle & c_j=1\\ 0 & c_j=0,\end{cases}
  \qquad
  \hat c_j^\dagger|c\rangle =
    \begin{cases}\varepsilon_j(c)\,|c^{j\mapsto1}\rangle & c_j=0\\ 0 & c_j=1.\end{cases}
\]
(\texttt{jwSign}, \texttt{fermionMultiAnnihilation\_mulVec\_basisVec},
\texttt{fermionMultiCreation\_mulVec\_basisVec} in
\texttt{Fermion/JordanWigner/AnnihilationCreationBasisVec.lean}.)
The base case on the vacuum (all-empty, string sign \(1\)) is
\(\hat c_j^\dagger|\mathrm{vac}\rangle = |0\dots0,\,j\mapsto1\rangle\)
(\texttt{jwSign\_zero\_config},
\texttt{fermionMultiCreation\_mulVec\_vacuum\_eq\_basisVec} in
\texttt{Fermion/JordanWigner/VacuumCreationBasisVec.lean}), the starting point
for the ordered creation-operator construction of the (11.2.3) basis states.

\medskip
\noindent\textbf{A single hopping term.}
Composing the two, the hopping term \(\hat c_p^\dagger \hat c_q\) moves the
electron from \(q\) to \(p\):
\[
  \hat c_p^\dagger \hat c_q\,|c\rangle =
    \begin{cases}
      \varepsilon_q(c)\,\varepsilon_p(c^{q\mapsto0})\,
        |c^{q\mapsto0,\,p\mapsto1}\rangle
        & c_q = 1 \text{ and } (c^{q\mapsto0})_p = 0,\\
      0 & \text{otherwise.}
    \end{cases}
\]
(\texttt{fermionMultiCreation\_mul\_Annihilation\_mulVec\_basisVec} in
\texttt{Fermion/JordanWigner/HopBasisVec.lean}.) This is the operator-level
content of (11.2.4); the fermion signs \(\varepsilon\) combine into the
sign of the hopping matrix element (11.2.5).

\medskip
\noindent\textbf{Spatial content of (11.2.4) for the one-hole sector.}
When an electron of spin \(s\) hops from an occupied site \(y\) into the hole
\(x\) of \(|\Phi_{x,\sigma}\rangle\), the resulting occupation configuration is
exactly that of \(|\Phi_{y,\sigma_{y\to x}}\rangle\), where the new hole is at
\(y\) and \(\sigma_{y\to x}\) carries the spin of the moved electron to \(x\):
\(\sigma_{y\to x}\) agrees with \(\sigma\) off \(x\) and equals \(\sigma_y\) at
\(x\). The hypothesis that the \(s\)-orbital at \(y\) is occupied forces \(s\)
to be the spin of the electron at \(y\). (\texttt{hubbardSpinMove},
\texttt{hubbardOneHoleConfig\_hop} in
\texttt{Fermion/JordanWigner/Hubbard/HopConfig.lean}.)

\medskip
\noindent\textbf{Operator action of the hopping term (11.2.4).}
Assembling the configuration identity with the basis-state hop computation, the
hole-filling hopping term acts on a one-hole basis state by
\[
  \hat c^\dagger_{(x,s)}\,\hat c_{(y,s)}\,|\Phi_{x,\sigma}\rangle
    = \varepsilon_{(y,s)}(c_{x,\sigma})\,
      \varepsilon_{(x,s)}\!\bigl(c_{x,\sigma}^{(y,s)\mapsto0}\bigr)\,
      |\Phi_{y,\sigma_{y\to x}}\rangle,
\]
when the \(s\)-orbital at \(y\) is occupied (and \(x\neq y\)). The scalar is the
honest computational-basis fermion sign in the sign-free \(\mathtt{basisVec}\)
convention; Tasaki's uniform \(-t_{x,y}\) of (11.2.5) arises in his ordered
creation-operator basis (11.2.3), recovered in a later layer.
(\texttt{hubbardHop\_mulVec\_hardcoreBasisState} in
\texttt{Fermion/JordanWigner/Hubbard/HopAction.lean}.)

\medskip
\noindent\textbf{Per-term matrix element (11.2.5).}
Pairing with another one-hole basis state and using orthonormality, the matrix
element of the hopping term between basis states is the sign product times the
indicator that the target is the moved configuration:
\[
  \langle \Phi_{y,\tau} |\, \hat c^\dagger_{(x,s)}\hat c_{(y,s)}\,|
    \Phi_{x,\sigma}\rangle
    = \varepsilon_{(y,s)}\,\varepsilon_{(x,s)}\,
      [\,\tau = \sigma_{y\to x}\,].
\]
The off-diagonal support is exactly \(\tau = \sigma_{y\to x}\), matching the
structure of (11.2.5). (\texttt{hubbardHopTerm\_inner\_hardcoreBasisState} in
\texttt{Fermion/JordanWigner/Hubbard/HopMatrixElement.lean}.)

%======================================================================
\subsection{Span of the one-hole hard-core sector (Tasaki §11.2, footnote 8)}
\label{subsec:hubbard-hardcore-span}
%======================================================================

\noindent\textbf{Statement.}
Tasaki's footnote 8 (p.~382) records that the one-hole hard-core sector
\(\mathcal H_{\rm hc}^{N}\) (with \(N=|\Lambda|-1\) electrons and no doubly
occupied site) is \emph{spanned} by the basis states \(|\Phi_{x,\sigma}\rangle\)
of (11.2.3). We call a computational basis configuration \emph{one-hole
hard-core} when it has no doubly occupied site and exactly one empty site (the
hole).

\medskip
\noindent\textbf{Argument.}
The combinatorial content is the surjectivity of the parametrization
\((x,\sigma)\mapsto c_{x,\sigma}\) onto the one-hole hard-core configurations:
given such a configuration, its unique empty site is the hole \(x\), and reading
off the spin of each singly occupied site recovers \(\sigma\). Concretely, at a
site \(i\neq x\), no double occupancy together with \(i\) not being the (unique)
hole forces exactly one orbital to be occupied, which fixes \(\sigma_i\). Hence
the basis states \(|\Phi_{x,\sigma}\rangle = \mathtt{basisVec}\,(c_{x,\sigma})\)
are precisely the computational basis vectors of the one-hole hard-core
configurations, so
\[
  \mathcal H_{\rm hc}^{N}
    = \operatorname{span}\bigl\{\,|\Phi_{x,\sigma}\rangle \ \big|\ x\in\Lambda,\
      \sigma\,\bigr\}.
\]

\medskip
\noindent\textbf{Lean declarations}
(file \texttt{Fermion/JordanWigner/Hubbard/HardcoreSpan.lean}):
\begin{itemize}
  \item \texttt{IsOneHoleHardcoreConfig N c}: no double occupancy and a unique
    empty site.
  \item \texttt{hubbardOneHoleConfig\_isOneHoleHardcore}: \(c_{x,\sigma}\) is
    one-hole hard-core.
  \item \texttt{exists\_eq\_hubbardOneHoleConfig\_of\_isOneHoleHardcore}:
    surjectivity of the parametrization.
  \item \texttt{hubbardOneHoleHardcoreSector N}: the sector
    \(\mathcal H_{\rm hc}^{N}\).
  \item \texttt{hubbardOneHoleHardcoreSector\_eq\_span\_basisState}: footnote 8
    spanning statement.
  \item \texttt{hubbardHardcoreBasisState\_mem\_sector}: each basis state lies in
    the sector.
\end{itemize}

%======================================================================
\subsection{Effective Hamiltonian on the hard-core sector (Tasaki §11.2)}
\label{subsec:hubbard-effective-hamiltonian}
%======================================================================

\noindent\textbf{Definition.}
The effective Hamiltonian of the infinite-\(U\) / one-hole Nagaoka sector is
the hard-core compression of the full Hubbard Hamiltonian
\(H = H_{\rm hop} + H_{\rm int}\),
\[
  \hat H_{\rm eff} := \hat P_{\rm hc}\,H\,\hat P_{\rm hc}.
\]

\medskip
\noindent\textbf{Properties} (unnumbered effective-Hamiltonian
infrastructure for Tasaki Theorems 11.5 and 11.7; 1st edition, §11.2,
pp.~381--388).
\begin{enumerate}
  \item \emph{Hermiticity}: for a Hermitian hopping matrix
    (\(\overline{t_{ij}} = t_{ji}\)) and real \(U\), \(\hat H_{\rm eff}\) is
    Hermitian, being the compression \(\hat P_{\rm hc}^\dagger H
    \hat P_{\rm hc}\) of a Hermitian operator by the Hermitian projection.
  \item \emph{\(U\uparrow\infty\) reduction}: on a hard-core vector
    \(\psi\in\mathcal H_{\rm hc}\), the on-site interaction vanishes
    (\(H_{\rm int}\psi = 0\)) and the inner projection fixes \(\psi\)
    (\(\hat P_{\rm hc}\psi = \psi\)), so
    \[
      \hat H_{\rm eff}\psi = \hat P_{\rm hc}\,(H_{\rm hop}\,\psi).
    \]
    The effective dynamics on the hard-core sector is therefore the projected
    hopping; the explicit hopping matrix elements (Tasaki eq.~(11.2.5)),
    which carry Jordan--Wigner fermion signs, are built on this reduction in
    a follow-up layer.
  \item \emph{Range}: \(\hat H_{\rm eff}\psi\in\mathcal H_{\rm hc}\) for every
    \(\psi\) (outermost factor is \(\hat P_{\rm hc}\)), and
    \(\hat P_{\rm hc}\,\hat H_{\rm eff}\psi = \hat H_{\rm eff}\psi\).
\end{enumerate}

\medskip
\noindent\textbf{Lean declarations}
(file \texttt{Fermion/JordanWigner/Hubbard/EffectiveHamiltonian.lean}):
\begin{itemize}
  \item \texttt{hubbardEffectiveHamiltonian N t U}:
    \(\hat H_{\rm eff}\).
  \item \texttt{hubbardEffectiveHamiltonian\_isHermitian}: Hermiticity.
  \item \texttt{hubbardEffectiveHamiltonian\_mulVec\_eq\_projected\_kinetic\_of\_mem}:
    \(\hat H_{\rm eff}\psi = \hat P_{\rm hc}(H_{\rm hop}\psi)\) on
    \(\mathcal H_{\rm hc}\).
  \item \texttt{hubbardEffectiveHamiltonian\_mulVec\_mem}:
    \(\hat H_{\rm eff}\psi\in\mathcal H_{\rm hc}\).
  \item \texttt{hubbardHardcoreProjection\_mulVec\_effectiveHamiltonian}:
    \(\hat P_{\rm hc}\hat H_{\rm eff}\psi = \hat H_{\rm eff}\psi\).
\end{itemize}

%======================================================================
\subsection{Tasaki ordered-creation basis (Tasaki §11.2, eq.~(11.2.3))}
\label{subsec:hubbard-tasaki-basis}
%======================================================================

\noindent\textbf{Definition (11.2.3).}
Tasaki's basis state for the one-hole hard-core sector is
\[
  |\Phi_{x,\sigma}\rangle
    := \hat c_{x,\uparrow}
       \Big(\prod_{y\in\Lambda} \hat c^\dagger_{y,\bar\sigma_y}\Big)
       |\Phi_{\rm vac}\rangle,
  \qquad \bar\sigma_x = \uparrow,
\]
where the creation product runs over all sites in a fixed order and
\(\bar\sigma\) agrees with \(\sigma\) off the hole site \(x\). Every site is
filled, then the spin-\(\uparrow\) electron at the hole \(x\) is annihilated.
This ``redundant'' double action is Tasaki's device for fixing the precise
fermion signs that make the hopping matrix elements (11.2.5) uniformly signed.

\medskip
\noindent\textbf{Reduction to a signed basis vector.}
We take the creation product in \emph{increasing} Jordan--Wigner index order.
The rightmost (largest-index) operator acts on the vacuum first, and each
successive creation operator acts at an index smaller than all currently
occupied modes; hence \emph{every} Jordan--Wigner string sign along the product
equals \(1\), and the ordered product on the vacuum is exactly the
computational basis vector of the all-filled configuration. The single
annihilation at the hole then contributes one nontrivial string sign
\(\varepsilon\in\{\pm 1\}\), so
\[
  |\Phi_{x,\sigma}\rangle
    = \varepsilon \cdot \mathrm{basisVec}\big(\texttt{hubbardOneHoleConfig}\,
      N\,x\,\sigma\big).
\]
Orthonormality of the \(|\Phi_{x,\sigma}\rangle\) is then inherited from the
computational basis vectors, since \(\varepsilon^2 = 1\).

\medskip
\noindent\textbf{Lean declarations}
(file \texttt{Fermion/JordanWigner/Hubbard/TasakiBasis.lean}):
\begin{itemize}
  \item \texttt{occupationOf js}: the configuration occupied exactly on the
    indices listed in \texttt{js}.
  \item \texttt{prod\_creation\_mulVec\_vacuum}: the ordered-creation fold
    lemma --- a strictly index-sorted product of creation operators on the
    vacuum equals \(\mathrm{basisVec}(\texttt{occupationOf js})\), every string
    sign being \(1\).
  \item \texttt{tasakiIndexList}, \texttt{tasakiIndexList\_sorted},
    \texttt{mem\_tasakiIndexList\_iff}: the increasing index list of the
    all-filled spin configuration, its strict sortedness, and its membership
    characterization.
  \item \texttt{hubbardTasakiBasisState N x }\(\sigma\): the (11.2.3) basis
    state.
  \item \texttt{hubbardTasakiBasisState\_eq\_smul\_basisVec}:
    \(|\Phi_{x,\sigma}\rangle = \varepsilon\cdot\mathrm{basisVec}\).
  \item \texttt{hubbardTasakiBasisSign\_mul\_self}: \(\varepsilon^2 = 1\).
  \item \texttt{hubbardTasakiBasisState\_inner} and
    \texttt{hubbardTasakiBasisState\_self\_inner}: orthonormality and
    normalization.
\end{itemize}

%======================================================================
\subsection{Uniform-sign hole-filling action (Tasaki §11.2, eq.~(11.2.4))}
\label{subsec:hubbard-tasaki-hop}
%======================================================================

\noindent\textbf{Statement (11.2.4).}
In the Tasaki basis the hole-filling hopping term acts with a \emph{uniform}
sign:
\[
  \hat c^\dagger_{(x,s)}\,\hat c_{(z,s)}\,|\Phi_{x,\sigma}\rangle
    = -\,|\Phi_{z,\sigma_{z\to x}}\rangle,
\]
where `\(\sigma_{z\to x}\)` is the spin configuration obtained by moving the
electron at the occupied source site `\(z\)` (of spin `\(s\)`) into the hole
`\(x\)`.

\medskip
\noindent\textbf{The fermion sign.}
On the sign-free `basisVec` basis the same term carries a configuration-dependent
Jordan--Wigner sign `\(\mathrm{jwSign}\cdot\mathrm{jwSign}\)`
(file \texttt{HopAction.lean}). Converting both basis states to the Tasaki basis
multiplies this by the two basis signs `\(\varepsilon\)`, and all four signs
combine to the constant `\(-1\)`. Writing
`\(\mathrm{jwSign} = (-1)^{\#\{\text{occupied modes below}\}}\)`, the explicit
counts are
\[
  \varepsilon(x) = (-1)^x,\quad \varepsilon(z) = (-1)^z,\quad
  E_{\rm src} = z - [x<z],\quad E_{\rm tgt} = x - [z<x],
\]
so the total parity exponent is
\[
  x + z + (z - [x<z]) + (x - [z<x]) = 2(x+z) - 1,
\]
which is odd because `\(x \neq z\)` forces exactly one of `\([x<z]\)`,
`\([z<x]\)` to equal `\(1\)`. Hence the combined sign is `\(-1\)`. The basis
sign is `\(\varepsilon = (-1)^x\)`, \emph{independent of} `\(\sigma\)`: there are
exactly `\(x\)` electrons (one per site below `\(x\)`) beneath the hole's
up-orbital `\(2x\)`.

\medskip
\noindent\textbf{Lean declarations}
(file \texttt{Fermion/JordanWigner/Hubbard/TasakiHopAction.lean}):
\begin{itemize}
  \item \texttt{jwSign\_eq\_neg\_one\_pow}: the string sign as
    `\((-1)^{\sum_{k<j}(c\,k)}\)`.
  \item \texttt{sum\_spinful\_reindex}: reindex a sum over the `\(2N+2\)` modes
    into a double site/spin sum.
  \item \texttt{hubbardTasakiBasisSign\_eq}: `\(\varepsilon = (-1)^x\)`.
  \item \texttt{hop\_jwSign\_source}, \texttt{hop\_jwSign\_target}: the two hop
    string signs as parities of occupied-site counts.
  \item \texttt{hubbardTasakiHop\_mulVec}: the uniform-sign action (11.2.4).
\end{itemize}

%======================================================================
\subsection{Effective-Hamiltonian matrix element (Tasaki §11.2, eq.~(11.2.5))}
\label{subsec:hubbard-matrix-element}
%======================================================================

\noindent\textbf{Statement (11.2.5), off-diagonal.}
For `\(x \neq y\)`,
\[
  \langle \Phi_{y,\tau} |\, \hat H_{\rm eff}\, | \Phi_{x,\sigma} \rangle
    = \begin{cases} -t_{x,y} & \text{if } \tau = \sigma_{y\to x},\\ 0 & \text{otherwise.}\end{cases}
\]

\medskip
\noindent\textbf{Proof.}
Since `\(|\Phi_{x,\sigma}\rangle\)` is hard-core, `\(\hat H_{\rm eff}\)` reduces
to the projected kinetic operator `\(\hat P_{\rm hc} H_{\rm hop}\)`. The hard-core
projection `\(\hat P_{\rm hc}\)` is diagonal and Hermitian and fixes the basis
states, so it acts as the identity on the `\(C_{y,\tau}\)` component and drops
out of the pairing, leaving `\(\langle \Phi_{y,\tau} | H_{\rm hop} | \Phi_{x,\sigma}\rangle\)`.
Expanding `\(H_{\rm hop} = \sum_{s,i,j} t_{ij}\, \hat c^\dagger_{(i,s)}\hat c_{(j,s)}\)`,
only the single channel `\((i,j,s) = (x,y,\sigma_y)\)` survives: evaluating the
hopped configuration at the occupied orbitals of `\(y\)` and `\(x\)` forces the
annihilation to occur at `\(y\)` and the creation at the hole `\(x\)`, so every
other channel either annihilates an empty orbital, produces a doubly-occupied
(non-hard-core) configuration, or merely returns the hole to `\(x\)` — none of
which overlaps `\(|\Phi_{y,\tau}\rangle\)`. The surviving channel is the
uniform-sign hole-filling action (11.2.4), `\(\hat c^\dagger_{(x,\sigma_y)}
\hat c_{(y,\sigma_y)} |\Phi_{x,\sigma}\rangle = -|\Phi_{y,\sigma_{y\to x}}\rangle\)`,
so no fermion-sign bookkeeping is repeated and the basis signs `\(\varepsilon =
(-1)^y\)` cancel (`\(\varepsilon^2 = 1\)`), giving `\(-t_{x,y}\)` times the
configuration-matching indicator. (The codebase kinetic uses `\(+t\)`; the
uniform `\(-1\)` of (11.2.4) supplies the minus, matching Tasaki's `\(-t_{x,y}\)`.)

\medskip
\noindent\textbf{Lean declarations}
(file \texttt{Fermion/JordanWigner/Hubbard/EffectiveHamiltonianMatrix.lean}):
\begin{itemize}
  \item \texttt{hardcore\_proj\_apply}: `\(\hat P_{\rm hc}\)` acts as the identity
    on a hard-core component.
  \item \texttt{hop\_config\_forces}: the hole-filling channel is the only one
    whose hop reaches `\(C_{y,\tau}\)`.
  \item \texttt{hubbardEffective\_tasaki\_matrixElement}: the (11.2.5) matrix
    element (off-diagonal, \(x \neq y\)).
  \item \texttt{hubbardEffective\_tasaki\_matrixElement\_diag}: the diagonal
    (\(x = y\)) case under no self-hopping (\(\forall i,\ t_{ii} = 0\)), which
    vanishes since `\(\hat H_{\rm eff}\)` moves the hole off `\(x\)`.
\end{itemize}

\medskip
\noindent\textbf{Diagonal vanishing.}
With no self-hopping the kinetic operator has no diagonal number term, so every
surviving hop moves the hole from \(x\) to a distinct site; the image of
\(|\Phi_{x,\sigma}\rangle\) therefore has no component on the hole-at-\(x\)
configuration \(C_{x,\tau}\). Concretely, a hop \(i \neq j\) empties the
source \(j\) without refilling it (\texttt{hop\_diag\_forces}), so the result is
not a one-hole-at-\(x\) configuration; and the \(i = j\) on-site term carries the
coefficient \(t_{ii} = 0\). Hence \(\langle \Phi_{x,\tau} | \hat H_{\rm eff}
| \Phi_{x,\sigma}\rangle = 0\).

%======================================================================
\subsection{Cauchy--Schwarz energy bound (Tasaki §11.2, eq.~(11.2.9))}
\label{subsec:hubbard-weak-nagaoka}
%======================================================================

\noindent\textbf{Statement (11.2.9).}
For non-negative hopping (\(t_{x,y} \geq 0\)) and no self-hopping
(\(t_{i,i}=0\)), the effective-Hamiltonian energy of any real Tasaki-basis
expansion \(|\Phi\rangle = \sum_{x,\sigma}\phi(x,\sigma)|\Phi_{x,\sigma}\rangle\)
is bounded below by that of the ferromagnetic state
\[
  |\Phi_\uparrow\rangle = \sum_x \xi_x\, |\Phi_{x,(\uparrow)}\rangle,
  \qquad \xi_x = \Big(\sum_\sigma \phi(x,\sigma)^2\Big)^{1/2},
\]
namely \(\langle \Phi_\uparrow | \hat H_{\rm eff} | \Phi_\uparrow\rangle
\leq \langle \Phi | \hat H_{\rm eff} | \Phi\rangle\). Hence the ferromagnetic
state is also a ground state.

\medskip
\noindent\textbf{Proof.}
Expanding the energy bilinearly and substituting the matrix element (11.2.5)
(diagonal terms vanish by no self-hopping) gives the quadratic form
\[
  \langle \Phi | \hat H_{\rm eff} | \Phi\rangle
    = -\sum_{x \neq y} t_{y,x} \sum_\sigma \phi(x,\sigma)\,
      \phi(y,\sigma_{x\to y}),
\]
where the inner indicator sum collapses because, among the canonical hole-spin
configurations \(\tau\) at \(y\), exactly the hole-move image \(\sigma_{x\to y}\)
matches the moved configuration (\texttt{oneHoleConfig\_move\_eq\_iff}). The
spin configurations on \(\Lambda\setminus\{x\}\) are encoded as the subtype
\texttt{HoleSpin}, on which \(\sigma\mapsto\sigma_{x\to y}\) is the bijection
\texttt{holeSpinMoveEquiv}; this reindexing yields
\(\sum_\sigma \phi(y,\sigma_{x\to y})^2 = \xi_y^2\). Cauchy--Schwarz
(\texttt{Finset.sum\_mul\_sq\_le\_sq\_mul\_sq}) then gives
\(\sum_\sigma \phi(x,\sigma)\phi(y,\sigma_{x\to y}) \leq \xi_x \xi_y\), and since
\(t_{y,x}\geq 0\) each off-diagonal term of \(|\Phi\rangle\) dominates the
corresponding term of \(|\Phi_\uparrow\rangle\) (whose quadratic form is
\(-\sum_{x\neq y} t_{y,x}\xi_x\xi_y\)).

\medskip
\noindent\textbf{Lean declarations}
(file \texttt{Fermion/JordanWigner/Hubbard/WeakNagaoka.lean}):
\begin{itemize}
  \item \texttt{HoleSpin}, \texttt{holeSpinMoveEquiv}: canonical hole-spin
    configurations and the hole-move bijection.
  \item \texttt{hubbardEffEnergy\_tasaki\_quadratic}: the quadratic-form
    expansion of the energy.
  \item \texttt{oneHoleConfig\_move\_eq\_iff}: the hole-move configuration
    bridge.
  \item \texttt{tasakiQuadForm\_ferro\_le}: the Cauchy--Schwarz bound on the
    real quadratic form.
  \item \texttt{hubbardWeakNagaoka\_energy\_bound}: the (11.2.9) energy
    inequality.
\end{itemize}

\medskip
\noindent\textbf{Norm identity (eqs.~(11.2.7)--(11.2.8)).}
The Tasaki basis states (indexed by hole site and canonical hole-spin) are
orthonormal (\texttt{tasakiState\_orthonormal}), so the squared norm of a
real expansion is the sum of squared coefficients,
\(\|\sum_p \phi_p \Phi_p\|^2 = \sum_p \phi_p^2\)
(\texttt{tasakiExpansion\_normSq}). Since
\(\sum_p (\text{ferro coeff})_p^2 = \sum_x \xi_x^2 = \sum_p \phi_p^2\)
(\texttt{ferroCoeff\_normSq\_eq}), the ferromagnetic state has the same norm,
\(\|\Phi_\uparrow\| = \|\Phi\|\) (\texttt{tasakiFerro\_normSq\_eq}). Together
with the (11.2.9) energy bound, this gives the quantitative weak-Nagaoka
content: \(|\Phi_\uparrow\rangle\) is a normalized state of the same norm with
no greater energy, hence is also a ground state.

%======================================================================
\subsection{SU(2) symmetry of the effective Hamiltonian (Tasaki §11.2)}
\label{subsec:hubbard-effham-su2}
%======================================================================

\noindent\textbf{Statement.}
The effective Hamiltonian inherits the full SU(2) spin symmetry of the Hubbard
model: \([\hat H_{\rm eff}, \hat S^\pm_{\rm tot}] = 0\).

\medskip
\noindent\textbf{Proof.}
Since \(\hat H_{\rm eff} = \hat P_{\rm hc}\hat H\hat P_{\rm hc}\) and
\([\hat H, \hat S^+_{\rm tot}] = 0\) (Tasaki §9.3.3), it suffices that
\(\hat S^+_{\rm tot}\) commutes with the hard-core projection. Each raising term
\(\hat c^\dagger_{k,\uparrow}\hat c_{k,\downarrow}\) commutes with every
double-occupancy operator \(n_{i,\uparrow}n_{i,\downarrow}\) (same site: both
products vanish by Pauli exclusion; distinct sites: the constituents commute),
hence with each hard-core factor \(1 - n_{i,\uparrow}n_{i,\downarrow}\), hence
with their non-commutative product \(\hat P_{\rm hc}\). The lowering symmetry
follows by taking adjoints, using that \(\hat H_{\rm eff}\) is Hermitian and
\((\hat S^+_{\rm tot})^\dagger = \hat S^-_{\rm tot}\). This SU(2) symmetry is the
structural reason for the \((2S_{\max}+1)\)-fold degeneracy of the weak Nagaoka
ground states (Theorem 11.5): \(\hat S^-_{\rm tot}\) maps a ferromagnetic ground
state to the full degenerate spin multiplet.

\medskip
\noindent\textbf{Lean declarations}
(file \texttt{Fermion/JordanWigner/Hubbard/EffectiveHamiltonianSpinSymmetry.lean}):
\texttt{fermionTotalSpinPlus\_commute\_hubbardHardcoreProjection},
\texttt{fermionTotalSpinPlus\_commute\_hubbardEffectiveHamiltonian},
\texttt{fermionTotalSpinMinus\_commute\_hubbardEffectiveHamiltonian}.

\medskip
\noindent\textbf{Conservation of total spin.}
The same argument applied to \(\hat S^z_{\rm tot}\) (which commutes with every
double-occupancy operator, all number operators commuting) gives
\([\hat H_{\rm eff}, \hat S^z_{\rm tot}] = 0\), and combining with the
raising/lowering versions yields \([\hat H_{\rm eff}, (\hat S_{\rm tot})^2] = 0\):
the effective Hamiltonian conserves total spin, so its eigenspaces split into
fixed-\(S_{\rm tot}\) sectors — the conservation law behind the
\(S_{\rm tot} = S_{\max}\) statement of Theorem 11.5 and the per-sector
Perron--Frobenius analysis of Theorem 11.7. (\texttt{fermionTotalSpinZ\_commute\_hubbardEffectiveHamiltonian},
\texttt{fermionTotalSpinSquared\_commute\_hubbardEffectiveHamiltonian}.)

%======================================================================
\subsection{Weak Nagaoka spin multiplet (Tasaki §11.2.1, Theorem 11.5 core)}
\label{subsec:weak-nagaoka-multiplet}
%======================================================================

\noindent\textbf{Reference.} Tasaki, \emph{Physics and Mathematics of Quantum
Many-Body Systems}, 1st edition, §11.2.1, Theorem 11.5, pp.~382--385; the SU(2)
ladder algebra is §9.3.3, p.~332.

\medskip
\noindent\textbf{SU(2) ladder algebra.}
On the spinful Jordan--Wigner chain we prove the standard \(\mathfrak{su}(2)\)
relations for the total-spin operators directly from the per-site fermion
commutators \([\hat S^+_{\rm tot}, c^\dagger_{j,\downarrow}] = c^\dagger_{j,\uparrow}\)
and \([\hat S^+_{\rm tot}, c_{j,\uparrow}] = -c_{j,\downarrow}\):
\[
  [\hat S^+_{\rm tot}, \hat S^-_{\rm tot}] = 2\,\hat S^z_{\rm tot},
  \qquad
  [(\hat S_{\rm tot})^2, \hat S^-_{\rm tot}] = 0,
  \qquad
  \hat S^+_{\rm tot}\hat S^-_{\rm tot}
    = (\hat S_{\rm tot})^2 - \hat S^z_{\rm tot}(\hat S^z_{\rm tot} - 1).
\]
The first is summed from the per-site contribution
\([\hat S^+_{\rm tot}, c^\dagger_{j,\downarrow}c_{j,\uparrow}]
   = n_{j,\uparrow} - n_{j,\downarrow}\); the second follows by expanding
\([(\hat S_{\rm tot})^2, \hat S^-_{\rm tot}]\) with the Casimir definition
\((\hat S_{\rm tot})^2 = \hat S^-_{\rm tot}\hat S^+_{\rm tot}
   + \hat S^z_{\rm tot}(\hat S^z_{\rm tot}+1)\), both sides reducing to
\(\hat S^-\hat S^-\hat S^+ + \hat S^-\hat S^z\hat S^z + \hat S^-\hat S^z\); the
third combines the Casimir definition with the ladder commutator.

\medskip
\noindent\textbf{The spin-lowering tower.}
Let \(v\) be a maximal-\(\hat S^z\) highest-weight state, i.e.
\(\hat S^+_{\rm tot} v = 0\) and \(\hat S^z_{\rm tot} v = (N/2)\,v\); then
\((\hat S_{\rm tot})^2 v = \tfrac N2(\tfrac N2 + 1)\,v\) automatically (the all-up
sector is entirely \(S_{\rm tot} = S_{\max} = N/2\)). For the descendants
\(v_k := (\hat S^-_{\rm tot})^k v\) the ladder algebra gives, for all
\(k = 0, \dots, N\),
\[
  \hat S^z_{\rm tot}\, v_k = \bigl(\tfrac N2 - k\bigr) v_k,
  \qquad
  (\hat S_{\rm tot})^2\, v_k = \tfrac N2(\tfrac N2+1)\, v_k,
  \qquad
  \hat S^+_{\rm tot}\hat S^-_{\rm tot}\, v_k = (k+1)(N-k)\, v_k.
\]
The last identity is purely algebraic and forces \(v_k \neq 0\) for \(k \le N\):
if \(v_{k+1} = \hat S^-_{\rm tot} v_k\) vanished then
\(\hat S^+_{\rm tot} v_{k+1} = (k+1)(N-k)\, v_k\) would vanish too, impossible
since \((k+1)(N-k) \neq 0\) for \(k < N\) and \(v_k \neq 0\) by induction (no
inner product is needed). Having distinct \(\hat S^z\) eigenvalues \(N/2 - k\),
the \(N+1\) nonzero states \(v_0, \dots, v_N\) are linearly independent.

\medskip
\noindent\textbf{Theorem 11.5 core.}
If in addition \(v\) is an \(\hat H_{\rm eff}\)-eigenvector at energy \(E\) (e.g.\
a ferromagnetic ground state), then \([\hat H_{\rm eff}, \hat S^-_{\rm tot}] = 0\)
makes every \(v_k\) an \(\hat H_{\rm eff}\)-eigenvector at the \emph{same} energy
\(E\). Thus a single ferromagnetic ground state generates
\(N+1 = 2S_{\max} + 1\) linearly independent degenerate ground states, all with
\(S_{\rm tot} = S_{\max} = N/2\) — the weak Nagaoka conclusion (eq.~(11.2.9)).

\medskip
\noindent\textbf{Existence of the ground state (the all-up block).}
To produce such a highest-weight eigenvector we represent \(\hat H_{\rm eff}\) on
the orthonormal one-hole Tasaki basis by the finite Hermitian matrix
\(M = T^\dagger \hat H_{\rm eff} T\) (columns of \(T\) are the basis states
\(|\Phi_p\rangle\)). The \emph{operator lift}
\(\hat H_{\rm eff}|\Phi_p\rangle = \sum_q \langle\Phi_q|\hat H_{\rm eff}|\Phi_p\rangle\,
|\Phi_q\rangle\) holds because \(\hat H_{\rm eff}|\Phi_p\rangle\) is hard-core
(\(P_{\rm hc}\)) and an \(N\)-electron eigenstate (\([\hat H_{\rm eff},\hat N]=0\)),
hence supported on one-hole hard-core configurations where the Tasaki basis is
complete. Consequently \(M\)-eigenvectors lift to \(\hat H_{\rm eff}\)-eigenvectors.
Since hopping the all-up hole keeps the configuration all-up, \(M\) leaves the
all-up block \(M_\uparrow\) invariant; a minimum eigenvector \(\xi\) of the
Hermitian \(M_\uparrow\) gives the all-up state
\(|\Phi_\uparrow\rangle = \sum_x \xi_x |\Phi_{x,\uparrow}\rangle\), which is a
nonzero highest-weight \(\hat H_{\rm eff}\)-eigenvector at the all-up sector
minimum. Feeding it to the spin-multiplet argument above yields
\(N+1 = 2S_{\max}+1\) linearly independent degenerate eigenstates with
\(S_{\rm tot} = S_{\max}\) — \textbf{Tasaki Theorem 11.5}. (No Perron--Frobenius is
used; that this all-up minimum is the global ground energy is the Schwarz bound
(11.2.9), and PF uniqueness belongs to Theorem 11.7.) A general Rayleigh--Ritz
equality lemma — a nonzero vector attaining the minimum-eigenvalue Rayleigh bound
is a minimum eigenvector — is provided as
\texttt{mulVec\_eq\_smul\_of\_rayleighOnVec\_eq\_min}.

\medskip
\noindent\textbf{The all-up minimum is the global minimum (Schwarz bound).}
The multiplet above sits at the all-up sector minimum \(\min(M_\uparrow)\). For
\(t \ge 0\) this equals the global one-hole minimum \(\min(M)\), so the multiplet
consists of genuine ground states. Indeed \(\min(M) \le \min(M_\uparrow)\) is the
principal-submatrix inequality (embed a unit all-up eigenvector); and
\(\min(M_\uparrow) \le \min(M)\) is Tasaki's Schwarz bound (11.2.9): the Tasaki
matrix is real (\(M = M_{\mathbb R}\,\)cast, since its entries are
\(-t_{x,y}\) times an indicator), so a global minimum eigenvector can be taken
real, and ferromagnetizing it (\texttt{ferroCoeff}, \(\sum\xi_x^2 = \sum\phi_p^2\))
does not raise the energy (\texttt{tasakiQuadForm\_ferro\_le}). Hence
\(\min(M_\uparrow) = \min(M)\), giving the global statement
\texttt{weakNagaoka\_theorem\_11\_5\_global}
(file \texttt{Fermion/JordanWigner/Hubbard/WeakNagaokaGlobalMin.lean}).

\medskip
\noindent The Lean development
(files \texttt{Fermion/JordanWigner/Hubbard/WeakNagaokaTheorem.lean},
\texttt{Fermion/JordanWigner/Hubbard/WeakNagaokaGroundState.lean}):
\texttt{fermionTotalSpinPlus\_commutator\_fermionTotalSpinMinus},
\texttt{fermionTotalSpinSquared\_commute\_fermionTotalSpinMinus},
\texttt{fermionTotalSpinPlus\_mul\_fermionTotalSpinMinus},
\texttt{spinMinusPow\_ne\_zero}, \texttt{spinMinusPow\_linearIndependent},
the capstone \texttt{weakNagaoka\_spinMultiplet}; and for the ground state
\texttt{tasakiEffMatrix}, \texttt{tasakiEffMatrixUp}, the operator lift
\texttt{hubbardEffectiveHamiltonian\_mulVec\_tasakiState}, \texttt{tasaki\_completeness},
and the final theorem \texttt{weakNagaoka\_theorem\_11\_5}.

%======================================================================
\subsection{Nagaoka's theorem and the connectivity condition
  (Tasaki §11.2.2, Definition 11.6 and Theorem 11.7)}
\label{subsec:nagaoka-theorem-11-7}
%======================================================================

\noindent This formalizes Tasaki, \emph{Physics and Mathematics of Quantum
Many-Body Systems}, 1st edition, §11.2.2, Definition 11.6 and Theorem 11.7,
pp.~385--388 (file
\texttt{Fermion/JordanWigner/Hubbard/NagaokaConnectivity.lean}).

\medskip
\noindent\textbf{Magnetization sectors.}
The one-hole Tasaki basis index \(I = \bigsqcup_x \{\sigma : \sigma_x = \uparrow\}\)
is graded by the (doubled) total \(S_z^{(3)} = \sum_{y \ne x}\sigma_y\), read off
the occupation configuration as \texttt{configMag}/\texttt{holeSpinMag}. Hole
hopping (11.2.4) preserves this grading, so the Tasaki matrix \(M\) (the
(11.2.5) matrix of \(\hat H_{\rm eff}\)) is \textbf{block-diagonal} across
magnetization sectors: \(M_{q,p} = 0\) when \(q,p\) have different magnetization
(\texttt{tasakiEffReMatrix\_eq\_zero\_of\_holeSpinMag\_ne}). There are exactly
\(N+1\) non-empty sectors, since \(\text{mag} = 2(\#\!\uparrow) - N\) with
\(\#\!\uparrow \in \{0,\dots,N\}\) (\texttt{holeSpinMag\_image\_card\_le}).

\medskip
\noindent\textbf{Definition 11.6 (connectivity condition).}
\texttt{nagaokaConnectivity} \(:= \forall m,\ (-M_m)\) is irreducible (in the
Perron--Frobenius sense), where \(M_m\) is \(M\) restricted to the fixed-\(m\)
sector. Since \(M_{q,p} = -t_{x,y}\cdot[\text{indicator}] \le 0\) off the
diagonal (for \(t \ge 0\)) and \(M\) vanishes on the diagonal, \(-M_m\) is
entrywise non-negative (condition (i) of Perron--Frobenius), and irreducibility
is condition (ii).

\medskip
\noindent\textbf{Theorem 11.7 (Nagaoka).}
For \(t \ge 0\), \(N = |\Lambda|-1\), \(U = \infty\) and the connectivity
condition, the one-hole \(\hat H_{\rm eff}\) ground eigenspace has dimension
exactly \(N+1 = 2S_{\max}+1\) and every ground state has \(S_{\rm tot} = S_{\max}\).
The proof matches the book's: apply Perron--Frobenius to \(-M_m\) on each fixed
\(S_z^{(3)}\) sector to get a \emph{non-degenerate} sector ground state, then
combine with Theorem 11.5.

\emph{Upper bound} \((\le N+1)\): on a non-empty connected sector,
\texttt{exists\_pos\_eigenvec\_max} gives a strictly positive Perron eigenvector
of \(-M_m\) and \texttt{eigenspace\_finrank\_le\_one\_of\_pos\_eigenvec} a
one-dimensional eigenspace; Collatz--Wielandt
(\texttt{hermitianMinEigenvalue\_lift\_eq\_sub\_pf}, shift \(c=0\)) identifies the
sector minimum as \(-\mu\), and the global minimum is \(\le\) every sector minimum
(a zero-extended sector eigenvector is a genuine full eigenvector). So at the
global minimum each sector contributes \(\le 1\); restricting a ground vector to
its sectors embeds the ground eigenspace into the product over the \(\le N+1\)
sectors (\texttt{tasakiEffMatrix\_ground\_finrank\_le\_N\_add\_one}).

\emph{Lower bound} \((\ge N+1)\): the SU(2) ferromagnetic tower
\((\hat S^-_{\rm tot})^k|\Phi_\uparrow\rangle\) (\(k = 0,\dots,N\)) of Theorem 11.5
consists of \(N+1\) linearly independent ground states. \(\hat S^-_{\rm tot}\)
conserves the particle number and preserves the hard-core subspace, so the tower
stays one-hole supported (\texttt{spinMinusPow\_pfFerroState\_oneHole\_supported});
its Tasaki coefficients are then \(N+1\) linearly independent \(M\)-eigenvectors
(\texttt{tasakiEffMatrix\_ground\_finrank\_ge}).

Combining gives \texttt{nagaoka\_theorem\_11\_7\_degeneracy} (degeneracy
\(= N+1\)) and, since the tower spans the ground eigenspace, every ground state
is a tower combination, hence a \((\hat S_{\rm tot})^2\)-eigenstate at
\(S_{\max}(S_{\max}+1)\) (\texttt{nagaoka\_theorem\_11\_7}). The backward bridge
\texttt{tasakiCoeff\_mulVec\_eigen\_of\_full} (a one-hole-supported full
eigenvector at \(E\) yields a coefficient \(M\)-eigenvector at \(E\)) is the
technical crux. The development is sorry-free.

\medskip
\noindent\emph{Definition 11.6 is here taken as per-sector irreducibility of
\(-M\); the book's equivalent ``connected through non-vanishing matrix elements
of \(\hat H_{\rm eff}\)'' is the same Perron--Frobenius notion.}

\medskip
\noindent\textbf{Theorem 11.8 and Lemma 11.9 (axiomatized, deferred).}
Tasaki §11.2.2 then gives two criteria for \emph{verifying} the connectivity
condition on a concrete bond graph \((\Lambda, B)\), \(B = \{\{x,y\} : t_{x,y}
\ne 0\}\): \textbf{Theorem 11.8} (Bobrow--Stubis--Li) — connectivity holds
\emph{iff} \((\Lambda, B)\) is biconnected and is not a simple loop (periodic
chain) with more than four sites; and \textbf{Lemma 11.9} — if \((\Lambda, B)\)
is connected by \emph{exchange bonds} (pairs lying on a common length-3/4 loop
whose removal keeps \(\Lambda\) connected), connectivity holds. In the Lean
development (file
\texttt{Fermion/JordanWigner/Hubbard/NagaokaConnectivityClassification.lean})
these are recorded as \textbf{\texttt{axiom}s}
(\texttt{nagaoka\_theorem\_11\_8}, \texttt{nagaoka\_lemma\_11\_9}), with the
bond-graph and biconnected / simple-loop / exchange-bond predicates defined
(\texttt{nagaokaBondGraph}, \texttt{IsBiconnected}, \texttt{IsSimpleLoopGTFour},
\texttt{IsExchangeBond}, \texttt{ConnectedByExchangeBonds}), and their proofs
\emph{deferred to a future implementation}. The rationale: Tasaki
\emph{explicitly leaves the proof of Theorem 11.8 to the original papers}
(Bobrow--Stubis--Li; Wilson's ``15-puzzle'' theorem, p.~387 refs [4], [81]) — it
is a cited external classification theorem; and a faithful Lean proof of Lemma
11.9 must turn its hole-motion (``15-puzzle'') argument into strong connectivity
of the per-sector quiver of \(-M\), a multi-PR effort in the present
matrix/configuration framework. \textbf{Theorem 11.7 above is \texttt{sorry}-free
and does not depend on these axioms}; they are isolated in a separate file so the
proven core stays axiom-clean.

%======================================================================
\subsection{Tasaki's flat-band ferromagnetism: model and Lemma 11.10
  (Tasaki §11.3.1)}
\label{subsec:tasaki-flatband}
%======================================================================

\noindent This formalizes the model setup of Tasaki §11.3.1 (the $d=1$
decorated/Delta chain) and Lemma 11.10 (files
\texttt{Fermion/JordanWigner/Hubbard/TasakiFlatBandModel.lean} and
\texttt{TasakiFlatBandBasis.lean}).

\medskip
\noindent\textbf{The model.}
The decorated chain $\Lambda = E \cup I$ has $K+1$ external sites $E$ (the
periodic chain) and $K+1$ internal sites $I$ (bond midpoints); it is realized in
the existing spinful Hubbard framework on $2K+2$ physical sites by interleaving
external $i \mapsto 2i$ (\texttt{deltaExternalSite}) and internal $i \mapsto 2i+1$
(\texttt{deltaInternalSite}) into $\mathrm{Fin}(2K+2)$. The localized
single-particle states are $\alpha_p$ (\texttt{flatBandAlpha}, eq.~(11.3.1):
$1$ at external $p$, $-\nu$ at the incident internal sites $p, p-1$) and $\beta_u$
(\texttt{flatBandBeta}, eq.~(11.3.2): $1$ at internal $u$, $+\nu$ at external
$u, u+1$); the fermion operators $\hat a_{p,\sigma} = \hat C_\sigma(\alpha_p)$,
$\hat b_{u,\sigma} = \hat C_\sigma(\beta_u)$ (eqs.~(11.3.3)/(11.3.4)) are the
corresponding combinations of $\hat c$'s. The Hamiltonian is
\texttt{flatBandHamiltonian} $= t\sum_{u,\sigma}\hat b^\dagger_{u,\sigma}
\hat b_{u,\sigma} + U\sum_x \hat n_{x\uparrow}\hat n_{x\downarrow}$
(eqs.~(11.3.5)/(11.3.6)), shown Hermitian.

\medskip
\noindent\textbf{Lemma 11.10.}
$\{\alpha_p\}_{p\in E} \cup \{\beta_u\}_{u\in I}$ is a basis of the
single-particle Hilbert space $\mathrm{Fin}(2K+2)\to\mathbb C$
(\texttt{flatBand\_linearIndependent}, \texttt{flatBandBasis}). Following Tasaki:
the diagonal evaluations $\alpha_p(2q) = [p=q]$, $\beta_u(2v{+}1) = [u=v]$ give the
individual linear independence; the even/odd site split
(\texttt{deltaSiteEquiv}) and the cross-orthogonality
$\langle\alpha_p,\beta_u\rangle = 0$ (\texttt{flatBandAlpha\_dot\_beta} --- the
$+\nu$ contribution at the shared external site cancels the $-\nu$ at the shared
internal site, via the cyclic incidence equivalence $p\in\{u,u{+}1\}
\Leftrightarrow u\in\{p,p{-}1\}$) make the two spans orthogonal, hence disjoint;
\texttt{linearIndependent\_sum} then gives the combined independence, and the
basis follows from $|E|+|I| = |\Lambda| = \dim$. The development is sorry-free.

\medskip
\noindent\textbf{The anticommutator $\{\hat b_{u,\sigma}, \hat a^\dagger_{p,\tau}\}
= 0$ (eq.~(11.3.7)).}
The $\hat b$ and $\hat a^\dagger$ operators anticommute
(\texttt{flatBandBAnnihilation\_ACreation\_anticomm}): expanding the
anticommutator of the two linear combinations bilinearly and using the spinful
canonical anticommutation relation $\{\hat c_{x,\sigma}, \hat c^\dagger_{y,\tau}\}
= \delta_{x,y}\,\delta_{\sigma,\tau}$
(\texttt{spinful\_annihilation\_creation\_anticomm}, from
\texttt{spinfulIndex} injectivity and the single-species CAR) reduces it, for
$\sigma=\tau$, to the cross-orthogonality
$\langle\alpha_p,\beta_u\rangle = 0$, and for $\sigma\ne\tau$ every term vanishes.
This is the algebraic input to the frustration-free analysis of Theorem 11.11.

\medskip
\noindent\textbf{The all-up $\alpha$ state and $\hat H_{\rm hop}|\Phi_{\alpha,
\mathrm{all}\uparrow}\rangle = 0$ (eqs.~(11.3.8)/(11.3.9)).}
The candidate ferromagnetic ground state is the all-up $\alpha$ Slater state
$|\Phi_{\alpha,\mathrm{all}\uparrow}\rangle = \bigl(\prod_p \hat a^\dagger_{p,
\uparrow}\bigr)|\mathrm{vac}\rangle$ (\texttt{flatBandAlphaAllUpState}, eq.~(11.3.9)).
A move-through lemma (\texttt{flatBand\_anticomm\_listProd\_mulVec\_vacuum}, by
list induction) shows that an operator $B$ that annihilates the vacuum and
anticommutes with each factor of an ordered product also annihilates the product
applied to the vacuum. Applied with $B = \hat b_{u,\sigma}$ (which annihilates
the vacuum and, by eq.~(11.3.7), anticommutes with every $\hat a^\dagger_{p,
\uparrow}$) it gives $\hat b_{u,\sigma}|\Phi_{\alpha,\mathrm{all}\uparrow}\rangle
= 0$ (the zero-energy condition (11.3.11) at the all-up state). Hence the hopping
Hamiltonian $\hat H_{\rm hop} = t\sum_{u,\sigma}\hat b^\dagger_{u,\sigma}
\hat b_{u,\sigma}$ (\texttt{flatBandHopping}) annihilates it:
$\hat H_{\rm hop}|\Phi_{\alpha,\mathrm{all}\uparrow}\rangle = 0$
(\texttt{flatBandHopping\_mulVec\_alphaAllUpState}, eq.~(11.3.8)) --- it is a
zero-energy state of the (positive-semidefinite) kinetic term.

\medskip
\noindent\textbf{$\hat H |\Phi_{\alpha,\mathrm{all}\uparrow}\rangle = 0$.}
The all-up state has no down-electrons, so $\hat c_{x,\downarrow}
|\Phi_{\alpha,\mathrm{all}\uparrow}\rangle = 0$ (the same move-through argument
with a down annihilation, which anticommutes with every $\hat a^\dagger_{p,
\uparrow}$); hence the on-site interaction $\hat H_{\rm int} = U\sum_x
\hat n_{x\uparrow}\hat n_{x\downarrow}$ annihilates it
(\texttt{hubbardDoubleOccupancy\_mulVec\_alphaAllUpState}). Together with
$\hat H_{\rm hop}|\Phi_{\alpha,\mathrm{all}\uparrow}\rangle = 0$ this gives
$\hat H|\Phi_{\alpha,\mathrm{all}\uparrow}\rangle = 0$
(\texttt{flatBandHamiltonian\_mulVec\_alphaAllUpState}): since
$\hat H = \hat H_{\rm hop} + \hat H_{\rm int}$ is a sum of positive-semidefinite
terms, $|\Phi_{\alpha,\mathrm{all}\uparrow}\rangle$ is a zero-energy ground state
(the existence half of Theorem 11.11).

\medskip
\noindent\textbf{General highest-weight SU(2) lowering tower
(\texttt{SpinLoweringTowerGeneral.lean}).}
To turn the single zero-energy state $|\Phi_{\alpha,\mathrm{all}\uparrow}\rangle$
into the full ferromagnetic multiplet we lower it with $\hat S^-_{\rm tot}$.
The SU(2) ladder lemmas already formalized for weak Nagaoka
(§\ref{subsec:weak-nagaoka-multiplet}) hardcode the highest weight to the \emph{chain
maximum} $N/2$, i.e.\ they assume $\hat S^z_{\rm tot}v = (N/2)v$. Tasaki's
flat-band ferromagnet has only $|E| = K+1$ electrons on $2K+2$ sites
($N = 2K+1$), so its highest-weight state carries
$\hat S^z_{\rm tot} = (K+1)/2 < N/2$; the $N/2$-specialized lemmas do not apply.
We therefore re-derive the same ladder algebra at an \emph{arbitrary} highest
weight $m\in\mathbb C$ (Casimir eigenvalue $\lambda\in\mathbb C$):
\[
  \hat S^z_{\rm tot}(\hat S^-_{\rm tot})^k v = (m-k)(\hat S^-_{\rm tot})^k v,
  \qquad
  \hat S^+_{\rm tot}\hat S^-_{\rm tot}(\hat S^-_{\rm tot})^k v
    = \bigl(\lambda-(m-k)(m-k-1)\bigr)(\hat S^-_{\rm tot})^k v,
\]
and, for a highest-weight state ($\hat S^+_{\rm tot}v = 0$,
$\hat S^z_{\rm tot}v = mv$), the Casimir value
$(\hat S_{\rm tot})^2 v = m(m+1)v$
(\texttt{fermionTotalSpinSquared\_mulVec\_of\_isTop\_general}). Specializing to a
step count $L\in\mathbb N$ with highest weight $m = L/2$, the factorization
$\lambda - (m-k)(m-k-1) = (k+1)(L-k)$ shows the ladder terminates only after $L$
steps: the $L+1$ lowered states $(\hat S^-_{\rm tot})^k v$ ($k = 0,\dots,L$) are
nonzero (\texttt{spinMinusPow\_ne\_zero\_general}) and linearly independent
(\texttt{spinMinusPow\_linearIndependent\_general}, distinct $\hat S^z$
eigenvalues $L/2-k$). The packaged
\texttt{highestWeight\_spinMultiplet\_general} states that a nonzero
highest-weight state generates an $(L+1)$-dimensional multiplet, all members in
the maximal-spin sector $(\hat S_{\rm tot})^2 = (L/2)(L/2+1)$ --- the SU(2) input
to the existence half of Theorem 11.11 at the correct electron-count highest
weight $L = K+1$.

\medskip
\noindent\textbf{The all-up $\alpha$ state is a highest-weight maximal-spin state
(\texttt{TasakiFlatBandHighestWeight.lean}, eq.~(11.3.10)).}
To feed $|\Phi_{\alpha,\mathrm{all}\uparrow}\rangle$ into the general tower we
verify its highest-weight data. The raising operator
$\hat S^+_{\rm tot} = \sum_i \hat c^\dagger_{i,\uparrow}\hat c_{i,\downarrow}$
annihilates it because each down annihilation does
($\hat c_{i,\downarrow}|\Phi_\alpha\rangle = 0$), giving
$\hat S^+_{\rm tot}|\Phi_\alpha\rangle = 0$
(\texttt{flatBandTotalSpinPlus\_mulVec\_alphaAllUpState}). For the $\hat S^z$
weight we count particles: a \emph{charge move-through} lemma
(\texttt{flatBand\_charge\_listProd\_mulVec\_vacuum}, list induction) shows that
if $Q$ annihilates the vacuum and each factor raises its charge by one
($Q\,A_p = A_p Q + A_p$), then $Q$ acts on $\bigl(\prod_p A_p\bigr)|\mathrm{vac}
\rangle$ with eigenvalue equal to the number of factors. Applied with
$Q = \hat N_\uparrow$ and the commutator
$[\hat N_\uparrow, \hat a^\dagger_{p,\uparrow}] = \hat a^\dagger_{p,\uparrow}$
(each $\hat a^\dagger_{p,\uparrow}$ being a linear combination of up creations) it
gives $\hat N_\uparrow|\Phi_\alpha\rangle = (K+1)|\Phi_\alpha\rangle$
(\texttt{flatBandTotalUpNumber\_mulVec\_alphaAllUpState}); and
$\hat N_\downarrow|\Phi_\alpha\rangle = 0$ since there are no down electrons.
Hence
\[
  \hat S^z_{\rm tot}|\Phi_\alpha\rangle
    = \tfrac12(\hat N_\uparrow - \hat N_\downarrow)|\Phi_\alpha\rangle
    = \tfrac{K+1}{2}\,|\Phi_\alpha\rangle
\]
(\texttt{flatBandTotalSpinZ\_mulVec\_alphaAllUpState}, eq.~(11.3.10)): the highest
weight is $m = (K+1)/2 = |E|/2$, the half-filled-band ferromagnetic weight,
strictly below the chain maximum $N/2 = (2K+1)/2$. Together with
$\hat S^+_{\rm tot}|\Phi_\alpha\rangle = 0$ these are exactly the hypotheses of
\texttt{highestWeight\_spinMultiplet\_general} (at $L = K+1$), so once
$|\Phi_\alpha\rangle\neq 0$ is established the full $(K+2)$-dimensional
maximal-spin multiplet follows.

\medskip
\noindent\textbf{Nonvanishing $|\Phi_{\alpha,\mathrm{all}\uparrow}\rangle\neq 0$
(\texttt{TasakiFlatBandNonvanishing.lean}).}
The remaining input is that the candidate ground state does not vanish. We avoid
all Slater-determinant / Gram machinery by using the \emph{triangular} structure
of the orbitals: on the external sites the $\alpha_p$ are exactly the unit
vectors, $\alpha_p(\texttt{deltaExternalSite}\ q) = \delta_{pq}$. Hence the
external up-site annihilation $\hat c_{2q,\uparrow}$ is the \emph{canonical dual}
of $\hat a^\dagger_{q,\uparrow}$:
\[
  \{\hat c_{2q,\uparrow},\, \hat a^\dagger_{p,\uparrow}\}
    = \alpha_p(2q) = \delta_{pq}
\]
(\texttt{flatBandExtUpAnnihilation\_ACreation\_anticomm}, from spinful CAR). A
\emph{peel} lemma (\texttt{flatBand\_extAnnihilation\_peel}): if $c$ kills the
vacuum, anticommutes to $1$ with $A$ and to $0$ with every factor of a product
$P$, then $c\,(A\,P)|\mathrm{vac}\rangle = P|\mathrm{vac}\rangle$ --- the dual
annihilation removes the leading creation. Peeling one factor at a time
(induction over the duplicate-free index list,
\texttt{flatBandAlpha\_listProd\_exists\_collapse}) produces an operator $D$ with
$D|\Phi_\alpha\rangle = |\mathrm{vac}\rangle$; since $|\mathrm{vac}\rangle\neq 0$
(\texttt{fermionMultiVacuum\_ne\_zero}, value $1$ at the empty configuration),
also $|\Phi_\alpha\rangle\neq 0$
(\texttt{flatBandAlphaAllUpState\_ne\_zero}). This is the identity-overlap special
case of the Slater determinant, formalized combinatorially without a determinant.

\medskip
\noindent\textbf{The ferromagnetic spin multiplet
(\texttt{TasakiFlatBandMultiplet.lean}, Theorem 11.11 existence half, spin
content).}
The three prerequisites just established --- the general highest-weight tower,
the highest-weight data of $|\Phi_\alpha\rangle$
($\hat S^+_{\rm tot}|\Phi_\alpha\rangle = 0$,
$\hat S^z_{\rm tot}|\Phi_\alpha\rangle = \tfrac{K+1}{2}|\Phi_\alpha\rangle$), and
$|\Phi_\alpha\rangle\neq 0$ --- are exactly the hypotheses of
\texttt{highestWeight\_spinMultiplet\_general} at $L = K+1 = |E|$. Feeding them in
yields \texttt{flatBand\_ferromagnetic\_multiplet}: the $K+2 = 2S_{\max}+1$ lowered
states $(\hat S^-_{\rm tot})^k|\Phi_\alpha\rangle$ ($k = 0,\dots,K+1$) are linearly
independent and all carry total spin
\[
  (\hat S_{\rm tot})^2\,(\hat S^-_{\rm tot})^k|\Phi_\alpha\rangle
    = S_{\max}(S_{\max}+1)\,(\hat S^-_{\rm tot})^k|\Phi_\alpha\rangle,
  \qquad S_{\max} = \frac{K+1}{2} = \frac{N_e}{2}.
\]
This is the maximal-spin multiplet structure of Tasaki's flat-band ferromagnet.
That every member is moreover a zero-energy \emph{ground} state requires the
energy tower $\hat H(\hat S^-_{\rm tot})^k|\Phi_\alpha\rangle = 0$ (flat-band
SU(2) lowering symmetry) and $\hat H\geq 0$, handled in subsequent steps.

\medskip
\noindent\textbf{Flat-band SU(2) lowering symmetry and the energy tower
(\texttt{TasakiFlatBandEnergyTower.lean}).}
The flat-band Hamiltonian is SU(2) symmetric, $[\hat S^\pm_{\rm tot},\hat H] = 0$.
The interaction $\hat H_{\rm int} = U\sum_x \hat n_{x\uparrow}\hat n_{x\downarrow}$
is the standard Hubbard one and reuses
\texttt{fermionTotalSpinPlus\_commute\_hubbardDoubleOccupancy}. For the kinetic
term $\hat H_{\rm hop} = t\sum_{u,\sigma}\hat b^\dagger_{u,\sigma}\hat b_{u,\sigma}$
the key point is that the $\hat b$ operators form a \emph{spin doublet} (both
spins are built from the same spatial $\beta_u$): lifting the single-particle
covariances $[\hat c_{x,\uparrow},\hat S^+] = \hat c_{x,\downarrow}$,
$[\hat c^\dagger_{x,\downarrow},\hat S^+] = -\hat c^\dagger_{x,\uparrow}$ termwise
gives
\[
  [\hat b_{u,\uparrow},\hat S^+] = \hat b_{u,\downarrow},\qquad
  [\hat b^\dagger_{u,\downarrow},\hat S^+] = -\hat b^\dagger_{u,\uparrow},
\]
while $\hat b^\dagger_{u,\uparrow}$ and $\hat b_{u,\downarrow}$ commute with
$\hat S^+$. Hence the spin-summed mode number
$\hat b^\dagger_{u,\uparrow}\hat b_{u,\uparrow} +
\hat b^\dagger_{u,\downarrow}\hat b_{u,\downarrow}$ commutes with $\hat S^+$ --- the
two off-diagonal pieces $\mp\hat b^\dagger_{u,\uparrow}\hat b_{u,\downarrow}$ cancel
(\texttt{commute\_two\_component\_kinetic\_of\_spinPlus\_relations}). Summing over
modes and adding the interaction gives $[\hat S^+_{\rm tot},\hat H] = 0$, and
$[\hat S^-_{\rm tot},\hat H] = 0$ follows by adjoint ($\hat H$ Hermitian,
$(\hat S^+_{\rm tot})^\dagger = \hat S^-_{\rm tot}$). Since $\hat H$ commutes with
$\hat S^-_{\rm tot}$ and $\hat H|\Phi_\alpha\rangle = 0$, the whole tower is
zero-energy:
\[
  \hat H\,(\hat S^-_{\rm tot})^k|\Phi_\alpha\rangle
    = (\hat S^-_{\rm tot})^k\,\hat H|\Phi_\alpha\rangle = 0
\]
(\texttt{flatBandHamiltonian\_mulVec\_spinMinusPow\_alphaAllUpState}). Combined with
the multiplet, the $K+2$ states are linearly independent, share total spin
$S_{\max} = (K+1)/2$, and are all zero-energy.

\medskip
\noindent\textbf{Positive-semidefiniteness and the ground states
(\texttt{TasakiFlatBandPosSemidef.lean}, Theorem 11.11 existence half).}
For $t, U\geq 0$ the flat-band Hamiltonian is positive-semidefinite,
$\hat H\geq 0$ (\texttt{flatBandHamiltonian\_posSemidef}): it is a nonnegative
combination of positive-semidefinite terms. Each kinetic term factors as
$\hat b^\dagger_{u,\sigma}\hat b_{u,\sigma} = (\hat b_{u,\sigma})^\dagger
\hat b_{u,\sigma}$, manifestly PSD; each interaction term $\hat n_{x\uparrow}
\hat n_{x\downarrow}$ is a Hermitian idempotent projection $P = P^\dagger P$, hence
PSD; and PSD is preserved by sums and nonnegative scalar multiples. Consequently
the energy (Rayleigh quotient) is nonnegative everywhere,
$\langle\psi|\hat H|\psi\rangle\geq 0$
(\texttt{flatBandHamiltonian\_rayleighOnVec\_nonneg}), while the tower states attain
$0$ (\texttt{flatBandHamiltonian\_rayleighOnVec\_spinMinusPow\_alphaAllUpState}).
Therefore every member $(\hat S^-_{\rm tot})^k|\Phi_\alpha\rangle$ of the multiplet
minimizes the energy --- it is a \emph{ground state}
(\texttt{flatBand\_alphaTower\_isGroundState}). Together with
\texttt{flatBand\_ferromagnetic\_multiplet}, this is the existence half of
Theorem 11.11: the $K+2 = 2S_{\max}+1$ states form a maximal-spin
($S_{\rm tot} = S_{\max} = (K+1)/2 = N_e/2$) degenerate \emph{ground-state}
multiplet of Tasaki's flat-band ferromagnet.

\medskip
\noindent\textbf{Frustration-free conditions on a ground state
(\texttt{TasakiFlatBandFrustrationFree.lean}, toward Theorem 11.11 uniqueness).}
The uniqueness half starts from the observation that $\hat H$ is a nonnegative
combination of positive-semidefinite terms with ground energy $0$, so a ground
state must be annihilated by \emph{each} term separately (Tasaki App.~A
Lemmas A.10/A.11). Concretely, the energy quadratic form decomposes as
\[
  \langle v|\hat H|v\rangle
    = t\sum_{u,\sigma}\langle v|\hat b^\dagger_{u,\sigma}\hat b_{u,\sigma}|v\rangle
      + U\sum_x \langle v|\hat n_{x\uparrow}\hat n_{x\downarrow}|v\rangle,
\]
a sum of nonnegative contributions (\texttt{flatBandHamiltonian\_rayleighOnVec\_%
decompose}). For a ground state ($\langle v|\hat H|v\rangle = 0$) with
$t, U > 0$ every contribution vanishes; since each term is positive-semidefinite,
a zero of its quadratic form is a kernel vector
(\texttt{posSemidef\_mulVec\_eq\_zero\_of\_rayleighOnVec\_zero}, from
\texttt{Matrix.PosSemidef.dotProduct\_mulVec\_zero\_iff}), and
$(\hat b_{u,\sigma})^\dagger \hat b_{u,\sigma} v = 0$ forces
$\hat b_{u,\sigma} v = 0$ (\texttt{conjTranspose\_mul\_self\_mulVec\_eq\_zero}).
This yields
\[
  \hat b_{u,\sigma}|v\rangle = 0\ \ (\forall u,\sigma) \quad (11.3.11),
  \qquad
  \hat n_{x\uparrow}\hat n_{x\downarrow}|v\rangle = 0\ \ (\forall x) \quad (11.3.12),
\]
(\texttt{flatBand\_groundState\_BAnnihilation\_mulVec\_eq\_zero},
\texttt{flatBand\_groundState\_doubleOccupancy\_mulVec\_eq\_zero}). From these
structural constraints Tasaki deduces that a ground state lives in the $\alpha$
band with no double occupancy, hence is spatially symmetric and carries maximal
spin --- the remaining steps of the uniqueness proof.

\medskip
\noindent\textbf{Number conservation and hard-core membership
(\texttt{TasakiFlatBandNumberConservation.lean}).}
Two structural facts organise the uniqueness analysis. First, $\hat H$ conserves
the total particle number, $[\hat H, \hat N] = 0$
(\texttt{flatBandHamiltonian\_commute\_fermionTotalNumber}): each kinetic mode
$\hat b^\dagger_{u,\sigma}\hat b_{u,\sigma}$ is number-conserving (lifting
$[\hat N, \hat c_j] = -\hat c_j$ to $\hat b_{u,\sigma}\hat N - \hat N\hat b_{u,
\sigma} = \hat b_{u,\sigma}$ and dually for $\hat b^\dagger$, so $\hat b^\dagger$
raises and $\hat b$ lowers $\hat N$ by one and the product commutes), and the
interaction $\hat n_{x\uparrow}\hat n_{x\downarrow}$ commutes with $\hat N$ as a
product of number operators. Hence ground states decompose into fixed-$N$ sectors,
and Tasaki works in the physical sector $N_e = |E| = K+1$. Second, the
no-double-occupancy condition (11.3.12) places any ground state in the Hubbard
hard-core subspace (\texttt{flatBand\_groundState\_mem\_hardcoreSubspace}).

\medskip
\noindent\textbf{The two subspaces of the uniqueness argument
(\texttt{TasakiFlatBandSubspaces.lean}).}
Tasaki's uniqueness statement is naturally phrased through two subspaces. The
\emph{$\alpha$-Fock subspace} \texttt{flatBandAlphaFockSubmodule} is the span of all
$\alpha$-band Slater states $\bigl(\prod_q \hat a^\dagger_q\bigr)|\mathrm{vac}
\rangle$ (\texttt{flatBandAlphaSlaterState}); the ferromagnetic state
$|\Phi_{\alpha,\mathrm{all}\uparrow}\rangle$ is the all-up such Slater state and so
lies in it (\texttt{flatBandAlphaAllUpState\_mem\_alphaFockSubmodule}). The
\emph{$\hat b$-kernel subspace} \texttt{flatBandBKernelSubmodule}
$= \bigcap_{u,\sigma}\ker \hat b_{u,\sigma}$ collects the states with no
$\beta$-band occupation; every ground state lies in it by (11.3.11)
(\texttt{flatBand\_groundState\_mem\_BKernelSubmodule}). With these in place, the
uniqueness half reduces to the inclusion
$\texttt{flatBandBKernelSubmodule}\subseteq\texttt{flatBandAlphaFockSubmodule}$ (a
ground state has no $\beta$-occupation, hence is built from the flat band) followed
by the symmetric/maximal-spin classification of the resulting $\alpha$-band states.

\medskip
\noindent\textbf{The easy inclusion $\alpha$-Fock $\subseteq$ $\hat b$-kernel
(\texttt{TasakiFlatBandAlphaFockKernel.lean}).}
The reverse-trivial inclusion
$\texttt{flatBandAlphaFockSubmodule}\subseteq\texttt{flatBandBKernelSubmodule}$
(\texttt{flatBandAlphaFockSubmodule\_le\_BKernelSubmodule}) holds because every
$\alpha$-Slater state is annihilated by every $\hat b_{u,\sigma}$: since
$\hat b_{u,\sigma}$ anticommutes with all $\hat a^\dagger_{p,\tau}$ (eq.~(11.3.7),
all spins) and kills the vacuum, moving it rightward through the ordered product
$\prod_q \hat a^\dagger_q$ (by induction on the product list) lands on
$\hat b_{u,\sigma}|\mathrm{vac}\rangle = 0$
(\texttt{flatBandBAnnihilation\_mulVec\_alphaSlaterState}). The genuinely hard
inclusion is the other one, $\hat b$-kernel $\subseteq$ $\alpha$-Fock, which
requires the Fock-space factorisation of the non-orthogonal $\{\alpha\}\cup\{\beta\}$
basis and is deferred.

\medskip
\noindent\textbf{Theorem 11.11 capstone (\texttt{TasakiFlatBandTheorem11\_11.lean}).}
The full theorem is assembled for the half-filled sector $N_e = K+1$. Define the
ground subspace as the zero-energy, fixed-number subspace
\[
  \texttt{flatBandHalfFilledGroundSubmodule}
    := \ker\hat H \;\sqcap\; \mathrm{eigenspace}(\hat N, K+1),
\]
and the ferromagnetic multiplet subspace as the span of the $K+2 = 2S_{\max}+1$
lowered states. Then
\[
  \texttt{flatBandHalfFilledGroundSubmodule}
    = \texttt{flatBandFerromagneticMultipletSubmodule}
\]
(\texttt{flatBand\_theorem\_11\_11\_groundSubmodule\_eq\_multipletSpan}), and every
ground state carries total spin $S_{\rm tot} = S_{\max} = (K+1)/2$
(\texttt{flatBand\_theorem\_11\_11\_groundState\_maximalSpin}). The $\geq$
(existence) direction is fully proven: each $(\hat S^-_{\rm tot})^k|\Phi_\alpha
\rangle$ is zero-energy ($\hat H \geq 0$) and carries $N_e = K+1$
($\hat N|\Phi_\alpha\rangle = (K+1)|\Phi_\alpha\rangle$, preserved by the
number-conserving tower $[\hat S^-_{\rm tot}, \hat N] = 0$). The $\leq$ (uniqueness)
direction uses a \emph{single documented axiom}
\texttt{flatBand\_zeroEnergy\_halfFilled\_mem\_ferromagneticMultipletSpan}: Tasaki's
Appendix~A classification (no $\beta$-occupation + no double occupancy $\Rightarrow$
spatially symmetric $\Rightarrow$ maximal spin), deferred under the same policy as
Theorem 11.8 / Lemma 11.9, because a faithful proof needs the Fock-space
factorisation of the non-orthogonal $\{\alpha\}\cup\{\beta\}$ basis and the
symmetric-tensor/SU(2) classification. The existence half does \emph{not} depend on
this axiom.

%======================================================================
\subsection{Mielke's flat-band ferromagnetism: the model (Tasaki §11.3.2)}
%======================================================================

\noindent Mielke's flat-band ferromagnetism (Tasaki §11.3.2) is the Hubbard model
on the \emph{line graph} of a base lattice, with uniform hopping. Its Hamiltonian
(eqs.~(11.3.31)/(11.3.6)) is
\[
  \hat H = t\sum_{\{x,y\}\in B}\hat c^\dagger_{x,\sigma}\hat c_{y,\sigma}
    + 2t\sum_x \hat n_x + U\sum_x \hat n_{x\uparrow}\hat n_{x\downarrow},
\]
i.e.\ the uniform-hopping graph Hubbard Hamiltonian plus a $2t\,\hat N$ shift placing
the ground-state energy at $0$. We formalize it on an arbitrary graph $G$ as
\texttt{mielkeHamiltonian} $= \texttt{hubbardHamiltonianOnGraph}\,G\,t\,U +
2t\cdot\hat N$ (\texttt{MielkeHamiltonian.lean}), and record its symmetries:
Hermiticity (\texttt{mielkeHamiltonian\_isHermitian}), particle-number conservation
$[\hat H,\hat N]=0$ (\texttt{mielkeHamiltonian\_commute\_fermionTotalNumber}), and
SU(2) invariance $[\hat S^-_{\rm tot},\hat H]=0$
(\texttt{fermionTotalSpinMinus\_commute\_mielkeHamiltonian}). The line graph is
mathlib's \texttt{SimpleGraph.lineGraph}. The two main results are Theorem 11.12
(a line graph has $D(\tilde\Lambda,B)$ zero-energy single-electron states) and
Theorem 11.13 (Mielke's ferromagnetism: for a biconnected base at half-filling
$N=D$, the ground states are the maximal-spin multiplet). Tasaki defers Theorem 11.12's
proof to §11.3.3 and omits Theorem 11.13's entirely. \textbf{Theorem 11.12 is now
\emph{proved}} (§11.3.3, \texttt{MielkeIncidenceMatrix.lean}, below); only Theorem 11.13
remains a documented axiom (\texttt{MielkeTheorems.lean}), matching the policy for
Theorem 11.8 / Lemma 11.9.
The line graph is mathlib's \texttt{SimpleGraph.lineGraph}, realised on
$\mathrm{Fin}\,(M{+}1)$ via a graph isomorphism
$\texttt{SimpleGraph.Iso}\,G\,(\texttt{lineGraph}\,G_{\rm base})$; the flat-band
count is $D(\tilde\Lambda,B) = |B|-|\tilde\Lambda|+1$ (bipartite base) or
$|B|-|\tilde\Lambda|$ (\texttt{mielkeFlatBandDim}). \texttt{mielke\_theorem\_11\_12}
(now a \emph{proved theorem}, §11.3.3) states that, for a \emph{connected} base lattice
with $|\tilde\Lambda|\le|B|$ (a cycle present), the single-electron operator
\texttt{mielkeSingleElectronOp} ($t\cdot\text{adjacency}+2t\,I$, eq.~(11.3.32)) has
kernel of \texttt{finrank} $D$;
\texttt{mielke\_theorem\_11\_13} states that for a biconnected base at half-filling
$N=D$, the ground subspace of \texttt{mielkeHamiltonian} has \texttt{finrank}
$N{+}1=2S_{\max}{+}1$ and every ground state is an $(\hat S_{\rm tot})^2$
eigenvector at $S_{\max}(S_{\max}{+}1)$, $S_{\max}=N/2$. The Mielke Hamiltonian
model and its symmetries are axiom-free; only these two classification statements
are deferred.

\medskip\noindent\textbf{Towards discharging Theorem 11.12: the incidence-matrix
factorisation (Tasaki §11.3.3, eqs.~(11.3.36)--(11.3.39)).} The proof of the flat
band realises the line-graph single-electron operator as a Gram matrix. Let
$S=\sqrt t\,\cdot$(\,mathlib \texttt{incMatrix}\,) be the \emph{incidence matrix} of
the base lattice $(\tilde\Lambda,B)$, the $\tilde\Lambda\times B$ matrix with
$S_{x,b}=\sqrt t$ when the vertex $x$ is an endpoint of the edge $b$ and $0$
otherwise (\texttt{mielkeIncidence}). Then its Gram matrix $T=S^{\mathsf H}S$, a
$B\times B$ matrix indexed by the edges $=$ the line-graph vertices, satisfies
\[
  (S^{\mathsf H}S)_{b,b'} = t\bigl(2\,\delta_{b,b'} + [\,b\neq b',\ b,b'\text{ share a vertex}\,]\bigr)
  = 2t\,\delta_{b,b'} + t\,(\text{adjacency of the line graph}),
\]
i.e.\ exactly \texttt{mielkeSingleElectronOpOn}$\,(\texttt{lineGraph}\,G)\,t$
(\texttt{mielkeIncidence\_conjTranspose\_mul\_self}, for $t\ge 0$). The diagonal
entry $2t$ is the degree of an edge ($2$ endpoints) and the off-diagonal $t$
records adjacent edges; the proof uses mathlib's
\texttt{sum\_incMatrix\_apply\_of\_mem\_edgeSet} for the diagonal and
\texttt{Sym2.mem\_and\_mem\_iff} (two distinct edges meet in at most one vertex)
for the off-diagonal count. This presents the line-graph operator as a manifestly
positive-semidefinite Gram matrix whose kernel is $\ker S$, the algebraic core for
the rank--nullity and bipartite zero-mode count (eq.~(11.3.41)) that discharge the
Theorem 11.12 axiom in subsequent steps. The single-electron operator is therefore
generalised to an arbitrary finite vertex type (\texttt{mielkeSingleElectronOpOn}),
with the original $\mathrm{Fin}\,(M{+}1)$ form a thin wrapper, since $S^{\mathsf H}S$
naturally lives over $G.\texttt{edgeSet}$ rather than $\mathrm{Fin}\,(M{+}1)$.
The rank step of Lemma 11.14 then follows: mathlib's
\texttt{ker\_mulVecLin\_conjTranspose\_mul\_self} gives $\ker(S^{\mathsf H}S)=\ker S$,
so by rank--nullity for $S$,
\[
  \dim\ker T = |B| - \operatorname{rank} S
\]
(\texttt{mielke\_lineGraph\_ker\_finrank\_eq}, for $t\ge 0$). For the companion
$\tilde T = SS^{\mathsf H}$, rather than build the (mathlib-absent) signless Laplacian we work
directly with $\ker S^{\mathsf H}$: the entry of $S^{\mathsf H}x$ at an edge $b=\{u,v\}$
is $\sqrt t\,(x_u + x_v)$, so for $t>0$ the kernel is
$\{x : x_u + x_v = 0 \text{ for all edges}\}$.  For a \emph{connected} base this kernel is
the one-dimensional span of the alternating $\pm1$ sign vector of a $2$-colouring when the
graph is bipartite, and is trivial otherwise (any nonzero kernel vector would $2$-colour
the graph) --- eq.~(11.3.41):
\[
  \dim\ker(SS^{\mathsf H}) = \dim\ker S^{\mathsf H}
    = \begin{cases} 1 & \text{bipartite},\\ 0 & \text{otherwise}\end{cases}
\]
(\texttt{mielke\_conjTranspose\_ker\_finrank}, for $t>0$ and a connected base).  Combining
with $\operatorname{rank}S = \operatorname{rank}(SS^{\mathsf H}) = |\tilde\Lambda| -
\dim\ker(SS^{\mathsf H})$ gives the flat-band dimension
$\dim\ker T = |B| - (|\tilde\Lambda| - (\text{bipartite}\,?\,1:0)) = D(\tilde\Lambda,B)$
(\texttt{mielke\_lineGraph\_ker\_finrank\_eq\_dim}, the general-base form of Theorem 11.12;
written with the inner subtraction to stay sound under truncated $\mathbb N$ subtraction,
agreeing with Tasaki's $|B|-|\tilde\Lambda|+1$ once $|\tilde\Lambda|\le|B|$).  The
final step (\texttt{mielke\_theorem\_11\_12}) transports this along the line-graph
isomorphism to the $\mathrm{Fin}\,(M{+}1)$ statement: reindexing a matrix by an
equivalence preserves its rank (\texttt{Matrix.rank\_submatrix}) and hence its kernel
dimension, so the flat-band count carries over verbatim, \textbf{discharging the former
\texttt{mielke\_theorem\_11\_12} axiom}.  The hypothesis $|\tilde\Lambda|\le|B|$ is exactly
where Tasaki's $|B|-|\tilde\Lambda|+1$ form coincides with the truncation-safe
$|B|-(|\tilde\Lambda|-\mathrm{bip})$ proved here.

\medskip\noindent\textbf{General theory of flat-band ferromagnetism (Tasaki §11.3.4,
Theorem 11.15, axiomatized).}  Mielke's general theory gives a compact necessary-and-%
sufficient condition for saturated ferromagnetism in any Hubbard model with a flat lowest
band.  Let $T$ be a Hermitian positive-semidefinite hopping matrix on $\Lambda=\{0,\dots,M\}$,
$h_0=\ker T$ the (assumed nonempty) flat band, $D_0=\dim h_0$, and $P_0$ the orthogonal
projection onto $h_0$.  Put $\Lambda_0=\{x\mid (P_0)_{x,x}\ne0\}$ and take the Hubbard model
with $U>0$ at filling $N=D_0$.  \textbf{Theorem 11.15} states that the model is
saturated-ferromagnetic ($N{+}1$-fold degenerate ground states with $S_{\rm tot}=N/2$) if and
only if the $\Lambda_0\times\Lambda_0$ submatrix $((P_0)_{x,y})_{x,y\in\Lambda_0}$ is
\emph{irreducible} (not block-decomposable).  In Lean (\texttt{GeneralFlatBand.lean}),
$P_0$ (\texttt{generalFlatBandProjectionMatrix}) is the matrix, in the standard orthonormal
basis, of \texttt{Submodule.starProjection} onto $\ker(\texttt{Matrix.toEuclideanLin}\,T)$;
irreducibility (\texttt{generalFlatBandProjectionIrreducible}) is \texttt{Matrix.IsIrreducible}
of the real nonnegative support matrix $\mathrm{normSq}((P_0)_{x,y})$ on $\Lambda_0$ (sound
because $P_0$ is Hermitian, so the support quiver is symmetric and strong connectivity equals
Tasaki's irreducibility); the ferromagnetism conclusion
(\texttt{generalFlatBandFerromagnetic}) reuses the Theorem 11.13 ground-subspace form
($\texttt{finrank}=D_0{+}1$, all states $(\hat S_{\rm tot})^2$-eigenvectors at $S_{\max}=D_0/2$).
The result is deep; Tasaki proves it via Lemma 11.16 (a special basis $\{\mu_z\}_{z\in I}$ of
$h_0$, $|I|=D_0$, site-localised so that $\mu_z(z)\ne0$ and $\mu_z(z')=0$ for $z'\in I\setminus\{z\}$)
and Theorem 11.17 (the model is ferromagnetic iff that basis is \emph{connected}, where
$\mu_z\sim\mu_{z'}$ when they share a nonzero component).  All three are recorded as documented
axioms (\texttt{tasaki\_theorem\_11\_15}, \texttt{generalFlatBand\_lemma\_11\_16},
\texttt{generalFlatBand\_theorem\_11\_17}; the basis connectivity is
\texttt{generalFlatBandBasisGraph}/\texttt{generalFlatBandBasisConnected}), matching the policy
for Theorems 11.8 / 11.13; their proofs are deferred to later work.

\medskip\noindent\textbf{Sector minimum energy and the ferromagnetism criterion (Tasaki
§11.4, eq.~(11.4.26)).}  Tasaki's non-singular Hubbard theory (Theorems 11.18--11.20) rests
on the precise definition of \emph{exhibiting ferromagnetism}.  For a many-body Hamiltonian
$\hat H$ the minimum energy in the total-spin-$S$ sector is
\[
  E_{\min}(S) = \min\{\langle\Phi|\hat H|\Phi\rangle : (\hat S_{\rm tot})^2\Phi = S(S{+}1)\Phi,\ \|\Phi\|=1\},
\]
and the model exhibits (saturated) ferromagnetism iff $E_{\min}(S_{\max}) < E_{\min}(S)$ for
every $S\ne S_{\max}$.  In Lean (\texttt{SectorMinEnergy.lean}) sectors are indexed by
$\mathtt{twoS}=2S\in\mathbb N$ (eigenvalue $(\mathtt{twoS}/2)(\mathtt{twoS}/2+1)$, allowing
half-integer spin); \texttt{sectorMinEnergy}~$\hat H~\mathtt{twoS}$ is the infimum of
\texttt{rayleighOnVec}~$\hat H$ over the unit (L²/\texttt{EuclideanSpace}) vectors of the
$(\hat S_{\rm tot})^2$-eigensector, and \texttt{exhibitsFerromagnetism}~$\hat H~\mathtt{twoSmax}$
is $\forall\,\mathtt{twoS}<\mathtt{twoSmax},\ E_{\min}(S_{\max})<E_{\min}(\mathtt{twoS})$.  These
definitions are the foundation of the §11.4 theorems (axiomatized in later work; Issue \#4189).
The non-singular model itself (eq.~(11.4.23)) is \texttt{nonsingularHubbardHamiltonian} $=
\texttt{flatBandHamiltonian} + (\zeta:\mathbb C)\cdot\texttt{hubbardKinetic}$, i.e.\ the flat-band
model $t\sum\hat b^\dagger\hat b + U\sum\hat n_\uparrow\hat n_\downarrow$ perturbed by
$\zeta\sum_{x,y,\sigma} t_{x,y}\,\hat c^\dagger_{x,\sigma}\hat c_{x,\sigma}$ (recovering the flat
band at $\zeta=0$).  Its Hermiticity, particle-number conservation, and SU(2) invariance
$[\hat S^\pm_{\rm tot},\hat H]=0$ are proved (axiom-free) by combining the flat-band and kinetic
symmetry lemmas (the lowering case via the adjoint of the raising case and Hermiticity).

\medskip\noindent\textbf{Local stability (Tasaki §11.4, Theorem 11.18, axiomatized).}  Tasaki's
first non-singular theorem is the local stability of ferromagnetism: there are constants
$\nu_0,\eta_0,\xi_0>0$ (depending only on $d=1$ and the hopping range $R$, uniformly in the
system size) such that, under $0<\nu\le\nu_0$, $|\zeta|\le\nu^3\eta_0$, and
$U\ge\xi_0 t|\zeta|/\nu^2$, the maximal-spin sector lies strictly below the once-flipped sector,
$E_{\min}(S_{\max})<E_{\min}(S_{\max}-1)$.  With $S_{\max}=(K{+}1)/2$ this is
\texttt{sectorMinEnergy}~$\hat H~(K{+}1)<\texttt{sectorMinEnergy}~\hat H~(K{-}1)$.  The
perturbation regularity (eqs.~(11.4.24)--(11.4.25): cyclic translation invariance and range-$R$
summability) is \texttt{IsNonsingularHopping}.  Tasaki's perturbation-theory proof is deferred;
the statement is the documented axiom \texttt{nonsingular\_theorem\_11\_18} (Issue \#4189).

\medskip\noindent\textbf{The §11.4.1 example model (eq.~(11.4.1)).}  Section 11.4.1 (Wannier
state perturbation theory) is heuristic and contains no numbered theorem; its one
formalisable object is the illustrative example --- the one-dimensional flat-band model with
the internal on-site potential shifted by $\gamma$.  Since the only $\gamma$-dependence is the
internal on-site term $(1{+}\gamma)t$ vs $t$, this is the non-singular model
(\texttt{nonsingularHubbardHamiltonian}) whose perturbation is $\gamma t$ times the
internal-site number operator: \texttt{wannierExampleModel}~$K\,\nu\,t\,\gamma\,U =
\texttt{nonsingularHubbardHamiltonian}~K\,\nu\,t\,(\gamma t)\,(\texttt{wannierExamplePerturbation}~K)\,U$.
It is Hermitian, and \texttt{wannierExampleModel\_gamma\_zero} proves it reduces to
\texttt{flatBandHamiltonian} at $\gamma=0$ (eq.~(11.4.1) reducing to (11.3.22)), explicitly
covering §11.4.1 in book order.

\medskip\noindent\textbf{Lattice translation (towards Theorem 11.19).}  Theorem 11.19 is
stated via the fermionic translation operator $\hat\tau_z$ (eq.~(11.4.30)).  Its underlying
combinatorial datum is built in \texttt{TranslationOperator.lean}: the spinful mode space
$\mathrm{Fin}(2(2K{+}1){+}2)=\mathrm{Fin}((2K{+}2)\cdot 2)$ factors as (physical site)
$\times$ (spin) (mode $2p{+}\sigma\leftrightarrow(p,\sigma)$, \texttt{modeSiteSpinEquiv}), and
\texttt{siteTranslationPerm}~$K\,z$ is the \texttt{Equiv.Perm} cyclically shifting the
physical site by $2z$ (one cell $=$ two physical sites), spin fixed.  This is step 1 toward
$\hat\tau_z$.  Step 2 (\texttt{FermionicTranslation.lean}) is the operator itself: the
Jordan--Wigner sign \texttt{translationJwSign}~$\pi\,\sigma = (-1)^{\#\{(i,j):i<j,\ \sigma_i=\sigma_j=1,\ \pi(j)<\pi(i)\}}$
(the parity of occupied inversions, $\pm1$), and \texttt{translationOperator}~$K\,z$, the
signed permutation $\hat\tau_z|\sigma\rangle = \varepsilon(\pi,\sigma)\,|\sigma\circ\pi^{-1}\rangle$
(\texttt{translationOperator\_mulVec\_basisVec} is exactly eq.~(11.4.30) on basis states;
\texttt{translationOperator\_zero} gives $\hat\tau_0=1$).  Finally
(\texttt{SpinWaveExcitation.lean}, completing §11.4.2): \texttt{momentumPhase}~$K\,k =
e^{-2\pi i k/(K+1)}$, the spin-wave energy \texttt{spinWaveEnergy}~$K\,\hat H\,k = E_{\rm SW}(k)$
is the infimum of \texttt{rayleighOnVec}~$\hat H$ over unit states in the joint eigenspace
$\hat S^z_{\rm tot}=(K{-}1)/2\;(=S_{\max}{-}1)$ and $\hat\tau=\texttt{momentumPhase}$, and
\textbf{Theorem 11.19} (\texttt{nonsingular\_theorem\_11\_19}, documented axiom) gives constants
$\nu_1,\eta_1,\xi_1,\xi_2,a_1,\dots,b_3>0$ such that, under (11.4.31)/(11.4.32), the dispersion
$E_{\rm SW}(k)-E_{\min}(S_{\max})$ is sandwiched between $F_2$ and $F_1$ times
$2\nu^4 U(1-\cos k)$ (eqs.~(11.4.33)--(11.4.35)).

\medskip\noindent\textbf{Ferromagnetism in the non-singular model (Tasaki §11.4.3,
Theorem 11.20, axiomatized).}  The culmination of §11.4 is a Hubbard model with a
\emph{nearly-flat} (non-flat) lowest band that nonetheless exhibits saturated ferromagnetism
--- the first such result with no singularity.  The model (eq.~(11.4.38)) is
\texttt{tasakiNonsingularHamiltonian}~$K\,\nu\,t\,s\,U =
\texttt{flatBandHamiltonian}~K\,\nu\,t\,U - s\cdot\sum_{p,\sigma}\hat a^\dagger_{p,\sigma}\hat a_{p,\sigma}$,
which reduces to the flat-band model at $s=0$ (\texttt{tasakiNonsingularHamiltonian\_zero\_s}).
\textbf{Theorem 11.20} (\texttt{tasaki\_theorem\_11\_20}, documented axiom, $d=1$): for every
$\nu>0$ there are thresholds, uniform in the system size, such that once $t/s$ and $U/s$ are
large enough the model satisfies \texttt{exhibitsFerromagnetism} (saturated ferromagnetism,
$S_{\max}=N/2$).  Tasaki's proof builds frustration-free local Hamiltonians $\hat h_p$ on
$4d+1$ sites (Lemmas 11.21--11.22) and reduces to Theorem 11.11; it is deferred.  The
proof-internal numbered lemmas are themselves recorded (\texttt{NonsingularLocalHamiltonian.lean}):
\texttt{nonsingularLocalHamiltonian}~$K\,\nu\,s\,t\,U\,\mathtt{lam}\,\kappa\,i$ is the local
Hamiltonian $\hat h_p$ (eq.~(11.4.48), $d=1$, on $p=i$ and its neighbours), and
\textbf{Lemma 11.21} (\texttt{nonsingular\_lemma\_11\_21}: $\hat h_p\ge0\ \forall p \Rightarrow
\texttt{exhibitsFerromagnetism}$, via Theorem 11.11), \textbf{Lemma 11.22}
(\texttt{nonsingular\_lemma\_11\_22}: parameter conditions $\Rightarrow \hat h_p\ge0$), and
\textbf{Lemma 11.23} (\texttt{nonsingular\_lemma\_11\_23}: the $t,U\uparrow\infty$ limit
positivity below $S_{\max}$) are documented axioms discharging Theorem 11.20's proof.

\medskip\noindent\textbf{Per-site fermionic spin operators (towards §11.5).}  The
ferromagnetic t-J model (§11.5.2, eq.~(11.5.4)) couples spins by a Heisenberg term
$\sum_{\langle x,y\rangle}(\hat n_x\hat n_y/4 - \hat{\mathbf S}_x\cdot\hat{\mathbf S}_y)$, which
needs the per-site spin operators of the Jordan--Wigner fermions
(\texttt{FermionSiteSpin.lean}): \texttt{fermionSiteSpinPlus/Minus/Z}~$N\,i$
($\hat S^+_x=\hat c^\dagger_{x\uparrow}\hat c_{x\downarrow}$, $\hat S^-_x=\hat c^\dagger_{x\downarrow}\hat c_{x\uparrow}$,
$\hat S^z_x=(\hat n_{x\uparrow}-\hat n_{x\downarrow})/2$) and the two-site dot
\texttt{fermionSpinDot}~$N\,i\,j = \tfrac12(\hat S^+_x\hat S^-_y+\hat S^-_x\hat S^+_y)+\hat S^z_x\hat S^z_y$.
They sum to the total spin operators (\texttt{fermionTotalSpinPlus\_eq\_sum\_siteSpinPlus}, …).

\medskip\noindent\textbf{The ferromagnetic t-J model and Proposition 11.24 (Tasaki §11.5.2,
axiomatized).}  On the no-double-occupancy (``hard-core'') subspace $H_N^{\rm hc}$ with
projection $\hat P_{\rm hc}$, the \emph{ferromagnetic t-J model} (eq.~(11.5.4)) is
\[
  \hat H_{tJ} \;=\; -\tau\,\hat P_{\rm hc}\!\!\sum_{\langle x,y\rangle,\sigma}\!\!
    \hat c^\dagger_{x\sigma}\hat c_{y\sigma}\,\hat P_{\rm hc}
  \;+\; J\!\!\sum_{\langle x,y\rangle}\!\Big(\tfrac{\hat n_x\hat n_y}{4}
    - \hat{\mathbf S}_x\cdot\hat{\mathbf S}_y\Big),
  \qquad \tau,J>0,
\]
the strong-coupling effective theory of the nearly-flat-band Hubbard model with electron
number $N$ strictly below the number of sites $L=|E|$.  We formalise it
(\texttt{TJModel.lean}) as \texttt{tJHamiltonian}~$N\,G\,\tau\,J$: the hopping is
\texttt{hubbardKineticOnGraph}~$N\,G\,1$ sandwiched by \texttt{hubbardHardcoreProjection}~$N$
(making it Hermitian, equal to Tasaki's $-\tau\hat P_{\rm hc}\sum\hat c^\dagger\hat c$ on
$H_N^{\rm hc}$), and the exchange term is $J\sum$ over the ordered adjacent pairs of $G$ of
$(\hat n_x\hat n_y)/4 - \texttt{fermionSpinDot}~N\,x\,y$, with $\hat n_x=\texttt{fermionSiteNumber}~N\,x$.
Since the maximal total spin at fixed electron number $N_e$ (with no double occupancy) is
$S_{\max}=N_e/2$, the statement requires the electron number to be \emph{fixed} and the
hard-core constraint imposed; we phrase it on the generic \emph{ground subspace} at filling
$N_e$ (\texttt{GroundSubspaceAtFilling.lean}): \texttt{fillingHardcoreStates} (unit
$N_e$-electron states in the no-double-occupancy subspace $H_{N_e}^{\rm hc}$),
\texttt{groundEnergyAtFilling} (their minimum energy), and
\texttt{groundSubmoduleAtFilling} ($H$-eigenspace at the ground energy $\sqcap$ the
$N_e$-electron sector $\sqcap$ $H_{N_e}^{\rm hc}$).  \textbf{Proposition
11.24} (\texttt{proposition\_11\_24}, documented axiom): for the one-dimensional periodic chain
($\texttt{SimpleGraph.cycleGraph}~(N+1)$), if the electron number $N_e<L$ is \emph{odd}, the
ground states have total spin $S_{\rm tot}=N_e/2$ and are non-degenerate apart from the trivial
$2S_{\rm tot}+1=N_e+1$-fold $SU(2)$ multiplet --- both halves captured at once by
\texttt{IsMaximalSpinMultipletSubmodule}~$N\,(\texttt{groundSubmoduleAtFilling}\dots)\,N_e$ (the
shared predicate reused from Mielke's Theorem 11.13 and the general flat-band Theorem 11.15).
Tasaki's proof is a Perron--Frobenius /
``spin-charge separation'' argument (Theorem A.18); the oddness of $N_e$ is exactly what makes
the wrap-around $1\leftrightarrow L$ hop sign-free under periodic boundary conditions.  The proof
is deferred (Issue \#4198).

\medskip\noindent\textbf{Proposition 11.24 discharge (Tasaki §11.5.2, Issue \#4230, in
progress).}  The axiom is being replaced by a faithful Lean proof along Tasaki's spin--charge
separation route.  \emph{(i) SU(2) invariance.}  $\hat H_{\rm tJ}$ is Hermitian
(\texttt{tJHamiltonian\_isHermitian}; the exchange term is self-adjoint only \emph{as a sum},
since $(\hat S_x\!\cdot\!\hat S_y)^{\mathsf H}=\hat S_y\!\cdot\!\hat S_x$) and commutes with
$\hat S^z_{\rm tot},\hat S^+_{\rm tot},\hat S^-_{\rm tot}$, so by Theorem A.17 every energy level
has an eigenstate of total $\hat S^z\in\{0,\tfrac12\}$.  \emph{(ii) The $\hat
S^z=\tfrac12$, $\hat N=N_e$ sector basis.}  A site-state $s:\mathrm{Fin}(N+1)\to\mathrm{Fin}\,3$
($0/1/2=\varnothing/\!\uparrow/\!\downarrow$) gives the spinful occupation $\texttt{tJConfigOf}\,s$,
which is hard-core, injective and orthonormal; its electron-number and spin-$z$ eigenvalues are the
site counts ($\hat N|\Phi_s\rangle=(\#\!\uparrow+\#\!\downarrow)|\Phi_s\rangle$, $\hat
S^z|\Phi_s\rangle=\tfrac12(\#\!\uparrow-\#\!\downarrow)|\Phi_s\rangle$), so under the count
hypotheses $\#\!\uparrow+\#\!\downarrow=N_e$ and $\#\!\uparrow=\#\!\downarrow+1$ it lies in the
$N_e$-electron, $\hat S^z=\tfrac12$ sector.  \emph{(iii) Sign-free hopping.}  For an \emph{allowed}
hop --- source $a$ carrying spin $\sigma$ and empty target $b$ --- acting
$\hat c^\dagger_{b\sigma}\hat c_{a\sigma}$ on $|\Phi_s\rangle$ produces, up to a Jordan--Wigner
sign, $|\Phi_{\texttt{tJSiteHop}\,s\,a\,b}\rangle$ (for disallowed source/target occupations the
operator term is simply $0$); the forward parity
$\mathrm{jwSign}(q,c)\,\mathrm{jwSign}(p,\texttt{update}\,c\,q\,0)=(-1)^{\sum_{q<k<p}c_k}$ (the
$2E_q$ from the modes below the source cancels) and its backward analogue collapse to $+1$ on a
nearest-neighbour bond (the single intervening orbital is empty) and to $(-1)^{N_e-1}=+1$ on the
periodic wrap bond $\{0,N\}$ for \emph{odd} $N_e$ (a three-way mode split shows the strictly-between
occupation is $N_e-1$).  Hence for an allowed nearest-neighbour or wrap hop the kinetic matrix
element $\langle\Phi_{s'}|\hat c^\dagger_{b\sigma}\hat c_{a\sigma}|\Phi_s\rangle=[\,s'=
\texttt{tJSiteHop}\,s\,a\,b\,]\in\{0,1\}$.  (Occupation-forbidden terms --- empty source or occupied
target --- already vanish at the operator level, $\hat c_{a\sigma}|\Phi_s\rangle=0$ or the creation
hitting a filled mode; the remaining non-adjacent hops, whose per-term element need not vanish, are
removed only by the coupling coefficient $\texttt{couplingOf}\,G\,1\,i\,j=0$ inside the kinetic
sum $\langle\Phi_{s'}|\hat K|\Phi_s\rangle=\sum_{\sigma,i,j}(\texttt{couplingOf}\,G\,1\,i\,j)\,
\langle\Phi_{s'}|\hat c^\dagger_{i\sigma}\hat c_{j\sigma}|\Phi_s\rangle$.)  \emph{(iv) Sign-free exchange.}  The ladder $\hat S^+_i\hat S^-_j=\hat
c^\dagger_{i\uparrow}\hat c_{i\downarrow}\hat c^\dagger_{j\downarrow}\hat c_{j\uparrow}$ maps an
antiparallel pair to $|\Phi_{\texttt{tJSpinSwap}\,s\,i\,j}\rangle$ with net sign $+1$ (each same-site
creation/annihilation pair contributes a squared string sign $=1$), giving the off-diagonal exchange
entry $-J/2$.  \emph{(v) Diagonal terms.}  $\hat n_x\hat n_y$ and $\hat S^z_x\hat S^z_y$ act as
scalars on $|\Phi_s\rangle$, so they vanish off the diagonal.  \emph{(vi) The effective matrix.}
$M=T^{\mathsf H}\hat H_{\rm tJ}T$ (Hermitian, $M_{s',s}=\langle\Phi_{s'}|\hat H_{\rm
tJ}|\Phi_s\rangle$); the hard-core bra drops both projections
($\hat P_{\rm hc}\hat K\hat P_{\rm hc}|\Phi_s\rangle=\hat P_{\rm hc}(\hat K|\Phi_s\rangle)$,
$(\hat P_{\rm hc}u)(\texttt{tJConfigOf}\,s')=u(\texttt{tJConfigOf}\,s')$), and the kinetic element
expands as $\langle\Phi_{s'}|\hat K|\Phi_s\rangle=\sum_{\sigma,i,j}(\text{adj}\,i\,j)\,
\langle\Phi_{s'}|\hat c^\dagger_{i\sigma}\hat c_{j\sigma}|\Phi_s\rangle$.  The \emph{general}
single-hop element (\texttt{tJ\_hop\_matrixElement\_apply}, arbitrary $i,j,\sigma$ before any
adjacency assumption) reads off the single-mode action at the bra configuration: it is the product
of the two Jordan--Wigner string signs times the indicator $[\texttt{tJConfigOf}\,s'=\text{hopped
config}]$ when the source $(j,\sigma)$ is occupied and the target $(i,\sigma)$ empty, and vanishes
otherwise (\texttt{tJ\_hop\_matrixElement\_eq\_zero\_of\_source}/\texttt{\_of\_target}, and
\texttt{\_of\_target\_other} when the target site carries the opposite spin, so the hopped
configuration would be doubly-occupied --- non-hard-core, hence unequal to any hard-core bra) ---
the foundation for the per-term $\langle\hat K\rangle\geq0$ non-negativity, which is \emph{now
proven}: \texttt{tJ\_kinetic\_summand\_zero\_or\_one} shows each summand $\texttt{couplingOf}\,i\,j
\cdot\langle\Phi_{s'}|\hat c^\dagger_{i\sigma}\hat c_{j\sigma}|\Phi_s\rangle$ is $0$ or $1$
(non-adjacent pairs killed by the coupling; an adjacent pair dispatched, via the four oriented
$\texttt{Fin}$-value cases \texttt{cycleGraph\_adj\_val\_cases}, to the rightward/leftward
nearest-neighbour and wrap hop matrix elements, or to the source/target vanishing), whence
\texttt{tJKinetic\_matrixElement\_nonneg} gives $0\leq\langle\Phi_{s'}|\hat K|\Phi_s\rangle$ (real
part non-negative, imaginary part zero).  The \emph{exchange} ladder is likewise $\{0,1\}$-valued
(\texttt{TJExchangeNonneg.lean}): $\hat S^-_j=\hat c^\dagger_{j\downarrow}\hat c_{j\uparrow}$
annihilates $|\Phi_s\rangle$ unless $s\,j=\uparrow$
(\texttt{fermionSpinFlip\_matrixElement\_eq\_zero\_of\_source}), and with $s\,j=\uparrow$ but
$s\,i\neq\downarrow$ ($i\neq j$) the subsequent $\hat S^+_i$ raising hits an empty down-orbital
(\texttt{\_eq\_zero\_of\_target}); hence \texttt{tJ\_exchange\_summand\_zero\_or\_one}: each summand
$\texttt{couplingOf}\,i\,j\cdot\langle\Phi_{s'}|\hat S^+_i\hat S^-_j|\Phi_s\rangle$ is $0$ or $1$
(non-adjacent pairs killed by the coupling; an adjacent antiparallel pair gives the sign-free
indicator $[\,s'=\texttt{tJSpinSwap}\,s\,i\,j\,]$), so the exchange off-diagonal entry
$-(J/2)\cdot(\geq0)\leq0$.  \emph{(vii) The full off-diagonal $M_{s',s}\leq0$ is now assembled}
(\texttt{TJExchangeBondSum.lean}).  Different-site ladders commute,
$\hat S^-_x\hat S^+_y=\hat S^+_y\hat S^-_x$ for $x\neq y$
(\texttt{fermionSiteSpinMinus\_mul\_Plus\_comm}; the two even bilinears act on disjoint mode sets,
so four cross-site anticommutations combine to $(-1)^4=+1$, through the generic ring helpers
\texttt{anticomm\_commute\_mul} / \texttt{anticomm\_commute\_mul\_mul} of
\texttt{Math/AnticommuteCommute.lean}: an element anticommuting with both factors of a product
commutes with the even product, and two even products of pairwise-anticommuting factors commute),
whence the reversed ladder is also
$\{0,1\}$-valued (\texttt{tJ\_exchange\_swap\_summand\_zero\_or\_one}, reusing the $(y,x)$ summand
and the symmetry of \texttt{couplingOf}).  Through a matrix-element linearity helper
$\texttt{tJME}(op)=\langle\Phi_{s'}|op|\Phi_s\rangle$ (additive, scalar-homogeneous, summable), the
interaction element expands (\texttt{tJInteraction\_matrixElement\_eq}) to
$-\tfrac12\sum_{x,y}(\texttt{couplingOf}\,x\,y\cdot\langle\hat S^+_x\hat S^-_y\rangle
+\texttt{couplingOf}\,x\,y\cdot\langle\hat S^-_x\hat S^+_y\rangle)$ once the density and $\hat S^3$
products are dropped (PR-B6) and $\hat S_x\!\cdot\!\hat S_y=\tfrac12(\hat S^+_x\hat S^-_y+\hat
S^-_x\hat S^+_y)+\hat S^3_x\hat S^3_y$ is split; each weighted summand lies in $\{0,1\}$, so the
bond-sum is a non-negative real and \texttt{tJInteraction\_matrixElement\_nonpos} gives the
exchange entry $-(J/2)\cdot(\geq0)\leq0$ (real part $\leq0$, imaginary part $0$).  Combined with the
kinetic $-\tau\cdot(\geq0)$, \texttt{tJEffMatrix\_offdiag\_nonpos} concludes: for $\tau,J\geq0$, odd
$N_e$, and $s'\neq s$, the entry $M_{s',s}=\langle\Phi_{s'}|\hat H_{\rm tJ}|\Phi_s\rangle$ is a
non-positive real.  Every off-diagonal
contribution is thus $-\tau$ (hopping), $-J/2$ (exchange), or $0$ (diagonal) --- all
\emph{non-positive}, the hypothesis of the Perron--Frobenius Theorem A.18.  \emph{(viii) The
real-symmetric sector matrix} (\texttt{TJSectorMatrix.lean}).  The Perron--Frobenius step acts on the
fixed $\hat N=N_e$, $\hat S^3=\tfrac12$ sector $\texttt{TJSpinHalfFillingSector}\,N\,N_e$ (site-states
$s$ with $\#\!\uparrow=\#\!\downarrow+1$ and $\#\!\uparrow+\#\!\downarrow=N_e$), where the wrap-hop is
sign-free and the entries are real.  The real form
$\texttt{tJEffReMatrix}\,s'\,s=(\texttt{tJEffMatrix}\,s'\,s).\mathrm{re}$ restricts to that sector by
$\texttt{submatrix}\,\texttt{Subtype.val}$ as $\texttt{tJEffReMatrixOnSector}$, and
$\texttt{tJEffReMatrixOnSector\_isSymm}$ proves symmetry directly from the Hermiticity of
$\texttt{tJEffMatrix}$: $M_{q,p}=\overline{M_{p,q}}$, so their real parts agree (no global realness of
the full matrix is needed --- it is only sector-real, since the wrap sign $(-1)^{N_e-1}$ depends on
the electron parity).  This is the $\texttt{Matrix}\ \cdots\ \mathbb R$ that feeds Theorem A.18.  \emph{(ix) The
connectivity step relation} (\texttt{TJStepRelation.lean}).  The elementary move
$\texttt{TJStep}\,s\,s'$ is one off-diagonal transition of the t-J Hamiltonian: a nearest-neighbour
or wrap hop of an electron into an adjacent empty site ($s'=\texttt{tJSiteHop}\,s\,a\,b$), or an
adjacent antiparallel spin exchange ($s'=\texttt{tJSpinSwap}\,s\,i\,j$).  Both are precomposition
with a transposition ($\texttt{tJSpinSwap}\,s\,i\,j=s\circ\texttt{Equiv.swap}\,i\,j$), hence site
bijections; so \texttt{tJSpinSwap\_count} / \texttt{tJSiteHop\_count} give
$\texttt{tJStep\_preserves\_sector}$: a step preserves $\#\!\uparrow$ and $\#\!\downarrow$, so the
reachability closure of \texttt{TJStep} stays inside one $\hat N=N_e$, $\hat S^3=\tfrac12$ sector.
This is the step relation whose positive matrix entries will feed
\texttt{matrix\_pow\_succ\_pos\_of\_path} and \texttt{Matrix.isIrreducible\_iff\_exists\_pow\_pos}.
\emph{(x) Per-step matrix-entry positivity} (\texttt{TJStepMatrixEntry.lean}).  For a hop $a\to b$,
the moved-electron kinetic summand $\texttt{couplingOf}\,b\,a\cdot\langle\Phi_{s'}|\hat
c^\dagger_{b\sigma}\hat c_{a\sigma}|\Phi_s\rangle=1$ (\texttt{tJ\_cycle\_hop\_kinetic\_summand\_eq\_one},
dispatching the forward / backward / wrap directional matrix elements via
\texttt{cycleGraph\_adj\_val\_cases}); since every kinetic summand is $\{0,1\}$-valued,
\texttt{Finset.single\_le\_sum} gives $\langle\Phi_{s'}|\hat K|\Phi_s\rangle.\mathrm{re}\geq1$
(\texttt{tJKinetic\_matrixElement\_re\_ge\_one\_of\_hop}).  For an exchange, the exchanged-pair
ladder summand equals $1$, so the bond-sum has real part $\geq1$ and the interaction entry satisfies
$\mathrm{re}\leq-J/2$.  Combining (\texttt{tJEffMatrix\_re\_neg\_of\_step}), for $\tau,J>0$ every
$\texttt{TJStep}\,s\,s'$ gives $M_{s',s}.\mathrm{re}<0$ (a hop $\leq-\tau$, an exchange $\leq-J/2$),
i.e. $B_{s',s}=-M_{s',s}>0$.  \emph{(xi) The adjacent-swap bridge} (\texttt{TJAdjacentSwap.lean}).
An \texttt{AdjacentSwapStep} exchanges the values at an adjacent pair of distinct sites
($s'=s\circ\texttt{Equiv.swap}\,a\,b$); it is always a \texttt{TJStep}
(\texttt{adjacentSwapStep\_to\_TJStep}: a hop if one site is empty, an exchange if the two carry
opposite spins), so \texttt{AdjacentSwapReachable} lifts to \texttt{TJStep}-reachability through
\texttt{ReflTransGen.mono} (\texttt{tjReachable\_of\_adjacentSwapReachable}).  The sector connectivity
thus reduces to the model-agnostic combinatorial fact that two $\texttt{Fin}\,3$-configurations with
equal value-counts are adjacent-swap reachable.  \emph{(xii) Reachability basics}
(\texttt{TJSwapReachableBasics.lean}): \texttt{adjacentSwapReachable\_swap} (the value-swapped
$s\circ\texttt{Equiv.swap}\,a\,b$ is reachable --- a single step if the values differ, \texttt{refl}
if equal), \texttt{adjacentSwapReachable\_count} (reachability preserves every value-count, via the
bijection \texttt{comp\_swap\_count}, so it stays in-sector), and
\texttt{exists\_needed\_value\_right\_of\_same\_counts} (equal counts $+$ prefix $[0,p)$ agreement
$+$ disagreement at $p$ $\Rightarrow$ the value $s'\,p$ occurs in $s$ strictly right of $p$ --- the
bubble target).
\emph{(xiii) Sector connectivity} (\texttt{TJSwapReachable.lean}): \texttt{bubble\_reachable} moves
the value at $q=p+g$ leftward to $p$ by $g$ nearest-neighbour swaps (prefix-preserving, induction on
$g$), and \texttt{reach\_of\_agree\_aux} is the selection sort (bubble the needed value into the
leftmost disagreement, recurse on the unfixed positions), assembling into
\texttt{adjacentSwapReachable\_of\_same\_counts}: any two $\texttt{Fin}\,3$-configurations with equal
value-counts are adjacent-swap reachable, hence (via the bridge) \texttt{TJStep}-reachable within a
sector.  \emph{(xiv) The shifted PF matrix} (\texttt{TJSectorShifted.lean}):
$\texttt{tJSectorShifted}=c\cdot 1-M$ with \texttt{tJSectorShifted\_isSymm}, \texttt{\_nonneg}
(entrywise $\geq 0$ for $c$ above the diagonal: diagonal $c-M_{qq}$, off-diagonal $-M_{q,p}$ from
\texttt{tJEffMatrix\_offdiag\_nonpos}), \texttt{tJStep\_ne} (a step changes the configuration), and
\texttt{tJSectorShifted\_pos\_of\_step} (a \texttt{TJStep} between distinct sector states gives a
strictly positive off-diagonal $B$-entry, from \texttt{tJEffMatrix\_re\_neg\_of\_step} $+$ symmetry).
\emph{(xv) Irreducibility} (\texttt{TJSectorIrreducible.lean}): the value-counts partition the sites
(\texttt{tJ\_value\_count\_total}, $\#\varnothing+\#\!\uparrow+\#\!\downarrow=N+1$) so distinct
sector states share all counts (\texttt{tJSector\_same\_counts}); a \texttt{TJStep}-reachable chain
lifts to a positive matrix-power entry of $B$ (\texttt{tJSectorShifted\_pow\_pos\_of\_reachable},
\texttt{matrix\_pow\_succ\_pos\_of\_pow\_pos\_step} along the chain, membership carried by
\texttt{tJStep\_preserves\_sector}); hence \texttt{tJSectorShifted\_isIrreducible} via
\texttt{isIrreducible\_iff\_exists\_pow\_pos} (positive diagonal $+$ connectivity from
\texttt{adjacentSwapReachable\_of\_same\_counts}).
\emph{(xvi) Perron--Frobenius ground state} (\texttt{TJSectorGroundState.lean}): a
\texttt{Nonempty}-witness (the $\hat S^3=\tfrac12$ filling with the $\uparrow$'s left of the
$\downarrow$'s; counts via \texttt{card\_filter\_val\_lt} and the total) supplies A.18's
\texttt{[Nonempty]}, and \texttt{tJEffReMatrixOnSector\_perronFrobenius} applies Theorem A.18 to
$M=\texttt{tJEffReMatrixOnSector}$: $M$ has a strictly positive eigenvector at its lowest eigenvalue
$\mu$ (the sector ground energy), with $\mathrm{finrank}\leq1$ (non-degeneracy) --- the unique
positive Perron--Frobenius ground state of the spin-charge-separated sector.
\emph{(xvii) Sector$\leftrightarrow$full bridge} (\texttt{TJExpansion.lean}):
$\texttt{tJExpansion}\,v=\sum_s v_s\,|\Phi_s\rangle$, the coefficient functional
$\texttt{tJExpansionCoeff}$, and the left inverse \texttt{tJExpansionCoeff\_tJExpansion} (the
orthonormal sector basis makes the expansion injective), mirroring the Nagaoka
\texttt{tasakiCoeff}/\texttt{tasakiCoeff\_expansion} bridge.
\emph{(xviii) Sector-basis completeness} (\texttt{TJCompleteness.lean}): \texttt{tJ\_completeness} (a
vector supported on the sector configurations equals its own sector expansion
$\texttt{tJExpansion}\,(\texttt{tJExpansionCoeff}\,v)$, by direct orthogonality of the sign-free
$\texttt{basisVec}(\texttt{tJConfigOf}\,s)$), with the inverse \texttt{tJSiteStateOf} and
\texttt{tJConfigOf\_tJSiteStateOf\_of\_hardcore} ($\texttt{tJConfigOf}\circ\texttt{tJSiteStateOf}=
\mathrm{id}$ on hard-core configurations).
\emph{(xix) Charge conservation} (\texttt{TJNumberCommute.lean}): $[\hat H_{\rm tJ},\hat N]=0$
(\texttt{fermionTotalNumber\_commute\_tJHamiltonian}) --- the kinetic sandwich, $\hat n_x\hat n_y$,
and $\hat S_x\!\cdot\!\hat S_y$ are all number-conserving (the $\hat S^+\hat S^-$ ladders via
\texttt{fermionTotalNumber\_commute\_hopping}), completing the conservation laws (the SU(2)
$[\hat H_{\rm tJ},\hat S^\alpha]=0$ were already established).

\emph{(xx) Hard-core preservation} (\texttt{TJHardcorePreserve.lean}): $\hat H_{\rm tJ}$ maps the
no-double-occupancy subspace into itself (\texttt{tJHamiltonian\_mulVec\_mem\_hardcore}), via
$[\hat H_{\rm tJ},\hat P_{\rm hc}]=0$ (\texttt{tJHamiltonian\_commute\_hubbardHardcoreProjection})
and $\hat P_{\rm hc}(\hat H_{\rm tJ}v)=\hat H_{\rm tJ}(\hat P_{\rm hc}v)=\hat H_{\rm tJ}v$.  Each
piece commutes with $\hat P_{\rm hc}=\prod_i(1-\hat n_{i\uparrow}\hat n_{i\downarrow})$: the kinetic
sandwich $\hat P_{\rm hc}\hat K\hat P_{\rm hc}$ by idempotency $\hat P_{\rm hc}^2=\hat P_{\rm hc}$
(\texttt{tJKinetic\_commute\_hubbardHardcoreProjection}), while the per-site
$\hat S^\pm_x,\hat S^3_x,\hat n_x$ commute with every same-site double occupancy
$\hat n_{i\uparrow}\hat n_{i\downarrow}$ (products of number operators, so simultaneously diagonal;
the ladders by the Pauli-exclusion interaction-term commutator) --- hence with every hard-core
factor and with their non-commutative product $\hat P_{\rm hc}$ --- so the density product
$\hat n_x\hat n_y$ and the spin dot $\hat S_x\!\cdot\!\hat S_y$
(\texttt{fermionSpinDot\_commute\_hubbardHardcoreProjection}) commute with $\hat P_{\rm hc}$.  This is
the operator-side input to the sector-matrix lift, mirroring
\texttt{fermionTotalSpinMinus\_mulVec\_mem\_hardcore}.

\emph{(xxi) Operator lift on a sector basis state} (\texttt{TJOperatorLift.lean}):
$\hat H_{\rm tJ}|\Phi_s\rangle=\sum_{s'}\langle\Phi_{s'}|\hat H_{\rm tJ}|\Phi_s\rangle\,|\Phi_{s'}\rangle$
(\texttt{tJHamiltonian\_mulVec\_tJConfigOf}) reassembles its sector-matrix column.  Since
$\hat H_{\rm tJ}|\Phi_s\rangle$ stays hard-core and keeps the $\hat N=N_e$ / $\hat S^3=\tfrac12$
eigenvalues (\texttt{tJHamiltonian\_mulVec\_preserves\_number}/\_spinZ, from
$[\hat H_{\rm tJ},\hat N]=[\hat H_{\rm tJ},\hat S^3]=0$), it is supported on the sector
configurations, where the basis is complete.  The key new input is the support restriction
\texttt{tJ\_mulVec\_apply\_eq\_zero\_of\_not\_sector}: a hard-core $\hat N=N_e$, $\hat S^3=\tfrac12$
eigenstate vanishes off the sector configurations, because a hard-core configuration with electron
count $N_e$ and $\hat S^3$ count $\tfrac12$ is exactly a sector basis index (the diagonal actions
\texttt{fermionTotalUpNumber/DownNumber/SpinZ\_mulVec\_apply} + the $\hat S^3$-shell vanishing
\texttt{mulVec\_apply\_eq\_zero\_of\_spinZ\_ne}).  Mirrors the Nagaoka
\texttt{hubbardEffectiveHamiltonian\_mulVec\_tasakiState}.

\emph{(xxii) Eigenvector lift} (\texttt{TJEigenvectorLift.lean}): a real Perron--Frobenius ground
eigenvector $c$ of $\texttt{tJEffReMatrixOnSector}$ at eigenvalue $\mu$ lifts to an
$\hat H_{\rm tJ}$-eigenvector at the same $\mu$,
$\hat H_{\rm tJ}(\texttt{tJExpansion}\,c)=\mu\cdot\texttt{tJExpansion}\,c$
(\texttt{tJHamiltonian\_mulVec\_tJExpansion\_ofReal}).  Two ingredients:
\texttt{tJHamiltonian\_mulVec\_tJExpansion} ($\hat H_{\rm tJ}$ acts on a sector expansion as the
effective matrix $M=\texttt{tJEffMatrix}$, from the raw-sum per-basis-state lift
\texttt{tJHamiltonian\_mulVec\_tJConfigOf'} and the identification
$\langle\Phi_{s'}|\hat H_{\rm tJ}|\Phi_s\rangle=M_{s',s}$), and the sector realness
\texttt{tJEffMatrix\_sector\_im\_zero} (off-diagonals by \texttt{tJEffMatrix\_offdiag\_nonpos}, the
diagonal by Hermiticity) which upgrades the real eigen-equation to the complex one.  Mirrors the
Nagaoka \texttt{hubbardEffectiveHamiltonian\_mulVec\_tasakiExpansion}.

\emph{(xxiii) Variational bound on the ground energy} (\texttt{TJGroundEnergy.lean}, E2 ``$\le$''
direction): a nonzero real sector eigenvector $c$ of $\texttt{tJEffReMatrixOnSector}$ at eigenvalue
$\mu$ gives $\texttt{groundEnergyAtFilling}\,\hat H_{\rm tJ}\,N_e \le \mu$
(\texttt{tJHamiltonian\_groundEnergyAtFilling\_le\_of\_sectorEigen}).  The lifted state
$\texttt{tJExpansion}\,c$ is admissible (a hard-core $\hat N=N_e$ eigenvector,
\texttt{tJExpansion\_mem\_hardcore} / \texttt{fermionTotalNumber\_mulVec\_tJExpansion} /
\texttt{tJExpansion\_ne\_zero\_of\_ne\_zero}) with energy $\mu$ (the (xxii) lift), so the
variational infimum is bounded above by $\mu$.  Built on a reusable generic lemma
\texttt{groundEnergyAtFilling\_le\_of\_eigenvector} (any nonzero $\hat N=N_e$ hard-core eigenvector
at real $\mu$ bounds the ground energy, via \texttt{ciInf\_le\_of\_le} and the \texttt{BddBelow}
brick \texttt{groundEnergyAtFilling\_bddBelow}).  Reverse direction ($\mu \le
\texttt{groundEnergyAtFilling}$, hence equality) by A.17 + odd $N_e$ next.

\emph{(xxiv) Reverse-bound crux} (\texttt{TJGroundEnergyReverse.lean}, E2 ``$\ge$''): a nonzero
$\hat S^3=\tfrac12$, $\hat N=N_e$, hard-core ($W$) eigenvector of $\hat H_{\rm tJ}$ at a real
eigenvalue $E$ yields a nonzero \emph{real} eigenvector of $\texttt{tJEffReMatrixOnSector}$ at $E$
(\texttt{tJ\_spinHalf\_W\_eigenvector\_to\_sector}).  It is sector-supported (support restriction),
hence equals its sector expansion; the operator action then becomes the complex sector matrix acting
on the coefficient vector, which is real on the sector (\texttt{tJEffMatrix\_sector\_im\_zero}), and
a complex eigenvector of a real-matrix cast at a real eigenvalue has a nonzero real or imaginary part
(\texttt{matrix\_eigenvec\_re/im\_of\_complex}).  Corollary
\texttt{tJ\_sectorMin\_le\_of\_spinHalf\_W\_eigenvalue}: under the Perron--Frobenius minimality of
$\mu$, every such $E$ satisfies $\mu \le E$.  Remaining for E2-$\ge$: a $W$-restricted A.17 placing
the global-minimum $W$-eigenvector into the $\hat S^3=\tfrac12$ block (A.17 + odd $N_e$), completing
$\texttt{groundEnergyAtFilling}=\mu$.

\emph{(xxv) Full $\hat N=N_e$ filling basis} (\texttt{TJFillingBasis.lean}, the $W$-compression
foundation): the all-$\hat S^3$ analog of the sector basis, indexing $W=(\hat N=N_e)\cap H^{\rm hc}$
by $\texttt{TJFillingSector}\,N\,N_e=\{s : \mathrm{Fin}(N{+}1)\to\mathrm{Fin}\,3 \mid \#{\uparrow}+\#{\downarrow}=N_e\}$.
With $\texttt{tJFillingExpansion}$/$\texttt{tJFillingExpansionCoeff}$ and the orthonormal left inverse,
the support restriction $\texttt{tJ\_mulVec\_apply\_eq\_zero\_of\_not\_filling}$ (a hard-core
$\hat N=N_e$ vector vanishes off the filling configurations — no $\hat S^3$ constraint), and
completeness $\texttt{tJ\_filling\_completeness}$.  This is the orthonormal basis of $W$ in which
$\hat H_{\rm tJ}$ compresses to the filling effective matrix $\hat H_W$ (toward
$\texttt{groundEnergyAtFilling}=\texttt{hermitianMinEigenvalue}\,\hat H_W$).

\emph{(xxvi) $W$-projection + compression homomorphism} (\texttt{TJFillingCompress.lean}): the
embedding $T=\texttt{tJFillingEmbedding}$ (columns $|\Phi_s\rangle$ over the filling index), the
submodule $W=\texttt{tJFillingWSubmodule}=(\hat N=N_e)\text{-eigenspace}\sqcap H^{\rm hc}$, and the
preservation predicate $\texttt{PreservesTJFillingW}\,B$ ($B$ maps $W\to W$).  The projection identity
$\texttt{tJFillingProjection\_mulVec\_eq\_of\_mem}$ ($T T^{\mathsf H}$ fixes $W$, via
$\texttt{Matrix.mulVec\_mulVec}$ + $\texttt{tJ\_filling\_completeness}$) yields the compression
$\texttt{tJFillingCompress}\,A = T^{\mathsf H} A T$ and the \textbf{homomorphism}
$\texttt{tJFillingCompress\_mul\_of\_right\_preserves}$: $\texttt{compress}(A)\,\texttt{compress}(B)
=\texttt{compress}(AB)$ when $B$ preserves $W$ (the intermediate $T T^{\mathsf H}$ is invisible).
$\texttt{preservesTJFillingW\_tJHamiltonian}$ records that $\hat H_{\rm tJ}$ preserves $W$ ($[\hat
H_{\rm tJ},\hat N]=0$ + hard-core preservation).  This homomorphism will transfer the $\hat S^\alpha$
su(2)/Hermitian/commute relations to their $W$-compressions, enabling the matrix A.17 at the filling
index.

\emph{(xxvii) A.17 operators preserve $W$} (\texttt{TJFillingSpinCompress.lean}):
$\texttt{preservesTJFillingW\_of\_commute}$ (an operator commuting with $\hat N$ and $\hat P_{\rm hc}$
preserves $W$) and the submodule closure $\texttt{preservesTJFillingW\_smul/\_add/\_sub}$ give
$\texttt{PreservesTJFillingW}$ for $\hat S^{(3)}$ ($\texttt{fermionTotalSpinZ}$), $\hat S^+$,
$\hat S^-$, and the Cartesian $\hat S^{(1)}=\tfrac12(\hat S^++\hat S^-)$,
$\hat S^{(2)}=-\tfrac{i}{2}(\hat S^+-\hat S^-)$ ($\texttt{tJTotalSpinOne}$/$\texttt{Two}$), plus
$\texttt{fermionTotalSpinPlus\_commute\_fermionTotalNumber}$.  Together with
$\texttt{preservesTJFillingW\_tJHamiltonian}$ these feed the compression homomorphism for the
$W$-restricted A.17.

\emph{(xxviii) Compressed $\hat H_W$, $\hat S^{(\alpha)}_W$ satisfy the A.17 hypotheses}
(\texttt{TJFillingCompressSpinAlgebra.lean}): via the compression homomorphism and $\texttt{compress}$
linearity ($\texttt{tJFillingCompress\_smul}$/$\texttt{\_sub}$/$\texttt{\_isHermitian}$), the
$W$-compressions inherit Hermiticity, the su(2) relations
$\texttt{tJFillingCompress\_su2\_12/\_23/\_31}$ ($[\hat S^{(1)}_W,\hat S^{(2)}_W]=i\hat S^3_W$ etc.)
and $\hat H_W$-commutativity $\texttt{tJFillingCompress\_tJHamiltonian\_commute\_one/\_two/\_three}$.
These are exactly the inputs of the matrix Theorem A.17
($\texttt{exists\_joint\_su2\_energy\_eigenstate}$), to be applied at the filling index next.

\emph{(xxix) Filling-embedding isometry and the Rayleigh bridge}
(\texttt{TJFillingRayleighBridge.lean}): the embedding has orthonormal columns,
$T^{\mathsf H} T = 1$ ($\texttt{tJFillingEmbedding\_conjTranspose\_mul\_self}$, from
$\texttt{tJConfigOf\_basisVec\_inner}$), hence is an isometry $\langle Tc,Tc\rangle=\langle c,c\rangle$
($\texttt{tJFillingExpansion\_dotProduct\_self}$).  The \emph{Rayleigh bridge}
$\texttt{rayleighOnVec\_tJFillingCompress}$ then gives
$\texttt{rayleighOnVec}\,\hat H_{\rm tJ}\,(\texttt{tJFillingExpansion}\,c)
=\texttt{rayleighOnVec}\,\hat H_W\,c$ (operator Rayleigh on a lifted state $=$ matrix Rayleigh of
the compression, via the rectangular adjoint identity $\langle c, T^{\mathsf H}y\rangle=\langle
Tc,y\rangle$).  This connects the operator-side $\texttt{groundEnergyAtFilling}$ to the finite
matrix $\hat H_W$, whose minimum eigenvalue is identified with $\mu$ via the $W$-restricted A.17.

\emph{(xxx) Eigenvector lift} (\texttt{TJFillingEigenLift.lean}):
$\texttt{mulVec\_tJFillingExpansion\_of\_compress\_eigen}$ --- if $A$ preserves $W$ and $\Phi$ is an
eigenvector of $\texttt{compress}(A)$ at $E$, then $\texttt{tJFillingExpansion}\,\Phi$ is an
eigenvector of $A$ at $E$ on $W$ (via the projection identity $T T^{\mathsf H}=\mathrm{id}$ on $W$),
together with $\texttt{tJFillingExpansion\_mem\_tJFillingWSubmodule}$.  This turns the filling-index
eigenstate produced by the matrix A.17 into a genuine $\hat H_{\rm tJ}$/$\hat S^3$-eigenstate inside
$W$, the final input to $\texttt{groundEnergyAtFilling}=\mu$.

\emph{(xxxi) $\hat S^3$ diagonal on filling expansions; odd $N_e$ has no $\hat S^3=0$ state}
(\texttt{TJFillingSpinZDiag.lean}): $\hat S^3_{\rm tot}$ scales the coefficient at $s$ by
$\tfrac12(\#{\uparrow}-\#{\downarrow})$ ($\texttt{fermionTotalSpinZ\_mulVec\_tJFillingExpansion}$), and
since for odd $N_e$ every filling state has $\#{\uparrow}\ne\#{\downarrow}$, the only $\hat S^3=0$
filling state is $0$ ($\texttt{tJFillingExpansion\_eq\_zero\_of\_spinZ\_mulVec\_eq\_zero}$).  This
kills the $\hat S^3=0$ branch of the $W$-restricted A.17 for odd $N_e$, forcing the
$\hat S^3=\tfrac12$ sector.

\emph{(xxxii) E2 capstone: $\texttt{groundEnergyAtFilling}=\mu$} (\texttt{TJGroundEnergyGe.lean}):
$\texttt{tJ\_perronFrobeniusMin\_le\_hermitianMinEigenvalue}$ applies the matrix Theorem A.17
($\texttt{ham\_eigenstate\_spin\_zero\_or\_half}$) to the minimum eigenvector of $\hat H_W$ (with the
proven compressed Hermitian/su(2)/commute hypotheses), lifts it to a $W$-eigenstate, and uses odd
$N_e$ to force $\hat S^3=\tfrac12$, giving $\mu\le\texttt{hermitianMinEigenvalue}\,\hat H_W$.  Then
$\texttt{tJ\_groundEnergyAtFilling\_ge\_of\_sectorMin}$ ($\mu\le\texttt{groundEnergyAtFilling}$, via
$\texttt{le\_ciInf}$ + the variational lower bound + the Rayleigh bridge + the isometry), and the
capstone $\texttt{tJHamiltonian\_groundEnergyAtFilling\_eq\_perronFrobeniusMin}$: the d=1
ferromagnetic t-J ground energy at odd filling $N_e$ equals the Perron--Frobenius sector minimum
$\mu$ ($\le$ from the lifted PF eigenvector, $\ge$ from the A.17), packaged with the strictly-positive
PF eigenvector $v$.  \textbf{E2 (ground-energy identification) is complete}, resting only on Theorem
A.17's documented \texttt{exists\_joint\_su2\_energy\_eigenstate}.

\emph{(xxxiii) The $\hat S^3=\tfrac12$ ground block is at most one-dimensional}
(\texttt{TJGroundSpinHalfFinrank.lean}, the SU(2) upper-bound seed):
$\texttt{tJ\_groundSubmodule\_spinHalf\_finrank\_le\_one}$ proves
$\dim_{\mathbb C}\bigl(\texttt{groundSubmoduleAtFilling}\,\hat H_{\rm tJ}\,N_e \sqcap (\hat S^3=\tfrac12)\bigr)\le 1$.
The argument fixes the Perron--Frobenius datum once (so $\texttt{groundEnergyAtFilling}=\mu$ and the
\emph{real} sector eigenspace at $\mu$ has $\dim_{\mathbb R}\le 1$); the real$\to$complex bridge
$\texttt{matrix\_complex\_eigenspace\_finrank\_le\_one\_of\_real}$ promotes this to
$\dim_{\mathbb C}\le 1$ for the complexified sector matrix
$(\texttt{tJEffReMatrixOnSector}).\mathrm{map}\,(\cdot:\mathbb R\to\mathbb C)$ at $\mu$; and the block
embeds into that eigenspace by the injective $\mathbb C$-linear map $\texttt{tJExpansionCoeff}_\ell$
($\Phi\mapsto\texttt{tJExpansionCoeff}\,\Phi$, injective because every block element is
sector-supported by $\texttt{tJ\_completeness}$), so
$\texttt{LinearMap.finrank\_le\_finrank\_of\_injective}$ finishes.  The supporting lemma
$\texttt{tJ\_spinHalf\_W\_complexSectorEigen}$ shows a hard-core $\hat N=N_e$, $\hat S^3=\tfrac12$
eigenvector of $\hat H_{\rm tJ}$ at real $E$ has, as its coefficient vector, a complex eigenvector of
the complexified sector matrix at $E$ (re-deriving the matrix action + sector-realness, injecting via
$\texttt{tJExpansionCoeff\_tJExpansion}$); no new axioms.

\emph{(xxxiv) The single-site spin-raising operator is sign-free on the sector basis}
(\texttt{TJSectorRaise.lean}, E3b PR1, toward lifting the PF vector to a highest weight):
$\texttt{fermionSiteSpinPlus\_mulVec\_tJConfigOf\_of\_down}$ proves $\hat S^+_x|\Phi_s\rangle =
|\Phi_{s,\,x\mapsto\uparrow}\rangle$ (for $s\,x=\downarrow$) with net Jordan--Wigner sign $+1$:
$\hat S^+_x=\hat c^\dagger_{x\uparrow}\hat c_{x\downarrow}$ acts on the \emph{adjacent} orbitals
$(x\uparrow,x\downarrow)=(2x,2x+1)$ (since $\texttt{spinfulIndex}\,N\,x\,\sigma=2x+\sigma$), and the
up-orbital $2x$ is empty before the move, so the two strings collapse
($\texttt{jwSign\_succ\_cancel\_high}$); the companion
$\texttt{fermionSiteSpinPlus\_mulVec\_tJConfigOf\_of\_not\_down}$ gives $\hat S^+_x|\Phi_s\rangle=0$
for $s\,x\neq\downarrow$, and the config identity $\texttt{tJConfigOf\_update\_raise}$ records the
orbital update.  This Marshall sign-freeness makes the iterated $(\hat S^+)^m$ have nonnegative
configuration coefficients, so the strictly-positive Perron--Frobenius coordinates ($v_q>0$) yield a
strictly-positive target coordinate, giving the non-vanishing $(\hat S^+)^{(N_e-1)/2}\Phi_0\neq 0$;
fully axiom-free.

\emph{(xxxv) Eigenvalue tracking along the spin-raising tower} (\texttt{TJRaisingTower.lean}, E3b
PR2): $\texttt{fermionTotalSpinZ\_mulVec\_spinPlusPow}$ ($\hat S^3 (\hat S^+)^k v = (m+k)(\hat
S^+)^k v$ --- each $\hat S^+$ step raises the $\hat S^3$ eigenvalue by one, via $[\hat S^3,\hat
S^+]=\hat S^+$; the raising mirror of the lowering tower),
$\texttt{tJHamiltonian\_mulVec\_spinPlusPow}$ ($\hat H_{\rm tJ}(\hat S^+)^k v=\mu(\hat S^+)^k v$, the
energy preserved since $[\hat H_{\rm tJ},\hat S^+]=0$), and
$\texttt{fermionTotalNumber\_mulVec\_spinPlusPow}$ ($\hat N(\hat S^+)^k v=N_e(\hat S^+)^k v$, the
electron number preserved since $[\hat N,\hat S^+]=0$).  So $(\hat S^+)^k\Phi_0$, when nonzero, is a
fixed-$N_e$, energy-$\mu$, $\hat S^3=\tfrac12+k$ eigenvector --- the building blocks for raising the
strictly-positive Perron--Frobenius ground vector to a nonzero highest weight; fully axiom-free.

\emph{(xxxvi) The spin-raising tower terminates} (\texttt{TJRaisingTermination.lean}, E3b PR3):
$\texttt{fermionTotalDownNumber\_mul\_fermionTotalSpinPlus}$ ($[\hat N_\downarrow,\hat S^+]=-\hat
S^+$, summed from the per-site relation), $\texttt{fermionTotalDownNumber\_mulVec\_spinPlusPow}$
($\hat N_\downarrow(\hat S^+)^k v=(m-k)(\hat S^+)^k v$ --- each raising removes one down-spin), and
$\texttt{spinPlusPow\_succ\_eq\_zero\_of\_downNumber}$: a vector with $\hat N_\downarrow v=m v$ is
annihilated by $m+1$ raisings ($(\hat S^+)^{m+1}v=0$), since the would-be $\hat N_\downarrow$
eigenvalue $-1$ forces $(\#{\downarrow}(w)+1)\,\psi(w)=0$ for every configuration $w$ (via the
diagonal $\texttt{fermionTotalDownNumber\_mulVec\_apply}$, $\#{\downarrow}\ge 0$).  Applied to a
$\hat S^3=\tfrac12$, $\hat N=N_e$ ground state ($\hat N_\downarrow=(N_e-1)/2$), the top
$\Omega=(\hat S^+)^{(N_e-1)/2}\Phi$ satisfies $\hat S^+\Omega=0$ --- a highest weight; fully
axiom-free.

\emph{(xxxvii) The raised sector ground state is a highest weight} (\texttt{TJHighestWeight.lean},
E3b PR4): $\texttt{tJ\_raised\_highestWeight}$ combines the tower tracking and termination --- a
sector ground state $\Phi$ with $\hat S^3\Phi=\tfrac12\Phi$, $\hat N_\downarrow\Phi=m\Phi$, $\hat
H_{\rm tJ}\Phi=\mu\Phi$ produces at the top $\Omega=(\hat S^+)^m\Phi$ a highest-weight ground state:
$\hat S^+\Omega=0$, $\hat S^3\Omega=(m+\tfrac12)\Omega$, $\hat H_{\rm tJ}\Omega=\mu\Omega$.  For the
d=1 ferromagnetic t-J ground state ($m=(N_e-1)/2$) this gives $\hat S^3\Omega=(N_e/2)\Omega$ --- the
maximal-spin highest weight feeding $\texttt{highestWeight\_spinMultiplet\_general}$; the remaining
input is $\Omega\neq 0$ (Marshall positivity).  Fully axiom-free.

\emph{(xxxviii) Sector eigenvalues of the lifted ground vector} (\texttt{TJExpansionSpinEigen.lean},
E3b PR5a): $\texttt{tJSpinHalfFillingSector\_down\_count}$ (every $\hat S^3=\tfrac12$, $\hat N=N_e$
sector state has $\#{\downarrow}=(N_e-1)/2$, by \texttt{omega} from $\#{\uparrow}=\#{\downarrow}+1$
and $\#{\uparrow}+\#{\downarrow}=N_e$), $\texttt{fermionTotalSpinZ\_mulVec\_tJExpansion}$ ($\hat
S^3(\texttt{tJExpansion}\,v)=\tfrac12(\texttt{tJExpansion}\,v)$), and
$\texttt{fermionTotalDownNumber\_mulVec\_tJExpansion}$ ($\hat N_\downarrow(\texttt{tJExpansion}\,v)
=((N_e-1)/2)(\texttt{tJExpansion}\,v)$).  These let the highest-weight extraction apply to the lifted
Perron--Frobenius ground vector $\Phi_0=\texttt{tJExpansion}(\mathbb C\circ v)$ with $m=(N_e-1)/2$;
fully axiom-free.

\emph{(xxxix) The total raising operator on a sector basis state} (\texttt{TJTotalRaiseAction.lean},
E3b PR5b-1): $\texttt{fermionTotalSpinPlus\_mulVec\_tJConfigOf}$ --- $\hat S^+_{\rm tot}|\Phi_s\rangle
= \sum_{x : s\,x=\downarrow} |\Phi_{s,\,x\mapsto\uparrow}\rangle$, the sign-free expansion (every term
has coefficient $+1$, summing the single-site action over sites).  This is the keystone for the
config-nonnegativity argument: the iterated raising preserves nonnegative configuration coefficients,
so $(\hat S^+)^{(N_e-1)/2}\Phi_0\neq 0$ follows from $v>0$.  Fully axiom-free.

\emph{(xl) Coefficient sum of the raising action} (\texttt{TJRaiseCoeffSum.lean}, E3b PR5b-2): the
linear functional $\texttt{coeffSum}\,\psi=\sum_c\psi_c$ with $\texttt{coeffSum\_basisVec}$
($\texttt{coeffSum}\,|c\rangle=1$) and $\texttt{coeffSum\_fermionTotalSpinPlus\_tJConfigOf}$
($\texttt{coeffSum}(\hat S^+_{\rm tot}|\Phi_s\rangle)=\#{\downarrow}(s)$, each raised basis state
contributing $1$).  With the uniform sector down-count $\#{\downarrow}=(N_e-1)/2$, iterating gives
$\texttt{coeffSum}((\hat S^+)^{(N_e-1)/2}\Phi_0)=((N_e-1)/2)!\cdot\sum_q v_q>0$ (from $v>0$), so the
raised vector is nonzero.  Fully axiom-free.

\emph{(xli) Coefficient sum of a filling expansion} (\texttt{TJFillingCoeffSum.lean}, E3b PR5b-3a):
$\texttt{coeffSum\_tJFillingExpansion}$ ($\texttt{coeffSum}(\texttt{tJFillingExpansion}\,v)=\sum_s
v_s$) and $\texttt{coeffSum\_fermionTotalSpinPlus\_tJFillingExpansion}$
($\texttt{coeffSum}(\hat S^+_{\rm tot}(\texttt{tJFillingExpansion}\,v))=\sum_s v_s\,\#{\downarrow}(s)$).
Together they give the recursion $\texttt{coeffSum}(\hat S^+\psi)=d\cdot\texttt{coeffSum}\,\psi$ for a
hard-core $\hat N_\downarrow$-eigenvector $\psi$ at $d$ (both sides equal $\sum_s v_s\#{\downarrow}(s)$),
the multiplicative step for the iterated-raising positivity.  Fully axiom-free.

\emph{(xlii) The coefficient-sum recursion step} (\texttt{TJRaisePositivityStep.lean}, E3b PR5b-3b):
$\texttt{fermionTotalDownNumber\_mulVec\_tJFillingExpansion}$ ($\hat N_\downarrow$ diagonal on a filling
expansion) and $\texttt{coeffSum\_fermionTotalSpinPlus\_eq\_of\_downEigen}$ --- a hard-core $\hat
N=N_e$ vector $\psi$ with $\hat N_\downarrow\psi=d\psi$ satisfies $\texttt{coeffSum}(\hat
S^+_{\rm tot}\psi)=d\cdot\texttt{coeffSum}\,\psi$ (both $\hat S^+\psi$ and $\hat N_\downarrow\psi$ have
coefficient sum $\sum_s\text{coeff}_s\,\#{\downarrow}(s)$ via $\texttt{tJ\_filling\_completeness}$, and
$\texttt{coeffSum}(\hat N_\downarrow\psi)=d\cdot\texttt{coeffSum}\,\psi$ by the eigen-equation).  The
multiplicative step driving $\texttt{coeffSum}((\hat S^+)^k\Phi_0)$ to $((N_e-1)/2)!\cdot\sum_q v_q>0$.
Fully axiom-free.

Remaining: the total-spin
identification ($S_{\mathrm{tot}}=N_e/2$) via the highest weight raised this way, the SU(2)-tower
lower bound (paired with this $\hat S^3=\tfrac12$ block
bound for the upper bound), and the capstone that
replaces the axiom.
The discharge lemmas add no
new t-J axioms; the one inherited dependency is the A.17 sector reduction
\texttt{tJHamiltonian\_eigenstate\_spin\_zero\_or\_half}, which rests on Theorem A.17's single
documented axiom \texttt{exists\_joint\_su2\_energy\_eigenstate}.  Every other layer (the
hop/exchange/diagonal sign computations and the matrix-element machinery) is fully axiom-free.

\medskip\noindent\textbf{The d=1 decorated Hubbard model, Lemma 11.25 and Theorem 11.26
(Tasaki §11.5.3, axiomatized).}  Tasaki obtains a rigorous example of \emph{metallic}
ferromagnetism by realising the t-J model as the strong-coupling limit of a Hubbard model on
the \emph{duplicated-internal-site} decorated lattice $\Lambda=E\cup(I\times\{1,2\})$.  We
formalise it (\texttt{MetallicFerroModel.lean}, $d=1$, $K+1$ cells) by packing the $3(K+1)$
sites into $\mathrm{Fin}(3K+3)$ (\texttt{decExternalSite}, \texttt{decInternalSite1},
\texttt{decInternalSite2}); the localised states \texttt{decAlpha}, \texttt{decBeta1},
\texttt{decBeta2} (eqs.~(11.5.7)--(11.5.9)) give fermion operators $\hat a,\hat b$ via
$\hat C_\sigma(\varphi)=\sum_x\varphi(x)\hat c_{x,\sigma}$, and the Hamiltonian is
$\hat H=\hat H_{\rm hop}+\hat H_{\rm int}$ with
$\hat H_{\rm hop}=t\sum_{(c,\zeta),\sigma}\hat b^\dagger_{(c,\zeta),\sigma}\hat b_{(c,\zeta),\sigma}
-s\sum_{\langle p,q\rangle\in E,\sigma}\hat a^\dagger_{p,\sigma}\hat a_{q,\sigma}$
(eq.~(11.5.14)) and $\hat H_{\rm int}=U\sum_x\hat n_{x\uparrow}\hat n_{x\downarrow}$
(eq.~(11.5.13)); the off-diagonal $\hat a^\dagger_p\hat a_q$ makes the lowest band dispersive,
and the internal duplication forces the standard normalised anticommutation
$\{\hat a_{p,\sigma},\hat a^\dagger_{q,\tau}\}=(1+4d\nu^2)\delta_{pq}\delta_{\sigma\tau}$
(eq.~(11.5.11)).  \textbf{Lemma 11.25} (\texttt{lemma\_11\_25}, documented axiom): the
$t,U\uparrow\infty$ limit of this Hubbard model at filling $N_e$ is equivalent to the
$J\uparrow\infty$ limit of the t-J model on $E=\texttt{cycleGraph}(K+1)$ with
$\tau=(1+4\nu^2)s$; Tasaki identifies both finite-energy subspaces with hard-core electrons on
$E$ carrying the same effective Hamiltonian, so the total-spin structure transfers.  We render
this faithfully: in the limits the Hubbard ground subspace at filling $N_e$ is the maximal-spin
$(N_e+1)$-fold multiplet \emph{iff} the t-J one is.  \textbf{Theorem 11.26}
(\texttt{theorem\_11\_26}, documented axiom): for $d=1$, $N_e\le K+1=|E|$ odd (Tasaki's
$N\le L$), in the limit $t,U\uparrow\infty$ the ground states are
$\texttt{IsMaximalSpinMultipletSubmodule}~(3K+2)\,(\texttt{groundSubmoduleAtFilling}\dots)\,N_e$
--- total spin $S_{\rm tot}=N_e/2$ and $N_e+1$-fold; a metallic ferromagnet when $N_e<K+1$ (the
lowest band is partially filled), reducing to the insulating Heisenberg ferromagnet at
$N_e=K+1$.  Tasaki derives it from Proposition 11.24 and Lemma 11.25 (Perron--Frobenius); the
proof is deferred (Issue \#4198).

\medskip\noindent\textbf{The Tanaka--Tasaki model and Theorem 11.27 (Tasaki §11.5.4,
axiomatized) --- the Chapter 11 capstone.}  The only rigorous example of metallic
ferromagnetism in a Hubbard model in \emph{arbitrary} dimension is the Tanaka--Tasaki model on
the heavily decorated lattice $\Lambda=E\times\{1,2,3\}\cup I\times\{1,2\}$ (externals
triplicated, internals duplicated).  We formalise the $d=1$ realisation
(\texttt{TanakaTasakiModel.lean}, $K+1$ cells) by packing the $5(K+1)$ sites into
$\mathrm{Fin}(5K+5)$ (\texttt{ttExtSite}, \texttt{ttIntSite}).  The special single-particle
states \texttt{ttAlpha}, \texttt{ttBeta}, \texttt{ttDeltaP}, \texttt{ttDeltaI}
(eqs.~(11.5.19)--(11.5.22); $\hat a_p$ carries the normalisation $1/\sqrt{3+4\nu^2}$ ensuring
the standard CAR (11.5.23)) give fermion operators through the shared helper
$\hat C_\sigma(\varphi)=\sum_x\varphi(x)\hat c_{x,\sigma}$ (\texttt{spinfulCreationFromVector},
\texttt{SpinfulVectorOperator.lean}).  The Hamiltonian is $\hat H=\hat H_{\rm hop}+\hat H_{\rm
int}$ with (eq.~(11.5.24))
$\hat H_{\rm hop}=\sum_{\langle p,q\rangle\in E,\sigma}(-s\,\hat a^\dagger_p\hat a_q
-t\,\hat b^\dagger_p\hat b_q)+u_1\sum_{p,\sigma}\hat b^\dagger_p\hat b_p
+u_2(\sum_{p,\sigma}\hat d^\dagger_p\hat d_p+\sum_{(u,\zeta),\sigma}\hat d^\dagger_{(u,\zeta)}\hat
d_{(u,\zeta)})$ and $\hat H_{\rm int}=U\sum_x\hat n_{x\uparrow}\hat n_{x\downarrow}$; the
$\hat d$-bands are pushed to infinity in the $u_2\uparrow\infty$ limit.  We encode that limit
\emph{faithfully} (\emph{not} as finite ``large enough'' thresholds --- Tasaki proves
Theorem 11.27 only in the limit and explicitly warns it is not expected to hold at finite
$u_2,U$, unlike the §11.4.3 insulating model): the $u_2,U=\infty$ finite-energy subspace is
\texttt{ttDKernel} ($\hat d\,\Phi=0$ for every $\hat d$, via Theorem A.12 / Lemma A.11) and the
$u_2,U=\infty$ effective Hamiltonian \texttt{ttEffectiveHamiltonian} is the surviving
$\hat a$/$\hat b$ hopping $\sum_{\langle p,q\rangle,\sigma}(-s\,\hat a^\dagger_p\hat a_q
-t\,\hat b^\dagger_p\hat b_q)+u_1\sum_{p,\sigma}\hat b^\dagger_p\hat b_p$.  \textbf{Theorem
11.27} (\texttt{theorem\_11\_27}, documented axiom): for $u_1>2(|s|+2|t|)$ (Tasaki's
$u_1>2d(|s|+2|t|)$) and $K+1\le N_e\le 2(K+1)$ (Tasaki's $L^d\le N\le 2L^d$), in the limit
every ground state in $\texttt{groundSubmoduleAtFilling}~(\texttt{ttEffectiveHamiltonian}\dots)\,N_e
\sqcap\texttt{ttDKernel}$ has the maximum possible total spin $S_{\rm tot}=N_e/2$ (an
$(\hat S_{\rm tot})^2$ eigenvector at $(N_e/2)(N_e/2+1)$; weaker than the full multiplet
predicate, as Tasaki claims only the spin).  The ground states are metallic when $1<N_e/L<2$.
Tasaki cites the original paper [63] for the technical proof; deferred (Issue \#4198).  This
completes the numbered results of §11.5 and of Chapter 11.

%======================================================================
\subsection{Total-spin Casimir and Definition 11.1 (Tasaki §11.1.1)}
\label{subsec:hubbard-saturated-ferro}
%======================================================================

\noindent\textbf{Definition.}
The total-spin Casimir is
\[
  (\hat{S}^2_{\rm tot})
  := \hat{S}^-_{\rm tot}\hat{S}^+_{\rm tot}
     + \hat{S}^z_{\rm tot}(\hat{S}^z_{\rm tot}+1).
\]
From $[\hat{S}^+,\hat{S}^-]=2\hat{S}_z$ one derives
$\hat{S}^2 = \hat{S}^-\hat{S}^+ + \hat{S}_z(\hat{S}_z+1)$.

\noindent\textbf{Eigenvalue on the all-up state.}
With $S_{\max} = (N+1)/2$:
\begin{align*}
  \hat{S}^z_{\rm tot}\,|\!\uparrow\cdots\uparrow\rangle
    &= \tfrac{N+1}{2}\,|\!\uparrow\cdots\uparrow\rangle,\\
  \hat{S}^+_{\rm tot}\,|\!\uparrow\cdots\uparrow\rangle
    &= 0 \quad\text{(highest-weight state)},\\
  (\hat{S}^2_{\rm tot})\,|\!\uparrow\cdots\uparrow\rangle
    &= S_{\max}(S_{\max}+1)\,|\!\uparrow\cdots\uparrow\rangle.
\end{align*}

\noindent\textbf{Theorem} (\textit{Tasaki §11.1.1, p.~372}).
$[(\hat{S}^2_{\rm tot}),\hat{H}] = 0$
(from SU(2) invariance; requires Hermiticity conditions on $t$, $U$).

\noindent\textbf{Definition 11.1} (Tasaki §11.1.1, p.~372; saturated ferromagnetism).
The Hubbard model exhibits saturated ferromagnetism if there exists a
ground-state energy $E_0$ such that every nonzero eigenvector $v$
with $\hat H v = E_0 v$ satisfies
$(\hat{S}^2_{\rm tot})v = S_{\max}(S_{\max}+1)v$.

\noindent\textbf{Lean names}
(file \texttt{Fermion/JordanWigner/Hubbard/SaturatedFerromagnetism.lean}):
\begin{itemize}
  \item \texttt{fermionTotalSpinSquared N}: the Casimir operator.
  \item \texttt{fermionTotalUpNumber\_mulVec\_allUpState}:
    $N_\uparrow\cdot|\!\uparrow\cdots\uparrow\rangle = (N+1)\,|\!\uparrow\cdots\uparrow\rangle$.
  \item \texttt{fermionTotalDownNumber\_mulVec\_allUpState}:
    $N_\downarrow\cdot|\!\uparrow\cdots\uparrow\rangle = 0$.
  \item \texttt{fermionTotalSpinZ\_mulVec\_allUpState}:
    $\hat{S}^z_{\rm tot}\cdot|\!\uparrow\cdots\uparrow\rangle
      = \frac{N+1}{2}\,|\!\uparrow\cdots\uparrow\rangle$.
  \item \texttt{fermionTotalSpinPlus\_mulVec\_allUpState}:
    $\hat{S}^+_{\rm tot}\cdot|\!\uparrow\cdots\uparrow\rangle = 0$.
  \item \texttt{fermionTotalSpinSquared\_mulVec\_allUpState}:
    $(\hat{S}^2_{\rm tot})\cdot|\!\uparrow\cdots\uparrow\rangle
      = S_{\max}(S_{\max}+1)\,|\!\uparrow\cdots\uparrow\rangle$.
  \item \texttt{fermionTotalSpinSquared\_commute\_hubbardHamiltonian}:
    $[(\hat{S}^2_{\rm tot}),\hat H] = 0$.
  \item \texttt{isSaturatedFerromagnet N t U}: the Lean predicate for
    Definition 11.1.
\end{itemize}

%======================================================================
\subsection{SU(2) algebra: \([\hat S^z,\hat S^-]=-\hat S^-\) and eigenvalue stepping (Tasaki §9.3.3, §11.1.1)}
\label{subsec:hubbard-su2-algebra}
%======================================================================

\noindent\textbf{Site-wise commutator.}
For each site $i$, let $A_i = c^\dagger_{i,\downarrow}c_{i,\uparrow}$.
One computes
\[
  [\hat S^z_{\rm tot},\,A_i]
  = \tfrac{1}{2}[N_\uparrow - N_\downarrow,\,A_i]
  = \tfrac{1}{2}(-A_i - A_i) = -A_i.
\]
Here $[N_\uparrow, A_i] = -A_i$ (since $N_\uparrow$ commutes with
$c^\dagger_{i,\downarrow}$ but not with $c_{i,\uparrow}$) and
$[N_\downarrow, A_i] = A_i$ (since $[N_\downarrow, c^\dagger_{i,\downarrow}] = c^\dagger_{i,\downarrow}$
and $N_\downarrow$ commutes with $c_{i,\uparrow}$).

\noindent\textbf{Global commutator.}
Summing over sites:
\[
  [\hat S^z_{\rm tot},\,\hat S^-_{\rm tot}] = -\hat S^-_{\rm tot}.
\]

\noindent\textbf{Consequences.}
\begin{enumerate}
\item (Eigenvalue preservation.)
  If $\hat H v = E_0 v$ and $[H, \hat S^-] = 0$, then
  $\hat H(\hat S^- v) = E_0\,(\hat S^- v)$.
\item (Eigenvalue step.)
  If $\hat S^z v = m\,v$, then
  $\hat S^z(\hat S^- v) = (m-1)(\hat S^- v)$.
  Proof: $\hat S^z \hat S^- = \hat S^- \hat S^z - \hat S^-$, so
  $\hat S^z(\hat S^- v) = \hat S^-(mv) - \hat S^- v = (m-1)(\hat S^- v)$.
\end{enumerate}

\noindent\textbf{Lean names}
(file \texttt{Fermion/JordanWigner/Hubbard/SaturatedFerromagnetism.lean}):
\begin{itemize}
  \item \texttt{fermionTotalSpinZ\_commutator\_fermionTotalSpinMinus}:
    $[\hat S^z_{\rm tot},\hat S^-_{\rm tot}] = -\hat S^-_{\rm tot}$
    (Tasaki §9.3.3, p.~332).
  \item \texttt{fermionTotalSpinMinus\_mulVec\_preserves\_hamiltonian\_eigenvalue}:
    if $\hat H v = E_0 v$ then $\hat H(\hat S^- v) = E_0(\hat S^- v)$
    (Tasaki §11.1.1, p.~373).
  \item \texttt{fermionTotalSpinZ\_mulVec\_spinMinus\_step}:
    if $\hat S^z v = m\,v$ then $\hat S^z(\hat S^- v) = (m-1)(\hat S^- v)$
    (Tasaki §2.4 eq.~(2.4.9); §11.1.1, p.~373).
\end{itemize}

%======================================================================
\section{Time-reversal map for a single \(S = 1/2\) spin (Tasaki §2.3)}
\label{sec:time-reversal-spin-half}
%======================================================================

\noindent\textbf{Definition} (Tasaki §2.3 eq.~(2.3.6), \(S = 1/2\)).
The time-reversal map on a single spin-\(1/2\) is
\(\hat{\Theta} := \hat{u}_2 \cdot \hat{K}\), where \(\hat{u}_2\) is
the \(\pi\)-rotation about the \(2\)-axis and \(\hat{K}\) is complex
conjugation on coefficients in the standard basis. Explicitly,
\[
  \hat{\Theta}\!\begin{pmatrix} a \\ b \end{pmatrix}
  \;=\; \begin{pmatrix} -b^{*} \\ a^{*} \end{pmatrix},
  \qquad a, b \in \mathbb{C}.
\]
In Lean this is \texttt{timeReversalSpinHalf} as a plain
\((\texttt{Fin}\,2 \to \mathbb{C}) \to (\texttt{Fin}\,2 \to \mathbb{C})\)
function (an antilinear map cannot be a complex matrix).

\noindent\textbf{Properties.}
\begin{itemize}
\item \(\hat{\Theta}|\psi^{\uparrow}\rangle = |\psi^{\downarrow}\rangle\)
and \(\hat{\Theta}|\psi^{\downarrow}\rangle = -|\psi^{\uparrow}\rangle\)
(\texttt{timeReversalSpinHalf\_spinHalfUp},
\texttt{\_spinHalfDown}).
\item Additivity \(\hat{\Theta}(v + w) = \hat{\Theta}(v) + \hat{\Theta}(w)\)
(\texttt{timeReversalSpinHalf\_add}).
\item \emph{Antilinearity in the scalar}
\(\hat{\Theta}(c \cdot v) = \overline{c} \cdot \hat{\Theta}(v)\)
(\texttt{timeReversalSpinHalf\_smul}). This is precisely Tasaki's
warning that \(\hat{\Theta}\) is not a linear operator.
\end{itemize}

\noindent\textbf{Theorem (Kramers degeneracy at \(S = 1/2\),
Tasaki eq.~(2.3.8) for half-odd-integer spin).}
\[
  \hat{\Theta}^{2} \;=\; -\hat{1}.
\]
\begin{proof}
By double application of the explicit formula:
\(
  \hat{\Theta}\!\bigl((-b^{*}, a^{*})^{\mathrm T}\bigr)
  = (-(\overline{a^{*}}), \overline{-b^{*}})^{\mathrm T}
  = (-a, -b)^{\mathrm T}
  = -(a, b)^{\mathrm T}.
\)
The double conjugation \(\overline{\overline{x}} = x\) cancels and
the explicit minus sign in the first component survives.
\end{proof}

\noindent\textbf{Theorem (spin sign flip,
Tasaki eq.~(2.3.14)).} For \(\alpha \in \{1, 2, 3\}\),
\[
  \hat{\Theta}\bigl(\hat{S}^{(\alpha)}\,v\bigr)
  \;=\; (-\hat{S}^{(\alpha)})\bigl(\hat{\Theta}\,v\bigr)
  \qquad \text{for every } v : \texttt{Fin}\,2 \to \mathbb{C}.
\]
This is the equivariance form of
\(\hat{\Theta}\,\hat{S}^{(\alpha)}\,\hat{\Theta}^{-1}
   = -\hat{S}^{(\alpha)}\).
The matrix-action form is needed because \(\hat{\Theta}\) is
antilinear and the conjugation
\(\hat{\Theta}^{-1} \cdot \hat{S}^{(\alpha)} \cdot \hat{\Theta}\)
cannot be written as a complex matrix in the conventional sense.
In Lean this is \texttt{timeReversalSpinHalf\_spinHalfOp\{1,2,3\}\_mulVec};
the proof is a direct \texttt{fin\_cases} computation using the
explicit formulae for \(\sigma^{x}, \sigma^{y}, \sigma^{z}\) and
the conjugate ring homomorphism on the entries.

\noindent\textbf{Decomposition} (Tasaki §2.3, p.~26).
The factorisation \(\hat{\Theta} = \hat{u}_2 \cdot \hat{K}\) is
recorded explicitly: \texttt{complexConjugationSpinHalf} defines
the antilinear conjugation
\(\hat{K}(v)_i := \overline{v_i}\) (which is involutive,
additive, and antilinear in the scalar), and
\texttt{timeReversalSpinHalf\_eq\_spinHalfRot2\_pi\_mulVec}
proves
\[
  \hat{\Theta}\,v
  \;=\;
  \hat{u}_2 \cdot \hat{K}(v)
  \;=\;
  \bigl(\texttt{spinHalfRot2}\ \pi\bigr).\texttt{mulVec}\,
    \bigl(\hat{K}(v)\bigr)
\]
for every \(v : \texttt{Fin}\,2 \to \mathbb{C}\). The proof
expands the explicit matrix entries of
\(\hat{u}_2 = \texttt{spinHalfRot2}\,\pi\) and uses
\(I^{2} = -1\) (\texttt{Complex.I\_sq}) to close the residual
arithmetic.

\noindent\textbf{Multi-spin extension and lattice Kramers.}
For finite \(\Lambda\) with \([\texttt{Fintype}\,\Lambda]
[\texttt{DecidableEq}\,\Lambda]\) the tensor product
\(\hat{\Theta}_{\text{tot}} :=
\bigotimes_{x \in \Lambda} \hat{\Theta}_x\) acts on the many-body
Hilbert space \((\Lambda \to \texttt{Fin}\,2) \to \mathbb{C}\) by
\[
  (\hat{\Theta}_{\text{tot}}\,v)(\tau)
  \;:=\;
  \Bigl(\prod_{x \in \Lambda} \varepsilon(\tau\,x)\Bigr)\,
  \overline{v(\texttt{flip}\,\tau)},
\]
where \(\texttt{flip}\,\tau\,x := 1 - \tau\,x\) is the spin-flipped
configuration and \(\varepsilon(0) := 1\),
\(\varepsilon(1) := -1\) is the per-site sign factor coming from
\(\hat{\Theta}|\!\uparrow\rangle = |\!\downarrow\rangle\),
\(\hat{\Theta}|\!\downarrow\rangle = -|\!\uparrow\rangle\). In Lean
this is \texttt{timeReversalSpinHalfMulti}; the helpers
\texttt{flipConfig}, \texttt{timeReversalSign} package the
configuration and sign data.

\noindent\textbf{Theorem (multi-spin Kramers degeneracy at
\(S = 1/2\), Tasaki §2.3 half-odd-integer extension).}
For any finite \(\Lambda\),
\[
  \hat{\Theta}_{\text{tot}}^{2} \;=\; (-1)^{|\Lambda|} \cdot \hat{1}.
\]
\begin{proof}
A direct computation: applying the formula twice and using
\(\overline{\overline{x}} = x\) and the per-site identity
\(\varepsilon(\tau\,x) \cdot \varepsilon((\texttt{flip}\,\tau)(x))
   = -1\) (from \texttt{timeReversalSign\_mul\_flip}) gives
\(\bigl(\prod_x \varepsilon(\tau\,x)\bigr)
  \cdot \bigl(\prod_x \varepsilon((\texttt{flip}\,\tau)(x))\bigr)
  = \prod_x (-1) = (-1)^{|\Lambda|}\),
which factors out to give \((-1)^{|\Lambda|} \cdot v(\tau)\).
\end{proof}

In particular \(\hat{\Theta}_{\text{tot}}^{2} = -\hat{1}\) when
\(|\Lambda|\) is odd (the Kramers minus sign survives) and
\(\hat{\Theta}_{\text{tot}}^{2} = +\hat{1}\) when \(|\Lambda|\)
is even (the minus signs cancel pairwise).

\noindent\textbf{Status / scope.} The \(S = 1/2\) single-spin and
multi-spin cases are closed
(\texttt{LatticeSystem/Quantum/TimeReversalSpinHalf.lean},
\texttt{LatticeSystem/Quantum/TimeReversalMulti.lean}). The
integer-spin case (\(\hat{\Theta}^{2} = +\hat{1}\) for any
\(|\Lambda|\)) is deferred.

\section{Continuum-limit roadmap (Phase A scaffold)}
\label{sec:continuum-limit-roadmap}
%======================================================================

The project's long-term goals include the $\varphi^4$ / Ising
continuum limit and lattice-QCD-style formalisations, both of which
are defined as limits $a \to 0$ of families of finite-spacing lattice
systems. The gap between the current finite-volume matrix framework
and what the continuum limit actually requires was surveyed during
Phase A scoping (consulted codex twice on scope and design choices)
and proposes the following four phases (mirrored on the project
page \href{https://phasetr.github.io/lattice-system/\#continuum-limit-roadmap}%
{\textsl{Continuum-limit roadmap}}):

\begin{description}
\item[Phase A (now, scaffold).] Add a \emph{thin} type-level tag
\begin{verbatim}
class LatticeWithSpacing (L : Type*) where
  spacing : NNReal
\end{verbatim}
attaching a lattice spacing $a : \mathbb{R}_{\geq 0}$ to the abstract
vertex type $\Lambda$. A default instance supplies
$\texttt{spacing} = 1$ on $\texttt{Fin}(N+1)$, so every pre-existing
Hamiltonian is \texttt{rfl}-equivalent to its unit-spacing
specialisation. No geometry, no rescaling, no continuum object.
\item[Phase B (deferred).] Lattice sequences $\Lambda_n$ with
$a_n \to 0$, rescaling of coupling constants
$J_n = \hat h \cdot a_n^{-2+d}$, and lattice-point embeddings in
$\mathbb{R}^d$.
\item[Phase C (deferred).] Operator-valued distribution / GNS /
Hilbert-space infrastructure. Per codex consultation (2026-04-22),
we do \emph{not} generalise
$\texttt{ManyBodyOp}\,\Lambda = \texttt{Matrix}\,\_\,\_\,\mathbb{C}$
to a type class preemptively: existing proofs depend on
Matrix-specific API (\texttt{conjTranspose}, \texttt{exp},
\texttt{trace}, \texttt{mulVec}, entry formulas). Refactor only when
a second concrete backend (infinite-volume Hilbert space, quasi-local
$C^*$-algebra) materialises.
\item[Phase D (deferred).] Running couplings $g : \mathbb{R}_{>0} \to
\mathbb{R}$ and renormalisation-group transformations.
\end{description}

The Phase A module \texttt{LatticeSystem/Lattice/Scale.lean} exports
the class \texttt{LatticeWithSpacing}, the accessor
\texttt{spacingOf\,}$\Lambda$, the default instance
\texttt{instLatticeWithSpacingFinSucc}, and the simp lemmas
\texttt{spacing\_fin\_succ}, \texttt{spacingOf\_fin\_succ}. Tests
pinning these values are in \texttt{LatticeSystem/Tests/Scale.lean}.

\section{Tasaki Appendix A: Mathematical Appendices}
\label{sec:appendix-a}

The book-order content following Chapter 11.  Many results here are the foundations of the
deferred Chapter-11 proof discharges (frustration-free Hamiltonians A.9/A.10 for §11.3/§11.4,
the strong-coupling limit A.11/A.12 for §11.5, and the Perron--Frobenius theorem A.17/A.18 for
§11.2/§11.5, the last largely already formalised in \texttt{Math/PerronFrobenius*} and
\texttt{Math/CollatzWielandt*}).  We prove results where mathlib supports them and axiomatize
the heavy analytic ones, in strict book order (Issue \#4205).

\medskip\noindent\textbf{Theorem A.1 (Lie product formula, axiomatized).}  For finite complex
matrices $A,B$ --- the operator algebra of a finite-volume quantum system --- the exponential
of a sum is the limit of the symmetric split product (Tasaki eq.~(A.2.21)),
\[
  e^{A+B} \;=\; \lim_{N\uparrow\infty}\bigl(e^{A/N}\,e^{B/N}\bigr)^{N}.
\]
Mathlib provides only the commuting case (\texttt{Matrix.exp\_add\_of\_commute}); the general
non-commuting Lie/Trotter formula for bounded operators is a standard but non-trivial
real-analysis result (Tasaki's footnote: ``nineteenth century mathematics''), recorded as the
documented axiom \texttt{lieProductFormula} (\texttt{Math/TasakiAppendixA/LieProduct.lean}) via
\texttt{NormedSpace.exp} and \texttt{Filter.Tendsto}.  It is not on the critical path of any
Chapter-11 proof; to be discharged later.

\medskip\noindent\textbf{Lemmas A.4--A.6 (positive-semidefinite basics, proved).}  Unlike the
Lie product formula, these are available in mathlib and are \emph{proved} (axiom-free,
\texttt{Math/TasakiAppendixA/PosSemidefBasics.lean}).  \textbf{Lemma A.4}
(\texttt{posSemidef\_iff\_eigenvalues\_nonneg}): a Hermitian matrix is positive-semidefinite
iff every eigenvalue is nonnegative.  \textbf{Lemma A.5} (\texttt{PosSemidef.add}): the sum of
two positive-semidefinite matrices is positive-semidefinite,
$\langle\Phi|(\hat A+\hat B)|\Phi\rangle=\langle\Phi|\hat A|\Phi\rangle+\langle\Phi|\hat B|\Phi\rangle\ge0$.
\textbf{Lemma A.6}: $\hat B^\dagger\hat B\ge0$ for any $\hat B$
(\texttt{posSemidef\_conjTranspose\_mul\_self}), and conversely every $\hat A\ge0$ is the square
$\hat A=\hat C^2$ of a positive-semidefinite operator $\hat C=\sqrt{\hat A}$
(\texttt{exists\_posSemidef\_sq\_eq\_of\_posSemidef}, the square root built from the Hermitian
continuous functional calculus \texttt{cfc Real.sqrt}); this square root is \emph{unique}
(\texttt{posSemidef\_sq\_eq\_unique}).

\medskip\noindent\textbf{Theorem A.7 (eigenvalue monotonicity, axiomatized).}  If self-adjoint
matrices satisfy $\hat A\le\hat B$ (the Loewner order), then every sorted eigenvalue of $\hat A$
is at most the corresponding one of $\hat B$: $\forall i,\ \hat A.\texttt{eigenvalues}_0\,i\le\hat
B.\texttt{eigenvalues}_0\,i$ (largest-first; equivalent to Tasaki's $a_j\le b_j$).  This is the
operator-monotonicity content of the min--max principle (eq.~(A.2.30)); mathlib has the
variational principle only for the extreme eigenvalues, so it is recorded as the documented
axiom \texttt{hermitian\_eigenvalues}$_0$\texttt{\_monotone} (\texttt{Math/TasakiAppendixA/EigenvalueMonotone.lean}),
to be discharged later.

\medskip\noindent\textbf{Corollary A.8 (trace of the exponential, axiomatized).}  An
immediate corollary of Theorem A.7, useful in quantum statistical mechanics: $\hat A\le\hat B$
implies $\mathrm{Tr}[e^{\hat A}]\le\mathrm{Tr}[e^{\hat B}]$ (eq.~(A.2.31)), and more generally
$\mathrm{Tr}[f(\hat A)]\le\mathrm{Tr}[f(\hat B)]$ for non-decreasing $f$ (eq.~(A.2.32)).  Recorded as
the documented axiom \texttt{hermitian\_trace\_exp\_monotone}
(\texttt{Math/TasakiAppendixA/TraceExpMonotone.lean}), resting on A.7.

\medskip\noindent\textbf{Lemmas A.9--A.10 (frustration-free Hamiltonians, proved).}  For
$\hat H=\sum_j\hat h_j$ with each local term bounded below, $\hat h_j\ge\varepsilon_j$ (i.e.
$\hat h_j-\varepsilon_j\ge0$): \textbf{Lemma A.9} (\texttt{frustration\_free\_isGroundState}) --- a
simultaneous eigenstate $\hat h_j\Phi=\varepsilon_j\Phi$ is a ground state, with
$\hat H-(\sum_j\varepsilon_j)\ge0$ (lower bound) and $\hat H\Phi=(\sum_j\varepsilon_j)\Phi$
(attained); \textbf{Lemma A.10} (\texttt{frustration\_free\_local\_eigen}, converse) --- if
$\hat H\Phi=(\sum_j\varepsilon_j)\Phi$ then $\hat h_j\Phi=\varepsilon_j\Phi$ for every $j$ (each
$\langle\Phi|\hat h_j-\varepsilon_j|\Phi\rangle\ge0$ summing to $0$, then Lemma A.11).  Both proved
(axiom-free, \texttt{Math/TasakiAppendixA/FrustrationFree.lean}).

\medskip\noindent\textbf{Lemma A.11 (zeros of the energy form, proved).}  If
$\hat A\ge0$ and $\langle\Phi|\hat A|\Phi\rangle=0$ then $\hat A|\Phi\rangle=0$
(\texttt{posSemidef\_mulVec\_eq\_zero\_of\_inner\_zero}); and if $\hat A=\hat B^\dagger\hat B$ then
$\langle\Phi|\hat A|\Phi\rangle=0$ implies $\hat B|\Phi\rangle=0$
(\texttt{conjTranspose\_mul\_self\_mulVec\_eq\_zero\_of\_inner\_zero}).  Stated over the existing
\texttt{Math/RayleighPosSemidefKernel}; we also record Tasaki's angular-momentum corollary
($\hat J^2=\sum_\alpha(\hat J^{(\alpha)})^2$, $\hat J^2|\Phi\rangle=0\Rightarrow\hat J^{(\alpha)}|\Phi\rangle=0$,
\texttt{mulVec\_eq\_zero\_of\_sq\_sum\_inner\_zero}).  All proved (axiom-free,
\texttt{Math/TasakiAppendixA/PosSemidefKernel.lean}).

\medskip\noindent\textbf{Theorem A.12 (strong-coupling effective Hamiltonian, axiomatized).}
For $\hat H_v=\hat H_0+v\hat V$ with $\hat H_0$ self-adjoint, $\hat V\ge0$, $v\ge0$, the
eigenstates whose energy stays finite as $v\uparrow\infty$ are governed by the effective
Schr\"odinger equation $\hat P_0\hat H_0|\Phi\rangle=E|\Phi\rangle$ on $H_0=\ker\hat V$.  We render the
effective equation in weak form ($\langle\psi|\hat H_0|\Phi\rangle=E\langle\psi|\Phi\rangle$ for all
$\psi\in H_0$, avoiding an explicit projection matrix); the axiom
(\texttt{effectiveHamiltonian\_strongCoupling\_limit}, \texttt{Math/TasakiAppendixA/EffectiveLimit.lean})
asserts, in the non-vacuous finite-energy${}\Rightarrow{}$effective direction, that a family
$\Phi_v$ of $\hat H_v$-eigenstates converging to a nonzero $\Phi$ with energy $E_v\to E$ has
$\Phi\in H_0$ and $\Phi$ solving the effective equation.  This underlies the §11.5 effective-theory limits
(\texttt{ttDKernel}/\texttt{ttEffectiveHamiltonian}).  This completes Appendix~A.2.

\medskip\noindent\textbf{§A.3 angular momentum --- Lemma A.14 (su(2) ladder, proved).}  For
angular-momentum operators with $[\hat J^{(1)},\hat J^{(2)}]=i\hat J^{(3)}$, the operator identity
$\hat J^-\hat J^+=\hat J^2-\hat J^{(3)}(\hat J^{(3)}+1)$ (eq.~(A.3.7), \texttt{angLower\_mul\_angRaise})
and the ladder non-vanishing (\texttt{angRaise\_mulVec\_ne\_zero}): on the joint eigenspace
$H_{J,M}$ with $-J\le M<J$ and $\Phi\ne0$, the raised state $\hat J^+\Phi$ is nonzero, since
$\hat J^-\hat J^+\Phi=(J-M)(J+M+1)\Phi$ with $(J-M)(J+M+1)>0$ (Tasaki eq.~(A.3.9)).  The raised state also lies in $H_{J,M+1}$
(\texttt{angRaise\_mem\_eigenspace}): $\hat J^2(\hat J^+\Phi)=J(J+1)(\hat J^+\Phi)$ since
$[\hat J^2,\hat J^+]=0$ (eq.~(A.3.5)) and $\hat J^{(3)}(\hat J^+\Phi)=(M+1)(\hat J^+\Phi)$ since
$[\hat J^{(3)},\hat J^+]=\hat J^+$ (eq.~(A.3.6)).  The mirror lowering direction
(\texttt{angRaise\_mul\_angLower}: $\hat J^+\hat J^-=\hat J^2-\hat J^{(3)}(\hat J^{(3)}-1)$, eq.~(A.3.8);
\texttt{angLower\_mulVec\_ne\_zero}: $\hat J^-\Phi\ne0$ for $-J<M\le J$; \texttt{angLower\_mem\_eigenspace}:
$\hat J^-\Phi\in H_{J,M-1}$ via $[\hat J^2,\hat J^-]=0$ and $[\hat J^{(3)},\hat J^-]=-\hat J^-$) completes
Lemma A.14.  Proved (axiom-free, \texttt{Math/TasakiAppendixA/AngularLadder.lean}).

\medskip\noindent\textbf{Lemma A.15 (spin bound, part 1, proved).}  For self-adjoint
$\hat J^{(\alpha)}$ the norm identities $\|\hat J^\pm\Phi\|^2=\{J(J+1)-M(M\pm1)\}\|\Phi\|^2$
(eq.~(A.3.9), \texttt{angRaise\_normSq}/\texttt{angLower\_normSq}) are nonnegative, so on a nonzero
$H_{J,M}$ state $J(J+1)-M(M+1)\ge0$ and $J(J+1)-M(M-1)\ge0$, giving the spin bound
$-J\le M\le J$ (i.e. $J-M\ge0$ and $J+M\ge0$, \texttt{angMom\_abs\_le\_J}).

\medskip\noindent\textbf{Lemma A.15 (integrality) and Theorem A.13 ($J=n/2$, proved).}  Iterating the
raising operator, \texttt{raiseIter} ($\hat J^+$ applied $k$ times), stays inside $H_{J,M+k}$
(\texttt{raiseIter\_eigenspace}) and stays nonzero while every intermediate magnetic number is
below $J$ (\texttt{raiseIter\_ne\_zero}).  Were $J-M$ non-integral, $\lfloor J-M\rfloor+1$ raisings
would produce a nonzero state with magnetic number ${>}J$, contradicting the spin bound; hence
$J-M\in\mathbb Z_{\ge0}$ (\texttt{angMom\_sub\_mem\_nat}).  Applying the same lemma to the
gauge-reflected operators $(\hat J^{(1)},-\hat J^{(2)},-\hat J^{(3)})$ (same $su(2)$ relations,
$M\mapsto-M$) gives $J+M\in\mathbb Z_{\ge0}$, so $2J=(J-M)+(J+M)\in\mathbb Z_{\ge0}$, i.e. $J=n/2$
(\texttt{angMom\_J\_eq\_half\_nat}, Theorem~A.13).  Proved (axiom-free,
\texttt{Math/TasakiAppendixA/AngularQuantization.lean}).

\medskip\noindent\textbf{Theorem A.16 (SU(2)-multiplet degeneracy, proved).}  For an
$su(2)$-invariant Hamiltonian ($\hat H\hat J^{(\alpha)}=\hat J^{(\alpha)}\hat H$), a joint energy
eigenstate $\Phi\in H_{J,M_0}$ ($\hat H\Phi=E\Phi$) has nonzero same-energy companions in every
$H_{J,J-k}$ for $k\le 2J$ (\texttt{ham\_su2\_multiplet}); as $k$ runs $0,\dots,2J$ this realizes the
$2J+1$ magnetic numbers, so $E$ is at least $(2J+1)$-fold degenerate (eq.~(A.3.14)).  The proof
raises $\Phi$ to the top $H_{J,J}$ ($J-M_0\in\mathbb Z_{\ge0}$ by Lemma~A.15) and lowers $k$ times
via the reflected raising; $\hat H$ preserves each rung because it commutes with $\hat J^\pm$
(\texttt{ham\_mulVec\_raiseIter}).  Proved (axiom-free,
\texttt{Math/TasakiAppendixA/AngularMultiplet.lean}).

\medskip\noindent\textbf{Theorem A.17 (spin-0/spin-1/2 sector suffices).}  For an
$su(2)$-invariant self-adjoint $\hat H$, every energy eigenvalue $E$ has a corresponding eigenstate
$\Phi$ with $\hat J^{(3)}\Phi=0$ or $\hat J^{(3)}\Phi=\tfrac12\Phi$ (\texttt{ham\_eigenstate\_spin\_zero\_or\_half}).
From a joint eigenstate of spin $J=n/2$ (Theorem~A.13) the multiplet (Theorem~A.16) contains the
magnetic number $M=J-\lfloor J\rfloor$, equal to $0$ for even $n$ and $1/2$ for odd $n$ (take
$k=\lfloor J\rfloor$).  The single ingredient beyond finite-matrix algebra---the existence of a
simultaneous eigenstate of the commuting self-adjoint family $\hat H,\hat{\boldsymbol J}^2,\hat J^{(3)}$
in the $E$-eigenspace (\texttt{exists\_joint\_su2\_energy\_eigenstate})---is a documented axiom
(\texttt{Math/TasakiAppendixA/AngularSpinHalfSector.lean}).

\medskip\noindent\textbf{Theorem A.18 (Perron--Frobenius for a real symmetric matrix, proved).}
For a real symmetric $M$ whose shifted matrix $\hat A=c\cdot1-M$ is irreducible (entrywise
nonnegative with strongly connected support---Tasaki's $M_{ij}\le0$ off-diagonal plus
connectivity), there is an eigenvalue $\mu$ of $M$ with a strictly positive eigenvector $v$
($M v=\mu v$, $v_i>0$); $\mu$ is the \emph{lowest} eigenvalue ($\mu\le\lambda$ for every eigenvalue
$\lambda$) and is \emph{nondegenerate} (\texttt{perronFrobenius\_real\_symmetric}).  The proof reuses
the project's Collatz--Wielandt Perron--Frobenius (positive eigenvector) and the eigenspace
simplicity (\texttt{eigenspace\_finrank\_le\_one\_of\_pos\_eigenvec}); the lowest-eigenvalue claim is
the variational $|w|$-argument (eq.~(A.4.1)): for any eigenpair $(\lambda,w)$, $\hat A|w|\ge|c-\lambda|\,|w|$
entrywise, and pairing with $v>0$ gives $c-\mu\ge c-\lambda$.  Proved (axiom-free,
\texttt{Math/TasakiAppendixA/PerronFrobeniusSymmetric.lean}).

\medskip\noindent\textbf{Theorems A.19 (polar) and A.20 (SVD).}  Any square complex matrix
$A$ admits the polar decomposition $A=WC$ with $W$ unitary and $C$ positive-semidefinite
(\texttt{matrix\_polar\_decomposition}, eq.~(A.4.2)) and the singular value decomposition
$A=UDV^\dagger$ with $U,V$ unitary and $D=\operatorname{diagonal}d$ ($d_i\ge0$) the diagonal of
$\sqrt{\lambda_i}$, the $\lambda_i\ge0$ being the eigenvalues of $A^\dagger A$
(\texttt{matrix\_singular\_value\_decomposition}, eqs.~(A.4.3)--(A.4.6)).  \texttt{mathlib} provides
the singular \emph{values} of a linear map but not the constructive matrix factorizations, so both
are recorded as documented axioms (statements kept in book order;
\texttt{Math/TasakiAppendixA/MatrixDecomposition.lean}).

\medskip\noindent\textbf{Theorems A.21 and A.22 (Wigner's theorem).}  Over a
finite-dimensional Hilbert space $H=$ \texttt{Fin\,D}$\to\mathbb C$ (operators $M(H)=$
\texttt{Matrix D D}$\,\mathbb C$): a linear $*$-automorphism $\Gamma$ ((A1)--(A4)) is implemented by a
unitary, $\Gamma(\hat A)=\hat U^\dagger\hat A\hat U$ (\texttt{wignerAutomorphism\_unitary}, eq.~(A.6.1)),
and an antilinear $*$-automorphism ((A4$'$)) by an antiunitary, in matrix form
$\Gamma(\hat A)=\hat U^\dagger\bar{\hat A}\hat U$ via entrywise conjugation
(\texttt{wignerAutomorphism\_antiunitary}, eq.~(A.6.2)).  A map on rank-one projections preserving the
transition probabilities $\operatorname{Tr}[P P']$ is implemented by a unitary or antiunitary
(\texttt{wignerProjection}, Theorem~A.22, eq.~(A.6.10)).  Recorded as documented axioms
(\texttt{Math/TasakiAppendixA/WignerTheorem.lean}).

\medskip\noindent\textbf{Definition A.23 (state) and Theorem A.24 (Banach--Alaoglu).}  Over an
abstract unital complex C*-algebra $A$, a \emph{state} is a weak-$*$-continuous linear functional
$\varphi\in$ \texttt{WeakDual}$\,\mathbb C\,A$ with $\varphi(1)=1$ and $\varphi(\hat A^\dagger\hat A)\ge0$
(\texttt{IsState}, Definition~A.23).  The set of all states is compact in the weak-$*$ topology
(\texttt{stateSpace\_isCompact}, Theorem~A.24, eq.~(A.7.3)) --- the operator-algebraic input for
infinite-volume limits of states; recorded as a documented axiom
(\texttt{Math/TasakiAppendixA/CStarState.lean}).

\medskip\noindent\textbf{Ground states of infinite systems (Definitions A.25/A.27, Theorem A.26).}
Modelling the Hamiltonian commutator $[\hat H,\cdot]$ by a derivation $\delta:A\to A$: a state $\omega$
is a \emph{ground state} iff $0\le\omega(\hat A^\dagger[\hat H,\hat A])$ for all $\hat A$
(\texttt{IsGroundState}, Def~A.25); it has a \emph{nonzero gap} iff some $\gamma>0$ has
$\omega(\hat A^\dagger[\hat H,\hat A])\ge\gamma\,\omega(\hat A^\dagger\hat A)$ whenever $\omega(\hat A)=0$
(\texttt{HasNonzeroGap}, Def~A.27).  A state is a ground state iff for every $L$ its energy
$\omega(\hat H_L)$ is the least over all states agreeing with $\omega$ outside $\Lambda_L$
(\texttt{groundState\_variational}, Theorem~A.26, eq.~(A.7.7), documented axiom;
\texttt{Math/TasakiAppendixA/CStarGroundState.lean}).

\medskip\noindent\textbf{Theorem A.28 (GNS construction).}  Every state $\rho$ on a
C*-algebra $A$ admits a GNS triple $(H_\rho,\pi_\rho,\Omega_\rho)$: a Hilbert space, a
$*$-representation $\pi_\rho:A\to B(H_\rho)$ (a \texttt{StarAlgHom} into $H\to_L H$), and a cyclic
vector with $\rho(\hat A)=\langle\Omega_\rho,\pi_\rho(\hat A)\Omega_\rho\rangle$ (eq.~(A.7.11)) and
$\{\pi_\rho(\hat A)\Omega_\rho\}$ dense (\texttt{gns\_construction}).  Every state is a vector state in
its GNS space; recorded as a documented axiom (the construction is in \texttt{mathlib}'s
\texttt{GelfandNaimarkSegal}; \texttt{Math/TasakiAppendixA/CStarGNS.lean}).

\medskip\noindent\textbf{Appendix A: status and axiomatization policy.}  Tasaki's Appendix A is
fully formalized in book order (A.1--A.28).  The linear-algebra and angular-momentum core that
\texttt{mathlib} supports is proved axiom-free: A.4--A.6, A.9--A.11, A.13--A.16, and A.18
(Perron--Frobenius for a real symmetric matrix).  The remaining items---A.1, A.7/A.8, A.12, the
simultaneous-diagonalization step of A.17, A.19/A.20, and the operator-algebraic A.21--A.28 (Wigner,
states, Banach--Alaoglu, ground states, GNS)---are recorded as faithful documented axioms.  These
axioms are \emph{not} active proof targets here: the heavy functional-analytic and operator-algebraic
theory has its natural home in a dedicated operator-algebra development (possibly contributed to
\texttt{mathlib}), and the axioms wait for that implementation.  The project's policy is to
axiomatize only the appendix results that Tasaki's formalized main text actually uses, to prove the
rest where \texttt{mathlib} allows, and otherwise to keep a faithful book-order axiom; every
theorem's \texttt{\#print axioms} keeps the dependency on these axioms auditable.

\section{Open items / TODO and axioms}
\label{sec:open-items}
%======================================================================

The following Tasaki §2.1 / §2.2 items are \emph{not yet fully proved}
in the Lean development. They are documented here so that future PRs
can pick them up and replace each axiom by a proof (or fill in the
deferred construction).

\subsection*{DONE: Problem 2.1.a for general \texorpdfstring{$S \geq 1$}{S >= 1} (P1d''')}

\textbf{Statement} (Tasaki p.15, solution S.1 p.493): For any spin $S$,
every operator on the single-site Hilbert space $h_0 = \mathbb{C}^{2S+1}$
can be written as a polynomial in $\hat 1, \hat S^{(1)}, \hat S^{(2)},
\hat S^{(3)}$.

\textbf{Status}: \emph{Done in general spin-$S$ form (Issue~\#458 closed
in PR~\#490).} The headline theorem
\texttt{LatticeSystem.Quantum.spinS\_adjoin\_eq\_top}
in
\texttt{Quantum/SpinS/SpanningTheorem.lean} states
\[
  \mathrm{Algebra.adjoin}\,\mathbb{C}\,\{\hat S^{(1)}_N, \hat S^{(2)}_N,
    \hat S^{(3)}_N\}
  \;=\; \top
  \;\subseteq\; \mathrm{Subalgebra}\,\mathbb{C}\,
    (\mathrm{Matrix}\,(\mathrm{Fin}\,(N+1))\,(\mathrm{Fin}\,(N+1))\,\mathbb{C}).
\]
The proof follows Tasaki's solution S.1: diagonal projectors $P_k$
are Lagrange-interpolation polynomials in $\hat S^{(3)}$
(\texttt{spinSDiagProj\_eq\_lagrange\_aeval}); off-diagonal matrix
units $E_{i,j}$ are products of stride-1 ladder units
(\texttt{single\_offset\_succ\_\{,swap\_\}mem\_adjoin}); the
entry-wise decomposition $M = \sum_{i,j} M_{i,j} \cdot E_{i,j}$
(\texttt{Matrix.matrix\_eq\_sum\_single}) closes the spanning. The
earlier concrete-case modules \texttt{pauliBasis} ($S = 1/2$) and
\texttt{spinOne\_decomposition} ($S = 1$) remain as illustrative
specialisations.

\subsection*{DONE / partial: Tasaki §2.5 antiferromagnetic status (issues \#240, \#412)}

The antiferromagnetic Heisenberg / Néel state machinery in Tasaki
§2.5 is largely formalised (chain / 2D / 3D Néel states + per-bond
expectations $-1/4$ + generic graph-centric \texttt{neelStateOf} +
Marshall sign machinery + time-reversal action; see
§\ref{subsec:neel-state-of}). Status of the substructure items:

\begin{itemize}
\item \textbf{DONE} Marshall--Lieb--Mattis Theorem 2.2 (uniqueness + sign
  structure of the AFM ground state) for general spin~$S$ on the
  magnetization-$M$ sector AND on the full Hilbert space, via a
  Perron--Frobenius argument on the Marshall-rotated dressed
  Heisenberg sector matrix. Tracked in Issue~\#412; assembled in
  PRs~\#794--\#869.
  \begin{itemize}
  \item \emph{Sector form (real)}: existence (\#853), eigenvector
    uniqueness (\#854), eigenvalue uniqueness (\#856), bundled
    same-eigenvalue (\#855), bundled forced-eigenvalue (\#859).
  \item \emph{Sector form (complex)}: complex sector matrix (\#857,
    Hermiticity~\#858), real-$\leftrightarrow$-complex
    correspondence (\#858, \#861), complex existence (\#860),
    complex Marshall-positive uniqueness (\#862), bundled complex
    sector ground-state theorem (\#865).
  \item \emph{Full-Hilbert-space form (textbook unit)}: sector
    embedding/restriction (\#867), Heisenberg commutes with
    embedding (\#868), and the full bundled theorem
    \texttt{marshallLiebMattis\_spinS\_heisenbergHamiltonianS\_groundState\_full}
    (\#869) — Tasaki §2.5 Theorem~2.2 ground-state form on the
    actual quantum Heisenberg Hamiltonian: existence + Marshall sign
    structure + same-eigenvalue uniqueness up to positive real
    scalar in the real part, sector-restricted to the
    magnetization-$M$ sector.
  \end{itemize}
\item \textbf{DONE} Theorem~2.3.  The current public statement is
  \texttt{tasaki\_2\_5\_theorem\_2\_3}, with structural proof witnesses
  \texttt{tasaki\_2\_5\_theorem\_2\_3\_bipartiteToy} and
  \texttt{tasaki\_2\_5\_theorem\_2\_3\_of\_bipartiteCompletePositive}.
  PR~\#4082 synchronized the public status rows with these canonical names.
\item \textbf{DONE} Problem 2.5.a (single-cluster ground-state energy
  $-S(1+zS)$ for general spin $S$ and coordination $z$).  The final
  equality wrapper
  \texttt{singleClusterHamiltonianS\_minEigenvalue\_eq\_gs\_of\_predicted\_joint\_witness}
  identifies the Hermitian minimum with the predicted ground-state energy for
  $1\le z$, under the explicit \texttt{[IsAlgClosed \(\mathbb C\)]}
  hypothesis.
\item \textbf{DONE} Problem 2.5.b (lower bound on $E_{\mathrm{GS}}$ via
  2.5.a).  The graph-local chain reaches the closed-form degree wrappers
  \texttt{tasaki25b\_heisenbergHamiltonianOnGraphS\_half\_lower\_bound\_closed\_form}
  and
  \texttt{tasaki25b\_heisenbergHamiltonianOnGraphS\_half\_lower\_bound\_degree\_closed\_form}.
\item \textbf{DONE} Problem 2.5.c (single-site squared expectation
  $\langle\Phi_{\mathrm{GS}},(\hat S_x^{(\alpha)})^2
  \Phi_{\mathrm{GS}}\rangle = S(S+1)/3$ in the AFM ground state).  The
  balanced structural wrapper
  \texttt{singleSiteSpinSquareExpectationS\_all\_axes\_eq\_of\_balanced\_bipartiteCompletePositive}
  proves the all-axis value $N(N+2)/12$ for normalized non-zero Heisenberg
  ground states under the standard balanced bipartite hypotheses.
\item \textbf{DONE} Problem 2.5.d (two-spin correlation under MLM).  The
  endpoint
  \texttt{twoSpinCorrelationS\_re\_neg\_of\_tasaki23\_balanced\_pf\_cross}
  extracts the concrete cross-sublattice negative real correlation from the
  balanced Perron--Frobenius package.
\item \textbf{Theorem~2.4 status}.  The spin-$1/2$ case-(i) target uniqueness
  and zero-magnetization wrappers are live for
  \(-1<\lambda<1\), \(D\ge 0\) as
  \texttt{spinHalf\_anisotropicHeisenbergS\_target\_finrank\_le\_one\_of\_MLM\_casimir\_ladder\_t23\_pf}
  and
  \texttt{spinHalf\_anisotropicHeisenbergS\_target\_groundState\_zero\_magnetization\_of\_MLM\_casimir\_ladder\_t23\_pf}
  in the strict-interior \(D>0\) form, and as the suffixed
  \texttt{\_D\_nonneg} declarations for the \(D=0\) boundary.  The
  \(\lambda=1\) scalar-shift boundary is live as
  \texttt{spinHalf\_anisotropicHeisenbergS\_lambda\_one\_finrank\_le\_one\_of\_MLM\_casimir\_ladder\_t23\_pf}
  and
  \texttt{spinHalf\_anisotropicHeisenbergS\_lambda\_one\_groundState\_zero\_magnetization\_of\_MLM\_casimir\_ladder\_t23\_pf}.
  The general spin-$S$ axis-swap unitary instance is now discharged by
  \texttt{axisSwapUnitarySSpinS} and
  \texttt{anisotropicHeisenbergS\_eigenspace\_finrank\_le\_two\_unconditional\_general};
  the bond-only \(D\ge0\) parity route gives the corresponding
  \texttt{\_D\_nonneg\_general} wrappers.  The new MLM/Casimir endpoint
  wrappers construct the SU(2)-endpoint uniqueness input from the general
  Theorem~2.3 Perron--Frobenius endpoint, so the general spin-$S$ case~(i)
  target uniqueness and zero-magnetization endpoint is live for
  \(-1<\lambda<1\), \(D\ge0\), and at the SU(2) point
  \((\lambda,D)=(1,0)\); the ion-only parity route covers the
  \(\lambda=1\), \(D>0\), \(2\le N\) edge.  Tasaki's case~(ii) now has the
  conditional path/target bridge and a strict-gap target bridge; the full
  \(\texttt{finrank}\le2\) replacement remains explicit.
\end{itemize}

The generic graph-centric \texttt{neelStateOf} (PR \#331) remains the
foundation for the older Néel-state examples, while the later spin-$S$
Marshall--Lieb--Mattis work is tracked by Issue~\#412.

\subsection*{Marshall--Lieb--Mattis Theorem 2.2 (Tasaki §2.5)
              — formalisation outline}
\label{subsec:mlm-spin-s-formalisation}

\textbf{Statement} (Tasaki §2.5 Theorem 2.2, p.~39, generic spin $S$
form on a magnetization sector). Let $G = (V, E)$ be a connected
bipartite graph with sublattice indicator $A: V \to \{0,1\}$, and
let $J: V \times V \to \mathbb{R}_{\ge 0}$ be a real, symmetric,
non-negative coupling supported on bipartite bonds (i.e.\
$J(x,y) = 0$ when $A(x) = A(y)$) with $J(x,y) > 0$ on every edge.
For a fixed magnetization quantum number $M$, the ground state of
the Heisenberg Hamiltonian
\[
  \hat H = \sum_{x,y \in V} J(x,y)\,\hat{\boldsymbol S}_x \cdot
  \hat{\boldsymbol S}_y
\]
restricted to the $M$-sector subspace
$\mathcal H_M = \{\Psi : (V \to \mathrm{Fin}(N+1)) \to \mathbb C \mid
  \Psi$ supported on $\{\sigma : \mathrm{magSum}(\sigma) = M\}\}$
exists, is unique up to a positive real scalar in its real part, and
has the Marshall sign structure
\[
  \Psi_{GS}(\sigma)
  \;=\;
  \mathrm{sign}_A(\sigma)\,v(\sigma) + 0i,
  \qquad v(\sigma) > 0 \quad \text{for } \sigma \in \mathcal H_M,
  \qquad \Psi_{GS}(\sigma) = 0 \quad \text{otherwise}.
\]

\textbf{Proof outline (formalised)}.
\begin{enumerate}
\item \emph{Marshall sign trick} (Tasaki Property~(ii), p.~41).
  The dressed matrix
  $\widetilde H_{\sigma\sigma'} = \mathrm{sign}_A(\sigma)\,
   \mathrm{sign}_A(\sigma')\,H_{\sigma\sigma'}$ has all
  off-diagonal entries non-positive: this is the sign-flip identity
  $\mathrm{sign}_A(\sigma)\,\mathrm{sign}_A(\mathrm{swap}(\sigma)) = -1$
  combined with the explicit two-site $\hat S \cdot \hat S$ matrix
  element on raise/lower bonds.
\item \emph{Bipartite raise/lower reachability} (Tasaki Property~(iii),
  p.~41–42). Any two configurations $\sigma, \sigma'$ in the same
  $M$-sector are connected by a sequence of bipartite raise/lower
  steps in the bipartite complete graph (with the
  intermediate-existence hypothesis on the configuration). On the
  $M$-sector subtype \texttt{magConfigS V N M} this lifts to a chain
  of \texttt{RaiseLowerStepSMagSector} steps.
\item \emph{Perron--Frobenius irreducibility} on the shifted
  dressed sector matrix
  $B := c\,\mathbf 1 - \widetilde H_{\mathrm{sec}}$, $c$ chosen
  larger than every diagonal entry of $\widetilde H_{\mathrm{sec}}$.
  Properties~(i)--(iii) plus reachability give $B$ symmetric,
  non-negative, and irreducible on the sector subtype.
\item \emph{Spectral conclusion}. PF gives a unique strictly positive
  Perron eigenvector $u$ of $B$ at eigenvalue $r > 0$, hence a
  positive eigenvector of $\widetilde H_{\mathrm{sec}}$ at
  $\mu = c - r$. Marshall-conjugation
  $\Phi(\sigma) := \mathrm{sign}_A(\sigma)\,u(\sigma)$ converts this
  to a (real-valued, sign-structured) eigenvector of the
  un-dressed Heisenberg sector matrix at the same $\mu$.
\item \emph{Eigenvalue uniqueness} on the dressed sector matrix
  (symmetry + Rayleigh identity): any two strictly positive
  eigenvectors share the same eigenvalue.
\item \emph{Real $\leftrightarrow$ complex correspondence} on the
  sector via real-part extraction (\#861) and \texttt{ofReal}
  embedding (\#858) — the complex sector eigenvector and the
  real-form sector eigenvector match on real coupling.
\item \emph{Sector $\leftrightarrow$ full-Hilbert-space lift}
  (\#868, \#869): the full Heisenberg Hamiltonian commutes with the
  zero-extension embedding from the sector, since
  $H_{\sigma\sigma'} = 0$ between configurations of different
  $\mathrm{magSum}$ (matrix-level form of $[H, \hat S^z_{\rm tot}] = 0$).
  Inversely, any full-Hilbert-space eigenvector at real $\mu$
  supported on the sector restricts to a sector eigenvector at $\mu$.
  These give a sector $\leftrightarrow$ full bijection on
  Marshall-positive sector-supported eigenvectors at the same $\mu$.
\end{enumerate}

\textbf{Lean entry point}:
\texttt{LatticeSystem/Quantum/SpinS/MagSectorEmbedding.lean},
theorem
\texttt{marshallLiebMattis\_spinS\_heisenbergHamiltonianS\_groundState\_full}
(\#869).

\subsection*{In progress: Tasaki §2.5 Theorem 2.3 ($\gamma$-4, $|A| \neq |B|$ case)}
\label{subsec:mlm-theorem-2-3}

\paragraph{Route correction (2026-05-26).} An earlier
``saturated-ladder-iterate'' route toward
\texttt{tasaki\_2\_5\_theorem\_2\_3} was found \emph{unsound}: it reduced
the theorem to a leaf hypothesis (lowerable attach-sum dominance) that is
false, because the ferromagnetic ladder iterate
$(\hat S^-_{\rm tot})^M\lvert\uparrow\cdots\uparrow\rangle$ is not the
Marshall-positive antiferromagnetic ground state (counterexample:
spin-$1/2$, $V=\{a,b\}$, $A=\{a\}$, $B=\{b\}$, step $M=0\to1$). Its $40$
conditional-wrapper Lean modules were removed. The \emph{sound} route keeps
the per-sector Perron--Frobenius ground state
(\texttt{tasaki\_2\_5\_theorem\_2\_3\_sector\_existence}, sorry-free) and
chains sectors by SU(2) invariance $[\hat H,\hat S^-_{\rm tot}]=0$; the
outside-sector lower bound comes from the toy-Hamiltonian Casimir argument.
See \texttt{.self-local/docs/tasaki-2-5-pf-route-design.md} and Issue \#3542.
Some passages below describe the toy/Casimir prediction machinery that is
retained; references to the removed conditional wrappers are obsolete.

\paragraph{Sound adjacent-sector ladder link
(\texttt{tasaki23\_pf\_ladder\_link\_succ}).} The first reconstructed
step. If $\Psi=\mathtt{magSectorEmbedding}\,\Phi$ is a Heisenberg
eigenvector at energy $\mu$ in magnetization sector $M$ and a
total-Casimir eigenvector at a value $\gamma$ away from the
lowering-kernel value
$\bigl(\tfrac{|V|N}{2}-M\bigr)\bigl(\tfrac{|V|N}{2}-M-1\bigr)$, then
$\hat S^-_{\rm tot}\Psi$ is again a Heisenberg eigenvector at the
\emph{same} $\mu$ (because $[\hat H,\hat S^-_{\rm tot}]=0$, via
\texttt{mulVec\_preserves\_eigenvalue\_of\_commuteS}), is non-zero (by
the total-Casimir non-vanishing criterion
\texttt{tasaki23\_totalSpinSOpMinus\_mulVec\_magSectorEmbedding\_ne\_zero\_of\_casimir\_ne\_kernel\_value}),
and lies in sector $M+1$. Crucially this uses no Marshall positivity of
$\hat S^-_{\rm tot}\Psi$, so it sidesteps the false leaf of the removed
route. Iterating it across the admissible range gives the constancy of
the sector ground energies (the common-energy chain).

\paragraph{Discharging the Casimir input
(\texttt{tasaki23\_pf\_ladder\_link\_succ\_of\_mem\_predictedGS}).} For a
sector vector lying in the predicted toy ground-state subspace
$\mathtt{bipartiteToyGroundStateSubspacePredicted}\,A\,N$, the two Casimir
hypotheses of the ladder link are discharged: in the canonical orientation
$|\neg A|\le|A|$, predicted-GS membership pins the total-Casimir eigenvalue
to $\mathtt{tasaki23PredictedCasimirValue}\,A\,N$, and that value differs
from the lowering-kernel value of every admissible sector strictly below
the right endpoint. So a predicted-GS Heisenberg eigenvector at $\mu$ in
such a sector has $\hat S^-_{\rm tot}\Psi$ a non-zero eigenvector at the
same $\mu$ in the next sector, with no remaining abstract Casimir input.
The remaining physics input is that the per-sector Perron--Frobenius ground
state itself lies in the predicted subspace (the total-spin determination,
to be supplied from the toy-Hamiltonian Casimir minimisation).

\paragraph{Adjacent-sector ground-energy bound
(\texttt{tasaki23\_pf\_sector\_energy\_succ\_le}).} Combining the ladder
link with the per-sector spectral lower bound for Marshall-positive ground
states
(\texttt{heisenbergHamiltonianSReMatrixOnMagSector\_eigenvalue\_ge\_of\_marshallPositive}):
if a predicted-GS eigenvector sits at $\mu$ in admissible sector $M$ below
the right endpoint and sector $M+1$ carries a Marshall-positive ground
state at $\mu'$, then $\mu'\le\mu$. The non-zero lowered image
$\hat S^-_{\rm tot}\Psi$ restricts to a non-zero complex sector eigenvector
at $\mu$, hence a non-zero real-form sector eigenvector at $\mu$ (real or
imaginary part), which the Marshall-positive ground state bounds below by
$\mu'$. With the raising companion this gives the constancy of the sector
ground energies (the common-energy chain), using no Marshall positivity of
the lowered vector.

\paragraph{Raising companion
(\texttt{tasaki23\_pf\_ladder\_link\_pred(\_of\_mem\_predictedGS)},
\texttt{tasaki23\_pf\_sector\_energy\_pred\_le}).} The symmetric raising
versions: $\hat S^+_{\rm tot}$ maps a sector vector (total-Casimir away
from the raising-kernel value) to a non-zero same-energy eigenvector in the
previous sector, with the Casimir hypotheses discharged by predicted-GS
membership strictly above the left endpoint. Hence a predicted-GS
eigenvector at $\mu$ in admissible sector $M+1$ above the left endpoint and
a Marshall-positive ground state at $\mu'$ in sector $M$ give $\mu'\le\mu$.
Combined with the lowering bound this yields $\mu_M=\mu_{M+1}$, the
constancy of the sector ground energies across adjacent admissible sectors.

\paragraph{Constancy step (\texttt{tasaki23\_pf\_sector\_energy\_eq}).}
The two bounds combine by \texttt{le\_antisymm} into the equality
$\mu_M=\mu_{M+1}$ for adjacent admissible sectors interior to the range,
each sector's Marshall-positive ground state being given both as a
full-space Heisenberg eigenvector and a real-form sector-matrix
eigenvector with predicted-GS membership.  This is the inductive
constancy step of the common-energy chain in the sound route, still
modulo the predicted-GS membership of the per-sector ground states (the
total-spin determination).

\paragraph{Overlap argument toward total-spin determination
(\texttt{tasaki23\_marshallPositive\_casimir\_eigenvalue\_eq}).} Two
Marshall-positive vectors in the same magnetization sector have a strictly
positive conjugate overlap: the Marshall signs cancel termwise and what
remains is $\sum_\sigma v_\sigma w_\sigma>0$. Since total-Casimir
eigenvectors at distinct eigenvalues are orthogonal (Hermiticity of
$(\hat S_{\rm tot})^2$), two such Marshall-positive eigenvectors must share
their $(\hat S_{\rm tot})^2$ eigenvalue. This is the bridge that will
transfer the predicted total spin from the toy-Hamiltonian ground state to
the antiferromagnetic Perron--Frobenius ground state, completing the
total-spin determination.

\paragraph{One-dimensionality of the sector ground eigenspace
(\texttt{tasaki23\_heis\_sector\_eigenvec\_proportional\_of\_marshallPositive}).}
The shifted dressed Heisenberg sector matrix is irreducible, so by geometric
simplicity of the Perron eigenvalue any real eigenvector at the Perron value
is a scalar multiple of the positive Perron eigenvector. Marshall-conjugating,
a Marshall-positive Heisenberg sector eigenvector $\varphi$ at $\mu$ has the
property that \emph{every} real Heisenberg sector eigenvector at $\mu$ is
$s\cdot\varphi$. Hence the per-sector Marshall-positive ground state is the
unique eigenvector (up to scale) at its energy; applying $(\hat S_{\rm tot})^2$
to it (which commutes with $\hat H$) therefore yields a scalar multiple of it,
i.e.\ it is a $(\hat S_{\rm tot})^2$ eigenvector — the step that pins its total
spin.

\paragraph{The Perron ground state is a total-Casimir eigenvector
(\texttt{tasaki23\_pf\_groundState\_casimir\_eigenvector}).} Concretely:
$(\hat S_{\rm tot})^2\Phi$ is, by $[\hat H,(\hat S_{\rm tot})^2]=0$, a Heisenberg
eigenvector at the same $\mu$ in the same magnetization sector; its real and
imaginary parts are real Heisenberg sector eigenvectors at $\mu$, hence scalar
multiples $a\Phi$, $b\Phi$ of the Marshall-positive ground state by the
one-dimensionality above. Recombining, $(\hat S_{\rm tot})^2\Phi=(a+bi)\Phi$.
With the overlap transfer (against the toy-Hamiltonian ground state) this pins
the total spin to the predicted value.

\paragraph{Ladder link for the ground state from Casimir non-vanishing
(\texttt{tasaki23\_pf\_groundState\_ladder\_link\_of\_casimir\_ne\_kernel}).}
Combining the previous two facts, the per-sector ground state satisfies the
adjacent-sector ladder link as soon as its total-Casimir eigenvalue is away
from the lowering-kernel value, with no predicted-GS-membership hypothesis:
$\hat S^-_{\rm tot}\Phi$ is a non-zero Heisenberg eigenvector at the same
$\mu$ in the next sector. This is the form in which the sound chain applies
directly to the actual Perron--Frobenius ground state; the single remaining
input is that the ground-state Casimir value equals the predicted value
(hence $\neq$ the kernel value), supplied by the toy-Hamiltonian overlap.

\paragraph{Overlap pin
(\texttt{tasaki23\_pf\_groundState\_casimir\_eq\_predicted\_of\_witness}).}
Tasaki's overlap step (eq.~2.5.12): given a Marshall-positive total-Casimir
eigenvector at the predicted value in sector $M$ (the toy-Hamiltonian ground
state), the per-sector Perron--Frobenius ground state is a total-Casimir
eigenvector at exactly the predicted value. The ground state is a Casimir
eigenvector at some \emph{real} $\gamma$ (the eigenvalue is real because
$(\hat S_{\rm tot})^2$ is Hermitian, lemma
\texttt{isHermitian\_eigenvalue\_star\_eq}); the Marshall-positive overlap with
the witness is non-zero, so the total-Casimir transfer forces
$\gamma=\gamma_{\rm pred}$. This pins the ground-state total spin, leaving only
the existence of the predicted-Casimir witness (the toy ground state) as input.

\paragraph{Toy ground state is a joint Casimir eigenvector
(\texttt{tasaki23\_toy\_groundState\_joint\_casimir\_eigenvector}).} Toward
constructing that witness: generalising the Casimir-eigenvector result to any
operator commuting with $\hat H$ and $\hat S_{\rm tot}^{(3)}$
(\texttt{tasaki23\_pf\_groundState\_commuting\_eigenvector}), and noting that all
three Casimirs $(\hat S_{\rm tot})^2$, $(\hat S_A)^2$, $(\hat S_{\neg A})^2$
commute with the toy Hamiltonian $\hat H_{\rm toy}=\hat H[\,J_{\rm bip}\,]$ (it is
their linear combination), the per-sector Marshall-positive toy ground state is a
\emph{simultaneous} eigenvector of all three Casimirs. Pinning the three
eigenvalues to the predicted values (max sublattice, min total) is the remaining
variational obligation.

\textbf{Goal} (Tasaki §2.5 Theorem 2.3, p.~42). For a connected
bipartite lattice $(\Lambda, B)$ with $|A| \ge 1$ and $|B| \ge 1$,
the AFM Heisenberg ground states have total spin $S_{\rm tot} =
\bigl||A| - |B|\bigr|\,S$ and are $2 S_{\rm tot} + 1$-fold
degenerate. The proof modifies the proof of Theorem 2.2 by
working in a non-zero magnetisation sector $M$ and using the
toy Hamiltonian (2.5.10) with the spin-$S$ Casimir identity
$\hat H_{\rm toy} = \frac{1}{2|\Lambda|}\bigl((\hat S_{\rm tot})^2
- (\hat S_A)^2 - (\hat S_B)^2\bigr)$.

\textbf{Step 1 (PR~\#1042)}: spin-$S$ sublattice spin operators
$\hat S_A^{(\alpha)} := \sum_{x : A x} \mathrm{onSiteS}\,x\,
\hat S_x^{(\alpha)}$ for $\alpha \in \{1, 2, 3\}$, plus
decomposition $\hat S_{\rm tot}^{(\alpha)} = \hat S_A^{(\alpha)}
+ \hat S_{\neg A}^{(\alpha)}$. Defined in
\texttt{LatticeSystem/Quantum/SpinS/SublatticeSpin.lean};
mirrors \texttt{Quantum/MarshallLiebMattis/SublatticeSpin.lean}
(spin-$1/2$ version).

Subsequent steps (planned, multi-PR): sublattice Casimir
$(\hat S_X)^2$, mixed-axes cross-commute, sublattice SU(2)
algebra, toy Hamiltonian Casimir identity, and the eigenvalue
analysis on $|A| \neq |B|$ leading to $S_{\rm tot} =
\bigl||A| - |B|\bigr|\,S$.

\textbf{Step 2 (PR~\#2773, foundation)}: bipartite imbalance
weight in
\texttt{LatticeSystem/Quantum/SpinS/NeelBipartiteWeight.lean}.
Defines \texttt{bipartiteImbalanceWeight A N}~:=~$(|A| - |\neg A|)
\cdot N/2$ — the \emph{signed} $\hat S_{\rm tot}^{(3)}$
magnetization eigenvalue of the Néel orientation
\texttt{neelStateOfS A N}. Its absolute value
$\bigl||A| - |\neg A|\bigr|\cdot N/2 = \bigl||A| - |\neg A|\bigr|\cdot S$
agrees with the Theorem 2.3 prediction under $S = N/2$; the signed
weight equals the predicted total spin in the $|A| \ge |\neg A|$
orientation and is its negative in the opposite orientation.
Packages: bridge with the existing
\texttt{magEigenvalueS\_neelConfigOfS} (the Néel configuration's
$\hat S_{\rm tot}^{(3)}$ eigenvalue equals the imbalance weight);
membership
\texttt{neelStateOfS A N $\in$ magSubspaceS $\Lambda$ N (bipartiteImbalanceWeight A N)}
via \texttt{basisVecS\_mem\_magSubspaceS}; two edge cases
(balanced $|A|=|\neg A| \Rightarrow$ weight $=0$, recovering the
$\gamma$-3 sector; saturated $|\neg A| = 0 \Rightarrow$ weight
$= m_{\max}$, recovering the all-up sector); and
real-part non-negativity when $|A| \ge |\neg A|$. This is the
first foundational step toward formalising the §2.5 Theorem 2.3
prediction $S_{\rm tot} = \bigl||A| - |B|\bigr|\,S$ at the
magnetisation level. Subsequent steps will address the lower
bound via the toy-Hamiltonian Casimir decomposition (using
PR~\#1042's sublattice operators) and the variational argument.

\textbf{Step 3 (PR~\#2774)}: toy-Hamiltonian eigenvalue on
simultaneous Casimir eigenvectors, in
\texttt{LatticeSystem/Quantum/SpinS/ToyHamiltonianJointEigenvalue.lean}.
The closed form
$\hat H_{\rm toy_S} = (\hat S_{\rm tot})^2 - (\hat S_A)^2 -
(\hat S_{\neg A})^2$ (PR \#1056) gives: for any joint eigenvector
$v$ of $(\hat S_{\rm tot})^2, (\hat S_A)^2, (\hat S_{\neg A})^2$ at
eigenvalues $(\alpha, \beta_A, \beta_B)$, $\hat H_{\rm toy_S}\cdot
v = (\alpha - \beta_A - \beta_B)\cdot v$
(\texttt{heisenbergToyHamiltonianS\_mulVec\_of\_jointCasimirEigenvector})
and $\langle v\,|\,\hat H_{\rm toy_S}\,v\rangle = (\alpha - \beta_A
- \beta_B)\,\langle v, v\rangle$
(\texttt{neelStateOfS\_heisenbergToyHamiltonianS\_expectation\_via\_joint}).
This is the eigenvalue identity used in Tasaki's Theorem 2.3
proof: minimising the toy energy over admissible joint
$(\alpha, \beta_A, \beta_B)$ subject to the angular-momentum
addition constraints gives the predicted ground state at
$S_{\rm tot} = \bigl||A| - |\neg A|\bigr|\,S$. Subsequent $\gamma$-4 PRs
will package the explicit Casimir-eigenvalue upper bounds
$\beta_A \le s_A(s_A+1)$, $\beta_B \le s_B(s_B+1)$ with
$s_A = |A|\cdot N/2$, $s_B = |\neg A|\cdot N/2$, the triangle-
inequality lower bound on $\alpha$, and the variational extraction
of the Theorem 2.3 ground-state energy and total spin.

\textbf{Step 4 (PR~\#2775)}: subspace-level form of Step 3 in
\texttt{LatticeSystem/Quantum/SpinS/ToyHamiltonianJointEigenspace.lean}.
Lifts PR \#2774's pointwise identity to the
\texttt{Submodule}-inclusion
$\texttt{eigenspace}((\hat S_{\rm tot})^2)_\alpha \sqcap
\texttt{eigenspace}((\hat S_A)^2)_{\beta_A} \sqcap
\texttt{eigenspace}((\hat S_{\neg A})^2)_{\beta_B} \le
\texttt{eigenspace}(\hat H_{\rm toy_S})_{\alpha - \beta_A - \beta_B}$
via \texttt{Module.End.mem\_eigenspace\_iff}. This is the natural
operator-level form that the upcoming $\gamma$-4 spectrum analysis will
chain through.

\textbf{Step 5 (PR~\#2776)}: Néel-state toy-Hamiltonian
expectation via the Casimir decomposition in
\texttt{LatticeSystem/Quantum/SpinS/NeelToyHamiltonianViaCasimir.lean}.
The Néel state $|\Phi_{\rm N\acute eel}\rangle$ is an eigenvector
of the sublattice Casimirs $(\hat S_A)^2$ and $(\hat S_{\neg A})^2$
at their maximum eigenvalues $s_A(s_A+1)$, $s_B(s_B+1)$ (PR \#1067)
but \emph{not} of $(\hat S_{\rm tot})^2$. Using PR \#1056's closed
form, the expectation splits as
$\langle\Phi_{\rm N\acute eel}|\hat H_{\rm toy_S}|\Phi_{\rm N\acute eel}\rangle
= \langle\Phi_{\rm N\acute eel}|(\hat S_{\rm tot})^2|\Phi_{\rm N\acute eel}\rangle
- s_A(s_A+1)\,\langle\Phi_{\rm N\acute eel},\Phi_{\rm N\acute eel}\rangle
- s_B(s_B+1)\,\langle\Phi_{\rm N\acute eel},\Phi_{\rm N\acute eel}\rangle$.
Combining with PR \#1093's $(\hat S_{\rm tot})^2$ expectation
$(s_A-s_B)^2 + (s_A+s_B)$ gives the known value $-2 s_A s_B =
-|A|\cdot|\neg A|\cdot N^2/2$, matching the direct PR \#1070
computation. Useful for subsequent variational comparisons
against the predicted toy-Hamiltonian ground state.

\textbf{Step 6 (PR~\#2777)}: complement-orientation Néel toy-
Hamiltonian expectation in
\texttt{LatticeSystem/Quantum/SpinS/NeelToyComplementExpectation.lean}.
$\langle\Phi_{\rm N\acute eel}(\neg A)|\hat H_{\rm toy_S}(A)|
\Phi_{\rm N\acute eel}(\neg A)\rangle = -|A|\cdot|\neg A|\cdot N^2/2$,
same value as the original-orientation expectation (PR \#1093).
The toy Hamiltonian is sublattice-swap invariant
($\hat H_{\rm toy_S}(A) = \hat H_{\rm toy_S}(\neg A)$ via
\texttt{sublatticeSpinSDot\_complement\_comm}), and the complement
Néel basis configuration produces the same per-bond
$m_x\cdot m_y = -N^2/4$ contribution. Completes the
orientation-pair (use $\Phi_{\rm N\acute eel}(A)$ for $|A|
\ge |\neg A|$, $\Phi_{\rm N\acute eel}(\neg A)$ for $|\neg A|
\ge |A|$) at the same toy-Hamiltonian value, reflecting the
choice-of-orientation freedom in Tasaki §2.5 Theorem 2.3.

\textbf{Step 7 (PR~\#2778)}: predicted toy-Hamiltonian ground-
state subspace and minimum energy in
\texttt{LatticeSystem/Quantum/SpinS/BipartiteToyMinEnergy.lean}.
For $|A| \ge |\neg A|$:
\begin{itemize}
\item \texttt{bipartiteToyMinEnergyPredicted A N}~:=~
$(s_A - s_B)(s_A - s_B + 1) - s_A(s_A+1) - s_B(s_B+1)$, with
$s_A = |A|\cdot N/2, s_B = |\neg A|\cdot N/2$. Ring computation
simplifies to $-|A|\cdot|\neg A|\cdot N^2/2 - |\neg A|\cdot N$.
\item \texttt{bipartiteToyGroundStateSubspacePredicted A N}: the
joint eigenspace of $((\hat S_{\rm tot})^2, (\hat S_A)^2,
(\hat S_{\neg A})^2)$ at the three target eigenvalues
(sublattice Casimirs at maximum, total Casimir at triangle-
inequality minimum).
\item Inclusion: predicted GS subspace $\le$ \texttt{eigenspace}
$(\hat H_{\rm toy_S})_{\texttt{bipartiteToyMinEnergyPredicted}}$,
direct application of PR \#2775's
\texttt{heisenbergToyHamiltonianS\_jointCasimirEigenspace\_le\_eigenspace}.
\end{itemize}
Minimality of \texttt{bipartiteToyMinEnergyPredicted} (variational
lower bound + non-emptiness of the joint subspace) is the next
$\gamma$-4 milestone, not yet established.

\textbf{Step 8 (PR~\#2779)}: structural properties of
\texttt{bipartiteToyMinEnergyPredicted} in
\texttt{LatticeSystem/Quantum/SpinS/BipartiteToyMinEnergyProperties.lean}.
Three edge cases / structural identities:
\begin{itemize}
\item \emph{Saturated edge case} $|\neg A| = 0$: predicted energy
$= 0$ (trivial empty-cross-bond limit).
\item \emph{Néel-state gap}:
\texttt{bipartiteToyMinEnergyPredicted}~$=$
$\langle\Phi_{\rm N\acute eel}|\hat H_{\rm toy_S}|\Phi_{\rm N\acute eel}\rangle
- |\neg A|\cdot N$. The Néel expectation
$-|A|\cdot|\neg A|\cdot N^2/2$ exceeds the predicted minimum by
$|\neg A|\cdot N \ge 0$ — Néel is \emph{not} the toy-H ground
state in the non-degenerate regime.
\item \emph{Balanced edge case} $|A| = |\neg A|$: imbalance term
vanishes, predicted energy $= -2 s_A(s_A+1)$ (the
``anti-ferromagnetic-saturation'' value).
\end{itemize}

\textbf{Step 9 (PR~\#2780)}: complement-orientation gap identity
in
\texttt{LatticeSystem/Quantum/SpinS/BipartiteToyMinEnergyComplement.lean}.
\texttt{bipartiteToyMinEnergyPredicted A N -
bipartiteToyMinEnergyPredicted ($\neg$A) N $=$
$(|A| - |\neg A|)\cdot N$}. The signed-difference imbalance in
PR \#2778 is orientation-dependent; the actual Tasaki prediction
$\min(E(A), E(\neg A))$ uses the orientation with the larger
sublattice. Direct ring computation from PR \#2779's simplified
form via filter congruence $\neg\neg A = A$.

\textbf{Step 10 (PR~\#2781)}: sublattice-swap symmetric predicted
min energy in
\texttt{LatticeSystem/Quantum/SpinS/BipartiteToyMinEnergySymm.lean}.
\texttt{bipartiteToyMinEnergyPredictedSymm A N}~:=~
$-|A|\cdot|\neg A|\cdot N^2/2 - \min(|A|, |\neg A|)\cdot N$, the
true Tasaki §2.5 Theorem 2.3 prediction independent of
orientation. Bridges:
\texttt{bipartiteToyMinEnergyPredictedSymm}~$=$~
\texttt{bipartiteToyMinEnergyPredicted A N} when
$|\neg A| \le |A|$ (via \texttt{Nat.min\_eq\_right}); equals
\texttt{bipartiteToyMinEnergyPredicted ($\neg$A) N} when
$|A| \le |\neg A|$ (via \texttt{Nat.min\_eq\_left} + filter
congruence). Sublattice-swap invariant:
\texttt{Symm ($\neg$A) N $=$ Symm A N} via
\texttt{Nat.min\_comm}.

\textbf{Step 11 (PR~\#2782)}: non-emptiness of the predicted GS
subspace in the saturated edge case ($|\neg A| = 0$) in
\texttt{LatticeSystem/Quantum/SpinS/BipartiteToyGroundStateSaturated.lean}.
For $|\neg A| = 0$, the all-up state $|\sigma_\top\rangle$ is in
\texttt{bipartiteToyGroundStateSubspacePredicted A N}, the first
concrete realisation of PR \#2778's predicted ground-state
subspace. Three eigenvector witnesses:
\begin{itemize}
\item $(\hat S_{\rm tot})^2\cdot|\sigma_\top\rangle =
s_A(s_A+1)\cdot|\sigma_\top\rangle$ (PR \#878 with $|V| = |A|$).
\item $(\hat S_A)^2\cdot|\sigma_\top\rangle =
s_A(s_A+1)\cdot|\sigma_\top\rangle$ (PR \#1067).
\item $(\hat S_{\neg A})^2\cdot|\sigma_\top\rangle =
0\cdot|\sigma_\top\rangle$ (PR \#1067 with $|\neg A| = 0$
collapsing the eigenvalue formula).
\end{itemize}

\textbf{Step 12 (PR~\#2783)}: all-down version of Step 11 in
\texttt{LatticeSystem/Quantum/SpinS/BipartiteToyGroundStateSaturatedLast.lean}.
For $|\neg A| = 0$, $|\sigma_\bot\rangle$ is also in
\texttt{bipartiteToyGroundStateSubspacePredicted A N}, via the
`\_last` variants of the same three eigenvector lemmas (PR \#879
total Casimir on all-down, PR \#1067 sublattice Casimir on
all-down). Together with PR \#2782, gives two distinct elements
of the predicted GS subspace for $N \ge 1$ (where
$|\sigma_\top\rangle \ne |\sigma_\bot\rangle$; the $N = 0$ spin-$0$
case collapses to a 1-dim space). Via the
saturated-ferromagnet ladder of PRs \#896 / \#2768, the entire
$(2 s_A + 1)$-dim irrep sits inside for any $N$, confirming the
subspace has the correct $(2 m_{\max} + 1)$ dimension in the
saturated case.

\textbf{Step 13 (PR~\#2784)}: sublattice $\leftrightarrow$ total
spin identities at the saturated edge case in
\texttt{LatticeSystem/Quantum/SpinS/SublatticeEqTotalOfCardZero.lean}.
When $|\neg A| = 0$:
\begin{itemize}
\item $\hat S_{\neg A}^{(\alpha)} = 0$ and $\hat S_{\neg A}^2 = 0$
(empty-sum vanishing).
\item $\hat S_A^{(\alpha)} = \hat S_{\rm tot}^{(\alpha)}$ for each
$\alpha$ (via the decomposition
$\hat S_{\rm tot}^{(\alpha)} = \hat S_A^{(\alpha)} +
\hat S_{\neg A}^{(\alpha)}$ + empty-sum vanishing).
\item $\hat S_A^2 = \hat S_{\rm tot}^2$.
\end{itemize}
Foundation for the dimension bound on the predicted GS subspace
in the saturated case:
$\texttt{eigenspace}(\hat S_{\rm tot}^2)_\alpha \le
\texttt{bipartiteToyGroundStateSubspacePredicted A N}$ whenever
$|\neg A| = 0$ and $\alpha = s_A(s_A+1)$.

\textbf{Step 14 (PR~\#2785)}: eigenspace bridge for the
saturated case in
\texttt{LatticeSystem/Quantum/SpinS/BipartiteToyGSEigenspaceBridge.lean}.
\texttt{totalSpinSSquaredEigenspace\_le\_bipartiteToyGroundStateSubspacePredicted\_of\_cardNotA\_zero}:
$\texttt{eigenspace}((\hat S_{\rm tot})^2)_{s_A(s_A+1)} \le
\texttt{bipartiteToyGroundStateSubspacePredicted A N}$ when
$|\neg A| = 0$. Each of the three joint Casimir conditions is
verified: (i) the $(\hat S_{\rm tot})^2$ condition is by
hypothesis; (ii) the $(\hat S_A)^2$ condition reduces to the
$(\hat S_{\rm tot})^2$ condition via PR \#2784's
$\hat S_A^2 = \hat S_{\rm tot}^2$; (iii) the $(\hat S_{\neg A})^2$
condition is trivial since the operator is the zero matrix.
Combined with PR \#2768's
\texttt{saturatedFerromagnetJointEigenspace = span(ladderIterateUp)}
and the finrank $2 m_{\max}+1$, this gives the dimension bound
$\dim(\texttt{predicted GS subspace}) \ge 2 s_A + 1$ at the
saturated edge case, matching the Theorem 2.3 prediction
$2 S_{\rm tot} + 1$.

\textbf{Step 15 (PR~\#2786)}: §2.4 $\leftrightarrow$ $\gamma$-4 bridge
at the saturated case in
\texttt{LatticeSystem/Quantum/SpinS/SaturatedFMJointLePredictedGS.lean}.
For $|\neg A| = 0$ and any coupling $J$,
$\texttt{saturatedFerromagnetJointEigenspace } J N \le
\texttt{bipartiteToyGroundStateSubspacePredicted A N}$.
Direct corollary of PR \#2785 applied to the
$(\hat S_{\rm tot})^2$-eigenspace component of the §2.4 joint
eigenspace, with $m_{\max} = s_A$ (forced by $|V| = |A|$).
Together with PRs \#2768 (joint = span ladder) and \#2769
($\dim = 2 m_{\max}+1$), this realises the dimension lower bound
$\dim(\texttt{predicted GS subspace}) \ge 2 s_A + 1$ concretely
at the saturated edge case.

\textbf{Step 16 (PR~\#2787)}: predicted GS subspace finrank lower
bound at the saturated case in
\texttt{LatticeSystem/Quantum/SpinS/BipartiteToyGSFinrankSaturated.lean}.
\texttt{bipartiteToyGroundStateSubspacePredicted\_finrank\_ge\_succ\_card\_mul\_N\_of\_cardNotA\_zero}:
when $|\neg A| = 0$,
$|V|\cdot N + 1 \le
\dim_\mathbb{C}(\texttt{bipartiteToyGroundStateSubspacePredicted A N})$.
Direct corollary of PR \#2786 + PR \#2769's
\texttt{saturatedFerromagnetJointEigenspace\_finrank\_eq} via
\texttt{Submodule.finrank\_mono}. This realises the Tasaki §2.5
Theorem 2.3 degeneracy lower bound $2 S_{\rm tot} + 1$ at the
saturated edge case where $S_{\rm tot} = m_{\max} = |V|\cdot S$.

\textbf{Step 17 (PR~\#2788)}: opposite saturated edge case
($|A| = 0$) in
\texttt{LatticeSystem/Quantum/SpinS/BipartiteToyGSFinrankOppositeSaturated.lean}.
\texttt{bipartiteToyGroundStateSubspacePredicted\_complement\_finrank\_ge\_succ\_card\_mul\_N\_of\_cardA\_zero}:
when $|A| = 0$, the complement-orientation predicted GS subspace
$\texttt{bipartiteToyGroundStateSubspacePredicted ($\neg$A) N}$
has finrank $\ge |V|\cdot N + 1$. Application of PR \#2787 to the
flipped sublattice, reducing $|\neg\neg A| = 0$ to $|A| = 0$ via
filter congruence. Together with PR \#2787, both saturated edge
cases now have the dimension lower bound, with the natural
orientation chosen automatically.

\textbf{Step 18 (PR~\#2789)}: $\hat H_{\rm toy_S}$-eigenspace
finrank lower bound at the saturated case in
\texttt{LatticeSystem/Quantum/SpinS/BipartiteToyEigenspaceFinrankSaturated.lean}.
When $|\neg A| = 0$, the $\hat H_{\rm toy_S}$-eigenspace at
\texttt{bipartiteToyMinEnergyPredicted A N} has finrank
$\ge |V|\cdot N + 1$. Direct chain of PR \#2775
(\texttt{predicted GS} $\subseteq$ \texttt{eigenspace}) and
PR \#2787 (\texttt{predicted GS} finrank $\ge |V|\cdot N + 1$)
via \texttt{Submodule.finrank\_mono}. This is the operator-level
version of the Theorem 2.3 degeneracy $2 S_{\rm tot} + 1$ lower
bound, applied directly to the $\hat H_{\rm toy_S}$-eigenspace.

\textbf{Step 19 (PR~\#2790)}: opposite saturated case
$\hat H_{\rm toy_S}(\neg A)$-eigenspace finrank lower bound in
\texttt{LatticeSystem/Quantum/SpinS/BipartiteToyEigenspaceFinrankOppositeSaturated.lean}.
Mirror of PR \#2789 via PR \#2788. When $|A| = 0$, the
$\hat H_{\rm toy_S}(\neg A)$-eigenspace at
\texttt{bipartiteToyMinEnergyPredicted ($\neg$A) N} has finrank
$\ge |V|\cdot N + 1$. Since the toy Hamiltonian is sublattice-
swap invariant, this also applies to the original-orientation
toy Hamiltonian.

\textbf{Step 20 (PR~\#2791)}: non-positivity of the predicted
min energy in
\texttt{LatticeSystem/Quantum/SpinS/BipartiteToyMinEnergyNonpos.lean}.
$(\texttt{bipartiteToyMinEnergyPredicted A N}).\texttt{re} \le 0$
for any sublattice $A$ and any $N$. Direct from the simplified
form $-|A|\cdot|\neg A|\cdot N^2/2 - |\neg A|\cdot N$ (PR \#2779),
both terms non-positive. Sanity check consistent with
antiferromagnetic energy lowering.

\textbf{Step 21 (PR~\#2792)}: predicted GS subspace finrank
\emph{equality} at saturated case in
\texttt{LatticeSystem/Quantum/SpinS/BipartiteToyGSFinrankEqSaturated.lean}.
Two theorems:
\begin{itemize}
\item Reverse bridge:
\texttt{bipartiteToyGroundStateSubspacePredicted A N $\le$
saturatedFerromagnetJointEigenspace 0 N} when $|\neg A| = 0$,
since $H_0 = 0$ and the eigenvalues match.
\item Finrank equality:
$\dim(\texttt{bipartiteToyGroundStateSubspacePredicted A N}) =
|V|\cdot N + 1$ when $|\neg A| = 0$. Combines the reverse bridge
(with PR \#2769's finrank equality on the saturated FM joint)
with PR \#2787's lower bound via \texttt{le\_antisymm}.
\end{itemize}
Operator-level confirmation that the predicted GS subspace is
\emph{exactly} the $(2 m_{\max}+1)$-dim spin-$m_{\max}$
irreducible representation at the saturated edge case, with
equality matching the Tasaki §2.5 Theorem 2.3 degeneracy
$2 S_{\rm tot} + 1$.

\textbf{Step 22 (PR~\#2793)}: opposite saturated case finrank
equality in
\texttt{LatticeSystem/Quantum/SpinS/BipartiteToyGSFinrankEqOppositeSaturated.lean}.
$\dim(\texttt{bipartiteToyGroundStateSubspacePredicted ($\neg$A) N})
= |V|\cdot N + 1$ when $|A| = 0$. Mirror of PR \#2792 via
orientation flip + $\neg\neg A = A$ filter congruence. Both
saturated edge cases now have the equality form, confirming the
predicted GS subspace dimension matches the Theorem 2.3 prediction
with equality in both orientations.

\textbf{Step 23 (PR~\#2794)}: $\hat H_{\rm toy_S}$-invariance of
the predicted GS subspace in
\texttt{LatticeSystem/Quantum/SpinS/BipartiteToyGSToyInvariant.lean}.
$\texttt{Submodule.map}\,\hat H_{\rm toy_S}.\texttt{mulVecLin}\,
(\texttt{predicted GS}) \le \texttt{predicted GS}$. Direct from
PR \#2775 (\texttt{predicted GS} $\subseteq$ \texttt{eigenspace} at
predicted min) + scalar-multiplication closure of the submodule.
Structural sanity check: the predicted GS subspace is
$\hat H_{\rm toy_S}$-invariant.

\textbf{Step 24 (PR~\#2795)}: Casimir-invariances of the
predicted GS subspace in
\texttt{LatticeSystem/Quantum/SpinS/BipartiteToyGSCasimirInvariant.lean}.
Trivial invariance under each of $(\hat S_{\rm tot})^2$,
$(\hat S_A)^2$, $(\hat S_{\neg A})^2$ from the corresponding
eigenspace membership components + scalar closure. Confirms the
predicted GS is closed under all three Casimirs, consistent with
its definition as their joint eigenspace.

\textbf{Step 25 (PR~\#2796)}: submodule equality
\texttt{predicted GS = saturated FM joint at $J = 0$} in the
saturated case in
\texttt{LatticeSystem/Quantum/SpinS/BipartiteToyGSEqSaturatedJointZero.lean}.
When $|\neg A| = 0$,
\texttt{bipartiteToyGroundStateSubspacePredicted A N $=$
saturatedFerromagnetJointEigenspace 0 N}. Combines PR \#2786
(saturated FM joint $\subseteq$ predicted GS) with PR \#2792's
reverse bridge via \texttt{le\_antisymm}. Identifies the $\gamma$-4
predicted GS subspace with the §2.4 saturated-ferromagnet
ground-state subspace at the saturated edge case (specialised to
$J = 0$).

\textbf{Step 26 (PR~\#2797)}: opposite-orientation submodule
equality in
\texttt{LatticeSystem/Quantum/SpinS/BipartiteToyGSEqOppositeSaturatedJointZero.lean}.
When $|A| = 0$, \texttt{bipartiteToyGroundStateSubspacePredicted
($\neg$A) N $=$ saturatedFerromagnetJointEigenspace 0 N}. Mirror of
PR \#2796 via the orientation flip and $\neg\neg A = A$ filter
congruence.

\textbf{Step 27 (PR~\#2798)}: $\hat S^z_{\rm tot}$-invariance of
the predicted GS subspace in
\texttt{LatticeSystem/Quantum/SpinS/BipartiteToyGSTotalSpinSOp3Invariant.lean}.
$\texttt{Submodule.map}\,\hat S^z_{\rm tot}.\texttt{mulVecLin}\,
(\texttt{predicted GS}) \le \texttt{predicted GS}$. The
$\hat S^z_{\rm tot}$ operator commutes with all three Casimirs
($(\hat S_{\rm tot})^2$, $(\hat S_A)^2$, $(\hat S_{\neg A})^2$), so
it preserves each Casimir eigenspace and hence their meet.
Standard ``magnetization-preserving operator acts within each
Casimir subspace'' picture.

\textbf{Step 28 (PR~\#2799)}: $\hat S^{(1)}_{\rm tot}$ and
$\hat S^{(2)}_{\rm tot}$ invariance of the predicted GS subspace
in
\texttt{LatticeSystem/Quantum/SpinS/BipartiteToyGSCartesianInvariant.lean}.
Mirror of PR \#2798 for the other two Cartesian axes. Each
$\hat S^{(\alpha)}_{\rm tot}$ commutes with all three Casimirs, so
it preserves each Casimir eigenspace and hence the predicted GS.
Together with PR \#2798, the predicted GS is closed under all three
Cartesian generators — the full SU(2) symmetry at the operator
level, a necessary structural condition for the irreducible-
representation interpretation of Theorem 2.3.

\textbf{Step 29 (PR~\#2800)}: $\hat S^+_{\rm tot}$ /
$\hat S^-_{\rm tot}$ invariance of the predicted GS subspace in
\texttt{LatticeSystem/Quantum/SpinS/BipartiteToyGSLadderInvariant.lean}.
Derived commutations
$[(\hat S_A)^2, \hat S^\pm_{\rm tot}] = 0$ via
$\hat S^\pm_{\rm tot} = \hat S^{(1)}_{\rm tot} \pm i \hat S^{(2)}_{\rm tot}$
and the Cartesian commutations of PR \#2799. Same proof pattern as
PR \#2798 gives the predicted-GS invariance under each ladder
operator. Together with PRs \#2798 / \#2799, full $su(2)$ algebra
closure: predicted GS is invariant under Cartan and ladders, hence
under the entire Lie algebra.

\textbf{Step 30 (PR~\#2801)}: simplification
\texttt{saturatedFerromagnetJointEigenspace 0 N $=$
($\hat S_{\rm tot})^2$ eigenspace} in
\texttt{LatticeSystem/Quantum/SpinS/SaturatedFMJointZeroEqTotalSpinSSquaredEigenspace.lean}.
With $J=0$, the H-eigenspace component trivialises (every vector
satisfies $0\cdot v = 0\cdot v$), so the joint reduces to the
Casimir eigenspace alone. Combined with PR \#2796, identifies the
saturated-case predicted GS subspace as the maximum-Casimir
eigenspace.

\textbf{Step 31 (PR~\#2802)}: composition of PRs \#2796 and \#2801
in
\texttt{LatticeSystem/Quantum/SpinS/BipartiteToyGSEqCasimirEigenspaceSaturated.lean}.
When $|\neg A| = 0$,
\texttt{bipartiteToyGroundStateSubspacePredicted A N $=$
Module.End.eigenspace ($\hat S_{\rm tot})^2$.mulVecLin
saturatedFerromagnetCasimirEigenvalueS}. The saturated-case
predicted GS is \emph{just} the maximum-Casimir
$(\hat S_{\rm tot})^2$-eigenspace, free from any auxiliary
Hamiltonian.

\textbf{Step 32 (PR~\#2803)}: opposite saturated case mirror of
PR \#2802 in
\texttt{LatticeSystem/Quantum/SpinS/BipartiteToyGSEqCasimirEigenspaceOppositeSaturated.lean}.
When $|A| = 0$, \texttt{bipartiteToyGroundStateSubspacePredicted
($\neg$A) N $=$ Module.End.eigenspace ($\hat S_{\rm tot})^2$.mulVecLin
saturatedFerromagnetCasimirEigenvalueS}. Apply PR \#2802 to
$\neg A$ with $|\neg\neg A| = 0$ reducing to $|A| = 0$.

\textbf{Step 33 (PR~\#2804)}: concrete basis identification in
\texttt{LatticeSystem/Quantum/SpinS/BipartiteToyGSEqSpanLadderSaturated.lean}.
When $|\neg A| = 0$,
\texttt{bipartiteToyGroundStateSubspacePredicted A N $=$ Submodule.span
$\mathbb{C}$ (Set.range (ladderIterateUp $\Lambda$ N))}. Composition
of PR \#2796 + PR \#2768. The saturated-case predicted GS subspace
is precisely the linear span of the $(2 m_{\max}+1)$ ladder
iterates $\{(\hat S^-_{\rm tot})^k\,|\sigma_\top\rangle\}_k$.

\textbf{Step 34 (PR~\#2805)}: opposite saturated case mirror of
PR \#2804 in
\texttt{LatticeSystem/Quantum/SpinS/BipartiteToyGSEqSpanLadderOppositeSaturated.lean}.
When $|A| = 0$,
\texttt{bipartiteToyGroundStateSubspacePredicted ($\neg$A) N $=$
Submodule.span $\mathbb{C}$ (Set.range (ladderIterateUp $\Lambda$ N))}.

\textbf{Step 35 (PR~\#2806)}: toy Hamiltonian vanishes at the
saturated case in
\texttt{LatticeSystem/Quantum/SpinS/HeisenbergToyHamiltonianZeroSaturated.lean}.
$\hat H_{\rm toy_S}(A, N) = 0$ when $|\neg A| = 0$ — no cross-
sublattice bonds. Direct from
\texttt{heisenbergToyHamiltonianS\_eq\_two\_sublatticeSpinSDot} +
the three axis-vanishing lemmas of PR \#2784.

\textbf{Step 36 (PR~\#2807)}: mirror of PR \#2806 in
\texttt{LatticeSystem/Quantum/SpinS/HeisenbergToyHamiltonianZeroCardA.lean}.
$\hat H_{\rm toy_S}(A, N) = 0$ when $|A| = 0$. Apply the
axis-vanishing lemmas of PR \#2784 to $A$ directly.

\textbf{Step 37 (PR~\#2808)}: predicted GS at saturated case is
invariant under any Heisenberg Hamiltonian in
\texttt{LatticeSystem/Quantum/SpinS/BipartiteToyGSSaturatedAnyJInvariant.lean}.
For any coupling $J$, at $|\neg A| = 0$,
\texttt{Submodule.map H\_J.mulVecLin (predicted GS) $\le$
predicted GS}. Combines PR \#2802 (predicted GS $=$
$(\hat S_{\rm tot})^2$ eigenspace) with the SU(2) invariance of
any Heisenberg Hamiltonian. Generalises PR \#2794 (toy-H
invariance) to any coupling at the saturated case.

\textbf{Step 38 (PR~\#2809)}: opposite saturated mirror of
PR \#2808 in
\texttt{LatticeSystem/Quantum/SpinS/BipartiteToyGSOppositeSaturatedAnyJInvariant.lean}.
At $|A| = 0$, for any $J$, the complement-orientation predicted
GS is invariant under
\texttt{heisenbergHamiltonianS J N}. Apply PR \#2808 to $\neg A$
via filter congruence.

\textbf{Step 39 (PR~\#2810)}: max-Casimir $(\hat S_{\rm tot})^2$
eigenspace dimension in
\texttt{LatticeSystem/Quantum/SpinS/TotalSpinSSquaredMaxEigenspaceFinrank.lean}.
$\dim_\mathbb{C} \texttt{Module.End.eigenspace}((\hat S_{\rm tot})^2)
_{\rm maxCasimir} = |V|\cdot N + 1 = 2 m_{\max}+1$. Operator-level
dimension of the spin-$m_{\max}$ irreducible representation at the
top weight, A-independent and J-independent. Composition of PR
\#2801 ($=$ saturated FM joint at $J=0$) and PR \#2769 ($=$
finrank).

\textbf{Step 40 (PR~\#2811)}: max-Casimir $(\hat S_{\rm tot})^2$
eigenspace $=$ span of ladder iterates in
\texttt{LatticeSystem/Quantum/SpinS/TotalSpinSSquaredMaxEigenspaceEqSpanLadder.lean}.
\texttt{Module.End.eigenspace} $(\hat S_{\rm tot})^2_{\rm maxCasimir}
= \texttt{Submodule.span}\,\mathbb{C}\,
(\texttt{Set.range (ladderIterateUp V N)})$. Concrete basis at
the operator level, A-independent and J-independent. Composition
of PR \#2801 ($=$ saturated FM joint at $J=0$) and PR \#2768 ($=$
span ladder).

\textbf{Step 41 (PR~\#2812)}: max-Casimir $(\hat S_{\rm tot})^2$
eigenspace invariance under any $H_J$ in
\texttt{LatticeSystem/Quantum/SpinS/TotalSpinSSquaredMaxEigenspaceHeisenbergInvariant.lean}.
For any $J$, $\texttt{Submodule.map H\_J.mulVecLin}\,(\text{max-
Casimir eigenspace}) \le \text{max-Casimir eigenspace}$. Direct
from $[\hat H_J, (\hat S_{\rm tot})^2] = 0$ (SU(2) invariance).
A-independent, J-independent invariance underlying PRs \#2808 /
\#2809.

\textbf{Step 42 (PR~\#2813)}: max-Casimir $(\hat S_{\rm tot})^2$
eigenspace is closed under all $su(2)$ generators in
\texttt{LatticeSystem/Quantum/SpinS/TotalSpinSSquaredMaxEigenspaceSU2Invariant.lean}.
Each of $\hat S^{(\alpha)}_{\rm tot}$ ($\alpha = 1, 2, 3$) and
$\hat S^\pm_{\rm tot}$ preserves the max-Casimir eigenspace, via
$[\hat T, (\hat S_{\rm tot})^2] = 0$. A general helper
\texttt{totalSpinSSquaredEigenspace\_invariant\_of\_commute}
packages ``any commuting operator preserves any
$(\hat S_{\rm tot})^2$-eigenspace''.

\textbf{Step 43 (PR~\#2814)}: all-aligned states are in the
max-Casimir $(\hat S_{\rm tot})^2$ eigenspace in
\texttt{LatticeSystem/Quantum/SpinS/AllAlignedStateMemTotalSpinSSquaredMaxEigenspace.lean}.
$|\sigma_\top\rangle, |\sigma_\bot\rangle \in
\texttt{Module.End.eigenspace}\,(\hat S_{\rm tot})^2_{\rm
maxCasimir}$. Direct membership form of PRs \#878 / \#879.

\textbf{Step 44 (PR~\#2815)}: predicted GS $\subseteq$ each
individual Casimir eigenspace at the target (any $A$, any $N$) in
\texttt{LatticeSystem/Quantum/SpinS/BipartiteToyGSLeTotalSpinSSquaredEigenspace.lean}.
Three trivial inclusions packaging the joint Casimir definition's
projections as standalone $\le$ statements. No saturation
hypothesis.

\textbf{Step 45 (PR~\#2816)}: $(\hat S_{\rm tot})^2$ eigenspace at
ANY eigenvalue $\alpha$ is $su(2)$-invariant in
\texttt{LatticeSystem/Quantum/SpinS/TotalSpinSSquaredAnyEigenspaceSU2Invariant.lean}.
Public helper \texttt{totalSpinSSquaredEigenspace\_invariant\_of\_commute\_at}
plus 5 specialised invariance theorems for each $su(2)$ generator.
Generalises PR \#2813 from the max-Casimir eigenvalue to any
$\alpha$. A-, J-, $\alpha$-independent.

\textbf{Step 46 (PR~\#2817)}: strict negativity of the predicted
min energy in the non-degenerate case in
\texttt{LatticeSystem/Quantum/SpinS/BipartiteToyMinEnergyStrictNeg.lean}.
$(\texttt{bipartiteToyMinEnergyPredicted A N}).\texttt{re} < 0$
when $|A| \ge 1$, $|\neg A| \ge 1$, $N \ge 1$. Strictly tightens
PR \#2791 in the non-degenerate regime.

\textbf{Step 47 (PR~\#2818)}: max-Casimir $(\hat S_{\rm tot})^2$
eigenspace is non-trivial in
\texttt{LatticeSystem/Quantum/SpinS/TotalSpinSSquaredMaxEigenspaceNeBot.lean}.
$\texttt{Module.End.eigenspace}\,(\hat S_{\rm tot})^2_{\rm
maxCasimir} \ne \bot$. Contains $|\sigma_\top\rangle$ (PR \#2814)
which is non-zero. A-independent, J-independent existence.

\textbf{Step 48 (PR~\#2820)}: predicted GS subspace at the
saturated case is non-trivial in
\texttt{LatticeSystem/Quantum/SpinS/BipartiteToyGSSaturatedNeBot.lean}.
\texttt{bipartiteToyGroundStateSubspacePredicted A N} $\ne \bot$
when $|\neg A| = 0$, and the mirror statement when $|A| = 0$ for
the complement orientation. Direct rewrite via the saturated
collapse (PR \#2802 / PR \#2803) into PR \#2818's max-Casimir
non-emptiness. First non-emptiness statement of the predicted
$\gamma$-4 ground-state subspace.

\textbf{Step 49 (PR~\#2821)}: saturated-FM joint eigenspace at
$J = 0$ is non-trivial in
\texttt{LatticeSystem/Quantum/SpinS/SaturatedFMJointZeroNeBot.lean}.
\texttt{saturatedFerromagnetJointEigenspace} $(0 : V \to V \to
\mathbb C)$ $N \ne \bot$. Direct rewrite via PR \#2801 into PR
\#2818's max-Casimir non-emptiness. J-, A-independent existence at
the trivial-coupling endpoint.

\textbf{Step 50 (PR~\#2823)}: predicted GS subspace at either
saturated edge is non-trivial in
\texttt{LatticeSystem/Quantum/SpinS/BipartiteToyGSSaturatedOrComplementNeBot.lean}.
If $|A| = 0 \lor |\neg A| = 0$, then either
\texttt{bipartiteToyGroundStateSubspacePredicted A N} $\ne \bot$
(when $|\neg A| = 0$) or its complement orientation
\texttt{(fun x => ! A x)} version is $\ne \bot$ (when $|A| = 0$).
Disjunctive unification of PR \#2820.

\textbf{Step 51 (PR~\#2824)}: complement mirror of PR \#2782's
saturated-edge witness in
\texttt{LatticeSystem/Quantum/SpinS/AllAlignedStateMemBipartiteToyGSSaturated.lean}.
At $|A| = 0$,
\texttt{allAlignedStateS}\,$\Lambda$\,$N$\,$0 \in$
\texttt{bipartiteToyGroundStateSubspacePredicted (fun x => ! A x)
N}. Applies PR \#2782's original-orientation theorem to the
complement orientation via the $\neg\neg A = A$ filter
congruence. Pairs with PR \#2820's complement ne\_bot to give the
explicit witness $|\sigma_\top\rangle$ at the $|A| = 0$ edge.

\textbf{Step 52 (PR~\#2825)}: \texttt{bipartiteImbalanceWeight}
is real in
\texttt{LatticeSystem/Quantum/SpinS/BipartiteImbalanceWeightImZero.lean}.
Realises \texttt{bipartiteImbalanceWeight A N} as the lift of the
real-axis imbalance times $N/2$:
$\texttt{bipartiteImbalanceWeight A N} =
((|A| - |\neg A|) \cdot N / 2 : \mathbb R)$.
Corollaries: imaginary part is $0$ and real part is
$(|A| - |\neg A|) \cdot N / 2 : \mathbb R$. Cleans up
Tasaki §2.5 Theorem 2.3's real-axis comparisons by eliminating
the implicit $\mathbb R \to \mathbb C$ coercion.

\textbf{Step 53 (PR~\#2826)}: norm of
\texttt{bipartiteImbalanceWeight} in
\texttt{LatticeSystem/Quantum/SpinS/BipartiteImbalanceWeightAbs.lean}.
$\|\texttt{bipartiteImbalanceWeight A N}\| =
||A| - |\neg A|| \cdot N / 2$. This is exactly the
$||A| - |\neg A|| \cdot S$ predicted total spin appearing in
Tasaki §2.5 Theorem 2.3 (recall $S = N/2$). Direct from
PR \#2825's ofReal realisation and \texttt{Complex.norm\_real}.

\textbf{Step 54 (PR~\#2840)}: integer-imbalance form of
\texttt{bipartiteImbalanceWeight} in
\texttt{LatticeSystem/Quantum/SpinS/BipartiteImbalanceWeightInt.lean}.
\texttt{bipartiteImbalanceWeight A N} $=
(((|A| - |\neg A|) : \mathbb Z) : \mathbb C) \cdot N / 2$. Lifts
the signed sublattice imbalance through $\mathbb Z \to \mathbb C$;
useful for orientation-flip arithmetic without going through the
real-axis representation.

\textbf{Step 55 (PR~\#2841)}: orientation flip negates
\texttt{bipartiteImbalanceWeight} in
\texttt{LatticeSystem/Quantum/SpinS/BipartiteImbalanceWeightComplement.lean}.
\texttt{bipartiteImbalanceWeight (fun x => ! A x) N}
$= -$\texttt{bipartiteImbalanceWeight A N}. Algebraic statement
that the complement orientation has the opposite sign; building
block for orientation independence of the Theorem 2.3 predicted
total-spin magnitude $||A| - |\neg A|| \cdot S$.

\textbf{Step 56 (PR~\#2842)}: orientation independence of the
imbalance-weight norm in
\texttt{LatticeSystem/Quantum/SpinS/BipartiteImbalanceWeightNormComplement.lean}.
$\|\texttt{bipartiteImbalanceWeight (}\neg A\texttt{) N}\| =
\|\texttt{bipartiteImbalanceWeight A N}\|$. Direct corollary of
PR \#2841 and \texttt{norm\_neg}; confirms the orientation
independence of the Theorem 2.3 predicted total-spin magnitude
$||A| - |\neg A|| \cdot S$.

\textbf{Step 57 (PR~\#2843)}: \texttt{bipartiteToyMinEnergyPredictedSymm}
is real in
\texttt{LatticeSystem/Quantum/SpinS/BipartiteToyMinEnergySymmReal.lean}.
\texttt{im = 0} and real part equals $-|A| \cdot |\neg A| \cdot N^2 / 2
- \min(|A|, |\neg A|) \cdot N$ on $\mathbb R$ (with \texttt{min}
taken in $\mathbb R$ after \texttt{Nat.cast\_min}). Mirrors PR
\#2825 for the asymmetric form.

\textbf{Step 58 (PR~\#2844)}: non-positivity of the predicted
symmetric min energy in
\texttt{LatticeSystem/Quantum/SpinS/BipartiteToyMinEnergySymmNonpos.lean}.
$(\texttt{bipartiteToyMinEnergyPredictedSymm A N}).\texttt{re} \le 0$.
Both contributions $-|A| \cdot |\neg A| \cdot N^2 / 2$ and
$-\min(|A|, |\neg A|) \cdot N$ are non-positive. Symmetric-form
mirror of PR \#2791.

\textbf{Step 59 (PR~\#2845)}: strict negativity of the predicted
symmetric min energy (non-degenerate case) in
\texttt{LatticeSystem/Quantum/SpinS/BipartiteToyMinEnergySymmStrictNeg.lean}.
$(\texttt{bipartiteToyMinEnergyPredictedSymm A N}).\texttt{re} < 0$
when $|A| \ge 1$, $|\neg A| \ge 1$, $N \ge 1$. Symmetric-form
mirror of PR \#2817.

\textbf{Step 60 (PR~\#2846)}: explicit witness of the saturated-FM
joint eigenspace at $J = 0$ in
\texttt{LatticeSystem/Quantum/SpinS/AllAlignedStateMemSaturatedFMJointZero.lean}.
\texttt{allAlignedStateS V N 0} $\in$
\texttt{saturatedFerromagnetJointEigenspace} $(0 : V \to V \to
\mathbb C)$ $N$. Strengthens PR \#2821's ne\_bot into a concrete
witness via PR \#2801 and PR \#2814.

\textbf{Step 61 (PR~\#2847)}: all-down mirror of PR \#2846 in
\texttt{LatticeSystem/Quantum/SpinS/AllAlignedStateLastMemSaturatedFMJointZero.lean}.
\texttt{allAlignedStateS V N (Fin.last N)} $\in$
\texttt{saturatedFerromagnetJointEigenspace} $(0 : V \to V \to
\mathbb C)$ $N$. Same chain (PR \#2801 + PR \#2814 all-down
branch).

\textbf{Step 62 (PR~\#2848)}: norm of the predicted symmetric
min energy in
\texttt{LatticeSystem/Quantum/SpinS/BipartiteToyMinEnergySymmNorm.lean}.
$\|\texttt{bipartiteToyMinEnergyPredictedSymm A N}\| = |A| \cdot
|\neg A| \cdot N^2 / 2 + \min(|A|, |\neg A|) \cdot N$. Combines PR
\#2843 (\texttt{im = 0} + \texttt{re} explicit) with PR \#2844
(\texttt{re} non-positive): norm $= |\texttt{re}| = -\texttt{re}$
for a non-positive real, then expand.

\textbf{Step 63 (PR~\#2849)}: symmetric predicted min energy
vanishes at saturated edges in
\texttt{LatticeSystem/Quantum/SpinS/BipartiteToyMinEnergySymmSaturated.lean}.
At $|\neg A| = 0$ or $|A| = 0$,
\texttt{bipartiteToyMinEnergyPredictedSymm A N} $= 0$. Symmetric-
form mirror of PR \#2779.

\textbf{Step 64 (PR~\#2850)}: imbalance-weight norm at saturated
edges equals $|\Lambda| \cdot N / 2$ in
\texttt{LatticeSystem/Quantum/SpinS/BipartiteImbalanceWeightNormSaturated.lean}.
At $|\neg A| = 0$ or $|A| = 0$,
$\|\texttt{bipartiteImbalanceWeight A N}\| = |\Lambda| \cdot N / 2
= m_{\rm max}$. Matches the maximum total spin attained at the
all-aligned ferromagnetic state.

\textbf{Step 65 (PR~\#2851)}: symm-energy norm vanishes at
saturated edges in
\texttt{LatticeSystem/Quantum/SpinS/BipartiteToyMinEnergySymmNormSaturated.lean}.
At $|\neg A| = 0$ or $|A| = 0$,
$\|\texttt{bipartiteToyMinEnergyPredictedSymm A N}\| = 0$.
Direct from PR \#2849 and \texttt{norm\_zero}.

\textbf{Step 66 (PR~\#2852)}: balanced-case symm-predicted min
energy in
\texttt{LatticeSystem/Quantum/SpinS/BipartiteToyMinEnergySymmBalanced.lean}.
When $|A| = |\neg A|$,
\texttt{bipartiteToyMinEnergyPredictedSymm A N} $= -|A|^2 \cdot
N^2 / 2 - |A| \cdot N$. Cross term collapses to $|A|^2$ and
$\min(|A|, |\neg A|) = |A|$. Symmetric-form mirror of the
asymmetric balanced form.

\textbf{Step 67 (PR~\#2853)}: balanced-case symm-energy norm in
\texttt{LatticeSystem/Quantum/SpinS/BipartiteToyMinEnergySymmNormBalanced.lean}.
When $|A| = |\neg A|$,
$\|\texttt{bipartiteToyMinEnergyPredictedSymm A N}\| = |A|^2 \cdot
N^2 / 2 + |A| \cdot N$. Direct from PR \#2848 balanced
specialisation.

\textbf{Step 68 (PR~\#2854)}: characterisation of
\texttt{bipartiteImbalanceWeight} $= 0$ in
\texttt{LatticeSystem/Quantum/SpinS/BipartiteImbalanceWeightEqZeroIff.lean}.
At $N \ge 1$, \texttt{bipartiteImbalanceWeight A N} $= 0
\Leftrightarrow |A| = |\neg A|$. The forward direction is PR
\#2773; the reverse uses the real-part formula and the strict
positivity $N / 2 > 0$ at $N \ge 1$.

\textbf{Step 69 (PR~\#2855)}: norm-zero characterisation in
\texttt{LatticeSystem/Quantum/SpinS/BipartiteImbalanceWeightNormEqZero.lean}.
At $N \ge 1$, $\|\texttt{bipartiteImbalanceWeight A N}\| = 0
\Leftrightarrow |A| = |\neg A|$. Direct from \texttt{norm\_eq\_zero}
and PR \#2854.

\textbf{Step 70 (PR~\#2856)}: strict sign of the imbalance-weight
real part in
\texttt{LatticeSystem/Quantum/SpinS/BipartiteImbalanceWeightReStrictSign.lean}.
At $N \ge 1$: $|A| > |\neg A|$ implies $(\texttt{bipartiteImbalanceWeight
A N}).\texttt{re} > 0$; $|A| < |\neg A|$ implies $< 0$. Strict-
inequality variants of PR \#2773.

\textbf{Step 71 (PR~\#2857)}: Néel-state magnetization norm =
Theorem 2.3 predicted spin in
\texttt{LatticeSystem/Quantum/SpinS/NeelStateMagNormEqPredicted.lean}.
$\|\texttt{magEigenvalueS (neelConfigOfS A N)}\| = ||A| - |\neg A||
\cdot N / 2 = ||A| - |\neg A|| \cdot S$. Quantitative connection
from the Néel variational input to Theorem 2.3's predicted total
spin magnitude.

\textbf{Step 72 (PR~\#2858)}: complement Néel-state magnetization
norm = same predicted spin in
\texttt{LatticeSystem/Quantum/SpinS/NeelStateComplementMagNormEqPredicted.lean}.
$\|\texttt{magEigenvalueS (neelConfigOfS (}\neg A\texttt{) N)}\| =
||A| - |\neg A|| \cdot N / 2$. Complement orientation has the
negated magnetization eigenvalue, so its norm is the same.

\textbf{Step 73 (PR~\#2859)}: Néel-state mag norm at saturated
edges $= m_{\rm max}$ in
\texttt{LatticeSystem/Quantum/SpinS/NeelStateMagNormSaturated.lean}.
At $|\neg A| = 0$ or $|A| = 0$,
$\|\texttt{magEigenvalueS (neelConfigOfS A N)}\| = |\Lambda|
\cdot N / 2 = m_{\rm max}$. Saturated-limit specialisation of
PR \#2857.

\textbf{Step 74 (PR~\#2860)}: strict positivity of the symm-energy
norm (non-degenerate case) in
\texttt{LatticeSystem/Quantum/SpinS/BipartiteToyMinEnergySymmNormPos.lean}.
$0 < \|\texttt{bipartiteToyMinEnergyPredictedSymm A N}\|$ when
$|A| \ge 1$, $|\neg A| \ge 1$, $N \ge 1$. Direct from PR \#2845's
strict negativity of the real part.

\textbf{Step 75 (PR~\#2861)}: strict positivity of
$\|\texttt{bipartiteImbalanceWeight A N}\|$ when $|A| \ne |\neg A|$
and $N \ge 1$ in
\texttt{LatticeSystem/Quantum/SpinS/BipartiteImbalanceWeightNormPos.lean}.
Contrapositive of PR \#2855.

\textbf{Step 76 (PR~\#2862)}: norm = real part when $|A| \ge |\neg A|$
in
\texttt{LatticeSystem/Quantum/SpinS/BipartiteImbalanceWeightNormEqRe.lean}.
$\|\texttt{bipartiteImbalanceWeight A N}\| =
(\texttt{bipartiteImbalanceWeight A N}).\texttt{re}$. The value
is real (PR \#2825) and non-negative when $|A| \ge |\neg A|$ (PR
\#2773), so $\texttt{norm} = \texttt{.re}$ directly.

\textbf{Step 77 (PR~\#2863)}: norm $= -$ real part when
$|A| \le |\neg A|$ in
\texttt{LatticeSystem/Quantum/SpinS/BipartiteImbalanceWeightNormEqNegRe.lean}.
$\|\texttt{bipartiteImbalanceWeight A N}\| =
-(\texttt{bipartiteImbalanceWeight A N}).\texttt{re}$.
Opposite-orientation mirror of PR \#2862 via \texttt{abs\_of\_nonpos}.

\textbf{Step 78 (PR~\#2864)}: iff characterisation of vanishing
symm-predicted min energy at $N \ge 1$ in
\texttt{LatticeSystem/Quantum/SpinS/BipartiteToyMinEnergySymmEqZeroIff.lean}.
$\texttt{bipartiteToyMinEnergyPredictedSymm A N} = 0
\Leftrightarrow |A| = 0 \lor |\neg A| = 0$. Reverse direction
(saturated edge $\to$ energy 0) from PR \#2849; forward direction
(energy 0 $\to$ saturated edge) via the contrapositive of PR
\#2845's strict negativity.

\textbf{Step 79 (PR~\#2865)}: norm-zero characterisation at
$N \ge 1$ in
\texttt{LatticeSystem/Quantum/SpinS/BipartiteToyMinEnergySymmNormEqZeroIff.lean}.
$\|\texttt{bipartiteToyMinEnergyPredictedSymm A N}\| = 0
\Leftrightarrow |A| = 0 \lor |\neg A| = 0$. Direct from
\texttt{norm\_eq\_zero} and PR \#2864.

\textbf{Step 80 (PR~\#2866)}: alternative form of symm-energy
norm positivity in
\texttt{LatticeSystem/Quantum/SpinS/BipartiteToyMinEnergySymmNormPosOfCardNeZero.lean}.
$0 < \|\texttt{bipartiteToyMinEnergyPredictedSymm A N}\|$ when
$N \ge 1$, $|A| \ne 0$, $|\neg A| \ne 0$. Contrapositive of
PR \#2865.

\textbf{Step 81 (PR~\#2867)}: positive forms of the imbalance-
weight norm in
\texttt{LatticeSystem/Quantum/SpinS/BipartiteImbalanceWeightNormPositive.lean}.
At $|\neg A| \le |A|$:
$\|\texttt{bipartiteImbalanceWeight A N}\| = (|A| - |\neg A|)
\cdot N / 2$. At $|A| \le |\neg A|$:
$\|\texttt{bipartiteImbalanceWeight A N}\| = (|\neg A| - |A|)
\cdot N / 2$. Drops the absolute-value bars case-specifically.

\textbf{Step 82 (PR~\#2868)}: case-specific symm-energy norm in
\texttt{LatticeSystem/Quantum/SpinS/BipartiteToyMinEnergySymmNormCardA.lean}.
At $|A| \le |\neg A|$: $\|\texttt{symm}\| = |A| \cdot |\neg A|
\cdot N^2 / 2 + |A| \cdot N$. At $|\neg A| \le |A|$:
$\|\texttt{symm}\| = |A| \cdot |\neg A| \cdot N^2 / 2 + |\neg A|
\cdot N$. Drops the $\min$ from PR \#2848.

\textbf{Step 83 (PR~\#2869)}: negated-form of the symm-energy real
part in
\texttt{LatticeSystem/Quantum/SpinS/BipartiteToyMinEnergySymmReNonpos.lean}.
$(\texttt{symm}).\texttt{re} = -(|A| \cdot |\neg A| \cdot N^2 / 2
+ \min(|A|, |\neg A|) \cdot N)$. Sign-normalised expression of
PR \#2843's \texttt{\_re\_eq}.

\textbf{Step 84 (PR~\#2870)}: bipartition cardinality sum in
\texttt{LatticeSystem/Quantum/SpinS/CardSumEqFintype.lean}.
$|A| + |\neg A| = \texttt{Fintype.card}\,\Lambda$ via
\texttt{Finset.card\_filter\_add\_card\_filter\_not}. Foundational
identity used across the $\gamma$-4 chain.

\textbf{Step 85 (PR~\#2871)}: saturated-edge cardinality
identification in
\texttt{LatticeSystem/Quantum/SpinS/CardSaturatedEdge.lean}.
$|\neg A| = 0 \to |A| = |\Lambda|$ and $|A| = 0 \to |\neg A| =
|\Lambda|$. Direct corollaries of PR \#2870 via \texttt{omega}.

\textbf{Step 86 (PR~\#2872)}: both saturated edges $\Leftrightarrow$
empty lattice in
\texttt{LatticeSystem/Quantum/SpinS/CardBothZeroIffEmpty.lean}.
$(|A| = 0 \land |\neg A| = 0) \Leftrightarrow
\texttt{Fintype.card}\,\Lambda = 0$. Direct corollary of PR \#2870.

\textbf{Step 87 (PR~\#2873)}: at least one sublattice non-empty in
\texttt{LatticeSystem/Quantum/SpinS/CardPosOfNonempty.lean}.
$[\texttt{Nonempty}\,\Lambda] \to 0 < |A| \lor 0 < |\neg A|$.
Direct corollary of PR \#2872.

\textbf{Step 88 (PR~\#2874)}: upper bound by $m_{\rm max}$ in
\texttt{LatticeSystem/Quantum/SpinS/BipartiteImbalanceWeightNormLeMMax.lean}.
$\|\texttt{bipartiteImbalanceWeight A N}\| \le |\Lambda| \cdot N / 2
= m_{\rm max}$. The predicted Theorem 2.3 spin magnitude is
bounded above by the maximum total spin. Direct from
$||A| - |\neg A|| \le |A| + |\neg A| = |\Lambda|$ (PR \#2870).

\textbf{Step 89 (PR~\#2875)}: symm-energy norm upper bound in
\texttt{LatticeSystem/Quantum/SpinS/BipartiteToyMinEnergySymmNormLeBound.lean}.
$\|\texttt{bipartiteToyMinEnergyPredictedSymm A N}\| \le
|\Lambda|^2 \cdot N^2 / 8 + |\Lambda| \cdot N / 2$. Via AM-GM
$|A| \cdot |\neg A| \le (|\Lambda|/2)^2$ and $\min(|A|, |\neg A|)
\le |\Lambda|/2$.

\textbf{Step 90 (PR~\#2876)}: bridge identity in
\texttt{LatticeSystem/Quantum/SpinS/BipartiteToyMinEnergySymmViaImbalance.lean}.
$8 \cdot ((\texttt{bipartiteToyMinEnergyPredictedSymm A N}).\mathrm{re}
   + \min(|A|, |\neg A|) \cdot N)
   + |\Lambda|^2 \cdot N^2
   = 4 \cdot ((\texttt{bipartiteImbalanceWeight A N}).\mathrm{re})^2$.
Algebraic consequence of $4 \cdot |A| \cdot |\neg A| =
(|A|+|\neg A|)^2 - (|A|-|\neg A|)^2$ combined with the
cardinality identity $|A| + |\neg A| = |\Lambda|$ (PR \#2870).
The predicted symmetric ground-state energy and the predicted
spin magnitude $||A| - |\neg A|| \cdot N / 2$ are linked by a
simple quadratic constraint involving only the lattice
cardinality.

\textbf{Step 91 (PR~\#2877)}: norm-squared = real-part squared
helper
(\texttt{bipartiteImbalanceWeight\_norm\_mul\_self\_eq\_re\_mul\_self}:
$\|\mathrm{biw}\|\cdot\|\mathrm{biw}\| =
\mathrm{biw}.\mathrm{re}\cdot\mathrm{biw}.\mathrm{re}$, since
$\mathrm{biw}$ is real, PR \#2825) and the resulting bridge
identity in $\|\mathrm{biw}\|^2$ form in
\texttt{LatticeSystem/Quantum/SpinS/BipartiteToyMinEnergySymmReViaImbalanceNormSq.lean}.
$8 \cdot ((\texttt{bipartiteToyMinEnergyPredictedSymm A N}).\mathrm{re}
   + \min(|A|, |\neg A|) \cdot N)
   + |\Lambda|^2 \cdot N^2
   = 4 \cdot \|\texttt{bipartiteImbalanceWeight A N}\|^2$.
Substitutes $(\mathrm{biw}.\mathrm{re})^2 = \|\mathrm{biw}\|^2$
into Step 90; the natural form for Tasaki §2.5 Theorem 2.3 where
the predicted spin magnitude
$\|\mathrm{biw}\| = ||A| - |\neg A||\cdot N/2$ appears directly.

\textbf{Step 92 (PR~\#2878)}: explicit closed form of the
symm-energy real part in
\texttt{LatticeSystem/Quantum/SpinS/BipartiteToyMinEnergySymmReEqHalfImbalanceNormSq.lean}.
$(\texttt{bipartiteToyMinEnergyPredictedSymm A N}).\mathrm{re}
   = \|\texttt{bipartiteImbalanceWeight A N}\|^2 / 2
   - |\Lambda|^2 \cdot N^2 / 8
   - \min(|A|, |\neg A|) \cdot N$.
Direct \texttt{linarith} rearrangement of Step 91. Expresses the
predicted symmetric ground-state energy as a function of the spin
magnitude squared, the lattice cardinality, and the minor
sublattice count.

\textbf{Step 93 (PR~\#2880)}: lower bound on the symm-energy real
part in
\texttt{LatticeSystem/Quantum/SpinS/BipartiteToyMinEnergySymmReGeNegBalanced.lean}.
$-|\Lambda|^2\cdot N^2/8 - \min(|A|, |\neg A|)\cdot N \le
   (\texttt{bipartiteToyMinEnergyPredictedSymm A N}).\mathrm{re}$.
Direct from Step 92 by dropping the non-negative
$\|\mathrm{biw}\|^2/2$ term. Combined with the existing upper bound
$(\mathrm{symm}).\mathrm{re} \le 0$, gives the sandwich
$-|\Lambda|^2\cdot N^2/8 - \min\cdot N \le \mathrm{symm.re} \le 0$;
equality at the lower bound iff $\|\mathrm{biw}\| = 0$
(for $N \ge 1$ this characterises balanced sublattices).

\textbf{Step 94 (PR~\#2881)}: balanced-case symm-energy real part in
$|\Lambda|$ form in
\texttt{LatticeSystem/Quantum/SpinS/BipartiteToyMinEnergySymmBalancedReEqCardinality.lean}.
When $|A| = |\neg A|$,
$(\texttt{bipartiteToyMinEnergyPredictedSymm A N}).\mathrm{re}
   = -|\Lambda|^2\cdot N^2/8 - |\Lambda|\cdot N/2$.
Combines the existing balanced complex form
($= -|A|^2\cdot N^2/2 - |A|\cdot N$, PR \#2842) with the cardinality
identity $2\cdot|A| = |\Lambda|$ (PR \#2870). Saturates the
Step 93 lower bound at balanced sublattices.

\textbf{Step 95 (PR~\#2882)}: configuration-independent global lower
bound on the symm-energy real part in
\texttt{LatticeSystem/Quantum/SpinS/BipartiteToyMinEnergySymmReGeGlobalMin.lean}.
For any $A$:
$-|\Lambda|^2\cdot N^2/8 - |\Lambda|\cdot N/2 \le
   (\texttt{bipartiteToyMinEnergyPredictedSymm A N}).\mathrm{re}$.
Combines Step 93 ($A$-dependent lower bound
$-|\Lambda|^2\cdot N^2/8 - \min(|A|, |\neg A|)\cdot N$) with the
elementary inequality $\min \le (|A|+|\neg A|)/2 = |\Lambda|/2$
(from PR \#2870). Saturated at balanced sublattices (Step 94):
balanced configurations realise the global minimum of the
symm-predicted energy real part.

\textbf{Step 96 (PR~\#2883)}: refined upper bound on the symm-energy
real part in
\texttt{LatticeSystem/Quantum/SpinS/BipartiteToyMinEnergySymmReLeNegMin.lean}.
$(\texttt{bipartiteToyMinEnergyPredictedSymm A N}).\mathrm{re}
   \le -\min(|A|, |\neg A|)\cdot N$.
Direct from Step 92 closed form together with the squared form of
the imbalance-norm $m_{\rm max}$ upper bound
$\|\mathrm{biw}\| \le |\Lambda|\cdot N/2$ (giving
$\|\mathrm{biw}\|^2/2 \le |\Lambda|^2\cdot N^2/8$, hence the first
two terms of the closed form contribute non-positively).
Refines the existing $\le 0$ bound when $\min > 0$; combined with
Step 93 gives the tight $\min$-aware sandwich
$-|\Lambda|^2\cdot N^2/8 - \min\cdot N \le \mathrm{symm.re} \le
  -\min\cdot N$.

\textbf{Step 97 (PR~\#2884)}: Néel-to-predicted-min energy gap in
\texttt{LatticeSystem/Quantum/SpinS/NeelStateOfSPredictedSymmGap.lean}.
$\langle\Phi_{\rm N\acute{e}el}(A,N) |\, \widehat H_{\rm toy}^S\,
   | \Phi_{\rm N\acute{e}el}(A,N)\rangle
   - \texttt{bipartiteToyMinEnergyPredictedSymm A N}
   = \min(|A|, |\neg A|)\cdot N$.
The N\'eel state's energy expectation is above the symmetric
predicted minimum by exactly $\min(|A|,|\neg A|)\cdot N$ (a
non-negative gap that vanishes at saturated edges and is maximised
at balanced sublattices). Consistent with Tasaki §2.5 Theorem 2.3:
the N\'eel state is a candidate ground state but the true GS energy
is lower by the sublattice-Casimir diagonal contribution.

\textbf{Step 98 (PR~\#2885)}: strict refined upper bound at
non-degenerate configurations in
\texttt{LatticeSystem/Quantum/SpinS/BipartiteToyMinEnergySymmReLtNegMinStrict.lean}.
When $|A| \ge 1$, $|\neg A| \ge 1$, $N \ge 1$:
$(\texttt{bipartiteToyMinEnergyPredictedSymm A N}).\mathrm{re}
   < -\min(|A|, |\neg A|)\cdot N$.
Strict version of Step 96; immediate from Step 92's
$\mathrm{re} = -|A|\cdot|\neg A|\cdot N^2/2 - \min\cdot N$ since
$|A|\cdot|\neg A|\cdot N^2 > 0$ under non-degeneracy.

\textbf{Step 99 (PR~\#2886)}: Néel-state expectation real part in
$\|\mathrm{biw}\|^2$ form in
\texttt{LatticeSystem/Quantum/SpinS/NeelStateOfSExpectationReViaImbalanceNormSq.lean}.
$(\langle\Phi_{\rm N\acute{e}el} |\, \widehat H_{\rm toy}^S\,
   | \Phi_{\rm N\acute{e}el}\rangle).\mathrm{re}
   = \|\texttt{bipartiteImbalanceWeight A N}\|^2/2
   - |\Lambda|^2\cdot N^2/8$.
Combines Step 97 (Néel--predicted-min gap $= \min\cdot N$) with
Step 92 (predicted-symm closed form): the $\min\cdot N$ terms cancel,
leaving a clean form depending only on the predicted spin magnitude
squared and the lattice cardinality.

\textbf{Step 100 (PR~\#2887)}: linear identity bridging $\min\cdot N$
and $\|\mathrm{biw}\|$ in
\texttt{LatticeSystem/Quantum/SpinS/MinNEqHalfCardTimesNMinusImbalanceNorm.lean}.
$\min(|A|, |\neg A|)\cdot N = |\Lambda|\cdot N/2 -
   \|\texttt{bipartiteImbalanceWeight A N}\|$.
Direct from $\|\mathrm{biw}\| = ||A|{-}|\neg A||\cdot N/2$
(PR \#2826) and $|A|+|\neg A| = |\Lambda|$ (PR \#2870), via the
elementary identity $2\cdot\min = (|A|+|\neg A|) - ||A|{-}|\neg A||$.
Enables $\min\cdot N$-free reformulations of all Tasaki §2.5
Theorem 2.3 ($\gamma$-4) energy bounds.

\textbf{Step 101 (PR~\#2888)}: univariate $\|\mathrm{biw}\|$-parameterisation
of the symm-energy real part in
\texttt{LatticeSystem/Quantum/SpinS/BipartiteToyMinEnergySymmReViaImbalanceNormOnly.lean}.
$(\texttt{bipartiteToyMinEnergyPredictedSymm A N}).\mathrm{re}
   = \|\mathrm{biw}\|^2/2 + \|\mathrm{biw}\|
   - |\Lambda|^2\cdot N^2/8 - |\Lambda|\cdot N/2$.
Direct from Step 92 (closed form with $\min\cdot N$) + Step 100
($\min\cdot N$ identity): eliminates $\min$-dependence,
parameterising the predicted GS energy as a univariate quadratic
in the predicted spin magnitude $\|\mathrm{biw}\|$.

\textbf{Step 102 (PR~\#2889)}: strict $m_{\rm max}$ upper bound on
imbalance norm at non-degenerate configurations in
\texttt{LatticeSystem/Quantum/SpinS/BipartiteImbalanceWeightNormLtMMaxNondegenerate.lean}.
When $|A| \ge 1$, $|\neg A| \ge 1$, $N \ge 1$:
$\|\texttt{bipartiteImbalanceWeight A N}\| < |\Lambda|\cdot N/2$.
Strict version of the $m_{\rm max}$ bound; via Step 100's identity
$\min\cdot N = |\Lambda|\cdot N/2 - \|\mathrm{biw}\|$, strictness
reduces to $\min\cdot N > 0$ at non-degeneracy.

\textbf{Step 103 (PR~\#2890)}: monotonicity of $\mathrm{symm.re}$ in
$\|\mathrm{biw}\|$ in
\texttt{LatticeSystem/Quantum/SpinS/BipartiteToyMinEnergySymmReMonoInImbalanceNorm.lean}.
For $A_1, A_2 : \Lambda \to \mathrm{Bool}$,
$\|\mathrm{biw}(A_1, N)\| \le \|\mathrm{biw}(A_2, N)\|
   \implies (\mathrm{symm}(A_1, N)).\mathrm{re} \le
            (\mathrm{symm}(A_2, N)).\mathrm{re}$.
Direct from Step 101 univariate form + monotonicity of
$x \mapsto x^2/2 + x$ on $[0, \infty)$. Characterises balanced as
the global energy minimiser and saturated as the maximiser.

\textbf{Step 104 (PR~\#2891)}: strict monotonicity of
$\mathrm{symm.re}$ in $\|\mathrm{biw}\|$ in
\texttt{LatticeSystem/Quantum/SpinS/BipartiteToyMinEnergySymmReStrictMonoInImbalanceNorm.lean}.
$\|\mathrm{biw}(A_1, N)\| < \|\mathrm{biw}(A_2, N)\|
   \implies (\mathrm{symm}(A_1, N)).\mathrm{re} <
            (\mathrm{symm}(A_2, N)).\mathrm{re}$.
Strict version of Step 103.

\textbf{Step 105 (PR~\#2892)}: AFM cross-term bridge in
\texttt{LatticeSystem/Quantum/SpinS/CardATimesCardNotATimesNSqViaImbalanceNormSq.lean}.
$|A|\cdot|\neg A|\cdot N^2 = |\Lambda|^2\cdot N^2/4
   - \|\texttt{bipartiteImbalanceWeight A N}\|^2$.
Direct from Step 90 (bridge $8\cdot(\mathrm{symm.re}+\min\cdot N)
+ |\Lambda|^2 N^2 = 4\cdot(\mathrm{biw.re})^2$) + Step 91
($\|\mathrm{biw}\|^2 = (\mathrm{biw.re})^2$) + the closed form
$\mathrm{symm.re} = -|A|\cdot|\neg A|\cdot N^2/2 - \min\cdot N$.
Expresses the Néel expectation magnitude in imbalance-norm form.

\textbf{Step 106 (PR~\#2893)}: strict positive Néel-vs-predicted-min
gap at non-degenerate in
\texttt{LatticeSystem/Quantum/SpinS/NeelStateOfSPredictedSymmGapPosNondegenerate.lean}.
When $|A| \ge 1$, $|\neg A| \ge 1$, $N \ge 1$:
$0 < (\langle\Phi_{\rm N\acute{e}el} |\,\widehat H_{\rm toy}^S\,
   | \Phi_{\rm N\acute{e}el}\rangle
   - \texttt{bipartiteToyMinEnergyPredictedSymm A N}).\mathrm{re}$.
Strict version of Step 97 (gap $= \min\cdot N$) under non-degeneracy.

\textbf{Step 107 (PR~\#2894)}: real-part iff form of
$\mathrm{symm.re} = 0$ at $N \ge 1$ in
\texttt{LatticeSystem/Quantum/SpinS/BipartiteToyMinEnergySymmReEqZeroIff.lean}.
$(\mathrm{symm}(A, N)).\mathrm{re} = 0 \iff |A| = 0 \vee |\neg A| = 0$.
Real-part mirror of the complex iff form. Forward direction is the
contrapositive of strict negativity at non-degenerate; reverse is
the saturated-edge zero.

\textbf{Step 108 (PR~\#2895)}: strict global lower bound at
unbalanced configurations in
\texttt{LatticeSystem/Quantum/SpinS/BipartiteToyMinEnergySymmReGtGlobalMinOfUnbalanced.lean}.
$-|\Lambda|^2\cdot N^2/8 - |\Lambda|\cdot N/2 <
   (\mathrm{symm}(A, N)).\mathrm{re}$ when $\|\mathrm{biw}\| > 0$.
Direct from Step 101 univariate form: $\|\mathrm{biw}\|^2/2
+ \|\mathrm{biw}\| > 0$ at $\|\mathrm{biw}\| > 0$. Strict version
of Step 95 (the global lower bound).

\textbf{Step 109 (PR~\#2896)}: Néel-state $(\widehat S_{\rm tot})^2$
expectation real part in $\|\mathrm{biw}\|^2$ form in
\texttt{LatticeSystem/Quantum/SpinS/NeelStateOfSTotalSpinSSquaredReViaImbalanceNormSq.lean}.
$(\langle\Phi_{\rm N\acute{e}el} |\, (\widehat S_{\rm tot})^2\,
   | \Phi_{\rm N\acute{e}el}\rangle).\mathrm{re}
   = \|\mathrm{biw}\|^2 + |\Lambda|\cdot N/2$.
Connects the N\'eel-state Casimir expectation to the predicted
spin magnitude squared $\|\mathrm{biw}\|^2 = (\mathrm{predicted}\,
S_{\rm tot})^2$ plus the lattice-size offset
$m_{\rm max} = |\Lambda|\cdot N/2$.

\textbf{Step 110 (PR~\#2897)}: balanced-case N\'eel
$(\widehat S_{\rm tot})^2$ expectation in
\texttt{LatticeSystem/Quantum/SpinS/NeelStateOfSTotalSpinSSquaredBalanced.lean}.
At $|A| = |\neg A|$ and $N \ge 1$:
$(\langle\Phi_{\rm N\acute{e}el} |\, (\widehat S_{\rm tot})^2\,
   | \Phi_{\rm N\acute{e}el}\rangle).\mathrm{re}
   = |\Lambda|\cdot N/2$.
Direct from Step 109 + $\|\mathrm{biw}\| = 0$ at balanced
(PR \#2854). Minimum value of the $(\widehat S_{\rm tot})^2$
expectation across all bipartitions.

\textbf{Step 111 (PR~\#2898)}: saturated-case N\'eel
$(\widehat S_{\rm tot})^2$ expectation in
\texttt{LatticeSystem/Quantum/SpinS/NeelStateOfSTotalSpinSSquaredSaturated.lean}.
At $|\neg A| = 0$ (or $|A| = 0$):
$(\langle\Phi_{\rm N\acute{e}el} |\, (\widehat S_{\rm tot})^2\,
   | \Phi_{\rm N\acute{e}el}\rangle).\mathrm{re}
   = (|\Lambda|\cdot N/2)\cdot(|\Lambda|\cdot N/2 + 1)
   = m_{\rm max}\cdot(m_{\rm max} + 1)$.
Direct from Step 109 + $\|\mathrm{biw}\| = |\Lambda|\cdot N/2$ at
saturated edges (PR \#2851/2852). The N\'eel state at saturated
edges is a $(\widehat S_{\rm tot})^2$ eigenvector at the
fully-polarised ferromagnetic maximum-spin eigenvalue.

\textbf{Step 112 (PR~\#2899)}: quantitative N\'eel-vs-predictedSymm
gap lower bound at non-degenerate in
\texttt{LatticeSystem/Quantum/SpinS/NeelStateOfSPredictedSymmGapGeN.lean}.
At $|A| \ge 1$, $|\neg A| \ge 1$:
$N \le (\langle\Phi_{\rm N\acute{e}el} |\, \widehat H_{\rm toy}^S\,
   | \Phi_{\rm N\acute{e}el}\rangle
   - \texttt{bipartiteToyMinEnergyPredictedSymm A N}).\mathrm{re}$.
Refines Step 106 strict positivity: since $\min(|A|, |\neg A|) \ge 1$
at non-degenerate, gap $= \min\cdot N \ge N$.

\textbf{Step 113 (PR~\#2901)}: N\'eel-vs-predicted $(\widehat S_{\rm tot})^2$
gap $= \min\cdot N$ in
\texttt{LatticeSystem/Quantum/SpinS/NeelStateOfSTotalSpinSSquaredVsPredictedGap.lean}.
$(\langle\Phi_{\rm N\acute{e}el} |\,(\widehat S_{\rm tot})^2\,
   | \Phi_{\rm N\acute{e}el}\rangle).\mathrm{re}
   - (\|\mathrm{biw}\|^2 + \|\mathrm{biw}\|)
   = \min(|A|,|\neg A|)\cdot N$.
Direct parallel of energy gap (Step 97) to the $(\widehat S_{\rm tot})^2$
Casimir: the N\'eel state's expectation is exactly $\min\cdot N$ above the
predicted Tasaki $S_{\rm tot}\cdot(S_{\rm tot}+1) =
\|\mathrm{biw}\|\cdot(\|\mathrm{biw}\|+1)$ eigenvalue. Gap vanishes at
saturated edges.

\textbf{Step 114 (PR~\#2902)}: strict positive N\'eel-vs-predicted
$(\widehat S_{\rm tot})^2$ gap at non-degenerate in
\texttt{LatticeSystem/Quantum/SpinS/NeelStateOfSTotalSpinSSquaredVsPredictedGapPosNondegenerate.lean}.
At $|A| \ge 1$, $|\neg A| \ge 1$, $N \ge 1$:
$\|\mathrm{biw}\|^2 + \|\mathrm{biw}\| <
   (\langle\Phi_{\rm N\acute{e}el} |\,(\widehat S_{\rm tot})^2\,
   | \Phi_{\rm N\acute{e}el}\rangle).\mathrm{re}$.
Strict version of Step 113; demonstrates the N\'eel state spans multiple
$S_{\rm tot}$ sectors (not the GS) at non-degenerate.

\textbf{Step 115 (PR~\#2903)}: quantitative N\'eel-vs-predicted
$(\widehat S_{\rm tot})^2$ gap lower bound at non-degenerate in
\texttt{LatticeSystem/Quantum/SpinS/NeelStateOfSTotalSpinSSquaredVsPredictedGapGeN.lean}.
At $|A| \ge 1$, $|\neg A| \ge 1$:
$N \le (\langle\Phi_{\rm N\acute{e}el} |\,(\widehat S_{\rm tot})^2\,
   | \Phi_{\rm N\acute{e}el}\rangle).\mathrm{re}
   - (\|\mathrm{biw}\|^2 + \|\mathrm{biw}\|)$.
Parallel of Step 112 (energy quantitative gap) for the
$(\widehat S_{\rm tot})^2$ Casimir: $\min \ge 1$ at non-degenerate,
gap $= \min\cdot N \ge N$.

\textbf{Step 116 (PR~\#2904)}: N\'eel-state $(\widehat S_{\rm tot}^{(3)})^2$
expectation real part in $\|\mathrm{biw}\|^2$ form in
\texttt{LatticeSystem/Quantum/SpinS/NeelStateOfSTotalSpinSOp3SqReViaImbalanceNormSq.lean}.
$(\langle\Phi_{\rm N\acute{e}el} |\,(\widehat S_{\rm tot}^{(3)})^2\,
   | \Phi_{\rm N\acute{e}el}\rangle).\mathrm{re} = \|\mathrm{biw}\|^2$.
Direct from $\gamma$-4 step 128 ($= ((|A|-|\neg A|)\cdot N/2)^2$) +
Step 91 ($\|\mathrm{biw}\|^2 = (\mathrm{biw.re})^2$). Identifies the
z-axis Casimir expectation with the predicted $S_{\rm tot}^2$;
transverse Casimir expectation
$(\widehat S_{\rm tot}^{(1)})^2 + (\widehat S_{\rm tot}^{(2)})^2 =
|\Lambda|\cdot N/2$ (bipartition-independent) by subtraction.

\textbf{Step 117 (PR~\#2905)}: linear N\'eel-state
$\widehat S_{\rm tot}^{(3)}$ expectation real part = $\mathrm{biw.re}$
in
\texttt{LatticeSystem/Quantum/SpinS/NeelStateOfSTotalSpinSOp3ReEqBiwRe.lean}.
$(\langle\Phi_{\rm N\acute{e}el} |\,\widehat S_{\rm tot}^{(3)}\,
   | \Phi_{\rm N\acute{e}el}\rangle).\mathrm{re} =
   (\texttt{bipartiteImbalanceWeight A N}).\mathrm{re}$.
Both equal $(|A|-|\neg A|)\cdot N/2$. The signed analog of the
squared form (Step 116) $\|\mathrm{biw}\|^2$.

\textbf{Step 118 (PR~\#2906)}: transverse Casimir N\'eel
expectation real part in
\texttt{LatticeSystem/Quantum/SpinS/NeelStateOfSTotalSpinSOp12SqReEqHalfCardTimesN.lean}.
$(\langle\Phi_{\rm N\acute{e}el} |\,(\widehat S_{\rm tot}^{(1)})^2
  + (\widehat S_{\rm tot}^{(2)})^2\,| \Phi_{\rm N\acute{e}el}\rangle
   ).\mathrm{re} = |\Lambda|\cdot N/2$.
Bipartition-independent. Together with Steps 109 + 116 yields the
operator-level Casimir decomposition on N\'eel: the bipartition
dependence of $\langle(\widehat S_{\rm tot})^2\rangle_{\rm N\acute{e}el}$
is fully concentrated in the z-axis $(\widehat S_{\rm tot}^{(3)})^2$
component.

\textbf{Step 119 (PR~\#2907)}: linear transverse N\'eel expectations
vanish in
\texttt{LatticeSystem/Quantum/SpinS/NeelStateOfSTotalSpinSOp1Op2ReEqZero.lean}.
$(\langle\Phi_{\rm N\acute{e}el} |\,\widehat S_{\rm tot}^{(1)}\,
   | \Phi_{\rm N\acute{e}el}\rangle).\mathrm{re} = 0$ and
$(\langle\Phi_{\rm N\acute{e}el} |\,\widehat S_{\rm tot}^{(2)}\,
   | \Phi_{\rm N\acute{e}el}\rangle).\mathrm{re} = 0$.
Direct from $\gamma$-4 step 241 (complex form $= 0$) +
\texttt{Complex.zero\_re}. Completes the linear part of the N\'eel
operator-level decomposition: z-axis = $\mathrm{biw.re}$, transverse = 0.

\textbf{Step 120 (PR~\#2908)}: zero $\widehat S_{\rm tot}^{(3)}$
variance on N\'eel in
\texttt{LatticeSystem/Quantum/SpinS/NeelStateOfSTotalSpinSOp3ZeroVariance.lean}.
$\langle(\widehat S_{\rm tot}^{(3)})^2\rangle_{\rm N\acute{e}el}.\mathrm{re}
   = (\langle\widehat S_{\rm tot}^{(3)}\rangle_{\rm N\acute{e}el}.\mathrm{re})^2$.
Combines Step 116 ($\langle(\widehat S_{\rm tot}^{(3)})^2\rangle =
\|\mathrm{biw}\|^2$) with Step 117 ($\langle\widehat S_{\rm tot}^{(3)}\rangle =
\mathrm{biw.re}$) + Step 91 ($\|\mathrm{biw}\|^2 = (\mathrm{biw.re})^2$).
Formally confirms the N\'eel state is an exact $\widehat S_{\rm tot}^{(3)}$
eigenvector at eigenvalue $\mathrm{biw.re} = (|A|-|\neg A|)\cdot N/2$.

\textbf{Step 121 (PR~\#2909)}: strict bound on
$\langle(\widehat S_{\rm tot}^{(3)})^2\rangle_{\rm N\acute{e}el}.\mathrm{re}$
at non-degenerate in
\texttt{LatticeSystem/Quantum/SpinS/NeelStateOfSTotalSpinSOp3SqLtMMaxSqNondegenerate.lean}.
At $|A| \ge 1$, $|\neg A| \ge 1$, $N \ge 1$:
$\langle(\widehat S_{\rm tot}^{(3)})^2\rangle_{\rm N\acute{e}el}.\mathrm{re}
   < |\Lambda|^2\cdot N^2/4 = m_{\rm max}^2$.
Combines Step 116 ($= \|\mathrm{biw}\|^2$) + Step 102
($\|\mathrm{biw}\| < m_{\rm max}$) squared.

\textbf{Step 122 (PR~\#2910)}: strict bound on
$\langle(\widehat S_{\rm tot})^2\rangle_{\rm N\acute{e}el}.\mathrm{re}$
at non-degenerate in
\texttt{LatticeSystem/Quantum/SpinS/NeelStateOfSTotalSpinSSquaredLtMMaxTimesMMaxPlusOneNondegenerate.lean}.
At $|A| \ge 1$, $|\neg A| \ge 1$, $N \ge 1$:
$\langle(\widehat S_{\rm tot})^2\rangle_{\rm N\acute{e}el}.\mathrm{re}
   < m_{\rm max}\cdot(m_{\rm max} + 1) = (|\Lambda|\cdot N/2)\cdot
   (|\Lambda|\cdot N/2 + 1)$.
Combines Step 109 + Step 102 squared. Strictly below the
fully-polarised ferromagnetic eigenvalue at non-degenerate.

\textbf{Step 123 (PR~\#2911)}: strict negativity of
$\langle\widehat H_{\rm toy}^S\rangle_{\rm N\acute{e}el}.\mathrm{re}$
at non-degenerate in
\texttt{LatticeSystem/Quantum/SpinS/NeelStateOfSHeisenbergToyExpectationReNegNondegenerate.lean}.
At $|A| \ge 1$, $|\neg A| \ge 1$, $N \ge 1$:
$(\langle\Phi_{\rm N\acute{e}el} |\,\widehat H_{\rm toy}^S\,
   | \Phi_{\rm N\acute{e}el}\rangle).\mathrm{re} < 0$.
Direct from $\langle\widehat H_{\rm toy}^S\rangle_{\rm N\acute{e}el} =
-|A|\cdot|\neg A|\cdot N^2/2$ ($\gamma$-4 step 131) +
$|A|\cdot|\neg A|\cdot N^2 > 0$ at non-degenerate.

\textbf{Step 124 (PR~\#2912)}: N\'eel sublattice Casimir expectations
in
\texttt{LatticeSystem/Quantum/SpinS/NeelStateOfSSublatticeSpinSquaredSExpectation.lean}.
$\langle\Phi_{\rm N\acute{e}el} |\,(\widehat S_A)^2\,
   | \Phi_{\rm N\acute{e}el}\rangle = s_A\cdot(s_A+1)$ and
$\langle\Phi_{\rm N\acute{e}el} |\,(\widehat S_{\neg A})^2\,
   | \Phi_{\rm N\acute{e}el}\rangle = s_B\cdot(s_B+1)$
where $s_A = |A|\cdot N/2$, $s_B = |\neg A|\cdot N/2$. Direct from
existing eigenvector relations + $\langle\Phi_{\rm N\acute{e}el}|
\Phi_{\rm N\acute{e}el}\rangle = 1$. Confirms N\'eel has maximum
spin on both sublattices --- a key input to Tasaki §2.5
Theorem 2.3 ($\gamma$-4): the N\'eel state matches the predicted GS
subspace's sublattice Casimir constraints.

\textbf{Step 125 (PR~\#2913)}: N\'eel state in joint sublattice-Casimir
eigenspace in
\texttt{LatticeSystem/Quantum/SpinS/NeelStateOfSMemJointSublatticeCasimirEigenspace.lean}.
$\Phi_{\rm N\acute{e}el}(A, N) \in
   \mathrm{eigenspace}((\widehat S_A)^2, s_A\cdot(s_A+1)) \cap
   \mathrm{eigenspace}((\widehat S_{\neg A})^2, s_B\cdot(s_B+1))$.
Two-thirds of the predicted GS subspace conditions. The N\'eel
state is NOT in the full predicted GS subspace because its
$(\widehat S_{\rm tot})^2$ expectation ($\|\mathrm{biw}\|^2 +
|\Lambda|\cdot N/2$, Step 109) exceeds the predicted eigenvalue
$\|\mathrm{biw}\|\cdot(\|\mathrm{biw}\|+1)$ by $\min\cdot N$
(Step 113).

\textbf{Step 126 (PR~\#2914)}: N\'eel = all-aligned + GS membership
at saturated $|\neg A| = 0$ in
\texttt{LatticeSystem/Quantum/SpinS/NeelStateOfSEqAllAlignedStateAtSaturated.lean}.
At $|\neg A| = 0$:
$\mathtt{neelConfigOfS}\,A\,N = (\lambda \_, 0)$, hence
$\Phi_{\rm N\acute{e}el}(A, N) = \mathtt{allAlignedStateS}\,\Lambda\,N\,0$.
Combined with the existing all-up GS witness (PR \#2782) yields
$\Phi_{\rm N\acute{e}el}(A, N) \in
   \mathtt{bipartiteToyGroundStateSubspacePredicted}\,A\,N$ at
saturated. The N\'eel state IS the predicted GS at saturated edges.

\textbf{Step 127 (PR~\#2915)}: N\'eel = all-down at saturated
complement $|A| = 0$ in
\texttt{LatticeSystem/Quantum/SpinS/NeelStateOfSEqAllAlignedStateAtSaturatedComplement.lean}.
At $|A| = 0$:
$\mathtt{neelConfigOfS}\,A\,N = (\lambda \_, \mathrm{Fin.last}\,N)$,
hence $\Phi_{\rm N\acute{e}el}(A, N) =
\mathtt{allAlignedStateS}\,\Lambda\,N\,(\mathrm{Fin.last}\,N)$.
Companion of Step 126.

\textbf{Step 128 (PR~\#2916)}: N\'eel $\in$ predicted GS subspace
(complement orientation) at $|A| = 0$ in
\texttt{LatticeSystem/Quantum/SpinS/NeelStateOfSMemBipartiteToyGSPredictedComplementSaturated.lean}.
$\Phi_{\rm N\acute{e}el}(A, N) \in
\mathtt{bipartiteToyGroundStateSubspacePredicted}\,(\lambda x, !A\,x)\,N$
at $|A| = 0$. Combines Step 127 (N\'eel = all-down) + the existing
all-down GS witness (PR \#2783). The N\'eel state IS the predicted
GS at the saturated complement edge under the complement
orientation.

\textbf{Step 129 (PR~\#2917)}: excess of symm-energy above balanced
minimum in
\texttt{LatticeSystem/Quantum/SpinS/BipartiteToyMinEnergySymmReExcessAboveBalanced.lean}.
$(\mathrm{symm}).\mathrm{re} + |\Lambda|^2\cdot N^2/8 + |\Lambda|\cdot N/2
   = \|\mathrm{biw}\|^2/2 + \|\mathrm{biw}\|$.
Quadratic in $\|\mathrm{biw}\|$. Vanishes at balanced; equals
$m_{\rm max}\cdot(m_{\rm max}+2)/2$ at saturated.

\textbf{Step 130 (PR~\#2918)}: squared $m_{\rm max}$ upper bound on
imbalance norm in
\texttt{LatticeSystem/Quantum/SpinS/BipartiteImbalanceWeightNormSqLeMMaxSq.lean}.
$\|\mathtt{bipartiteImbalanceWeight}\,A\,N\|^2 \le |\Lambda|^2\cdot N^2/4
= m_{\rm max}^2$. Direct from the linear $m_{\rm max}$ bound squared.

\textbf{Step 131 (PR~\#2919)}: N\'eel $\notin$ predicted GS at
non-degenerate in
\texttt{LatticeSystem/Quantum/SpinS/NeelStateOfSNotMemBipartiteToyGSPredictedNondegenerate.lean}.
At $|A| \ge 1$, $|\neg A| \ge 1$, $N \ge 1$, $|\neg A| \le |A|$:
$\Phi_{\rm N\acute{e}el}(A, N) \notin
\mathtt{bipartiteToyGroundStateSubspacePredicted}\,A\,N$. By
contradiction: GS membership would force
$\langle(\widehat S_{\rm tot})^2\rangle_{\rm N\acute{e}el} =
\|\mathrm{biw}\|\cdot(\|\mathrm{biw}\|+1)$, but Step 114 shows
strict $>$. Complements Step 126 (saturated case: N\'eel $\in$ GS).

\textbf{Step 132 (PR~\#2921)}: balanced-case N\'eel toy-Hamiltonian
expectation in
\texttt{LatticeSystem/Quantum/SpinS/NeelStateOfSHeisenbergToyExpectationReBalanced.lean}.
At $|A| = |\neg A|$, $N \ge 1$:
$(\langle\Phi_{\rm N\acute{e}el} |\,\widehat H_{\rm toy}^S\,
   | \Phi_{\rm N\acute{e}el}\rangle).\mathrm{re} = -|\Lambda|^2\cdot N^2/8$.
Direct from Step 99 ($\|\mathrm{biw}\|^2/2 - |\Lambda|^2\cdot N^2/8$)
+ $\|\mathrm{biw}\| = 0$ at balanced. Variational benchmark for
the actual GS energy.

\textbf{Step 133 (PR~\#2922)}: saturated-edge N\'eel toy-Hamiltonian
expectation $= 0$ in
\texttt{LatticeSystem/Quantum/SpinS/NeelStateOfSHeisenbergToyExpectationReSaturated.lean}.
At $|\neg A| = 0$ (or $|A| = 0$):
$(\langle\Phi_{\rm N\acute{e}el} |\,\widehat H_{\rm toy}^S\,
   | \Phi_{\rm N\acute{e}el}\rangle).\mathrm{re} = 0$.
N\'eel achieves zero energy at saturated edges, matching the
predicted GS energy --- confirming N\'eel IS a GS at saturated
edges.

\textbf{Step 134 (PR~\#2923)}: joint sublattice-Casimir eigenspace
strictly contains predicted GS at non-degenerate in
\texttt{LatticeSystem/Quantum/SpinS/JointSublatticeCasimirStrictSupsetPredictedGS.lean}.
At $|\neg A| \le |A|$ with both positive and $N \ge 1$:
$\mathtt{bipartiteToyGroundStateSubspacePredicted}\,A\,N \subsetneq
   \mathrm{eigenspace}((\widehat S_A)^2, s_A(s_A+1))
   \cap \mathrm{eigenspace}((\widehat S_{\neg A})^2, s_B(s_B+1))$.
The N\'eel state witnesses the strict inclusion: in the joint
Casimir eigenspace (Step 125) but not in the predicted GS
(Step 131). The "extra" states have maximum sublattice spins but
wrong total spin.

\textbf{Step 135 (PR~\#2924)}: predicted GS finrank $<$ joint Casimir
finrank at non-deg in
\texttt{LatticeSystem/Quantum/SpinS/PredictedGSFinrankLtJointCasimirFinrank.lean}.
$\mathrm{finrank}(\mathtt{bipartiteToyGroundStateSubspacePredicted}\,A\,N)
   < \mathrm{finrank}(\mathtt{jointSublatticeCasimirEigenspace}\,A\,N)$.
Direct from Step 134's strict subspace inclusion in finite-dim
space. Quantifies the "extra" subspace beyond the predicted GS.

\textbf{Step 136 (PR~\#2925)}: predicted GS $=$ joint Casimir
eigenspace at saturated in
\texttt{LatticeSystem/Quantum/SpinS/PredictedGSEqJointCasimirEigenspaceSaturated.lean}.
At $|\neg A| = 0$:
$\mathtt{bipartiteToyGroundStateSubspacePredicted}\,A\,N =
   \mathtt{jointSublatticeCasimirEigenspace}\,A\,N$.
Complement of Step 134's strict inclusion at non-deg. At saturated,
$(\widehat S_{\neg A})^2 = 0$ makes that eigenspace condition vacuous,
and $(\widehat S_A)^2 = (\widehat S_{\rm tot})^2$ makes the
$(\widehat S_A)^2$ condition imply the $(\widehat S_{\rm tot})^2$
condition.

\textbf{Step 137 (PR~\#2926)}: complement-orientation analog of
Step 136 in
\texttt{LatticeSystem/Quantum/SpinS/PredictedGSEqJointCasimirEigenspaceSaturatedComplement.lean}.
At $|A| = 0$:
$\mathtt{bipartiteToyGroundStateSubspacePredicted}\,(\lambda x, !A\,x)\,N =
\mathtt{jointSublatticeCasimirEigenspace}\,(\lambda x, !A\,x)\,N$.
Direct from Step 136 via $|\neg\neg A| = |A| = 0$.

\textbf{Step 138 (PR~\#2927)}: joint Casimir eigenspace finrank at
saturated in
\texttt{LatticeSystem/Quantum/SpinS/JointCasimirEigenspaceFinrankSaturated.lean}.
At $|\neg A| = 0$:
$\mathrm{finrank}(\mathtt{jointSublatticeCasimirEigenspace}\,A\,N) =
|\Lambda|\cdot N + 1$. Direct from Step 136 (predicted GS = joint
Casimir at saturated) + existing saturated finrank result. The
dimension matches the SU(2) multiplet $2m_{\rm max} + 1$ at
saturated edges.

\textbf{Step 139 (PR~\#2928)}: complement saturated joint Casimir
finrank in
\texttt{LatticeSystem/Quantum/SpinS/JointCasimirEigenspaceFinrankSaturatedComplement.lean}.
At $|A| = 0$:
$\mathrm{finrank}(\mathtt{jointSublatticeCasimirEigenspace}\,
   (\lambda x, !A\,x)\,N) = |\Lambda|\cdot N + 1$.
Mirror of Step 138 via $|\neg\neg A| = |A| = 0$.

\textbf{Step 140 (PR~\#2929)}: joint Casimir $=$ $(\widehat S_{\rm tot})^2$
eigenspace at saturated in
\texttt{LatticeSystem/Quantum/SpinS/JointCasimirEqTotalSpinSSquaredEigenspaceSaturated.lean}.
At $|\neg A| = 0$:
$\mathtt{jointSublatticeCasimirEigenspace}\,A\,N =
\mathrm{eigenspace}(\widehat S_{\rm tot}^2, \mathtt{saturatedFerromagnetCasimirEigenvalueS}\,\Lambda\,N)$.
Composes Step 136 (predicted GS = joint Casimir at saturated) +
the existing predicted GS = $(\widehat S_{\rm tot})^2$ eigenspace at
max Casimir at saturated. The joint Casimir IS the fully-polarised
ferromagnetic SU(2) multiplet at saturated.

\textbf{Step 141 (PR~\#2930)}: predicted $(\widehat S_{\rm tot})^2$
eigenvalue in $\|\mathrm{biw}\|$ form in
\texttt{LatticeSystem/Quantum/SpinS/PredictedTotalSpinSquaredEigenvalueViaImbalanceNorm.lean}.
At $|\neg A| \le |A|$:
$((s_A - s_B)\cdot((s_A - s_B) + 1)).\mathrm{re} =
\|\mathrm{biw}\|\cdot(\|\mathrm{biw}\| + 1)$.
Identifies Tasaki's predicted $S_{\rm tot}(S_{\rm tot}+1)$ with the
imbalance-norm quadratic, with $\|\mathrm{biw}\| = ||A|-|\neg A||\cdot N/2$
playing the role of $S_{\rm tot}$.

\textbf{Step 142 (PR~\#2931)}: complement-orientation predicted
$(\widehat S_{\rm tot})^2$ eigenvalue in
\texttt{LatticeSystem/Quantum/SpinS/PredictedTotalSpinSquaredEigenvalueViaImbalanceNormComplement.lean}.
At $|A| \le |\neg A|$:
$((s_B - s_A)\cdot((s_B - s_A) + 1)).\mathrm{re} =
\|\mathrm{biw}\|\cdot(\|\mathrm{biw}\| + 1)$.
Same eigenvalue as Step 141 via complement-norm invariance.

\textbf{Step 143 (PR~\#2932)}: unified signed-form predicted
$(\widehat S_{\rm tot})^2$ eigenvalue (no orientation hypothesis) in
\texttt{LatticeSystem/Quantum/SpinS/PredictedTotalSpinSquaredEigenvalueViaImbalanceRe.lean}.
$((s_A - s_B)\cdot((s_A - s_B) + 1)).\mathrm{re} =
\mathrm{biw.re}\cdot(\mathrm{biw.re} + 1)$.
Generalises Steps 141 and 142 via $\mathrm{biw.re} = \pm\|\mathrm{biw}\|$.

\textbf{Step 144 (PR~\#2933)}: N\'eel-vs-predicted $(\widehat S_{\rm tot})^2$
gap $= 0$ iff saturated edge at $N \ge 1$ in
\texttt{LatticeSystem/Quantum/SpinS/NeelTotalSpinSSqGapEqZeroIffSaturated.lean}.
$\mathrm{gap} = 0 \iff |A| = 0 \vee |\neg A| = 0$.
Direct from Step 113 (gap $= \min\cdot N$) + $\min\cdot N = 0 \iff
\min = 0$ at $N \ge 1$.

\textbf{Step 145 (PR~\#2934)}: energy gap iff characterisation at
$N \ge 1$ in
\texttt{LatticeSystem/Quantum/SpinS/NeelToyEnergyGapEqZeroIffSaturated.lean}.
$(\langle\Phi_{\rm N\acute{e}el}|\widehat H_{\rm toy}^S|\Phi_{\rm N\acute{e}el}\rangle
   - \texttt{bipartiteToyMinEnergyPredictedSymm}\,A\,N).\mathrm{re} = 0
   \iff |A| = 0 \vee |\neg A| = 0$.
Energy parallel of Step 144; both gaps vanish simultaneously at
saturated edges.

\textbf{Step 146 (PR~\#2935)}: non-negativity of predicted
$(\widehat S_{\rm tot})^2$ eigenvalue at $|\neg A| \le |A|$ in
\texttt{LatticeSystem/Quantum/SpinS/PredictedTotalSpinSquaredEigenvalueReNonneg.lean}.
$0 \le ((s_A - s_B)\cdot((s_A - s_B) + 1)).\mathrm{re}$.
Direct from Step 141 ($= \|\mathrm{biw}\|\cdot(\|\mathrm{biw}\|+1)$)
$+ \|\mathrm{biw}\| \ge 0$. The predicted Casimir eigenvalue is
physical under the proper orientation.

\textbf{Step 147 (PR~\#2936)}: predicted $(\widehat S_{\rm tot})^2$
eigenvalue $= 0$ iff balanced at $|\neg A| \le |A|$, $N \ge 1$ in
\texttt{LatticeSystem/Quantum/SpinS/PredictedTotalSpinSquaredEigenvalueEqZeroIffBalanced.lean}.
$((s_A - s_B)\cdot((s_A - s_B) + 1)).\mathrm{re} = 0 \iff |A| = |\neg A|$.
The predicted GS has singlet $S_{\rm tot} = 0$ exactly at balanced.

\textbf{Step 148 (PR~\#2937)}: strict positivity of predicted
$(\widehat S_{\rm tot})^2$ eigenvalue at strict unbalanced in
\texttt{LatticeSystem/Quantum/SpinS/PredictedTotalSpinSquaredEigenvaluePosUnbalanced.lean}.
At $|\neg A| < |A|$, $N \ge 1$:
$0 < ((s_A - s_B)\cdot((s_A - s_B) + 1)).\mathrm{re}$.
The predicted GS has strictly positive $S_{\rm tot}(S_{\rm tot}+1)$
at strict unbalanced --- non-singlet.

\textbf{Step 149 (PR~\#2938)}: saturated-edge predicted
$(\widehat S_{\rm tot})^2$ eigenvalue $= m_{\rm max}(m_{\rm max}+1)$ in
\texttt{LatticeSystem/Quantum/SpinS/PredictedEigenvalueSaturatedEqMMaxQuad.lean}.
At $|\neg A| = 0$: predicted eigenvalue $=
(|\Lambda|\cdot N/2)\cdot(|\Lambda|\cdot N/2 + 1)$. Direct from
Step 141 $+ \|\mathrm{biw}\| = m_{\rm max}$ at saturated.

\textbf{Step 150 (PR~\#2939)}: balanced-case predicted
$(\widehat S_{\rm tot})^2$ eigenvalue $= 0$ (singlet) in
\texttt{LatticeSystem/Quantum/SpinS/PredictedEigenvalueBalancedEqZero.lean}.
At $|A| = |\neg A|$: predicted eigenvalue $= 0$. Direct from
Step 143 (signed unified form) $+\ \mathrm{biw.re} = 0$ at balanced.
The predicted GS at balanced has $S_{\rm tot} = 0$ (singlet), matching
Tasaki's balanced-case prediction.

\textbf{Step 151 (PR~\#2941)}: complement saturated-edge predicted
$(\widehat S_{\rm tot})^2$ eigenvalue $= m_{\rm max}(m_{\rm max}+1)$ in
\texttt{LatticeSystem/Quantum/SpinS/PredictedEigenvalueComplementSaturatedEqMMaxQuad.lean}.
At $|A| = 0$: complement-orientation predicted eigenvalue $=
(|\Lambda|\cdot N/2)(|\Lambda|\cdot N/2 + 1)$. Mirror of Step 149.

\textbf{Step 152 (PR~\#2942)}: upper bound on predicted
$(\widehat S_{\rm tot})^2$ eigenvalue $\le m_{\rm max}(m_{\rm max}+1)$
at $|\neg A| \le |A|$ in
\texttt{LatticeSystem/Quantum/SpinS/PredictedEigenvalueLeMMaxQuad.lean}.
$((s_A-s_B)\cdot((s_A-s_B)+1)).\mathrm{re} \le
(|\Lambda|\cdot N/2)(|\Lambda|\cdot N/2 + 1)$.
Direct from Step 141 + $\|\mathrm{biw}\| \le m_{\rm max}$ +
monotonicity of $x\cdot(x+1)$ on $[0, \infty)$. Equality is attained
at saturated edges (Step 149).

\textbf{Step 153 (PR~\#2943)}: upper bound on complement predicted
$(\widehat S_{\rm tot})^2$ eigenvalue $\le m_{\rm max}(m_{\rm max}+1)$
at $|A| \le |\neg A|$ in
\texttt{LatticeSystem/Quantum/SpinS/PredictedEigenvalueComplementLeMMaxQuad.lean}.
$((s_B-s_A)\cdot((s_B-s_A)+1)).\mathrm{re} \le
(|\Lambda|\cdot N/2)(|\Lambda|\cdot N/2 + 1)$. Mirror of Step 152 via
Step 142 + $\|\mathrm{biw}\| \le m_{\rm max}$ + monotonicity of
$x\cdot(x+1)$ on $[0, \infty)$. Equality is attained at saturated
complement (Step 151).

\textbf{Step 154 (PR~\#2944)}: unified upper bound on predicted
$(\widehat S_{\rm tot})^2$ eigenvalue $\le m_{\rm max}(m_{\rm max}+1)$
\emph{without orientation hypothesis} in
\texttt{LatticeSystem/Quantum/SpinS/PredictedEigenvalueLeMMaxQuadUnified.lean}.
$((s_A-s_B)\cdot((s_A-s_B)+1)).\mathrm{re} \le
(|\Lambda|\cdot N/2)(|\Lambda|\cdot N/2 + 1)$ in the
$(s_A - s_B)$ form. Direct from Step 143 (signed unified form
$= \mathrm{biw.re}\cdot(\mathrm{biw.re}+1)$) +
$\|\mathrm{biw}\| \le m_{\rm max}$ + $\mathrm{biw.im} = 0$ giving
$|\mathrm{biw.re}| \le m_{\rm max}$, then \texttt{nlinarith} on
$x(x+1)$ over $[-m_{\rm max}, m_{\rm max}]$. Unifies Step 152
($|\neg A| \le |A|$) and Step 153 ($|A| \le |\neg A|$ complement).

\textbf{Step 155 (PR~\#2945)}: strict upper bound on predicted
$(\widehat S_{\rm tot})^2$ eigenvalue $< m_{\rm max}(m_{\rm max}+1)$
\emph{at non-degenerate} ($|A| \ge 1$, $|\neg A| \ge 1$, $N \ge 1$) in
\texttt{LatticeSystem/Quantum/SpinS/PredictedEigenvalueLtMMaxQuadNondegenerate.lean}.
Strengthens Step 154 from $\le$ to $<$. Direct from Step 143 + the
strict imbalance-norm bound $\|\mathrm{biw}\| < m_{\rm max}$
(non-degenerate) + $\mathrm{biw.im} = 0$ giving
$|\mathrm{biw.re}| < m_{\rm max}$, then \texttt{nlinarith} on
$x(x+1)$ over $(-m_{\rm max}, m_{\rm max})$. At non-degenerate, the
predicted GS Casimir is strictly below the fully-polarised
ferromagnetic value; this completes the "saturated iff equality"
characterisation.

\textbf{Step 156 (PR~\#2946)}: strict upper bound on complement
predicted $(\widehat S_{\rm tot})^2$ eigenvalue
$< m_{\rm max}(m_{\rm max}+1)$ \emph{at non-degenerate}
($|A| \le |\neg A|$, $|A| \ge 1$, $|\neg A| \ge 1$, $N \ge 1$) in
\texttt{LatticeSystem/Quantum/SpinS/PredictedEigenvalueComplementLtMMaxQuadNondegenerate.lean}.
Mirror of Step 155 strengthening Step 153 from $\le$ to $<$. Direct
from Step 142 (complement $= \|\mathrm{biw}\|\cdot(\|\mathrm{biw}\|+1)$)
+ Step 144 ($\|\mathrm{biw}\| < m_{\rm max}$ at non-degenerate) +
monotonicity of $x(x+1)$ on $[0, m_{\rm max})$ via \texttt{nlinarith}.

\textbf{Step 157 (PR~\#2947)}: generalised strict predicted
$(\widehat S_{\rm tot})^2$ eigenvalue $< m_{\rm max}(m_{\rm max}+1)$
at $|\neg A| \ge 1$, $N \ge 1$ (\emph{without} $|A| \ge 1$) in
\texttt{LatticeSystem/Quantum/SpinS/PredictedEigenvalueLtMMaxQuadOfCardNotAPos.lean}.
Generalises Step 155 by removing the $|A| \ge 1$ hypothesis. Direct
from Step 143 + strict $\mathrm{biw.re} < m_{\rm max}$ (since
$|A| - |\neg A| < |\Lambda|$ when $|\neg A| \ge 1$, scaled by
$N \ge 1$) + unconditional $-m_{\rm max} \le \mathrm{biw.re}$ (from
$\|\mathrm{biw}\| \le m_{\rm max}$ + $\mathrm{biw.im} = 0$), then
\texttt{nlinarith}. Sets up the equality characterisation
$\mathrm{eigenvalue} = m_{\rm max}(m_{\rm max}+1) \Leftrightarrow
|\neg A| = 0$.

\textbf{Step 158 (PR~\#2948)}: iff characterisation predicted
$(\widehat S_{\rm tot})^2$ eigenvalue
$= m_{\rm max}(m_{\rm max}+1) \Leftrightarrow |\neg A| = 0$ (at
$N \ge 1$) in
\texttt{LatticeSystem/Quantum/SpinS/PredictedEigenvalueEqMMaxQuadIffCardNotAZero.lean}.
Forward direction = contrapositive of Step 157; backward = Step 149.
The predicted Casimir hits the fully-polarised ferromagnetic value
\emph{iff} all sites are in $A$ (saturated edge in the
$(s_A - s_B)$ orientation).

\textbf{Step 159 (PR~\#2949)}: iff characterisation (complement)
complement predicted $(\widehat S_{\rm tot})^2$ eigenvalue
$= m_{\rm max}(m_{\rm max}+1) \Leftrightarrow |A| = 0$ at
$|A| \le |\neg A|$, $N \ge 1$ in
\texttt{LatticeSystem/Quantum/SpinS/PredictedEigenvalueComplementEqMMaxQuadIffCardAZero.lean}.
Mirror of Step 158. Forward = contrapositive of Step 156 (combined
with $|A| \le |\neg A|$ to derive $|\neg A| \ge 1$ from $|A| \ge 1$);
backward = Step 151. Completes the complement equality
characterisation.

\textbf{Step 160 (PR~\#2950)}: range characterisation predicted
$(\widehat S_{\rm tot})^2$ eigenvalue
$\in [0, m_{\rm max}(m_{\rm max}+1)]$ at proper orientation
$|\neg A| \le |A|$ in
\texttt{LatticeSystem/Quantum/SpinS/PredictedEigenvalueRangeIcc.lean}.
Combines Step 145 ($\ge 0$) and Step 152 ($\le m_{\rm max}(m_{\rm max}+1)$)
into a single closed-interval statement. Endpoints attained at
balanced (lower) and saturated (upper).

\textbf{Step 161 (PR~\#2951)}: non-negativity of complement predicted
$(\widehat S_{\rm tot})^2$ eigenvalue at $|A| \le |\neg A|$ in
\texttt{LatticeSystem/Quantum/SpinS/PredictedTotalSpinSquaredEigenvalueReNonnegComplement.lean}.
Mirror of Step 145. Direct from Step 142
(complement $= \|\mathrm{biw}\|\cdot(\|\mathrm{biw}\|+1)$) +
$\|\mathrm{biw}\| \ge 0$.

\textbf{Step 162 (PR~\#2952)}: range characterisation (complement)
complement predicted $(\widehat S_{\rm tot})^2$ eigenvalue
$\in [0, m_{\rm max}(m_{\rm max}+1)]$ at complement orientation
$|A| \le |\neg A|$ in
\texttt{LatticeSystem/Quantum/SpinS/PredictedEigenvalueComplementRangeIcc.lean}.
Mirror of Step 160. Combines Step 161 ($\ge 0$) and Step 153
($\le m_{\rm max}(m_{\rm max}+1)$).

\textbf{Step 163 (PR~\#2953)}: complement predicted $(\widehat S_{\rm tot})^2$
eigenvalue $= 0 \Leftrightarrow |A| = |\neg A|$ at $|A| \le |\neg A|$,
$N \ge 1$ in
\texttt{LatticeSystem/Quantum/SpinS/PredictedEigenvalueComplementEqZeroIffBalanced.lean}.
Mirror of Step 146. Direct from Step 142 + Step 81
($\|\mathrm{biw}\| = 0 \Leftrightarrow |A|=|\neg A|$ at $N \ge 1$) +
$\|\mathrm{biw}\| + 1 \ge 1 > 0$.

\textbf{Step 164 (PR~\#2954)}: strict positivity of complement
predicted $(\widehat S_{\rm tot})^2$ eigenvalue at strict unbalanced
$|A| < |\neg A|$, $N \ge 1$ in
\texttt{LatticeSystem/Quantum/SpinS/PredictedEigenvalueComplementPosUnbalanced.lean}.
Mirror of Step 147. Direct from Step 142 + $\|\mathrm{biw}\| > 0$
(strict unbalanced norm bound) via \texttt{nlinarith}.

\textbf{Step 165 (PR~\#2955)}: unified signed-form complement predicted
$(\widehat S_{\rm tot})^2$ eigenvalue
$((s_B - s_A)\cdot((s_B - s_A) + 1)).\mathrm{re} =
\mathrm{biw.re}\cdot(\mathrm{biw.re} - 1)$ \emph{without orientation
hypothesis} in
\texttt{LatticeSystem/Quantum/SpinS/PredictedTotalSpinSquaredEigenvalueViaImbalanceReComplement.lean}.
Mirror of Step 143. Step 142 ($\|\mathrm{biw}\|\cdot(\|\mathrm{biw}\|+1)$)
is an orientation-specialised consequence at $|A| \le |\neg A|$ where
$\mathrm{biw.re} \le 0$ so $-\mathrm{biw.re} = \|\mathrm{biw}\|$. Allows
orientation-free reasoning about the complement form.

\textbf{Step 166 (PR~\#2956)}: unified iff complement predicted
$(\widehat S_{\rm tot})^2$ eigenvalue
$= m_{\rm max}(m_{\rm max}+1) \Leftrightarrow |A| = 0$ at $N \ge 1$
(\emph{no} orientation hypothesis) in
\texttt{LatticeSystem/Quantum/SpinS/PredictedEigenvalueComplementEqMMaxQuadIffCardAZeroUnified.lean}.
Generalises Step 159. Direct from Step 165 (signed complement form
$\mathrm{biw.re}\cdot(\mathrm{biw.re}-1)$) + $|\mathrm{biw.re}| \le m_{\rm max}$:
factoring $(\mathrm{biw.re}-(m_{\rm max}+1))(\mathrm{biw.re}+m_{\rm max}) = 0$
and using $\mathrm{biw.re} \le m_{\rm max} < m_{\rm max}+1$ forces
$\mathrm{biw.re} = -m_{\rm max}$, hence $|A| = 0$ at $N \ge 1$.

\textbf{Step 167 (PR~\#2957)}: generalised strict complement predicted
$(\widehat S_{\rm tot})^2$ eigenvalue $< m_{\rm max}(m_{\rm max}+1)$
at $|A| \ge 1$, $N \ge 1$ (\emph{no} $|\neg A| \ge 1$, no orientation)
in
\texttt{LatticeSystem/Quantum/SpinS/PredictedEigenvalueComplementLtMMaxQuadOfCardAPos.lean}.
Mirror of Step 157. Direct from Step 165 + strict
$\mathrm{biw.re} > -m_{\rm max}$ (since $|A| - |\neg A| > -|\Lambda|$
when $|A| \ge 1$, scaled by $N \ge 1$) + unconditional
$\mathrm{biw.re} \le m_{\rm max}$, then \texttt{nlinarith}.

\textbf{Step 168 (PR~\#2958)}: canonical predicted $(\widehat S_{\rm tot})^2$
eigenvalue value at complement-saturated $|A| = 0$:
$\mathrm{eigenvalue.re} = m_{\rm max}(m_{\rm max} - 1)$ (NOT the
ferromagnetic $m_{\rm max}(m_{\rm max} + 1)$) in
\texttt{LatticeSystem/Quantum/SpinS/PredictedEigenvalueCanonicalAtComplementSaturated.lean}.
Mirror of Step 149 in the wrong orientation. Direct from Step 143
(signed canonical $\mathrm{biw.re}\cdot(\mathrm{biw.re}+1)$) +
$\mathrm{biw.re} = -m_{\rm max}$ at $|A| = 0$. The canonical form at
complement-saturated has Casimir strictly below the ferromagnetic
value — wrong orientation; the complement form (Step 151) gives the
correct value.

\textbf{Step 169 (PR~\#2959)}: complement predicted $(\widehat S_{\rm tot})^2$
eigenvalue value at canonical-saturated $|\neg A| = 0$:
$\mathrm{complement\_eigenvalue.re} = m_{\rm max}(m_{\rm max} - 1)$
(NOT the ferromagnetic $m_{\rm max}(m_{\rm max} + 1)$) in
\texttt{LatticeSystem/Quantum/SpinS/PredictedEigenvalueComplementAtCanonicalSaturated.lean}.
Mirror of Step 168 — complement form at the wrong orientation edge.
Direct from Step 165 (signed complement form
$\mathrm{biw.re}\cdot(\mathrm{biw.re}-1)$) + $\mathrm{biw.re} = m_{\rm max}$
at $|\neg A| = 0$.

\textbf{Step 170 (PR~\#2960)}: canonical-minus-complement predicted
$(\widehat S_{\rm tot})^2$ eigenvalue difference $= 2 \cdot \mathrm{biw.re}$
\emph{unconditionally} in
\texttt{LatticeSystem/Quantum/SpinS/PredictedEigenvalueCanonicalComplementDiff.lean}.
Direct from Step 143 (canonical $= \mathrm{biw.re}\cdot(\mathrm{biw.re}+1)$)
$-$ Step 165 (complement $= \mathrm{biw.re}\cdot(\mathrm{biw.re}-1)$)
$= 2\cdot\mathrm{biw.re}$ via \texttt{ring}. The canonical-complement
gap measures twice the signed imbalance; vanishes at balanced.

\textbf{Step 171 (PR~\#2962)}: complement saturated joint Casimir
eigenspace $=$ $(\widehat S_{\rm tot})^2$ eigenspace at max Casimir
in
\texttt{LatticeSystem/Quantum/SpinS/JointCasimirEqTotalSpinSSquaredEigenspaceSaturatedComplement.lean}.
Mirror of the canonical saturated theorem
(\texttt{jointSublatticeCasimirEigenspace\_eq\_totalSpinSSquaredEigenspace\_of\_cardNotA\_zero})
for the complement sublattice function via $|\neg\neg A| = |A| = 0$
reduction.

\textbf{Step 172 (PR~\#2963)}: canonical $+$ complement predicted
$(\widehat S_{\rm tot})^2$ eigenvalue sum $= 2 \cdot \mathrm{biw.re}^2$
\emph{unconditionally} in
\texttt{LatticeSystem/Quantum/SpinS/PredictedEigenvalueCanonicalComplementSum.lean}.
Companion of Step 170 (difference $= 2 \cdot \mathrm{biw.re}$).
Direct from Step 143 $+$ Step 165 via \texttt{ring}. Together gives the
affine parameterisation $\mathrm{canonical} = \mathrm{biw.re}^2 + \mathrm{biw.re}$,
$\mathrm{complement} = \mathrm{biw.re}^2 - \mathrm{biw.re}$.

\textbf{Step 173 (PR~\#2964)}: canonical predicted $(\widehat S_{\rm tot})^2$
eigenvalue $= \mathrm{biw.re}^2 + \mathrm{biw.re}$ \emph{unconditionally}
in
\texttt{LatticeSystem/Quantum/SpinS/PredictedEigenvalueCanonicalEqSquareAdd.lean}.
Expanded polynomial form of Step 143. Convenient for \texttt{nlinarith}
and arithmetic on the canonical signed form.

\textbf{Step 174 (PR~\#2965)}: complement predicted $(\widehat S_{\rm tot})^2$
eigenvalue $= \mathrm{biw.re}^2 - \mathrm{biw.re}$ \emph{unconditionally}
in
\texttt{LatticeSystem/Quantum/SpinS/PredictedEigenvalueComplementEqSquareSub.lean}.
Expanded polynomial form of Step 165. Companion of Step 173 — together
gives the symmetric affine parameterisation
$\mathrm{canonical} = \mathrm{biw.re}^2 + \mathrm{biw.re}$,
$\mathrm{complement} = \mathrm{biw.re}^2 - \mathrm{biw.re}$.

\textbf{Step 175 (PR~\#2966)}: $\max(\mathrm{canonical}, \mathrm{complement})$
predicted $(\widehat S_{\rm tot})^2$ eigenvalue
$= \|\mathrm{biw}\|\cdot(\|\mathrm{biw}\|+1)$ \emph{unconditionally} in
\texttt{LatticeSystem/Quantum/SpinS/PredictedEigenvalueMaxCanonicalComplement.lean}.
The \emph{physical} predicted $(\widehat S_{\rm tot})^2$ eigenvalue is the
max over both orientations. Direct from Step 173 + Step 174 by case-split
on sign of $\mathrm{biw.re}$ and $\mathrm{biw.im} = 0 \Rightarrow
\|\mathrm{biw}\| = |\mathrm{biw.re}|$. Unifies Step 141 (canonical at
$|\neg A| \le |A|$) and Step 142 (complement at $|A| \le |\neg A|$) into a
single orientation-free statement.

\textbf{Step 176 (PR~\#2967)}: $\min(\mathrm{canonical}, \mathrm{complement})$
predicted $(\widehat S_{\rm tot})^2$ eigenvalue
$= \|\mathrm{biw}\|\cdot(\|\mathrm{biw}\|-1)$ \emph{unconditionally} in
\texttt{LatticeSystem/Quantum/SpinS/PredictedEigenvalueMinCanonicalComplement.lean}.
Companion of Step 175 (max). The $\min$ is the "wrong"-orientation
eigenvalue, strictly below the physical predicted value. At saturated
edges where $\|\mathrm{biw}\| = m_{\rm max}$, this equals
$m_{\rm max}(m_{\rm max}-1)$ (matching Steps 168/169
alternative-orientation values).

\textbf{Step 177 (PR~\#2968)}: $\max - \min$ of canonical/complement
predicted $(\widehat S_{\rm tot})^2$ eigenvalues $= 2\cdot\|\mathrm{biw}\|$
\emph{unconditionally} in
\texttt{LatticeSystem/Quantum/SpinS/PredictedEigenvalueMaxMinDiff.lean}.
Companion of Step 170 (canonical $-$ complement $= 2\cdot\mathrm{biw.re}$)
using $\|\mathrm{biw}\|$ instead of signed $\mathrm{biw.re}$. Direct
from Step 175 $-$ Step 176 $= \|\mathrm{biw}\|\cdot(\|\mathrm{biw}\|+1) -
\|\mathrm{biw}\|\cdot(\|\mathrm{biw}\|-1) = 2\cdot\|\mathrm{biw}\|$ via
\texttt{ring}. The max-min gap is non-negative since $\|\mathrm{biw}\| \ge 0$.

\textbf{Step 178 (PR~\#2969)}: iff $\|\mathrm{biw}\| = m_{\rm max}
\Leftrightarrow$ saturated edge at $N \ge 1$ in
\texttt{LatticeSystem/Quantum/SpinS/BipartiteImbalanceWeightNormEqMMaxIff.lean}.
The imbalance norm hits its maximum exactly at saturated edges
($|A| = 0 \vee |\neg A| = 0$). Backward via Step 90/91; forward via
contrapositive of Step 144 (strict $\|\mathrm{biw}\| < m_{\rm max}$ at
non-degenerate).

\textbf{Step 179 (PR~\#2970)}: unified max iff
$\max(\mathrm{canonical}, \mathrm{complement}) = m_{\rm max}(m_{\rm max}+1)
\Leftrightarrow$ saturated edge at $N \ge 1$ in
\texttt{LatticeSystem/Quantum/SpinS/PredictedEigenvalueMaxEqMMaxQuadIff.lean}.
The \emph{physical} predicted (max over orientations) $(\widehat S_{\rm tot})^2$
eigenvalue attains the fully-polarised ferromagnetic value iff at a
saturated edge. Direct from Step 175 ($\max = \|\mathrm{biw}\|\cdot(\|\mathrm{biw}\|+1)$)
$+$ Step 178 ($\|\mathrm{biw}\| = m_{\rm max} \Leftrightarrow$ saturated)
using polynomial factoring. Unifies Steps 158/166 via the physical
$\max$ quantity.

\textbf{Step 180 (PR~\#2971)}: unified max iff
$\max(\mathrm{canonical}, \mathrm{complement}) > 0 \Leftrightarrow$
unbalanced at $N \ge 1$ in
\texttt{LatticeSystem/Quantum/SpinS/PredictedEigenvalueMaxPosIffUnbalanced.lean}.
The \emph{physical} predicted $(\widehat S_{\rm tot})^2$ eigenvalue is
strictly positive iff $|A| \ne |\neg A|$. Direct from Step 175
($\max = \|\mathrm{biw}\|\cdot(\|\mathrm{biw}\|+1)$) + Step 81
($\|\mathrm{biw}\| = 0 \Leftrightarrow$ balanced at $N \ge 1$). Unifies
Step 147 (canonical) and Step 164 (complement) via the
physical-quantity iff.

\textbf{Step 181 (PR~\#2972)}: range Icc for the physical predicted
$(\widehat S_{\rm tot})^2$ eigenvalue:
$\max(\mathrm{canonical}, \mathrm{complement}) \in [0, m_{\rm max}(m_{\rm max}+1)]$
\emph{unconditionally} in
\texttt{LatticeSystem/Quantum/SpinS/PredictedEigenvalueMaxMemIcc.lean}.
Direct from Step 175 ($\max = \|\mathrm{biw}\|\cdot(\|\mathrm{biw}\|+1)$)
+ $0 \le \|\mathrm{biw}\| \le m_{\rm max}$ + monotonicity. Unifies the
orientation-specialised range Icc (Steps 160/162) via the physical
$\max$ quantity.

\textbf{Step 182 (PR~\#2973)}: lower bound
$\min(\mathrm{canonical}, \mathrm{complement}) \ge -1/4$ \emph{unconditionally}
in
\texttt{LatticeSystem/Quantum/SpinS/PredictedEigenvalueMinGeNegOneFourth.lean}.
The parabola $x(x-1)$ has minimum $-1/4$ at $x = 1/2$. Direct from
Step 176 (min $= \|\mathrm{biw}\|(\|\mathrm{biw}\|-1)$) + $(x - 1/2)^2 \ge 0$.

\textbf{Step 183 (PR~\#2974)}: unified max iff
$\max(\mathrm{canonical}, \mathrm{complement}) = 0 \Leftrightarrow$ balanced
at $N \ge 1$ in
\texttt{LatticeSystem/Quantum/SpinS/PredictedEigenvalueMaxEqZeroIffBalanced.lean}.
The \emph{physical} predicted $(\widehat S_{\rm tot})^2$ eigenvalue vanishes
(singlet $S_{\rm tot} = 0$) exactly at balanced sublattices. Complements
Step 180 ($\max > 0 \Leftrightarrow$ unbalanced). Unifies Steps 146/163
(orientation-specialised balanced iffs) via the physical $\max$ quantity.

\textbf{Step 184 (PR~\#2975)}: unified strict max bound
$\max(\mathrm{canonical}, \mathrm{complement}) < m_{\rm max}(m_{\rm max}+1)$
at non-saturated ($|A| \ge 1$, $|\neg A| \ge 1$, $N \ge 1$) in
\texttt{LatticeSystem/Quantum/SpinS/PredictedEigenvalueMaxLtMMaxQuad.lean}.
Contrapositive of Step 179 combined with Step 181 upper bound. The
physical predicted $(\widehat S_{\rm tot})^2$ eigenvalue is strictly
below the fully-polarised ferromagnetic value at non-saturated. Unifies
Steps 155/156 (orientation-specialised strict bounds) via the physical
$\max$ quantity.

\textbf{Step 185 (PR~\#2976)}: $\max + \min$ of canonical/complement
predicted $(\widehat S_{\rm tot})^2$ eigenvalues $= 2\cdot\|\mathrm{biw}\|^2$
\emph{unconditionally} in
\texttt{LatticeSystem/Quantum/SpinS/PredictedEigenvalueMaxAddMin.lean}.
Companion of Step 177 (max $-$ min $= 2\cdot\|\mathrm{biw}\|$). Direct
from Step 175 $+$ Step 176 via \texttt{ring}.

\textbf{Step 186 (PR~\#2977)}: iff $\max - \min > 0 \Leftrightarrow$
unbalanced at $N \ge 1$ in
\texttt{LatticeSystem/Quantum/SpinS/PredictedEigenvalueMaxSubMinPosIffUnbalanced.lean}.
The canonical-complement eigenvalue gap (proportional to $\|\mathrm{biw}\|$)
vanishes iff balanced. Direct from Step 177 (gap $= 2\cdot\|\mathrm{biw}\|$)
+ Step 81 ($\|\mathrm{biw}\| = 0 \Leftrightarrow$ balanced). Companion
of Step 180 ($\max > 0 \Leftrightarrow$ unbalanced).

\textbf{Step 187 (PR~\#2978)}: $|\mathrm{canonical} - \mathrm{complement}|
= 2\cdot\|\mathrm{biw}\|$ \emph{unconditionally} in
\texttt{LatticeSystem/Quantum/SpinS/PredictedEigenvalueAbsDiff.lean}.
Unsigned version of Step 170. Direct from
$|2\cdot\mathrm{biw.re}| = 2\cdot|\mathrm{biw.re}| = 2\cdot\|\mathrm{biw}\|$
(using $\mathrm{biw.im} = 0$). Equivalent to Step 177 ($\max - \min$)
since $\max - \min = |\mathrm{canonical} - \mathrm{complement}|$.

\textbf{Step 188 (PR~\#2979)}: $\max(\mathrm{canonical}, \mathrm{complement})
= \|\mathrm{biw}\|^2 + \|\mathrm{biw}\|$ \emph{unconditionally} in
\texttt{LatticeSystem/Quantum/SpinS/PredictedEigenvalueMaxEqNormSqAdd.lean}.
Expanded polynomial form of Step 175. Convenient for \texttt{nlinarith}
and arithmetic.

\textbf{Step 189 (PR~\#2980)}: $\min(\mathrm{canonical}, \mathrm{complement})
= \|\mathrm{biw}\|^2 - \|\mathrm{biw}\|$ \emph{unconditionally} in
\texttt{LatticeSystem/Quantum/SpinS/PredictedEigenvalueMinEqNormSqSub.lean}.
Expanded polynomial form of Step 176. Companion of Step 188 — together
makes the affine parameterisation
$\max = \|\mathrm{biw}\|^2 + \|\mathrm{biw}\|$,
$\min = \|\mathrm{biw}\|^2 - \|\mathrm{biw}\|$ explicit.

\textbf{Step 190 (PR~\#2981)}: $\max(\mathrm{canonical}, \mathrm{complement})
= \mathrm{canonical}$ at proper orientation $|\neg A| \le |A|$ in
\texttt{LatticeSystem/Quantum/SpinS/PredictedEigenvalueMaxEqCanonical.lean}.
The canonical orientation gives the physical predicted eigenvalue
exactly at proper orientation. Direct from Step 173 + Step 174 +
$\mathrm{biw.re} \ge 0$ at proper orientation.

\textbf{Step 191 (PR~\#2983)}: iff
$\max(\mathrm{canonical}, \mathrm{complement}) < m_{\rm max}(m_{\rm max}+1)
\Leftrightarrow$ non-saturated at $N \ge 1$ in
\texttt{LatticeSystem/Quantum/SpinS/PredictedEigenvalueMaxLtMMaxQuadIff.lean}.
Iff form of Step 184. Direct combine of Step 179
($\max = m_{\rm max}(m_{\rm max}+1) \Leftrightarrow$ saturated) +
Step 181 ($\max \le m_{\rm max}(m_{\rm max}+1)$).

\textbf{Step 192 (PR~\#2984)}: $\min(\mathrm{canonical}, \mathrm{complement})
= \mathrm{complement}$ at proper orientation $|\neg A| \le |A|$ in
\texttt{LatticeSystem/Quantum/SpinS/PredictedEigenvalueMinEqComplement.lean}.
Companion of Step 190 ($\max = \mathrm{canonical}$ at proper
orientation). Direct from Step 173 + Step 174 + $\mathrm{biw.re} \ge 0$
at proper orientation.

\textbf{Step 193 (PR~\#2985)}: $\max(\mathrm{canonical}, \mathrm{complement})
= \mathrm{complement}$ at complement orientation $|A| \le |\neg A|$ in
\texttt{LatticeSystem/Quantum/SpinS/PredictedEigenvalueMaxEqComplement.lean}.
Mirror of Step 190. The complement orientation gives the physical
predicted eigenvalue at complement orientation. Direct from Step 173 +
Step 174 + $\mathrm{biw.re} \le 0$ at complement orientation.

\textbf{Step 194 (PR~\#2986)}: $\min(\mathrm{canonical}, \mathrm{complement})
= \mathrm{canonical}$ at complement orientation $|A| \le |\neg A|$ in
\texttt{LatticeSystem/Quantum/SpinS/PredictedEigenvalueMinEqCanonical.lean}.
Companion of Step 193 ($\max = \mathrm{complement}$ at complement
orientation). Direct from Step 173 + Step 174 + $\mathrm{biw.re} \le 0$
at complement orientation.

\textbf{Step 195 (PR~\#2987)}:
$\min(\mathrm{canonical}, \mathrm{complement}) \le
\max(\mathrm{canonical}, \mathrm{complement})$ \emph{unconditionally}
in
\texttt{LatticeSystem/Quantum/SpinS/PredictedEigenvalueMinLeMax.lean}.
Trivial corollary of \texttt{min\_le\_max}; convenient for combining
with other inequalities.

\textbf{Step 196 (PR~\#2988)}: iff $\|\mathrm{biw}\| > 0 \Leftrightarrow$
unbalanced at $N \ge 1$ in
\texttt{LatticeSystem/Quantum/SpinS/BipartiteImbalanceWeightNormPosIff.lean}.
Iff form combining $\|\mathrm{biw}\| \ge 0$ and Step 81 ($\|\mathrm{biw}\|
= 0 \Leftrightarrow$ balanced). Convenient for switching between strict
positivity and balanced/unbalanced predicates.

\textbf{Step 197 (PR~\#2989)}: $\mathrm{biw.re} \le 0$ at complement
orientation $|A| \le |\neg A|$ in
\texttt{LatticeSystem/Quantum/SpinS/BipartiteImbalanceWeightReNonpos.lean}.
Mirror of the existing `\_re\_nonneg\_of\_cardA\_ge\_cardNotA` lemma.
Direct from $\mathrm{biw.re}$-eq + orientation + $N \ge 0$.

\textbf{Step 198 (PR~\#2990)}: $\mathrm{canonical} \le
\max(\mathrm{canonical}, \mathrm{complement})$ \emph{unconditionally}
in
\texttt{LatticeSystem/Quantum/SpinS/PredictedEigenvalueCanonicalLeMax.lean}.
Trivial corollary of \texttt{le\_max\_left}. Useful for chaining
inequalities.

\textbf{Step 199 (PR~\#2991)}: $\mathrm{complement} \le
\max(\mathrm{canonical}, \mathrm{complement})$ \emph{unconditionally}
in
\texttt{LatticeSystem/Quantum/SpinS/PredictedEigenvalueComplementLeMax.lean}.
Companion of Step 198. Trivial corollary of \texttt{le\_max\_right}.

\textbf{Step 200 (PR~\#2992)}:
$\min(\mathrm{canonical}, \mathrm{complement}) \le \mathrm{canonical}$
\emph{unconditionally} in
\texttt{LatticeSystem/Quantum/SpinS/PredictedEigenvalueMinLeCanonical.lean}.
Trivial corollary of \texttt{min\_le\_left}.

\textbf{Step 201 (PR~\#2993)}:
$\min(\mathrm{canonical}, \mathrm{complement}) \le \mathrm{complement}$
\emph{unconditionally} in
\texttt{LatticeSystem/Quantum/SpinS/PredictedEigenvalueMinLeComplement.lean}.
Companion of Step 200. Trivial corollary of \texttt{min\_le\_right}.

\textbf{Step 202 (PR~\#2994)}: predicted min energy decomposition
$(bipartiteToyMinEnergyPredicted\ A\ N).\mathrm{re} =
\mathrm{canonical\_eigenvalue.re} -
\mathrm{Re}(s_A\cdot(s_A+1)) - \mathrm{Re}(s_B\cdot(s_B+1))$ in
\texttt{LatticeSystem/Quantum/SpinS/BipartiteToyMinEnergyPredictedReViaEigenvalue.lean}.
Direct from the definition via \texttt{Complex.sub\_re}. Makes
explicit the relationship between predicted min energy of $H_{\rm toy}$
and the canonical predicted $(\widehat S_{\rm tot})^2$ eigenvalue.

\textbf{Step 203 (PR~\#2995)}: predicted min energy via max (physical)
eigenvalue at proper orientation: at $|\neg A| \le |A|$,
$(bipartiteToyMinEnergyPredicted\ A\ N).\mathrm{re} =
\max(\mathrm{canonical}, \mathrm{complement}) -
\mathrm{Re}(s_A\cdot(s_A+1)) - \mathrm{Re}(s_B\cdot(s_B+1))$ in
\texttt{LatticeSystem/Quantum/SpinS/BipartiteToyMinEnergyPredictedReAtProperOrientation.lean}.
Combines Step 202 + Step 190 (proper orientation $\max = \mathrm{canonical}$).
Uses the physical predicted $(\widehat S_{\rm tot})^2$ eigenvalue.

\textbf{Step 204 (PR~\#2996)}: predicted min energy via min
(wrong-orientation) eigenvalue at complement orientation: at
$|A| \le |\neg A|$,
$(bipartiteToyMinEnergyPredicted\ A\ N).\mathrm{re} =
\min(\mathrm{canonical}, \mathrm{complement}) -
\mathrm{Re}(s_A\cdot(s_A+1)) - \mathrm{Re}(s_B\cdot(s_B+1))$ in
\texttt{LatticeSystem/Quantum/SpinS/BipartiteToyMinEnergyPredictedReAtComplementOrientation.lean}.
Companion of Step 203. Combines Step 202 + Step 194. At complement
orientation, the canonical-formula uses the smaller $\min$ eigenvalue,
underestimating the physical min energy.

\textbf{Step 205 (PR~\#2997)}: complement-orientation sum identity
$bipartiteToyMinEnergyPredicted\ A\ N + bipartiteToyMinEnergyPredicted\ (\neg A)\ N
= -|A|\cdot|\neg A|\cdot N^2 - |\Lambda|\cdot N$ in
\texttt{LatticeSystem/Quantum/SpinS/BipartiteToyMinEnergyPredictedAddComplement.lean}.
Companion of the existing `\_sub\_complement` gap identity. Together
gives the affine sandwich: $predicted\_min\ A = -|A|\cdot|\neg A|\cdot N^2/2
- |\neg A|\cdot N$ and $predicted\_min\ (\neg A) = -|A|\cdot|\neg A|\cdot N^2/2
- |A|\cdot N$.

\textbf{Step 206 (PR~\#2998)}: iff $predicted\_min\_energy\ A =
predicted\_min\_energy\ (\neg A) \Leftrightarrow$ balanced at $N \ge 1$
in
\texttt{LatticeSystem/Quantum/SpinS/BipartiteToyMinEnergyPredictedEqComplementIffBalanced.lean}.
The two orientations give the same predicted min energy \emph{iff}
balanced sublattices. Direct from the existing `\_sub\_complement`
gap identity $= (|A| - |\neg A|)\cdot N$ + $N \ge 1$ cancellation.

\textbf{Step 207 (PR~\#2999)}: $(predicted\_min\ A).\mathrm{re} \le
(predictedSymm\ A).\mathrm{re}$ \emph{unconditionally} in
\texttt{LatticeSystem/Quantum/SpinS/BipartiteToyMinEnergyPredictedReLeSymmRe.lean}.
The canonical-orientation predicted min energy is bounded above by
the symmetric (min-aware) form. Difference $= (|\neg A| -
\min(|A|, |\neg A|))\cdot N \ge 0$. Equality at $|\neg A| \le |A|$;
strict at $|A| < |\neg A|$ and $N \ge 1$.

\textbf{Step 208 (PR~\#3000)}: $(predicted\_min\ \neg A).\mathrm{re} \le
(predictedSymm\ A).\mathrm{re}$ \emph{unconditionally} in
\texttt{LatticeSystem/Quantum/SpinS/BipartiteToyMinEnergyPredictedComplementReLeSymmRe.lean}.
Mirror of Step 207. The complement-orientation predicted min energy
is bounded above by the symmetric form. Difference $= (|A| -
\min(|A|, |\neg A|))\cdot N \ge 0$. Equality at $|A| \le |\neg A|$;
strict at $|\neg A| < |A|$ and $N \ge 1$. Together with Step 207
confirms $predictedSymm$ as upper bound on both orientations.

\textbf{Step 209 (PR~\#3001)}: $\max((predicted\_min\ A).\mathrm{re},
(predicted\_min\ \neg A).\mathrm{re}) = (predictedSymm\ A).\mathrm{re}$
\emph{unconditionally} in
\texttt{LatticeSystem/Quantum/SpinS/BipartiteToyMinEnergyPredictedMaxEqPredictedSymm.lean}.
Identifies $bipartiteToyMinEnergyPredictedSymm$ as the \emph{max} of
orientation-specific predicted min energies. Direct from Steps 207/208
(upper bounds) plus equality at the orientation matching
$\min(|A|, |\neg A|)$. The symmetric form is the physical Tasaki §2.5
Theorem 2.3 prediction.

\textbf{Step 210 (PR~\#3003)}: at balanced
$(predicted\_min\ A).\mathrm{re} = (predictedSymm\ A).\mathrm{re}$ in
\texttt{LatticeSystem/Quantum/SpinS/BipartiteToyMinEnergyPredictedReEqSymmAtBalanced.lean}.
At $|A| = |\neg A|$, the canonical orientation's predicted min energy
equals the symmetric form. Direct from the simplified forms + balanced
cardinality.

\textbf{Step 211 (PR~\#3004)}: at balanced
$(predicted\_min\ \neg A).\mathrm{re} = (predictedSymm\ A).\mathrm{re}$
in
\texttt{LatticeSystem/Quantum/SpinS/BipartiteToyMinEnergyPredictedComplementReEqSymmAtBalanced.lean}.
Mirror of Step 210. Combined: at balanced, both orientations coincide
and equal the symmetric form.

\textbf{Step 212 (PR~\#3005)}: at canonical-saturated $|\neg A| = 0$,
$(predicted\_min\ A).\mathrm{re} = (predictedSymm\ A).\mathrm{re} = 0$
in
\texttt{LatticeSystem/Quantum/SpinS/BipartiteToyMinEnergyPredictedReEqSymmAtSaturated.lean}.
Both formulas vanish at the fully-polarised saturated edge.

\textbf{Step 213 (PR~\#3006)}: at complement-saturated $|A| = 0$,
$(predicted\_min\ \neg A).\mathrm{re} = (predictedSymm\ A).\mathrm{re} = 0$
in
\texttt{LatticeSystem/Quantum/SpinS/BipartiteToyMinEnergyPredictedComplementReEqSymmAtCardAZero.lean}.
Mirror of Step 212. The complement-orientation predicted min energy
vanishes at complement-saturated.

\textbf{Step 214 (PR~\#3007)}: iff strict
$(predicted\_min\ A).\mathrm{re} < (predictedSymm\ A).\mathrm{re}
\Leftrightarrow |A| < |\neg A|$ at $N \ge 1$ in
\texttt{LatticeSystem/Quantum/SpinS/BipartiteToyMinEnergyPredictedLtSymmIff.lean}.
Companion strict version of Step 207. Via case-split on
$|A| < |\neg A|$ vs $|\neg A| \le |A|$.

\textbf{Step 215 (PR~\#3008)}: iff strict (complement)
$(predicted\_min\ \neg A).\mathrm{re} < (predictedSymm\ A).\mathrm{re}
\Leftrightarrow |\neg A| < |A|$ at $N \ge 1$ in
\texttt{LatticeSystem/Quantum/SpinS/BipartiteToyMinEnergyPredictedComplementLtSymmIff.lean}.
Mirror of Step 214. Companion strict version of Step 208.

\textbf{Step 216 (PR~\#3009)}: iff equality
$(predicted\_min\ A).\mathrm{re} = (predictedSymm\ A).\mathrm{re}
\Leftrightarrow |\neg A| \le |A|$ at $N \ge 1$ in
\texttt{LatticeSystem/Quantum/SpinS/BipartiteToyMinEnergyPredictedEqSymmIff.lean}.
Trichotomy companion of Step 214. Direct from Step 207 ($\le$) +
Step 214 ($<$ iff).

\textbf{Step 217 (PR~\#3010)}: iff equality (complement)
$(predicted\_min\ \neg A).\mathrm{re} = (predictedSymm\ A).\mathrm{re}
\Leftrightarrow |A| \le |\neg A|$ at $N \ge 1$ in
\texttt{LatticeSystem/Quantum/SpinS/BipartiteToyMinEnergyPredictedComplementEqSymmIff.lean}.
Mirror of Step 216. Direct from Step 208 ($\le$) + Step 215 ($<$ iff).

\textbf{Step 218 (PR~\#3011)}: min of orientation-specific predicted
min energies $= -|A|\cdot|\neg A|\cdot N^2/2 - \max(|A|, |\neg A|)\cdot N$
\emph{unconditionally} in
\texttt{LatticeSystem/Quantum/SpinS/BipartiteToyMinEnergyPredictedMinComplementRe.lean}.
Companion of Step 209 ($\max = predictedSymm$). The $\min$ is the
"unphysical" lower predicted min energy, strictly below $predictedSymm$
at unbalanced.

\textbf{Steps 219--261 (PRs~\#3012--\#3054)}: orientation-sandwich
algebraic closure + operator-algebraic bridges.

PRs~\#3012--\#3050 systematically complete the algebraic framework
for the four orientation-pair points
\begin{equation*}
  \min \;\le\; \mathrm{avg} \;\le\; predictedSymm = \max,
\end{equation*}
all linked by the half-spread $\frac{1}{2}\bigl||A| - |\neg A|\bigr|\,N$.
The framework establishes:
\begin{itemize}
  \item the explicit sum identity $\max + \min = -|A|\cdot|\neg A|\cdot N^2 - |\Lambda|\cdot N$ (PR~\#3031);
  \item the spread formula $\max - \min = \bigl||A| - |\neg A|\bigr|\,N$ (PR~\#3012)
        together with its sign trichotomy at $N \ge 1$ (PRs~\#3013/\#3014);
  \item the half-spread separations $predictedSymm - \mathrm{avg} = \mathrm{avg} - \min
        = \frac{1}{2}\bigl||A| - |\neg A|\bigr|\,N$ (PRs~\#3034/\#3039);
  \item the three endpoint formulas in terms of midpoint $\pm$ half-spread
        (PRs~\#3032/\#3043/\#3045);
  \item equality and strict iffs at balanced / unbalanced for every pair
        of sandwich points (PRs~\#3037/\#3038/\#3040/\#3041/\#3046--\#3049);
  \item non-positivity of $\min$, $\max$ and avg, plus strict negativity
        of avg at non-degenerate (PRs~\#3029/\#3030/\#3042/\#3050).
\end{itemize}

PRs~\#3051--\#3054 add four operator-algebraic bridges connecting the
Néel-state toy-Hamiltonian expectation $\langle \Phi_{\rm N\acute eel}|
\hat H_{\rm toy}|\Phi_{\rm N\acute eel}\rangle = -|A|\cdot|\neg A|\cdot
N^2/2$ to the algebraic orientation-pair:
\begin{align*}
  \langle \Phi_{\rm N\acute eel}|\hat H_{\rm toy}|\Phi_{\rm N\acute eel}\rangle.\mathrm{re} - \mathrm{avg} &= \tfrac{|\Lambda|\cdot N}{2}, \\
  \langle \Phi_{\rm N\acute eel}|\hat H_{\rm toy}|\Phi_{\rm N\acute eel}\rangle.\mathrm{re} - \min &= \max(|A|, |\neg A|)\cdot N, \\
  \langle \Phi_{\rm N\acute eel}|\hat H_{\rm toy}|\Phi_{\rm N\acute eel}\rangle.\mathrm{re} - (predicted\_min\ A).\mathrm{re} &= |\neg A|\cdot N, \\
  \langle \Phi_{\rm N\acute eel}|\hat H_{\rm toy}|\Phi_{\rm N\acute eel}\rangle.\mathrm{re} - (predicted\_min\ \neg A).\mathrm{re} &= |A|\cdot N.
\end{align*}
Together with the existing gap to $predictedSymm$ (Tasaki §2.5
Theorem 2.3 prediction = $\min(|A|, |\neg A|)\cdot N$, PR~\#3000-area),
these bridges pin down the Néel-state expectation against every
sandwich point. Individual files are in
\texttt{LatticeSystem/Quantum/SpinS/BipartiteToyMinEnergyPredicted*.lean}
and
\texttt{LatticeSystem/Quantum/SpinS/NeelStateOfSExpectation*.lean};
see \texttt{docs/index.md} for the full list with per-PR provenance.

\textbf{Final statement (PR \#3337, consolidated in PRs \#3917/\#3922/\#3919)}: the Tasaki \S2.5 Theorem 2.3 conclusion is packaged as a \texttt{Prop}-valued definition \texttt{tasaki\_2\_5\_theorem\_2\_3} in \texttt{LatticeSystem/Quantum/SpinS/Theorem23SectorExistence.lean}. The detailed module-split structure described in earlier versions of this guide was massively consolidated by the Theorem 2.3 capstone consolidation refactor (Issue \#3920): PR \#3917 removed the two vacuous-at-\(N=1\) capstones \texttt{tasaki\_2\_5\_theorem\_2\_3\_bipartiteToy} and \texttt{tasaki\_2\_5\_theorem\_2\_3\_of\_bipartiteCompletePositive}; PR \#3922 pruned the \texttt{h\_intermediate}-threaded middle chain (the original sector-existence / dominance / predicted-Casimir / common-energy chain modules); PR \#3919 swept 525 newly-orphan modules across the whole \texttt{LatticeSystem/} tree (\textasciitilde 35{,}000 lines). The current sound proof witness for any \(N \ge 1\) is the structural variant \texttt{tasaki\_2\_5\_theorem\_2\_3\_structural} (PR \#3891 toy coupling, PR \#3893 general bipartite \(J\)). See \texttt{LatticeSystem.lean} for the current canonical import graph.

\textbf{Conditional final wrapper}: in the canonical orientation
$|\neg A| \le |A|$, Lean now proves
\texttt{tasaki\_2\_5\_theorem\_2\_3\_of\_left\_endpoint\_predictedGS\_of\_onA\_neg\_lt\_offA}.
It first invokes the predicted-GS left-endpoint interval-chain theorem
to obtain one common energy $\mu$ across all admissible sectors, then
uses the per-sector Theorem 2.2 wrapper to replace each interval-chain
witness by a sector ground-state witness carrying the final
within-sector uniqueness clause.  The statement keeps exactly the
remaining mathematical inputs explicit: predicted-GS membership for
the source vector, the lowered off-$A$ dominance estimate, and global
minimality of the common value.  Once these callbacks are supplied, the
wrapper proves the final \texttt{tasaki\_2\_5\_theorem\_2\_3} proposition
for Tasaki, \emph{Physics and Mathematics of Quantum Many-Body Systems},
Springer 2020, \S2.5 Theorem 2.3, p.~42.

\textbf{Admissible-sector interval combinatorics}: the sector range is
now formalised as a closed finite interval.  Lean proves the membership
characterisation
\[
  M \in \mathrm{tasaki23GroundStateSectors}\,A\,N
  \quad\Longleftrightarrow\quad
  \min(|A|,|\neg A|)N \le M \le \max(|A|,|\neg A|)N
\]
as \texttt{tasaki23GroundStateSectors\_mem\_iff}.  The endpoint lemmas
\texttt{tasaki23GroundStateSectors\_left\_mem} and
\texttt{tasaki23GroundStateSectors\_right\_mem} put the two boundary
sectors into the range, while
\texttt{tasaki23GroundStateSectors\_succ\_mem\_of\_mem\_of\_lt\_right}
and
\texttt{tasaki23GroundStateSectors\_pred\_mem\_of\_left\_lt\_of\_mem}
record the one-step successor/predecessor closure inside the interval.
Finally,
\[
  |\mathrm{tasaki23GroundStateSectors}\,A\,N|
  =
  \mathrm{tasaki23PredictedDegeneracy}\,A\,N
  =
  \bigl||A|-|\neg A|\bigr|\,N+1
\]
is formalised as \texttt{tasaki23GroundStateSectors\_card}.  These
lemmas now live in
\texttt{LatticeSystem/Quantum/SpinS/Theorem23Sectors.lean} and provide
the finite-range bookkeeping used to iterate the lowering and raising
adjacent-sector energy comparisons across the Theorem 2.3 multiplet.

\textbf{Ladder eigenvalue preservation}: the next proof step packages
the SU(2) ladder comparison used to identify neighbouring sector
eigenvalues.  If
\[
  \hat H\Psi = \mu\Psi,
\]
then the existing commutators
$[\hat H,\hat S^-_{\rm tot}]=0$ and
$[\hat H,\hat S^+_{\rm tot}]=0$ give
\[
  \hat H(\hat S^-_{\rm tot}\Psi) = \mu\,\hat S^-_{\rm tot}\Psi,
  \qquad
  \hat H(\hat S^+_{\rm tot}\Psi) = \mu\,\hat S^+_{\rm tot}\Psi .
\]
Lean names:
\texttt{heisenbergHamiltonianS\_mulVec\_totalSpinSOpMinus\_of\_eigenvec}
and
\texttt{heisenbergHamiltonianS\_mulVec\_totalSpinSOpPlus\_of\_eigenvec}.
These are the one-step sector-linking identities for the Theorem 2.3
multiplet; the remaining proof work is the non-vanishing and
Marshall-sign compatibility needed to feed the laddered vector into
the per-sector uniqueness theorem.

\textbf{Ladder-kernel Casimir consequences}: the complementary
SU(2) obstruction is now recorded.  Let \(\Psi\) lie in the
\(\hat S^{(3)}_{\rm tot}\)-eigenspace with eigenvalue \(m\).  The
total-spin Casimir rearrangements give
\[
  \hat S^+_{\rm tot}\hat S^-_{\rm tot}
  =
  \hat{\boldsymbol S}_{\rm tot}^2
  -(\hat S^{(3)}_{\rm tot})^2
  +\hat S^{(3)}_{\rm tot}
\]
and
\[
  \hat S^-_{\rm tot}\hat S^+_{\rm tot}
  =
  \hat{\boldsymbol S}_{\rm tot}^2
  -(\hat S^{(3)}_{\rm tot})^2
  -\hat S^{(3)}_{\rm tot}.
\]
Hence, if \(\hat S^-_{\rm tot}\Psi=0\), then
\[
  \hat{\boldsymbol S}_{\rm tot}^2\Psi=m(m-1)\Psi,
\]
while if \(\hat S^+_{\rm tot}\Psi=0\), then
\[
  \hat{\boldsymbol S}_{\rm tot}^2\Psi=m(m+1)\Psi.
\]
Lean names these two consequences
\[
\begin{gathered}
\texttt{tasaki23\_totalSpinSSquared\_mulVec\_of}\\
\texttt{\_totalSpinSOpMinus\_eq\_zero\_of\_mem\_magSubspaceS}
\end{gathered}
\]
and
\[
\begin{gathered}
\texttt{tasaki23\_totalSpinSSquared\_mulVec\_of}\\
\texttt{\_totalSpinSOpPlus\_eq\_zero\_of\_mem\_magSubspaceS}.
\end{gathered}
\]
These statements identify the total-spin representation forced by a
vanished adjacent-sector ladder move, and are the structural
alternative to treating the still-open local dominance inequality as a
stand-alone positivity fact.

\textbf{Casimir-based ladder non-vanishing}: the previous kernel
alternative is also packaged in the contrapositive form needed for the
sector interval argument.  Suppose \(\Psi\ne0\) lies in the
\(\hat S^{(3)}_{\rm tot}\)-eigenspace of value \(m\), and
\[
  \hat{\boldsymbol S}_{\rm tot}^2\Psi=\gamma\Psi.
\]
If \(\gamma\ne m(m-1)\), then
\[
  \hat S^-_{\rm tot}\Psi\ne0.
\]
Otherwise the lowering kernel theorem would force the same Casimir
action to equal \(m(m-1)\Psi\), so
\((\gamma-m(m-1))\Psi=0\), contradicting \(\Psi\ne0\) and
\(\gamma\ne m(m-1)\).  The raising companion is identical: if
\(\gamma\ne m(m+1)\), then
\[
  \hat S^+_{\rm tot}\Psi\ne0.
\]
Lean names these two non-vanishing criteria
\[
\begin{gathered}
\texttt{tasaki23\_totalSpinSOpMinus\_mulVec\_ne\_zero}\\
\texttt{\_of\_casimir\_ne\_kernel\_value}
\end{gathered}
\]
and
\[
\begin{gathered}
\texttt{tasaki23\_totalSpinSOpPlus\_mulVec\_ne\_zero}\\
\texttt{\_of\_casimir\_ne\_kernel\_value}.
\end{gathered}
\]
They are the direct bridge from the SU(2) Casimir obstruction to the
claim that an adjacent-sector ladder move can vanish only at the
corresponding representation endpoint.  The Casimir obstruction and
contrapositive non-vanishing criteria are collected in
\texttt{LatticeSystem/Quantum/SpinS/Theorem23Casimir.lean}.
The same file also contains the sector-embedded wrappers
\[
\begin{gathered}
\texttt{tasaki23\_totalSpinSOpMinus\_mulVec\_magSectorEmbedding}\\
\texttt{\_ne\_zero\_of\_casimir\_ne\_kernel\_value},\\
\texttt{tasaki23\_totalSpinSOpPlus\_mulVec\_magSectorEmbedding}\\
\texttt{\_ne\_zero\_of\_casimir\_ne\_kernel\_value}.
\end{gathered}
\]
They combine the abstract Casimir obstruction with
\texttt{magSectorEmbedding\_mem\_magSubspaceS}, so a sector vector
\(\Phi\) can be used directly after zero extension to the full
Hilbert space.
The follow-up wrappers discharge the nonzero-vector hypothesis for the
actual Theorem 2.2 sector vector.  If \(v_\tau>0\) for every sector
configuration \(\tau\), then
\[
  \mathrm{magSectorEmbedding}
    \left(\tau\mapsto
      \bigl((\mathrm{marshallSignS}(A,\tau))_{\rm re} v_\tau\bigr)
    \right)\ne0,
\]
because the Marshall sign has real part \(\pm1\).  Lean records this as
\texttt{tasaki23\_marshallPositive\_magSectorEmbedding\_ne\_zero} and
then packages the lowering and raising Casimir criteria for this
Marshall-positive vector as
\[
\begin{gathered}
\texttt{tasaki23\_totalSpinSOpMinus\_mulVec\_marshallPositive}\\
\texttt{\_magSectorEmbedding\_ne\_zero\_of\_casimir\_ne\_kernel\_value},\\
\texttt{tasaki23\_totalSpinSOpPlus\_mulVec\_marshallPositive}\\
\texttt{\_magSectorEmbedding\_ne\_zero\_of\_casimir\_ne\_kernel\_value}.
\end{gathered}
\]
The adjacent-sector energy comparison now has a Casimir-nonvanishing
entry point.  For a Marshall-positive source vector
\[
  \Psi_M=
  \mathrm{magSectorEmbedding}
    \left(\tau\mapsto
      \bigl((\mathrm{marshallSignS}(A,\tau))_{\rm re} v_\tau\bigr)
    \right)
\]
with \(v_\tau>0\), suppose
\[
  \hat S_{\rm tot}^2\Psi_M=\gamma\Psi_M
  \qquad\text{and}\qquad
  \gamma\ne m(m-1),
  \quad m=|V|N/2-M .
\]
Then the lowering image \(\hat S^-_{\rm tot}\Psi_M\) is nonzero by
the Marshall-positive Casimir criterion.  If the same lowered vector
also satisfies the already isolated Marshall-positivity hypothesis in
sector \(M+1\), Theorem~2.2 uniqueness in that sector identifies the
successor ground-state energy with the source energy.  Lean records
this as
\[
\begin{gathered}
\texttt{tasaki23\_lowering\_identifies\_adjacent\_sector\_energy}\\
\texttt{\_with\_casimir\_nonzero}.
\end{gathered}
\]
The raising companion uses the endpoint value \(m(m+1)\) and is named
\[
\begin{gathered}
\texttt{tasaki23\_raising\_identifies\_adjacent\_sector\_energy}\\
\texttt{\_with\_casimir\_nonzero}.
\end{gathered}
\]
Finally,
\[
\begin{gathered}
\texttt{tasaki23\_successor\_sector\_common\_energy}\\
\texttt{\_of\_site\_sum\_pos\_of\_casimir\_ne\_kernel\_value}
\end{gathered}
\]
packages the lowering direction with the admissible-sector successor
lemma.  Thus the one-step chain can now carry both pieces of data
needed later: the successor sector belongs to the predicted interval,
and the ladder vector is nonzero for the SU(2) Casimir reason rather
than merely because strict target-sector positivity implies nonzero.
The raising-direction interval wrapper is now recorded as well.  If
the source sector \(M+1\) belongs to the admissible interval and is not
the left endpoint, then
\[
\begin{gathered}
\texttt{tasaki23\_predecessor\_sector\_common\_energy}\\
\texttt{\_of\_site\_sum\_pos}
\end{gathered}
\]
uses \texttt{tasaki23GroundStateSectors\_pred\_mem\_of\_left\_lt\_of\_mem}
and the already isolated raised-vector site-sum positivity package to
produce a Marshall-positive eigenvector in sector \(M\) at the same
energy.  Its Casimir endpoint-mismatch form,
\[
\begin{gathered}
\texttt{tasaki23\_predecessor\_sector\_common\_energy}\\
\texttt{\_of\_site\_sum\_pos\_of\_casimir\_ne\_kernel\_value},
\end{gathered}
\]
adds the hypothesis that the source Casimir value is not
\[
  \left(|V|N/2-(M+1)\right)
  \left(|V|N/2-(M+1)+1\right),
\]
so that the nonzero raised-vector conclusion comes from the same
SU(2) obstruction rather than from the target-sector positivity
hypothesis alone.

\textbf{Predicted-Casimir endpoint mismatch}: the predicted total
spin in Tasaki's theorem is now connected to the concrete Casimir
kernel values that occur in the adjacent-sector wrappers.  Lean names
\[
  S_{\rm pred}(A,N)
    = \bigl||A|-|\neg A|\bigr|\,N/2
\]
as \texttt{tasaki23PredictedTotalSpin} and records
\[
  \Gamma_{\rm pred}(A,N)
    = S_{\rm pred}(S_{\rm pred}+1)
\]
as \texttt{tasaki23PredictedCasimirValue}.  The half-width identity
\[
  S_{\rm pred}
  = \frac{\max(|A|,|\neg A|)N-\min(|A|,|\neg A|)N}{2}
\]
is formalised as
\texttt{tasaki23PredictedTotalSpin\_eq\_sector\_half\_width}; the
partition identity \(|A|+|\neg A|=|V|\) is recorded as
\texttt{tasaki23\_card\_filter\_A\_add\_card\_notA}.  If
\(M\) lies in the admissible interval and is not the right endpoint,
then \(m=|V|N/2-M\) satisfies \(-S_{\rm pred}<m\le S_{\rm pred}\).
Hence
\[
  m(m-1)<S_{\rm pred}(S_{\rm pred}+1),
\]
because
\[
  S_{\rm pred}(S_{\rm pred}+1)-m(m-1)
  =(S_{\rm pred}+m)(S_{\rm pred}-m+1)>0.
\]
Lean first records the corresponding real strict inequality as
\texttt{tasaki23\_lowering\_kernel\_value\_lt\_predictedCasimirValue\_of\_mem\_of\_lt\_right}
and then the complex non-equality as
\texttt{tasaki23\_predictedCasimirValue\_ne\_lowering\_kernel\_value\_of\_mem\_of\_lt\_right}.
The raising direction is dual: if the source sector \(M+1\) lies in
the admissible interval and is not the left endpoint, then
\(-S_{\rm pred}\le m<S_{\rm pred}\), and
\[
  S_{\rm pred}(S_{\rm pred}+1)-m(m+1)
  =(S_{\rm pred}-m)(S_{\rm pred}+m+1)>0.
\]
This yields the real strict inequality
\texttt{tasaki23\_raising\_kernel\_value\_lt\_predictedCasimirValue\_of\_mem\_of\_left\_lt}
and its complex non-equality form
\texttt{tasaki23\_predictedCasimirValue\_ne\_raising\_kernel\_value\_of\_mem\_of\_left\_lt}.
These endpoint mismatch lemmas are precisely the arithmetic input
needed to turn a predicted total-Casimir eigenvalue into the
\texttt{h\(\gamma\)\_ne} hypotheses used by the successor and
predecessor Casimir non-vanishing chain links.  Reference: H.~Tasaki,
\emph{Physics and Mathematics of Quantum Many-Body Systems}, Springer
2020, \S2.5 Theorem 2.3, p.~42.

\textbf{Sector-support bridge and lowering-sector shift}: the same PR
also records the support side of the adjacent-sector comparison.  First,
\texttt{totalSpinSOp3\_mulVec\_apply\_eq\_magEigenvalueS\_mul} gives the
pointwise diagonal action
\[
  (\hat S^{(3)}_{\rm tot}\Psi)(\sigma)
  = \mathrm{magEigenvalueS}(\sigma)\,\Psi(\sigma)
\]
for arbitrary full spin-$S$ vectors.  Hence
\texttt{magSectorEmbedding\_mem\_magSubspaceS} says that a zero-extended
\(\mathrm{magSumS}=M\) sector vector lies in the
\(\hat S^{(3)}_{\rm tot}\)-eigenspace with eigenvalue
\(|V|N/2 - M\), and
\texttt{magSubspaceS\_apply\_eq\_zero\_of\_magSumS\_ne} gives the converse
support statement.  The re-embedding bridge
\texttt{magSectorEmbedding\_magSectorRestriction\_of\_mem\_magSubspaceS}
then says that any full vector in this eigenspace is exactly the
zero-extension of its sector restriction, so later component arguments
can reuse the sector-embedding formulas for vectors obtained from
successor \texttt{magSubspaceS} support.  Combining these with the
existing commutator shift lemma for \(\hat S^-_{\rm tot}\) proves
\[
  \hat S^-_{\rm tot}\,
    \mathrm{magSectorEmbedding}(\Phi_M)
  \quad\text{is supported on}\quad
  \mathrm{magSumS}=M+1,
\]
formalised as
\texttt{totalSpinSOpMinus\_mulVec\_magSectorEmbedding\_supported\_succ}.
Together with the eigenvalue-preservation lemma above, this shows that
the lowered vector, when nonzero, is an eigenvector in the neighbouring
sector at the same Heisenberg eigenvalue.  This is the next direct step
toward identifying the sector energies in Tasaki §2.5 Theorem 2.3.

\textbf{Adjacent-sector energy identification, conditional lowering step}:
the next bridge isolates exactly how the lowered vector will be used once
its Marshall positivity is proved.  Suppose
\[
  \hat H\,\mathrm{magSectorEmbedding}(\Phi_M)
    = \mu\,\mathrm{magSectorEmbedding}(\Phi_M)
\]
and suppose the lowered vector
\(\hat S^-_{\rm tot}\mathrm{magSectorEmbedding}(\Phi_M)\) satisfies the
Marshall-positive hypothesis in the adjacent sector \(M+1\).  The previous
support lemma gives sector \(M+1\) support, and the ladder-commutation lemma
gives
\[
  \hat H\bigl(\hat S^-_{\rm tot}\mathrm{magSectorEmbedding}(\Phi_M)\bigr)
    =
  \mu\,\hat S^-_{\rm tot}\mathrm{magSectorEmbedding}(\Phi_M).
\]
Applying the full-Hilbert-space Theorem 2.2 uniqueness clause in sector
\(M+1\) then identifies the target sector ground-state eigenvalue with
\(\mu\).  This is formalised as
\texttt{tasaki23\_lowering\_identifies\_adjacent\_sector\_energy} in
\texttt{LatticeSystem/Quantum/SpinS/Theorem23Local.lean}.  The next
bookkeeping step is to separate the non-vanishing consequence of strict
Marshall positivity from the still-missing proof of that positivity.

\textbf{Lowered-vector non-vanishing from strict Marshall positivity}:
the next Lean step makes the non-vanishing part explicit.  If the
adjacent sector \(M+1\) is non-empty and the lowered vector
\[
  \Psi^- =
  \hat S^-_{\rm tot}\,\mathrm{magSectorEmbedding}(\Phi_M)
\]
satisfies
\[
  0 <
  (\mathrm{marshallSignS}\,A\,\tau).\mathrm{re}\,
  (\Psi^-(\tau)).\mathrm{re}
  \qquad
  \text{for every }\tau\in\mathrm{magConfigS}(V,N,M+1),
\]
then choosing any \(\tau\) in the non-empty sector shows
\(\Psi^-(\tau)\neq 0\), hence \(\Psi^-\neq 0\) as a full-space vector.
This is formalised as
\texttt{tasaki23\_lowered\_ne\_zero\_of\_marshall\_pos}.  The wrapper
\texttt{tasaki23\_lowering\_identifies\_adjacent\_sector\_energy\_with\_nonzero}
packages this non-zero conclusion together with the conditional
adjacent-sector energy identification above.  The remaining
critical-path input is now the proof that the lowered sector vector
indeed satisfies the Marshall-positive hypothesis.

\textbf{Lowered-vector site-sum positivity form}: the next Lean step
opens the total lowering operator at a target configuration:
\[
  \bigl(\hat S^-_{\rm tot}\,\mathrm{magSectorEmbedding}(\Phi_M)\bigr)(\tau)
  =
  \sum_{x\in V}
  \bigl(\hat S^-_x\,\mathrm{magSectorEmbedding}(\Phi_M)\bigr)(\tau).
\]
This is formalised as
\texttt{totalSpinSOpMinus\_mulVec\_magSectorEmbedding\_apply\_eq\_site\_sum}.
It converts the remaining Marshall-positivity obligation into the
single-site sum form
\[
  0 <
  (\mathrm{marshallSignS}\,A\,\tau).\mathrm{re}\,
  \sum_{x\in V}
  \mathrm{Re}\,
  \bigl(\hat S^-_x\,\mathrm{magSectorEmbedding}(\Phi_M)\bigr)(\tau).
\]
The bridge
\texttt{tasaki23\_lowered\_marshall\_pos\_of\_site\_sum\_pos} turns this
local site-sum hypothesis back into the global lowered-vector
Marshall-positivity hypothesis, and
\texttt{tasaki23\_lowering\_identifies\_adjacent\_sector\_energy\_of\_site\_sum\_pos}
feeds it directly into the non-vanishing plus adjacent-sector energy
identification package.  The remaining critical-path proof can now
focus on the predecessor configurations appearing in each single-site
lowering term.

\textbf{Single-site Marshall sign for lowering predecessors}: after the
site-sum expansion, each non-zero single-site lowering contribution
comes from a predecessor configuration \(\rho\) which agrees with the
target \(\tau\) away from one site \(x\) and satisfies
\[
  \rho_x + 1 = \tau_x .
\]
For such one-site differences, the Marshall-sign product factors as
\[
  \mathrm{marshallSignS}(A,\tau)\,
  \mathrm{marshallSignS}(A,\rho)
  =
  \begin{cases}
    (-1)^{\tau_x+\rho_x}, & A(x),\\
    1, & \neg A(x).
  \end{cases}
\]
The predecessor relation makes \(\tau_x+\rho_x\) odd, hence the product
is \(-1\) on \(A\)-sites and \(1\) off \(A\).  Since the Marshall signs
are real, the same identities hold for the products of their real
parts.  The Lean names are
\texttt{marshallSignS\_mul\_of\_agree\_off\_site},
\texttt{marshallSignS\_mul\_of\_agree\_off\_site\_A\_true\_lower},
\texttt{marshallSignS\_mul\_of\_agree\_off\_site\_A\_false\_lower}, and
their real-part variants
\texttt{marshallSignS\_re\_mul\_re\_of\_agree\_off\_site\_A\_true\_lower}
and
\texttt{marshallSignS\_re\_mul\_re\_of\_agree\_off\_site\_A\_false\_lower}.
These are the local sign facts needed before proving positivity of the
single-site lowering summands.

\textbf{Single-site lowering component formula}: the local lowering
summand has an explicit predecessor description.  Let
\(\tau\in\mathrm{magConfigS}(V,N,M+1)\) and fix \(x\in V\).  If
\(\tau_x=0\), then there is no source configuration whose value at
\(x\) is lowered into \(\tau_x\), and Lean proves
\[
  \bigl(\hat S^-_x\,\mathrm{magSectorEmbedding}(\Phi_M)\bigr)(\tau)=0.
\]
This is
\[
\begin{gathered}
\texttt{onSiteS\_spinSOpMinus\_mulVec\_magSectorEmbedding\_apply}\\
\texttt{\_eq\_zero\_of\_target\_zero}.
\end{gathered}
\]
If instead \(\tau_x>0\), define the predecessor
\[
  \rho = \tau[x \mapsto \tau_x-1].
\]
The private bookkeeping lemma
\[
\begin{gathered}
\texttt{magSumS\_single\_site}\\
\texttt{\_lowering\_predecessor}
\end{gathered}
\]
proves
\(\mathrm{magSumS}(\rho)=M\), and the matrix-vector product collapses to
the single non-zero column
\[
  \bigl(\hat S^-_x\,\mathrm{magSectorEmbedding}(\Phi_M)\bigr)(\tau)
  =
  (\hat S^-_N)_{\tau_x,\rho_x}\,\Phi_M(\rho).
\]
This is formalised as
\[
\begin{gathered}
\texttt{onSiteS\_spinSOpMinus\_mulVec\_magSectorEmbedding\_apply}\\
\texttt{\_single\_site\_pred}.
\end{gathered}
\]
Together with the preceding Marshall-sign predecessor identities, this
reduces the remaining lowered-vector positivity proof to signed sums of
strictly positive source-sector coefficients and positive local lowering
matrix entries.

\textbf{Local signed single-site contribution split}: Lean also records
the real-part form of the predecessor component formula.  Since every
\(\hat S^-_N\) matrix entry is real, the non-zero local contribution
satisfies
\[
  \mathrm{Re}\,
  \bigl(\hat S^-_x\,\mathrm{magSectorEmbedding}(\Phi_M)\bigr)(\tau)
  =
  \mathrm{Re}\,(\hat S^-_N)_{\tau_x,\rho_x}\,
  \mathrm{Re}\,\Phi_M(\rho).
\]
This is formalised as
\[
\begin{gathered}
\texttt{onSiteS\_spinSOpMinus\_mulVec\_magSectorEmbedding\_apply}\\
\texttt{\_single\_site\_pred\_re}.
\end{gathered}
\]
Combining this identity with
\texttt{spinSOpMinus\_apply\_re\_pos\_of\_lower}, source-sector
Marshall positivity, and the one-site Marshall sign product gives the
two local signed alternatives:
\[
  A(x)=\mathrm{false}
  \quad\Rightarrow\quad
  0 <
  \mathrm{Re}\,\mathrm{marshallSignS}(A,\tau)\,
  \mathrm{Re}\,
  \bigl(\hat S^-_x\,\mathrm{magSectorEmbedding}(\Phi_M)\bigr)(\tau),
\]
whereas
\[
  A(x)=\mathrm{true}
  \quad\Rightarrow\quad
  \mathrm{Re}\,\mathrm{marshallSignS}(A,\tau)\,
  \mathrm{Re}\,
  \bigl(\hat S^-_x\,\mathrm{magSectorEmbedding}(\Phi_M)\bigr)(\tau)
  < 0 .
\]
The Lean names are
\[
\begin{gathered}
\texttt{tasaki23\_signed\_single\_site\_lowering\_component}\\
\texttt{\_pos\_of\_A\_false}
\end{gathered}
\]
and
\[
\begin{gathered}
\texttt{tasaki23\_signed\_single\_site\_lowering\_component}\\
\texttt{\_neg\_of\_A\_true}.
\end{gathered}
\]
Thus the formal local analysis now exposes the sign obstruction that
the full site-sum positivity proof must control: lowered sites off
\(A\) contribute with the Marshall-positive sign, while lowered sites
on \(A\) contribute with the opposite sign.

\textbf{Boundary-inclusive lowered sign-sum bounds}: the local split is
now extended to the boundary case \(\tau_x=0\).  In that case the
single-site lowering contribution vanishes, while for \(\tau_x>0\) the
strict local alternatives above apply.  Hence Lean proves the weak
single-site bounds
\[
  A(x)=\mathrm{false}
  \quad\Rightarrow\quad
  0 \le
  \mathrm{Re}\,\mathrm{marshallSignS}(A,\tau)\,
  \mathrm{Re}\,
  \bigl(\hat S^-_x\,\mathrm{magSectorEmbedding}(\Phi_M)\bigr)(\tau),
\]
and
\[
  A(x)=\mathrm{true}
  \quad\Rightarrow\quad
  \mathrm{Re}\,\mathrm{marshallSignS}(A,\tau)\,
  \mathrm{Re}\,
  \bigl(\hat S^-_x\,\mathrm{magSectorEmbedding}(\Phi_M)\bigr)(\tau)
  \le 0 .
\]
The Lean names are
\[
\begin{gathered}
\texttt{tasaki23\_signed\_single\_site\_lowering\_component}\\
\texttt{\_nonneg\_of\_A\_false}
\end{gathered}
\]
and
\[
\begin{gathered}
\texttt{tasaki23\_signed\_single\_site\_lowering\_component}\\
\texttt{\_nonpos\_of\_A\_true}.
\end{gathered}
\]
Summing over the two filtered sublattices gives
\[
  0 \le
  \sum_{x:A(x)=\mathrm{false}}
  \mathrm{Re}\,\mathrm{marshallSignS}(A,\tau)\,
  \mathrm{Re}\,
  \bigl(\hat S^-_x\,\mathrm{magSectorEmbedding}(\Phi_M)\bigr)(\tau),
\]
and
\[
  \sum_{x:A(x)=\mathrm{true}}
  \mathrm{Re}\,\mathrm{marshallSignS}(A,\tau)\,
  \mathrm{Re}\,
  \bigl(\hat S^-_x\,\mathrm{magSectorEmbedding}(\Phi_M)\bigr)(\tau)
  \le 0 .
\]
These are formalised as
\texttt{tasaki23\_signed\_lowering\_offA\_sum\_nonneg} and
\texttt{tasaki23\_signed\_lowering\_onA\_sum\_nonpos}.  Lean also
records the strict off-\(A\) witness form
\[
\texttt{tasaki23\_signed\_lowering\_offA\_sum\_pos\_of\_exists\_lowerable\_offA}.
\]
It uses \texttt{Finset.sum\_pos'}: every off-\(A\) term is
non-negative, and any site \(x\notin A\) with \(\tau_x>0\) supplies a
strictly positive summand.  The aggregate version
\[
\begin{gathered}
\texttt{tasaki23\_exists\_lowerable\_offA\_of\_offA\_occupation\_pos},\\
\texttt{tasaki23\_signed\_lowering\_offA\_sum\_pos\_of\_offA\_occupation\_pos}
\end{gathered}
\]
extracts this witness from a positive off-\(A\) occupation sum.  These
bounds decompose the remaining strict site-sum positivity problem into
the non-negative off-\(A\) part and the non-positive on-\(A\) part,
without yet proving that the former
dominates the latter.

\textbf{Lowered site-sum dominance bridge}: Lean now names the signed
local contribution
\[
  T_x(\tau)=
  \mathrm{Re}\,\mathrm{marshallSignS}(A,\tau)\,
  \mathrm{Re}\,
  \bigl(\hat S^-_x\,\mathrm{magSectorEmbedding}(\Phi_M)\bigr)(\tau)
\]
as \texttt{tasaki23SignedLoweringSiteContribution} in
\texttt{LatticeSystem/Quantum/SpinS/Theorem23LocalCoefficient.lean}.
Since \(A(x)\) is
Boolean, the full signed site-sum splits exactly as
\[
  \mathrm{Re}\,\mathrm{marshallSignS}(A,\tau)
  \sum_{x\in V}
  \mathrm{Re}\,
  \bigl(\hat S^-_x\,\mathrm{magSectorEmbedding}(\Phi_M)\bigr)(\tau)
  =
  \sum_{x:A(x)=\mathrm{false}} T_x(\tau)
  +
  \sum_{x:A(x)=\mathrm{true}} T_x(\tau).
\]
This is
\texttt{tasaki23\_signed\_lowering\_site\_sum\_eq\_offA\_add\_onA}.
Consequently the pointwise dominance condition
\[
  -\sum_{x:A(x)=\mathrm{true}} T_x(\tau)
  <
  \sum_{x:A(x)=\mathrm{false}} T_x(\tau)
\]
implies the strict site-sum positivity hypothesis.  Lean records this
as
\texttt{tasaki23\_signed\_lowering\_site\_sum\_pos\_of\_onA\_neg\_lt\_offA}
and wraps it into lowered-vector Marshall positivity as
\texttt{tasaki23\_lowered\_marshall\_pos\_of\_onA\_neg\_lt\_offA}.  This
still leaves the dominance inequality itself as the remaining
substantive proof obligation.

\textbf{Dominance-form successor common-energy link}: the same
dominance condition is now wired directly into the adjacent-sector
successor package.  Lean records this as
\texttt{tasaki23\_successor\_sector\_common\_energy\_of\_onA\_neg\_lt\_offA}.
Given a source-sector Marshall-positive eigenvector in an admissible
sector \(M\), the interval hypothesis \(M < \max(|A|,|\neg A|)N\),
and the pointwise condition
\[
  -\sum_{x:A(x)=\mathrm{true}} T_x(\tau)
  <
  \sum_{x:A(x)=\mathrm{false}} T_x(\tau),
\]
the theorem first invokes the dominance-to-site-sum bridge above and
then applies
\texttt{tasaki23\_successor\_sector\_common\_energy\_of\_site\_sum\_pos}.
The conclusion is the same one-step common-energy result: \(M+1\) is
admissible, \(S^-_{\mathrm{tot}}\Psi_M\ne 0\), and the successor sector
has a Marshall-positive eigenvector at the same eigenvalue \(\mu\).
Thus the remaining local task is sharply isolated as the proof of the
displayed dominance inequality.

\textbf{Raising-direction local component package}: the same
adjacent-sector bookkeeping is now available for the raising ladder.
First, the support bridge
\[
\begin{gathered}
\texttt{totalSpinSOpPlus\_mulVec}\\
\texttt{\_magSectorEmbedding\_supported\_pred}
\end{gathered}
\]
states that if \(\Phi_{M+1}\) is embedded from the
\(\mathrm{magSumS}=M+1\) sector, then
\(\hat S^+_{\rm tot}\,\mathrm{magSectorEmbedding}(\Phi_{M+1})\) is
supported on the adjacent sector \(\mathrm{magSumS}=M\).  Componentwise,
Lean expands
\[
  \bigl(\hat S^+_{\rm tot}\,\mathrm{magSectorEmbedding}(\Phi_{M+1})\bigr)(\tau)
  =
  \sum_{x\in V}
  \bigl(\hat S^+_x\,\mathrm{magSectorEmbedding}(\Phi_{M+1})\bigr)(\tau)
\]
as
\[
\begin{gathered}
\texttt{totalSpinSOpPlus\_mulVec\_magSectorEmbedding\_apply}\\
\texttt{\_eq\_site\_sum}.
\end{gathered}
\]

For a target \(\tau\in\mathrm{magConfigS}(V,N,M)\), the local
raising summand vanishes at a saturated site \(\tau_x=N\):
\[
\begin{gathered}
\texttt{onSiteS\_spinSOpPlus\_mulVec\_magSectorEmbedding\_apply}\\
\texttt{\_eq\_zero\_of\_target\_top}.
\end{gathered}
\]
If \(\tau_x<N\), the unique source successor is
\[
  \sigma = \tau[x \mapsto \tau_x+1],
  \qquad
  \mathrm{magSumS}(\sigma)=M+1,
\]
and Lean proves
\[
  \bigl(\hat S^+_x\,\mathrm{magSectorEmbedding}(\Phi_{M+1})\bigr)(\tau)
  =
  (\hat S^+_N)_{\tau_x,\sigma_x}\,\Phi_{M+1}(\sigma).
\]
The theorem names are
\[
\begin{gathered}
\texttt{onSiteS\_spinSOpPlus\_mulVec\_magSectorEmbedding\_apply}\\
\texttt{\_single\_site\_succ}
\end{gathered}
\]
and its real-part form
\[
\begin{gathered}
\texttt{onSiteS\_spinSOpPlus\_mulVec\_magSectorEmbedding\_apply}\\
\texttt{\_single\_site\_succ\_re}.
\end{gathered}
\]
Since the raising matrix coefficient has positive real part on this
successor pair, source-sector Marshall positivity gives the same local
sign split:
\[
  A(x)=\mathrm{false}
  \Rightarrow
  0 <
  \mathrm{Re}\,\mathrm{marshallSignS}(A,\tau)\,
  \mathrm{Re}\,
  \bigl(\hat S^+_x\,\mathrm{magSectorEmbedding}(\Phi_{M+1})\bigr)(\tau),
\]
and
\[
  A(x)=\mathrm{true}
  \Rightarrow
  \mathrm{Re}\,\mathrm{marshallSignS}(A,\tau)\,
  \mathrm{Re}\,
  \bigl(\hat S^+_x\,\mathrm{magSectorEmbedding}(\Phi_{M+1})\bigr)(\tau)
  < 0 .
\]
These are formalised as
\[
\begin{gathered}
\texttt{tasaki23\_signed\_single\_site\_raising\_component}\\
\texttt{\_pos\_of\_A\_false}
\end{gathered}
\]
and
\[
\begin{gathered}
\texttt{tasaki23\_signed\_single\_site\_raising\_component}\\
\texttt{\_neg\_of\_A\_true}.
\end{gathered}
\]

\textbf{Raising-direction adjacent-sector energy package}: the
raising local component package now feeds the same sector-uniqueness
machinery as the lowering direction.  If
\(\Phi_{M+1}\) is an embedded eigenvector with eigenvalue \(\mu\), then
commutation of the Hamiltonian with \(\hat S^+_{\rm tot}\) gives
\[
  H\,\hat S^+_{\rm tot}\Psi_{M+1}
  = \mu\,\hat S^+_{\rm tot}\Psi_{M+1}.
\]
The support bridge places this raised vector in sector \(M\).  Hence,
under the Marshall-positive hypothesis
\[
  0 <
  \mathrm{Re}\,\mathrm{marshallSignS}(A,\tau)\,
  \mathrm{Re}\,
  \bigl(\hat S^+_{\rm tot}\Psi_{M+1}\bigr)(\tau)
  \qquad(\tau\in\mathrm{magConfigS}(V,N,M)),
\]
the sector-\(M\) Theorem~2.2 uniqueness theorem identifies the
sector-\(M\) eigenvalue with \(\mu\).  The non-zero conclusion follows
by evaluating the strict positive inequality at any target sector
configuration.  Lean packages these facts as
\[
\begin{gathered}
\texttt{tasaki23\_raised\_ne\_zero\_of\_marshall\_pos},\\
\texttt{tasaki23\_raising\_identifies\_adjacent\_sector\_energy},\\
\texttt{tasaki23\_raising\_identifies\_adjacent\_sector\_energy}\\
\texttt{\_with\_nonzero}.
\end{gathered}
\]
The site-sum bridge
\[
\begin{gathered}
\texttt{tasaki23\_raised\_marshall\_pos\_of\_site\_sum\_pos},\\
\texttt{tasaki23\_raising\_identifies\_adjacent\_sector\_energy}\\
\texttt{\_of\_site\_sum\_pos}
\end{gathered}
\]
allows the same comparison to be invoked from the local hypothesis
obtained by expanding \(\hat S^+_{\rm tot}\) as
\(\sum_x \hat S^+_x\).

\textbf{Raised sign-sum bounds and dominance bridge}: the raising
direction now has the same boundary-inclusive local sign-sum package as
the lowering direction.  Away from the top boundary, the existing
single-site raising formula shows that the signed local contribution is
strictly positive off \(A\) and strictly negative on \(A\).  At the top
boundary, the one-site raising summand is zero.  Lean records the
boundary-inclusive local forms as
\[
\begin{gathered}
\texttt{tasaki23\_signed\_single\_site\_raising\_component}\\
\texttt{\_nonneg\_of\_A\_false},\\
\texttt{tasaki23\_signed\_single\_site\_raising\_component}\\
\texttt{\_nonpos\_of\_A\_true}.
\end{gathered}
\]
Summing over the two Boolean sublattices gives
\[
\begin{gathered}
\texttt{tasaki23\_signed\_raising\_offA\_sum\_nonneg},\\
\texttt{tasaki23\_signed\_raising\_onA\_sum\_nonpos}.
\end{gathered}
\]
The strict off-\(A\) witness companion is
\[
\texttt{tasaki23\_signed\_raising\_offA\_sum\_pos\_of\_exists\_raiseable\_offA}.
\]
If \(x\notin A\) and \(\tau_x<N\), the strict single-site raising
positivity theorem supplies one positive summand, while all other
off-\(A\) summands remain non-negative.
The aggregate vacancy form is
\[
\begin{gathered}
\texttt{tasaki23\_exists\_raiseable\_offA\_of\_offA\_vacancy\_pos},\\
\texttt{tasaki23\_signed\_raising\_offA\_sum\_pos\_of\_offA\_vacancy\_pos}.
\end{gathered}
\]
It extracts a raiseable witness from a positive off-\(A\) vacancy sum
\(\sum_{x\notin A}(N-\tau_x)>0\).
Lean names the local raised contribution as
\texttt{tasaki23SignedRaisingSiteContribution} and decomposes the full
signed site-sum as
\[
  \mathrm{Re}\,\mathrm{marshallSignS}(A,\tau)
  \sum_{x\in V}
  \mathrm{Re}\,
  \bigl(\hat S^+_x\,\mathrm{magSectorEmbedding}(\Phi_{M+1})\bigr)(\tau)
  =
  \sum_{x:A(x)=\mathrm{false}} R_x(\tau)
  +
  \sum_{x:A(x)=\mathrm{true}} R_x(\tau),
\]
where \(R_x(\tau)\) denotes the named raised contribution.  This is
\texttt{tasaki23\_signed\_raising\_site\_sum\_eq\_offA\_add\_onA}.
Consequently the dominance condition
\[
  -\sum_{x:A(x)=\mathrm{true}} R_x(\tau)
  <
  \sum_{x:A(x)=\mathrm{false}} R_x(\tau)
\]
implies the strict raised site-sum positivity hypothesis:
\texttt{tasaki23\_signed\_raising\_site\_sum\_pos\_of\_onA\_neg\_lt\_offA}.
It is then wrapped into raised-vector Marshall positivity as
\texttt{tasaki23\_raised\_marshall\_pos\_of\_onA\_neg\_lt\_offA}.
The predecessor-sector common-energy step can now be invoked directly
from the same dominance hypothesis through
\texttt{tasaki23\_predecessor\_sector\_common\_energy\_of\_onA\_neg\_lt\_offA}.

\textbf{Adjacent common-energy successor link}: the interval
combinatorics and the lowered site-sum comparison are now packaged in
the Lean theorem
\[
\texttt{tasaki23\_successor\_sector\_common\_energy\_of\_site\_sum\_pos}.
\]
If \(M\) belongs to the admissible interval
\[
  [\min(|A|,|B|)N,\ \max(|A|,|B|)N]
\]
and is not its right endpoint, then
\[
  M+1 \in
  [\min(|A|,|B|)N,\ \max(|A|,|B|)N]
\]
by \texttt{tasaki23GroundStateSectors\_succ\_mem\_of\_mem\_of\_lt\_right}.
Given a Marshall-positive source-sector eigenvector
\[
  \Psi_M(\tau)=\operatorname{sgn}_A(\tau)v_M(\tau), \qquad v_M(\tau)>0,
\]
with eigenvalue \(\mu<c\), and the strict site-sum positivity hypothesis
for \(\hat S^-_{\rm tot}\Psi_M\) in sector \(M+1\), the theorem returns
non-vanishing of the lowered vector and a Marshall-positive vector in
sector \(M+1\) at the same eigenvalue \(\mu\).  This is the one-step
successor link needed to iterate Theorem~2.2 across the Theorem~2.3
multiplet.  The remaining substantive input is still the proof of the
strict site-sum positivity hypothesis for the actual sector ground
states.

\textbf{Predicted-Casimir threading layer (PR \#3917-\#3919 consolidation)}: an earlier version of this guide listed ~10 sub-paragraphs describing the modules \texttt{Theorem23PredictedLadderCasimirTransfer}, \texttt{Theorem23PredictedCasimirEnergy}, \texttt{Theorem23Dominance}, \texttt{Theorem23DominancePredictedGS}, etc., that threaded the predicted-Casimir hypothesis through the adjacent-sector chain. The Theorem 2.3 capstone consolidation refactor (Issue \#3920) removed those modules and the entire threading chain (PR \#3922 pass-through pruning, PR \#3919 525-file orphan sweep) because the structural-variant chain \texttt{tasaki\_2\_5\_theorem\_2\_3\_structural} (PR \#3891 toy coupling, PR \#3893 general bipartite \(J\)) provides a sound, truly-unconditional proof for any \(N \ge 1\) without the predicted-Casimir threading. The SU(2) Casimir commutation \([\hat S_{\mathrm{tot}}^2,\hat S^\pm_{\mathrm{tot}}]=0\) (formalised as \texttt{totalSpinSSquared\_commute\_totalSpinSOpMinus} and \texttt{...\_Plus}) is still part of the sound chain; the abandoned wrappers around it are not.

The predecessor-difference route has also been exposed at this same
site-sum boundary.  Lean names the callback-level adapter
\[
\begin{gathered}
\texttt{tasaki23\_lowered\_site\_sum\_pos\_callback}\\
\texttt{\_of\_unpacked\_reembedded\_real\_source\_weight}\\
\texttt{\_predecessor\_difference\_pos}.
\end{gathered}
\]
It is the site-sum API form of the already proved bridge from the
off-\(A\) minus on-\(A\) predecessor raising-source difference to
strict positivity of the lowered single-site sum.  The final wrapper now
uses this API directly:
\[
\begin{gathered}
\texttt{tasaki23\_energy\_interval\_chain\_with\_predictedGS}\\
\texttt{\_of\_left\_endpoint\_predictedGS}\\
\texttt{\_of\_unpacked\_reembedded\_real\_source\_weight}\\
\texttt{\_predecessor\_difference\_pos},\\
\texttt{tasaki\_2\_5\_theorem\_2\_3}\\
\texttt{\_of\_left\_endpoint\_threaded\_predictedGS}\\
\texttt{\_of\_unpacked\_reembedded\_real\_source\_weight}\\
\texttt{\_predecessor\_difference\_pos}\\
\texttt{\_via\_lowered\_site\_sum\_of\_sector\_minimality}.
\end{gathered}
\]
At each successor step, predicted-GS membership supplies the
re-embedded source-weight identity and the lowered sublattice
Casimir/support facts; the predecessor-difference callback is then
converted to strict lowered site-sum positivity and passed directly to
the site-sum successor chain.  This avoids first restating the same local
boundary as a raising-source dominance final callback.

The same direct site-sum route now has an outside-interval real-sector
minimality wrapper:
\[
\begin{gathered}
\texttt{tasaki\_2\_5\_theorem\_2\_3}\\
\texttt{\_of\_left\_endpoint\_threaded\_predictedGS}\\
\texttt{\_of\_unpacked\_reembedded\_real\_source\_weight}\\
\texttt{\_predecessor\_difference\_pos}\\
\texttt{\_via\_lowered\_site\_sum}\\
\texttt{\_of\_outside\_real\_sector\_minimality}.
\end{gathered}
\]
It uses the common Marshall-positive energy chain together with
\(\texttt{tasaki23\_real\_sector\_minimality\_on\_groundStateSectors}\)
\(\texttt{\_of\_common\_energy\_chain}\) to prove the real-sector
lower bound on every admissible sector.  The remaining spectral
callback therefore only covers sectors outside
\(\texttt{tasaki23GroundStateSectors}\).  This is the same
Perron--Frobenius comparison used above for the shifted dressed real
sector matrix; see Seneta, \emph{Non-negative Matrices and Markov
Chains}, 3rd ed., Springer 2006, \S1.2, pp.~27--28.

The outside-sector callback is further reduced by
\[
\begin{gathered}
\texttt{tasaki23\_outside\_real\_sector\_minimality}\\
\texttt{\_of\_outside\_sector\_ground\_energy\_lower\_bound}.
\end{gathered}
\]
For an outside sector \(M\), the per-sector Theorem 2.2 wrapper supplies
a Marshall-positive sector ground state at energy \(\mu_M<c\).  The
same shifted dressed real-sector Perron--Frobenius comparison gives
\(\mu_M\le \mu'\) for every nonzero real-form eigenvector in that
sector.  Hence the final outside-sector task may focus only on proving
\(\mu\le\mu_M\) for these Theorem 2.2 ground-state representatives,
rather than for arbitrary real eigenvectors directly.

The quantified callback shapes used in this final bridge are now named:
\[
\begin{gathered}
\texttt{tasaki23CommonEnergyChain},\\
\texttt{tasaki23OutsideGroundEnergyLowerCallback},\\
\texttt{tasaki23OutsideGroundEnergyLowerFamilyCallback},\\
\texttt{tasaki23LeftEndpointPredictedGSCallback},\\
\texttt{tasaki23SourcePredictedGSCallback},\\
\texttt{tasaki23PredecessorDifferenceCallback}.
\end{gathered}
\]
The module-split structure of the Theorem 2.3 chain that this paragraph used to describe at length (each `\texttt{LatticeSystem.Quantum.SpinS.Theorem23*}` module) was largely deleted in the consolidation refactor (Issue \#3920); see the consolidated-statement paragraph above for context. The current canonical structural variants live in the `\texttt{Theorem23Structural*}` modules.
The predecessor raising-source comparison is also exposed as the reusable
callback bridge
\[
\begin{gathered}
\texttt{tasaki23PredecessorDifferenceCallback}\\
\texttt{\_of\_unpacked\_reembedded\_real\_source\_weight}\\
\texttt{\_predecessor\_raising\_source\_sum\_lt}.
\end{gathered}
\]
It takes the fully threaded strict inequality saying that the on-\(A\)
predecessor raising-source sum is less than the off-\(A\) sum, and returns
\(\texttt{tasaki23PredecessorDifferenceCallback}\) by rewriting the same
comparison as positivity of the off-\(A\) minus on-\(A\) difference.  The
common-energy theorem with predecessor raising-source dominance now calls this
named bridge instead of carrying that `linarith` conversion anonymously.
The explicit lowerable coefficient comparison has the analogous direct bridge
\[
\begin{gathered}
\texttt{tasaki23PredecessorDifferenceCallback}\\
\texttt{\_of\_unpacked\_reembedded\_real\_source\_weight}\\
\texttt{\_predecessor\_explicit\_lowerable}\\
\texttt{\_positive\_source\_coefficient\_lt}.
\end{gathered}
\]
It first uses the local coefficient mirror
\[
\begin{gathered}
\texttt{tasaki23\_raising\_predecessor\_source\_sum\_lt}\\
\texttt{\_of\_lowerable\_positive\_source\_attach\_sum\_lt}
\end{gathered}
\]
to convert attached sums of
\(\texttt{tasaki23LoweringPredecessorPositiveSourceLowerableCoefficient}\)
to predecessor raising-source dominance, then returns the same
predecessor-difference callback.  The common-energy theorem from explicit
lowerable coefficients now enters the predecessor-difference common-energy
route through this named bridge.

The short final boundary
\[
\texttt{tasaki\_2\_5\_theorem\_2\_3\_of\_named\_callbacks}
\]
collects the same discharged outside-ground theorem behind exactly the
three remaining named callbacks:
\[
\begin{gathered}
\texttt{tasaki23SourcePredictedGSCallback},\\
\texttt{tasaki23PredecessorDifferenceCallback},\\
\texttt{tasaki23OutsideGroundEnergyLowerFamilyCallback}.
\end{gathered}
\]
It is an alias for the current threaded predicted-GS, predecessor-
difference, and outside-sector lower-bound route; no extra mathematical
hypothesis is added.

The variant
\[
\texttt{tasaki\_2\_5\_theorem\_2\_3\_of\_left\_endpoint\_named\_callbacks}
\]
uses the weaker left-endpoint predicted-GS callback instead of the
uniform source-sector callback.  Its other two visible inputs are still
\[
\begin{gathered}
\texttt{tasaki23PredecessorDifferenceCallback},\\
\texttt{tasaki23OutsideGroundEnergyLowerFamilyCallback}.
\end{gathered}
\]
The admissible-sector non-emptiness input is discharged by the
physical-range construction
\(\texttt{magConfigS\_nonempty\_of\_le\_card\_mul}\).

For the older direct lowered-site-sum route, Lean also names the two
remaining anonymous callbacks:
\[
\begin{gathered}
\texttt{tasaki23LoweredSiteSumCallback},\\
\texttt{tasaki23SectorMinimalityCallback}.
\end{gathered}
\]
The public aliases
\[
\begin{gathered}
\texttt{tasaki\_2\_5\_theorem\_2\_3\_of\_lowered\_site\_sum\_named\_callbacks},\\
\texttt{tasaki\_2\_5\_theorem\_2\_3\_of\_left\_endpoint\_lowered\_site\_sum}\\
\texttt{\_named\_callbacks}
\end{gathered}
\]
therefore expose the direct interval-chain route through named
predicted-GS, site-sum, and sector-minimality inputs without changing
the mathematical hypotheses.  The outside-ground callback layer also records
\[
\texttt{tasaki23OutsideGroundEnergyLowerFamilyCallback\_of\_sector\_minimality}.
\]
This bridge converts a sector-minimality callback into the
\texttt{tasaki23OutsideGroundEnergyLowerFamilyCallback}: the
Marshall-positive outside-sector representative is restricted to its
magnetization sector, non-zeroness follows from positivity and
\((\operatorname{Re}\mathrm{MarshallSign})^2=1\), and sector minimality
then gives the required \(\mu\le\mu_M\).  A second bridge records the
remaining ladder-reach form as
\[
\begin{gathered}
\texttt{tasaki23OutsideGroundAdmissibleReachCallback},\\
\texttt{tasaki23OutsideGroundEnergyLowerFamilyCallback}\\
\texttt{\_of\_admissible\_reach}.
\end{gathered}
\]
Here the local input is no longer a sectorwise lower bound: for each
Marshall-positive Theorem 2.2 representative outside the admissible
interval, it supplies a nonzero eigenvector at the same eigenvalue
inside some admissible sector.  The common-energy chain already gives
the real-sector lower bound on admissible sectors.  The reached complex
eigenvector is converted to that real-form statement by taking its real
part, or, if the real part vanishes, by taking its nonzero imaginary
part.  This proves \(\mu\le\mu_M\) and hence the same outside-ground
lower-family callback.  The outside reach task is then split into its
full-space and directional forms.  The full-space ladder form is named by
\[
\begin{gathered}
\texttt{tasaki23OutsideGroundAdmissibleFullReachCallback},\\
\texttt{tasaki23OutsideGroundAdmissibleReachCallback\_of\_full\_reach},\\
\texttt{tasaki23OutsideGroundEnergyLowerFamilyCallback}\\
\texttt{\_of\_full\_admissible\_reach}.
\end{gathered}
\]
This form asks the ladder argument for a nonzero full-space eigenvector
whose magnetization lies in an admissible sector.  Lean then restricts it
to sector coordinates with \(\texttt{magSectorRestriction}\).  The support
round-trip lemma
\(\texttt{magSectorEmbedding\_magSectorRestriction\_of\_mem\_magSubspaceS}\)
preserves non-zeroness, and
\(\texttt{heisenbergHamiltonianSMatrixOnMagSector\_mulVec\_magSectorRestriction\_of\_full\_eigen}\)
transfers the full eigen equation to the admissible sector matrix.  The
discharged final boundary
\[
\begin{gathered}
\texttt{tasaki\_2\_5\_theorem\_2\_3}\\
\texttt{\_of\_threaded\_predictedGS}\\
\texttt{\_of\_unpacked\_reembedded\_real\_source\_weight}\\
\texttt{\_predecessor\_difference\_pos}\\
\texttt{\_of\_full\_admissible\_reach}\\
\texttt{\_discharge\_nonempty}
\end{gathered}
\]
therefore exposes the remaining outside-reach input in the same
full-space form produced by total-spin ladder operators.
Lean further narrows this full-space reach input to non-zeroness of
iterated total-spin ladder outputs.  The local ladder layer records
\[
\begin{gathered}
\texttt{heisenbergHamiltonianS\_mulVec\_totalSpinSOpMinus}\\
\texttt{\_pow\_of\_eigenvec},\\
\texttt{heisenbergHamiltonianS\_mulVec\_totalSpinSOpPlus}\\
\texttt{\_pow\_of\_eigenvec},\\
\texttt{totalSpinSOpMinus\_pow\_mulVec\_mem\_magSubspaceS\_of\_mem},\\
\texttt{totalSpinSOpPlus\_pow\_mulVec\_mem\_magSubspaceS\_of\_mem}.
\end{gathered}
\]
Thus \((\hat S^-_{\rm tot})^k\) and \((\hat S^+_{\rm tot})^k\)
preserve the Heisenberg eigenvalue and shift the magnetization
subspace by \(-k\) and \(+k\), respectively.  The remaining directional
inputs can therefore be stated as
\[
\begin{gathered}
\texttt{tasaki23OutsideGroundLeftIteratedLadderFullReachCallback},\\
\texttt{tasaki23OutsideGroundRightIteratedLadderFullReachCallback},\\
\texttt{tasaki23OutsideGroundAdmissibleFullReachCallback}\\
\texttt{\_of\_iterated\_ladder\_callbacks},\\
\texttt{tasaki23OutsideGroundEnergyLowerFamilyCallback}\\
\texttt{\_of\_iterated\_ladder\_full\_reach}.
\end{gathered}
\]
The left callback chooses an admissible sector \(K=M+k\) and proves
that \((\hat S^-_{\rm tot})^k\Psi_M\ne0\); the right callback chooses
an admissible sector with \(M=K+k\) and proves
\((\hat S^+_{\rm tot})^k\Psi_M\ne0\).  The discharged final boundary
\[
\begin{gathered}
\texttt{tasaki\_2\_5\_theorem\_2\_3}\\
\texttt{\_of\_threaded\_predictedGS}\\
\texttt{\_of\_unpacked\_reembedded\_real\_source\_weight}\\
\texttt{\_predecessor\_difference\_pos}\\
\texttt{\_of\_iterated\_ladder\_full\_reach}\\
\texttt{\_discharge\_nonempty}
\end{gathered}
\]
then obtains the outside-sector lower family through
\(\texttt{tasaki23OutsideGroundEnergyLowerFamilyCallback\_of\_iterated\_ladder\_full\_reach}\)
and passes that lower family directly to the source predicted-GS common-energy
final boundary.
The non-zeroness input has a Casimir-based discharge:
\[
\begin{gathered}
\texttt{totalSpinSOpMinus\_pow\_mulVec\_ne\_zero}\\
\texttt{\_of\_casimir\_ne\_kernel\_values},\\
\texttt{totalSpinSOpPlus\_pow\_mulVec\_ne\_zero}\\
\texttt{\_of\_casimir\_ne\_kernel\_values},\\
\texttt{tasaki23OutsideGroundAdmissibleFullReachCallback}\\
\texttt{\_of\_iterated\_ladder\_casimir\_callbacks}.
\end{gathered}
\]
For each intermediate ladder step, these lemmas compare the preserved
total-Casimir eigenvalue with the one-step kernel value
\(m(m-1)\) in the lowering direction, respectively \(m(m+1)\) in the
raising direction.  The final predecessor-difference boundary
\[
\begin{gathered}
\texttt{tasaki\_2\_5\_theorem\_2\_3}\\
\texttt{\_of\_threaded\_predictedGS}\\
\texttt{\_of\_unpacked\_reembedded\_real\_source\_weight}\\
\texttt{\_predecessor\_difference\_pos}\\
\texttt{\_of\_iterated\_ladder\_casimir\_full\_reach}\\
\texttt{\_discharge\_nonempty}
\end{gathered}
\]
therefore replaces the raw iterated-ladder non-zeroness hypotheses by
left and right Casimir-avoidance callbacks, obtains the outside-sector lower
family through
\(\texttt{tasaki23OutsideGroundEnergyLowerFamilyCallback\_of\_iterated\_ladder\_casimir\_full\_reach}\),
and enters the same source predicted-GS common-energy final boundary directly.
The saturated-ferromagnet total-Casimir value gives a further endpoint
bridge.  Lean records the elementary mismatches
\[
\begin{gathered}
\texttt{saturatedFerromagnetCasimirEigenvalueS\_ne}\\
\texttt{\_lowering\_kernel\_value\_of\_lt\_card\_mul},\\
\texttt{saturatedFerromagnetCasimirEigenvalueS\_ne}\\
\texttt{\_raising\_kernel\_value\_of\_pos\_of\_le\_card\_mul}.
\end{gathered}
\]
In magnetization coordinate \(n\), these are the two facts that
\[
  m_{\max}(m_{\max}+1)
  \ne (m_{\max}-n)(m_{\max}-n-1)
  \quad (n<|V|N),
\]
and
\[
  m_{\max}(m_{\max}+1)
  \ne (m_{\max}-n)(m_{\max}-n+1)
  \quad (0<n\le |V|N).
\]
Thus a source vector already known to have the saturated total-Casimir
eigenvalue automatically avoids every intermediate kernel value on the
way to the left or right endpoint of the admissible interval.  The
callbacks
\[
\begin{gathered}
\texttt{tasaki23OutsideGroundLeftSaturatedCasimirSourceCallback},\\
\texttt{tasaki23OutsideGroundRightSaturatedCasimirSourceCallback}
\end{gathered}
\]
package this remaining max-Casimir source input, and
\[
\begin{gathered}
\texttt{tasaki23OutsideGroundLeftIteratedLadderCasimirFullReachCallback}\\
\texttt{\_of\_saturated\_casimir\_source},\\
\texttt{tasaki23OutsideGroundRightIteratedLadderCasimirFullReachCallback}\\
\texttt{\_of\_saturated\_casimir\_source}
\end{gathered}
\]
choose the left or right interval endpoint and discharge all
intermediate Casimir-avoidance obligations.  The corresponding final
wrapper is
\[
\begin{gathered}
\texttt{tasaki\_2\_5\_theorem\_2\_3}\\
\texttt{\_of\_threaded\_predictedGS}\\
\texttt{\_of\_unpacked\_reembedded\_real\_source\_weight}\\
\texttt{\_predecessor\_difference\_pos}\\
\texttt{\_of\_saturated\_casimir\_sources}\\
\texttt{\_discharge\_nonempty}.
\end{gathered}
\]
It now obtains the outside-sector lower family through
\(\texttt{tasaki23OutsideGroundEnergyLowerFamilyCallback\_of\_saturated\_casimir\_sources}\)
and passes that lower family directly to the source predicted-GS common-energy
final boundary.
The next wrapper replaces the saturated-Casimir source equation by the
more concrete ladder-span input.  Lean uses
\[
\begin{gathered}
\texttt{tasaki23OutsideGroundLeftSaturatedLadderSpanSourceCallback},\\
\texttt{tasaki23OutsideGroundRightSaturatedLadderSpanSourceCallback}
\end{gathered}
\]
together with the maximum-Casimir eigenspace identification
\[
\texttt{totalSpinSSquaredEigenspace\_max\_eq\_span\_ladderIterateUp}
\]
to prove
\[
\begin{gathered}
\texttt{tasaki23OutsideGroundLeftSaturatedCasimirSourceCallback}\\
\texttt{\_of\_ladder\_span\_source},\\
\texttt{tasaki23OutsideGroundRightSaturatedCasimirSourceCallback}\\
\texttt{\_of\_ladder\_span\_source}.
\end{gathered}
\]
Thus it is enough to show that the outside-sector source vector lies in
the span of the saturated ferromagnetic total-spin ladder.  The
corresponding final boundary is
\[
\begin{gathered}
\texttt{tasaki\_2\_5\_theorem\_2\_3}\\
\texttt{\_of\_threaded\_predictedGS}\\
\texttt{\_of\_unpacked\_reembedded\_real\_source\_weight}\\
\texttt{\_predecessor\_difference\_pos}\\
\texttt{\_of\_saturated\_ladder\_span\_sources}\\
\texttt{\_discharge\_nonempty}.
\end{gathered}
\]
It now obtains the outside-sector lower family through
\(\texttt{tasaki23OutsideGroundEnergyLowerFamilyCallback\_of\_saturated\_ladder\_span\_sources}\)
and passes that lower family directly to the source predicted-GS common-energy
final boundary.
Lean also records the same input in the saturated-ferromagnet joint
eigenspace form.  The callbacks
\[
\begin{gathered}
\texttt{tasaki23OutsideGroundLeftSaturatedJointSourceCallback},\\
\texttt{tasaki23OutsideGroundRightSaturatedJointSourceCallback}
\end{gathered}
\]
ask the source vector to lie in
\(\texttt{saturatedFerromagnetJointEigenspace}\ J\ N\).  The Tasaki
\S2.4 ladder-span identification
\[
\texttt{saturatedFerromagnetJointEigenspace\_eq\_span\_ladderIterateUp}
\]
then proves
\[
\begin{gathered}
\texttt{tasaki23OutsideGroundLeftSaturatedLadderSpanSourceCallback}\\
\texttt{\_of\_saturated\_joint\_source},\\
\texttt{tasaki23OutsideGroundRightSaturatedLadderSpanSourceCallback}\\
\texttt{\_of\_saturated\_joint\_source}.
\end{gathered}
\]
The corresponding final boundary is
\[
\begin{gathered}
\texttt{tasaki\_2\_5\_theorem\_2\_3}\\
\texttt{\_of\_threaded\_predictedGS}\\
\texttt{\_of\_unpacked\_reembedded\_real\_source\_weight}\\
\texttt{\_predecessor\_difference\_pos}\\
\texttt{\_of\_saturated\_joint\_sources}\\
\texttt{\_discharge\_nonempty}.
\end{gathered}
\]
It now obtains the outside-sector lower family through
\(\texttt{tasaki23OutsideGroundEnergyLowerFamilyCallback\_of\_saturated\_joint\_sources}\)
and passes that lower family directly to the source predicted-GS common-energy
final boundary.
The next boundary splits this joint membership into its two defining
eigenvalue equations.  Lean names the Heisenberg half by
\[
\begin{gathered}
\texttt{tasaki23OutsideGroundLeftSaturatedHeisenbergSourceCallback},\\
\texttt{tasaki23OutsideGroundRightSaturatedHeisenbergSourceCallback},
\end{gathered}
\]
and combines it with the existing saturated-Casimir source callbacks via
\[
\begin{gathered}
\texttt{tasaki23OutsideGroundLeftSaturatedJointSourceCallback}\\
\texttt{\_of\_saturated\_eigen\_sources},\\
\texttt{tasaki23OutsideGroundRightSaturatedJointSourceCallback}\\
\texttt{\_of\_saturated\_eigen\_sources}.
\end{gathered}
\]
Thus the final boundary
\[
\begin{gathered}
\texttt{tasaki\_2\_5\_theorem\_2\_3}\\
\texttt{\_of\_threaded\_predictedGS}\\
\texttt{\_of\_unpacked\_reembedded\_real\_source\_weight}\\
\texttt{\_predecessor\_difference\_pos}\\
\texttt{\_of\_saturated\_eigen\_sources}\\
\texttt{\_discharge\_nonempty}
\end{gathered}
\]
exposes the remaining source-vector task as the two saturated eigenvalue
equations for \(H_J\) and \((\hat S_{\mathrm{tot}})^2\).
The Heisenberg half can be reduced one step further.  The source
eigenvector equation already has the form
\[
  H_J\Phi = \mu_M\Phi.
\]
Lean packages the remaining scalar identification
\((\mu_M:\mathbb C)=E_{\mathrm{sat}}(J,N)\), where
\(E_{\mathrm{sat}}\) is
\texttt{saturatedFerromagnetEigenvalueS}, as
\[
\begin{gathered}
\texttt{tasaki23OutsideGroundLeftSaturatedEnergySourceCallback},\\
\texttt{tasaki23OutsideGroundRightSaturatedEnergySourceCallback}.
\end{gathered}
\]
Rewriting the source eigenvector equation along this scalar equality proves
\[
\begin{gathered}
\texttt{tasaki23OutsideGroundLeftSaturatedHeisenbergSourceCallback}\\
\texttt{\_of\_saturated\_energy\_source},\\
\texttt{tasaki23OutsideGroundRightSaturatedHeisenbergSourceCallback}\\
\texttt{\_of\_saturated\_energy\_source}.
\end{gathered}
\]
The corresponding final boundary is
\[
\begin{gathered}
\texttt{tasaki\_2\_5\_theorem\_2\_3}\\
\texttt{\_of\_threaded\_predictedGS}\\
\texttt{\_of\_unpacked\_reembedded\_real\_source\_weight}\\
\texttt{\_predecessor\_difference\_pos}\\
\texttt{\_of\_saturated\_source\_energy}\\
\texttt{\_discharge\_nonempty}.
\end{gathered}
\]
Conversely, the saturated joint source-vector input already contains both
eigenvalue components.  Lean records this extraction by projecting
membership in
\(\texttt{saturatedFerromagnetJointEigenspace}\) to the
\(H_J\)- and \((\widehat S_{\rm tot})^2\)-eigenspace factors.  The
Casimir projection directly gives
\[
\begin{gathered}
\texttt{tasaki23OutsideGroundLeftSaturatedCasimirSourceCallback}\\
\texttt{\_of\_saturated\_joint\_source},\\
\texttt{tasaki23OutsideGroundRightSaturatedCasimirSourceCallback}\\
\texttt{\_of\_saturated\_joint\_source}.
\end{gathered}
\]
For the scalar source-energy callback, the saturated \(H_J\) equation is
compared with the already available source equation
\(H_J\Phi=\mu_M\Phi\).  Since strict Marshall positivity implies
\(\Phi\neq0\), Lean cancels the nonzero vector from
\((\mu_M-E_{\rm sat})\Phi=0\) and obtains
\[
\begin{gathered}
\texttt{tasaki23OutsideGroundLeftSaturatedEnergySourceCallback}\\
\texttt{\_of\_saturated\_joint\_source},\\
\texttt{tasaki23OutsideGroundRightSaturatedEnergySourceCallback}\\
\texttt{\_of\_saturated\_joint\_source}.
\end{gathered}
\]
The corresponding final boundary is
\[
\begin{gathered}
\texttt{tasaki\_2\_5\_theorem\_2\_3}\\
\texttt{\_of\_threaded\_predictedGS}\\
\texttt{\_of\_unpacked\_reembedded\_real\_source\_weight}\\
\texttt{\_predecessor\_difference\_pos}\\
\texttt{\_of\_saturated\_joint\_sources}\\
\texttt{\_extract\_source\_eigen\_callbacks}\\
\texttt{\_discharge\_nonempty}.
\end{gathered}
\]
At the callback level, the same projection gives the lower-family bridge
\[
\begin{gathered}
\texttt{tasaki23OutsideGroundEnergyLowerFamilyCallback}\\
\texttt{\_of\_saturated\_joint\_sources}.
\end{gathered}
\]
It projects the left and right saturated joint source-vector callbacks to
their saturated total-Casimir source components, then reuses
\(\texttt{tasaki23OutsideGroundEnergyLowerFamilyCallback}\)
\(\texttt{\_of\_saturated\_casimir\_sources}\).  Thus the joint
source-vector API now supplies the outside-sector lower-family callback
directly, not only the final theorem boundary.
The public named-callback final theorem also consumes this bridge through
\[
\begin{gathered}
\texttt{tasaki\_2\_5\_theorem\_2\_3}\\
\texttt{\_of\_saturated\_joint\_sources}\\
\texttt{\_named\_callbacks}.
\end{gathered}
\]
It keeps the common-energy side in the existing source predicted-GS and
predecessor-difference callbacks, while the outside-sector input is supplied
directly by the left and right saturated joint source-vector callbacks.
The same uniform named-callback boundary is also exposed through the full
saturated-source extraction stack:
\[
\begin{gathered}
\texttt{tasaki\_2\_5\_theorem\_2\_3}\\
\texttt{\_of\_saturated\_casimir\_sources}\\
\texttt{\_named\_callbacks},\\
\texttt{tasaki\_2\_5\_theorem\_2\_3}\\
\texttt{\_of\_saturated\_ladder\_span\_sources}\\
\texttt{\_named\_callbacks},\\
\texttt{tasaki\_2\_5\_theorem\_2\_3}\\
\texttt{\_of\_saturated\_eigen\_sources}\\
\texttt{\_named\_callbacks},\\
\texttt{tasaki\_2\_5\_theorem\_2\_3}\\
\texttt{\_of\_saturated\_source\_energy}\\
\texttt{\_named\_callbacks},\\
\texttt{tasaki\_2\_5\_theorem\_2\_3}\\
\texttt{\_of\_saturated\_joint\_sources}\\
\texttt{\_extract\_source\_eigen\_callbacks}\\
\texttt{\_named\_callbacks},\\
\texttt{tasaki\_2\_5\_theorem\_2\_3}\\
\texttt{\_of\_saturated\_joint\_references}\\
\texttt{\_extract\_source\_eigen\_callbacks}\\
\texttt{\_named\_callbacks}.
\end{gathered}
\]
These wrappers match the later left-endpoint saturated-source stack, but keep
the common-energy input as the uniform source predicted-GS callback.
The companion
\[
\begin{gathered}
\texttt{tasaki\_2\_5\_theorem\_2\_3}\\
\texttt{\_of\_left\_endpoint}\\
\texttt{\_saturated\_joint\_sources}\\
\texttt{\_named\_callbacks}
\end{gathered}
\]
does the same for the weaker left-endpoint predicted-GS boundary.
The left-endpoint named-callback API is also exposed through the saturated
source stack:
\[
\begin{gathered}
\texttt{tasaki\_2\_5\_theorem\_2\_3}\\
\texttt{\_of\_left\_endpoint}\\
\texttt{\_saturated\_casimir\_sources}\\
\texttt{\_named\_callbacks},\\
\texttt{tasaki\_2\_5\_theorem\_2\_3}\\
\texttt{\_of\_left\_endpoint}\\
\texttt{\_saturated\_ladder\_span\_sources}\\
\texttt{\_named\_callbacks},\\
\texttt{tasaki\_2\_5\_theorem\_2\_3}\\
\texttt{\_of\_left\_endpoint}\\
\texttt{\_saturated\_eigen\_sources}\\
\texttt{\_named\_callbacks},\\
\texttt{tasaki\_2\_5\_theorem\_2\_3}\\
\texttt{\_of\_left\_endpoint}\\
\texttt{\_saturated\_source\_energy}\\
\texttt{\_named\_callbacks},\\
\texttt{tasaki\_2\_5\_theorem\_2\_3}\\
\texttt{\_of\_left\_endpoint}\\
\texttt{\_saturated\_joint\_sources}\\
\texttt{\_extract\_source\_eigen\_callbacks}\\
\texttt{\_named\_callbacks}.
\end{gathered}
\]
These variants compose the already established saturated-Casimir,
ladder-span, saturated-Heisenberg, scalar source-energy, and joint-source
extraction bridges with the left-endpoint final boundary.  Thus the weaker
left-endpoint predicted-GS hypothesis can be used at every saturated-source
outside-ground layer.
The source-vector side of the saturated joint callback is now connected to a
sectorwise saturated reference-vector API:
\[
\begin{gathered}
\texttt{tasaki23OutsideGroundLeftSaturatedJointReferenceCallback},\\
\texttt{tasaki23OutsideGroundRightSaturatedJointReferenceCallback},\\
\texttt{tasaki23OutsideGroundLeftSaturatedJointSourceCallback}\\
\texttt{\_of\_saturated\_joint\_reference},\\
\texttt{tasaki23OutsideGroundRightSaturatedJointSourceCallback}\\
\texttt{\_of\_saturated\_joint\_reference}.
\end{gathered}
\]
For each outside sector, the reference callback supplies a strictly
Marshall-positive real sector vector whose full embedding lies in
\(\texttt{saturatedFerromagnetJointEigenspace}\), together with a real
representative of the saturated Heisenberg eigenvalue.  The bridge restricts
the source and reference full eigenvectors to the sector, uses the
Marshall-positive eigenvalue-uniqueness theorem to identify the source energy
with the saturated energy, and then uses sector eigenvector uniqueness to show
that the source embedding is a positive scalar multiple of the saturated joint
reference.
Composing this reference-vector bridge with the existing saturated
joint-source outside-ground route gives the lower-family and named final
boundaries:
\[
\begin{gathered}
\texttt{tasaki23OutsideGroundEnergyLowerFamilyCallback}\\
\texttt{\_of\_saturated\_joint\_references},\\
\texttt{tasaki\_2\_5\_theorem\_2\_3}\\
\texttt{\_of\_threaded\_predictedGS}\\
\texttt{\_of\_unpacked\_reembedded\_real\_source\_weight}\\
\texttt{\_predecessor\_difference\_pos}\\
\texttt{\_of\_saturated\_joint\_references}\\
\texttt{\_discharge\_nonempty},\\
\texttt{tasaki\_2\_5\_theorem\_2\_3}\\
\texttt{\_of\_saturated\_joint\_references}\\
\texttt{\_named\_callbacks},\\
\texttt{tasaki\_2\_5\_theorem\_2\_3}\\
\texttt{\_of\_left\_endpoint}\\
\texttt{\_saturated\_joint\_references}\\
\texttt{\_named\_callbacks}.
\end{gathered}
\]
Thus the reference-vector API can now be used directly at the discharged
predecessor-difference final boundary and at both public named-callback final
boundaries.  In the discharged route, Lean passes the resulting lower family
directly to the source predicted-GS common-energy final boundary instead of
calling the saturated joint-source final wrapper.
The reference-vector input is also exposed in the more concrete saturated
ladder-span form.  Lean names the callbacks
\[
\begin{gathered}
\texttt{tasaki23OutsideGroundLeftSaturatedLadderReferenceCallback},\\
\texttt{tasaki23OutsideGroundRightSaturatedLadderReferenceCallback}.
\end{gathered}
\]
They keep the real saturated Heisenberg eigenvalue witness and strict
Marshall positivity from the joint-reference callback, but replace
membership in
\(\texttt{saturatedFerromagnetJointEigenspace}\) by membership in
\[
  \operatorname{span}_{\mathbb C}
    \{\texttt{ladderIterateUp}\ V\ N\ k\}.
\]
The Tasaki \S2.4 identification
\[
\texttt{saturatedFerromagnetJointEigenspace\_eq\_span\_ladderIterateUp}
\]
proves
\[
\begin{gathered}
\texttt{tasaki23OutsideGroundLeftSaturatedJointReferenceCallback}\\
\texttt{\_of\_saturated\_ladder\_reference},\\
\texttt{tasaki23OutsideGroundRightSaturatedJointReferenceCallback}\\
\texttt{\_of\_saturated\_ladder\_reference}.
\end{gathered}
\]
Composing with the joint-reference route gives
\[
\begin{gathered}
\texttt{tasaki23OutsideGroundEnergyLowerFamilyCallback}\\
\texttt{\_of\_saturated\_ladder\_references},\\
\texttt{tasaki\_2\_5\_theorem\_2\_3}\\
\texttt{\_of\_saturated\_ladder\_references}\\
\texttt{\_named\_callbacks},\\
\texttt{tasaki\_2\_5\_theorem\_2\_3}\\
\texttt{\_of\_left\_endpoint}\\
\texttt{\_saturated\_ladder\_references}\\
\texttt{\_named\_callbacks}.
\end{gathered}
\]
The same ladder-reference inputs are also exposed at the discharged and
source-eigen extraction final boundaries:
\[
\begin{gathered}
\texttt{tasaki\_2\_5\_theorem\_2\_3}\\
\texttt{\_of\_threaded\_predictedGS}\\
\texttt{\_of\_unpacked\_reembedded\_real\_source\_weight}\\
\texttt{\_predecessor\_difference\_pos}\\
\texttt{\_of\_saturated\_ladder\_references}\\
\texttt{\_discharge\_nonempty},\\
\texttt{tasaki\_2\_5\_theorem\_2\_3}\\
\texttt{\_of\_threaded\_predictedGS}\\
\texttt{\_of\_unpacked\_reembedded\_real\_source\_weight}\\
\texttt{\_predecessor\_difference\_pos}\\
\texttt{\_of\_saturated\_ladder\_references}\\
\texttt{\_extract\_source\_eigen\_callbacks}\\
\texttt{\_discharge\_nonempty},\\
\texttt{tasaki\_2\_5\_theorem\_2\_3}\\
\texttt{\_of\_saturated\_ladder\_references}\\
\texttt{\_extract\_source\_eigen\_callbacks}\\
\texttt{\_named\_callbacks},\\
\texttt{tasaki\_2\_5\_theorem\_2\_3}\\
\texttt{\_of\_left\_endpoint}\\
\texttt{\_saturated\_ladder\_references}\\
\texttt{\_extract\_source\_eigen\_callbacks}\\
\texttt{\_named\_callbacks}.
\end{gathered}
\]
These wrappers first convert ladder references to saturated joint references.
The non-extraction discharged wrapper passes the resulting lower family
directly to the source predicted-GS common-energy final boundary; the
source-eigen extraction wrappers still reuse the corresponding saturated
joint-reference extraction boundary.
This pushes the remaining reference construction toward explicit total-spin
ladder vectors rather than abstract joint-eigenspace membership.
Lean also records the sectorwise singleton-ladder version of this input.
The callbacks
\[
\begin{gathered}
\texttt{tasaki23OutsideGroundLeftSaturatedLadderIterateReferenceCallback},\\
\texttt{tasaki23OutsideGroundRightSaturatedLadderIterateReferenceCallback}
\end{gathered}
\]
ask only that the Marshall-positive reference embedding lie in the singleton
span of the sector iterate
\(\texttt{ladderIterateUp}\ V\ N\ \langle M,\_\rangle\).  Since this
singleton is contained in the full range of
\(\texttt{ladderIterateUp}\ V\ N\), Lean first records the span-inclusion
helper in \texttt{SaturatedLadderJointEigenspace.lean},
\[
\texttt{ladderIterateUp\_singleton\_span\_le\_span\_range}.
\]
It then proves
\[
\begin{gathered}
\texttt{tasaki23OutsideGroundLeftSaturatedLadderReferenceCallback}\\
\texttt{\_of\_saturated\_ladder\_iterate\_reference},\\
\texttt{tasaki23OutsideGroundRightSaturatedLadderReferenceCallback}\\
\texttt{\_of\_saturated\_ladder\_iterate\_reference}.
\end{gathered}
\]
Composing these inclusions with the saturated ladder-reference route gives
\[
\begin{gathered}
\texttt{tasaki23OutsideGroundEnergyLowerFamilyCallback}\\
\texttt{\_of\_saturated\_ladder\_iterate\_references},\\
\texttt{tasaki\_2\_5\_theorem\_2\_3}\\
\texttt{\_of\_saturated\_ladder\_iterate\_references}\\
\texttt{\_named\_callbacks},\\
\texttt{tasaki\_2\_5\_theorem\_2\_3}\\
\texttt{\_of\_left\_endpoint}\\
\texttt{\_saturated\_ladder\_iterate\_references}\\
\texttt{\_named\_callbacks}.
\end{gathered}
\]
Thus the public named-callback boundary can now be stated using a single
explicit sector ladder iterate, with the span inclusion handling the
conversion to the earlier ladder-span API.
The same singleton-ladder input is now exposed at the discharged and
source-eigen extraction final boundaries:
\[
\begin{gathered}
\texttt{tasaki\_2\_5\_theorem\_2\_3}\\
\texttt{\_of\_threaded\_predictedGS}\\
\texttt{\_of\_unpacked\_reembedded\_real\_source\_weight}\\
\texttt{\_predecessor\_difference\_pos}\\
\texttt{\_of\_saturated\_ladder\_iterate\_references}\\
\texttt{\_discharge\_nonempty},\\
\texttt{tasaki\_2\_5\_theorem\_2\_3}\\
\texttt{\_of\_threaded\_predictedGS}\\
\texttt{\_of\_unpacked\_reembedded\_real\_source\_weight}\\
\texttt{\_predecessor\_difference\_pos}\\
\texttt{\_of\_saturated\_ladder\_iterate\_references}\\
\texttt{\_extract\_source\_eigen\_callbacks}\\
\texttt{\_discharge\_nonempty},\\
\texttt{tasaki\_2\_5\_theorem\_2\_3}\\
\texttt{\_of\_saturated\_ladder\_iterate\_references}\\
\texttt{\_extract\_source\_eigen\_callbacks}\\
\texttt{\_named\_callbacks},\\
\texttt{tasaki\_2\_5\_theorem\_2\_3}\\
\texttt{\_of\_left\_endpoint}\\
\texttt{\_saturated\_ladder\_iterate\_references}\\
\texttt{\_extract\_source\_eigen\_callbacks}\\
\texttt{\_named\_callbacks}.
\end{gathered}
\]
These wrappers first convert singleton ladder-iterate references to saturated
ladder references.  The non-extraction discharged wrapper passes the resulting
lower family directly to the source predicted-GS common-energy final boundary;
the source-eigen extraction wrappers still reuse the corresponding saturated
ladder-reference extraction boundary.
The next boundary makes the remaining coefficient-level construction explicit.
Lean introduces
\[
\begin{gathered}
\texttt{tasaki23OutsideGroundLeftSaturatedLadderIterate}\\
\texttt{MarshallPositiveReferenceCallback},\\
\texttt{tasaki23OutsideGroundRightSaturatedLadderIterate}\\
\texttt{MarshallPositiveReferenceCallback}.
\end{gathered}
\]
These callbacks replace singleton-span membership by the stronger equality
\[
\texttt{magSectorEmbedding}\bigl(
  ((\texttt{marshallSignS}\ A\ \tau.1).\texttt{re})\,w(\tau)
\bigr)
= \texttt{ladderIterateUp}\ V\ N\ \langle M,\_\rangle ,
\]
with \(w(\tau)>0\) for every sector configuration \(\tau\).  Equality with
the ladder iterate immediately gives membership in the singleton span, so Lean
proves
\[
\begin{gathered}
\texttt{tasaki23OutsideGroundLeftSaturatedLadderIterateReferenceCallback}\\
\texttt{\_of\_marshall\_positive\_reference},\\
\texttt{tasaki23OutsideGroundRightSaturatedLadderIterateReferenceCallback}\\
\texttt{\_of\_marshall\_positive\_reference}
\end{gathered}
\]
by rewriting with that equality and applying
\(\texttt{Submodule.subset\_span}\) to the singleton member.  Composing this
with the saturated ladder-iterate reference route yields
\[
\begin{gathered}
\texttt{tasaki23OutsideGroundEnergyLowerFamilyCallback}\\
\texttt{\_of\_saturated\_ladder\_iterate}\\
\texttt{\_marshall\_positive\_references}.
\end{gathered}
\]
This is still a Tasaki §2.5 Theorem~2.3 boundary (Tasaki, Springer 2020,
§2.5 Theorem~2.3, p.~42): it isolates the remaining task as a concrete
strict-positivity formula for the sector total-spin ladder iterate.
Lean then exposes that formula as a pointwise sector-coefficient boundary.
The callbacks
\[
\begin{gathered}
\texttt{tasaki23OutsideGroundLeftSaturatedLadderIterate}\\
\texttt{MarshallPositiveCoefficientCallback},\\
\texttt{tasaki23OutsideGroundRightSaturatedLadderIterate}\\
\texttt{MarshallPositiveCoefficientCallback}
\end{gathered}
\]
require, for every sector configuration \(\tau\),
\[
\texttt{ladderIterateUp}\ V\ N\ \langle M,\_\rangle\ \tau.1
=
((\texttt{marshallSignS}\ A\ \tau.1).\texttt{re})\,w(\tau),
\qquad w(\tau)>0.
\]
The full-vector equality follows because the ladder iterate belongs to the
magnetization subspace with label
\((|V|N/2)-M\):
\[
\texttt{totalSpinSOpMinus\_pow\_allAlignedStateS\_zero\_mem\_magSubspaceS}.
\]
Using the zero-extension inverse
\[
\texttt{magSectorEmbedding\_magSectorRestriction\_of\_mem\_magSubspaceS},
\]
Lean proves
\[
\begin{gathered}
\texttt{tasaki23OutsideGroundLeftSaturatedLadderIterate}\\
\texttt{MarshallPositiveReferenceCallback}\\
\texttt{\_of\_marshall\_positive\_coefficients},\\
\texttt{tasaki23OutsideGroundRightSaturatedLadderIterate}\\
\texttt{MarshallPositiveReferenceCallback}\\
\texttt{\_of\_marshall\_positive\_coefficients}.
\end{gathered}
\]
Composing once more with the reference route gives
\[
\begin{gathered}
\texttt{tasaki23OutsideGroundEnergyLowerFamilyCallback}\\
\texttt{\_of\_saturated\_ladder\_iterate}\\
\texttt{\_marshall\_positive\_coefficients}.
\end{gathered}
\]
Thus the remaining positivity task is now purely coefficient-wise on the
sector restriction of the unnormalised total-spin ladder iterate.
The first recursive coefficient boundary is the successor site-sum formula.
Assume the \(M\)-sector iterate has the pointwise form
\[
\texttt{ladderIterateUp}\ V\ N\ \langle M,\_\rangle\ \sigma.1
=
((\texttt{marshallSignS}\ A\ \sigma.1).\texttt{re})\,w(\sigma).
\]
If \(M+1\) is still in the ladder range, then
\[
\begin{aligned}
&\texttt{ladderIterateUp}\ V\ N\ \langle M+1,\_\rangle\ \tau.1 \\
&\quad =
\sum_{x\in V}
\Bigl(
\texttt{onSiteS}\ x\ (\texttt{spinSOpMinus}\ N)
\Bigr)
\Bigl(
\texttt{magSectorEmbedding}
  (\lambda \sigma.\,
   ((\texttt{marshallSignS}\ A\ \sigma.1).\texttt{re})w(\sigma))
\Bigr)(\tau.1).
\end{aligned}
\]
Lean names this as
\[
\begin{gathered}
\texttt{tasaki23OutsideGroundSaturatedLadderIterate}\\
\texttt{MarshallPositiveCoefficientSuccSiteSum}.
\end{gathered}
\]
The proof uses the iterate identity
\[
\texttt{totalSpinSOpMinus\_mulVec\_ladderIterateUp\_interior}
\]
to rewrite the successor iterate as \( \hat S^-_{\mathrm{tot}}\) applied to
the previous iterate, reuses the coefficient-to-zero-extension bridge above,
and finally applies
\[
\texttt{totalSpinSOpMinus\_mulVec\_magSectorEmbedding\_apply\_eq\_site\_sum}.
\]
Thus the remaining strict-positivity step is localized to the single-site
lowering summands in Tasaki §2.5 Theorem~2.3 (Tasaki, Springer 2020,
§2.5 Theorem~2.3, p.~42).
The next boundary packages that localized positivity into the successor
coefficient form.  If the lowerable positive-source predecessor coefficient
sum over \(A\) is strictly dominated by the corresponding sum over
\(\neg A\), then the successor coefficients admit strictly positive weights
\(w_{\mathrm{succ}}\):
\[
\begin{gathered}
\texttt{tasaki23OutsideGroundSaturatedLadderIterate}\\
\texttt{MarshallPositiveCoefficientSucc}\\
\texttt{\_of\_positive\_source\_lowerable\_coefficient\_lt}.
\end{gathered}
\]
Here \(w_{\mathrm{succ}}(\tau)\) is the Marshall sign times the real part of
the successor coefficient.  The proof uses
\[
\texttt{totalSpinSOpMinus\_mulVec\_marshallSignedEmbedding\_im\_zero}
\]
for realness and
\[
\texttt{tasaki23\_lowered\_marshall\_pos}\\
\texttt{\_of\_positive\_source\_lowerable\_coefficient\_lt}
\]
for strict positivity.
The base of the coefficient chain is also explicit:
\[
\begin{gathered}
\texttt{tasaki23OutsideGroundSaturatedLadderIterate}\\
\texttt{MarshallPositiveCoefficientZero}.
\end{gathered}
\]
It uses the uniqueness of the \(M=0\) sector,
\(\texttt{magConfigS\_zero\_eq\_allAlignedConfigS}\), and the identity
\(\texttt{marshallSignS\_const\_zero}\) to give the constant positive weight
\(w_0=1\) for \(\texttt{ladderIterateUp}\ V\ N\ \langle0,\_\rangle\).
Lean then packages this base and the one-step result into the full coefficient
chain.  The intermediate induction boundary
\[
\begin{gathered}
\texttt{tasaki23OutsideGroundSaturatedLadderIterate}\\
\texttt{MarshallPositiveCoefficientFamily}\\
\texttt{\_of\_zero\_and\_successor\_dominance}
\end{gathered}
\]
takes a positive coefficient representation at \(M=0\), plus the same
lowerable positive-source dominance input at each successor step, and returns
strictly positive Marshall-signed coefficients for every
\(\texttt{ladderIterateUp}\ V\ N\ \langle M,\_\rangle\) in the saturated range.
The closed form
\[
\begin{gathered}
\texttt{tasaki23OutsideGroundSaturatedLadderIterate}\\
\texttt{MarshallPositiveCoefficientFamily}\\
\texttt{\_of\_successor\_dominance}
\end{gathered}
\]
discharges the \(M=0\) input with the zero-sector theorem, so successor
dominance is the only remaining coefficient input.
This is the explicit induction closure of the coefficient recursion used at
the Tasaki §2.5 Theorem~2.3 outside-ground boundary.  Lean also records the
left and right callback forms
\[
\begin{gathered}
\texttt{tasaki23OutsideGroundLeftSaturatedLadderIterate}\\
\texttt{MarshallPositiveCoefficientCallback}\\
\texttt{\_of\_successor\_dominance},\\
\texttt{tasaki23OutsideGroundRightSaturatedLadderIterate}\\
\texttt{MarshallPositiveCoefficientCallback}\\
\texttt{\_of\_successor\_dominance}.
\end{gathered}
\]
These wrappers feed the closed coefficient family into the existing
outside-sector coefficient callbacks once the saturated Heisenberg eigenvalue
has been represented by a real scalar.  On the left, the endpoint hypothesis
gives \(M<|V|N+1\) by the cardinality bound on the smaller sublattice.  On the
right, nonemptiness of \(\texttt{magConfigS}\ V\ N\ M\) gives a configuration
with magnetization \(M\), and \(\texttt{magSumS\_le}\) yields
\(M\le |V|N\), hence \(M<|V|N+1\).  This is the coefficient-level discharge
needed by the Marshall-positive reference bridge (Tasaki, Springer 2020,
§2.5 Theorem~2.3, p.~42).
Lean also packages this coefficient-level discharge as an outside-sector
lower-family theorem:
\[
\begin{gathered}
\texttt{tasaki23OutsideGroundEnergyLowerFamilyCallback}\\
\texttt{\_of\_saturated\_ladder\_iterate}\\
\texttt{\_successor\_dominance}.
\end{gathered}
\]
It supplies the left and right successor-dominance coefficient callbacks to the
Marshall-positive coefficient bridge, thereby producing the full
\(\texttt{tasaki23OutsideGroundEnergyLowerFamilyCallback}\) without requiring
callers to package the two sides separately.  The real-coupling variant uses
\(\texttt{saturatedFerromagnetEigenvalueS\_exists\_real}\) to discharge the
saturated-energy real-scalar input.  The discharged final wrappers with suffixes
\[
\begin{gathered}
\texttt{\_of\_saturated\_ladder\_iterate}\\
\texttt{\_successor\_dominance\_discharge\_nonempty},\\
\texttt{\_of\_saturated\_ladder\_iterate}\\
\texttt{\_successor\_dominance}\\
\texttt{\_of\_real\_couplings\_discharge\_nonempty}
\end{gathered}
\]
feed the same lower-family route directly into the source predicted-GS
common-energy final boundary, while preserving the public hypotheses of the
threaded predicted-GS and predecessor-difference discharged API.
The uniform and left-endpoint named-callback final boundaries expose the same
successor-dominance route through suffixes
\[
\begin{gathered}
\texttt{\_of\_saturated\_ladder\_iterate}\\
\texttt{\_successor\_dominance\_named\_callbacks},\\
\texttt{\_of\_left\_endpoint\_saturated\_ladder\_iterate}\\
\texttt{\_successor\_dominance\_named\_callbacks},
\end{gathered}
\]
together with real-coupling variants.  These wrappers pass
\(\texttt{tasaki23OutsideGroundEnergyLowerFamilyCallback}\)
\(\texttt{\_of\_saturated\_ladder\_iterate}\)
\(\texttt{\_successor\_dominance}\) to
\(\texttt{tasaki\_2\_5\_theorem\_2\_3\_of\_named\_callbacks}\) or to the
left-endpoint analogue.  As in the discharged final boundary, the
real-coupling variants expose no separate saturated-energy real-scalar
hypothesis.
The remaining dominance input can also be stated in the explicit lowerable
attached-sum form already used in the local coefficient API.  The theorem
\[
\begin{gathered}
\texttt{tasaki23OutsideGroundSaturatedLadderIterate}\\
\texttt{MarshallPositiveCoefficientFamily}\\
\texttt{\_of\_lowerable\_attach\_sum\_dominance}
\end{gathered}
\]
uses
\(\texttt{tasaki23\_positive\_source\_lowerable\_coefficient\_lt\_of\_attach\_sum\_lt}\)
to turn dominance of attached sums of
\(\texttt{tasaki23LoweringPredecessorPositiveSourceLowerableCoefficient}\)
into the filtered positive-source coefficient dominance required by the
successor induction.  The left and right coefficient callbacks have the same
attached-sum variants, with suffix
\(\texttt{\_of\_lowerable\_attach\_sum\_dominance}\), so the outside-ground
boundary no longer needs to mention the boundary-inclusive filtered sum
directly.  Their real-coupling variants, with suffix
\[
\texttt{\_of\_real\_couplings},
\]
use \(\texttt{saturatedFerromagnetEigenvalueS\_exists\_real}\) to discharge
the saturated-energy real-scalar input at the left and right coefficient
callback level.  Finally,
\[
\begin{gathered}
\texttt{tasaki23OutsideGroundEnergyLowerFamilyCallback}\\
\texttt{\_of\_saturated\_ladder\_iterate}\\
\texttt{\_lowerable\_attach\_sum\_dominance}
\end{gathered}
\]
combines the left and right attached-sum coefficient callbacks with the
existing saturated ladder-iterate coefficient bridge.  Thus, once the
saturated ferromagnetic energy is represented by a real scalar, the same
explicit lowerable dominance input supplies the outside-sector lower-energy
family used by the final Tasaki Theorem~2.3 boundary.  The real-coupling
variant of this lower-family theorem discharges the same saturated-energy
real-scalar input internally from the realness of the couplings.  The
uniform and left-endpoint named final wrappers with suffix
\[
\texttt{\_of\_saturated\_ladder\_iterate\_lowerable\_attach\_sum\_dominance}
\]
consume this explicit lowerable attach-sum input directly as their
outside-sector lower-family source; their real-coupling variants use the
real-coupling lower-family theorem and expose no separate saturated-energy
real-scalar hypothesis.  The corresponding discharged non-empty final
wrappers feed the same lower-family theorem directly into the source
predicted-GS common-energy final boundary, so the discharged route can also be
stated directly from the explicit lowerable attach-sum dominance input.
The companion theorem
\[
\begin{gathered}
\texttt{tasaki23OutsideGroundEnergyLowerFamilyCallback}\\
\texttt{\_of\_saturated\_ladder\_iterate}\\
\texttt{\_predecessor\_difference\_pos}
\end{gathered}
\]
uses
\(\texttt{tasaki23\_lowerable\_positive\_source\_attach\_sum\_lt\_callback\_of\_offA\_sub\_onA\_pos}\)
to obtain that attached-sum dominance from the off-\(A\) minus on-\(A\)
predecessor raising-source positive-difference form.  This is the same local
comparison shape used by the predecessor-difference boundary, now connected
directly to the saturated ladder-iterate outside-ground route.  Lean packages
that pointwise input as
\[
\texttt{tasaki23SaturatedLadderIteratePredecessorDifferencePosCallback}
\]
and also records the real-coupling variant
\[
\begin{gathered}
\texttt{tasaki23OutsideGroundEnergyLowerFamilyCallback}\\
\texttt{\_of\_saturated\_ladder\_iterate}\\
\texttt{\_predecessor\_difference\_pos}\\
\texttt{\_of\_real\_couplings}.
\end{gathered}
\]
Here the separate real-scalar hypothesis for the saturated ferromagnetic
energy is supplied by
\(\texttt{saturatedFerromagnetEigenvalueS\_exists\_real}\), whose proof uses
the explicit saturated-energy formula and the realness of every \(J(x,y)\).
Lean also exposes a discharged final wrapper with the same suffix but without
\(\texttt{\_of\_real\_couplings}\).  That wrapper keeps the saturated-energy
real-scalar witness as an explicit input and feeds
\(\texttt{tasaki23OutsideGroundEnergyLowerFamilyCallback\_of\_saturated\_ladder\_iterate\_predecessor\_difference\_pos}\)
directly into the source predicted-GS common-energy final boundary.  The
real-coupling final API
wrapper
\[
\begin{gathered}
\texttt{tasaki\_2\_5\_theorem\_2\_3}\\
\texttt{\_of\_threaded\_predictedGS}\\
\texttt{\_of\_unpacked\_reembedded}\\
\texttt{\_real\_source\_weight}\\
\texttt{\_predecessor\_difference\_pos}\\
\texttt{\_of\_saturated\_ladder\_iterate}\\
\texttt{\_predecessor\_difference\_pos}\\
\texttt{\_of\_real\_couplings}\\
\texttt{\_discharge\_nonempty}
\end{gathered}
\]
uses the real-coupling lower-family callback directly and invokes the same
source predicted-GS common-energy final boundary, so the real-coupling
discharged route no longer exposes either the outside-sector lower-family
hypothesis or the saturated-energy real-scalar hypothesis.  The
shorter named wrapper
\[
\begin{gathered}
\texttt{tasaki\_2\_5\_theorem\_2\_3}\\
\texttt{\_of\_saturated\_ladder\_iterate}\\
\texttt{\_predecessor\_difference\_pos\_named\_callbacks}
\end{gathered}
\]
passes the same lower-family callback into
\(\texttt{tasaki\_2\_5\_theorem\_2\_3\_of\_named\_callbacks}\), keeping the
threaded predicted-GS and common-energy predecessor-difference inputs unchanged
while discharging the outside-sector lower bounds from the saturated
ladder-iterate predecessor-difference route.  Its left-endpoint analogue
\[
\begin{gathered}
\texttt{tasaki\_2\_5\_theorem\_2\_3}\\
\texttt{\_of\_left\_endpoint}\\
\texttt{\_saturated\_ladder\_iterate}\\
\texttt{\_predecessor\_difference\_pos\_named\_callbacks}
\end{gathered}
\]
does the same for
\(\texttt{tasaki\_2\_5\_theorem\_2\_3\_of\_left\_endpoint\_named\_callbacks}\).
The real-coupling named variants
\[
\begin{gathered}
\texttt{tasaki\_2\_5\_theorem\_2\_3}\\
\texttt{\_of\_saturated\_ladder\_iterate}\\
\texttt{\_predecessor\_difference\_pos}\\
\texttt{\_of\_real\_couplings\_named\_callbacks},\\
\texttt{tasaki\_2\_5\_theorem\_2\_3}\\
\texttt{\_of\_left\_endpoint}\\
\texttt{\_saturated\_ladder\_iterate}\\
\texttt{\_predecessor\_difference\_pos}\\
\texttt{\_of\_real\_couplings\_named\_callbacks}
\end{gathered}
\]
remove the visible saturated-energy real-scalar input from both public named
boundaries.  This is still a boundary for Tasaki, Springer 2020, §2.5
Theorem~2.3, p.~42; it only replaces an explicit scalar-realness premise by
the real-coupling hypothesis already present in the theorem.  The uniform
real-coupling named variant now also routes through
\(\texttt{tasaki\_2\_5\_theorem\_2\_3\_of\_named\_callbacks}\), so it uses
the same reduced common-energy plus outside-sector final path as the other
named routes; the non-negative-matrix comparison remains the Seneta,
3rd ed., §1.2, pp.~27--28 step described above.
The left-endpoint named boundary also has the extracted reference-vector
form
\[
\begin{gathered}
\texttt{tasaki\_2\_5\_theorem\_2\_3}\\
\texttt{\_of\_left\_endpoint}\\
\texttt{\_saturated\_joint\_references}\\
\texttt{\_extract\_source\_eigen\_callbacks}\\
\texttt{\_named\_callbacks}.
\end{gathered}
\]
It first applies the same reference-to-source conversion, then reuses the
left-endpoint saturated joint-source extraction route through the scalar
source-energy and saturated-Casimir callback layer.
The same reference-vector API is also exposed at the discharged
predecessor-difference final boundary by
\[
\begin{gathered}
\texttt{tasaki\_2\_5\_theorem\_2\_3}\\
\texttt{\_of\_threaded\_predictedGS}\\
\texttt{\_of\_unpacked\_reembedded\_real\_source\_weight}\\
\texttt{\_predecessor\_difference\_pos}\\
\texttt{\_of\_saturated\_joint\_references}\\
\texttt{\_discharge\_nonempty},\\
\texttt{tasaki\_2\_5\_theorem\_2\_3}\\
\texttt{\_of\_threaded\_predictedGS}\\
\texttt{\_of\_unpacked\_reembedded\_real\_source\_weight}\\
\texttt{\_predecessor\_difference\_pos}\\
\texttt{\_of\_saturated\_joint\_references}\\
\texttt{\_extract\_source\_eigen\_callbacks}\\
\texttt{\_discharge\_nonempty}.
\end{gathered}
\]
Both wrappers first turn the saturated joint references into saturated joint
source-vector callbacks.  The direct wrapper then converts those callbacks to
the outside-sector lower family and feeds that family directly into the
source predicted-GS common-energy final boundary, while the extracted wrapper
reuses the route through the scalar source-energy and saturated-Casimir
extraction layer.
Lean also exposes the same direct lower-family step one layer earlier:
\[
\begin{gathered}
\texttt{tasaki23OutsideGroundEnergyLowerFamilyCallback}\\
\texttt{\_of\_saturated\_eigen\_sources},\\
\texttt{tasaki23OutsideGroundEnergyLowerFamilyCallback}\\
\texttt{\_of\_saturated\_source\_energy}.
\end{gathered}
\]
The first theorem combines the saturated \(H_J\)-source and saturated-Casimir
source callbacks into saturated joint source-vector callbacks, then applies
the joint-source outside-ground bridge.  The second first turns scalar
source-energy callbacks into the saturated \(H_J\)-source equations and then
uses the same bridge.  Consequently the discharged saturated eigen-source and
scalar source-energy final wrappers now keep their public hypotheses but pass
the resulting outside-sector lower family directly to the source predicted-GS
common-energy final boundary.
The uniform and left-endpoint named-callback wrappers for the same saturated
eigen-source and scalar source-energy inputs now reuse these two lower-family
bridges as well before entering
\(\texttt{tasaki\_2\_5\_theorem\_2\_3\_of\_named\_callbacks}\) or its
left-endpoint analogue.
Likewise,
\(\texttt{tasaki23OutsideGroundEnergyLowerFamilyCallback\_of\_saturated\_ladder\_span\_sources}\)
packages the conversion from saturated ladder-span sources to saturated
Casimir sources as an outside-ground lower-family bridge, so the uniform and
left-endpoint ladder-span named wrappers can enter the same common-energy
final boundaries directly.
The outside reach task is also split into its
two endpoint directions by
\[
\begin{gathered}
\texttt{tasaki23GroundStateSectors\_not\_mem\_iff\_lt\_left\_or\_right\_lt},\\
\texttt{tasaki23OutsideGroundLeftAdmissibleReachCallback},\\
\texttt{tasaki23OutsideGroundRightAdmissibleReachCallback},\\
\texttt{tasaki23OutsideGroundAdmissibleReachCallback\_of\_side\_callbacks}.
\end{gathered}
\]
The first statement is the arithmetic fact that a sector outside the
closed interval
\([\min(|A|,|\neg A|)N,\max(|A|,|\neg A|)N]\) is either strictly left of
the interval or strictly right of it.  The two side callbacks then name
the corresponding lowering and raising ladder-reach tasks, and the
composition wrapper recovers the full outside admissible-reach callback
by case splitting on that arithmetic alternative.  Lean also records the
direct side-input bridge
\[
\begin{gathered}
\texttt{tasaki23OutsideGroundEnergyLowerFamilyCallback}\\
\texttt{\_of\_side\_admissible\_reach}
\end{gathered}
\]
which first recombines the side callbacks and then invokes the
admissible-reach lower-family bridge.  The discharged final boundary
\[
\begin{gathered}
\texttt{tasaki\_2\_5\_theorem\_2\_3}\\
\texttt{\_of\_threaded\_predictedGS}\\
\texttt{\_of\_unpacked\_reembedded\_real\_source\_weight}\\
\texttt{\_predecessor\_difference\_pos}\\
\texttt{\_of\_side\_admissible\_reach}\\
\texttt{\_discharge\_nonempty}
\end{gathered}
\]
therefore obtains the outside-sector lower family through
\(\texttt{tasaki23OutsideGroundEnergyLowerFamilyCallback\_of\_side\_admissible\_reach}\)
and exposes the final proof boundary directly in terms of the two directional
outside-sector reach tasks.  The final named predecessor-difference
boundaries
\[
\begin{gathered}
\texttt{tasaki\_2\_5\_theorem\_2\_3\_of\_predecessor\_difference}\\
\texttt{\_named\_callbacks},\\
\texttt{tasaki\_2\_5\_theorem\_2\_3\_of\_left\_endpoint}\\
\texttt{\_predecessor\_difference\_named\_callbacks}
\end{gathered}
\]
use the sector-minimality bridge above to replace the outside-ground
lower-family hypothesis by
\texttt{tasaki23SectorMinimalityCallback}, leaving only predicted-GS
membership and the predecessor-difference local comparison as the other
named inputs.  The corresponding admissible-reach final boundaries
\[
\begin{gathered}
\texttt{tasaki\_2\_5\_theorem\_2\_3\_of\_admissible\_reach}\\
\texttt{\_named\_callbacks},\\
\texttt{tasaki\_2\_5\_theorem\_2\_3\_of\_left\_endpoint}\\
\texttt{\_admissible\_reach\_named\_callbacks}
\end{gathered}
\]
use
\texttt{tasaki23OutsideGroundEnergyLowerFamilyCallback\_of\_admissible\_reach}
instead, so the final public boundary can consume the ladder-reach
callback together with the real-coupling and strict diagonal-shift
hypotheses.  The non-discharged threaded layer, where admissible-sector
non-emptiness is still passed explicitly, also exposes the stronger reach
inputs
\[
\begin{gathered}
\texttt{tasaki\_2\_5\_theorem\_2\_3\_of\_threaded\_predictedGS}\\
\texttt{\_of\_unpacked\_reembedded\_real\_source\_weight}\\
\texttt{\_predecessor\_difference\_pos\_of\_full\_admissible\_reach},\\
\texttt{tasaki\_2\_5\_theorem\_2\_3\_of\_threaded\_predictedGS}\\
\texttt{\_of\_unpacked\_reembedded\_real\_source\_weight}\\
\texttt{\_predecessor\_difference\_pos\_of\_iterated\_ladder}\\
\texttt{\_full\_reach},\\
\texttt{tasaki\_2\_5\_theorem\_2\_3\_of\_threaded\_predictedGS}\\
\texttt{\_of\_unpacked\_reembedded\_real\_source\_weight}\\
\texttt{\_predecessor\_difference\_pos\_of\_iterated\_ladder}\\
\texttt{\_casimir\_full\_reach},\\
\texttt{tasaki\_2\_5\_theorem\_2\_3\_of\_threaded\_predictedGS}\\
\texttt{\_of\_unpacked\_reembedded\_real\_source\_weight}\\
\texttt{\_predecessor\_difference\_pos\_of\_side\_admissible\_reach}.
\end{gathered}
\]
These wrappers route through the same direct outside-ground lower-family
bridges before invoking the source predicted-GS common-energy final boundary.
The same explicit-nonempty threaded layer can also consume saturated source
data directly:
\[
\begin{gathered}
\texttt{tasaki\_2\_5\_theorem\_2\_3\_of\_threaded\_predictedGS}\\
\texttt{\_of\_unpacked\_reembedded\_real\_source\_weight}\\
\texttt{\_predecessor\_difference\_pos\_of\_saturated\_casimir}\\
\texttt{\_sources},\\
\texttt{tasaki\_2\_5\_theorem\_2\_3\_of\_threaded\_predictedGS}\\
\texttt{\_of\_unpacked\_reembedded\_real\_source\_weight}\\
\texttt{\_predecessor\_difference\_pos\_of\_saturated\_ladder}\\
\texttt{\_span\_sources},\\
\texttt{tasaki\_2\_5\_theorem\_2\_3\_of\_threaded\_predictedGS}\\
\texttt{\_of\_unpacked\_reembedded\_real\_source\_weight}\\
\texttt{\_predecessor\_difference\_pos\_of\_saturated\_joint}\\
\texttt{\_sources}.
\end{gathered}
\]
These wrappers keep the sector non-emptiness callback explicit and then use
the saturated-Casimir, saturated-ladder-span, or saturated-joint source
outside-ground lower-family bridge before entering the same source
predicted-GS common-energy final boundary.
The saturated eigen/source-energy inputs are available at the same layer:
\[
\begin{gathered}
\texttt{tasaki\_2\_5\_theorem\_2\_3\_of\_threaded\_predictedGS}\\
\texttt{\_of\_unpacked\_reembedded\_real\_source\_weight}\\
\texttt{\_predecessor\_difference\_pos\_of\_saturated\_eigen}\\
\texttt{\_sources},\\
\texttt{tasaki\_2\_5\_theorem\_2\_3\_of\_threaded\_predictedGS}\\
\texttt{\_of\_unpacked\_reembedded\_real\_source\_weight}\\
\texttt{\_predecessor\_difference\_pos\_of\_saturated\_source}\\
\texttt{\_energy}.
\end{gathered}
\]
They keep the same explicit sector non-emptiness callback while replacing the
outside-ground lower-family input by the saturated eigen-source or scalar
saturated source-energy bridge.
Sectorwise saturated references are also exposed without discharging sector
non-emptiness:
\[
\begin{gathered}
\texttt{tasaki\_2\_5\_theorem\_2\_3\_of\_threaded\_predictedGS}\\
\texttt{\_of\_unpacked\_reembedded\_real\_source\_weight}\\
\texttt{\_predecessor\_difference\_pos\_of\_saturated\_joint}\\
\texttt{\_references},\\
\texttt{tasaki\_2\_5\_theorem\_2\_3\_of\_threaded\_predictedGS}\\
\texttt{\_of\_unpacked\_reembedded\_real\_source\_weight}\\
\texttt{\_predecessor\_difference\_pos\_of\_saturated\_ladder}\\
\texttt{\_references},\\
\texttt{tasaki\_2\_5\_theorem\_2\_3\_of\_threaded\_predictedGS}\\
\texttt{\_of\_unpacked\_reembedded\_real\_source\_weight}\\
\texttt{\_predecessor\_difference\_pos\_of\_saturated\_ladder}\\
\texttt{\_iterate\_references}.
\end{gathered}
\]
These wrappers keep the same explicit sector non-emptiness callback while
using the saturated joint-reference, ladder-reference, or singleton
ladder-iterate reference bridge to build the outside-ground lower family.
The same explicit sector non-emptiness layer also reaches the
Marshall-positive ladder-iterate inputs:
\[
\begin{gathered}
\texttt{tasaki\_2\_5\_theorem\_2\_3\_of\_threaded\_predictedGS}\\
\texttt{\_of\_unpacked\_reembedded\_real\_source\_weight}\\
\texttt{\_predecessor\_difference\_pos\_of\_saturated\_ladder}\\
\texttt{\_iterate\_marshall\_positive\_references},\\
\texttt{tasaki\_2\_5\_theorem\_2\_3\_of\_threaded\_predictedGS}\\
\texttt{\_of\_unpacked\_reembedded\_real\_source\_weight}\\
\texttt{\_predecessor\_difference\_pos\_of\_saturated\_ladder}\\
\texttt{\_iterate\_marshall\_positive\_coefficients}.
\end{gathered}
\]
These wrappers use the direct Marshall-positive reference or coefficient
bridge to build the outside-ground lower family before entering the same
source predicted-GS common-energy final boundary.
The discharged variants
\[
\begin{gathered}
\texttt{tasaki\_2\_5\_theorem\_2\_3\_of\_threaded\_predictedGS}\\
\texttt{\_of\_unpacked\_reembedded\_real\_source\_weight}\\
\texttt{\_predecessor\_difference\_pos\_of\_saturated\_ladder}\\
\texttt{\_iterate\_marshall\_positive\_references}\\
\texttt{\_discharge\_nonempty},\\
\texttt{tasaki\_2\_5\_theorem\_2\_3\_of\_threaded\_predictedGS}\\
\texttt{\_of\_unpacked\_reembedded\_real\_source\_weight}\\
\texttt{\_predecessor\_difference\_pos\_of\_saturated\_ladder}\\
\texttt{\_iterate\_marshall\_positive\_coefficients}\\
\texttt{\_discharge\_nonempty}
\end{gathered}
\]
use the same direct Marshall-positive bridges but enter the source
common-energy boundary with admissible-sector non-emptiness discharged
internally.
The source-eigen extraction variants
\[
\begin{gathered}
\texttt{tasaki\_2\_5\_theorem\_2\_3\_of\_threaded\_predictedGS}\\
\texttt{\_of\_unpacked\_reembedded\_real\_source\_weight}\\
\texttt{\_predecessor\_difference\_pos\_of\_saturated\_ladder}\\
\texttt{\_iterate\_marshall\_positive\_references}\\
\texttt{\_extract\_source\_eigen\_callbacks\_discharge\_nonempty},\\
\texttt{tasaki\_2\_5\_theorem\_2\_3}\\
\texttt{\_of\_saturated\_ladder\_iterate\_marshall\_positive}\\
\texttt{\_references\_extract\_source\_eigen\_callbacks}\\
\texttt{\_named\_callbacks},\\
\texttt{tasaki\_2\_5\_theorem\_2\_3}\\
\texttt{\_of\_left\_endpoint\_saturated\_ladder\_iterate}\\
\texttt{\_marshall\_positive\_references}\\
\texttt{\_extract\_source\_eigen\_callbacks\_named\_callbacks}
\end{gathered}
\]
and their coefficient analogues first convert the concrete
Marshall-positive inputs to singleton ladder-iterate references and then
reuse the source-eigen extraction route.  Thus the scalar saturated
source-energy and saturated-Casimir callbacks remain internal for the
discharged, uniform named, and left-endpoint named final surfaces.
Under real couplings this extraction layer is also exposed directly at the
saturated ladder-iterate dominance level.  The wrappers with suffixes
\[
\begin{gathered}
\texttt{\_saturated\_ladder\_iterate\_successor\_dominance}\\
\texttt{\_of\_real\_couplings\_extract\_source\_eigen\_callbacks},\\
\texttt{\_saturated\_ladder\_iterate\_lowerable\_attach\_sum}\\
\texttt{\_dominance\_of\_real\_couplings}\\
\texttt{\_extract\_source\_eigen\_callbacks},\\
\texttt{\_saturated\_ladder\_iterate\_predecessor\_difference\_pos}\\
\texttt{\_of\_real\_couplings\_extract\_source\_eigen\_callbacks}
\end{gathered}
\]
are provided for the discharged, uniform named, and left-endpoint named final
surfaces.  They first build the left and right pointwise
Marshall-positive coefficient callbacks from the corresponding dominance
input, using real couplings to discharge the saturated-energy real-scalar
witness, and then reuse the coefficient source-eigen extraction boundary.
Consequently the public final statements expose the dominance input rather
than the intermediate coefficient or scalar source-eigen callbacks.
The same explicit sector non-emptiness layer reaches the saturated
ladder-iterate dominance inputs:
\[
\begin{gathered}
\texttt{tasaki\_2\_5\_theorem\_2\_3\_of\_threaded\_predictedGS}\\
\texttt{\_of\_unpacked\_reembedded\_real\_source\_weight}\\
\texttt{\_predecessor\_difference\_pos\_of\_saturated\_ladder}\\
\texttt{\_iterate\_successor\_dominance},\\
\texttt{tasaki\_2\_5\_theorem\_2\_3\_of\_threaded\_predictedGS}\\
\texttt{\_of\_unpacked\_reembedded\_real\_source\_weight}\\
\texttt{\_predecessor\_difference\_pos\_of\_saturated\_ladder}\\
\texttt{\_iterate\_lowerable\_attach\_sum\_dominance},\\
\texttt{tasaki\_2\_5\_theorem\_2\_3\_of\_threaded\_predictedGS}\\
\texttt{\_of\_unpacked\_reembedded\_real\_source\_weight}\\
\texttt{\_predecessor\_difference\_pos\_of\_saturated\_ladder}\\
\texttt{\_iterate\_predecessor\_difference\_pos},
\end{gathered}
\]
together with their \texttt{\_of\_real\_couplings} variants.  These wrappers
use the direct successor-dominance, lowerable attach-sum dominance, or
predecessor-difference positivity bridge before entering the same source
predicted-GS common-energy final boundary; the real-coupling variants also
remove the explicit saturated-energy scalar.
The named-callback layer has the corresponding Marshall-positive
ladder-iterate reference and coefficient entry points:
\[
\begin{gathered}
\texttt{tasaki\_2\_5\_theorem\_2\_3\_of\_saturated\_ladder}\\
\texttt{\_iterate\_marshall\_positive\_references\_named\_callbacks},\\
\texttt{tasaki\_2\_5\_theorem\_2\_3\_of\_saturated\_ladder}\\
\texttt{\_iterate\_marshall\_positive\_coefficients\_named\_callbacks},\\
\texttt{tasaki\_2\_5\_theorem\_2\_3\_of\_left\_endpoint}\\
\texttt{\_saturated\_ladder\_iterate\_marshall\_positive}\\
\texttt{\_references\_named\_callbacks},\\
\texttt{tasaki\_2\_5\_theorem\_2\_3\_of\_left\_endpoint}\\
\texttt{\_saturated\_ladder\_iterate\_marshall\_positive}\\
\texttt{\_coefficients\_named\_callbacks}.
\end{gathered}
\]
These wrappers pass the direct Marshall-positive reference or coefficient
outside-ground lower-family bridge to the uniform or left-endpoint named final
boundary.
The same named-callback layer also exposes the stronger
reach inputs
\[
\begin{gathered}
\texttt{tasaki\_2\_5\_theorem\_2\_3\_of\_full\_admissible\_reach}\\
\texttt{\_named\_callbacks},\\
\texttt{tasaki\_2\_5\_theorem\_2\_3\_of\_iterated\_ladder\_full\_reach}\\
\texttt{\_named\_callbacks},\\
\texttt{tasaki\_2\_5\_theorem\_2\_3\_of\_iterated\_ladder\_casimir}\\
\texttt{\_full\_reach\_named\_callbacks},\\
\texttt{tasaki\_2\_5\_theorem\_2\_3\_of\_side\_admissible\_reach}\\
\texttt{\_named\_callbacks}
\end{gathered}
\]
and their left-endpoint analogues.  These wrappers use the direct bridges
\[
\begin{gathered}
\texttt{tasaki23OutsideGroundEnergyLowerFamilyCallback}\\
\texttt{\_of\_full\_admissible\_reach},\\
\texttt{tasaki23OutsideGroundEnergyLowerFamilyCallback}\\
\texttt{\_of\_iterated\_ladder\_full\_reach},\\
\texttt{tasaki23OutsideGroundEnergyLowerFamilyCallback}\\
\texttt{\_of\_iterated\_ladder\_casimir\_full\_reach},\\
\texttt{tasaki23OutsideGroundEnergyLowerFamilyCallback}\\
\texttt{\_of\_side\_admissible\_reach}
\end{gathered}
\]
before reusing the uniform or left-endpoint named final boundary.  Thus the
named public API can state the remaining outside-ground task directly as
full-space reach, iterated-ladder reach, Casimir-certified ladder reach, or
left/right side reach.  These final boundaries live in
\(\texttt{LatticeSystem.Quantum.SpinS.Theorem23Final}\).  The lowered-site-sum
public final-boundary aliases now
live in the Lean module
\(\texttt{LatticeSystem.Quantum.SpinS.Theorem23FinalLoweredSiteSumLeft}\)
(split from
\(\texttt{LatticeSystem.Quantum.SpinS.Theorem23FinalLoweredSiteSum}\)), with the
lowered-vector-Marshall final suffix split into
\(\texttt{LatticeSystem.Quantum.SpinS.Theorem23FinalLoweredMarshall}\), while
\(\texttt{Theorem23PredictedEndpoint}\) carries the predicted-Casimir
endpoint mismatch bridges,
\(\texttt{Theorem23Predicted}\) carries the predicted-GS and cross-ladder
bridges, and
\(\texttt{Theorem23}\) keeps the adjacent common-energy chain.

Lean packages the resulting sectorwise minimality statement as
\[
\begin{gathered}
\texttt{tasaki23\_sector\_minimality\_of\_common\_energy\_chain}\\
\texttt{\_and\_outside\_sector\_ground\_energy\_lower\_bound}.
\end{gathered}
\]
On admissible sectors it uses the common-energy chain to obtain the
real-sector lower bound.  On outside sectors it applies the preceding
outside-sector ground-energy bridge.  Finally, it invokes the real-to-
complex sector bridge, so the remaining spectral input for the final
global-minimality argument is only the estimate \(\mu\le\mu_M\) for the
Marshall-positive Theorem 2.2 ground state in each outside sector.

The corresponding common-energy final wrapper is
\[
\begin{gathered}
\texttt{tasaki\_2\_5\_theorem\_2\_3\_of\_common\_energy\_chain}\\
\texttt{\_and\_outside\_sector\_ground\_energy\_lower\_bound}.
\end{gathered}
\]
Once the adjacent-sector chain has produced a common energy \(\mu\)
on the admissible interval, this wrapper uses the sector-minimality
package above and then the sector-to-global minimality bridge.  Thus
the final full-space theorem can be reached from the common-energy
chain while leaving only outside-sector estimates of the form
\(\mu\le\mu_M\) for the Marshall-positive representatives supplied by
the sector Theorem 2.2 wrapper.

The predecessor-difference lowered-site-sum boundary can now use this
outside-sector ground-energy form directly:
\[
\begin{gathered}
\texttt{tasaki\_2\_5\_theorem\_2\_3}\\
\texttt{\_of\_left\_endpoint\_threaded\_predictedGS}\\
\texttt{\_of\_unpacked\_reembedded\_real\_source\_weight}\\
\texttt{\_predecessor\_difference\_pos}\\
\texttt{\_via\_lowered\_site\_sum}\\
\texttt{\_of\_outside\_sector\_ground\_energy\_lower\_bound}.
\end{gathered}
\]
It first runs the same predecessor-difference interval chain to obtain
the common Marshall-positive energy on
\(\texttt{tasaki23GroundStateSectors}\).  The remaining outside-sector
input is then the ground-representative estimate
\(\mu\le\mu_M\); the already proved outside-sector bridge turns this into
the outside-real-sector minimality callback, and the previous final
wrapper completes the Tasaki Theorem 2.3 package.  This is still the
same Tasaki §2.5 Theorem 2.3, p.~42 Perron--Frobenius endpoint, with
the non-negative-matrix comparison following Seneta, \emph{Non-negative
Matrices and Markov Chains}, 3rd ed., Springer 2006, \S1.2,
pp.~27--28.
The corresponding admissible-reach predecessor-difference boundary
\[
\begin{gathered}
\texttt{tasaki\_2\_5\_theorem\_2\_3}\\
\texttt{\_of\_left\_endpoint\_threaded\_predictedGS}\\
\texttt{\_of\_unpacked\_reembedded\_real\_source\_weight}\\
\texttt{\_predecessor\_difference\_pos}\\
\texttt{\_of\_admissible\_reach}
\end{gathered}
\]
uses
\texttt{tasaki23OutsideGroundEnergyLowerFamilyCallback\_of\_admissible\_reach}
to supply that outside-sector ground-energy input from the
ladder-reach callback, together with the real-coupling and strict
diagonal-shift hypotheses.
The discharged companions
\[
\begin{gathered}
\texttt{tasaki\_2\_5\_theorem\_2\_3}\\
\texttt{\_of\_left\_endpoint\_threaded\_predictedGS}\\
\texttt{\_of\_unpacked\_reembedded\_real\_source\_weight}\\
\texttt{\_predecessor\_difference\_pos}\\
\texttt{\_of\_outside\_sector\_ground\_energy\_lower\_bound}\\
\texttt{\_discharge\_nonempty}
\end{gathered}
\]
and
\[
\begin{gathered}
\texttt{tasaki\_2\_5\_theorem\_2\_3}\\
\texttt{\_of\_left\_endpoint\_threaded\_predictedGS}\\
\texttt{\_of\_unpacked\_reembedded\_real\_source\_weight}\\
\texttt{\_predecessor\_difference\_pos}\\
\texttt{\_of\_admissible\_reach}\\
\texttt{\_discharge\_nonempty}
\end{gathered}
\]
remove the explicit admissible-sector non-emptiness callback from these
left-endpoint predecessor-difference boundaries.  The proof uses
\(\texttt{tasaki23GroundStateSectors\_le\_card\_mul}\) to place each
admissible sector in the physical range and then applies
\(\texttt{magConfigS\_nonempty\_of\_le\_card\_mul}\).  The mathematical
content is unchanged: the local input is still the predecessor-difference
comparison, and the outside-sector input is either the ground-energy lower
family or the admissible-reach callback.  References are Tasaki,
Springer 2020, \S2.5 Theorem 2.3, p.~42, and Seneta,
\emph{Non-negative Matrices and Markov Chains}, 3rd ed., Springer 2006,
\S1.2, pp.~27--28 for the non-negative-matrix comparison.

For the next proof stage Lean also exposes the common-energy output
itself, separated from the final outside-sector estimate:
\[
\begin{gathered}
\texttt{tasaki23\_common\_energy\_chain}\\
\texttt{\_of\_left\_endpoint\_predictedGS}\\
\texttt{\_of\_unpacked\_reembedded\_real\_source\_weight}\\
\texttt{\_predecessor\_difference\_pos}.
\end{gathered}
\]
This theorem invokes the predecessor-difference interval chain, obtains
one real energy \(\mu\) and Marshall-positive representatives in every
admissible sector, and then forgets the extra predicted-GS membership
field carried by the interval induction.  Its conclusion is exactly
\(\exists \mu,\texttt{tasaki23CommonEnergyChain}~A~J~N~c~\mu\), the first
of the two remaining callback outputs needed by the common-energy plus
outside-sector final wrapper.
The source-callback variants
\[
\begin{gathered}
\texttt{tasaki23\_common\_energy\_chain}\\
\texttt{\_of\_source\_predictedGS}\\
\texttt{\_of\_unpacked\_reembedded\_real\_source\_weight}\\
\texttt{\_predecessor\_difference\_pos}
\end{gathered}
\]
and its \(\texttt{\_discharge\_nonempty}\) variant keep the same
common-energy construction but accept the uniform
\(\texttt{tasaki23SourcePredictedGSCallback}\).  The first wrapper
derives the left-endpoint predicted-GS input from
\(\texttt{tasaki23GroundStateSectors\_left\_mem}\).  The discharged
variant also supplies admissible-sector non-emptiness from
\(\texttt{tasaki23GroundStateSectors\_le\_card\_mul}\) and
\(\texttt{magConfigS\_nonempty\_of\_le\_card\_mul}\), so the visible
inputs are the uniform predicted-GS callback and the
predecessor-difference callback.
Conversely, the source-callback closure
\[
\begin{gathered}
\texttt{tasaki23\_source\_predictedGSCallback}\\
\texttt{\_of\_left\_endpoint\_predictedGS}\\
\texttt{\_of\_unpacked\_reembedded\_real\_source\_weight}\\
\texttt{\_predecessor\_difference\_pos}
\end{gathered}
\]
derives the uniform source callback from the left-endpoint route itself.
The interval chain first supplies one predicted-GS representative in each
admissible sector.  The per-sector Theorem 2.2 uniqueness statement then
identifies any other Marshall-positive source eigenvector in that sector as
a positive real scalar multiple of the threaded representative, so
submodule closure under scalar multiplication transfers membership in
\(\texttt{bipartiteToyGroundStateSubspacePredicted}\).
The final boundary
\[
\begin{gathered}
\texttt{tasaki\_2\_5\_theorem\_2\_3}\\
\texttt{\_of\_source\_predictedGS\_common\_energy\_chain}\\
\texttt{\_and\_outside\_sector\_ground\_energy\_lower\_bound}\\
\texttt{\_discharge\_nonempty}
\end{gathered}
\]
uses this discharged source common-energy wrapper as the first half of
the reduced proof route.  It obtains
\(\exists\mu,\texttt{tasaki23CommonEnergyChain}~A~J~N~c~\mu\), applies
the outside-sector lower-family callback at that \(\mu\), and then
invokes
\(\texttt{tasaki\_2\_5\_theorem\_2\_3\_of\_common\_energy\_chain\_and\_outside\_sector\_ground\_energy\_lower\_bound}\).
Thus the public source predicted-GS route now factors through the same
common-energy plus outside-sector theorem used as the final reduced
boundary.
The companion without the \(\texttt{\_discharge\_nonempty}\) suffix keeps the
admissible-sector non-emptiness callback explicit, and the older uniform
threaded outside-ground entry point now calls this companion directly while
preserving its public hypotheses.
The physical-range non-empty outside-ground entry point now discharges this
non-emptiness callback and calls the same companion directly, rather than
passing through the older threaded wrapper.
The left-endpoint analogue
\[
\begin{gathered}
\texttt{tasaki\_2\_5\_theorem\_2\_3}\\
\texttt{\_of\_left\_endpoint\_predictedGS\_common\_energy\_chain}\\
\texttt{\_and\_outside\_sector\_ground\_energy\_lower\_bound}\\
\texttt{\_discharge\_nonempty}
\end{gathered}
\]
uses the same reduced final theorem but constructs
\(\exists\mu,\texttt{tasaki23CommonEnergyChain}~A~J~N~c~\mu\) from the
weaker left-endpoint predicted-GS callback.  The non-empty admissible-sector
input to the common-energy chain is again supplied from
\(\texttt{tasaki23GroundStateSectors\_le\_card\_mul}\) and
\(\texttt{magConfigS\_nonempty\_of\_le\_card\_mul}\), so the visible route
matches the left-endpoint named-callback boundary while making the
common-energy plus outside-sector factorisation explicit.
The short named-callback entry points
\(\texttt{tasaki\_2\_5\_theorem\_2\_3\_of\_named\_callbacks}\) and
\(\texttt{tasaki\_2\_5\_theorem\_2\_3\_of\_left\_endpoint\_named\_callbacks}\)
now reuse these two explicit common-energy final boundaries internally.
Thus the public callback API remains the same, but both named routes
construct the common interval energy first and then invoke the
outside-sector lower-family callback at that energy before applying the
reduced final theorem.
The discharged uniform outside-ground entry point
\(\texttt{tasaki\_2\_5\_theorem\_2\_3\_of\_threaded\_predictedGS\_}\cdots
\texttt{\_outside\_sector\_ground\_energy\_lower\_bound\_discharge\_nonempty}\)
now keeps the same public hypotheses while calling the source predicted-GS
common-energy final boundary directly.  Thus downstream discharged wrappers
also share the reduced common-energy plus outside-sector final route.
The admissible-reach entry point follows the same pattern: it converts the
reach callback to the outside-sector lower-family callback and feeds that
family directly into the explicit source common-energy final boundary.  The
physical-range and discharged reach wrappers now use the same reduced
common-energy plus outside-sector route without changing their public
hypotheses.

The common-energy boundary has also been pushed one step closer to the
local comparison now used in the predecessor-source wrappers:
\[
\begin{gathered}
\texttt{tasaki23\_common\_energy\_chain}\\
\texttt{\_of\_left\_endpoint\_predictedGS}\\
\texttt{\_of\_unpacked\_reembedded\_real\_source\_weight}\\
\texttt{\_predecessor\_raising\_source\_sum\_lt}.
\end{gathered}
\]
Instead of assuming the named predecessor-difference callback, this
theorem assumes the fully threaded dominance statement saying that the
on-\(A\) predecessor raising-source sum is strictly smaller than the
off-\(A\) sum.  Lean converts this inequality to the difference form
\(0<\text{offA}-\text{onA}\) by linear arithmetic and then invokes the
predecessor-difference common-energy wrapper above.  Thus the first
callback output is now available directly from the same raising-source
dominance shape used by the predecessor final boundary.

The lowerable-coefficient mirror is now available in the converse
dominance direction as well:
\[
\begin{gathered}
\texttt{tasaki23\_raising\_predecessor\_source\_sum\_lt}\\
\texttt{\_of\_lowerable\_positive\_source\_attach\_sum\_lt}.
\end{gathered}
\]
It rewrites the attached explicit lowerable positive-source coefficient
sums to the matching predecessor raising-source sums on both sublattices.
Consequently the common-energy endpoint can also be invoked from the
coefficient form
\[
\begin{gathered}
\texttt{tasaki23\_common\_energy\_chain}\\
\texttt{\_of\_left\_endpoint\_predictedGS}\\
\texttt{\_of\_unpacked\_reembedded\_real\_source\_weight}\\
\texttt{\_predecessor\_explicit\_lowerable}\\
\texttt{\_positive\_source\_coefficient\_lt}.
\end{gathered}
\]
This does not yet prove the local coefficient dominance itself; it moves
the first remaining callback output from the raising-source dominance
shape to the explicit lowerable coefficient shape used by the predecessor
source-weight wrappers.

The same explicit lowerable coefficient boundary has also been connected
to the outside-sector lower-bound family:
\[
\begin{gathered}
\texttt{tasaki\_2\_5\_theorem\_2\_3}\\
\texttt{\_of\_left\_endpoint\_threaded\_predictedGS}\\
\texttt{\_of\_unpacked\_reembedded\_real\_source\_weight}\\
\texttt{\_predecessor\_explicit\_lowerable}\\
\texttt{\_positive\_source\_coefficient\_lt}\\
\texttt{\_of\_outside\_sector\_ground\_energy\_lower\_bound}.
\end{gathered}
\]
This wrapper uses the common-energy chain supplied by the explicit
lowerable coefficient callback for the admissible sectors and delegates
the complementary sectors to
\texttt{tasaki23OutsideGroundEnergyLowerFamilyCallback}.  Thus the
sector-minimality callback no longer appears at this boundary; the
remaining visible mathematical inputs are the explicit lowerable local
dominance and the outside-sector ground-energy lower family.

The predecessor-difference callback now also feeds this
explicit-lowerable outside-sector route.  Lean records
\[
\begin{gathered}
\texttt{tasaki\_2\_5\_theorem\_2\_3}\\
\texttt{\_of\_left\_endpoint\_threaded\_predictedGS}\\
\texttt{\_of\_unpacked\_reembedded\_real\_source\_weight}\\
\texttt{\_predecessor\_difference\_pos}\\
\texttt{\_via\_explicit\_lowerable}\\
\texttt{\_of\_outside\_sector\_ground\_energy\_lower\_bound}.
\end{gathered}
\]
It converts the predecessor positive-difference callback to the attached
explicit lowerable coefficient dominance callback before invoking the
outside-sector ground-energy wrapper above.

Lean also records the same endpoint under the public boundary name
\[
\begin{gathered}
\texttt{tasaki\_2\_5\_theorem\_2\_3}\\
\texttt{\_of\_left\_endpoint\_threaded\_predictedGS}\\
\texttt{\_of\_unpacked\_reembedded\_real\_source\_weight}\\
\texttt{\_predecessor\_difference\_pos}\\
\texttt{\_of\_outside\_sector\_ground\_energy\_lower\_bound}.
\end{gathered}
\]
This alias hides the internal explicit-lowerable coefficient route in the declaration
name.  From the outside, the remaining final inputs are exactly the
left-endpoint predicted-GS membership, the local predecessor-difference
comparison, and the outside-sector ground-representative estimates
\(\mu\le\mu_M\).

The same outside-sector mechanism is also recorded one layer earlier,
before passing through the predecessor-difference callback.  Lean records
\[
\begin{gathered}
\texttt{tasaki23\_common\_energy\_chain}\\
\texttt{\_of\_left\_endpoint\_predictedGS}\\
\texttt{\_of\_unpacked\_lowered\_joint\_cross\_ladder}\\
\texttt{\_component\_lt}
\end{gathered}
\]
by forgetting the predicted-ground-state membership carried by the
unpacked cross-ladder interval chain.  The corresponding final wrapper
\[
\begin{gathered}
\texttt{tasaki\_2\_5\_theorem\_2\_3}\\
\texttt{\_of\_left\_endpoint\_threaded\_predictedGS}\\
\texttt{\_of\_unpacked\_lowered\_joint\_cross\_ladder}\\
\texttt{\_component\_lt}\\
\texttt{\_of\_outside\_sector\_ground\_energy\_lower\_bound}
\end{gathered}
\]
combines that common-energy chain with
\texttt{tasaki23OutsideGroundEnergyLowerFamilyCallback}.  This boundary
therefore proves the final theorem from the unpacked cross-ladder local
component comparison and the outside-sector lower family, without using
\texttt{tasaki23PredecessorDifferenceCallback}.

The coefficient form of the same outside-sector boundary is also available:
\[
\begin{gathered}
\texttt{tasaki\_2\_5\_theorem\_2\_3}\\
\texttt{\_of\_left\_endpoint\_threaded\_predictedGS}\\
\texttt{\_of\_unpacked\_lowered\_joint\_cross\_ladder}\\
\texttt{\_coefficient\_lt}\\
\texttt{\_of\_outside\_sector\_ground\_energy\_lower\_bound}.
\end{gathered}
\]
It first uses the existing coefficient wrapper to rewrite the signed
lowered-component comparison into strict dominance between the on-\(A\) and
off-\(A\) predecessor coefficient sums.  The common-energy/outside-sector
bridge is then applied exactly as in the component boundary.

The next boundary removes the special left-endpoint shape from the
public API.  Lean records
\[
\begin{gathered}
\texttt{tasaki\_2\_5\_theorem\_2\_3}\\
\texttt{\_of\_threaded\_predictedGS}\\
\texttt{\_of\_unpacked\_reembedded\_real\_source\_weight}\\
\texttt{\_predecessor\_difference\_pos}\\
\texttt{\_of\_outside\_sector\_ground\_energy\_lower\_bound}.
\end{gathered}
\]
Its predicted-GS hypothesis is uniform in every source sector
\(M\in\mathcal G(A,N)\).  The wrapper simply evaluates that uniform
callback at the left endpoint
\(\min(|A|,|A^c|)N\), already known to belong to
\(\mathcal G(A,N)\), and then invokes the public outside-ground
predecessor-difference boundary above.  Thus the visible remaining local
work is no longer tied to a bespoke left-endpoint predicted-GS callback.

Lean also records the range step needed to remove the
admissible-interval-specific non-emptiness callback.  The lemma
\[
  \texttt{tasaki23GroundStateSectors\_le\_card\_mul}
\]
in \texttt{LatticeSystem/Quantum/SpinS/Theorem23Sectors.lean}
states that if \(M\in\mathcal G(A,N)\), then
\[
  M \le |V|N.
\]
Indeed \(M\le \max(|A|,|A^c|)N\), while
\(\max(|A|,|A^c|)\le |A|+|A^c|=|V|\).  The boundary
\[
\begin{gathered}
\texttt{tasaki\_2\_5\_theorem\_2\_3}\\
\texttt{\_of\_physical\_range\_nonempty\_threaded\_predictedGS}\\
\texttt{\_of\_unpacked\_reembedded\_real\_source\_weight}\\
\texttt{\_predecessor\_difference\_pos}\\
\texttt{\_of\_outside\_sector\_ground\_energy\_lower\_bound}
\end{gathered}
\]
therefore asks only for non-emptiness of \(\texttt{magConfigS}\) on
the physical range \(M\le |V|N\), and derives the previous
admissible-sector non-emptiness callback from this range lemma.

The physical-range non-emptiness input is itself discharged by
\[
  \texttt{magConfigS\_nonempty\_of\_le\_card\_mul}.
\]
For \(M\le |V|N\), choose a subset of cardinality \(M\) in the
elementary slot set \(V\times\{0,\ldots,N-1\}\).  At each site \(x\),
count the chosen slots over \(x\); this count is at most \(N\), hence
gives a value in \(\texttt{Fin}(N+1)\).  Summing these fiber
cardinalities over \(V\) recovers the cardinality \(M\) of the chosen
slot set.  The wrapper
\[
\begin{gathered}
\texttt{tasaki\_2\_5\_theorem\_2\_3}\\
\texttt{\_of\_threaded\_predictedGS}\\
\texttt{\_of\_unpacked\_reembedded\_real\_source\_weight}\\
\texttt{\_predecessor\_difference\_pos}\\
\texttt{\_of\_outside\_sector\_ground\_energy\_lower\_bound}\\
\texttt{\_discharge\_nonempty}
\end{gathered}
\]
therefore removes the explicit `\texttt{magConfigS}` non-emptiness
callback from this boundary.
The same three public boundary levels also have admissible-reach versions:
\[
\begin{gathered}
\texttt{tasaki\_2\_5\_theorem\_2\_3}\\
\texttt{\_of\_threaded\_predictedGS}\\
\texttt{\_of\_unpacked\_reembedded\_real\_source\_weight}\\
\texttt{\_predecessor\_difference\_pos}\\
\texttt{\_of\_admissible\_reach},\\
\texttt{tasaki\_2\_5\_theorem\_2\_3}\\
\texttt{\_of\_physical\_range\_nonempty\_threaded\_predictedGS}\\
\texttt{\_of\_unpacked\_reembedded\_real\_source\_weight}\\
\texttt{\_predecessor\_difference\_pos}\\
\texttt{\_of\_admissible\_reach},\\
\texttt{tasaki\_2\_5\_theorem\_2\_3}\\
\texttt{\_of\_threaded\_predictedGS}\\
\texttt{\_of\_unpacked\_reembedded\_real\_source\_weight}\\
\texttt{\_predecessor\_difference\_pos}\\
\texttt{\_of\_admissible\_reach}\\
\texttt{\_discharge\_nonempty}.
\end{gathered}
\]
Each wrapper obtains the outside-sector ground-energy lower family from
\texttt{tasaki23OutsideGroundEnergyLowerFamilyCallback\_of\_admissible\_reach}
using the real-coupling and strict diagonal-shift hypotheses, and then
passes it directly to the corresponding source common-energy final boundary.

\textbf{Lowered vector Marshall-positivity final wrappers}: the
site-sum formulation is only a technical expansion of the total lowering
operator.  Lean therefore also records the reverse conversion
\[
\begin{gathered}
\texttt{tasaki23\_lowered\_site\_sum\_pos\_of\_marshall\_pos},\\
\texttt{tasaki23\_raised\_site\_sum\_pos\_of\_marshall\_pos}.
\end{gathered}
\]
For example, if the lowered full-space vector satisfies
\[
  0 <
  \operatorname{Re}(\operatorname{MarshallSign}_A(\tau))\,
  \operatorname{Re}\bigl((S^-_{\mathrm{tot}}\Phi)(\tau)\bigr)
\]
for every target-sector configuration \(\tau\), then the expansion
\[
  (S^-_{\mathrm{tot}}\Phi)(\tau)
    = \sum_x (S^-_x\Phi)(\tau)
\]
immediately gives the strict lowered site-sum positivity used by the
successor step.  The interval chain
\[
\texttt{tasaki23\_energy\_interval\_chain\_of\_left\_endpoint}
\texttt{\_predictedGS\_of\_lowered\_marshall\_pos}
\]
and the final wrappers
\[
\begin{gathered}
\texttt{tasaki\_2\_5\_theorem\_2\_3}\\
\texttt{\_of\_predictedGS\_of\_lowered\_marshall\_pos},\\
\texttt{tasaki\_2\_5\_theorem\_2\_3}\\
\texttt{\_of\_predictedGS\_of\_lowered\_marshall\_pos}\\
\texttt{\_of\_sector\_minimality},\\
\texttt{tasaki\_2\_5\_theorem\_2\_3}\\
\texttt{\_of\_predictedGS\_of\_lowered\_marshall\_pos}\\
\texttt{\_of\_real\_sector\_minimality}
\end{gathered}
\]
therefore take the natural remaining mathematical input, namely
Marshall positivity of \(S^-_{\mathrm{tot}}\Phi\), and hide the
site-sum conversion internally.

\textbf{Left-endpoint threading for lowered Marshall positivity}: the
lowered site-sum successor step also transfers the predicted total-Casimir
identity to the successor representative.  Lean records this as
\[
\texttt{tasaki23\_successor\_sector\_common\_energy}
\texttt{\_with\_successor\_predictedCasimir\_of\_site\_sum\_pos}.
\]
The proof first obtains the usual successor package from the direct
site-sum step.  The lowered vector has the predicted Casimir by ladder
preservation, and the successor vector is a nonzero positive real scalar
multiple of that lowered vector in Marshall coordinates; scalar cancellation
therefore gives the successor predicted-Casimir identity.  Consequently the
interval chains
\[
\begin{gathered}
\texttt{tasaki23\_energy\_interval\_chain\_with\_predictedCasimir}\\
\texttt{\_of\_left\_endpoint\_predictedCasimir\_of\_lowered\_site\_sum\_pos},\\
\texttt{tasaki23\_energy\_interval\_chain\_with\_predictedCasimir}\\
\texttt{\_of\_left\_endpoint\_predictedCasimir\_of\_lowered\_marshall\_pos}
\end{gathered}
\]
need the predicted-Casimir callback only at the left endpoint.  Under
\(|A^c|\le |A|\), left-endpoint membership in
\(\mathcal G_{\mathrm{pred}}(A,N)\) supplies that Casimir identity, giving
the final wrappers
\[
\begin{gathered}
\texttt{tasaki\_2\_5\_theorem\_2\_3}\\
\texttt{\_of\_left\_endpoint\_threaded\_predictedGS\_of\_lowered\_marshall\_pos},\\
\texttt{tasaki\_2\_5\_theorem\_2\_3}\\
\texttt{\_of\_left\_endpoint\_threaded\_predictedGS\_of\_lowered\_marshall\_pos}\\
\texttt{\_of\_sector\_minimality},\\
\texttt{tasaki\_2\_5\_theorem\_2\_3}\\
\texttt{\_of\_left\_endpoint\_threaded\_predictedGS\_of\_lowered\_marshall\_pos}\\
\texttt{\_of\_real\_sector\_minimality}.
\end{gathered}
\]
These wrappers remove the all-source predicted-GS membership callback from
the lowered-Marshall formulation; the remaining local hypothesis is the
pointwise Marshall positivity of \(S^-_{\mathrm{tot}}\Phi\).

\textbf{Predicted-GS membership threading}: the same scalar relation also
threads actual membership in the predicted toy ground-state subspace.  If
\(\Psi_M\in\mathcal G_{\mathrm{pred}}(A,N)\), then the already formalized
ladder invariance gives
\[
  S^-_{\mathrm{tot}}\Psi_M\in\mathcal G_{\mathrm{pred}}(A,N).
\]
The adjacent-sector uniqueness step identifies this lowered image with a
positive real scalar multiple of the successor representative:
\[
  S^-_{\mathrm{tot}}\Psi_M = r\,\Psi_{M+1},\qquad r>0.
\]
Since \(\mathcal G_{\mathrm{pred}}(A,N)\) is a complex subspace,
scalar cancellation gives \(\Psi_{M+1}\in
\mathcal G_{\mathrm{pred}}(A,N)\).  Lean records the scalar-cancellation
lemma as
\[
\texttt{tasaki23\_mem\_bipartiteToyGroundStateSubspacePredicted}
\texttt{\_of\_real\_smul\_eq}
\]
and the successor package as
\[
\texttt{tasaki23\_successor\_sector\_common\_energy}
\texttt{\_with\_successor\_predictedGS\_of\_site\_sum\_pos}.
\]
Consequently the interval chains
\[
\begin{gathered}
\texttt{tasaki23\_energy\_interval\_chain\_with\_predictedGS}\\
\texttt{\_of\_left\_endpoint\_predictedGS\_of\_lowered\_site\_sum\_pos},\\
\texttt{tasaki23\_energy\_interval\_chain\_with\_predictedGS}\\
\texttt{\_of\_left\_endpoint\_predictedGS\_of\_lowered\_marshall\_pos}
\end{gathered}
\]
carry predicted-GS membership from the left endpoint through the whole
admissible interval.  This is still conditional on the local lowered
site-sum or lowered-Marshall positivity input, exactly as in Tasaki,
Springer 2020, Section 2.5, Theorem 2.3, p. 42.

\textbf{Real-sector interval lower bound}: the common-energy chain also
discharges the real-sector spectral lower bound on the admissible
interval.  The matrix input is the elementary Perron comparison: if a
non-negative symmetric real matrix has a strictly positive eigenvector
with eigenvalue \(r\), then every real eigenvalue \(s\) satisfies
\(s\le r\).  In Lean this is
\[
\texttt{eigenvalue\_le\_of\_nonnegative\_isSymm\_of\_positive\_eigenvec}.
\]
Applying it to the shifted dressed sector matrix
\(cI-H_{\mathrm{dress},M}\) gives
\[
\texttt{shiftedDressedSReMatrixOnMagSector\_eigenvalue\_le\_of\_positive\_eigenvec}
\]
and, after undoing the Marshall sign conjugation,
\[
\texttt{heisenbergHamiltonianSReMatrixOnMagSector\_eigenvalue\_ge\_of\_marshallPositive}.
\]
Thus any Marshall-positive representative at the common energy \(\mu\)
proves \(\mu\le\mu'\) for every real-form eigenvector in the same
admissible sector.  The theorem
\[
\texttt{tasaki23\_real\_sector\_minimality\_on\_groundStateSectors\_of\_common\_energy\_chain}
\]
packages this for all \(M\in\texttt{tasaki23GroundStateSectors}\).
The final lowered-Marshall wrapper
\[
\texttt{tasaki\_2\_5\_theorem\_2\_3\_of\_left\_endpoint\_threaded\_predictedGS}\\
\texttt{\_of\_lowered\_marshall\_pos\_of\_outside\_real\_sector\_minimality}
\]
therefore leaves only the sectors outside the admissible interval as the
remaining real-sector spectral input.  This is the finite-dimensional
Perron-Frobenius comparison used in Tasaki, Springer 2020,
Section 2.5, Theorem 2.3, p. 42; for the non-negative-matrix
background see Seneta, \emph{Non-negative Matrices and Markov Chains},
3rd ed., Springer 2006, Section 1.2, pp. 27--28.

\textbf{Lowered Marshall coefficient identities}: the remaining local
input is the strict Marshall positivity of the lowered ladder image.
For this, the single-site predecessor formula is now recorded in its
Marshall-signed coefficient form.  If \(\tau\in H_{M+1}\) and
\(\tau^{x-}\in H_M\) is obtained by lowering \(\tau_x\), then the real
single-site contribution satisfies
\[
  \operatorname{sgn}_A(\tau)\,
    \operatorname{Re}\bigl((S_x^-\Phi)_\tau\bigr)
  =
  \begin{cases}
    c_x(\tau)\,
      \operatorname{sgn}_A(\tau^{x-})
      \operatorname{Re}\Phi(\tau^{x-}), & A(x)=\mathrm{false},\\
    -c_x(\tau)\,
      \operatorname{sgn}_A(\tau^{x-})
      \operatorname{Re}\Phi(\tau^{x-}), & A(x)=\mathrm{true}.
  \end{cases}
\]
Here \(c_x(\tau)>0\) is the real lowering matrix coefficient.  Lean
records the two cases as
\[
\texttt{tasaki23\_signed\_single\_site\_lowering\_component\_eq\_of\_A\_false}
\]
and
\[
\texttt{tasaki23\_signed\_single\_site\_lowering\_component\_eq\_neg\_of\_A\_true}.
\]
These identities do not by themselves prove positivity; they isolate the
coefficient-level sublattice sign relation that the final strict
site-sum dominance argument has to sum.  This is the local ladder
coefficient analysis used in Tasaki, Springer 2020, Section 2.5,
Theorem 2.3, p. 42.

\textbf{Boundary-inclusive lowered coefficient split}: the same
coefficient relation is now packaged in a form suitable for finite sums.
Define \(C_x(\tau)\) to be the preceding coefficient when
\(\tau_x>0\), and \(0\) when \(\tau_x=0\).  Lean names this
\[
\texttt{tasaki23LoweringPredecessorSignedCoefficient}.
\]
Then
\[
\texttt{tasaki23\_signed\_lowering\_site\_contribution\_eq\_coefficient\_of\_A\_false}
\]
states that the signed \(x\)-contribution is \(C_x(\tau)\) for
\(x\notin A\), while
\[
\texttt{tasaki23\_signed\_lowering\_site\_contribution\_eq\_neg\_coefficient\_of\_A\_true}
\]
states that it is \(-C_x(\tau)\) for \(x\in A\).  Summing these
pointwise identities gives
\[
\texttt{tasaki23\_signed\_lowering\_offA\_sum\_eq\_coefficient\_sum}
\]
and
\[
\texttt{tasaki23\_signed\_lowering\_onA\_sum\_eq\_neg\_coefficient\_sum}.
\]
Thus the remaining lowered-site-sum positivity input has been reduced
to a coefficient-level comparison between the two sublattice sums,
still within the ladder step of Tasaki, Springer 2020, Section 2.5,
Theorem 2.3, p. 42.

For the Marshall-signed real source vector
\(\Phi_\sigma=\operatorname{sgn}_A(\sigma)v_\sigma\), the coefficient
\(C_x(\tau)\) has a sign-normalized form.  At a lowerable site the two
Marshall signs cancel, and Lean records
\[
\texttt{tasaki23\_lowering\_predecessor\_signed\_coefficient\_eq\_positive\_source}
\]
as
\[
  C_x(\tau)=c_x(\tau)v_{\tau^{x-}} .
\]
Since \(c_x(\tau)>0\), the companion lemmas
\[
\texttt{tasaki23\_lowering\_predecessor\_signed\_coefficient\_pos\_of\_lowerable}
\]
and
\[
\texttt{tasaki23\_lowering\_predecessor\_signed\_coefficient\_nonneg}
\]
show strict positivity at lowerable sites and non-negativity in the
boundary-inclusive convention.

Lean also packages this sign-normalized quantity as
\[
\texttt{tasaki23LoweringPredecessorPositiveSourceCoefficient}.
\]
The bridge
\[
\texttt{tasaki23\_lowering\_predecessor\_signed\_coefficient\_eq\_positive\_source\_coefficient}
\]
identifies the original signed predecessor coefficient for the
Marshall-signed source vector with this positive-source coefficient, and
\[
\texttt{tasaki23\_lowering\_predecessor\_coefficient\_sum\_eq\_positive\_source\_sum}
\]
does the same for finite site sums.  Thus
\[
\texttt{tasaki23\_signed\_lowering\_site\_sum\_pos\_of\_positive\_source\_coefficient\_lt}
\]
and
\[
\texttt{tasaki23\_lowered\_marshall\_pos\_of\_positive\_source\_coefficient\_lt}
\]
allow the remaining dominance callback to be stated directly in terms of
the positive real source coefficients \(v_{\tau^{x-}}\).
Lean first names the explicit lowerable branch and records that the
boundary-inclusive coefficient agrees with it at lowerable sites:
\[
\texttt{tasaki23\_positive\_source\_coefficient\_eq\_lowerable\_coefficient}.
\]
Since the boundary-inclusive coefficient is zero when \(\tau_x=0\), Lean
also records the finite-sum lowerable-site restriction
\[
\texttt{tasaki23\_positive\_source\_coefficient\_sum\_eq\_lowerable\_sum}.
\]
Together with
\[
\begin{gathered}
\texttt{tasaki23\_signed\_lowering\_site\_sum\_pos}\\
\texttt{\_of\_positive\_source\_lowerable\_coefficient\_lt}
\end{gathered}
\]
and
\[
\begin{gathered}
\texttt{tasaki23\_lowered\_marshall\_pos}\\
\texttt{\_of\_positive\_source\_lowerable\_coefficient\_lt},
\end{gathered}
\]
the remaining comparison may be stated after filtering both sublattice
sums to the lowerable sites of the successor configuration.
The final callback boundary is correspondingly threaded by
\[
\begin{gathered}
\texttt{tasaki\_2\_5\_theorem\_2\_3\_of\_left\_endpoint\_threaded\_predictedGS}\\
\texttt{\_of\_unpacked\_reembedded\_source\_weight\_predecessor}\\
\texttt{\_lowerable\_positive\_source\_coefficient\_lt\_of\_sector\_minimality}.
\end{gathered}
\]
This wrapper keeps the predecessor source-weight identity, the four
sublattice-Casimir equations, and the two successor-sector support facts,
but asks the remaining local estimate only for the lowerable-site
positive-source coefficient comparison.
Its outside-sector companion
\[
\begin{gathered}
\texttt{tasaki\_2\_5\_theorem\_2\_3\_of\_left\_endpoint\_threaded\_predictedGS}\\
\texttt{\_of\_unpacked\_reembedded\_source\_weight\_predecessor}\\
\texttt{\_lowerable\_positive\_source\_coefficient\_lt}\\
\texttt{\_of\_outside\_sector\_ground\_energy\_lower\_bound}
\end{gathered}
\]
combines the lowerable wrapper with
\texttt{tasaki23OutsideGroundEnergyLowerFamilyCallback}.  Sector minimality
is supplied by the same common-energy/outside-sector bridge, while the
remaining local comparison is now lowerable-site filtered.
Lean then removes the last boundary-inclusive notation from that
callback.  The finite-sum identity
\[
\texttt{tasaki23\_positive\_source\_lowerable\_filter\_sum\_eq\_lowerable\_attach\_sum}
\]
rewrites each lowerable-filtered coefficient sum as an attached sum of
the explicit branch
\(\texttt{tasaki23LoweringPredecessorPositiveSourceLowerableCoefficient}\).
The final wrapper
\[
\begin{gathered}
\texttt{tasaki\_2\_5\_theorem\_2\_3\_of\_left\_endpoint\_threaded\_predictedGS}\\
\texttt{\_of\_unpacked\_reembedded\_source\_weight\_predecessor}\\
\texttt{\_explicit\_lowerable\_positive\_source\_coefficient\_lt\_of\_sector\_minimality}
\end{gathered}
\]
therefore exposes the remaining comparison directly as a strict
inequality between explicit lowerable predecessor coefficient sums.
Its outside-sector companion
\[
\begin{gathered}
\texttt{tasaki\_2\_5\_theorem\_2\_3\_of\_left\_endpoint\_threaded\_predictedGS}\\
\texttt{\_of\_unpacked\_reembedded\_source\_weight\_predecessor}\\
\texttt{\_explicit\_lowerable\_positive\_source\_coefficient\_lt}\\
\texttt{\_of\_outside\_sector\_ground\_energy\_lower\_bound}
\end{gathered}
\]
uses the same outside-sector bridge after the explicit lowerable wrapper has
turned the attached coefficient comparison into the lowerable-filtered
positive-source comparison.  The remaining outside-sector input is again
\texttt{tasaki23OutsideGroundEnergyLowerFamilyCallback}.
After the real predecessor source-weight identity is available, Lean
also records the final boundary
\[
\begin{gathered}
\texttt{tasaki\_2\_5\_theorem\_2\_3\_of\_left\_endpoint\_threaded\_predictedGS}\\
\texttt{\_of\_unpacked\_reembedded\_real\_source\_weight\_predecessor}\\
\texttt{\_explicit\_lowerable\_positive\_source\_coefficient\_lt\_of\_sector\_minimality}.
\end{gathered}
\]
It converts the complex predecessor source-weight identity to the real
equality above before invoking the explicit lowerable callback.  Thus
the remaining local comparison receives the real filtered raising-sum
identity, the four maximum sublattice-Casimir equations, the two
successor-sector support facts, and the attached lowerable
positive-source coefficient sums directly.
The coefficient bridge
\[
\begin{gathered}
\texttt{tasaki23\_spinSOpPlus\_lowering\_predecessor}\\
\texttt{\_re\_eq\_spinSOpMinus}
\end{gathered}
\]
identifies the real raising coefficient from the lowering predecessor
back to the successor with the real lowering coefficient used in the
explicit lowerable predecessor coefficient.  Consequently
\[
\begin{gathered}
\texttt{tasaki23\_lowerable\_positive\_source\_coefficient}\\
\texttt{\_eq\_raising\_predecessor\_source}
\end{gathered}
\]
rewrites the attached coefficient summand as the matching predecessor
raising coefficient times the positive predecessor source value.  This
aligns the explicit lowerable sums with the real source-weight
raising sums.
Lean also records the attached finite-sum version
\[
\begin{gathered}
\texttt{tasaki23\_lowerable\_positive\_source\_attach\_sum}\\
\texttt{\_eq\_raising\_predecessor\_source\_sum}
\end{gathered}
\]
and its strict-dominance wrapper
\[
\begin{gathered}
\texttt{tasaki23\_lowerable\_positive\_source\_attach\_sum\_lt}\\
\texttt{\_of\_raising\_predecessor\_source\_sum\_lt}.
\end{gathered}
\]
These theorems move the remaining local estimate from explicit
lowerable coefficient notation to the predecessor raising-source sums
that occur in the real source-weight identity.
Lean also records the sign information for these sums:
\[
\begin{gathered}
\texttt{tasaki23\_raising\_predecessor\_source\_summand\_pos},\\
\texttt{tasaki23\_raising\_predecessor\_source\_attach\_sum\_nonneg}.
\end{gathered}
\]
For every lowerable successor site, the real \(S^+\) coefficient from
the lowering predecessor back to the successor is strictly positive, and
the source coefficient \(v(\tau^{x-})\) is strictly positive by the
Marshall-positive real-sector representative.  Therefore every
predecessor raising-source summand is positive, and the attached sums
appearing in the remaining dominance callback are non-negative.
For the finite-set partition step, Lean introduces the boundary-inclusive
coefficient
\[
\texttt{tasaki23RaisingPredecessorSourceCoefficient},
\]
which agrees with the lowerable predecessor raising-source summand when
\(\tau_x>0\) and is zero otherwise.  The bridge
\[
\begin{gathered}
\texttt{tasaki23\_raising\_predecessor\_source\_attach\_sum}\\
\texttt{\_eq\_boundary\_sum}
\end{gathered}
\]
rewrites the attached lowerable-site sum as an ambient finite-set sum of
this boundary coefficient.  The partition theorem
\[
\begin{gathered}
\texttt{tasaki23\_raising\_predecessor\_source\_univ\_attach\_sum}\\
\texttt{\_eq\_onA\_add\_offA}
\end{gathered}
\]
then splits the total predecessor raising-source sum into the on-\(A\)
and off-\(A\) attached sums.  This is the Boolean finite-set partition
needed before the real predecessor source-weight identity is applied to
the two sublattice raising-source sums.
The local comparison is also recorded in the difference form
\[
\begin{gathered}
\texttt{tasaki23\_raising\_predecessor\_source\_sum\_lt}\\
\texttt{\_of\_offA\_sub\_onA\_pos}.
\end{gathered}
\]
It converts positivity of the off-\(A\) attached predecessor
raising-source sum minus the on-\(A\) attached sum into the strict
dominance inequality required by the final callback.  Its quantified
callback form
\[
\begin{gathered}
\texttt{tasaki23\_raising\_predecessor\_source\_sum\_lt\_callback}\\
\texttt{\_of\_offA\_sub\_onA\_pos}
\end{gathered}
\]
passes this conversion pointwise for every successor configuration.
This is the form that matches the real predecessor source-weight
identity before the remaining sign estimate is discharged.
Composing this bridge with the lowerable coefficient mirror also gives
\[
\begin{gathered}
\texttt{tasaki23\_lowerable\_positive\_source\_attach\_sum\_lt}\\
\texttt{\_of\_offA\_sub\_onA\_pos}
\end{gathered}
\]
and its callback form.  Hence a positive off-\(A\) minus on-\(A\)
predecessor raising-source difference supplies the same strict dominance
of attached explicit lowerable positive-source coefficient sums that the
explicit-lowerable boundary expects.
Lean also identifies the boundary-sum version of this difference with
the actual full lowering component:
\[
\begin{gathered}
\texttt{tasaki23\_raising\_predecessor\_source\_difference}\\
\texttt{\_eq\_lowered\_marshall\_component}.
\end{gathered}
\]
For the Marshall-signed representative
\(\Psi_\sigma=\operatorname{sgn}_A(\sigma)v_\sigma\), the theorem
rewrites
\[
  \sum_{x\notin A} R_x(\tau)-\sum_{x\in A}R_x(\tau)
\]
as
\[
  \operatorname{sgn}_A(\tau)\,
  \operatorname{Re}\bigl((\hat S^-_{\rm tot}\Psi)_\tau\bigr).
\]
The proof combines the on/off sublattice lowering-component coefficient
splits with the positive-source and raising-source sum bridges.  Thus
the difference-form callback is now linked directly to the
operator-level lowered-vector positivity required in Tasaki's adjacent
sector step.
The companion bridge
\[
\begin{gathered}
\texttt{tasaki23\_lowered\_marshall\_pos}\\
\texttt{\_of\_raising\_predecessor\_source\_difference\_pos}
\end{gathered}
\]
uses the lowerable-to-boundary sum conversion on both sublattices.
Hence the already isolated positive difference callback directly proves
\[
  0<
  \operatorname{sgn}_A(\tau)\,
  \operatorname{Re}\bigl((\hat S^-_{\rm tot}\Psi)_\tau\bigr),
\]
which is exactly the lowered-sector Marshall-positivity hypothesis in
the adjacent-sector energy comparison.
Lean also records the fully threaded callback form
\[
\begin{gathered}
\texttt{tasaki23\_lowered\_marshall\_pos}\\
\texttt{\_of\_unpacked\_reembedded\_real\_source\_weight}\\
\texttt{\_predecessor\_difference\_pos}.
\end{gathered}
\]
This version keeps the predicted-GS membership, real predecessor
source-weight equality, the four lowered-component sublattice-Casimir
equations, and the two successor-sector support facts in the hypothesis.
Once those data produce the positive off-\(A\) minus on-\(A\)
predecessor raising-source difference, the callback immediately returns
the lowered-vector Marshall-positivity component.
The source-specialized bridge
\[
\texttt{tasaki23\_lowered\_site\_sum\_pos\_of\_source\_lowered\_marshall\_pos}
\]
then expands this lowered total-ladder positivity back into the
single-site lowered sum positivity used by the adjacent-sector chain.
The lowered-Marshall interval wrappers call this source-form bridge
directly for the Marshall-signed positive real representatives that occur
in the sector chain.
The bridge
\[
\texttt{tasaki23\_lowered\_site\_sum\_pos\_of\_raising\_predecessor\_source\_difference\_pos}
\]
packages the same composition directly from the predecessor
raising-source difference to the site-sum callback.
Its fully threaded counterpart
\[
\begin{gathered}
\texttt{tasaki23\_lowered\_site\_sum\_pos}\\
\texttt{\_of\_unpacked\_reembedded\_real\_source\_weight}\\
\texttt{\_predecessor\_difference\_pos}
\end{gathered}
\]
keeps the same predicted-GS, real source-weight, sublattice-Casimir, and
successor-support hypotheses as the lowered-Marshall callback, but
returns the strict single-site lowered sum positivity directly.
The corresponding final wrapper
\[
\begin{gathered}
\texttt{tasaki\_2\_5\_theorem\_2\_3\_of\_left\_endpoint\_threaded\_predictedGS}\\
\texttt{\_of\_unpacked\_reembedded\_real\_source\_weight\_predecessor}\\
\texttt{\_raising\_source\_sum\_lt\_of\_sector\_minimality}
\end{gathered}
\]
therefore asks the remaining callback directly for strict dominance of
the two predecessor raising-source sums.  It converts that dominance
back to the explicit lowerable coefficient comparison with
\(\texttt{tasaki23\_lowerable\_positive\_source\_attach\_sum\_lt\_of\_raising\_predecessor\_source\_sum\_lt}\),
while keeping the real predecessor source-weight equality and the
lowered-component Casimir/support data available to the callback.
The difference-form final wrapper
\[
\begin{gathered}
\texttt{tasaki\_2\_5\_theorem\_2\_3\_of\_left\_endpoint\_threaded\_predictedGS}\\
\texttt{\_of\_unpacked\_reembedded\_real\_source\_weight\_predecessor}\\
\texttt{\_difference\_pos\_of\_sector\_minimality}
\end{gathered}
\]
asks instead for positivity of the off-\(A\) minus on-\(A\) predecessor
raising-source difference.  It applies
\(\texttt{tasaki23\_raising\_predecessor\_source\_sum\_lt\_callback\_of\_offA\_sub\_onA\_pos}\)
to recover the strict-dominance callback, matching the form naturally
produced by the real predecessor source-weight identity.

\textbf{Lowered coefficient dominance bridge}: the coefficient split
also gives the direct callback form used by the lowered Marshall
positivity wrappers.  Since the on-\(A\) signed contribution is
\(-\sum_{x\in A}C_x(\tau)\) and the off-\(A\) contribution is
\(\sum_{x\notin A}C_x(\tau)\), the old condition
\[
  -\sum_{x\in A}\operatorname{contribution}_x(\tau)
  <
  \sum_{x\notin A}\operatorname{contribution}_x(\tau)
\]
is rewritten as the coefficient inequality
\[
  \sum_{x\in A} C_x(\tau)
  <
  \sum_{x\notin A} C_x(\tau).
\]
Lean packages this as
\[
\texttt{tasaki23\_signed\_lowering\_site\_sum\_pos\_of\_onA\_coefficient\_lt\_offA}
\]
and then as the lowered-vector Marshall-positivity wrapper
\[
\texttt{tasaki23\_lowered\_marshall\_pos\_of\_onA\_coefficient\_lt\_offA}.
\]
This does not prove the quantitative dominance estimate; it removes the
signed local contribution notation from the remaining local input and
leaves a pure predecessor-coefficient comparison.

\textbf{Sublattice lowered coefficient components}: the coefficient
comparison is now connected to the actual sublattice lowering
operators.  Lean first expands
\[
  \hat S_A^- \Phi
  \quad\hbox{and}\quad
  \hat S_{\bar A}^- \Phi
\]
as the filtered sums of the single-site lowering contributions over
\(A\) and \(\bar A\):
\[
\texttt{sublatticeSpinSOpMinus\_mulVec\_magSectorEmbedding\_apply\_eq\_onA\_site\_sum}
\]
and
\[
\texttt{sublatticeSpinSOpMinus\_complement\_mulVec\_magSectorEmbedding\_apply\_eq\_offA\_site\_sum}.
\]
Combining those expansions with the coefficient split gives
\[
  \operatorname{sgn}_A(\tau)
  \operatorname{Re}((\hat S_{\bar A}^- \Phi)_\tau)
  =
  \sum_{x\notin A} C_x(\tau),
\]
and
\[
  \operatorname{sgn}_A(\tau)
  \operatorname{Re}((\hat S_A^- \Phi)_\tau)
  =
  -\sum_{x\in A} C_x(\tau).
\]
The corresponding Lean names are
\[
\texttt{tasaki23\_signed\_lowering\_offA\_sublattice\_component\_eq\_coefficient\_sum}
\]
and
\[
\texttt{tasaki23\_signed\_lowering\_onA\_sublattice\_component\_eq\_neg\_coefficient\_sum}.
\]
This places the remaining dominance estimate at the operator level:
one has to compare the Marshall-signed real components of
\(\hat S_A^- \Phi\) and \(\hat S_{\bar A}^- \Phi\), which is the form
needed to use the predicted ground-state and sublattice-Casimir
structure in Tasaki, Springer 2020, Section 2.5, Theorem 2.3, p. 42.

\textbf{Sublattice lowered dominance bridge}: the coefficient
dominance hypothesis is now replaced by the operator-component
comparison
\[
  -\operatorname{sgn}_A(\tau)
  \operatorname{Re}((\hat S_A^- \Phi)_\tau)
  <
  \operatorname{sgn}_A(\tau)
  \operatorname{Re}((\hat S_{\bar A}^- \Phi)_\tau).
\]
Using the two sublattice coefficient identities above, this is exactly
the predecessor-coefficient dominance hypothesis.  Lean packages the
strict site-sum version as
\[
\texttt{tasaki23\_signed\_lowering\_site\_sum\_pos\_of\_sublattice\_component\_lt}
\]
and the corresponding lowered-vector Marshall-positivity wrapper as
\[
\texttt{tasaki23\_lowered\_marshall\_pos\_of\_sublattice\_component\_lt}.
\]
Thus the remaining local input for the lowered adjacent-sector step has
the same form as the sublattice-ladder comparison appearing in Tasaki,
Springer 2020, Section 2.5, Theorem 2.3, p. 42.

\textbf{Predicted-GS sublattice Casimir bridges}: the predicted toy
ground-state subspace is defined as the joint eigenspace for
\((\hat S_{\mathrm{tot}})^2\), \((\hat S_A)^2\), and
\((\hat S_{\bar A})^2\).  In addition to the already used total-Casimir
bridge, Lean now exposes the two sublattice components:
\[
\begin{gathered}
\texttt{tasaki23\_sublatticeSpinSquaredS\_mulVec}\\
\texttt{\_of\_mem\_bipartiteToyGroundStateSubspacePredicted},\\
\texttt{tasaki23\_sublatticeSpinSquaredS\_complement\_mulVec}\\
\texttt{\_of\_mem\_bipartiteToyGroundStateSubspacePredicted}.
\end{gathered}
\]
These state that a predicted-GS vector is an eigenvector of
\((\hat S_A)^2\) at \(s_A(s_A+1)\) and of
\((\hat S_{\bar A})^2\) at \(s_B(s_B+1)\), where
\(s_A=|A|N/2\) and \(s_B=|\bar A|N/2\).  Combining these with the
predicted-GS lowering closure gives the lowered-image companions
\[
\begin{gathered}
\texttt{tasaki23\_lowered\_sublatticeSpinSquaredS}\\
\texttt{\_of\_mem\_bipartiteToyGroundStateSubspacePredicted},\\
\texttt{tasaki23\_lowered\_sublatticeSpinSquaredS\_complement}\\
\texttt{\_of\_mem\_bipartiteToyGroundStateSubspacePredicted}.
\end{gathered}
\]
This is the Casimir-side input needed for the next comparison of
\(\hat S_A^- \Phi\) and \(\hat S_{\bar A}^- \Phi\) in the proof of
Tasaki, Springer 2020, Section 2.5, Theorem 2.3, p. 42.

\textbf{Sublattice-ladder Casimir preservation}: the Casimir identities
also descend to the two individual lowering components.  The support
lemmas
\[
\begin{gathered}
\texttt{sublatticeSpinSquaredS\_commute\_sublatticeSpinSOpMinus},\\
\texttt{sublatticeSpinSquaredS\_commute\_sublatticeSpinSOpMinus\_complement}
\end{gathered}
\]
rewrite the Cartesian commutation of \((\hat S_A)^2\) with
\(\hat S_A^{(1)}\) and \(\hat S_A^{(2)}\) into the ladder basis, and record
the corresponding cross-sublattice commutation with
\(\hat S_{\bar A}^-\).  Therefore the predicted-GS bridges
\[
\begin{gathered}
\texttt{tasaki23\_sublatticeSpinSquaredS\_sublatticeSpinSOpMinus}\\
\texttt{\_of\_mem\_bipartiteToyGroundStateSubspacePredicted},\\
\texttt{tasaki23\_sublatticeSpinSquaredS\_complement}\\
\texttt{\_sublatticeSpinSOpMinus\_complement}\\
\texttt{\_of\_mem\_bipartiteToyGroundStateSubspacePredicted}
\end{gathered}
\]
show that \(\hat S_A^-\Phi\) and \(\hat S_{\bar A}^-\Phi\) stay in the
maximum \(A\)- and \(\bar A\)-sublattice Casimir eigenspaces,
respectively.  This is the component-level Casimir input for the
remaining comparison between the two lowered sublattice pieces.

\textbf{Joint sublattice-Casimir preservation for lowered components}:
the same cross-sublattice commutation also gives the missing opposite
Casimir identities.  The Lean bridges
\[
\begin{gathered}
\texttt{tasaki23\_sublatticeSpinSquaredS\_complement}\\
\texttt{\_sublatticeSpinSOpMinus}\\
\texttt{\_of\_mem\_bipartiteToyGroundStateSubspacePredicted},\\
\texttt{tasaki23\_sublatticeSpinSquaredS}\\
\texttt{\_sublatticeSpinSOpMinus\_complement}\\
\texttt{\_of\_mem\_bipartiteToyGroundStateSubspacePredicted}
\end{gathered}
\]
show that \(\hat S_A^-\Phi\) also has the maximum
\(\bar A\)-sublattice Casimir eigenvalue, and that
\(\hat S_{\bar A}^-\Phi\) also has the maximum \(A\)-sublattice Casimir
eigenvalue.  Hence both individual lowered components lie in the joint
maximum sublattice-Casimir eigenspace.  This is the precise
component-Casimir package required before proving the strict
Marshall-signed comparison in Tasaki, Springer 2020, Section 2.5,
Theorem 2.3, p. 42.

\textbf{Lowered component joint-eigenspace membership}: the preceding
four component Casimir identities are also packaged in the submodule
language used by the later comparison step.  The Lean bridges
\[
\begin{gathered}
\texttt{tasaki23\_sublatticeSpinSOpMinus}\\
\texttt{\_mem\_joint\_sublattice\_casimir\_eigenspace}\\
\texttt{\_of\_mem\_bipartiteToyGroundStateSubspacePredicted},\\
\texttt{tasaki23\_sublatticeSpinSOpMinus\_complement}\\
\texttt{\_mem\_joint\_sublattice\_casimir\_eigenspace}\\
\texttt{\_of\_mem\_bipartiteToyGroundStateSubspacePredicted}
\end{gathered}
\]
state directly that both \(\hat S_A^-\Phi\) and
\(\hat S_{\bar A}^-\Phi\) belong to the intersection of the maximum
\((\hat S_A)^2\)- and \((\hat S_{\bar A})^2\)-eigenspaces.  The proof is
just \texttt{Submodule.mem\_inf} applied to the same-side and opposite-side
Casimir bridges above.  This removes an API layer before the remaining
pointwise comparison of the Marshall-signed real parts.

\textbf{Lowered component joint-Casimir successor-sector membership}: for
a sector-\(M\) representative, the lowered components also carry the
successor-sector support.  The named subspace
\[
\texttt{tasaki23LoweredJointMagSubspace}
\]
is the intersection of
\(\texttt{tasaki23JointSublatticeCasimirEigenspace}(A,N)\) with the
magnetisation subspace corresponding to \(\texttt{magSumS}=M+1\).  The
two bridges
\[
\begin{gathered}
\texttt{tasaki23\_sublatticeSpinSOpMinus}\\
\texttt{\_mem\_lowered\_joint\_magSubspace}\\
\texttt{\_of\_mem\_bipartiteToyGroundStateSubspacePredicted},\\
\texttt{tasaki23\_sublatticeSpinSOpMinus\_complement}\\
\texttt{\_mem\_lowered\_joint\_magSubspace}\\
\texttt{\_of\_mem\_bipartiteToyGroundStateSubspacePredicted}
\end{gathered}
\]
combine the joint-Casimir membership above with the standard
sublattice-lowering magnetisation shift.  Thus both lowered components
are available as successor-sector vectors with the joint-Casimir
structure needed for the next component comparison.

\textbf{Lowered joint-magnetisation callback threading}: the interval
chain now exposes the stronger lowered-component package directly to the
remaining local comparison callback.  The wrapper
\[
\begin{gathered}
\texttt{tasaki23\_energy\_interval\_chain\_with\_predictedGS}\\
\texttt{\_of\_left\_endpoint\_predictedGS}\\
\texttt{\_of\_lowered\_joint\_magSubspace\_component\_lt}\\
\texttt{\_of\_predictedGS}
\end{gathered}
\]
threads, for each source sector \(M\), the memberships
\[
  \hat S_A^-\Phi,\ \hat S_{\bar A}^-\Phi
    \in \texttt{tasaki23LoweredJointMagSubspace}(A,N,M).
\]
Thus the callback sees both the joint-Casimir structure and the
successor-sector support at once.  The final wrappers
\[
\begin{gathered}
\texttt{tasaki\_2\_5\_theorem\_2\_3\_of\_left\_endpoint\_threaded}\\
\texttt{\_predictedGS\_of\_lowered\_joint\_magSubspace}\\
\texttt{\_component\_lt\_of\_sector\_minimality},\\
\texttt{tasaki\_2\_5\_theorem\_2\_3\_of\_left\_endpoint\_threaded}\\
\texttt{\_predictedGS\_of\_lowered\_joint\_magSubspace}\\
\texttt{\_coefficient\_lt\_of\_sector\_minimality}
\end{gathered}
\]
make the same stronger callback available at the final theorem
boundary.  The coefficient form rewrites on/off predecessor coefficient
dominance into the strict Marshall-signed sublattice-component
comparison using the existing component formulas.

\textbf{Lowered joint cross-ladder callback threading}: the same
callback is further sharpened by exposing the predicted-GS
raised-lowered cross-ladder identity directly.  The interval wrapper
\[
\begin{gathered}
\texttt{tasaki23\_energy\_interval\_chain\_with\_predictedGS}\\
\texttt{\_of\_left\_endpoint\_predictedGS}\\
\texttt{\_of\_lowered\_joint\_cross\_ladder\_component\_lt}\\
\texttt{\_of\_predictedGS}
\end{gathered}
\]
passes the identity
\[
\begin{aligned}
  &\hat S_A^+(\hat S_{\bar A}^-\Phi)
    + \hat S_{\bar A}^+(\hat S_A^-\Phi) \\
  &\quad =
    E_{\mathrm{toy}}^{\mathrm{pred}}\Phi
    - 2\,\hat S_A^{(3)}\hat S_{\bar A}^{(3)}\Phi
\end{aligned}
\]
together with the two memberships in
\(\texttt{tasaki23LoweredJointMagSubspace}(A,N,M)\).  The final
component and coefficient wrappers expose the same cross-ladder-aware
callback at the Theorem 2.3 boundary.  Thus the remaining local
comparison is isolated under precisely the cross-ladder equation and
successor-sector joint-Casimir structure already formalised.

\textbf{Unpacked lowered joint cross-ladder callback}: the packed
membership
\[
  \hat S_A^-\Phi,\ \hat S_{\bar A}^-\Phi
  \in \texttt{tasaki23LoweredJointMagSubspace}(A,N,M)
\]
is now unpacked before the local comparison callback is invoked.  The
extraction lemmas expose, for each lowered component, the two maximum
sub\-lattice-Casimir eigenvector equations for \((\hat S_A)^2\) and
\((\hat S_{\bar A})^2\), together with membership in the successor
magnetization subspace
\[
  \texttt{magSubspaceS}\bigl(|\Lambda|N/2-(M+1)\bigr).
\]
The interval wrapper
\[
\begin{gathered}
\texttt{tasaki23\_energy\_interval\_chain\_with\_predictedGS}\\
\texttt{\_of\_left\_endpoint\_predictedGS}\\
\texttt{\_of\_unpacked\_lowered\_joint\_cross\_ladder}\\
\texttt{\_component\_lt\_of\_predictedGS}
\end{gathered}
\]
and the corresponding final wrapper pass these unpacked equations
alongside the raised-lowered cross-ladder identity.  The remaining
component comparison can therefore work directly with the equations and
sector support rather than reopening the packed submodule membership.
The companion coefficient wrapper rewrites predecessor coefficient
dominance through the existing on/off sublattice component formulas and
then invokes this unpacked component wrapper.

\textbf{Re-embedded cross-ladder source-sector site sums}: the
successor \texttt{magSubspaceS} support can now be used to replace each
lowered component by the zero-extension of its sector restriction.  The
theorem
\begin{Verbatim}
tasaki23_cross_ladder_reembedded_source_site_sum_eq_energy_sub_two_op3_of_predictedGS
\end{Verbatim}
evaluates the raised-lowered cross-ladder identity at a source-sector
configuration and rewrites the two raised terms as
\[
\begin{aligned}
  &\sum_{x\in A}
    \hat S_x^+\,\iota_{M+1}\bigl((\hat S_{\bar A}^-\Phi)|_{H_{M+1}}\bigr)
  +\sum_{x\notin A}
    \hat S_x^+\,\iota_{M+1}\bigl((\hat S_A^-\Phi)|_{H_{M+1}}\bigr).
\end{aligned}
\]
The supporting sublattice expansions are
\begin{Verbatim}
sublatticeSpinSOpPlus_mulVec_magSectorEmbedding_apply_eq_onA_site_sum
sublatticeSpinSOpPlus_complement_mulVec_magSectorEmbedding_apply_eq_offA_site_sum
\end{Verbatim}
This is the site-sum form needed before converting the cross-ladder
equation into the local Marshall-signed coefficient comparison.

The final-theorem callback
\begin{Verbatim}
tasaki_2_5_theorem_2_3_of_left_endpoint_threaded_predictedGS_of_reembedded_cross_ladder_coefficient_lt_of_sector_minimality
\end{Verbatim}
threads this pointwise source-sector identity directly to the remaining
coefficient-comparison hypothesis.  Thus the local input no longer has to
open the operator-level cross-ladder equation or redo the re-embedding
step; it can compare the two successor-sector restrictions using the
filtered \(A\) and \(\bar A\) raising sums and the already-unpacked
sublattice-Casimir equations.

\textbf{Sublattice \(S^3\) source weights}: the right-hand side of the
same source-sector identity also contains
\(\hat S_A^3\hat S_{\bar A}^3\Phi\).  The local diagonal bridge
\begin{Verbatim}
onSiteS_spinSOp3_mulVec_apply
\end{Verbatim}
states that \(\hat S_x^3\) acts on a configuration \(\sigma\) by the
scalar \(N/2-\sigma_x\).  Summing this identity over the two Boolean
sublattices gives
\begin{Verbatim}
sublatticeSpinSOp3_mulVec_apply_eq_onA_weight
sublatticeSpinSOp3_complement_mulVec_apply_eq_offA_weight
\end{Verbatim}
and applying the two diagonal actions in sequence gives
\begin{Verbatim}
sublatticeSpinSOp3_mul_complement_mulVec_apply_eq_onA_offA_weight
\end{Verbatim}
Thus the source-sector component of
\(\hat S_A^3\hat S_{\bar A}^3\Phi\) is
\[
  \left(\sum_{x\in A}\left(\frac N2-\sigma_x\right)\right)
  \left(\sum_{x\notin A}\left(\frac N2-\sigma_x\right)\right)\Phi_\sigma,
\]
which turns the non-ladder term in the re-embedded cross-ladder equation
into an explicit scalar multiple of the source coefficient.

The re-embedded site-sum identity is correspondingly sharpened by
\begin{Verbatim}
tasaki23_cross_ladder_reembedded_source_site_sum_eq_energy_sub_two_sublattice_weight_product_of_predictedGS
\end{Verbatim}
to the scalar right-hand side
\[
  \left(E_{\mathrm{toy}}
    -2
      \left(\sum_{x\in A}\left(\frac N2-\sigma_x\right)\right)
      \left(\sum_{x\notin A}\left(\frac N2-\sigma_x\right)\right)
  \right)\Phi_\sigma .
\]
This removes \(\hat S_A^3\hat S_{\bar A}^3\) from the local
cross-ladder comparison and leaves only the two re-embedded
successor-sector raising sums against an explicit source coefficient.
The local predecessor specialization
\begin{Verbatim}
tasaki23_cross_ladder_reembedded_source_weight_eq_lowering_predecessor_of_predictedGS
\end{Verbatim}
then evaluates this scalar identity at the configuration obtained from
a successor-sector \(\tau\in\mathcal H_{M+1}\) by lowering a chosen
site \(x\).  This is the same source configuration used in
\texttt{tasaki23LoweringPredecessorSignedCoefficient}, so the
source-weight equation is now aligned with the coefficient comparison
over successor configurations and sites.
For a Marshall-positive sector embedding
\(\Psi(\rho)=\operatorname{sgn}_A(\rho)v_\rho\), Lean also records the
real-part form of this predecessor right-hand side:
\[
\begin{gathered}
\texttt{tasaki23\_cross\_ladder\_reembedded\_source\_weight}\\
\texttt{\_lowering\_predecessor\_rhs\_re\_eq}.
\end{gathered}
\]
It rewrites
\[
  \operatorname{Re}\bigl((E_{\rm toy}-2W_A(\rho)W_{\bar A}(\rho))
    \Psi(\rho)\bigr)
\]
as
\[
  \bigl(E_{\rm toy}^{\rm re}-2W_A(\rho)W_{\bar A}(\rho)\bigr)
  \operatorname{sgn}_A(\rho)_{\rm re}v_\rho,
\]
where the two \(W\)-terms are the real filtered sums of
\(N/2-\rho_y\).  This is the real scalar form needed before comparing
positive-source predecessor coefficients.
The companion theorem
\[
\begin{gathered}
\texttt{tasaki23\_cross\_ladder\_reembedded\_source\_weight}\\
\texttt{\_lowering\_predecessor\_re\_eq}
\end{gathered}
\]
applies this RHS rewrite after taking real parts of the full
predecessor source-weight equality.  It exposes the real part of the
two filtered raising sums as the same real scalar coefficient times the
Marshall-signed positive source coefficient at \(\rho\).
The final callback boundary is updated by
\begin{Verbatim}
tasaki_2_5_theorem_2_3_of_left_endpoint_threaded_predictedGS_of_reembedded_source_weight_coefficient_lt_of_sector_minimality
\end{Verbatim}
so that the remaining local coefficient-comparison hypothesis receives
this scalar source-weight identity directly, together with the four
maximum sublattice-Casimir equations for the two lowered components.
Thus the callback no longer has to consume the raw
\(\hat S_A^3\hat S_{\bar A}^3\) operator term.
The companion boundary
\begin{Verbatim}
tasaki_2_5_theorem_2_3_of_left_endpoint_threaded_predictedGS_of_unpacked_reembedded_source_weight_coefficient_lt_of_sector_minimality
\end{Verbatim}
uses the unpacked lowered-joint callback instead.  It passes the same
scalar source-weight identity, the four maximum sublattice-Casimir
equations, and the two successor-sector
\(\mathcal H_{M+1}\)-support facts for
\(\hat S_A^-\Psi\) and \(\hat S_{\bar A}^-\Psi\).  Hence the remaining
local comparison can use both the explicit coefficient
\(E_{\rm toy}-2W_A(\sigma)W_{\bar A}(\sigma)\) and the fact that the two
lowered components are already restricted to the successor sector.
The outside-sector version
\begin{Verbatim}
tasaki_2_5_theorem_2_3_of_left_endpoint_threaded_predictedGS_of_unpacked_reembedded_source_weight_coefficient_lt_of_outside_sector_ground_energy_lower_bound
\end{Verbatim}
combines this unpacked source-weight wrapper with
\texttt{tasaki23OutsideGroundEnergyLowerFamilyCallback}.  It therefore
uses the scalar re-embedded source-weight callback as the remaining local
coefficient input, while admissible-sector minimality comes from the
common-energy chain and non-admissible sectors are handled by the
outside-sector lower family.
The predecessor-specialized boundary
\begin{Verbatim}
tasaki_2_5_theorem_2_3_of_left_endpoint_threaded_predictedGS_of_unpacked_reembedded_source_weight_predecessor_coefficient_lt_of_sector_minimality
\end{Verbatim}
pushes this one step further.  It instantiates the source-weight identity
at each predecessor obtained from a successor \(\tau\) by lowering a
site \(x\), and passes those predecessor identities directly to the
coefficient-comparison callback.  The callback therefore receives the
filtered raising sums and the scalar coefficient
\((E_{\rm toy}-2W_AW_{\bar A})\Psi\) already evaluated at the same
configuration used by the lowered predecessor coefficient.
Its outside-sector companion
\begin{Verbatim}
tasaki_2_5_theorem_2_3_of_left_endpoint_threaded_predictedGS_of_unpacked_reembedded_source_weight_predecessor_coefficient_lt_of_outside_sector_ground_energy_lower_bound
\end{Verbatim}
combines the predecessor-specialized wrapper with
\texttt{tasaki23OutsideGroundEnergyLowerFamilyCallback}.  This leaves the
remaining local coefficient comparison at lowering predecessor
configurations, while sector minimality is supplied by the common-energy
chain and the outside-sector lower family.
The final positive-source boundary
\begin{Verbatim}
tasaki_2_5_theorem_2_3_of_left_endpoint_threaded_predictedGS_of_unpacked_reembedded_source_weight_predecessor_positive_source_coefficient_lt_of_sector_minimality
\end{Verbatim}
combines this predecessor-specialized source-weight callback with the
sign-normalized coefficient bridge above.  The local callback still
receives the filtered raising sums, scalar source-weight right-hand side,
four maximum sublattice-Casimir equations, and two successor-sector
support facts, but it may state the dominance estimate directly as
\[
  \sum_{x\in A} C_x^{+}(\tau)
  <
  \sum_{x\notin A} C_x^{+}(\tau),
\]
where \(C_x^{+}\) is the positive-source predecessor coefficient.
Its outside-sector companion
\begin{Verbatim}
tasaki_2_5_theorem_2_3_of_left_endpoint_threaded_predictedGS_of_unpacked_reembedded_source_weight_predecessor_positive_source_coefficient_lt_of_outside_sector_ground_energy_lower_bound
\end{Verbatim}
combines the positive-source predecessor wrapper with
\texttt{tasaki23OutsideGroundEnergyLowerFamilyCallback}.  Thus the
outside-sector final boundary now accepts the sign-normalized
positive-source predecessor coefficient comparison directly, while sector
minimality is again supplied by the common-energy chain and the
outside-sector lower family.

\textbf{Joint-component callback}: the joint maximum sublattice-Casimir
eigenspace used above is now named as
\[
\texttt{tasaki23JointSublatticeCasimirEigenspace}.
\]
The interval-chain wrapper
\[
\begin{gathered}
\texttt{tasaki23\_energy\_interval\_chain\_with\_predictedGS}\\
\texttt{\_of\_left\_endpoint\_predictedGS}\\
\texttt{\_of\_joint\_sublattice\_component\_lt}\\
\texttt{\_of\_predictedGS}
\end{gathered}
\]
threads, at each successor step, the propagated membership
\(\Phi\in\mathcal G_{\mathrm{pred}}(A,N)\) together with
\[
  \hat S_A^-\Phi,\ \hat S_{\bar A}^-\Phi
    \in \texttt{tasaki23JointSublatticeCasimirEigenspace}(A,N)
\]
into the local comparison callback.  The remaining local task is
therefore exactly the strict pointwise inequality between the
Marshall-signed real parts of the two lowered sublattice components,
under the predicted-GS and joint-Casimir hypotheses already available.

\textbf{Joint-component final wrapper}: the Lean theorem
\[
\begin{gathered}
\texttt{tasaki\_2\_5\_theorem\_2\_3\_of\_left\_endpoint\_threaded}\\
\texttt{\_predictedGS\_of\_joint\_sublattice\_component\_lt}\\
\texttt{\_of\_sector\_minimality}
\end{gathered}
\]
exposes the same local callback at the final theorem boundary.  It
invokes the joint-component interval chain above, discards the threaded
predicted-GS field after extracting the common sector energy, and then
uses the existing sector-minimality bridge to supply the global
minimality clause in the final Theorem 2.3 statement.  Thus the final
wrapper can now assume the strict component inequality only in the
context where predicted-GS membership and the two joint-Casimir
memberships have already been provided.

\textbf{Joint-coefficient final wrapper}: the Lean theorem
\[
\begin{gathered}
\texttt{tasaki\_2\_5\_theorem\_2\_3\_of\_left\_endpoint\_threaded}\\
\texttt{\_predictedGS\_of\_joint\_sublattice\_coefficient\_lt}\\
\texttt{\_of\_sector\_minimality}
\end{gathered}
\]
lowers the same final callback from operator components to the
predecessor coefficient sums.  The callback still receives the
predicted-GS membership of the source vector and the two joint
sublattice-Casimir memberships for \(\hat S_A^-\Phi\) and
\(\hat S_{\bar A}^-\Phi\).  It then proves the strict inequality
between the on-\(A\) and off-\(A\) predecessor coefficient sums.  The
wrapper rewrites this coefficient inequality using
\texttt{tasaki23\_signed\_lowering\_onA\_sublattice\_component\_eq\_neg\_coefficient\_sum}
and
\texttt{tasaki23\_signed\_lowering\_offA\_sublattice\_component\_eq\_coefficient\_sum},
obtaining the joint-component comparison required by the preceding
final wrapper.

\textbf{Predicted-GS toy cross-dot bridge}: predicted toy ground-state
membership also gives the pointwise toy-Hamiltonian identity recorded in
\[
\begin{gathered}
\texttt{tasaki23\_heisenbergToyHamiltonianS\_mulVec}\\
\texttt{\_of\_mem\_bipartiteToyGroundStateSubspacePredicted}
\end{gathered}
\]
Combining this with the Casimir identity
\(\hat H_{\mathrm{toy},S}=2\,\hat S_A\cdot \hat S_{\bar A}\) gives
\[
\begin{gathered}
\texttt{tasaki23\_two\_sublatticeSpinSDot\_mulVec}\\
\texttt{\_of\_mem\_bipartiteToyGroundStateSubspacePredicted}
\end{gathered}
\]
The general identity
\[
\texttt{sublatticeSpinSDot\_eq\_op3\_add\_ladder}
\]
rewrites the dot product as
\(\hat S_A^3\hat S_{\bar A}^3+
\frac12(\hat S_A^+\hat S_{\bar A}^-+
\hat S_A^-\hat S_{\bar A}^+)\), yielding
\[
\begin{gathered}
\texttt{tasaki23\_two\_cross\_ladder\_mulVec}\\
\texttt{\_of\_mem\_bipartiteToyGroundStateSubspacePredicted}
\end{gathered}
\]
The follow-up corollary
\[
\begin{gathered}
\texttt{tasaki23\_cross\_ladder\_sum\_mulVec\_eq}\\
\texttt{\_energy\_sub\_two\_op3}\\
\texttt{\_of\_predictedGS}
\end{gathered}
\]
isolates
\(\hat S_A^+\hat S_{\bar A}^-+\hat S_A^-\hat S_{\bar A}^+\)
as the predicted toy-energy term minus
\(2\hat S_A^3\hat S_{\bar A}^3\).  The sequential form
\[
\begin{gathered}
\texttt{tasaki23\_cross\_ladder\_sequential\_mulVec\_eq}\\
\texttt{\_energy\_sub\_two\_op3\_of\_predictedGS}
\end{gathered}
\]
then rewrites the same left hand side as
\(\hat S_A^+(\hat S_{\bar A}^-\Phi)+
\hat S_A^-(\hat S_{\bar A}^+\Phi)\).  Thus the cross-sublattice
spin-dot constraint is available in the ladder-term form needed for the
remaining comparison between \(\hat S_A^-\Phi\) and
\(\hat S_{\bar A}^-\Phi\).  The raised-lowered component variant
\[
\begin{gathered}
\texttt{tasaki23\_cross\_ladder\_raised\_lowered}\\
\texttt{\_components\_eq\_energy\_sub\_two\_op3}\\
\texttt{\_of\_predictedGS}
\end{gathered}
\]
commutes the second term to
\(\hat S_{\bar A}^+(\hat S_A^-\Phi)\), so the two summands act directly
on the two lowered components.

\textbf{Sublattice-component final wrapper}: the final theorem wrapper can
now consume that sublattice comparison directly.  The Lean theorem
\[
\begin{gathered}
\texttt{tasaki\_2\_5\_theorem\_2\_3\_of\_left\_endpoint\_threaded}\\
\texttt{\_predictedGS\_of\_sublattice\_component\_lt}\\
\texttt{\_of\_outside\_real\_sector\_minimality}
\end{gathered}
\]
assumes, for each source-sector Marshall-positive representative \(\Phi\),
the pointwise inequality
\[
  -\operatorname{sgn}_A(\tau)
    \operatorname{Re}\bigl((\hat S_A^-\Phi)_\tau\bigr)
  <
  \operatorname{sgn}_A(\tau)
    \operatorname{Re}\bigl((\hat S_{\bar A}^-\Phi)_\tau\bigr).
\]
It converts this hypothesis to lowered-vector Marshall positivity using
\(\texttt{tasaki23\_lowered\_marshall\_pos\_of\_sublattice\_component\_lt}\)
and then invokes the existing left-endpoint predicted-GS final wrapper with
outside-interval real-sector minimality.  Thus the remaining local input is
exactly the sublattice-ladder comparison above, rather than the more
abstract positivity of the total lowered vector.

\textbf{Predicted-GS-aware sublattice-component interval chain}: the local
comparison can now be restricted to source vectors already known to lie in
the predicted toy ground-state subspace.  The Lean theorem
\[
\begin{gathered}
\texttt{tasaki23\_energy\_interval\_chain\_with\_predictedGS}\\
\texttt{\_of\_left\_endpoint\_predictedGS}\\
\texttt{\_of\_sublattice\_component\_lt\_of\_predictedGS}
\end{gathered}
\]
uses the same left-endpoint induction as the predicted-GS threading
wrapper, but in each successor step it passes the propagated membership
\(\Phi\in\texttt{bipartiteToyGroundStateSubspacePredicted}(A,N)\)
to the component comparison callback.  After applying
\(\texttt{tasaki23\_signed\_lowering\_site\_sum\_pos\_of\_sublattice\_component\_lt}\),
the ordinary strict site-sum successor step supplies the next
Marshall-positive representative and the existing scalar-cancellation
argument transports predicted-GS membership forward.  This leaves the
remaining Casimir calculation in the natural form: prove the inequality
between \(\hat S_A^-\Phi\) and \(\hat S_{\bar A}^-\Phi\) for predicted-GS
source vectors only.

The dominance-form sector-existence companions are
\[
\begin{gathered}
\texttt{tasaki23\_successor\_sector\_existence}\\
\texttt{\_with\_lowered\_predictedCasimir\_of\_onA\_neg\_lt\_offA},\\
\texttt{tasaki23\_predecessor\_sector\_existence}\\
\texttt{\_with\_raised\_predictedCasimir\_of\_onA\_neg\_lt\_offA}.
\end{gathered}
\]
They first choose the sector representative by the same per-sector
Theorem~2.2 existence wrapper, then convert the lowered or raised
off-\(A\) dominance hypothesis into the strict site-sum positivity input
using
\(\texttt{tasaki23\_signed\_\{lowering,raising\}\_site\_sum\_pos\_of\_onA\_neg\_lt\_offA}\).
The dominance estimate is still an explicit hypothesis; the wrappers
only connect it to the predicted-Casimir adjacent-chain machinery.

\textbf{Predicted-GS sector-existence dominance wrappers}: in the
canonical orientation \(|\neg A|\le |A|\), the explicit
predicted-Casimir callback in the preceding sector-existence wrappers
can be replaced by a predicted-GS membership callback.  The Lean
theorems
\[
\begin{gathered}
\texttt{tasaki23\_successor\_sector\_existence}\\
\texttt{\_with\_lowered\_predictedGS\_of\_onA\_neg\_lt\_offA},\\
\texttt{tasaki23\_predecessor\_sector\_existence}\\
\texttt{\_with\_raised\_predictedGS\_of\_onA\_neg\_lt\_offA}
\end{gathered}
\]
apply
\[
\texttt{tasaki23\_totalSpinSSquared\_mulVec\_of\_mem}\\
\texttt{\_bipartiteToyGroundStateSubspacePredicted}
\]
to the source-sector vector selected by Theorem~2.2.  This supplies the
\(C_{\mathrm{pred}}\) total-Casimir identity required by the
dominance-form predicted-Casimir sector-existence wrappers, while the
lowered or raised off-\(A\) dominance estimate remains explicit.

\textbf{Predicted-Casimir adjacent common-energy wrappers}: the
Casimir endpoint mismatch is now connected directly to the one-step
sector chain.  The Lean theorems
\[
\begin{gathered}
\texttt{tasaki23\_successor\_sector\_common\_energy}\\
\texttt{\_of\_site\_sum\_pos\_of\_predictedCasimirValue},\\
\texttt{tasaki23\_predecessor\_sector\_common\_energy}\\
\texttt{\_of\_site\_sum\_pos\_of\_predictedCasimirValue}
\end{gathered}
\]
specialize the preceding Casimir non-vanishing wrappers to
\[
  C_{\mathrm{pred}}
  =
  S_{\mathrm{tot}}(S_{\mathrm{tot}}+1),
  \qquad
  S_{\mathrm{tot}}
  =
  \frac{\bigl||A|-|\neg A|\bigr|N}{2}.
\]
For the lowering direction, if the source-sector vector in sector
\(M\) has total-Casimir eigenvalue \(C_{\mathrm{pred}}\) and
\(M\) is not the right endpoint, then
\[
  C_{\mathrm{pred}}
  \ne
  (|\Lambda|N/2-M)(|\Lambda|N/2-M-1)
\]
by
\texttt{tasaki23\_predictedCasimirValue\_ne\_lowering\_kernel\_value\_of\_mem\_of\_lt\_right}.
Thus the abstract endpoint-mismatch hypothesis in the lowering
Casimir wrapper is discharged automatically.  The raising theorem does
the same for a source vector in sector \(M+1\), using
\[
  C_{\mathrm{pred}}
  \ne
  (|\Lambda|N/2-(M+1))(|\Lambda|N/2-(M+1)+1)
\]
and
\texttt{tasaki23\_predictedCasimirValue\_ne\_raising\_kernel\_value\_of\_mem\_of\_left\_lt}.
These wrappers still assume the predicted-Casimir eigenvector
hypothesis and the local site-sum positivity input; their role is to
remove the now-solved endpoint mismatch from the final Theorem~2.3
chain assembly.

\textbf{Predicted-Casimir dominance common-energy wrappers}: the
local site-sum positivity input in the preceding wrappers can now be
supplied directly from the lowered or raised dominance hypotheses.  The
Lean theorems
\[
\begin{gathered}
\texttt{tasaki23\_successor\_sector\_common\_energy}\\
\texttt{\_of\_onA\_neg\_lt\_offA\_of\_predictedCasimirValue},\\
\texttt{tasaki23\_predecessor\_sector\_common\_energy}\\
\texttt{\_of\_onA\_neg\_lt\_offA\_of\_predictedCasimirValue}
\end{gathered}
\]
combine the predicted-Casimir common-energy wrappers with
\(\texttt{tasaki23\_signed\_\{lowering,raising\}\_site\_sum\_pos\_of\_onA\_neg\_lt\_offA}\).
For a source sector vector already known to have total-Casimir
eigenvalue \(C_{\mathrm{pred}}\), the remaining positivity input is
therefore the pointwise comparison
\[
  -\sum_{x\in A} L_x(\tau) < \sum_{x\notin A} L_x(\tau)
  \quad\text{or}\quad
  -\sum_{x\in A} R_x(\tau) < \sum_{x\notin A} R_x(\tau).
\]
The wrappers do not prove those dominance inequalities; they keep the
final chain API in the form needed once the dominance estimate is
available.

\textbf{Predicted-GS dominance common-energy wrappers}: in the
canonical orientation \(|\neg A|\le |A|\), membership in the predicted
toy ground-state subspace can now replace the explicit total-Casimir
hypothesis in the dominance common-energy step.  The Lean theorems
\[
\begin{gathered}
\texttt{tasaki23\_successor\_sector\_common\_energy}\\
\texttt{\_of\_onA\_neg\_lt\_offA\_of\_mem}\\
\texttt{\_bipartiteToyGroundStateSubspacePredicted},\\
\texttt{tasaki23\_predecessor\_sector\_common\_energy}\\
\texttt{\_of\_onA\_neg\_lt\_offA\_of\_mem}\\
\texttt{\_bipartiteToyGroundStateSubspacePredicted}
\end{gathered}
\]
first apply
\[
\texttt{tasaki23\_totalSpinSSquared\_mulVec\_of\_mem}\\
\texttt{\_bipartiteToyGroundStateSubspacePredicted}
\]
to obtain the \(C_{\mathrm{pred}}\) total-Casimir identity for the
source-sector zero-extension.  They then invoke the predicted-Casimir
dominance wrappers above.  Thus the remaining source-sector inputs are
membership in the predicted toy ground-state subspace and the same
explicit lowered or raised dominance inequality.

\textbf{Predicted-GS energy interval chain}: the canonical-orientation
lowering chain has now been assembled over the whole admissible
magnetization interval.  The Lean theorem
\[
\texttt{tasaki23\_energy\_interval\_chain\_of\_left\_endpoint}\\
\texttt{\_predictedGS\_of\_onA\_neg\_lt\_offA}
\]
starts at the left endpoint
\[
  L=\min(|A|,|\neg A|)\,N
\]
using the per-sector Theorem~2.2 existence wrapper, obtaining a
Marshall-positive eigenvector with energy \(\mu\).  For each
successor step \(M\mapsto M+1\) with \(M<R\), where
\[
  R=\max(|A|,|\neg A|)\,N,
\]
it invokes the predicted-GS lowered dominance common-energy wrapper.
Induction by \texttt{Nat.le\_induction} therefore yields, for every
\(M\in[L,R]\), a Marshall-positive sector vector with the same energy
\(\mu\).

This is still a conditional chain, not the final Theorem~2.3 proof:
for each currently chosen sector vector it assumes exactly the two
remaining source-sector callbacks, namely membership in
\(\texttt{bipartiteToyGroundStateSubspacePredicted}\) and the lowered
off-\(A\) dominance estimate.  The purpose is to separate the discrete
interval propagation from those remaining analytic inputs.

\textbf{Global minimality from sector minimality}: the remaining
global-minimality callback has been reduced to a sectorwise statement.
The Lean theorem
\[
\texttt{heisenbergHamiltonianSMatrixOnMagSector\_mulVec}\\
\texttt{\_magSectorRestriction\_of\_full\_eigen}
\]
states that the restriction of any full-space eigenvector to a
magnetization sector is a sector-matrix eigenvector at the same
eigenvalue, without assuming the full vector is supported in that
sector.  If the full eigenvector is nonzero, choosing a configuration
where it has a nonzero coefficient gives a nonzero restriction to the
corresponding \(\texttt{magSumS}\) sector.  Hence
\[
\texttt{tasaki23\_global\_minimality\_of\_sector\_minimality}
\]
proves the full global-minimality clause from lower bounds on all
sector eigenvalues.  The wrapper
\[
\begin{gathered}
\texttt{tasaki\_2\_5\_theorem\_2\_3}\\
\texttt{\_of\_left\_endpoint\_predictedGS}\\
\texttt{\_of\_onA\_neg\_lt\_offA}\\
\texttt{\_of\_sector\_minimality}
\end{gathered}
\]
therefore replaces the previous full-space global-minimality callback
by a sharper sector-minimality callback, while leaving the predicted-GS
membership and lowered off-\(A\) dominance inputs explicit.

The sector-minimality callback is reduced one step further to a real
sector statement.  For real coupling, the existing real-part extraction
lemma and the new Lean theorem
\[
\begin{gathered}
\texttt{heisenbergHamiltonianSReMatrix}\\
\texttt{OnMagSector\_mulVec\_im}\\
\texttt{\_of\_complex\_eigenvec}
\end{gathered}
\]
show that the real and imaginary parts of any complex sector eigenvector
at a real eigenvalue are real-form sector eigenvectors at the same
eigenvalue.  Since a nonzero complex vector has a nonzero real or
imaginary component,
\[
\texttt{tasaki23\_sector\_minimality\_of\_real\_sector\_minimality}
\]
proves complex sector minimality from lower bounds on nonzero real-form
sector eigenvectors.  The wrapper
\[
\begin{gathered}
\texttt{tasaki\_2\_5\_theorem\_2\_3}\\
\texttt{\_of\_left\_endpoint\_predictedGS}\\
\texttt{\_of\_onA\_neg\_lt\_offA}\\
\texttt{\_of\_real\_sector}\\
\texttt{\_minimality}
\end{gathered}
\]
therefore leaves the remaining spectral lower-bound input entirely in
the real sector matrix.

\textbf{Predicted-Casimir final chain}: the final interval chain also has a
form that no longer asks for predicted toy ground-state membership.  The
Lean theorem
\[
\begin{gathered}
\texttt{tasaki23\_energy\_interval\_chain}\\
\texttt{\_of\_left\_endpoint\_predictedCasimir}\\
\texttt{\_of\_onA\_neg\_lt\_offA}
\end{gathered}
\]
uses the same lowered off-\(A\) dominance step, but its source-vector input
is the predicted total-Casimir identity itself.  The helper
\[
\begin{gathered}
\texttt{tasaki\_2\_5\_theorem\_2\_3}\\
\texttt{\_of\_common\_energy\_chain}
\end{gathered}
\]
then factors the final uniqueness packaging common to all interval-chain
wrappers.  The intermediate wrapper
\[
\begin{gathered}
\texttt{tasaki\_2\_5\_theorem\_2\_3}\\
\texttt{\_of\_left\_endpoint\_predictedCasimir}\\
\texttt{\_of\_onA\_neg\_lt\_offA}\\
\texttt{\_of\_sector}\\
\texttt{\_minimality}
\end{gathered}
\]
uses the sectorwise minimality bridge for the global lower bound.
Consequently
\[
\begin{gathered}
\texttt{tasaki\_2\_5\_theorem\_2\_3}\\
\texttt{\_of\_left\_endpoint\_predictedCasimir}\\
\texttt{\_of\_onA\_neg\_lt\_offA}\\
\texttt{\_of\_real\_sector}\\
\texttt{\_minimality}
\end{gathered}
\]
reduces the remaining final inputs to the source predicted-Casimir identity,
the lowered off-\(A\) dominance estimate, and the real-sector spectral lower
bound.

\textbf{Predicted-GS total-Casimir bridge}: the total-Casimir component
of the predicted toy ground-state subspace has now been put into the
same notation as the adjacent-sector chain.  In the canonical
orientation \(|\neg A|\le |A|\), the Lean theorem
\[
\begin{gathered}
\texttt{tasaki23PredictedCasimirValue}\\
\texttt{\_eq\_canonical}\\
\texttt{\_of\_card\_notA\_le\_cardA}
\end{gathered}
\]
identifies \(C_{\mathrm{pred}}\) with the canonical joint-Casimir target
appearing in the predicted toy ground-state subspace.  Thus
\[
\begin{gathered}
\texttt{tasaki23\_totalSpinSSquared\_mulVec}\\
\texttt{\_of\_mem}\\
\texttt{\_bipartiteToyGroundStateSubspacePredicted}
\end{gathered}
\]
states that any vector in that predicted toy ground-state subspace
satisfies
\[
  \hat S_{\mathrm{tot}}^2\Phi
    =
  C_{\mathrm{pred}}\Phi .
\]
This supplies the exact source-sector Casimir hypothesis required by
the successor/predecessor predicted-Casimir chain once the sector
ground-state vector is known to lie in the predicted toy ground-state
subspace.

\textbf{Predicted-GS ladder closure}: the predicted toy ground-state
subspace is also closed under the total ladder operators in the exact
pointwise form needed by the adjacent-sector chain.  The Lean theorems
\[
\begin{gathered}
\texttt{tasaki23\_totalSpinSOpMinus\_mulVec\_mem}\\
\texttt{\_bipartiteToyGroundStateSubspacePredicted},\\
\texttt{tasaki23\_totalSpinSOpPlus\_mulVec\_mem}\\
\texttt{\_bipartiteToyGroundStateSubspacePredicted}
\end{gathered}
\]
state that
\[
  \Phi\in\mathcal G_{\mathrm{pred}}
  \quad\Longrightarrow\quad
  \hat S^-_{\mathrm{tot}}\Phi\in\mathcal G_{\mathrm{pred}},
  \qquad
  \hat S^+_{\mathrm{tot}}\Phi\in\mathcal G_{\mathrm{pred}} .
\]
These are direct uses of the previously formalized submodule
invariance of \(\mathcal G_{\mathrm{pred}}\).  Together with the
predicted-GS total-Casimir bridge, they carry both membership in the
predicted toy ground-state subspace and the corresponding
\(C_{\mathrm{pred}}\) total-Casimir identity through a ladder move,
without adding a separate membership assumption for the ladder image.
In the canonical orientation \(|\neg A|\le |A|\), the combined Lean
theorems
\[
\begin{gathered}
\texttt{tasaki23\_lowered\_predictedCasimirValue}\\
\texttt{\_of\_mem}\\
\texttt{\_bipartiteToyGroundStateSubspacePredicted},\\
\texttt{tasaki23\_raised\_predictedCasimirValue}\\
\texttt{\_of\_mem}\\
\texttt{\_bipartiteToyGroundStateSubspacePredicted}
\end{gathered}
\]
therefore give directly
\[
  \hat S_{\mathrm{tot}}^2(\hat S^\pm_{\mathrm{tot}}\Phi)
    =
  C_{\mathrm{pred}}\,\hat S^\pm_{\mathrm{tot}}\Phi
\]
from \(\Phi\in\mathcal G_{\mathrm{pred}}\), with no separate
predicted-Casimir hypothesis for the ladder image.

\subsection*{DONE: Tasaki §2.5 Theorem 2.3 for the standard bipartite
                coupling (sound Perron--Frobenius route, option (a))}
\label{subsec:tasaki-2-5-theorem-2-3-toy-done}

After the earlier saturated-ladder route was found unsound and removed,
the final statement \texttt{tasaki\_2\_5\_theorem\_2\_3} was closed for
the standard bipartite antiferromagnetic coupling
\(J=\texttt{bipartiteCoupling}\,A\) (canonical orientation
\(|\bar A|\le|A|\)) along the sound Perron--Frobenius route.  The proof
\texttt{tasaki\_2\_5\_theorem\_2\_3\_bipartiteToy} is sorry-free
(\texttt{\#print axioms} reports only \texttt{propext},
\texttt{Classical.choice}, \texttt{Quot.sound}).  Its four tiers are:

\begin{itemize}
\item \textbf{Per-sector existence} (Theorem 2.2, \#869): every
admissible magnetisation sector
\(M\in\texttt{tasaki23GroundStateSectors}\,A\,N
  =[\,\min(|A|,|\bar A|)N,\ \max(|A|,|\bar A|)N\,]\)
hosts a Marshall-positive Perron--Frobenius ground state
\(\sum_\sigma \varepsilon_\sigma v_\sigma\,|\sigma\rangle\),
\(v_\sigma>0\) (Tasaki (2.5.4)).

\item \textbf{TIER 4 --- common energy.}
\texttt{tasaki23\_common\_groundEnergy} fixes one energy
\(\mu\) realised in \emph{every} admissible sector.  Each sector ground
state carries the predicted total Casimir
\(C_{\mathrm{pred}}=|s_A-s_B|(|s_A-s_B|+1)\)
(per-sector overlap pin, \#3731), and the Casimir-only adjacent-sector
constancy \texttt{tasaki23\_pf\_sector\_energy\_eq\_of\_casimir} (\#3713)
propagates \(\mu\) across the interval by \(\le\)-induction.  This
consumes \emph{only} the predicted total Casimir, never the lowered
Marshall \texttt{site\_sum} positivity.

\item \textbf{TIER 5 --- global minimality.}
\texttt{tasaki23\_eigenvalue\_ge\_common} is a sector min--max engine: a
non-zero full-space eigenvector has a non-zero magnetisation-sector
restriction, whose real or imaginary part is a non-zero real sector
eigenvector; the Collatz--Wielandt per-sector bound
(\texttt{\_eigenvalue\_ge\_of\_marshallPositive}) closes the admissible
case, and the non-admissible bound \texttt{hOutside} the rest.  For the
toy coupling \texttt{hOutside} is the Lieb--Mattis lower bound
\texttt{tasaki23\_toy\_sector\_energy\_ge\_predicted} (\#3734): the
Marshall-positive sector ground state is a joint
\((\hat S^2_{\mathrm{tot}},\hat S^2_A,\hat S^2_{\bar A})\) eigenvector, so
the universal minimum-energy bound
\(E_{\mathrm{pred}}\le(\gamma_{\mathrm{tot}}-\gamma_A-\gamma_B)_{\mathrm{re}}\)
(\#3725) holds in \emph{every} sector.

\item \textbf{TIER 6 --- capstone.}
\(\mu=(\texttt{bipartiteToyMinEnergyPredicted}\,A\,N)_{\mathrm{re}}
   =E_{\mathrm{pred}}\); each sector energy is pinned to \(E_{\mathrm{pred}}\)
from below by the toy bound and from above by the predicted-energy
witness (\#3729 / \#3680), and the minimality engine closes the global
clause.
\end{itemize}

The only gap to the fully general coupling is the Lieb--Mattis energy
monotonicity \texttt{hOutside} for an arbitrary bipartite \(J\) --- the
single remaining hypothesis of the minimality engine.  The symmetric
orientation \(|A|<|\bar A|\) is the mirror image.

\textbf{General-\(J\) \texttt{hOutside} route (in progress).} Codex-validated,
this needs only the SU(2) \emph{inward ladder}, not reflection positivity:
an eigenvector in a non-admissible sector is moved to the admissible band
edge by \(\hat S^\mp_{\mathrm{tot}}\) (which commutes with \(H\), so the
energy is preserved), where the Collatz--Wielandt band-edge bound gives
\(\mu\le\mu_M\).  The non-vanishing of that ladder step is the magnitude
identity (\texttt{totalSpinSOpMinus\_mulVec\_normSq\_eq}), valid on any
weight-\(w\) vector \(\Phi\) (\(\hat S^3_{\mathrm{tot}}\Phi=w\Phi\)):
\[
  \|\hat S^-_{\mathrm{tot}}\Phi\|^2
    = \|\hat S^+_{\mathrm{tot}}\Phi\|^2 + 2\,w\,\|\Phi\|^2,
\]
from \(\hat S^+\hat S^- = \hat S^-\hat S^+ + 2\hat S^3\) and
\((\hat S^-)^\dagger=\hat S^+\).  Hence \(\hat S^-_{\mathrm{tot}}\Phi\neq0\)
when \(\Phi\neq0\) and \(w_{\mathrm{re}}>0\) (dually for \(\hat S^+\)).

Iterating this (\texttt{raise\_iterate\_ne\_zero},
\texttt{lower\_iterate\_ne\_zero}) carries a non-admissible-sector
eigenvector to the admissible band edge --- non-zero, weight-shifted, and
at the same energy --- and the Collatz--Wielandt band-edge bound
(\texttt{tasaki23\_admissible\_eigenvector\_ge}) then yields
\(\mu\le\mu_M\).  This is \texttt{tasaki23\_general\_hOutside}, which
discharges \texttt{hOutside} for an arbitrary bipartite \(J\); feeding it
to the minimality engine extends Theorem~2.3 from the standard coupling to
every connected bipartite antiferromagnetic coupling.

\textbf{General-\(J\) capstone (DONE).}
\texttt{tasaki\_2\_5\_theorem\_2\_3\_of\_bipartiteCompletePositive} assembles
the full statement for an arbitrary real symmetric non-negative bipartite
coupling positive on the complete bipartite graph: the common energy
\(\mu\) (TIER~4), the per-sector Theorem~2.2 ground state (\#869) with its
eigenvalue reconciled to \(\mu\) by the two-sided Collatz--Wielandt bound,
and global minimality from the engine plus the general-\(J\)
\texttt{hOutside}.  It is sorry-free (\texttt{\#print axioms} =
\texttt{propext}, \texttt{Classical.choice}, \texttt{Quot.sound}).  The only
residual restrictions are the canonical orientation \(|\bar A|\le|A|\) (the
mirror case is \(A\leftrightarrow\bar A\)) and \(s_B>0\) (non-degeneracy).

Reference: H.\ Tasaki, \emph{Physics and Mathematics of Quantum
Many-Body Systems}, Springer 2020, §2.5 Theorem 2.3, p.~42;
E.\ Lieb, D.\ Mattis, \emph{J. Math. Phys.}\ \textbf{3} (1962) 749.

\subsection*{Status: Tasaki §2.5 Theorem 2.4 (Mattis--Nishimori)}
\label{subsec:tasaki-2-5-theorem-2-4}

Tasaki's anisotropic Hamiltonian (2.5.14)
\[
  \hat H = \sum_{\{x,y\}\in B} \bigl(\hat S^{(1)}_x\hat S^{(1)}_y
    + \hat S^{(2)}_x\hat S^{(2)}_y + \lambda\,\hat S^{(3)}_x\hat S^{(3)}_y\bigr)
    + D\sum_{x}(\hat S^{(3)}_x)^2
\]
(Ising anisotropy \(\lambda\), crystal field \(D\)) is formalized as
\texttt{anisotropicHeisenbergS} (with two-site term \texttt{spinSDotXXZ} and
single-ion term \texttt{singleIonAnisotropyS}).  It is Hermitian for real
parameters and reduces to the isotropic Heisenberg Hamiltonian at
\(\lambda=1,\ D=0\).  Unlike the isotropic model it is only
\(U(1)\)-invariant: \texttt{anisotropicHeisenbergS\_commute\_totalSpinSOp3}
shows \([\hat H,\hat S^3_{\mathrm{tot}}]=0\) (each site \(\hat S^3\) commutes
with \(\hat S^3_{\mathrm{tot}}\), and
\(\texttt{spinSDotXXZ}=\texttt{spinSDot}+(\lambda-1)\hat S^3_x\hat S^3_y\)),
so \(\hat H\) preserves every magnetization sector
(\texttt{anisotropicHeisenbergS\_mulVec\_mem\_magSubspaceS\_of\_mem}).  It does
\emph{not} commute with \((\hat S_{\mathrm{tot}})^2\), so Marshall--Lieb--Mattis
cannot be applied directly; Theorem~2.4 instead argues sector-by-sector
(even/odd Perron--Frobenius for at-most-double degeneracy) plus a deformation
from the \(SU(2)\) point \((\lambda,D)=(1,0)\). (Issue \#3739, in progress.)

The single-site spin reversal \(F\) (permutation matrix of \texttt{Fin.rev},
\(k\mapsto N-k\); the \(\pi\)-rotation about axis~1) conjugates
\(F\hat S^3F=-\hat S^3\), \(F\hat S^\pm F=\hat S^\mp\)
(\texttt{spinReversalS\_conj\_*}).  The many-body product
\(\Theta=\bigotimes_xF\) (\texttt{manyBodyReversalS}, the configuration
reversal \(\sigma\mapsto\mathrm{rev}\circ\sigma\)) is an involution with
\(\Theta(\texttt{onSiteS}\,z\,A)\Theta=\texttt{onSiteS}\,z\,(FAF)\), giving the
\(M\leftrightarrow-M\) magnetization reflection
\(\Theta\hat S^3_{\mathrm{tot}}\Theta=-\hat S^3_{\mathrm{tot}}\)
(\texttt{manyBodyReversalS\_conj\_totalSpinSOp3}).  This reflection yields
\(E_M=E_{-M}\) for the crossing argument and, with the ground state in
\(\mathcal H_0\), the conclusion \(\hat S^3_{\mathrm{tot}}\Phi_{\mathrm{GS}}=0\).
The \(\hat S^3_{\mathrm{tot}}\Phi_{\mathrm{GS}}=0\) half is formalized as
\texttt{anisotropicHeisenbergS\_unique\_groundState\_has\_zero\_magnetization}
(PR \#3896): given \(\dim_{\mathbb C}(\text{GS-eigenspace})\le 1\), every nonzero
GS \(\Phi\) lies in a one-dimensional eigenspace, both \(\hat
S^3_{\mathrm{tot}}\Phi\) and \(\Theta\Phi\) lie there too, so
\(\hat S^3_{\mathrm{tot}}\Phi=\gamma\Phi\) and \(\Theta\Phi=\delta\Phi\) for some
scalars \(\gamma,\delta\in\mathbb C\); \(\Theta^2=1\) gives \(\delta^2=1\), and the
anticommutation \(\Theta\hat S^3_{\mathrm{tot}}=-\hat S^3_{\mathrm{tot}}\Theta\)
applied to \(\Phi\) gives \(\delta\gamma=-\gamma\delta\), i.e.\ \(2\delta\gamma=0\);
since \(\delta\ne 0\), \(\gamma=0\) and \(\hat S^3_{\mathrm{tot}}\Phi=0\).  The
uniqueness hypothesis itself (\(\dim\le 1\)) awaits the deformation argument from
\((\lambda,D)=(1,0)\) — separate forthcoming work.

The matrix-element form of the U(1) block-diagonality,
\texttt{anisotropicHeisenbergS\_apply\_eq\_zero\_of\_magSumS\_ne} (PR \#3898),
states \(\langle\sigma'|\hat H|\sigma\rangle=0\) whenever \texttt{magSumS}
\(\sigma'\ne\) \texttt{magSumS} \(\sigma\): evaluate the matrix commutation
\([\hat H,\hat S^3_{\mathrm{tot}}]=0\) at \((\sigma',\sigma)\), using diagonality
of \(\hat S^3_{\mathrm{tot}}\) with eigenvalues \texttt{magEigenvalueS},
yields \(\langle\sigma'|\hat H|\sigma\rangle(M_\sigma-M_{\sigma'})=0\); when
\texttt{magSumS} differs, \texttt{Nat} cast injectivity forces \(M_\sigma\ne M_{\sigma'}\)
and hence the matrix element vanishes.  This is the key building block for sector
decompositions in the eventual Theorem 2.4 obligation (2.a) uniqueness argument.

The vector-level corollary,
\texttt{anisotropicHeisenbergS\_magSectorProjection\_eigen} (PR \#3899), states
that for an \(\hat H\)-eigenvector \(v\) at \(\mu\), the sector projection
\(w := \texttt{magSectorEmbedding}(\texttt{magSectorRestriction}_{(M)}\,v)\) (\(v\)
on sector \(M\), zero elsewhere) is also an \(\hat H\)-eigenvector at \(\mu\).  In
the \texttt{mulVec} expansion of \(\hat H\,w\), the cross-sector matrix entries
vanish by PR \#3898, so only same-sector entries survive; combined with
\(w=v\) on sector \(M\) and \(\hat H\,v=\mu v\), this gives \(\hat H\,w=\mu w\)
both on sector \(M\) and off it (where both sides are zero).  This is the
vector-level enabler of sector-by-sector spectral analysis used in the
forthcoming uniqueness step of Theorem 2.4 obligation (2.a).

The companion linear-independence lemma
\texttt{linearIndependent\_pair\_of\_magSubspaceS\_distinct} (PR \#3900) states
that two nonzero vectors in distinct \(\hat S^3_{\mathrm{tot}}\) eigenspaces
\(M_1\ne M_2\) are linearly independent: from \(a\,\Phi_1 + b\,\Phi_2 = 0\), the
combination \(a\,\Phi_1\) lies in
\(\texttt{magSubspaceS}_{M_1}\cap\texttt{magSubspaceS}_{M_2}=\{0\}\)
(\texttt{magSubspaceS\_disjoint}), forcing each scalar to vanish.  The triple
extension \texttt{linearIndependent\_triple\_of\_magSubspaceS\_distinct} (PR
\#3901) applies mathlib's \texttt{Module.End.eigenvectors\_linearIndependent'} to
the linear endomorphism \texttt{Matrix.toLin'\,(totalSpinSOp3\,\(\Lambda\)\,N)}
to get a three-element linearly independent family for three pairwise distinct
sectors.  Combined with three nonzero vectors in three pairwise distinct
sectors, this contradicts Theorem 2.4 obligation (1)'s \(\dim\le2\) bound — the
contradiction lever in the sector decomposition + reflection argument.

The application
\texttt{anisotropicHeisenbergS\_threeLI\_of\_admis\_and\_nonadmis} (PR \#3902)
specialises the triple-LI to the symmetric \(|A|=|\neg A|\) configuration: given
an admissible-sector eigenvector \(\Phi\) (at \(\hat S^3_{\mathrm{tot}}=0\)) and
a non-admissible-sector eigenvector \(\Psi\) (at \(\hat S^3_{\mathrm{tot}}=M'\ne
0\)) at the same energy, the family \(\{\Phi,\Psi,\Theta\Psi\}\) is linearly
independent.  The reflected eigenvector \(\Theta\Psi\) lies in the
\(-M'\)-eigenspace via
\texttt{manyBodyReversalS\_mulVec\_mem\_magSubspaceS\_neg}; nonzeroness of
\(\Theta\Psi\) follows from \(\Theta^2=1\) and \(\Psi\ne0\); and the three
sectors \(0,M',-M'\) are pairwise distinct precisely because \(M'\ne0\).

The final contradiction step
\texttt{anisotropicHeisenbergS\_finrank\_le\_two\_no\_admis\_plus\_nonadmis} (PR
\#3903) packages this into the direct dichotomy: if the \(\mu\)-eigenspace has
\(\dim_{\mathbb C}\le 2\) and contains both an admissible-sector eigenvector
\(\Phi\) and a non-admissible-sector eigenvector \(\Psi\), then \texttt{False}.
The reflected vector \(\Theta\Psi\) is at the same energy via
\texttt{anisotropicHeisenbergS\_mul\_manyBodyReversalS}; the 3-LI family from
PR \#3902 lifts to the eigenspace as a subtype-valued LI family via
\texttt{LinearIndependent.of\_comp E.subtype}; mathlib's
\texttt{LinearIndependent.fintype\_card\_le\_finrank} then gives \(3\le\dim E\),
contradicting \(\dim E\le 2\).

The contrapositive
\texttt{anisotropicHeisenbergS\_eigvec\_nonadmis\_projection\_zero} (PR \#3904)
phrases this in projection form: under the same \(\dim E\le 2\) + admissible
eigenvector hypothesis, any other eigenvector \(\Psi\) at the same energy has
zero \texttt{magSectorEmbedding}\((\texttt{magSectorRestriction}_{(M)}\Psi)\) for
every non-admissible sector index \(M\) (i.e.\ \(2M\ne|\Lambda|\cdot N\)).  The
projection itself is again a \(\mu\)-eigenvector by PR \#3899, and lives in
\texttt{magSubspaceS}\(_\Lambda^N\) at eigenvalue \(M'=|\Lambda|\cdot N/2 - M\ne0\);
non-zeroness would activate PR \#3903 and give \texttt{False}, so the projection
vanishes.  Combined with the magnetization-sector decomposition
\texttt{eq\_sum\_magSectorEmbedding\_magSectorRestriction}, this places the
entire \(\mu\)-eigenspace inside the admissible sector — the last step before
the SU(2) symmetric \(\dim\le1\) conclusion.

The packaging
\texttt{anisotropicHeisenbergS\_eigvec\_in\_admis\_sector\_of\_finrank\_le\_two}
(PR \#3905) makes the inclusion explicit: under the same \(\dim E\le 2\) +
admissible eigenvector hypothesis, every eigenvector \(\Psi\) at \(\mu\) lies in
\texttt{magSubspaceS}\(_\Lambda^N\,0\).  Expanding
\(\Psi=\sum_{M\in\mathbb N}\texttt{magSectorEmbedding}(\texttt{magSectorRestriction}_{(M)}\Psi)\)
via \texttt{eq\_sum\_magSectorEmbedding\_magSectorRestriction}, the non-admissible
summands are killed by PR \#3904; the admissible summand (at \(2M=|\Lambda|\cdot
N\)) is in
\texttt{magSubspaceS}\(_\Lambda^N\,(|\Lambda|\cdot N/2-M)=\)\texttt{magSubspaceS}\(_\Lambda^N\,0\)
via \texttt{magSectorEmbedding\_mem\_magSubspaceS}; closure under sums gives
\(\Psi\) in the admissible sector.  Combined with the within-admissible-sector
Perron–Frobenius \(\dim=1\), this completes the SU(2) symmetric \(\dim\le1\)
chain.

The abstract finrank counting step
\texttt{anisotropicHeisenbergS\_finrank\_le\_one\_from\_admis\_pf} (PR \#3906)
finishes the dimension argument: given a within-admissible-sector hypothesis
\(\dim_{\mathbb C}(E\cap A)\le1\) (where
\(A=\texttt{magSubspaceS}\,\Lambda\,N\,0\) and \(E\) is the \(\mu\)-eigenspace
of \(\hat H\)), PR \#3905 gives \(E\subseteq A\), hence \(E\cap A=E\), and the
hypothesis transfers to \(\dim_{\mathbb C}E\le1\).  Once the within-admissible
\(\dim\le1\) is supplied (e.g.\ from the Perron–Frobenius argument on the
admissible sector matrix), the SU(2) symmetric \(\dim\le1\) capstone is
immediate.

The SU(2) specialisation
\texttt{heisenbergHamiltonianS\_finrank\_le\_one\_at\_SU2\_conditional} (PR
\#3907) bundles PR \#3906 with PR \#3897's eigenspace identification
\(\texttt{anisotropicHeisenbergS}\,J\,1\,0\,N=\texttt{heisenbergHamiltonianS}\,J\,N\):
under the same three hypotheses (\(\dim_{\mathbb C}E\le2\) at \(\mu\),
admissible eigenvector at \(\mu\), within-admissible \(\dim\le1\)) the
isotropic Heisenberg Hamiltonian's \(\mu\)-eigenspace has
\(\dim_{\mathbb C}\le1\).  This is the SU(2) endpoint of the deformation
argument: combined with PR \#3888's spin-\(1/2\) obligation (1) and Theorem 2.3's
within-admissible PF, the SU(2)-symmetric ground state is unique, and the
deformation can propagate.

The sector \texttt{LinearEquiv} \texttt{magSectorLinearEquiv} (PR \#3908) is the
algebraic bridge underlying the within-admissible PF input: it identifies
sector-indexed vectors \(\texttt{magConfigS}\,\Lambda\,N\,M\to\mathbb C\) with
the subspace
\(\texttt{magSubspaceS}\,\Lambda\,N\,(|\Lambda|\cdot N/2-M)\) of the full
Hilbert space via \texttt{magSectorEmbedding} (forward) and
\texttt{magSectorRestriction} (backward).  Round-trips are by
\texttt{magSectorRestriction\_magSectorEmbedding} and
\texttt{magSectorEmbedding\_magSectorRestriction\_of\_mem\_magSubspaceS};
linearity is by the existing \texttt{magSectorEmbedding\_add}/\texttt{\_smul}.
This bridge transfers within-sector-matrix Perron–Frobenius
\(\dim\le1\) results to the full-Hilbert-space \(\cap\)-sector form required by
PR \#3906.

The immediate finrank consequence
\texttt{magSubspaceS\_finrank\_eq\_sector} (PR \#3909) records the dimension
identity \(\dim_{\mathbb C}\,\texttt{magSubspaceS}\,\Lambda\,N\,(|\Lambda|\cdot
N/2 - M) = \dim_{\mathbb C}(\texttt{magConfigS}\,\Lambda\,N\,M\to\mathbb C)\),
obtained from \texttt{magSectorLinearEquiv} via
\texttt{LinearEquiv.finrank\_eq}.  This is the basic dimension-matching lemma
on which the eigenspace-level transfer of PF \(\dim\le1\) is built.

The eigenvalue side of the transfer requires the sector-restriction lemma at
arbitrary \(\mu\in\mathbb C\): the existing
\texttt{heisenbergHamiltonianSMatrixOnMagSector\_mulVec\_magSectorRestriction\_of\_full\_eigen}
takes \(\mu\in\mathbb R\) and lifts via the coercion \((\mu:\mathbb C)\); the
generalisation
\texttt{heisenbergHamiltonianSMatrixOnMagSector\_mulVec\_magSectorRestriction\_of\_full\_eigen\_complex}
(PR \#3910) is the \(\mu\in\mathbb C\) variant with the same proof structure
(the eigenvalue appears only as a scalar multiplier on \(\Psi\), and the matrix-entry
sector-zero lemma
\texttt{heisenbergHamiltonianS\_apply\_eq\_zero\_of\_magSumS\_ne} is
eigenvalue-independent).  This matches \texttt{End.eigenspace}'s
\(\mathbb C\)-valued eigenvalue convention used by the finrank infrastructure.

The companion embedding direction
\texttt{heisenbergHamiltonianS\_mulVec\_magSectorEmbedding\_complex} (PR \#3911)
is the dual \(\mu\in\mathbb C\) generalisation of
\texttt{heisenbergHamiltonianS\_mulVec\_magSectorEmbedding}: the embedding of a
sector \(\mathbb C\)-eigenvector is a full Hilbert-space
\(\mathbb C\)-eigenvector at the same \(\mu\).  Together these two PRs give the
two-sided \(\mu\in\mathbb C\) correspondence between the sector matrix and the
full Hilbert space, on which the within-admissible PF \(\dim\le1\) transfer
relies.

The eigenspace-level finrank equality
\texttt{heisenbergHamiltonianS\_sector\_matrix\_eigenspace\_finrank\_eq} (PR
\#3912) combines PRs \#3910/\#3911 with PR \#3908's
\texttt{magSectorLinearEquiv} into a \texttt{LinearEquiv} between
\(\texttt{End.eigenspace}(\hat{H}_{\text{sec}})\,\mu\) and
\((\texttt{End.eigenspace}(\hat{H})\,\mu)\cap\texttt{magSubspaceS}\,\Lambda\,N\,(|\Lambda|\cdot
N/2 - M)\); the forward map is \texttt{magSectorEmbedding} (lift, with
\(\mathbb C\)-eigvec property by \#3911), the backward map is
\texttt{magSectorRestriction} (restrict, with \(\mathbb C\)-eigvec property by
\#3910), and round-trips are by the existing
\texttt{magSectorRestriction\_magSectorEmbedding} and
\texttt{magSectorEmbedding\_magSectorRestriction\_of\_mem\_magSubspaceS}.
Applying \texttt{LinearEquiv.finrank\_eq} gives the dimension identity, the
within-admissible PF input transfer needed for the SU(2) symmetric
\(\dim\le1\) capstone.

The direct application
\texttt{heisenbergHamiltonianS\_eigenspace\_inf\_magSubspaceS\_finrank\_le\_one\_of\_sector}
(PR \#3913) lifts the inequality direction: if
\(\dim_{\mathbb C}\,\texttt{End.eigenspace}(\hat{H}_{\text{sec}})\,\mu\le1\),
then
\(\dim_{\mathbb C}\bigl((\texttt{End.eigenspace}(\hat{H})\,\mu)\cap
\texttt{magSubspaceS}\,\Lambda\,N\,(|\Lambda|\cdot N/2 - M)\bigr)\le1\).  This
connects the existing Theorem 2.3 chain's sector-matrix PF uniqueness to the
abstract PR \#3906 / PR \#3907 \(\dim\le1\) machinery for the SU(2) symmetric
capstone.

The bridge to the existing Theorem 2.3 chain output is supplied by
\texttt{heisenbergHamiltonianS\_finrank\_le\_one\_at\_SU2\_from\_sector\_pf}
(PR \#3915): under the hypotheses (i) \(\dim_{\mathbb C}E\le2\), (ii)
admissible eigenvector \(\Phi\), (iii) symmetric sector index \(M\) with
\(2M=|\Lambda|\cdot N\), (iv) within-admissible sector-matrix PF
\(\dim_{\mathbb C}\le1\), the isotropic Heisenberg \(\mu\)-eigenspace has
\(\dim_{\mathbb C}\le1\).  PR \#3913 turns (iv) into the form needed by PR
\#3907; the sector index equality
\((|\Lambda|\cdot N/2 - M : \mathbb C)=0\) follows from (iii) by
\texttt{push\_cast} + \texttt{linear\_combination}.

The natural physical packaging
\texttt{heisenbergHamiltonianS\_finrank\_le\_one\_at\_SU2\_symmetric\_card\_eq}
(PR \#3916) replaces the abstract sector-index hypothesis (iii) by the
sublattice card equality \(|A|=|\neg A|\): with \(M:=|A|\cdot N\), the
constraint \(2M=|\Lambda|\cdot N\) follows from
\texttt{tasaki23\_card\_filter\_A\_add\_card\_notA} and the card equality.
This is the standard physical statement form of the SU(2) symmetric
\(\dim\le1\) capstone.

The later deformation chain reduces spin-\(1/2\) obligation~(2) to a single
strict-gap input at the SU(2) point:
\[
  E_{M_{\mathrm{balanced}}}(1,0) < E_M(1,0)
  \qquad (M\ne M_{\mathrm{balanced}}).
\]
The capstone
\texttt{spinHalf\_anisotropicHeisenbergS\_obligation\_2\_single\_axiom}
packages the first-crossing/argmin argument and shows that this pointwise
strict gap is enough to contradict any proposed non-balanced ground-state
violation along the path from \((1,0)\) to \((\lambda',D')\).

The current structural refinement is
\texttt{strict\_gap\_at\_SU2\_of\_global\_unique} and its path-named form
\texttt{strict\_gap\_at\_path\_zero\_of\_global\_unique}.  Assume that the full
SU(2)-point ground eigenspace at \(E_{\rm full}\) has
\(\dim_{\mathbb C}\le1\), that the balanced sector attains the global minimum
\(E_{M_{\rm balanced}}=E_{\rm full}\), and that every non-balanced sector under
consideration has non-zero centered magnetization.  The general sector-minimum
bound gives \(E_{\rm full}\le E_M\).  If equality held for a non-balanced
sector, the sector-ground-vector construction would give two non-zero full
eigenvectors at \(E_{\rm full}\): one in centered magnetization \(0\), and one
in a non-zero centered-magnetization sector.  The existing two-sector/reflection
contradiction forbids this under \(\dim\le2\), hence certainly under
\(\dim\le1\).  Therefore \(E_{\rm full}<E_M\), and the balanced-sector equality
turns this into the strict gap.  This is implemented in
\texttt{AnisotropicHeisenbergStrictGapFromGlobalUniqueness.lean}.

Finally,
\texttt{spinHalf\_anisotropicHeisenbergS\_obligation\_2\_of\_SU2\_global\_unique}
uses this structural strict-gap theorem for every valid \(M\ne
M_{\rm balanced}\), then invokes the single-axiom capstone.  The remaining
mathematical input is therefore no longer the pointwise strict-gap family
itself, but the SU(2) global uniqueness plus balanced-ground-sector equality
expected from the Marshall--Lieb--Mattis/Casimir chain.  Reference: H.~Tasaki,
\emph{Physics and Mathematics of Quantum Many-Body Systems}, Springer 2020,
§2.5 Theorem 2.4, pp.~43--44.

The next refinement removes the balanced-ground-sector equality as an
independent input.  The theorem
\texttt{hermitianMinEigenvalue\_balanced\_eq\_full\_at\_SU2\_of\_global\_unique}
takes a full ground eigenvector \(\Phi\) of \(\hat H(1,0)\) at
\(E_{\rm full}\).  Global uniqueness and the already-formalized
zero-magnetization theorem imply
\(\hat S^3_{\rm tot}\Phi=0\).  If \(M_{\rm balanced}\) has zero centered
magnetization, then \(\Phi\) belongs to that magnetization sector, so the
sector restriction is non-zero and remains an eigenvector at \(E_{\rm full}\)
by the inverse-lift lemma.  Hence \(E_{M_{\rm balanced}}\le E_{\rm full}\).
The opposite inequality is the existing sector lower-bound
\texttt{hermitianMinEigenvalue\_anisotropicHeisenbergS\_full\_le\_sector},
so the two minima are equal.  This is implemented in
\texttt{AnisotropicHeisenbergBalancedFromGlobalUniqueness.lean}.

Composing this bridge with the PR~\#4018 capstone gives the one-input wrapper
shown below, implemented in
\path{AnisotropicHeisenbergSpinHalfObligation2SU2UniqueOnly.lean}.
\begin{Verbatim}
spinHalf_anisotropicHeisenbergS_obligation_2_of_SU2_global_unique_only
\end{Verbatim}
The spin-\(1/2\) obligation~(2) boundary now has one remaining SU(2)-endpoint
mathematical input: full ground eigenspace uniqueness at \(E_{\rm full}\).

The non-circular route to that endpoint uniqueness is started in the new
endpoint file.  Its main Lean entry points are:
\begin{Verbatim}
Theorem24SectorPFFromTheorem23.lean
Theorem24SU2GlobalUniquenessFromMLM.lean
marshallDiagonalOnMagSector
marshallDiagonalOnMagSector_mul_self
dressedHeisenbergSReMatrixOnMagSector_map_eq_marshall_conj_heisenberg
heisenbergHamiltonianSMatrixOnMagSector_finrank_le_one_of_marshall_positive
tasaki23GroundStateSectors_eq_singleton_of_card_eq
tasaki23PredictedTotalSpin_eq_zero_of_card_eq
tasaki23PredictedCasimirValue_eq_zero_of_card_eq
tasaki23_sector_lift_and_casimir_zero_of_card_eq
tasaki23_balanced_sector_matrix_finrank_le_one_of_common_min
heisenbergHamiltonianS_totalSpinSSquared_mulVec_eq_zero_of_sector_pf_zero_casimir
heisenbergHamiltonianS_totalSpinSSquared_mulVec_lower_landed_eq_zero_of_sector_pf
heisenbergHamiltonianS_totalSpinSSquared_mulVec_raise_landed_eq_zero_of_sector_pf
totalSpinSOpPlus_mulVec_eq_zero_of_totalSpinSSquared_mulVec_eq_zero_of_mem_zero_magSubspaceS
totalSpinSOpMinus_mulVec_eq_zero_of_totalSpinSSquared_mulVec_eq_zero_of_mem_zero_magSubspaceS
not_exists_totalSpinSOpMinus_image_of_zero_casimir_zero_magSubspaceS
not_exists_totalSpinSOpPlus_image_of_zero_casimir_zero_magSubspaceS
not_totalSpinSOpMinus_pow_mulVec_ne_zero_of_zero_casimir_zero_magSubspaceS
not_totalSpinSOpPlus_pow_mulVec_ne_zero_of_zero_casimir_zero_magSubspaceS
tasaki23_strict_hOutside_of_card_eq_zero_casimir_ladder_obstruction
hermitianMinEigenvalue_eq_common_of_eigenvector_and_global_lower
exists_tasaki23_common_energy_eq_hermitianMinEigenvalue
heisenbergHamiltonianS_full_eigenspace_finrank_le_one_of_outside_projection_zero
heisenbergHamiltonianS_outside_projection_zero_of_strict_sector_lower
heisenbergHamiltonianS_full_eigenspace_finrank_le_one_of_strict_sector_lower
exists_tasaki23_common_energy_and_heisenbergHamiltonianS_full_eigenspace_finrank_le_one
exists_tasaki23_common_energy_and_heisenbergHamiltonianS_full_eigenspace_finrank_le_one_of_casimir_ladder
exists_tasaki23_common_energy_and_heisenbergHamiltonianS_full_eigenspace_finrank_le_one_of_casimir_ladder_t23_pf
spinHalf_anisotropicHeisenbergS_obligation_2_of_MLM_casimir_ladder
spinHalf_anisotropicHeisenbergS_obligation_2_of_MLM_casimir_ladder_t23_pf
spinHalf_anisotropicHeisenbergS_obligation_2_strict_gap_of_MLM_casimir_ladder_t23_pf
\end{Verbatim}
The first group packages the sector Perron--Frobenius callback itself.  The
Marshall diagonal on a magnetization sector is involutive, and conjugating the
bare complex sector Hamiltonian by it gives the complexification of the
dressed real sector matrix.  A Theorem~2.3 Marshall-positive real sector
witness therefore becomes a positive eigenvector of the shifted dressed
non-negative irreducible sector matrix.  Perron--Frobenius geometric
simplicity, the shift correspondence, real-to-complex transfer, and the
Marshall similarity then give
\(\dim\ker(H_M-\mu)\le1\) for the bare complex sector matrix at the common
energy.  This is the new
\texttt{tasaki23\_balanced\_sector\_matrix\_finrank\_le\_one\_of\_common\_min}
bridge.

The next arithmetic name records the balanced-cardinality arithmetic: when
\(|A|=|\Lambda\setminus A|\), the Tasaki Theorem~2.3 admissible sector interval
\([\min(|A|,|B|)N,\max(|A|,|B|)N]\) is the singleton \(\{|A|N\}\).  The next
three names record the corresponding total-spin simplification:
the predicted total spin and total-Casimir value are both zero, and the
structural PF/Casimir lift of an admissible sector ground state specializes to
a full Heisenberg eigenvector with total-Casimir eigenvalue zero.  The next
bridge transfers this zero-Casimir property along any one-dimensional
sector-supported full eigenspace: if sector PF simplicity gives
\(\dim(E_\mu\cap\mathcal H_M)\le1\), and one non-zero vector on that line has
\((\hat S_{\rm tot})^2\Phi=0\), then every other full eigenvector at the same
\(\mu\) supported in the same sector also has zero total-Casimir image.  This is
the equality-case pin needed after the outside-sector ladder lands in the
balanced sector.  The lowering- and raising-landed variants combine this pin
with the existing non-vanishing inward-ladder iterates: if a non-zero
Heisenberg eigenvector reaches a sector whose ground line is already pinned by
sector PF simplicity to a zero-Casimir vector, then the landed vector is both
non-zero and killed by \((\hat S_{\rm tot})^2\).  These two bridges isolate the
left- and right-outside equality cases of \texttt{tasaki23\_general\_hOutside}.
The next six names formalize the representation-theoretic obstruction on the
balanced zero-magnetization line.  A vector with
\((\hat S_{\rm tot})^2\Phi=0\) and \(\hat S^3_{\rm tot}\Phi=0\) is killed by
both \(\hat S^+_{\rm tot}\) and \(\hat S^-_{\rm tot}\), using
\(\hat S^-_{\rm tot}\hat S^+_{\rm tot}
= (\hat S_{\rm tot})^2-(\hat S^3_{\rm tot})^2-\hat S^3_{\rm tot}\) and
\(\hat S^+_{\rm tot}\hat S^-_{\rm tot}
= (\hat S_{\rm tot})^2-(\hat S^3_{\rm tot})^2+\hat S^3_{\rm tot}\).  Since
\(\hat S^+_{\rm tot}\) and \(\hat S^-_{\rm tot}\) are adjoints, a non-zero
vector cannot simultaneously lie in the range of one total ladder and in the
kernel of the other.  Thus no positive total-ladder iterate can land as a
non-zero zero-Casimir vector in the zero magnetization sector, which is the
contradiction used after the landed vector is pinned to the balanced
zero-Casimir line.
The strict outside-sector wrapper then strengthens the Theorem~2.3
non-admissible-sector lower bound \(\mu\le\mu_M\) to \(\mu<\mu_M\).  If
equality held, the real sector eigenvector in the outside sector would lift to
a full Heisenberg eigenvector at \(\mu\).  Because the balanced-cardinality
case makes the admissible interval the singleton \(|A|N\), a positive
lowering or raising iterate moves that vector to the balanced zero
magnetization sector.  The landed-vector bridges pin the iterate to total
Casimir zero, and the singlet image obstruction rules out its non-zero
positive-step ladder-image origin.  This supplies the strict outside-sector
callback required by the full SU(2)-endpoint uniqueness bridge.
The following two theorems
identify the MLM common energy \(\mu\) with the Hermitian minimum of the full
Heisenberg Hamiltonian: the sector ground-state existence gives
\(\lambda_{\min}\le\mu\), while Theorem~2.3's global lower-bound callback
applied to a minimum eigenvector gives \(\mu\le\lambda_{\min}\).  The projection
bridge then isolates the remaining mathematical task.  Once every full ground
eigenvector has zero projection to sectors outside the singleton balanced
sector, the existing sector-matrix/sector-subspace finrank bridge transfers the
Perron--Frobenius sector simplicity to the full Hilbert-space eigenspace.

The strict-ordering interface says that if every real eigenvalue in every
sector \(M\ne |A|N\) is strictly above \(\mu\), then a full eigenvector at
\(\mu\) has zero outside projections: a non-zero outside projection would give
a real or imaginary sector eigenvector at the same \(\mu\), contradicting
strictness.  Consequently the first packaged theorem assembles the consumer
boundary: Theorem~2.3 supplies the common energy and full lower bound, balanced
cardinality supplies the singleton sector band, and the two remaining callbacks
are exactly balanced-sector PF simplicity and strict outside-sector ordering.
The final packaged theorem removes the strict-ordering callback by constructing
it from the balanced zero-Casimir vector and the total-ladder obstruction above.
The \(t23\_pf\) endpoint then supplies the balanced-sector PF callback from the
Theorem~2.3 witness described above, so the full SU(2)-endpoint uniqueness no
longer exposes either strict outside-sector ordering or balanced-sector PF
simplicity as a callback.  References:
H.~Tasaki, \emph{Physics and
Mathematics of Quantum Many-Body Systems}, Springer 2020, §2.5
Theorems~2.3--2.4, pp.~42--44; E.~Lieb and D.~Mattis, J. Math. Phys. 3 (1962),
749.

The spin-\(1/2\) wrapper
\texttt{spinHalf\_anisotropicHeisenbergS\_obligation\_2\_of\_MLM\_casimir\_ladder}
then composes this Heisenberg endpoint with the deformation capstone
\texttt{spinHalf\_anisotropicHeisenbergS\_obligation\_2\_of\_SU2\_global\_unique\_only}.
At the SU(2) point the anisotropic Hamiltonian is exactly the isotropic
Heisenberg Hamiltonian, formalized by
\texttt{anisotropicHeisenbergS\_one\_zero} and the eigenspace bridge
\texttt{anisotropicHeisenbergS\_at\_SU2\_eigenspace\_eq\_heisenbergHamiltonianS}.
Thus the Heisenberg Hermitian-minimum eigenspace bound
\(\dim E_{\min}\le1\) transports to \(\hat H(1,0)\), supplying the last
abstract SU(2)-global-uniqueness input of the spin-\(1/2\) obligation
boundary.  The \(t23\_pf\) wrapper also constructs that balanced-sector
Perron--Frobenius simplicity from the structural Theorem~2.3 witness, leaving
only the Theorem~2.3 predicate, the coupling and strict diagonal hypotheses,
balanced-cardinality arithmetic, and the deformation-path hypotheses at the
current spin-\(1/2\) obligation boundary.

The strict-gap wrapper
\texttt{spinHalf\_anisotropicHeisenbergS\_obligation\_2\_strict\_gap\_of\_MLM\_casimir\_ladder\_t23\_pf}
then exposes the same result in the direct Theorem~2.4 obligation form:
\[
  E_{M_{\rm balanced}}(\lambda',D')
  < E_M(\lambda',D')
  \qquad (M\ne M_{\rm balanced}).
\]
Its proof is only the final order-theoretic packaging step.  If this strict
inequality failed, then \(E_M(\lambda',D')\le
E_{M_{\rm balanced}}(\lambda',D')\), which is exactly the concrete violation
hypothesis ruled out by
\texttt{spinHalf\_anisotropicHeisenbergS\_obligation\_2\_of\_MLM\_casimir\_ladder\_t23\_pf}.
Thus the PR~\#4027 contradiction form and the strict-gap public form are
mathematically equivalent at this boundary; the latter is the cleaner consumer
surface for obligation~(2).

The next packaging step removes the last per-sector bookkeeping from that
consumer surface.  The theorem
\texttt{spinHalf\_anisotropicHeisenbergS\_strict\_gap\_all\_M\_of\_MLM\_casimir\_ladder\_t23\_pf}
states the same target strict gap uniformly for every non-empty
\(M\ne M_{\rm balanced}\).  A non-empty sector contains a configuration
\(\sigma\), hence \(M=\texttt{magSumS}\,\sigma\le |\Lambda|N\), so \(M\) lies in
the finite range required by the centered-nonzero hypothesis; the previous
strict-gap theorem then applies to this arbitrary \(M\).  The theorem
\texttt{spinHalf\_anisotropicHeisenbergS\_balanced\_eq\_full\_of\_MLM\_casimir\_ladder\_t23\_pf}
uses this uniform strict gap with
\texttt{hermitianMinEigenvalue\_balanced\_eq\_full\_of\_strict\_gap} and rewrites
the path endpoint \(t=1\) by
\texttt{anisotropicHeisenbergParametricPath\_one}.  It gives the target ground
energy identity
\[
  E_{M_{\rm balanced}}(\lambda',D')
  = \lambda_{\min}(\hat H(\lambda',D')).
\]
No new mathematical callback is introduced; the theorem is a consumer wrapper
around the Theorem~2.3/Perron--Frobenius strict-gap endpoint.

The next consumer theorem separates the last target-point PF input from the
linear-algebra packaging.  First,
\texttt{anisotropicHeisenbergS\_sector\_matrix\_eigenspace\_finrank\_eq} is the
anisotropic analogue of the Heisenberg sector-transfer theorem: zero-extension
and restriction give a linear equivalence between the eigenspace of the sector
matrix \(\hat H_M(\lambda,D)\) and the full eigenspace of
\(\hat H(\lambda,D)\) intersected with the corresponding magnetization
subspace.  Consequently
\texttt{anisotropicHeisenbergS\_eigenspace\_inf\_magSubspaceS\_finrank\_le\_one\_of\_sector}
transfers a sector-matrix \( \texttt{finrank}\le 1 \) hypothesis to the
admissible-intersection form used by the existing abstract theorem
\texttt{anisotropicHeisenbergS\_finrank\_le\_one\_from\_admis\_pf}.
At the target point, the theorem
\texttt{spinHalf\_anisotropicHeisenbergS\_target\_finrank\_le\_one\_of\_balanced\_sector\_pf\_and\_MLM\_casimir\_ladder\_t23\_pf}
combines this transfer with the global \( \texttt{finrank}\le 2 \) bound, the
balanced-sector ground vector, and the PR~\#4029 identity
\(E_{M_{\rm balanced}}(\lambda',D')=\lambda_{\min}(\hat H(\lambda',D'))\).
Thus balanced-sector Perron--Frobenius simplicity at the target point implies
full ground-state uniqueness.  Finally
\texttt{spinHalf\_anisotropicHeisenbergS\_target\_groundState\_zero\_magnetization\_of\_balanced\_sector\_pf\_and\_MLM\_casimir\_ladder\_t23\_pf}
composes this uniqueness bound with
\texttt{anisotropicHeisenbergS\_unique\_groundState\_has\_zero\_magnetization}
to obtain \(\hat S^3_{\rm tot}\Phi=0\) for every non-zero target ground state.

The target balanced-sector Perron--Frobenius input is then discharged at the
anisotropic point itself.  The file
\texttt{DressedAnisotropicMatrixOnMagSector.lean} defines the real
sector matrix, its Marshall-dressed conjugate, and the shifted matrix
\(B=cI-M_{\rm dressed}\).  Off the diagonal the anisotropic \(\lambda\)- and
\(D\)-terms agree with the Heisenberg exchange part, so the existing spin-\(S\)
bipartite reachability theorem gives positivity of a power of \(B\), hence
irreducibility.  The file \texttt{AnisotropicSectorPFAtMin.lean} applies
Perron--Frobenius to \(B\), uses the Collatz--Wielandt identification of the
PF eigenvalue with the spectral maximum of \(B\), and translates through the
strict shift, real-to-complex comparison, and Marshall similarity to prove
\[
  \dim_{\mathbb C}\ker
  \bigl(\hat H_M(\lambda,D)-\lambda_{\min}(\hat H_M(\lambda,D))I\bigr)
  \le 1 .
\]
For spin \(1/2\), the wrapper
\texttt{spinHalf\_anisotropicHeisenbergS\_balanced\_sector\_pf\_at\_target}
chooses the shift \(c\) strictly above all dressed diagonal entries and supplies
this sector PF bound at the balanced target sector.  The public wrappers
\texttt{spinHalf\_anisotropicHeisenbergS\_target\_finrank\_le\_one\_of\_MLM\_casimir\_ladder\_t23\_pf}
and
\texttt{spinHalf\_anisotropicHeisenbergS\_target\_groundState\_zero\_magnetization\_of\_MLM\_casimir\_ladder\_t23\_pf}
therefore remove the explicit balanced-sector PF callback from the current
spin-\(1/2\) Theorem~2.4 target boundary.

The \(D=0\) edge of this spin-\(1/2\) target boundary is handled in
\texttt{AnisotropicHeisenbergSpinHalfDNonnegBoundary.lean}.  The generic
spin-\(S\) parity-block Perron--Frobenius route assumes \(D>0\) because
single-ion \(\pm2\) parity moves need strict single-ion matrix entries.  At
spin \(1/2\), however, a configuration is indexed by
\(\mathrm{Fin}(1+1)\), so such a move is impossible:
\texttt{singleIonStepS\_spinHalf\_false}.  Thus every actual parity step is
either a transverse raise/lower step or a parity-bond step, and those
positivity inputs only require \(D\ge0\).  The resulting wrappers
\texttt{spinHalf\_anisotropicHeisenbergS\_eigenspace\_finrank\_le\_two\_at\_global\_min\_D\_nonneg},
\texttt{spinHalf\_anisotropicHeisenbergS\_obligation\_2\_of\_MLM\_casimir\_ladder\_t23\_pf\_D\_nonneg},
\texttt{spinHalf\_anisotropicHeisenbergS\_balanced\_eq\_full\_of\_MLM\_casimir\_ladder\_t23\_pf\_D\_nonneg},
\texttt{spinHalf\_anisotropicHeisenbergS\_target\_finrank\_le\_one\_of\_MLM\_casimir\_ladder\_t23\_pf\_D\_nonneg},
and
\texttt{spinHalf\_anisotropicHeisenbergS\_target\_groundState\_zero\_magnetization\_of\_MLM\_casimir\_ladder\_t23\_pf\_D\_nonneg}
extend the final spin-\(1/2\) target conclusion to
\(-1<\lambda<1\), \(D\ge0\).

The same \(D\ge0\) boundary is now available for general spin-\(S\) in
\texttt{AnisotropicHeisenbergSpinSDNonnegBoundary.lean}.  The new
\texttt{BondParityStepS} and \texttt{BondParityReachableS} relation removes
the single-ion strict-positive branch from the parity-block connectivity
argument.  Bond-only reachability is total on
\texttt{bipartiteCompleteGraphOf A}: the proof lowers both configurations to
a minimum magnetization representative using transverse transfers and
parity-bond lowering, then uses the existing equal-\(\texttt{magSumS}\)
raise/lower connectivity.  Consequently the shifted dressed parity-block
matrix is irreducible under \(D\ge0\), not merely \(D>0\).  The wrappers
\texttt{anisotropicHeisenbergS\_target\_finrank\_le\_one\_of\_MLM\_casimir\_ladder\_t23\_pf\_D\_nonneg\_general}
and
\texttt{anisotropicHeisenbergS\_target\_zero\_magnetization\_of\_MLM\_casimir\_ladder\_t23\_pf\_D\_nonneg\_general}
therefore extend the general spin-\(S\) case~(i) target endpoint to
\(-1<\lambda<1\), \(D\ge0\).

For general spin-\(S\), the \(\lambda=1\), \(D>0\) parity-block
Perron--Frobenius capstone is separated from the strict parity-bond route in
\texttt{ParityReachableNoParityBond.lean},
\texttt{ParityReachableNoParityBondTotal.lean},
\texttt{DressedAxisSwapIonParityLambdaOne.lean}, and
\texttt{DressedAxisSwapIonParityBlockIrreducibleLambdaOne.lean}.  The relation
\texttt{IonParityStepS} is generated only by transverse raise/lower moves and
single-ion \(\pm2\) moves.  The theorem \texttt{ionParityReachableS\_total}
proves totality on the bipartite complete graph for equal \(\texttt{magSumS}\)
parity when \(2\le N\): repeatedly lower \(\texttt{magSumS}\) by two using a
direct single-ion move, or first stack two unit occupations by transverse
transfers and then apply a single-ion lowering move.  The resulting
\texttt{shiftedDressedAxisSwappedReMatrixOnParityBlock\_isIrreducible\_lambda\_one\_D\_pos}
gives irreducibility at \(\lambda=1\), \(D>0\) without invoking the
same-direction parity-bond coefficient, which vanishes at \(\lambda=1\).

The spin-\(1/2\) \(\lambda=1\) edge of case~(i) is handled separately in
\texttt{AnisotropicHeisenbergSpinHalfLambdaOneBoundary.lean}.  This edge is
not a limit of the strict parity-bond Perron--Frobenius proof: the
same-direction parity-bond coefficient is \((1-\lambda)/4\), hence vanishes at
\(\lambda=1\).  Instead, at spin \(1/2\),
\[
  (S^{(3)})^2 = \frac14 I,
\]
so
\[
  \texttt{singleIonAnisotropyS}\;D\;1
    = \frac{D|\Lambda|}{4} I.
\]
The formal statements
\texttt{spinSOp3\_one\_sq\_eq\_quarter\_smul\_one},
\texttt{singleIonAnisotropyS\_spinHalf\_eq\_scalar}, and
\texttt{anisotropicHeisenbergS\_one\_D\_spinHalf\_eq\_heisenberg\_add\_scalar}
therefore identify \(\hat H(J,1,D)\) with the Heisenberg Hamiltonian plus a
scalar shift.  The generic scalar-shift lemmas
\texttt{eigenspace\_add\_smul\_one\_eq} and
\texttt{hermitianMinEigenvalue\_add\_smul\_one} transfer the Theorem~2.3/MLM
SU(2)-endpoint uniqueness to
\texttt{spinHalf\_anisotropicHeisenbergS\_lambda\_one\_finrank\_le\_one\_of\_MLM\_casimir\_ladder\_t23\_pf},
and the existing uniqueness-to-zero-magnetization theorem gives
\texttt{spinHalf\_anisotropicHeisenbergS\_lambda\_one\_groundState\_zero\_magnetization\_of\_MLM\_casimir\_ladder\_t23\_pf}.
For general spin-\(S\), the case~(i) target endpoint is now obtained on
\(-1<\lambda<1\), \(D\ge0\), by composing the \(D\ge0\) parity-block route with
the SU(2)-global-uniqueness wrappers and the MLM/Casimir endpoint below.  The
SU(2) point \((\lambda,D)=(1,0)\) is exposed by the direct public wrappers
\texttt{anisotropicHeisenbergS\_target\_finrank\_le\_one\_of\_MLM\_casimir\_ladder\_t23\_pf\_lambda\_one\_D\_zero\_general}
and
\texttt{anisotropicHeisenbergS\_target\_zero\_magnetization\_of\_MLM\_casimir\_ladder\_t23\_pf\_lambda\_one\_D\_zero\_general}.
The \(\lambda=1\), \(D>0\), \(2\le N\) edge is also live in
\texttt{AnisotropicHeisenbergSpinSLambdaOneBoundary.lean}: the ion-only
parity-block irreducibility replaces the vanished parity-bond coefficient,
and the path with first coordinate constantly \(1\) feeds the same
first-crossing argument.  The resulting endpoint wrappers are
\texttt{anisotropicHeisenbergS\_target\_finrank\_le\_one\_of\_MLM\_casimir\_ladder\_t23\_pf\_lambda\_one\_D\_pos\_general}
and
\texttt{anisotropicHeisenbergS\_target\_zero\_magnetization\_of\_MLM\_casimir\_ladder\_t23\_pf\_lambda\_one\_D\_pos\_general}.
Case~(ii) now has the conditional bridge below; the global
crossing/degeneracy replacement remains the next task.

For the degeneracy bound the Hamiltonian is gauged to the axis-swapped
\(\hat H'=U^{-1}\hat HU\) (the even/odd parity structure needed for
Perron--Frobenius).  Only the ground-state degeneracy count must be
transferred, so a \emph{similarity} suffices:
\texttt{matrix\_similar\_eigenspace\_finrank\_eq} shows that for invertible
\(U\), each \(\mu\)-eigenspace of \(U^{-1}\hat HU\) and of \(\hat H\) has equal
\texttt{finrank} (via \(v\mapsto U\,v\), a linear equivalence of eigenspaces).
This lets the explicit gauge unitary be isolated behind the
\texttt{AxisSwapUnitaryS} interface.  The concrete general spin-\(S\)
instance is now \texttt{axisSwapUnitarySSpinS}, the \(\pi/2\) rotation about
axis~1 built from \(\texttt{spinSRot1}\).  Specializing existing
\texttt{AxisSwapUnitaryS}-parameterized wrappers with this instance gives
\texttt{axisSwapUnitarySSpinS\_anisotropic\_axisSwapped\_eigenspace\_finrank\_eq},
the corresponding conditional block and minimum-eigenvalue wrappers, and the
strict case~(i) general-\(N\) degeneracy capstone
\texttt{anisotropicHeisenbergS\_eigenspace\_finrank\_le\_two\_unconditional\_general}.

The final target uniqueness step for general spin-\(S\) is now packaged as a
conditional bridge:
\texttt{anisotropicHeisenbergS\_target\_finrank\_le\_one\_of\_balanced\_sector\_pf}
combines the full target \(\texttt{finrank}\le2\) bound, equality of the
balanced sector minimum with the full ground energy, and balanced-sector
Perron--Frobenius simplicity to obtain full ground-eigenspace
\(\texttt{finrank}\le1\).  The companion
\texttt{anisotropicHeisenbergS\_target\_groundState\_zero\_magnetization\_of\_balanced\_sector\_pf}
then applies the existing uniqueness-implies-zero-magnetization theorem.  The
balanced-sector PF input itself is supplied in general spin-\(S\) form by
\texttt{anisotropicHeisenbergS\_balanced\_sector\_pf\_at\_target}, which
chooses a finite strict diagonal shift and applies the existing anisotropic
sector PF/minimum theorem.  Consequently
\texttt{anisotropicHeisenbergS\_target\_finrank\_le\_one\_of\_balanced\_eq\_full}
and
\texttt{anisotropicHeisenbergS\_target\_groundState\_zero\_magnetization\_of\_balanced\_eq\_full}
remove the explicit balanced-sector PF callback.  Finally,
\texttt{anisotropicHeisenbergS\_target\_finrank\_le\_one\_of\_strict\_gap}
and
\texttt{anisotropicHeisenbergS\_target\_groundState\_zero\_magnetization\_of\_strict\_gap}
compose the existing general-\(N\)
\texttt{hermitianMinEigenvalue\_balanced\_eq\_full\_of\_strict\_gap}
with those PF-free wrappers.  Thus the remaining general spin-\(S\) target
energy input is the strict gap
\(E_{M_{\mathrm{bal}}}<E_M\) over all non-balanced non-empty sectors.
The case~(ii) bridge in
\texttt{AnisotropicHeisenbergSpinSCaseIIConditionalBridge.lean} first records
that the same deformation path stays in the region \(\lambda\ge1\),
\(D\le0\):
\texttt{anisotropicHeisenbergParametricPath\_in\_case\_ii\_region}.  It then
repackages the same PF-free target step under case~(ii) hypotheses via
\texttt{anisotropicHeisenbergS\_case\_ii\_target\_finrank\_le\_one\_of\_balanced\_eq\_full}
and
\texttt{anisotropicHeisenbergS\_case\_ii\_target\_zero\_magnetization\_of\_balanced\_eq\_full}.
These first wrappers intentionally keep the missing global
crossing/degeneracy replacement explicit as balanced-sector/full-ground
equality plus full \(\texttt{finrank}\le2\); they do not reuse the case~(i)
axis-swap sign argument.  The follow-up strict-gap bridge in
\texttt{AnisotropicHeisenbergSpinSCaseIIStrictGapBridge.lean} replaces the
direct equality input by the strict sector gap
\(E_{M_{\mathrm{bal}}}<E_M\) over every non-balanced non-empty sector:
\texttt{anisotropicHeisenbergS\_case\_ii\_target\_finrank\_le\_one\_of\_strict\_gap}
and
\texttt{anisotropicHeisenbergS\_case\_ii\_target\_zero\_magnetization\_of\_strict\_gap}.
The follow-up no-full-finrank bridge in
\texttt{AnisotropicHeisenbergSpinSTargetFromStrictGapNoFullFinrank.lean}
removes that full \(\texttt{finrank}\le2\) input once the strict gap is
available.  The key lemma
\texttt{anisotropicHeisenbergS\_groundState\_mem\_balanced\_sector\_of\_strict\_gap}
shows directly that every full ground vector lies in the balanced
magnetization sector: each sector projection is again an eigenvector, and a
non-balanced projection would give a sector eigenvector below that sector's
Hermitian minimum.  Combining this containment with balanced-sector
Perron--Frobenius simplicity gives
\texttt{anisotropicHeisenbergS\_case\_ii\_target\_finrank\_le\_one\_of\_strict\_gap\_no\_full\_le\_two}
and
\texttt{anisotropicHeisenbergS\_case\_ii\_target\_zero\_magnetization\_of\_strict\_gap\_no\_full\_le\_two}.
The next case~(ii) bridge in
\texttt{AnisotropicHeisenbergSpinSCaseIIStrictGapFromCrossing.lean}
compresses that strict-gap derivation to a single target crossing
contradiction callback.  The theorem
\texttt{anisotropicHeisenbergS\_case\_ii\_strict\_gap\_all\_M\_of\_crossing\_contradiction}
says that if every non-balanced target crossing
\(E_M(\lambda,D)\le E_{M_{\mathrm{bal}}}(\lambda,D)\) is impossible under
\(1\le\lambda\), \(D\le0\), then \(E_{M_{\mathrm{bal}}}<E_M\) for every
non-balanced non-empty sector.  The wrappers
\texttt{anisotropicHeisenbergS\_case\_ii\_target\_finrank\_le\_one\_of\_crossing\_contradiction}
and
\texttt{anisotropicHeisenbergS\_case\_ii\_target\_zero\_magnetization\_of\_crossing\_contradiction}
then feed this strict gap to the no-full-finrank bridge above.  The next
bridge in
\texttt{AnisotropicHeisenbergSpinSCaseIICrossingFromPath.lean} evaluates the
case~(ii) deformation path at \(t=1\).  The theorem
\texttt{anisotropicHeisenbergS\_case\_ii\_crossing\_contradiction\_of\_path\_crossing\_contradiction}
turns a contradiction for every crossing along the path inside the region
\(1\le\lambda\), \(D\le0\) into the target crossing contradiction above.  Its
target wrappers
\texttt{anisotropicHeisenbergS\_case\_ii\_target\_finrank\_le\_one\_of\_path\_crossing\_contradiction}
and
\texttt{anisotropicHeisenbergS\_case\_ii\_target\_zero\_magnetization\_of\_path\_crossing\_contradiction}
therefore leave only the path crossing contradiction callback as the
remaining case~(ii) input.  The next bridge in
\texttt{AnisotropicHeisenbergSpinSCaseIIPathCrossingFromSet.lean} rewrites that
input in the first-crossing set language.  The theorem
\texttt{anisotropicHeisenbergS\_case\_ii\_path\_crossing\_contradiction\_of\_set\_contradiction}
observes that a pointwise path crossing is a witness that
\(\texttt{perMCrossingSet}\,M\cap[0,1]\) is non-empty.  Its target wrappers
\texttt{anisotropicHeisenbergS\_case\_ii\_target\_finrank\_le\_one\_of\_set\_contradiction}
and
\texttt{anisotropicHeisenbergS\_case\_ii\_target\_zero\_magnetization\_of\_set\_contradiction}
therefore leave only the crossing-set non-emptiness contradiction as the
remaining case~(ii) input.  The next bridge in
\texttt{AnisotropicHeisenbergSpinSCaseIICrossingSetFromFirst.lean} moves that
input to the achieved first-crossing point.  The theorem
\texttt{anisotropicHeisenbergS\_case\_ii\_crossing\_set\_contradiction\_of\_first\_crossing}
uses
\texttt{sInf\_perMCrossingSet\_inter\_Icc\_mem} to turn non-emptiness of
\(\texttt{perMCrossingSet}\,M\cap[0,1]\) into membership of its
\(\inf\).  Its target wrappers
\texttt{anisotropicHeisenbergS\_case\_ii\_target\_finrank\_le\_one\_of\_first\_crossing}
and
\texttt{anisotropicHeisenbergS\_case\_ii\_target\_zero\_magnetization\_of\_first\_crossing}
therefore leave only the first-crossing contradiction callback as the
remaining case~(ii) input.  The next bridge in
\texttt{AnisotropicHeisenbergSpinSCaseIIArgminFirstCrossing.lean} narrows
that input to the sector selected by the finite argmin extraction.  The
definition \texttt{caseIIArgminFirstCrossingContradiction} records the
available data: the chosen sector is in the valid range, is non-balanced, has
non-empty \(\texttt{perMCrossingSet}\cap[0,1]\), and minimises the
\(\inf\) of that set among all non-balanced sectors with non-empty crossing
set.  The theorem
\texttt{anisotropicHeisenbergS\_case\_ii\_crossing\_contradiction\_of\_argmin\_first\_crossing}
turns a target crossing at \(t=1\) into a non-empty crossing set, applies
\texttt{exists\_M\_chosen\_argmin\_per\_M\_first\_crossing}, converts the
total crossing-set argmin statement back to the sector-level set using
\texttt{perMCrossingSet\_total\_eq\_perMCrossingSet}, and invokes the argmin
first-crossing callback at
\(\texttt{sInf}(\texttt{perMCrossingSet}\,M_{\mathrm{chosen}}\cap[0,1])\).
Its target wrappers
\texttt{anisotropicHeisenbergS\_case\_ii\_target\_finrank\_le\_one\_of\_argmin\_first\_crossing}
and
\texttt{anisotropicHeisenbergS\_case\_ii\_target\_zero\_magnetization\_of\_argmin\_first\_crossing}
therefore leave only the argmin first-crossing contradiction callback as the
remaining case~(ii) input.  The bridge
\texttt{AnisotropicHeisenbergSpinSCaseIIArgminFinrank.lean} discharges that
callback from the standard first-crossing algebra once a full ground
eigenspace bound \(\texttt{finrank}\le2\) is available at the selected
first-crossing parameter.  The theorem
\texttt{caseIIArgminFirstCrossingContradiction\_of\_first\_crossing\_finrank\_bound}
uses the argmin property to obtain the balanced-sector ground-state equality
below and at the selected \(\inf\), uses the crossing equality at that
\(\inf\), and then applies the embedded two-sector contradiction to the
balanced and selected non-balanced sector ground vectors.  Its target wrappers
\texttt{anisotropicHeisenbergS\_case\_ii\_target\_finrank\_le\_one\_of\_first\_crossing\_finrank\_bound}
and
\texttt{anisotropicHeisenbergS\_case\_ii\_target\_zero\_magnetization\_of\_first\_crossing\_finrank\_bound}
therefore leave the first-crossing \(\texttt{finrank}\le2\) input, together
with the standard SU(2)-point strict-gap/ground-state data, as the next
case~(ii) boundary.  The subsequent path-global bridge in
\texttt{AnisotropicHeisenbergSpinSCaseIIPathGlobalFinrank.lean} removes the
selected-sector \(\inf\) expression from that input.  The theorem
\texttt{caseII\_first\_crossing\_finrank\_bound\_of\_path\_global\_finrank\_bound}
uses
\texttt{sInf\_perMCrossingSet\_inter\_Icc\_mem} to observe that each selected
first-crossing time lies in \([0,1]\).  Its target wrappers
\texttt{anisotropicHeisenbergS\_case\_ii\_target\_finrank\_le\_one\_of\_path\_global\_finrank\_bound}
and
\texttt{anisotropicHeisenbergS\_case\_ii\_target\_zero\_magnetization\_of\_path\_global\_finrank\_bound}
therefore leave a path-global full-ground-eigenspace
\(\texttt{finrank}\le2\) theorem as the remaining case~(ii) input.
The block-path bridge in
\texttt{AnisotropicHeisenbergSpinSCaseIIBlockPathFinrank.lean} rewrites this
input into the standard parity-block form.  The theorem
\texttt{caseII\_path\_global\_finrank\_bound\_of\_axisSwapped\_submatrix\_blocks\_le\_one}
uses
\texttt{anisotropicHeisenbergS\_eigenspace\_finrank\_le\_two\_of\_submatrix\_blocks\_le\_one\_general}
at the eigenvalue \(\mu\) equal to the full Hermitian minimum of
\(\hat H(\gamma(t))\).  Thus, for every \(t\in[0,1]\), two axis-swapped
parity-block submatrix \(\texttt{finrank}\le1\) bounds imply the required
full \(\texttt{finrank}\le2\) bound.  The wrappers
\texttt{anisotropicHeisenbergS\_case\_ii\_target\_finrank\_le\_one\_of\_axisSwapped\_submatrix\_blocks\_path}
and
\texttt{anisotropicHeisenbergS\_case\_ii\_target\_zero\_magnetization\_of\_axisSwapped\_submatrix\_blocks\_path}
therefore leave only the pathwise parity-block simplicity input.
\texttt{AnisotropicHeisenbergSpinSCaseIIBlockPFMin.lean} rewrites that
remaining input into the older PF/min conditional form.  The theorem
\texttt{axisSwappedAnisotropicHeisenbergS\_submatrix\_finrank\_le\_one\_at\_full\_min\_of\_pf\_min}
fixes a parity block \(p=0,1\).  If a real number \(\nu\) has
\(\texttt{finrank}\le1\) for that bare axis-swapped block and
\(\nu\) equals the block Hermitian minimum, then the same block has
\(\texttt{finrank}\le1\) at the full anisotropic ground energy.  Indeed,
\texttt{hermitianMinEigenvalue\_axisSwapped\_eq\_min\_block\_mins}
and
\texttt{axisSwapUnitarySSpinS\_hermitianMinEigenvalue\_anisotropic\_eq\_axisSwapped}
identify the full ground energy with the minimum of the two block minima.  If
this minimum is the chosen block minimum, the existing conditional theorem
\texttt{axisSwappedAnisotropicHeisenbergS\_submatrix\_finrank\_le\_one\_at\_min\_conditional}
applies; if it is strictly lower, then
\texttt{hermitian\_eigenspace\_eq\_bot\_of\_real\_lt\_min} makes the chosen
block eigenspace zero.  The path wrapper
\texttt{caseII\_axisSwapped\_submatrix\_blocks\_path\_of\_pf\_min} packages
this for both parity blocks along \(\gamma(t)\), and the target wrappers
\texttt{anisotropicHeisenbergS\_case\_ii\_target\_finrank\_le\_one\_of\_block\_pf\_min\_path}
and
\texttt{anisotropicHeisenbergS\_case\_ii\_target\_zero\_magnetization\_of\_block\_pf\_min\_path}
feed the result to the block-path bridge.  The same file introduces the
abbreviations \texttt{axisSwappedParityBlockPFMinAt} and
\texttt{axisSwappedParityBlockPFMinPath} for this repeated PF/min callback
shape, together with
\texttt{axisSwappedParityBlockStrictRawSupportPath},
\texttt{axisSwappedParityBlockLambdaOneRawSupportPath}, and
\texttt{axisSwappedParityBlockDZeroRawSupportPath} for the three
raw-support inputs along the path.  The theorem
\texttt{axisSwappedAnisotropicHeisenbergS\_submatrix\_pf\_min\_path\_of\_caseII\_raw\_support}
turns the fixed strict and zero-coefficient raw-support consumers into a
pathwise callback for one parity block.  The case~(ii) path-region inequalities
choose the strict branch, the \(\lambda=1\) ion-only branch, or the \(D=0\)
bond-only branch.  The corner \((\lambda,D)=(1,0)\) remains an explicit
PF/min callback because the strict coefficients used by those three consumers
all vanish there.  The wrapper
\texttt{caseII\_axisSwapped\_parityBlockPFMinPath\_of\_raw\_support} applies
this selector to the even and odd parity blocks, and
\texttt{caseII\_axisSwapped\_submatrix\_blocks\_path\_of\_raw\_support\_pf\_min}
feeds those two callbacks into the full-ground-energy block transfer.  The
target wrappers
\texttt{anisotropicHeisenbergS\_case\_ii\_target\_finrank\_le\_one\_of\_raw\_support\_pf\_min\_path}
and
\texttt{anisotropicHeisenbergS\_case\_ii\_target\_zero\_magnetization\_of\_raw\_support\_pf\_min\_path}
derive the raw-support zero-support hypothesis from the no-self-coupling and
bipartite-support assumptions, then expose the target statements directly from
the raw-support input surface.  The finite-block helper
\texttt{exists\_parityBlock\_dressed\_diag\_strict\_upper\_bound} supplies the
strict diagonal shift for any dressed real diagonal on a parity block.  Hence
\texttt{caseII\_axisSwapped\_parityBlockPFMinPath\_of\_reachability},
\texttt{caseII\_axisSwapped\_submatrix\_blocks\_path\_of\_reachability\_pf\_min},
and the target wrappers
\texttt{anisotropicHeisenbergS\_case\_ii\_target\_finrank\_le\_one\_of\_reachability\_pf\_min\_path}
and
\texttt{anisotropicHeisenbergS\_case\_ii\_target\_zero\_magnetization\_of\_reachability\_pf\_min\_path}
replace the lower-level raw-support path inputs by reachability totality plus
the explicit corner PF/min callbacks.  The reachability-totality bridge in
\texttt{AnisotropicHeisenbergSpinSCaseIIBlockReachability.lean} lifts
full-configuration \(\texttt{ParityReachableS}\),
\(\texttt{IonParityReachableS}\), and \(\texttt{BondParityReachableS}\) paths
to the corresponding fixed-parity-block relations by parity preservation.  Its
totality wrappers then supply the strict and bond-only block reachability
inputs from \(1\le N\), and the ion-only block reachability input from
\(2\le N\).  Consequently
\texttt{caseII\_axisSwapped\_parityBlockPFMinPath\_of\_total\_reachability}
and
\texttt{caseII\_axisSwapped\_submatrix\_blocks\_path\_of\_total\_reachability\_pf\_min}
leave only the two explicit \((\lambda,D)=(1,0)\) corner PF/min callbacks.
\texttt{AnisotropicHeisenbergSpinSCaseIIParityGauge.lean} begins the
case~(ii) PF sign layer.  In the region \(\lambda\ge1\), \(D\le0\), the
Marshall sign alone does not put the parity-bond and single-ion off-diagonal
terms into the case~(i) shifted-matrix sign convention: those moves change
\(\texttt{magSumS}\) by \(\pm2\).  On a fixed parity block the additional
gauge
\[
  \Xi(\sigma)=(-1)^{\texttt{magSumS}(\sigma)/2}
\]
agrees with Tasaki's block convention
\((-1)^{(\texttt{magSumS}(\sigma)-p)/2}\) up to a constant sign.  The lemmas
\texttt{caseIIParityGaugeSign\_mul\_eq\_one\_of\_magSumS\_eq},
\texttt{caseIIParityGaugeSign\_mul\_eq\_neg\_one\_of\_magSumS\_add\_two},
and
\texttt{caseIIParityGaugeSign\_mul\_eq\_neg\_one\_of\_add\_two\_magSumS}
record exactly the needed sign behavior: equal \(\texttt{magSumS}\) gives
product \(1\), while a difference of \(2\) gives product \(-1\).  The diagonal
matrix \texttt{caseIIParityGaugeDiagonalOnParity} squares to the identity by
\texttt{caseIIParityGaugeDiagonalOnParity\_mul\_self}, so it can be combined
with the Marshall diagonal in the next case~(ii) shifted PF matrix layer.
\texttt{AnisotropicHeisenbergSpinSCaseIIParityGaugedMatrix.lean} constructs
that structural matrix layer.  The real sign
\texttt{caseIIParityGaugeSignReal} is the real-valued version of \(\Xi\), and
\texttt{caseIICombinedGaugeSignOnParity} multiplies it with the Marshall sign
on each parity block.  The diagonal
\texttt{caseIICombinedGaugeDiagonalOnParity} squares to the identity by
\texttt{caseIICombinedGaugeDiagonalOnParity\_mul\_self}.  Rather than
duplicating the complex bare matrix, the Lean definition
\texttt{caseIIParityGaugedAxisSwappedReMatrixOnParityBlock} applies \(\Xi\) on
the left and right of the existing Marshall-dressed real matrix:
\[
  R^{(ii)}_{A,\lambda,D,p}(\sigma,\tau)
    = \Xi(\sigma)\,
      \operatorname{Re}\!\bigl(\widetilde H'_{A,\lambda,D}(\sigma,\tau)\bigr)\,
      \Xi(\tau).
\]
Its diagonal agrees with the dressed real matrix by
\texttt{caseIIParityGaugedAxisSwappedReMatrixOnParityBlock\_apply\_diag}, and
it is symmetric for real coupling data by
\texttt{caseIIParityGaugedAxisSwappedReMatrixOnParityBlock\_isSymm\_of\_real}.
The shifted PF candidate is
\[
  B^{(ii)}_{A,\lambda,D,p,c}
    = c I - R^{(ii)}_{A,\lambda,D,p},
\]
implemented as
\texttt{shiftedCaseIIParityGaugedAxisSwappedReMatrixOnParityBlock}.  The file
records its diagonal, off-diagonal, strict-shift diagonal positivity, and
symmetry facts.  The remaining case~(ii) PF work is to prove sign
non-negativity, reachability/irreducibility, and PF/min identification for
this shifted matrix in the region \(\lambda\ge1\), \(D\le0\).
\texttt{AnisotropicHeisenbergSpinSCaseIIParityGaugedSigns.lean} adds the
corresponding structural sign-transfer API.  The real gauge product lemmas
\texttt{caseIIParityGaugeSignReal\_mul\_eq\_one\_of\_magSumS\_eq},
\texttt{caseIIParityGaugeSignReal\_mul\_eq\_neg\_one\_of\_magSumS\_add\_two},
and
\texttt{caseIIParityGaugeSignReal\_mul\_eq\_neg\_one\_of\_add\_two\_magSumS}
are the real analogues of the complex sign layer.  They yield entry
identities saying that the parity-gauged real matrix equals the dressed real
entry on equal-\(\texttt{magSumS}\) pairs and equals its negative when the two
magnetization sums differ by \(2\).  Consequently dressed negativity on
transverse moves and dressed positivity on the \(\pm2\) moves both become
parity-gauged non-positivity.  The theorem
\texttt{shiftedCaseIIParityGaugedAxisSwappedReMatrixOnParityBlock\_nonneg}
packages the structural implication: a diagonal shift bound plus
off-diagonal non-positivity of \(R^{(ii)}\) makes \(cI-R^{(ii)}\) entrywise
non-negative.  The strict transfer lemmas
\texttt{shiftedCaseIIParityGaugedAxisSwappedReMatrixOnParityBlock\_pos\_of\_magSumS\_eq},
\texttt{shiftedCaseIIParityGaugedAxisSwappedReMatrixOnParityBlock\_pos\_of\_magSumS\_add\_two},
and
\texttt{shiftedCaseIIParityGaugedAxisSwappedReMatrixOnParityBlock\_pos\_of\_add\_two\_magSumS}
turn the corresponding dressed strict signs into strict positivity of the
shifted entry.  The local sign file
\texttt{AnisotropicHeisenbergSpinSCaseIILocalSigns.lean} supplies those
dressed strict signs and feeds them into this transfer API.  For transverse
\(\pm1,\mp1\) moves,
\texttt{shiftedCaseIIParityGaugedAxisSwappedReMatrixOnParityBlock\_pos\_of\_raiseLowerStepS}
uses the positive \((1+\lambda)/4\) coefficient and the Marshall sign flip;
for parity-bond \(\pm1,\pm1\) moves,
\texttt{shiftedCaseIIParityGaugedAxisSwappedReMatrixOnParityBlock\_pos\_of\_parityBondStepS}
uses the negative \((1-\lambda)/4\) coefficient together with the extra
case~(ii) parity gauge; and for same-site \(\pm2\) moves,
\texttt{shiftedCaseIIParityGaugedAxisSwappedReMatrixOnParityBlock\_pos\_of\_singleIonStepS}
uses \(D<0\) against the negative \((\hat S^2)^2\) off-diagonal entry.  Thus
\texttt{AnisotropicHeisenbergSpinSCaseIIBlockIrreducible.lean} defines the
block-level closure
\texttt{parityReachableSOnBlock} and packages the three local signs as
\texttt{shiftedCaseIIParityGaugedAxisSwappedReMatrixOnParityBlock\_pos\_of\_parityStepS}.
Entrywise non-negativity plus block reachability gives a positive matrix-power
entry, and reachability totality plus strict diagonal shift gives the
conditional irreducibility theorem
\texttt{shiftedCaseIIParityGaugedBlock\_isIrreducible\_of\_blockReachable\_total}.
\texttt{AnisotropicHeisenbergSpinSCaseIIBlockNonneg.lean} then lowers the
explicit entrywise non-negativity input to a diagonal shift bound and a
\(\texttt{magSumS}\) support split for off-diagonal block pairs.  Equal
\(\texttt{magSumS}\) pairs use dressed non-positivity; target/source
\(\texttt{magSumS}\) raised by two use dressed non-negativity together with the
case~(ii) parity-gauge sign flip; all other off-diagonal pairs are supplied as
zero-support.  The wrappers
\texttt{caseIIParityGaugedBlock\_offdiag\_nonpos\_of\_magSum\_support},
\texttt{shiftedCaseIIParityGaugedBlock\_nonneg\_of\_magSum\_support}, and
\texttt{shiftedCaseIIParityGaugedBlock\_irreducible\_of\_magSum\_support}
feed this support split through shifted entrywise non-negativity and the block
irreducibility layer.
\texttt{AnisotropicHeisenbergSpinSCaseIIStepSupport.lean} then supplies the
next consumer layer: it turns a step-or-zero split into those
\(\texttt{magSumS}\) support hypotheses.  Equal-\(\texttt{magSumS}\)
off-diagonal pairs are represented by a transverse
\(\texttt{RaiseLowerStepS}\) or zero support; target/source
\(\texttt{magSumS}\) raised by two are represented by
\(\texttt{ParityBondStepS}\), \(\texttt{SingleIonStepS}\), or zero support.
The step signs are exactly the local case~(ii) dressed signs already proved
above, weakened to non-positive/non-negative inequalities; dressed zero support
is transferred through the parity gauge by
\texttt{caseIIParityGaugedAxisSwappedReMatrixOnParityBlock\_eq\_zero\_of\_dressed\_zero}.
The wrappers
\texttt{shiftedCaseIIParityGaugedBlock\_nonneg\_of\_step\_support} and
\texttt{shiftedCaseIIParityGaugedBlock\_irreducible\_of\_step\_support}
therefore reduce the block PF input to a raw total-Hamiltonian witness/zero
classification.  The next paragraphs supply that raw classification and then
record the boundary move-set choices when a strict local coefficient vanishes;
the following PF/min layer then converts shifted non-negative irreducible
blocks into bare parity-block PF/min witnesses.

\texttt{AnisotropicHeisenbergSpinSCaseIIRawSupport.lean} proves that raw
classification under the support-zero assumption
\[
  J(x,y)=0 \quad\text{whenever}\quad
  (x,y)\notin E(\texttt{bipartiteCompleteGraphOf}~A).
\]
For a raw off-diagonal dressed real entry of the axis-swapped Hamiltonian,
the local zero tests show that a nonzero bond contribution must be one of the
two local \(S^1S^1+\lambda S^2S^2\) moves: a transverse
\texttt{RaiseLowerStepS} when \(\texttt{magSumS}\) is preserved, or a
\texttt{ParityBondStepS} when one configuration has \(\texttt{magSumS}\) two
larger than the other.  A nonzero single-ion contribution must be a
\texttt{SingleIonStepS}; all other off-diagonal entries vanish.  The wrappers
\texttt{dressedAxisSwappedReMatrix\_raiseLower\_or\_zero\_of\_magSum\_eq},
\texttt{dressedAxisSwappedReMatrix\_parity\_or\_single\_or\_zero\_of\_magSum\_add\_two},
and
\texttt{dressedAxisSwappedAnisotropicHeisenbergSReMatrix\_apply\_eq\_zero\_of\_not\_magSum\_step}
therefore supply the step-or-zero hypotheses directly.  Finally,
\texttt{shiftedCaseIIParityGaugedBlock\_nonneg\_of\_raw\_support} and
\texttt{shiftedCaseIIParityGaugedBlock\_irreducible\_of\_raw\_support} feed
that raw support classification into the case~(ii) shifted block
non-negativity and conditional irreducibility layers.  The next boundary layer
handles the zero-coefficient move-set variants, and the following PF/min layer
turns those shifted block inputs into bare parity-block PF/min witnesses.

\texttt{AnisotropicHeisenbergSpinSCaseIIBoundaryMoveSets.lean} records the
boundary move-set variants where one strict local coefficient vanishes.  At
\(\lambda=1\), the parity-bond coefficient is zero, so the boundary input uses
the ion-only relation generated by transverse moves and
\(\texttt{SingleIonStepS}\); the wrappers
\texttt{shiftedCaseIIBlock\_nonneg\_of\_ion\_step\_support\_lambda\_one} and
\texttt{shiftedCaseIIBlock\_irreducible\_of\_ion\_step\_support\_lambda\_one}
package entrywise non-negativity and conditional irreducibility from
\texttt{ionParityReachableSOnBlock}.  At \(D=0\), the single-ion term has zero
coefficient; the lemma
\texttt{dressedAxisSwappedReMatrix\_zero\_of\_singleIonStep\_D\_zero} removes
that branch from the raw classification, and
\texttt{shiftedCaseIIBlock\_nonneg\_of\_raw\_support\_D\_zero} plus
\texttt{shiftedCaseIIBlock\_irreducible\_of\_raw\_support\_D\_zero} feed the
result to the bond-only relation \texttt{bondParityReachableSOnBlock}.  The
raw \(\lambda=1\) endpoint is also discharged: the local lemma
\texttt{spinSDotXXZSwap\_apply\_eq\_zero\_of\_parityBondStepS\_witness\_lambda\_one}
uses the ladder form to kill parity-bond entries by the coefficient
\((1-\lambda)/4=0\), and
\texttt{dressedAxisSwappedReMatrix\_zero\_of\_parityBondStep\_lambda\_one}
removes that branch from the total dressed real entry.  Hence
\texttt{shiftedCaseIIBlock\_nonneg\_of\_raw\_support\_lambda\_one} and
\texttt{shiftedCaseIIBlock\_irreducible\_of\_raw\_support\_lambda\_one} feed
raw support directly to the ion-only relation.  The
\texttt{AnisotropicHeisenbergSpinSCaseIIParityBlockPFMin.lean} layer then
converts these shifted parity-gauged block inputs into bare parity-block
PF/min witnesses.  It proves that the real combined parity/Marshall gauge
diagonal complexifies to an involutive diagonal and that the complexified
case~(ii) parity-gauged block is similar to the bare axis-swapped parity
block:
\[
  R_{\mathrm{case~II},p}^{\mathbb C}
  = G_p\,\widehat H'_{p}\,G_p .
\]
For a fixed parity block, Perron--Frobenius supplies a positive eigenvector and
a one-dimensional eigenspace for the shifted non-negative irreducible matrix
\(B=cI-R_{\mathrm{case~II},p}\).  The shift changes the eigenvalue from
\(r\) to \(c-r\), and Collatz--Wielandt identifies \(c-r\) with the Hermitian
minimum of the complexified real block.  The similarity then transfers the
\(\texttt{finrank}\le1\) statement to the bare axis-swapped block at its
Hermitian minimum.  The wrappers
\texttt{axisSwappedAnisotropicHeisenbergS\_submatrix\_pf\_min\_of\_caseII\_raw\_support},
\texttt{axisSwappedAnisotropicHeisenbergS\_submatrix\_pf\_min\_of\_caseII\_raw\_support\_lambda\_one},
and
\texttt{axisSwappedAnisotropicHeisenbergS\_submatrix\_pf\_min\_of\_caseII\_raw\_support\_D\_zero}
consume the strict raw-support layer, the \(\lambda=1\) ion-only boundary,
and the \(D=0\) bond-only boundary, respectively.  The remaining pathwise
case~(ii) input is to supply the appropriate strict or zero-coefficient
reachability hypotheses along the deformation path.

The new general spin-\(S\) obligation-(2) wrapper supplies that strict-gap
callback from one SU(2)-endpoint input.  First,
\texttt{anisotropicHeisenbergS\_eigenspace\_finrank\_le\_two\_at\_global\_min\_general}
combines the axis-swapped parity-block capstone with the minimum-eigenvalue
axis-swap identity, and
\texttt{anisotropicHeisenbergS\_eigenspace\_finrank\_le\_two\_at\_global\_min\_path\_general}
specializes it to the deformation path \(\gamma(t)\), \(0<t\le1\).  The
first-crossing contradiction
\texttt{anisotropicHeisenbergS\_obligation\_2\_of\_SU2\_global\_unique\_only\_general}
then uses the existing general-\(N\)
\texttt{strict\_gap\_at\_path\_zero\_of\_global\_unique} and
\texttt{hermitianMinEigenvalue\_balanced\_eq\_full\_at\_SU2\_of\_global\_unique}
to discharge the SU(2)-endpoint strict-gap and balanced-ground inputs from
full SU(2)-point ground-eigenspace uniqueness.  Negating the target violation
gives
\texttt{anisotropicHeisenbergS\_strict\_gap\_all\_M\_of\_SU2\_global\_unique\_general},
and the final wrappers
\texttt{anisotropicHeisenbergS\_target\_finrank\_le\_one\_of\_SU2\_global\_unique\_general}
and
\texttt{anisotropicHeisenbergS\_target\_groundState\_zero\_magnetization\_of\_SU2\_global\_unique\_general}
feed this strict gap to the PF-free target endpoint.

The general spin-\(S\) MLM endpoint file removes this remaining SU(2)
uniqueness callback in strict case~(i).  The transport theorem
\texttt{anisotropicHeisenbergS\_SU2\_ground\_eigenspace\_finrank\_le\_one\_of\_heisenberg\_general}
rewrites \(\hat H(1,0)\) to the isotropic Heisenberg Hamiltonian and transfers
the Heisenberg Hermitian-minimum eigenspace bound to the anisotropic
SU(2)-point eigenspace.  The wrappers
\texttt{anisotropicHeisenbergS\_target\_finrank\_le\_one\_of\_MLM\_casimir\_ladder\_t23\_pf\_general}
and
\texttt{anisotropicHeisenbergS\_target\_zero\_magnetization\_of\_MLM\_casimir\_ladder\_t23\_pf\_general}
compose that transport with
\texttt{exists\_tasaki23\_common\_energy\_and\_heisenbergHamiltonianS\_full\_eigenspace\_finrank\_le\_one\_of\_casimir\_ladder\_t23\_pf}
and the SU(2)-global-uniqueness target wrappers above.  Thus the general
spin-\(S\) strict case~(i) target uniqueness and zero-magnetization endpoint no
longer exposes an abstract SU(2)-endpoint uniqueness hypothesis; the
\texttt{\_D\_nonneg\_general} wrappers above extend the same endpoint to
\(D\ge0\).  The SU(2)-point wrappers and the
\(\lambda=1\), \(D>0\), \(2\le N\) ion-only wrappers cover the general
spin-\(S\) \(\lambda=1\), \(D\ge0\) case~(i) boundary.  Tasaki case~(ii) now
has the conditional path/target bridge, the strict-gap bridge, the
no-full-finrank bridge, the crossing-callback bridge, and the path-callback
bridge, the crossing-set bridge, the first-crossing bridge, the
argmin-first-crossing bridge, the first-crossing finrank bridge, the
path-global finrank bridge, the block-path finrank bridge, and the
block-PF/min bridge described above, plus the first parity-gauge sign layer,
the parity-gauged shifted-matrix structural layer, and the shifted-matrix
sign-transfer layer.  The local sign layer now proves strict shifted-entry
positivity for the transverse, parity-bond, and single-ion elementary moves in
the corresponding strict local regimes: \(-1<\lambda\) for transverse moves,
\(1<\lambda\) for parity-bond moves, and \(D<0\) for single-ion moves.  The
block irreducibility layer now turns those step signs into block matrix-power
positivity from block-level parity reachability and into conditional
irreducibility from entrywise non-negativity, strict diagonal shift, and
reachability totality.  The block non-negativity bridge now lowers that
entrywise non-negativity input to a diagonal shift bound and a
\(\texttt{magSumS}\) off-diagonal support split, distinguishing equal
\(\texttt{magSumS}\), target raised by two, source raised by two, and
zero-support pairs.  The raw support layer now proves this split directly from
the total Hamiltonian under the bipartite support-zero assumption on \(J\).  The
boundary move-set layer now supplies the ion-only \(\lambda=1\) and bond-only
\(D=0\) variants, including raw-support wrappers at both zero-coefficient
boundaries.  The parity-block PF/min layer now transfers shifted case~(ii)
parity-gauged non-negative irreducible blocks to bare axis-swapped
parity-block PF/min witnesses in the strict raw-support region, at the
\(\lambda=1\) ion-only boundary, and at the \(D=0\) bond-only boundary.  The
block-PF/min layer now selects among these fixed consumers along the
case~(ii) deformation path, packages the even and odd parity-block callbacks,
feeds them to the full-ground-energy block transfer, and exposes the target
uniqueness and zero-magnetization wrappers directly from that raw-support
surface.  It now also supplies the strict diagonal shift from finite
parity-block boundedness, and the block reachability totality bridge supplies
the strict, ion-only, and bond-only reachability inputs for \(2\le N\).  The
corner bridge now separates the exact SU(2) corner from the non-corner path:
away from \((\lambda,D)=(1,0)\), the strict, ion-only, and bond-only PF/min
consumers are still selected pointwise from total reachability; at the corner,
the path-global full \(\texttt{finrank}\le2\) input is supplied directly from
a full SU(2)-point ground-eigenspace \(\texttt{finrank}\le1\) callback.  This
avoids asking each whole parity block to be PF-simple at its own block minimum
at the SU(2) point, which is stronger than the crossing argument needs.  The
total-reachability target layer now composes this corner bridge with the
MLM/Casimir/Theorem~2.3 endpoint: the endpoint supplies full SU(2)-point
uniqueness, and that same uniqueness input is reused for the exact-corner
\(\texttt{finrank}\le2\) branch, the balanced SU(2) sector/full-ground
equality, and the strict non-balanced sector gap at \(t=0\).  Consequently,
for \(2\le N\), total reachability on the bipartite complete graph feeds the
case~(ii) target uniqueness and zero-magnetization wrappers directly, modulo
the standard balanced-sector bookkeeping.

The current general spin-\(S\), \(2\le N\), public region wrapper
(\texttt{AnisotropicHeisenbergSpinSTheorem24.lean}) is now only a dispatch
layer.  For \(-1<\lambda<1\), \(D\ge0\), it invokes the
\(\texttt{\_D\_nonneg\_general}\) endpoint.  For \(\lambda=1\), \(D\ge0\), it
splits \(D=0\) from \(D>0\): the former is the direct SU(2) endpoint, while the
latter is the ion-only \(\lambda=1\), \(D>0\) endpoint.  For
\(\lambda\ge1\), \(D\le0\), it invokes the case~(ii) total-reachability
endpoint above.  This gives
\[
  \dim_{\mathbb C}\ker\bigl(\hat H_{\lambda,D}
    - E_{\rm GS}(\lambda,D)\bigr) \le 1
\]
throughout the packaged parameter region, with \(E_{\rm GS}\) represented in
Lean by the Hermitian minimum eigenvalue of the anisotropic Hamiltonian.  The
second wrapper then applies the existing uniqueness-implies-zero-magnetization
theorem, giving \(\hat S^3_{\mathrm{tot}}\Phi=0\) for every non-zero target
ground state \(\Phi\).  This is the composed Theorem~2.4 endpoint currently
available for the general spin-\(S\), \(2\le N\), input surface.

For spin-\(1/2\), the remaining exact \(N=1\) case~(ii) packaging is handled
in the spin-half case~(ii) module.  The only obstruction to
specialising the general total-reachability wrapper to \(N=1\) is the
\(\lambda=1\), \(D<0\) ion-only branch, since ion-only totality requires
\(2\le N\).  For a strict case~(ii) target \(1<\lambda\), the linear path
\(\gamma(t)=((1-t)+t\lambda,tD)\) has first coordinate \(1\) only at \(t=0\),
so the non-corner path never enters that ion-only branch.  The strict parity
and \(D=0\) bond-only reachability inputs both require only \(1\le N\), while
the exact corner is still discharged by the full SU(2)-point uniqueness
callback.  The resulting strict case~(ii) target wrappers are:
\begin{Verbatim}[breaklines,breakanywhere,fontsize=\scriptsize]
spinHalf_anisotropicHeisenbergS_case_ii_target_finrank_le_one_of_MLM_casimir_ladder_t23_pf
spinHalf_anisotropicHeisenbergS_case_ii_target_zero_magnetization_of_MLM_casimir_ladder_t23_pf
\end{Verbatim}
Together with the spin-\(1/2\) \(D\ge0\) case~(i) endpoint and the scalar
\(\lambda=1\) boundary, they feed the exact \(N=1\) region wrappers:
\begin{Verbatim}[breaklines,breakanywhere,fontsize=\scriptsize]
spinHalf_anisotropicHeisenbergS_tasaki24_target_finrank_le_one_of_MLM_casimir_ladder_t23_pf
spinHalf_anisotropicHeisenbergS_tasaki24_target_zero_magnetization_of_MLM_casimir_ladder_t23_pf
\end{Verbatim}

The axis-swapped bond rewrites in ladder form (2.5.16),
\[
\hat S^1_x\hat S^1_y+\lambda\hat S^2_x\hat S^2_y+\hat S^3_x\hat S^3_y
=\tfrac{1+\lambda}{4}\bigl(\hat S^+_x\hat S^-_y+\hat S^-_x\hat S^+_y\bigr)
+\tfrac{1-\lambda}{4}\bigl(\hat S^+_x\hat S^+_y+\hat S^-_x\hat S^-_y\bigr)
+\hat S^3_x\hat S^3_y,
\]
(\texttt{spinSDotXXZSwap\_ladder\_form}): the \(\tfrac{1-\lambda}{4}(\hat
S^+\hat S^++\hat S^-\hat S^-)\) term shifts the magnetization by \(\pm2\),
coupling the \(\hat S^3_{\mathrm{tot}}\) sectors into the two even/odd parity
blocks.  A complementary symmetry is the spin reversal \(\Theta\)
(\texttt{manyBodyReversalS}), which anticommutes with \(\hat S^3_{\mathrm{tot}}\)
and commutes with \(\hat H\), hence carries the magnetization sector \(M\) onto
\(-M\) while preserving every \(\hat H\)-eigenspace; the resulting
\(E_M=E_{-M}\) per-sector degeneracy symmetry
(\texttt{manyBodyReversalS\_anisotropic\_sector\_eigenspace\_finrank\_eq})
underlies the ``first crossing forces \(\ge3\)-fold degeneracy'' step.

Dressing \(\hat H'\) by the Marshall sign matrix \(\Theta_A=\operatorname{diag}(\text{marshallSignS}\,A)\)
(a similarity, since the signs are \(\pm1\)) makes the off-diagonal entries
sign-definite.  For a bipartite antiferromagnetic coupling (\(J\ge0\) real, no
self-bonds, cross-sublattice support) and case~(i) (\(-1\le\lambda\le1\),
\(D\ge0\)) every off-diagonal entry of \(\hat H'_{\mathrm{dressed}}\) is
\(\le0\) (\texttt{dressedAxisSwappedAnisotropicHeisenbergS\_offdiag\_re\_nonpos}):
on a bipartite bond any \(\pm1\) move on the \(A\)-site gives Marshall sign
\(-1\), flipping the real non-negative transverse/parity bond entries, while the
single-ion \(\pm2\) moves are even (Marshall \(+1\)) with coefficient
\(-D/4\le0\).  Thus \(-\hat H'_{\mathrm{dressed}}\) is entrywise real with
non-negative off-diagonal part --- the Perron--Frobenius input for the
parity-block degeneracy bound.

Reference: H.\ Tasaki, \emph{Physics and Mathematics of Quantum
Many-Body Systems}, Springer 2020, §2.5 Theorem 2.4, p.~43--44 (eq.~2.5.16).

\subsection*{DONE / status: Tasaki Problem 2.5.a ($\gamma$-5,
                single-cluster star-graph ground-state energy)}
\label{subsec:tasaki-problem-2-5-a}

\textbf{Goal} (Tasaki §2.5 Problem 2.5.a, p.~38). For a single
$(z+1)$-vertex star graph with one central spin and $z$ leaves, all of
spin $S = N/2$, with antiferromagnetic Heisenberg coupling
$\hat H = \sum_{j=1}^{z} \hat{\boldsymbol S}_0 \cdot
\hat{\boldsymbol S}_j$, the ground-state energy is
$E_{\rm GS} = -S(1 + zS)$.

\textbf{Strategy.} The Hamiltonian admits a Casimir decomposition
\[
  2 \hat H = (\hat S_{\rm tot})^2 - (\hat S_0)^2 - (\hat S_R)^2,
\]
where $\hat S_R^{(\alpha)} := \sum_{j=1}^{z}
\mathrm{onSiteS}\,j\,\hat S_j^{(\alpha)}$ is the leaf-spin operator.
The three Casimirs $(\hat S_{\rm tot})^2, (\hat S_0)^2,
(\hat S_R)^2$ pairwise commute (proven via on-site / cross-site
commutativity), giving a joint eigenbasis. On a joint
$(\alpha, \beta)$-eigenstate of $((\hat S_{\rm tot})^2,
(\hat S_R)^2)$ the $\hat H$-eigenvalue is
\[
  E(\alpha, \beta) = \tfrac{1}{2}\bigl(\alpha
    - \tfrac{N(N+2)}{4} - \beta\bigr),
\]
where the central Casimir $(\hat S_0)^2$ acts as the scalar
$N(N+2)/4 = S(S+1)$. Minimising $E$ corresponds to maximising
$\beta$ and minimising $\alpha$ subject to angular-momentum
addition constraints; the ground state lies at $s_R = z N/2$ (max
leaf spin) with $s_{\rm tot} = (z-1)N/2$, giving
$E_{\rm GS} = -S(1 + zS)$.

\textbf{Formalisation status} ($\gamma$-5, in progress).  The
single-cluster modules are live again in the default Lean import graph:
\texttt{LatticeSystem.Quantum.SpinS.SingleClusterHamiltonianConcreteClusters}
imports the abstract Hamiltonian and energy companion, and
\texttt{LatticeSystem.lean} imports that tip module.

\begin{itemize}
  \item Single-cluster Hamiltonian
        \texttt{singleClusterHamiltonianS z N} on $\mathrm{Fin}(z+1)$
        with Hermiticity, $z=0$ edge case, all-aligned
        expectation $\langle \Phi_c | \hat H | \Phi_c\rangle =
        z(N/2 - c)^2$ ($\gamma$-5 steps 243--248).
  \item Leaf operators $\hat S_R^{(\alpha)}$, leaf Casimir
        $(\hat S_R)^2 := \sum_\alpha (\hat S_R^{(\alpha)})^2$,
        and the $\hat S_0 \cdot \hat S_R$ decomposition
        ($\gamma$-5 steps 249--252).
  \item Center-leaf commutativity and the Casimir decomposition
        $2 \hat H = (\hat S_{\rm tot})^2 - (\hat S_0)^2
        - (\hat S_R)^2$ ($\gamma$-5 steps 253--256).
  \item Single-site Casimir scalar action and the joint-Casimir
        eigenvalue formula
        \texttt{singleClusterHamiltonianS\_eigenvalue\_of\_joint\_casimir\_eigenvec}
        giving $\hat H \cdot v = ((\alpha - N(N+2)/4 - \beta)/2) \cdot v$
        ($\gamma$-5 steps 257--261).
  \item Eigenvector forms on extremal aligned states
        $|\Phi_\top\rangle, |\Phi_\bot\rangle$:
        $\hat H \cdot |\Phi_\top\rangle =
        z(N/2)^2 \cdot |\Phi_\top\rangle$ and
        $(\hat S_R)^2 \cdot |\Phi_\top\rangle =
        (zN/2)(zN/2+1) \cdot |\Phi_\top\rangle$
        (analogues for $|\Phi_\bot\rangle$); these confirm the
        extremal states are joint eigenstates at the maximum
        Casimir sector $(s_R, s_{\rm tot}) = (zN/2, (z+1)N/2)$
        ($\gamma$-5 steps 264--267).
  \item Named ground-state and maximum-Casimir-sector energies
        \texttt{singleClusterGSEnergyS z N := -(N/2)(zN/2+1)} and
        \texttt{singleClusterMaxEnergyS z N := z (N/2)$^2$},
        with conditional eigenvalue identities
        \texttt{singleClusterHamiltonianS\_mulVec\_eq\_gs\_energy\_smul}
        and the corresponding maximum-sector identity
        ($\gamma$-5 steps 268--271).
  \item Energy hygiene: real-part sign / order / reality / strict
        gap / strict negativity ($\gamma$-5 steps 272--283).
  \item Spin-specific numerical values for dimer / trimer /
        quartet at $S \in \{1/2, 1, 3/2, 2, 5/2\}$
        ($\gamma$-5 steps 277, 279, 282, 284, 289, 290, 291,
        297, 298).
  \item Dimer ($z=1$) concrete attainability: the leaf-Casimir
        collapses to the scalar $N(N+2)/4 \cdot 1$, so any
        $(\hat S_{\rm tot})^2$-eigenvector at $\alpha = 0$
        (singlet) is automatically an $\hat H$-eigenvector at
        \texttt{singleClusterGSEnergyS 1 N} = $-S(S+1)$, and at
        $\alpha = N(N+1)$ (top spin $s_{\rm tot} = 2S$) at
        \texttt{singleClusterMaxEnergyS 1 N} = $S^2$. This is the
        strongest concrete realisation of Tasaki Problem 2.5.a in
        the dimer case ($\gamma$-5 steps 285--288).
  \item Trimer ($z=2$) analysis: explicit leaf-Casimir
        decomposition $(\hat S_R)^2 = (N(N+2)/2) \cdot 1
        + 2 \cdot (\hat{\boldsymbol S}_1 \cdot \hat{\boldsymbol S}_2)$,
        eigenvalue formula $\hat H \cdot v
        = ((\alpha - 3N(N+2)/4 - 2\beta)/2) \cdot v$ at joint
        $(\hat S_{\rm tot}^2, \hat{\boldsymbol S}_1 \cdot
        \hat{\boldsymbol S}_2)$-eigenvector with eigenvalues
        $(\alpha, \beta)$, and three concrete sector eigenvalues:
        GS sector ($s_{\rm tot} = N/2$, leaf triplet) at
        \texttt{singleClusterGSEnergyS 2 N} $= -N(N+1)/2$, top
        sector ($s_{\rm tot} = 3N/2$, leaf triplet) at
        \texttt{singleClusterMaxEnergyS 2 N} $= N^2/2$, and
        leaf-singlet decoupling sector at $0$. Together these give
        the complete $S=1/2$ trimer spectrum
        $\{-1, 0, 1/2\}$ analytically ($\gamma$-5 steps 292--296).
  \item Quartet ($z=3$) and pentamer ($z=4$) concrete formulas:
        \texttt{leafSpinSSquared\_three} and
        \texttt{leafSpinSSquared\_four} decompose the leaf Casimir
        into the scalar diagonal contribution plus twice the sum of
        unordered leaf--leaf couplings.  The corresponding
        \texttt{singleClusterHamiltonianS\_eigenvalue\_quartet}
        and \texttt{...\_pentamer} theorems turn joint eigenvectors
        of $(\hat S_{\rm tot})^2$ and the leaf--leaf sum into
        $\hat H$-eigenvectors.  The GS-sector and top-sector
        specialisations recover
        \texttt{singleClusterGSEnergyS 3 N},
        \texttt{singleClusterMaxEnergyS 3 N},
        \texttt{singleClusterGSEnergyS 4 N}, and
        \texttt{singleClusterMaxEnergyS 4 N}
        ($\gamma$-5 steps 303--306 and 312--315).
  \item Leaf-singlet decoupling:
        \texttt{singleClusterHamiltonianS\_eigenvalue\_trimer\_leaf\_singlet},
        \texttt{...\_quartet\_leaf\_singlet},
        \texttt{...\_pentamer\_leaf\_singlet}, and the generic
        \texttt{singleClusterHamiltonianS\_eigenvalue\_leaf\_singlet}
        prove that if the leaves have total spin zero and the total
        Casimir is the central-spin value $N(N+2)/4$, then
        $\hat H \cdot v = 0$ ($\gamma$-5 steps 296, 321--323).
  \item Minimum-eigenvalue upper-bound bridge:
        \texttt{singleClusterHamiltonianS\_hermitianMinEigenvalue\_le\_gs\_of\_gs\_sector}
        shows that any non-zero vector in the predicted GS Casimir
        sector gives
        $\lambda_{\min}(\hat H) \le \operatorname{Re}
        \texttt{singleClusterGSEnergyS z N}$.  The companion
        \texttt{...\_of\_exists\_gs\_sector} packages the
        Clebsch--Gordan vector construction as the remaining
        hypothesis ($\gamma$-5 step 324).
  \item Minimum-eigenvalue lower/equality bridge:
        \texttt{singleClusterGSEnergyS\_re\_le\_hermitianMinEigenvalue\_of\_global\_eigenvalue\_lower}
        says that a global lower-bound callback for every non-zero
        real-energy eigenvector of $\hat H$ gives
        $\operatorname{Re}\texttt{singleClusterGSEnergyS z N}
        \le \lambda_{\min}(\hat H)$.  The equality wrappers
        \texttt{singleClusterHamiltonianS\_hermitianMinEigenvalue\_eq\_gs\_of\_gs\_sector\_and\_global\_lower}
        and \texttt{...\_of\_exists\_gs\_sector\_and\_global\_lower}
        combine this lower bound with the GS-sector upper-bound
        witness to identify the Hermitian minimum eigenvalue with
        the predicted GS energy ($\gamma$-5 step 325).
  \item Joint-Casimir lower callback bridge:
        \texttt{singleCluster\_global\_eigenvalue\_lower\_of\_joint\_casimir\_energy\_lower}
        turns the concrete joint-Casimir spectral-exhaustion target
        into the global lower-bound callback above.  Its hypothesis says
        that every non-zero real-energy $\hat H$-eigenvector admits
        $\hat S_{\rm tot}^2$ and $\hat S_R^2$ eigenvalues
        $(\alpha,\beta)$ whose Casimir energy
        $(\alpha - N(N+2)/4 - \beta)/2$ lies above
        $\operatorname{Re}\texttt{singleClusterGSEnergyS z N}$.
        Consumer forms
        \texttt{singleClusterGSEnergyS\_re\_le\_hermitianMinEigenvalue\_of\_joint\_casimir\_energy\_lower}
        and the corresponding witness / existential equality wrappers
        reduce the remaining lower-bound work to proving that
        joint-Casimir callback and its sector arithmetic bound
        ($\gamma$-5 step 326).
  \item Joint-Casimir sector arithmetic lower bound:
        \texttt{singleClusterGSEnergyS\_re\_le\_casimir\_energy\_of\_joint\_bounds}
        proves the pure arithmetic implication
        \[
          \alpha_{\rm pred} \le \operatorname{Re}\alpha,\quad
          \operatorname{Re}\beta \le \beta_{\max}
          \quad\Longrightarrow\quad
          \operatorname{Re}\texttt{singleClusterGSEnergyS z N}
          \le
          \operatorname{Re}\frac{\alpha - N(N+2)/4 - \beta}{2},
        \]
        where
        $\alpha_{\rm pred} = ((z-1)N/2)(((z-1)N/2)+1)$ and
        $\beta_{\max} = (zN/2)(zN/2+1)$.  The callback forms
        \texttt{singleCluster\_global\_eigenvalue\_lower\_of\_joint\_casimir\_bounds}
        and
        \texttt{singleClusterGSEnergyS\_re\_le\_hermitianMinEigenvalue\_of\_joint\_casimir\_bounds}
        turn these sector bounds for every non-zero real-energy
        $\hat H$-eigenvector into the global lower callback and
        Hermitian-min lower bound ($\gamma$-5 step 327).
  \item Conditional equality from joint-Casimir sector bounds:
        \texttt{singleClusterHamiltonianS\_minEigenvalue\_eq\_gs\_of\_gs\_sector\_and\_joint\_casimir\_bounds}
        combines a concrete predicted GS-sector witness with the
        joint-Casimir sector-bounds callback to obtain
        $\lambda_{\min}(\hat H)=
        \operatorname{Re}\texttt{singleClusterGSEnergyS z N}$.
        The existential companion
        \texttt{...\_minEigenvalue\_eq\_gs\_of\_exists\_gs\_sector\_and\_joint\_casimir\_bounds}
        packages the Clebsch--Gordan witness as the remaining
        existence hypothesis ($\gamma$-5 step 328).
  \item Leaf-Casimir spectral maximum bound:
        \texttt{leafSpinSSquared\_eq\_sublatticeSpinSquaredS\_leaf}
        identifies the leaf Casimir with the sublattice Casimir for
        $\{x:\operatorname{Fin}(z+1)\mid x\ne 0\}$, and
        \texttt{leafSpinSSquared\_eigenvalue\_re\_le\_max}
        applies the general sublattice bound to get
        $\operatorname{Re}\beta \le (zN/2)(zN/2+1)$ for every
        non-zero $\hat S_R^2$ eigenvector.  The callback forms
        \texttt{singleCluster\_global\_eigenvalue\_lower\_of\_joint\_casimir\_total\_lower}
        and
        \texttt{singleClusterGSEnergyS\_re\_le\_hermitianMinEigenvalue\_of\_joint\_casimir\_total\_lower}
        remove the leaf upper-bound hypothesis from the lower-bound
        bridge, and the companion wrappers
        \texttt{singleClusterHamiltonianS\_minEigenvalue\_eq\_gs\_of\_gs\_sector\_and\_joint\_casimir\_total\_lower}
        and
        \texttt{singleClusterHamiltonianS\_minEigenvalue\_eq\_gs\_of\_exists\_gs\_sector\_and\_joint\_casimir\_total\_lower}
        package the corresponding conditional equality endpoints
        ($\gamma$-5 step 329).
  \item Coupled leaf/center sector energy bridge:
        \texttt{sublatticeSpinSquaredS\_compl\_leaf\_eq\_spinSDot\_zero\_zero}
        identifies the complement of the leaf sublattice with the
        singleton center Casimir, hence with the scalar single-site value.
        For $1\le z$,
        \texttt{singleClusterGSEnergyS\_re\_le\_casimir\_energy\_of\_coupled\_leaf\_sector}
        applies the existing Clebsch--Gordan coupled lower theorem to the
        leaf/center partition and divides the toy-sector energy by two.
        The callback forms
        \texttt{singleCluster\_global\_eigenvalue\_lower\_of\_coupled\_leaf\_sector}
        and
        \texttt{singleClusterGSEnergyS\_re\_le\_hermitianMinEigenvalue\_of\_coupled\_leaf\_sector}
        reduce the remaining lower-bound work to proving that real-energy
        eigenvectors decompose into magnetization and joint total/leaf
        Casimir sectors ($\gamma$-5 step 330).
  \item Single-cluster magnetization projection bridge:
        \texttt{mulVec\_magProjFn\_eq\_of\_mulVec\_basisVecS\_mem\_magSubspaceS}
        abstracts the existing matrix-entry argument: any operator that sends
        each computational basis vector to its own magnetization subspace
        commutes with the pointwise magnetization projection \texttt{magProjFn}.
        \texttt{singleClusterHamiltonianS\_commute\_totalSpinSOp3}
        follows by summing the two-site commutators
        $[\hat{\boldsymbol S}_0\cdot\hat{\boldsymbol S}_j,
        \hat S^3_{\rm tot}]=0$ over all leaves.  Therefore
        \texttt{singleClusterHamiltonianS\_mulVec\_mem\_magSubspaceS\_of\_mem}
        shows that the Hamiltonian preserves each magnetization subspace.
        The basis-vector form
        \texttt{singleClusterHamiltonianS\_mulVec\_basisVecS\_mem\_magSubspaceS}
        supplies the hypothesis of the generic projection helper.
        The pointwise magnetization projection theorem
        \texttt{singleClusterHamiltonianS\_mulVec\_magProjFn\_eq}
        then gives
        $\hat H P_M v = P_M \hat H v$.  Consequently
        \texttt{singleClusterHamiltonianS\_eigenvec\_exists\_weight\_component}
        extracts a non-zero magnetization component from any non-zero
        single-cluster eigenvector, preserving the same energy
        ($\gamma$-5 step 331).
  \item Invariant-submodule eigenvector bridge:
        \texttt{exists\_eigenvector\_in\_invariant\_submodule}
        is the finite-dimensional linear-algebra bridge for the next
        sector-decomposition layer.  If a non-zero submodule \(p\) is invariant
        under an endomorphism \(f\), then the restriction \(f|_p\) has an
        eigenvector by Mathlib's algebraically closed finite-dimensional
        eigenvalue theorem.  Forgetting the subtype gives \(0\ne v\in p\)
        with \(f v=\mu v\).  This is intended for the next joint total/leaf
        Casimir eigensector extraction after the single-cluster
        magnetization-component step ($\gamma$-5 step 332).
  \item Joint invariant-submodule eigenvector bridge:
        \texttt{exists\_joint\_eigenvector\_in\_invariant\_submodule}
        iterates the preceding bridge for two commuting endomorphisms.  If
        \(p\) is non-zero and invariant under \(f\) and \(g\), and \(fg=gf\),
        first choose a non-zero \(f\)-eigenvector in \(p\).  The intersection
        of \(p\) with that \(f\)-eigenspace is non-zero and is preserved by
        \(g\); commutativity gives \(f(gx)=g(fx)=\lambda gx\).  Applying the
        one-operator bridge again gives \(0\ne v\in p\) with
        \(f v=\lambda v\) and \(g v=\mu v\).  This packages the linear-algebra
        step needed to extract simultaneous total/leaf Casimir data inside the
        single-cluster magnetization/Hamiltonian component ($\gamma$-5 step
        333).
  \item Single-cluster Casimir invariance bridge:
        \texttt{singleClusterHamiltonianS\_commute\_totalSpinSSquared}
        and
        \texttt{singleClusterHamiltonianS\_commute\_leafSpinSSquared}
        turn the Casimir decomposition
        \[
          (2:\mathbb C)\,\hat H
          = \hat S_{\rm tot}^2 - \hat S_0^2 - \hat S_R^2
        \]
        into concrete commutation theorems.  The central-site term
        \(\hat S_0^2\) is scalar, and the total/leaf commutation follows from
        the general sublattice-Casimir API after identifying
        \(\hat S_R^2\) with the sublattice Casimir of the predicate \(x\ne 0\).
        Hence the submodule
        \[
          \operatorname{Eig}(\hat H,\mu) \cap \mathcal H_M
        \]
        is invariant under both \(\hat S_{\rm tot}^2\) and \(\hat S_R^2\):
        the eigenspace part is preserved by commutation with \(\hat H\), and
        the magnetization part is preserved by the existing total and
        sublattice Casimir magnetization-preservation theorems.  The theorem
        \texttt{exists\_joint\_casimir\_eigenvector\_in\_singleClusterHamiltonianMagEigenspaceSubmodule}
        then applies the joint invariant-submodule bridge to any non-zero
        such submodule, producing simultaneous total/leaf Casimir eigenvector
        data under the explicit algebraically closed field hypothesis
        \texttt{[IsAlgClosed \(\mathbb C\)]} ($\gamma$-5 step 334).
  \item Same-energy joint-Casimir lower bridge:
        \texttt{singleCluster\_global\_eigenvalue\_lower\_of\_exists\_joint\_casimir\_energy\_lower}
        weakens the earlier lower-bound callback.  The original non-zero
        real-energy eigenvector \(v\) of \(\hat H\) no longer needs to be a
        joint total/leaf Casimir eigenvector.  It is enough to find another
        non-zero vector \(w\) with the same Hamiltonian eigenvalue
        \(\mu\), together with
        \[
          \hat S_{\rm tot}^2 w=\alpha w,\qquad
          \hat S_R^2 w=\beta w
        \]
        and the Casimir-energy lower bound
        \[
          \operatorname{Re} E_{\rm GS}
          \le
          \operatorname{Re}\left(
            \frac{\alpha-N(N+2)/4-\beta}{2}\right).
        \]
        The existing joint-Casimir Hamiltonian eigenvalue theorem applied to
        \(w\) gives
        \[
          \hat H w =
            \frac{\alpha-N(N+2)/4-\beta}{2}\,w.
        \]
        Since \(w\ne0\), scalar cancellation identifies this Casimir energy
        with the original real eigenvalue \(\mu\), and taking real parts gives
        the global lower callback.  The companion declarations
        \texttt{singleClusterGSEnergyS\_re\_le\_hermitianMinEigenvalue\_of\_exists\_joint\_casimir\_energy\_lower},
        \texttt{singleClusterHamiltonianS\_hermitianMinEigenvalue\_eq\_gs\_of\_gs\_sector\_and\_exists\_joint\_lower},
        and
        \texttt{singleClusterHamiltonianS\_minEigenvalue\_eq\_gs\_of\_exists\_gs\_sector\_and\_exists\_joint\_lower}
        package the Hermitian-min lower bound and conditional equality
        endpoints ($\gamma$-5 step 335).
  \item Single-cluster joint-Casimir witness bridge:
        \texttt{singleCluster\_exists\_joint\_casimir\_energy\_lower\_of\_mag\_component\_joint\_lower}
        composes the magnetization projection bridge with the invariant
        submodule Casimir bridge.  Given a non-zero real-energy Hamiltonian
        eigenvector \(v\),
        \[
          \hat H v=\mu v,\qquad \mu\in\mathbb R ,
        \]
        the theorem
        \texttt{singleClusterHamiltonianS\_eigenvec\_exists\_weight\_component}
        provides a magnetization value \(M\) and a non-zero component
        \(v_M=\texttt{magProjFn}\,M\,v\) with the same Hamiltonian eigenvalue:
        \[
          v_M\ne0,\qquad \hat H v_M=\mu v_M,\qquad v_M\in\mathcal H_M .
        \]
        Thus \(v_M\) lies in
        \[
          \operatorname{Eig}(\hat H,\mu)\cap\mathcal H_M,
        \]
        proving that the corresponding
        \texttt{singleClusterHamiltonianMagEigenspaceSubmodule} is non-zero.
        Under the explicit hypothesis
        \texttt{[IsAlgClosed \(\mathbb C\)]},
        \texttt{exists\_joint\_casimir\_eigenvector\_in\_singleClusterHamiltonianMagEigenspaceSubmodule}
        then returns a non-zero vector \(w\) in the same submodule such that
        \[
          \hat S_{\rm tot}^2w=\alpha w,\qquad
          \hat S_R^2w=\beta w .
        \]
        The submodule membership also contains the Hamiltonian eigenspace
        membership, so Lean recovers \(\hat H w=\mu w\) by rewriting
        \texttt{Module.End.eigenspace} membership through
        \texttt{Matrix.mulVecLin\_apply}.  Therefore any remaining lower-bound
        callback for such extracted joint vectors supplies the same-energy
        witness callback introduced in the previous item.  The companion
        declarations
        \texttt{singleClusterGSEnergyS\_re\_le\_hermitianMinEigenvalue\_of\_mag\_component\_joint\_lower},
        \texttt{singleClusterHamiltonianS\_minEigenvalue\_eq\_gs\_of\_gs\_sector\_and\_mag\_joint\_lower},
        and
        \texttt{singleClusterHamiltonianS\_minEigenvalue\_eq\_gs\_of\_exists\_gs\_sector\_and\_mag\_joint\_lower}
        pass this witness callback into the PR~\#4044 lower/equality wrappers
        ($\gamma$-5 step 336).
  \item Extracted coupled leaf-sector lower bridge:
        \texttt{singleCluster\_mag\_component\_joint\_lower\_of\_coupled\_leaf\_sector}
        closes the lower-bound callback left by the previous item.  Let
        \(w\ne0\) lie in an extracted Hamiltonian/magnetization submodule:
        \[
          w\in \operatorname{Eig}(\hat H,\mu)\cap\mathcal H_M,\qquad
          \hat S_{\rm tot}^2w=\alpha w,\qquad
          \hat S_R^2w=\beta w .
        \]
        Since \(w\ne0\) in \(\mathcal H_M\), the magnetization value is
        attained by some configuration \(\sigma\):
        \[
          M=\operatorname{magEigenvalueS}(\sigma)
           =m_{\max}-\operatorname{magSumS}(\sigma),
          \qquad
          m_{\max}=|\operatorname{Fin}(z+1)|N/2 .
        \]
        The coupled leaf-sector theorem uses the opposite indexing convention
        \(k-m_{\max}\).  Lean sets
        \[
          k=|\operatorname{Fin}(z+1)|N-\operatorname{magSumS}(\sigma)
        \]
        and rewrites
        \[
          M=k-m_{\max}.
        \]
        Therefore the same vector \(w\) belongs to the magnetization sector
        required by
        \texttt{singleClusterGSEnergyS\_re\_le\_casimir\_energy\_of\_coupled\_leaf\_sector}.
        For \(1\le z\), that theorem gives
        \[
          \operatorname{Re}E_{\rm GS}
          \le
          \operatorname{Re}\left(
            \frac{\alpha-N(N+2)/4-\beta}{2}\right).
        \]
        Feeding this callback into the PR~\#4045 witness bridge yields
        \texttt{singleClusterGSEnergyS\_re\_le\_hermitianMinEigenvalue\_of\_extracted\_coupled\_leaf\_sector}.
        The companion declarations
        \texttt{singleClusterHamiltonianS\_minEigenvalue\_eq\_gs\_of\_gs\_sector\_and\_extracted\_coupled\_leaf\_sector}
        and
        \texttt{singleClusterHamiltonianS\_minEigenvalue\_eq\_gs\_of\_exists\_gs\_sector\_and\_extracted\_coupled}
        package the corresponding conditional equality endpoints under the
        explicit hypotheses \(1\le z\) and
        \texttt{[IsAlgClosed \(\mathbb C\)]}
        ($\gamma$-5 step 337).
  \item Single-cluster predicted GS joint-Casimir witness:
        \texttt{singleCluster\_exists\_gs\_joint\_casimir\_witness}
        specializes the bipartite predicted joint-Casimir eigenvector
        \texttt{exists\_jointPredictedCasimir\_eigenvector} to the leaf/center
        partition of \(\operatorname{Fin}(z+1)\).  The leaf predicate is
        \(A(x)=\operatorname{decide}(x\ne0)\).  For \(1\le z\),
        \[
          |\neg A|=1\le z=|A|,
        \]
        so the canonical orientation hypothesis of the bipartite construction
        is available.  The theorem then returns a non-zero vector \(v\) with
        \[
          \hat S_{\rm tot}^2v
            =
            \left(\frac{(z-1)N}{2}\right)
            \left(\frac{(z-1)N}{2}+1\right)v,
          \qquad
          \hat S_R^2v
            =
            \left(\frac{zN}{2}\right)
            \left(\frac{zN}{2}+1\right)v.
        \]
        In Lean, the total-Casimir scalar is simplified using
        \texttt{tasaki23PredictedCasimirValue\_eq\_sub},
        \texttt{singleCluster\_leaf\_bool\_card}, and
        \texttt{singleCluster\_compl\_leaf\_bool\_card}; the leaf-Casimir
        equation is obtained by rewriting the sublattice equation with
        \texttt{leafSpinSSquared\_eq\_sublatticeSpinSquaredS\_leaf}.
        Feeding this existential package into
        \texttt{singleClusterHamiltonianS\_minEigenvalue\_eq\_gs\_of}
        \texttt{\_exists\_gs\_sector\_and\_extracted\_coupled}
        gives
        \texttt{singleClusterHamiltonianS\_minEigenvalue\_eq\_gs\_of}
        \texttt{\_predicted\_joint\_witness}:
        for \(1\le z\),
        \[
          \lambda_{\min}(\hat H)
          =
          \operatorname{Re}E_{\rm GS}
        \]
        under the explicit hypothesis
        \texttt{[IsAlgClosed \(\mathbb C\)]}
        ($\gamma$-5 step 338).
  \item Problem 2.5.b finite-sum lower-bound bridge.  The declaration
        \texttt{sum\_lower\_bounds\_le}
        \texttt{\_hermitianMinEigenvalue\_sum}
        packages the Lemma~A.5 argument in minimum-eigenvalue form.  For a
        finite family of Hermitian matrices \(H_i\), if each local minimum
        eigenvalue is bounded below by \(\varepsilon_i\), then
        \[
          \sum_i \varepsilon_i
          \le
          \lambda_{\min}\!\left(\sum_i H_i\right).
        \]
        The proof chooses a unit eigenvector attaining the minimum eigenvalue
        of \(\sum_i H_i\).  Rayleigh--Ritz gives
        \(\lambda_{\min}(H_i)\le\langle v,H_i v\rangle\) for each summand,
        and \texttt{rayleighOnVec\_sum\_matrix} identifies the sum of the
        Rayleigh quotients with the Rayleigh quotient of the summed matrix.

        The two \texttt{tasaki25b\_local\_cluster} wrappers specialize this
        abstract sum bound to local cluster Hamiltonians whose
        minimum eigenvalues are the Problem~2.5.a values.  In closed form this
        is Tasaki's lower bound
        \[
          E_{\rm GS}
          \ge
          -\sum_{x\in A} S\bigl(1+|N(x)|S\bigr),
          \qquad S=N/2,
        \]
        once the global bipartite Hamiltonian is decomposed as
        \(\hat H=\sum_{x\in A}\hat h_x\) (Tasaki, Springer 2020,
        Problem~2.5.b, p.~38; solution p.~497; Lemma~A.5, p.~468).
  \item Problem 2.5.b graph-local decomposition.  The declaration
        \texttt{graphLocalClusterHamiltonianS} defines, on the full
        Hilbert space over the graph vertex set,
        \[
          h_x=\sum_{y\in N_G(x)}
            \hat{\mathbf S}_x\cdot\hat{\mathbf S}_y .
        \]
        The theorem
        \texttt{heisenbergHamiltonianOnGraphS\_one\_eq\_sum}
        \texttt{\_graphLocalClusterHamiltonianS} records the repository's
        ordered-pair convention:
        \[
          \hat H_G(1)=\sum_x h_x .
        \]
        The bipartite refinement
        \texttt{heisenbergHamiltonianOnGraphS\_half\_eq\_sum}
        \texttt{\_filter\_graphLocalClusterHamiltonianS} proves that, if every
        edge crosses the chosen side \(A\),
        \[
          \hat H_G(1/2)=\sum_{x\in A}h_x .
        \]
        The proof splits the all-vertex sum into \(A\) and its complement,
        swaps each oriented adjacent pair, and uses
        \(\hat{\mathbf S}_x\cdot\hat{\mathbf S}_y
        =\hat{\mathbf S}_y\cdot\hat{\mathbf S}_x\).  This fixes the coefficient
        convention for Tasaki's local-cluster decomposition (Tasaki, Springer
        2020, Problem~2.5.b, p.~38; solution p.~497).
  \item Problem 2.5.b single-cluster transport bridge.  The declaration
        \texttt{siteConfigEquiv} lifts a site equivalence \(e:\alpha\simeq\beta\)
        to the corresponding equivalence of spin-configuration spaces, with
        simp rules \texttt{siteConfigEquiv\_apply} and
        \texttt{siteConfigEquiv\_symm\_apply}.  The theorems
        \texttt{reindex\_onSiteS\_siteEquiv} and
        \texttt{reindex\_spinSDot\_siteEquiv} show that matrix reindexing along
        this configuration equivalence simply transports the site labels:
        \[
          \operatorname{reindex}_e
          (\hat{\mathbf S}_a\cdot\hat{\mathbf S}_b)
          =
          \hat{\mathbf S}_{e(a)}\cdot\hat{\mathbf S}_{e(b)} .
        \]
        Consequently the transported single-cluster Hamiltonian
        \texttt{transportedSingleClusterHamiltonianS} is exactly the sum of
        transported center-leaf dot products.  The canonical equivalence
        \(\mathrm{Fin}(|s|+1)\simeq\mathrm{Option}(s)\), sending \(0\) to
        \(\mathrm{none}\) and successors to \(\mathrm{some}\) leaves via
        \texttt{Finset.equivFin}, uses the finite-sum reindexing helper
        \texttt{sum\_fin\_erase\_zero\_eq\_sum\_succ}, and yields the named
        option-star identity
        \texttt{transportedSingleClusterHamiltonianS\_option\_eq}.
  \item Problem 2.5.b graph-local star block bridge.  After fixing an outside
        configuration \(\eta\), the new embedding sends a configuration on
        \(\mathrm{Option}(N_G(x))\) into the full graph configuration space.
        Its pointwise rules identify the center, neighbor, and outside
        coordinates.  Forward and converse agreement-transfer lemmas reduce
        the two-site support condition for
        \(\hat{\mathbf S}_x\cdot\hat{\mathbf S}_y\) to the option-star support
        condition.  The centered-edge matrix-element theorem then identifies
        each graph edge with the option-star edge.  Summing over
        \(y\in N_G(x)\) gives the off-block zero theorem, saying different
        outside restrictions do not couple, and the block-equality theorem,
        identifying each fixed-outside block with the option-star Hamiltonian.
        The exact Lean names are listed in \texttt{docs/index.md}.
  \item Problem 2.5.b graph-local product lower-bound bridge.  The full
        configuration space around \(x\) is reindexed as
        \[
          (\mathrm{Option}(N_G(x))\to\{0,\ldots,N\})
          \times
          (O_x\to\{0,\ldots,N\}),
        \]
        where \(O_x\) is the subtype of sites outside the center and its
        graph neighbours.  In these product coordinates, the graph-local star
        is block diagonal in the outside coordinate, and each diagonal block is
        the option-star Hamiltonian.  The product-coordinate multiplication
        formula gives a Rayleigh-numerator decomposition as a sum over outside
        blocks, and the corresponding norm decomposition gives the abstract
        lower-bound lift: any common Rayleigh lower bound for the option-star
        blocks is also a Hermitian minimum-eigenvalue lower bound for the
        product-coordinate graph-local star.  Reindex-invariance lemmas for
        matrix-vector multiplication, norms, and Rayleigh numerators transfer
        the same lower bound back to the original same-Hilbert-space
        graph-local star.  The exact Lean names are listed in
        \texttt{docs/index.md}.
  \item Problem 2.5.b option-star and graph-local sum wrapper.  The canonical
        option-star Hamiltonian is Hermitian because it is the transported
        single-cluster Hamiltonian along
        \(\mathrm{Fin}(|s|+1)\simeq\mathrm{Option}(s)\).  The forward
        Rayleigh reindexing lemmas identify the Rayleigh quotient of the
        option-star matrix with the Rayleigh quotient of the abstract
        single-cluster matrix after precomposition by the configuration
        equivalence.  Therefore the Problem~2.5.a equality
        \[
          \lambda_{\min}(\hat h_{\rm cl}(|s|))
          =
          \operatorname{Re}E_{\rm GS}(|s|)
        \]
        for \(1\le |s|\), under \texttt{[IsAlgClosed \(\mathbb C\)]}, gives
        the option-star Rayleigh lower bound with
        \(\operatorname{Re}E_{\rm GS}(|s|)\).

        Applying the graph-local product lower-bound bridge yields
        \[
          \operatorname{Re}E_{\rm GS}(|N_G(x)|)
          \le
          \lambda_{\min}(h_x)
        \]
        whenever the center has at least one graph neighbour.  Summing these
        local lower bounds and using the one-sided bipartite decomposition
        \(\hat H_G(1/2)=\sum_{x\in A}h_x\) gives the formal lower bound
        \[
          \sum_{x\in A}\operatorname{Re}E_{\rm GS}(|N_G(x)|)
          \le
          \lambda_{\min}(\hat H_G(1/2)).
        \]
        The closed-form wrappers substitute the Problem~2.5.a formula
        \[
          \operatorname{Re}E_{\rm GS}(z)
          =
          -\frac{N}{2}\left(\frac{zN}{2}+1\right)
        \]
        and use the graph identity \(|N_G(x)|=\deg_G(x)\).  Thus the
        one-sided graph-Hamiltonian theorem is stated in the expected
        degree form
        \[
          \sum_{x\in A}
          -\frac{N}{2}\left(\frac{\deg_G(x)N}{2}+1\right)
          \le
          \lambda_{\min}(\hat H_G(1/2)),
        \]
        still under the explicit positive-degree assumption on the selected
        side.  Isolated centers are intentionally excluded because their
        local star Hamiltonian is zero, whereas the closed formula would
        contribute a negative number.
        The exact Lean names are listed in \texttt{docs/index.md}.  This is
        the graph-local form of Tasaki's Problem~2.5.b lower-bound argument
        (Tasaki, Springer 2020, Problem~2.5.b, p.~38; solution p.~497).
\end{itemize}

\textbf{Strategic gaps.} The single-cluster Problem 2.5.a equality is now
available for \(1\le z\) under the explicit
\texttt{[IsAlgClosed \(\mathbb C\)]} hypothesis.  The remaining optional
cleanup is to discharge that hypothesis by replacing the algebraically closed
eigenvector-extraction bridge with a constructive finite-dimensional spectral
decomposition if that later becomes useful.  For Problem~2.5.b, the graph-local
finite-sum lower bound is now available for the half-coupling graph Hamiltonian
under the same algebraically closed-field hypothesis and a positive local-degree
assumption on the chosen bipartition side.

\subsection*{DONE: Spin-$S$ saturated ferromagnetic state (Tasaki §2.4 generalised)}
\label{subsec:spin-s-all-aligned}

\textbf{Statement} (spin-$S$ generalisation of Tasaki §2.4 P1i,
saturated ferromagnet). For a multi-site spin-$S$ system on a finite
vertex set $V$ with $S = N/2$, the all-aligned (constant-spin) basis
state $|\sigma_{\!c}\rangle = \otimes_x |c\rangle$ with
$\sigma\,x = c$ for all $x : V$ is:
\begin{enumerate}
\item A $\hat S^z_{\mathrm{tot}}$-eigenvector with eigenvalue
  $|V|\cdot N/2 - |V|\cdot c$ (any $c$).
\item For $c = 0$ (\emph{highest weight}, ``all up''): the unique
  configuration with $\mathrm{magSum} = 0$, hence by
  $[H,\hat S^z_{\rm tot}] = 0$ a Heisenberg eigenvector for ANY
  coupling $J$, with eigenvalue
  $\sum_x J(x,x)\cdot \tfrac{N(N+2)}{4} +
   \sum_{x\neq y} J(x,y)\cdot \tfrac{N^2}{4}$. Annihilated by
  $\hat S^+_{\mathrm{tot}}$.
  $(\hat S_{\mathrm{tot}})^2$-eigenvector with the highest-weight
  Casimir eigenvalue
  $(|V|\cdot N/2)\cdot(|V|\cdot N/2 + 1)$.
\item For $c = N$ (\emph{lowest weight}, ``all down''): symmetric —
  same eigenvalues (since $(N/2 - N)^2 = (N/2)^2$), annihilated by
  $\hat S^-_{\mathrm{tot}}$.
\end{enumerate}

These together identify the all-up and all-down basis states as the
extremal-weight vectors of the $J_{\mathrm{tot}} = |V|\cdot S
= |V|\cdot N/2$ irreducible SU(2) representation in the multi-site
Hilbert space.

\textbf{Status}: Formalized in
\texttt{LatticeSystem/Quantum/SpinS/AllAlignedState.lean} as PRs
\#875--\#879 (Issue~\#412). Theorems include
\texttt{allAlignedConfigS}, \texttt{allAlignedStateS},
\texttt{magSumS\_pos\_of\_ne\_allAlignedConfigS\_zero}
(uniqueness in $M=0$ sector),
\texttt{heisenbergHamiltonianS\_mulVec\_allAlignedStateS\_zero\_eigenvalue}
(Heisenberg eigenvalue),
\texttt{totalSpinSOpPlus\_mulVec\_allAlignedStateS\_zero}
(highest-weight annihilation),
\texttt{totalSpinSSquared\_mulVec\_allAlignedStateS\_zero\_eigenvalue}
(Casimir eigenvalue), and the c=N parallels.

\textbf{Iterated ladder (PRs~\#881--\#883)}: subsequently extended
with ladder commutation and iterated-ladder eigenvalue preservation:
\begin{itemize}
\item Heisenberg + Casimir + axis-total Commute conversions (\#881,
  \#882). The Heisenberg Hamiltonian and the Casimir
  $(\hat S_{\rm tot})^2$ both commute with each
  $\hat S^{(\alpha)}_{\rm tot}$, hence with $\hat S^\pm_{\rm tot}$,
  hence with $(\hat S^\pm_{\rm tot})^k$ by induction.
\item Iterated-ladder eigenvalue preservation
  \texttt{heisenbergHamiltonianS\_mulVec\_totalSpinSOpMinus\_pow\_allAlignedStateS\_zero}
  and the Casimir analogue: for any $k$, the ladder iterate
  $(\hat S^-_{\rm tot})^k\,|\sigma_\top\rangle$ is a
  Heisenberg + Casimir simultaneous eigenstate at the same
  saturated-ferromagnet eigenvalue (and dually for
  $(\hat S^+_{\rm tot})^k\,|\sigma_\bot\rangle$).
\item Multi-site Cartan relations
  \texttt{totalSpinSOp3\_commutator\_totalSpinSOpMinus} and
  \texttt{...\_OpPlus}: $[\hat S^z_{\rm tot}, \hat S^\pm_{\rm tot}] =
  \pm \hat S^\pm_{\rm tot}$ (\#883), the operator-algebraic input to
  the magnetic-quantum-number labelling along the ladder.
\end{itemize}

\textbf{Magnetic-quantum-number labelling along the ladder
(PRs~\#886--\#887)}: with the multi-site Cartan in hand, the
unnormalised iterates $(\hat S^\mp_{\rm tot})^k\,|\sigma_\top\rangle$
and $(\hat S^\pm_{\rm tot})^k\,|\sigma_\bot\rangle$ are simultaneous
$\hat S^z_{\rm tot}$-eigenvectors with eigenvalue $m_{\max}\mp k$ resp.
$-m_{\max}\pm k$, where $m_{\max} = |\Lambda|\cdot N/2$.
\begin{itemize}
\item Single-step shifts on the all-aligned states
  \texttt{totalSpinSOp3\_mulVec\_totalSpinSOpMinus\_mulVec\_allAlignedStateS\_zero}
  and \texttt{...\_OpPlus\_mulVec\_allAlignedStateS\_last}~(\#886):
  the once-lowered (resp. once-raised) all-up (resp. all-down)
  state is an $\hat S^z_{\rm tot}$ eigenvector at $m_{\max}-1$
  (resp. $-m_{\max}+1$).
\item Generic single-step shift on \emph{any} $\hat S^z_{\rm tot}$
  eigenvector
  \texttt{totalSpinSOp3\_mulVec\_totalSpinSOpMinus\_mulVec\_eigenvec}
  and \texttt{...\_OpPlus\_mulVec\_eigenvec}~(\#887): if
  $\hat S^z_{\rm tot}\psi = \lambda\psi$ then
  $\hat S^z_{\rm tot}(\hat S^\mp_{\rm tot}\psi) = (\lambda\mp 1)
   (\hat S^\mp_{\rm tot}\psi)$. Proof: rearrange
  $\hat S^z_{\rm tot}\hat S^\mp_{\rm tot} = \hat S^\mp_{\rm tot}
   \hat S^z_{\rm tot}\mp\hat S^\mp_{\rm tot}$ via
  \texttt{sub\_eq\_iff\_eq\_add} on the multi-site Cartan, lift to
  \texttt{mulVec}, substitute $\psi$ as an eigenvector.
\item Iterated form
  \texttt{totalSpinSOp3\_mulVec\_totalSpinSOpMinus\_pow\_allAlignedStateS\_zero}
  and \texttt{...\_OpPlus\_pow\_allAlignedStateS\_last}~(\#887):
  $\hat S^z_{\rm tot}\,(\hat S^\mp_{\rm tot})^k|\sigma_{\top/\bot}\rangle
   = (\pm m_{\max}\mp k)\,(\hat S^\mp_{\rm tot})^k
     |\sigma_{\top/\bot}\rangle$. Proof: induction on $k$ at the
  \emph{eigenvector} level, applying the generic single-step shift
  to the eigenvector supplied by the IH. This bypasses the
  technically delicate operator-level iterated Cartan
  $\hat S^z_{\rm tot}\,(\hat S^\mp_{\rm tot})^k =
   (\hat S^\mp_{\rm tot})^k\,\hat S^z_{\rm tot} \mp
   k\,(\hat S^\mp_{\rm tot})^k$.
\end{itemize}

\textbf{Extremal magnetization sectors are 1-dimensional
(PR~\#885)}: the combinatorial counterpart of the uniqueness of the
highest- and lowest-weight states of the $J_{\rm tot}=|\Lambda|\cdot S$
irreducible SU(2) representation. In the file
\texttt{Quantum/SpinS/MagConfigExtremalCardinality.lean} the
$M=0$ (resp. $M=|V|\cdot N$) magnetization sector subtype
\texttt{magConfigS V N M} is shown to contain exactly the
all-up (resp. all-down) configuration, hence
\texttt{Fintype.card (magConfigS V N 0) = 1} via
\texttt{Fintype.card\_eq\_one\_iff} and the
\texttt{magSumS\_pos\_of\_ne\_allAlignedConfigS\_zero}~/
\texttt{magSumS\_lt\_card\_mul\_of\_ne\_allAlignedConfigS\_last}
strict bounds; \texttt{Nonempty} instances follow.

\textbf{Magnetization subspaces form an internal direct sum
(PR~\#889)}: the spin-`S` analog of the spin-1/2
\texttt{magnetizationSubspace\_isInternal} (in
\texttt{Quantum/MagnetizationSubspace.lean}). The new file
\texttt{Quantum/SpinS/MagnetizationDirectSum.lean} proves
\texttt{magSubspaceS\_eq\_eigenspace} (bridge to
\texttt{Module.End.eigenspace} of
$(\hat S^z_{\rm tot}).\texttt{mulVecLin}$), which lets
\texttt{magSubspaceS\_iSupIndep} inherit from the mathlib
\texttt{Module.End.eigenspaces\_iSupIndep} (applicable because
$\mathbb{C}$ is a field, hence \texttt{IsDomain} and
\texttt{IsTorsionFree}). Combined with the existing
\texttt{iSup\_magSubspaceS\_eq\_top} (in
\texttt{Quantum/SpinS/Magnetization.lean}), this gives
\texttt{magSubspaceS\_isInternal}: the multi-site Hilbert space
$(\Lambda\to\textsf{Fin}(N+1))\to\mathbb{C}$ decomposes as the
internal direct sum of the magnetization subspaces. As corollaries,
the PR~\#887 ladder iterates
$(\hat S^\mp_{\rm tot})^k\,|\sigma_{\top/\bot}\rangle$ are placed in
their respective magnetization sectors via
\texttt{mem\_magSubspaceS\_iff}.

\textbf{Once-lowered/raised state is non-zero on the saturated
ferromagnet (PR~\#890)}: foundational non-vanishing for the
saturated-ferromagnet ladder. In the new file
\texttt{Quantum/SpinS/LadderIterateNonvanishing.lean} the
configuration \texttt{oneFlippedUpConfig V x\_0 hN} (requires
$0<N$) flips one site $x_0:V$ from $0$ to $1\in\textsf{Fin}(N+1)$
and leaves the rest at $0$. The explicit value identity
\texttt{totalSpinSOpMinus\_mulVec\_allAlignedStateS\_zero\_at\_oneFlippedUpConfig}
gives $((\hat S^-_{\rm tot}\,|\sigma_\top\rangle))(\texttt{oneFlipped})=\sqrt{N}$:
distribute $\hat S^-_{\rm tot}=\sum_x\texttt{onSiteS}\,x\,\hat S^-$
via \texttt{Matrix.sum\_mulVec}, isolate only the pivot $x=x_0$ via
\texttt{sum\_eq\_single} (off-site agreement fails for
$x\neq x_0$), reduce via \texttt{spinSOpMinus\_apply\_lower}
$(0+1=1)=\sqrt{N\cdot 1}$. The corollary
\texttt{totalSpinSOpMinus\_mulVec\_allAlignedStateS\_zero\_ne\_zero}
follows for $0<N$ and $[\texttt{Nonempty V}]$ since
$\sqrt{N}>0$. The symmetric raising side
\texttt{oneFlippedDownConfig\,/\,..\_at\_oneFlippedDownConfig\,/\,..\_ne\_zero}
covers $\hat S^+_{\rm tot}\,|\sigma_\bot\rangle$.

\textbf{Multi-site Cartan \(+\)/\(-\) relation (PR~\#893)}: completes the
SU(2) algebra of total operators with
$[\hat S^+_{\rm tot}, \hat S^-_{\rm tot}] = 2\,\hat S^z_{\rm tot}$
(and the antisymmetric corollary $-2\,\hat S^z_{\rm tot}$). The
new file \texttt{Quantum/SpinS/MultiSiteCartanPlusMinus.lean}
lifts the single-site \texttt{spinSOpPlus\_commutator\_spinSOpMinus}
to multi-site via \texttt{onSiteS\_commutator\_totalOnSiteS}
(off-site terms cancel via cross-site commutativity, on-site reduces
to \texttt{onSiteS x ((2:C) • spinSOp3 N)}), then applies linearity
of \texttt{onSiteS} in the matrix argument and \texttt{Finset.smul\_sum}.

\textbf{Casimir rearrangement (PR~\#894)}: the standard SU(2)
Casimir identity at the multi-site level
$\hat S^+_{\rm tot} \hat S^-_{\rm tot} =
 (\hat S_{\rm tot})^2 - (\hat S^z_{\rm tot})^2 + \hat S^z_{\rm tot}$
(and symmetric $-\hat S^z_{\rm tot}$ for MinusPlus order). The new
file \texttt{Quantum/SpinS/CasimirRearrangement.lean} first proves
the sum identity
$\hat S^+ \hat S^- + \hat S^- \hat S^+ = 2\,((\hat S^{(1)}_{\rm tot})
^2 + (\hat S^{(2)}_{\rm tot})^2)$ by expanding both
$(A \pm iB)(A \mp iB)$ products — the $\pm i [A, B]$ cross terms
cancel — then combines with the Cartan \(+\)/\(-\) (\#893) and divides by
2 via \texttt{smul\_right\_injective}. Substituting
$(\hat S^{(1)}_{\rm tot})^2 + (\hat S^{(2)}_{\rm tot})^2 =
 (\hat S_{\rm tot})^2 - (\hat S^{(3)}_{\rm tot})^2$
via \texttt{totalSpinSSquared\_def} closes the identity.

\textbf{Inductive non-vanishing of $(\hat S^-_{\rm tot})^k\,|\sigma_\top\rangle$
for $k\leq 2m_{\max}$ (PR~\#895)}: the technical hurdle for the
full ladder linear independence. In the new file
\texttt{Quantum/SpinS/IterateInductiveNonvanishing.lean} the
eigenvalue identity
$\hat S^+_{\rm tot}\,((\hat S^-_{\rm tot})^{k+1}|\sigma_\top\rangle)
 = (k+1)(|V|\cdot N - k)\,(\hat S^-_{\rm tot})^k|\sigma_\top\rangle$
is derived from the Casimir rearrangement (\#894) plus iterate
eigenvalue identities (\#882 Casimir, \#887 $\hat S^z$ — applied
twice for $(\hat S^z)^2$). Then induction: base $k=0$ is
\texttt{allAlignedStateS\_ne\_zero} (\#891); step assumes
$v_k \neq 0$ and $k < |V|\cdot N$, contradicts $v_{k+1} = 0$ via
the scalar-mul cancellation (the scalar $(k+1)(|V|\cdot N - k)$
is non-zero in $\mathbb{C}$ for $k < |V|\cdot N$).

\textbf{Full saturated-ferromagnet ladder LI (PR~\#896)}: the
$(2m_{\max} + 1)$-element ladder basis statement. In the new file
\texttt{Quantum/SpinS/SaturatedFullLadderLI.lean} the family
\texttt{ladderIterateUp V N k} (parameterised by $k:\textsf{Fin}
(|V|\cdot N+1)$) and the eigenvalue function
\texttt{ladderEigenvalueUp V N k = $m_{\max}-k$} are introduced.
Injectivity of the eigenvalue function is by
\texttt{linear\_combination -hij}; per-$k$
\texttt{Module.End.HasEigenvector} witnesses combine
\texttt{ladderIterateUp\_mem\_eigenspace} (PR \#887 + \#889 bridge)
with \texttt{ladderIterateUp\_hasEigenvector}'s non-vanishing
(PR \#895, with the bound \texttt{Nat.lt\_succ\_iff.mp k.isLt}).
The main theorem
\texttt{ladderIterateUp\_linearIndependent} applies
\texttt{Module.End.eigenvectors\_linearIndependent'} to assemble
the LI conclusion. This completes the Tasaki §2.4 saturated-
ferromagnet ground-state ladder basis identification: the
$(2m_{\max} + 1)$-vector family
$\{(\hat S^-_{\rm tot})^k\,|\sigma_\top\rangle :
 k = 0,\ldots,2m_{\max}\}$ is linearly independent, identifying
the ground-state sector with the $J_{\rm tot} = m_{\max}$
irreducible SU(2) representation.

\textbf{Pairwise orthogonality of the ladder iterates (PR~\#898)}:
strengthens the LI conclusion to orthogonality. The new file
\texttt{Quantum/SpinS/SaturatedFullLadderOrthogonality.lean}
proves the generic helper
\texttt{Matrix.IsHermitian.dotProduct\_eq\_zero\_of\_eigenvalues\_ne}:
for a Hermitian matrix $M$ and two eigenvectors at distinct real
eigenvalues, $\alpha\,\langle v,w\rangle = \langle Mv, w\rangle =
\langle v, Mw\rangle = \beta\,\langle v,w\rangle$, hence
$\langle v,w\rangle = 0$. The shift to the left side uses
\texttt{Matrix.dotProduct\_mulVec} and \texttt{Matrix.star\_mulVec}
combined with Hermiticity. Applied to $\hat S^z_{\rm tot}$
(Hermitian by \texttt{totalSpinSOp3\_isHermitian}) with the
real eigenvalues $m_{\max} - k$ (from \texttt{Complex.conj\_natCast})
and the injectivity of \texttt{ladderEigenvalueUp}.

\textbf{Ladder iterates lie in a single Heisenberg eigenspace
(PR~\#899)}: defines
\texttt{saturatedFerromagnetEigenvalueS J N} as the H diagonal
at the all-up configuration; bridges PR \#881 to
\texttt{Module.End.eigenspace} membership via
\texttt{Module.End.mem\_eigenspace\_iff} +
\texttt{Matrix.mulVecLin\_apply}. Combined with the LI of \#896
and \texttt{LinearIndependent.fintype\_card\_le\_finrank}, gives
the dimension lower bound
\texttt{heisenbergHamiltonianS\_eigenspace\_finrank\_ge\_succ\_card\_mul\_N}:
the H-eigenspace at $c_J$ has $\dim_{\mathbb{C}} \geq 2m_{\max}+1$.
The LI family is restricted to the eigenspace via the subtype
embedding \texttt{E.subtype} and
\texttt{LinearIndependent.of\_comp}.

\textbf{Ladder iterates lie in a single Casimir eigenspace
(PR~\#900)}: mirror of PR \#899 for the Casimir operator
$(\hat S_{\rm tot})^2$. Defines
\texttt{saturatedFerromagnetCasimirEigenvalueS V N} as
$m_{\max}(m_{\max}+1)$, bridges PR \#882 to
\texttt{Module.End.eigenspace}, and provides the parallel
dimension lower bound
\texttt{totalSpinSSquared\_eigenspace\_finrank\_ge\_succ\_card\_mul\_N}.

\textbf{Simultaneous (H, $\hat S_{\rm tot}^2$, $\hat S^z_{\rm tot}$)
eigenvector bundle (PR~\#901)}: packages PRs \#881, \#882, \#887,
\#895 into the operator-level form of Tasaki §2.4 (2.4.10):
each iterate $(\hat S^-_{\rm tot})^k\,|\sigma_\top\rangle$ is
non-zero and is simultaneously a Heisenberg eigenvector at
$c_J$, a Casimir eigenvector at $m_{\max}(m_{\max}+1)$, and an
$\hat S^z_{\rm tot}$ eigenvector at $m_{\max}-k$. The resulting
\texttt{ladderIterateUp\_simultaneous\_eigenvector} statement is
the structural identification of the unnormalised ferromagnetic
ground-state ladder iterates $|\Phi_M\rangle$ with the
$J_{\rm tot} = m_{\max}$ irreducible SU(2) representation
(eigenvalues $(c_J, m_{\max}(m_{\max}+1), M)$).

\textbf{Joint $(H, \hat S_{\rm tot}^2)$ eigenspace (PR~\#903)}: the
new file \texttt{Quantum/SpinS/SaturatedLadderJointEigenspace.lean}
defines \texttt{saturatedFerromagnetJointEigenspace J N} as the
meet ($\sqcap$) of the Heisenberg and Casimir eigenspaces at
$(c_J, m_{\max}(m_{\max}+1))$; combines \#899 and \#900 to show
each ladder iterate lies in the joint eigenspace, then derives
\texttt{saturatedFerromagnetJointEigenspace\_finrank\_ge\_succ\_card\_mul\_N}:
$\dim_\mathbb{C} \geq |V|\cdot N + 1 = 2m_{\max} + 1$.

\textbf{Extremal magnetisation sectors are 1-dimensional (PR~\#2759,
toward Tasaki §2.4 Theorem 2.1)}: in
\texttt{Quantum/SpinS/SaturatedLadderJointEigenspace.lean} two
companion theorems
\texttt{magSubspaceS\_mMax\_inf\_saturatedFerromagnetJointEigenspace}
and \texttt{magSubspaceS\_neg\_mMax\_inf\_saturatedFerromagnetJointEigenspace}
state
\[
  H_{\pm m_{\max}} \sqcap
  \texttt{saturatedFerromagnetJointEigenspace J N} =
  \mathbb{C}\,|\sigma_{\top/\bot}\rangle.
\]
Argument: the extremal magnetisation subspaces are 1-dimensional
(PR \#908, \texttt{magSubspaceS\_(neg\_)mMax\_eq\_span\_allAlignedStateS\_(zero/last)}),
and $|\sigma_\top\rangle = \texttt{ladderIterateUp V N 0}$ lies in
the joint eigenspace by the ladder-iterate membership lemma of
\#903; the corresponding $|\sigma_\bot\rangle$ membership is bridged
via \texttt{saturatedFerromagnetEigenvalueS\_explicit} (closed-form
eigenvalue identical at the two extremal configurations by SU(2)
symmetry). These are the first two concrete sector contributions
to the upper bound $\dim_\mathbb{C}\,\texttt{joint} \le 2m_{\max}+1$
that, combined with the lower bound of \#903 and the span-equals-joint
argument of \#904, closes Tasaki §2.4 Theorem 2.1 (saturated-
ferromagnet ground-state subspace = $(2m_{\max}+1)$-dim irreducible
SU(2) representation). The intermediate sectors $M$ with
$-m_{\max} < M < m_{\max}$ are handled by the $\hat S^\mp_{\rm tot}$
lowering-injectivity argument (PRs \#2761/\#2762/\#2763/\#2764) and
the magnetisation projector decomposition (PR \#2768).

\textbf{Raising-operator joint-magnetisation shift (PR~\#2760,
inductive step)}: the next theorem
\texttt{totalSpinSOpPlus\_mulVec\_mem\_saturatedFerromagnetJointEigenspace\_inf\_magSubspaceS}
states that, for any
$v \in \texttt{joint} \sqcap H_M$,
$\hat S^+_{\rm tot} \cdot v \in \texttt{joint} \sqcap H_{M+1}$.
Argument: $[\hat H, \hat S^+_{\rm tot}] = [\hat S_{\rm tot}^2,
\hat S^+_{\rm tot}] = 0$ (\#881, \#882) preserve the joint
eigenvalues, while
$[\hat S^z_{\rm tot}, \hat S^+_{\rm tot}] = +\hat S^+_{\rm tot}$
shifts the magnetic quantum number by $+1$
(\#887 \texttt{\_OpPlus\_mulVec\_eigenvec}). The companion lemma
\texttt{mulVec\_preserves\_eigenvalue\_of\_commuteS} is the spin-$S$
analogue of the generic eigenvalue propagation lemma in
\texttt{Quantum/TotalSpin/Casimir.lean}. Combined with the
$M = m_{\max}$ extremal 1-dim result of \#2759, this establishes the
induction skeleton that propagates the 1-dim bound to every sector,
once combined with the raising-operator injectivity on
$\texttt{joint} \sqcap H_{M}$ for $M < m_{\max}$ added in PR \#2761.

\textbf{Raising-operator injectivity (PR~\#2761, kernel-trivial
step)}: the next theorem
\texttt{totalSpinSOpPlus\_mulVec\_eq\_zero\_imp\_eq\_zero\_of\_mem\_saturatedFerromagnetJointEigenspace\_inf\_magSubspaceS}
states that, for any
$v \in \texttt{joint} \sqcap H_M$ with $M \neq m_{\max}$, if
$\hat S^+_{\rm tot} \cdot v = 0$ then $v = 0$. Argument: the
MinusPlus Casimir rearrangement
$\hat S^-_{\rm tot} \cdot \hat S^+_{\rm tot} =
\hat S_{\rm tot}^2 - (\hat S^z_{\rm tot})^2 - \hat S^z_{\rm tot}$
of \#894, applied to $v$ on the joint eigenspace with the
$\hat S^z_{\rm tot}$-eigenvalue $M$, gives
\[
  \hat S^-_{\rm tot} \cdot \hat S^+_{\rm tot} \cdot v
  = (m_{\max}(m_{\max}+1) - M^2 - M)\,v
  = (m_{\max} - M)(m_{\max} + M + 1)\,v.
\]
With $\hat S^+_{\rm tot} v = 0$, the LHS vanishes, so the scalar
factor annihilates $v$. The case $M = -m_{\max} - 1$ is ruled out by
PR~\#905 (\texttt{magSubspaceS V N (-m\_max - 1) = bottom}), and
$M \neq m_{\max}$ is the hypothesis; together they force $v = 0$.
Combined with the joint-magnetisation shift of \#2760, this means
$\hat S^+_{\rm tot}$ restricted to
$\texttt{joint} \sqcap H_M$ is an injective linear map into
$\texttt{joint} \sqcap H_{M+1}$ for every $M \neq m_{\max}$, hence
$\dim_\mathbb{C}(\texttt{joint} \sqcap H_M) \le
\dim_\mathbb{C}(\texttt{joint} \sqcap H_{M+1})$. Iterating from
$M = -m_{\max}$ upward propagates the 1-dim bound at $H_{m_{\max}}$
(from \#2759) down to every sector — the final summation step is
PR \#2768.

\textbf{Linear-map packaging and finrank chain (PR~\#2762)}: the
three theorems
\texttt{totalSpinSOpPlusJointMagShift J M},
\texttt{totalSpinSOpPlusJointMagShift\_injective}, and
\texttt{saturatedFerromagnetJointEigenspace\_inf\_magSubspaceS\_finrank\_le\_succ}
package PR \#2760 (joint-magnetisation shift) and PR \#2761
(kernel-trivial) into a single \texttt{LinearMap} between the
joint-magnetisation sectors and apply mathlib's
\texttt{LinearMap.finrank\_le\_finrank\_of\_injective} to derive
$\dim_\mathbb{C}(\texttt{joint} \sqcap H_M) \le
\dim_\mathbb{C}(\texttt{joint} \sqcap H_{M+1})$ for every
$M \neq m_{\max}$. PR \#2763 iterates this chain from $M = -m_{\max}$ upward, giving
the per-sector $\le 1$ bound at every $M$ in the spectrum; PR \#2768
sums the $(2m_{\max}+1)$ such bounds via the magnetisation projector
decomposition (\texttt{sum\_magProjFn\_eq}) and the joint-preservation
to derive $\texttt{joint} = \texttt{span}(\texttt{ladderIterateUp})$.

\textbf{Pointwise magnetisation projector (PR~\#2765, framework
for §2.4 Theorem 2.1)}: the definition
\texttt{magProjFn (M : $\mathbb{C}$) (v : (V $\to$ Fin (N+1)) $\to$ $\mathbb{C}$)}
keeps `v $\sigma$' when
$\texttt{magEigenvalueS}~\sigma = M$ and zeros it out otherwise.
The companion theorem
\texttt{magProjFn\_mem\_magSubspaceS} shows
$\texttt{magProjFn M v} \in \texttt{magSubspaceS V N M}$ by the
diagonal $\hat S^z_{\rm tot}$-eigenvalue equation.

\textbf{magProjFn linearity + commutation with H (PR~\#2766)}:
extends the framework with
(i) the magnetisation-subspace support property
\texttt{magSubspaceS\_apply\_eq\_zero\_of\_magEigenvalueS\_ne}
(vectors in $\texttt{magSubspaceS V N M}$ vanish at any
configuration with $\texttt{magEigenvalueS} \neq M$),
(ii) linearity \texttt{magProjFn\_add} / \texttt{magProjFn\_smul},
(iii) generic off-magnetisation matrix-entry vanishing
\texttt{matrix\_entry\_eq\_zero\_of\_mulVec\_basisVecS\_mem\_magSubspaceS},
and (iv) the Heisenberg commutation
\texttt{heisenbergHamiltonianS\_mulVec\_magProjFn\_eq}:
$\hat H \cdot \texttt{magProjFn M v} =
\texttt{magProjFn M} (\hat H \cdot v)$,
proved by case analysis on $\texttt{magEigenvalueS}~\sigma = M$
together with the matrix-entry vanishing for
$\hat H \cdot \texttt{basisVecS}~\tau$ (PR~\#1078). \textbf{Casimir commutation + joint preservation (PR~\#2767)}:
extends the framework with (v) the Casimir commutation
\texttt{totalSpinSSquared\_mulVec\_magProjFn\_eq}:
$(\hat S_{\rm tot})^2 \cdot \texttt{magProjFn M v} =
\texttt{magProjFn M} ((\hat S_{\rm tot})^2 \cdot v)$
(same matrix-entry-vanishing argument as the Heisenberg case),
and (vi) the joint preservation
\texttt{magProjFn\_mem\_saturatedFerromagnetJointEigenspace}:
$\texttt{magProjFn M v} \in \texttt{joint}$ for $v \in \texttt{joint}$,
by combining the two commutations with \texttt{magProjFn\_smul} —
the H- and Casimir-eigenvalues factor through the projector since
$\texttt{magProjFn M}(c \cdot v) = c \cdot \texttt{magProjFn M v}$.

\textbf{Tasaki §2.4 Theorem 2.1 closure (PR~\#2768)}:
the closure combines three theorems.
(i) \texttt{sum\_magProjFn\_eq}: every $v$ decomposes as
$v = \sum_{k : \mathrm{Fin}(|V|\cdot N + 1)}
\texttt{magProjFn}(m_{\max} - k.\textrm{val})~v$
(each $\sigma$ has unique $\texttt{magSumS}~\sigma \in \{0, \ldots,
|V|\cdot N\}$ selecting exactly one term).
(ii) \texttt{saturatedFerromagnetJointEigenspace\_le\_span\_ladderIterateUp}:
$\texttt{joint} \subseteq \texttt{span}(\mathrm{range}(\texttt{ladderIterateUp}))$.
For $v \in \texttt{joint}$, the joint preservation of \#2767 puts
each magnetisation component into
$\texttt{joint} \sqcap H_{m_{\max} - k.\textrm{val}} =
\mathbb{C}\,\texttt{ladderIterateUp V N k}$ (per the per-sector
identification of \#2764). Summing across the $2m_{\max}+1$ sectors
yields $v \in \texttt{span}(\mathrm{range}(\texttt{ladderIterateUp}))$.
(iii) \texttt{saturatedFerromagnetJointEigenspace\_eq\_span\_ladderIterateUp}:
the final equality $\texttt{joint} =
\texttt{span}(\mathrm{range}(\texttt{ladderIterateUp}))$ via
\texttt{le\_antisymm} with the reverse inclusion of PR~\#904.

Identifies the saturated-ferromagnet joint $(H, \hat S_{\rm tot}^2)$
eigenspace as exactly the $(2m_{\max}+1)$-dimensional irreducible
SU(2) representation generated by the lowering operator
$\hat S^-_{\rm tot}$ from the highest-weight state
$|\sigma_\top\rangle$. Combined with PRs \#896 (linear independence),
\#898 (orthogonality), \#903 (lower bound), and \#904
($\dim_\mathbb{C} = 2m_{\max}+1$), this is the operator-level form
of Tasaki §2.4 Theorem 2.1.

\textbf{Per-sector identification (PR~\#2764)}: the theorem
\texttt{saturatedFerromagnetJointEigenspace\_inf\_magSubspaceS\_eq\_span\_ladderIterateUp}
states that for every $k \in \mathrm{Fin}(|V|\cdot N + 1)$,
\[
  \texttt{joint} \sqcap H_{m_{\max} - k}
  = \mathbb{C}\,\texttt{ladderIterateUp V N k}.
\]
Argument: the LHS is $\le 1$-dim (\#2763) and contains the non-zero
ladder iterate $(\hat S^-_{\rm tot})^k |\sigma_\top\rangle$
(joint membership: \#903 \texttt{ladderIterateUp\_mem\_saturatedFerromagnetJointEigenspace};
magnetisation membership: \#889 via \#887 \texttt{totalSpinSOpMinus\_pow\_allAlignedStateS\_zero\_mem\_magSubspaceS};
non-vanishing: \#895). Conclude by
\texttt{Submodule.eq\_of\_le\_of\_finrank\_le}. Each ladder iterate
is thereby identified as the unique (up to scaling) non-zero
vector in its joint-magnetisation sector. Summed across the
$2m_{\max}+1$ values of $k$ in PR~\#2768 via the magnetisation
projector decomposition (\texttt{sum\_magProjFn\_eq}) and joint
preservation to identify
$\texttt{joint} = \texttt{span}(\texttt{ladderIterateUp})$ and
complete Tasaki §2.4 Theorem 2.1.

\textbf{Per-sector $\dim \le 1$ (PR~\#2763)}: the theorem
\texttt{saturatedFerromagnetJointEigenspace\_inf\_magSubspaceS\_finrank\_le\_one}
states that for every $k \in \mathbb{N}$,
$\dim_\mathbb{C}(\texttt{joint} \sqcap H_{m_{\max} - k}) \le 1$.
The proof is by induction on $k$: base case $k = 0$ uses \#2759
to identify $\texttt{joint} \sqcap H_{m_{\max}} = \mathbb{C}\,|\sigma_\top\rangle$
and \texttt{finrank\_span\_singleton} to conclude $\dim = 1$;
the step case uses the finrank chain of \#2762 to bound
$\dim_\mathbb{C}(\texttt{joint} \sqcap H_{m_{\max} - (k+1)}) \le
\dim_\mathbb{C}(\texttt{joint} \sqcap H_{m_{\max} - k}) \le 1$ by
the inductive hypothesis. The $2m_{\max}+1$ spectrum values
$M = m_{\max} - k$ for $k \in \{0, 1, \ldots, 2m_{\max}\}$ now each
contribute at most $1$ to $\dim_\mathbb{C}(\texttt{joint})$; the
final summation step (PR \#2768) uses the magnetisation projector
decomposition (\texttt{sum\_magProjFn\_eq}) together with the
joint-preservation (\#2767) to conclude
$\texttt{joint} = \texttt{span}(\texttt{ladderIterateUp})$ and hence
$\dim_\mathbb{C}(\texttt{joint}) = 2m_{\max} + 1$, completing
Tasaki §2.4 Theorem 2.1.

\textbf{Ladder span finrank (PR~\#904)}: the new file
\texttt{Quantum/SpinS/SaturatedLadderSpan.lean} applies mathlib's
\texttt{finrank\_span\_eq\_card} to the LI family of \#896,
yielding
\texttt{ladderIterateUp\_span\_finrank\_eq\_succ\_card\_mul\_N}:
$\dim_\mathbb{C} \texttt{Submodule.span} = |V|\cdot N + 1 =
2m_{\max} + 1$. The span is contained in the joint eigenspace
of \#903 via \texttt{Submodule.span\_le}.

\textbf{Finrank-equality corollary (PR~\#2769)}: in
\texttt{Quantum/SpinS/SaturatedLadderSpan.lean},
\texttt{saturatedFerromagnetJointEigenspace\_finrank\_eq} promotes
the lower bound of \#903 to the operator-level equality
$\dim_\mathbb{C}(\texttt{joint J N}) = |V|\cdot N + 1 = 2m_{\max}+1$
by rewriting through \#2768's \texttt{\_eq\_span\_ladderIterateUp}
and applying the span-finrank result of \#904. Identifies the
saturated-ferromagnet joint eigenspace as the unique
$(2m_{\max}+1)$-dimensional irreducible $SU(2)$ representation at
the operator level.

\textbf{$\hat S^-_{\rm tot}$ invariance (PR~\#2770)}: the new file
\texttt{Quantum/SpinS/SaturatedLadderLoweringInvariant.lean} adds
the operational lowering-side $su(2)$ statement.
Per-iterate action:
\texttt{totalSpinSOpMinus\_mulVec\_ladderIterateUp\_interior} gives
$\hat S^-_{\rm tot}\cdot\texttt{ladderIterateUp V N }k =
\texttt{ladderIterateUp V N }\langle k+1, \cdot\rangle$ for $k+1 <
|V|\cdot N + 1$ (definitionally from the iterate definition plus
$\texttt{pow\_succ'}$);
\texttt{totalSpinSOpMinus\_mulVec\_ladderIterateUp\_last} gives
$\hat S^-_{\rm tot}\cdot\texttt{ladderIterateUp V N }(\texttt{Fin.last})
= 0$ (boundary annihilation via \#905). Lifted via
\texttt{Submodule.map\_span} + \texttt{Submodule.span\_le} on
\#2768's $\texttt{joint} = \texttt{span}(\texttt{ladderIterateUp})$
to the global invariance
$\texttt{Submodule.map}\,\hat S^-_{\rm tot}.\texttt{mulVecLin}\,
(\texttt{joint J N}) \le \texttt{joint J N}$. Combined with the
finrank equality of \#2769 and the raising-side annihilation of
\#907, this identifies the joint eigenspace as the
$(2m_{\max}+1)$-dimensional spin-$m_{\max}$ representation of
$su(2)$ under the total-spin generators.

\textbf{$\hat S^+_{\rm tot}$ invariance (PR~\#2771)}: the new file
\texttt{Quantum/SpinS/SaturatedLadderRaisingInvariant.lean}
completes the operational $su(2)$ statement on the raising side.
Per-iterate action:
\texttt{totalSpinSOpPlus\_mulVec\_ladderIterateUp\_zero} gives
$\hat S^+_{\rm tot}\cdot\texttt{ladderIterateUp V N }\langle 0,
\cdot\rangle = 0$ (top-weight annihilation, via PR \#877's
\texttt{totalSpinSOpPlus\_mulVec\_allAlignedStateS\_zero});
\texttt{totalSpinSOpPlus\_mulVec\_ladderIterateUp\_succ} gives
$\hat S^+_{\rm tot}\cdot\texttt{ladderIterateUp V N }\langle j+1,
\cdot\rangle = ((j+1)\cdot(|V|\cdot N - j))\cdot
\texttt{ladderIterateUp V N }\langle j, \cdot\rangle$ (Casimir-
rearrangement scalar identity from PRs \#882, \#887, \#894 packaged
as
\texttt{totalSpinSOpPlus\_mulVec\_totalSpinSOpMinus\_pow\_succ\_allAlignedStateS\_zero}).
Lifted via the same
\texttt{Submodule.map\_span} + \texttt{Submodule.span\_le} pattern
of PR \#2770 to the global invariance
$\texttt{Submodule.map}\,\hat S^+_{\rm tot}.\texttt{mulVecLin}\,
(\texttt{joint J N}) \le \texttt{joint J N}$. Together with
\#2770, this establishes invariance of \texttt{joint J N} under
both $su(2)$ ladder generators, completing the operational
identification with the $(2m_{\max}+1)$-dim spin-$m_{\max}$
irreducible representation.

\textbf{Full Cartesian $su(2)$ closure (PR~\#2772)}: the new file
\texttt{Quantum/SpinS/SaturatedLadderCartesianInvariant.lean}
completes the operational $su(2)$-module structure of the
saturated joint eigenspace. Two total-operator identities lift the
single-site definitions of \texttt{spinSOp1} / \texttt{spinSOp2}:
$\hat S^{(1)}_{\rm tot} = (1/2)(\hat S^+_{\rm tot} +
\hat S^-_{\rm tot})$
(\texttt{totalSpinSOp1\_eq\_one\_half\_smul\_plus\_add\_minus}) and
$\hat S^{(2)}_{\rm tot} = (1/(2 i))(\hat S^+_{\rm tot} -
\hat S^-_{\rm tot})$
(\texttt{totalSpinSOp2\_eq\_one\_over\_two\_I\_smul\_plus\_sub\_minus}),
both via \texttt{onSiteS\_smul} + \texttt{onSiteS\_add/sub}.
The $\hat S^{(1)}_{\rm tot}$ and $\hat S^{(2)}_{\rm tot}$
invariances
(\texttt{saturatedFerromagnetJointEigenspace\_totalSpinSOp\{1,2\}\_invariant})
reduce via these decompositions to the ladder invariances of PRs
\#2770 / \#2771 plus submodule add/sub/smul closure. The
$\hat S^{(3)}_{\rm tot}$ case
(\texttt{saturatedFerromagnetJointEigenspace\_totalSpinSOp3\_invariant})
is independent: each ladder iterate is an
$\hat S^z_{\rm tot}$-eigenvector at eigenvalue $m_{\max}-k$
(\texttt{ladderIterateUp\_mem\_eigenspace}), so its
$\hat S^z_{\rm tot}$-image is a scalar multiple of the same
iterate, still in $\texttt{span}(\texttt{ladderIterateUp}) =
\texttt{joint J N}$ (\#2768). The three together establish
invariance under all $su(2)$ generators acting through the
total-spin operators, completing the operational identification of
the saturated joint eigenspace with the $(2m_{\max}+1)$-dim
irreducible representation of $su(2)$ at top weight $m_{\max}$.
This finishes the operational closure of Tasaki §2.4 Theorem 2.1
on top of \#2768 (joint = span ladder), \#2769 (finrank
$= 2m_{\max}+1$), \#2770 ($\hat S^-_{\rm tot}$ invariance), and
\#2771 ($\hat S^+_{\rm tot}$ invariance).

\textbf{Boundary annihilation (PR~\#905)}: the new file
\texttt{Quantum/SpinS/LadderBoundaryAnnihilation.lean} caps the
ladder at exactly $2m_{\max}+1$ non-zero terms by proving
\texttt{totalSpinSOpMinus\_pow\_succ\_card\_mul\_N\_allAlignedStateS\_zero}:
$(\hat S^-_{\rm tot})^{|V|\cdot N + 1}\,|\sigma_\top\rangle = 0$.
Argument: PR \#887 places the iterate in the
$\hat S^z_{\rm tot}$-eigenspace at $-m_{\max}-1$; the helper
\texttt{magSubspaceS\_eq\_bot\_of\_not\_in\_spectrum} shows that
this magnetization subspace is $\bot$ (since $\hat S^z_{\rm tot}$
is diagonal in the \texttt{basisVecS} basis with eigenvalues
\texttt{magEigenvalueS}~$\sigma \in [-m_{\max}, m_{\max}]$, never
$-m_{\max}-1$).

\textbf{Symmetric raising-side boundary annihilation (PR~\#907)}:
extends \texttt{Quantum/SpinS/LadderBoundaryAnnihilation.lean} with
\texttt{magEigenvalueS\_ne\_mMax\_add\_one} and
\texttt{totalSpinSOpPlus\_pow\_succ\_card\_mul\_N\_allAlignedStateS\_last}:
$(\hat S^+_{\rm tot})^{|V|\cdot N + 1}\,|\sigma_\bot\rangle = 0$
via $m_{\max}+1$ off-spectrum.

\textbf{Extremal magnetization subspaces are 1-dimensional
(PR~\#908)}: the new file
\texttt{Quantum/SpinS/MagSubspaceExtremalDim.lean} proves
\texttt{magEigenvalueS\_eq\_(neg\_)mMax\_iff\_allAlignedConfigS\_(last/zero)}:
the eigenvalue $\pm m_{\max}$ is achieved exactly at the all-up
(resp. all-down) constant configuration. Combined with the
diagonality of $\hat S^z_{\rm tot}$, this yields
\texttt{magSubspaceS\_(neg\_)mMax\_eq\_span\_allAlignedStateS\_(last/zero)}:
\texttt{magSubspaceS V N $(\pm m_{\max})$ = span $\mathbb{C}$
\{$|\sigma_{\top/\bot}\rangle$\}}. Analytic counterpart of PR \#885's
combinatorial cardinality.

\textbf{Maximal-lowering iterate identification (PR~\#909)}: the new
file \texttt{Quantum/SpinS/LadderBottom.lean} combines PR \#887's
eigenvalue identity at $k = |V|\cdot N$ ($-m_{\max}$ eigenvalue)
with PR \#908's extremal subspace identification, giving
\texttt{totalSpinSOpMinus\_pow\_card\_mul\_N\_allAlignedStateS\_zero\_eq\_smul\_last}:
there exists a non-zero $c\in\mathbb{C}$ with
$(\hat S^-_{\rm tot})^{|V|\cdot N}\,|\sigma_\top\rangle =
c\,|\sigma_\bot\rangle$. Non-vanishing of $c$ uses PR \#895.
A symmetric span-membership statement is provided for the raising
side; the non-zero scalar version is in PR \#910.

\textbf{Symmetric inductive non-vanishing for raising side
(PR~\#910)}: the new file
\texttt{Quantum/SpinS/IterateInductiveNonvanishingPlus.lean} mirrors
PR \#895 using PR \#894's MinusPlus Casimir rearrangement
$\hat S^- \hat S^+ = \hat S_{\rm tot}^2 - (\hat S^z_{\rm tot})^2 -
\hat S^z_{\rm tot}$. The same eigenvalue scalar $(k+1)(2m_{\max}-k)$
arises with the raising-side $\hat S^z_{\rm tot}$ eigenvalue
$-m_{\max}+k$ (PR \#887 raising side), giving
\texttt{totalSpinSOpPlus\_pow\_allAlignedStateS\_last\_ne\_zero} for
$k\leq|V|\cdot N$ and the non-zero scalar identification
\texttt{totalSpinSOpPlus\_pow\_card\_mul\_N\_allAlignedStateS\_last\_eq\_smul\_zero}.

\textbf{Round-trip identity (PR~\#912)}: the new file
\texttt{Quantum/SpinS/LadderRoundTrip.lean} composes \#909 (lowering
reaches $c_1\,|\sigma_\bot\rangle$) with \#910 (raising from
$|\sigma_\bot\rangle$ reaches $c_2\,|\sigma_\top\rangle$) to give
\texttt{totalSpinSOpPlus\_pow\_mulVec\_totalSpinSOpMinus\_pow\_allAlignedStateS\_zero\_eq\_smul}:
the round-trip $(\hat S^+_{\rm tot})^{|V|\cdot N}\,(\hat S^-_{\rm tot})^{|V|\cdot N}\,|\sigma_\top\rangle$
returns to a non-zero $c_1c_2\cdot|\sigma_\top\rangle$. The
saturated-ferromagnet ladder closes on itself.

\textbf{Expectation values on all-aligned states (PR~\#913)}: the new
file \texttt{Quantum/SpinS/AllAlignedStateExpectations.lean} computes
the inner-product expectation values
$\langle v, w\rangle := \texttt{dotProduct (star v) w}$ on the
all-aligned states: norm-squared 1
(\texttt{basisVecS\_inner\_self} / \texttt{allAlignedStateS\_inner\_self}),
$\hat S^z_{\rm tot}$ expectation $\pm m_{\max}$, and Casimir
$\hat S_{\rm tot}^2$ expectation $m_{\max}(m_{\max}+1)$. These are
the explicit $J = m_{\max}$ quantum numbers.

\textbf{basisVecS orthonormality (PR~\#914)}: the new file
\texttt{Quantum/SpinS/BasisVecSOrthonormal.lean} extends \#913 with
the off-diagonal case: distinct configurations $\sigma\neq\tau$
give orthogonal basis vectors. Bundled in Kronecker form
\texttt{basisVecS\_inner\_kronecker}.

\textbf{Ladder iterate expectation values (PR~\#915)}: the new file
\texttt{Quantum/SpinS/IterateExpectations.lean} packages the
eigenvalue identities (\#881 H, \#882 Casimir, \#887 $\hat S^z$)
into the inner-product form
$\langle v_k, A\,v_k\rangle = \alpha\,\langle v_k, v_k\rangle$ for
each iterate $v_k = (\hat S^-_{\rm tot})^k\,|\sigma_\top\rangle$ and
$A\in\{\hat S^z_{\rm tot}, \hat S_{\rm tot}^2, H\}$ with the
corresponding eigenvalue $\alpha$.

\textbf{basisVecS spans + linearly independent (PR~\#917)}: the new
file \texttt{Quantum/SpinS/BasisVecSBasis.lean} proves
\texttt{basisVecS\_span\_eq\_top} (via
\texttt{fun\_eq\_sum\_smul\_basisVecS}) and
\texttt{basisVecS\_linearIndependent} (via pointwise evaluation of
zero-sums). These structural facts identify
\texttt{basisVecS} as a basis of the multi-site Hilbert space.

\textbf{Multi-site dimension formula (PR~\#918)}: the new file
\texttt{Quantum/SpinS/MultiSiteFinrank.lean} proves
\texttt{multiSiteSpinS\_finrank}: $\dim_\mathbb{C} ((V \to
\textsf{Fin}\,(N+1)) \to \mathbb{C}) = (N+1)^{|V|}$. Direct from
mathlib's \texttt{Module.finrank\_fintype\_fun\_eq\_card} +
\texttt{Fintype.card\_fun} + \texttt{Fintype.card\_fin}. The
standard quantum-mechanical dimension formula $(2S+1)^{|\Lambda|}$.

\textbf{basisSpinS as Module.Basis (PR~\#919)}: the new file
\texttt{Quantum/SpinS/BasisSpinS.lean} bundles \#917 into a
\texttt{Module.Basis (V $\to$ Fin (N+1)) $\mathbb{C}$
((V $\to$ Fin (N+1)) $\to$ $\mathbb{C}$)} via
\texttt{Module.Basis.mk}, with the apply equality
\texttt{basisSpinS V N $\sigma$ = basisVecS $\sigma$}.

\textbf{Universal single-site Casimir expectation (PR~\#920)}: the
new file \texttt{Quantum/SpinS/SingleSiteCasimirExpectation.lean}
combines \texttt{spinSDot\_self} ($\hat{S}_x \cdot \hat{S}_x =
(N(N+2)/4)\,\mathbf{1}$) with \texttt{dotProduct\_smul} to give
\texttt{spinSDot\_self\_expectation\_normalized}: $\langle\Phi,
\hat{S}_x \cdot \hat{S}_x\,\Phi\rangle = S(S+1)$ for any normalized
$\Phi$, independent of state — the universal single-site Casimir
identity. Foundation for Tasaki Problem 2.5.c ($\gamma$-7).

\textbf{Problem 2.5.c single-site squared bridge (PR~\#4056)}: the
new Problem 2.5.c module packages the squared-axis expectation
\[
  E_\alpha(x,\Phi)=
  \langle \Phi,(\hat{S}^{(\alpha)}_x)^2\Phi\rangle .
\]
The theorem
\texttt{singleSiteSpinSquareExpectationS\_axis\_sum} expands the
definition of \(\hat{S}_x\cdot \hat{S}_x\) and proves
\[
  E_1(x,\Phi)+E_2(x,\Phi)+E_3(x,\Phi)
  =
  \langle\Phi,\hat{S}_x\cdot \hat{S}_x\,\Phi\rangle .
\]
Combining this with the normalized Casimir identity above gives the
sum \(N(N+2)/4=S(S+1)\).  The all-axis wrapper then records the
algebraic reduction needed for Tasaki Problem~2.5.c: if the later
AFM SU(2)-symmetry input supplies
\(E_1=E_2=E_3\), each axis has expectation
\[
  E_\alpha(x,\Phi)=\frac{N(N+2)}{12}=\frac{S(S+1)}{3}.
\]
This PR deliberately leaves the equal-axis ground-state symmetry as an
explicit hypothesis.  Reference: Tasaki, Springer 2020, Problem~2.5.c,
p.~43.

\textbf{Problem 2.5.c unitary symmetry input (PR~\#4057)}: the
next module records the generic symmetry mechanism behind the equal-axis
input.  If a many-body unitary \(T\) has inverse \(T^{-1}=T^\dagger\),
conjugates a single-site component \(\hat S_x^{(a)}\) to
\(\hat S_x^{(b)}\), and fixes the state through \(T^{-1}\Phi=\Phi\), then
\[
  \langle\Phi,(\hat S_x^{(b)})^2\Phi\rangle
  =
  \langle\Phi,(\hat S_x^{(a)})^2\Phi\rangle .
\]
The proof expands \((\hat S_x^{(b)})^2=T(\hat S_x^{(a)})^2T^{-1}\),
uses the state-invariance hypothesis on the ket side, and moves \(T\) to
the bra side as \(T^\dagger=T^{-1}\).  The existing axis-swap API can
supply such a conjugation hypothesis for axes \(2\) and \(3\); the remaining
work is the AFM ground-state invariance or fuller SU(2) rotation input.
Reference: Tasaki, Springer 2020, Problem~2.5.c, p.~43, and the
Theorem~2.4 symmetry context, pp.~43--44.

\textbf{Problem 2.5.c lifted axis-swap adjoint input (PR~\#4058)}:
the next module supplies the concrete adjoint data needed to instantiate the
unitary bridge with the spin-\(S\) axis swap from the Theorem~2.4
infrastructure.  Entrywise,
\[
  (\Theta_W^\dagger)_{\sigma',\sigma}
  =
  \overline{\prod_x W_x(\sigma_x,\sigma'_x)}
  =
  \prod_x W_x^\dagger(\sigma'_x,\sigma_x),
\]
so \((\otimes_x W_x)^\dagger=\otimes_x W_x^\dagger\).  Applying this to
the rotation \(U=\exp(-i\pi\hat S^{(1)}/2)\), together with the existing
single-site theorem \(U^\dagger=U^{-1}\), gives
\((\otimes_x U)^\dagger=\otimes_x U^{-1}\).  The module then combines this
with the existing conjugation identity
\((\otimes_x U)\hat S_x^{(2)}(\otimes_x U^{-1})=\hat S_x^{(3)}\) and the
PR~\#4057 bridge to prove
\[
  \langle\Phi,(\hat S_x^{(3)})^2\Phi\rangle
  =
  \langle\Phi,(\hat S_x^{(2)})^2\Phi\rangle
\]
under the explicit invariance hypothesis \((\otimes_x U^{-1})\Phi=\Phi\).
The AFM ground-state invariance and the additional rotation data needed to
equate all three axes remain separate inputs.  Reference: Tasaki, Springer
2020, Problem~2.5.c, p.~43, and the Theorem~2.4 symmetry context,
pp.~43--44.

\textbf{Problem 2.5.c axis-swap equal-axes wrapper (PR~\#4059)}:
the next module combines the preceding two Problem~2.5.c inputs.  Write
\[
  E_\alpha(x,\Phi)=
  \langle\Phi,(\hat S_x^{(\alpha)})^2\Phi\rangle .
\]
For normalized \(\Phi\), PR~\#4056 proves the desired value from
\(E_1=E_2=E_3\).  PR~\#4058 supplies \(E_3=E_2\) whenever the inverse
lifted axis-swap unitary fixes \(\Phi\).  Therefore the wrapper records that
\[
  (\otimes_x U^{-1})\Phi=\Phi,\qquad E_1(x,\Phi)=E_2(x,\Phi)
\]
together imply
\[
  E_1(x,\Phi)=E_2(x,\Phi)=E_3(x,\Phi)=\frac{N(N+2)}{12}.
\]
The remaining equality \(E_1=E_2\) is still explicit: current general
spin-\(S\) rotation infrastructure only supplies the axis-1 rotation used to
swap axes \(2\) and \(3\).  Reference: Tasaki, Springer 2020,
Problem~2.5.c, p.~43, and the Theorem~2.4 symmetry context, pp.~43--44.

\textbf{Problem 2.5.c two-symmetry axis input (PR~\#4060)}:
the next module replaces the bare equality \(E_1=E_2\) in PR~\#4059 by an
explicit second unitary-symmetry input.  Suppose a many-body unitary \(T\) and
its inverse \(T^{-1}\) satisfy
\[
  T^\dagger=T^{-1},\qquad T^{-1}T=1,\qquad T^{-1}\Phi=\Phi,
\]
and that the local spin components are conjugate at site \(x\):
\[
  T\hat S_x^{(2)}T^{-1}=\hat S_x^{(1)}.
\]
The PR~\#4057 unitary bridge gives \(E_1(x,\Phi)=E_2(x,\Phi)\).  Feeding this
into PR~\#4059, together with the lifted axis-swap invariance
\((\otimes_x U^{-1})\Phi=\Phi\), proves
\[
  E_1(x,\Phi)=E_2(x,\Phi)=E_3(x,\Phi)=\frac{N(N+2)}{12}.
\]
The concrete construction of \(T\), or a theorem saying that the AFM ground
state is fixed by such a symmetry, remains the open mathematical input.
Reference: Tasaki, Springer 2020, Problem~2.5.c, p.~43, and the
Theorem~2.4 symmetry context, pp.~43--44.

\textbf{Problem 2.5.c z-axis rotation input (PR~\#4061)}:
the next module instantiates the abstract second unitary of PR~\#4060 by a
concrete general spin-\(S\) rotation about axis \(3\).  Put
\[
  R_3(\theta)=\exp(-i\theta\hat S^{(3)}).
\]
Using the ladder identities
\[
  \hat S^+=\hat S^{(1)}+i\hat S^{(2)},\qquad
  \hat S^-=\hat S^{(1)}-i\hat S^{(2)}
\]
and the Cartan relations, the module proves
\[
  R_3(-\pi/2)\hat S^{(2)}R_3(\pi/2)=\hat S^{(1)}.
\]
Tensoring over all sites then gives
\[
  (\otimes_x R_3(-\pi/2))\hat S_x^{(2)}
  (\otimes_x R_3(\pi/2))=\hat S_x^{(1)}.
\]
Thus PR~\#4060 applies with this concrete rotation.  The remaining input is
that the AFM ground state is fixed by the inverse lifted z-axis rotation, along
with the existing lifted axis-swap invariance.  Reference: Tasaki, Springer
2020, Problem~2.5.c, p.~43, and the Theorem~2.4 symmetry context, pp.~43--44.

\textbf{Problem 2.5.c phase-invariant axis input (PR~\#4062)}:
the exact-invariance input in PRs~\#4057, \#4060, and \#4061 is stronger than
what one obtains from a one-dimensional ground-state line.  If
\[
  T A T^{-1}=B,\qquad T^{-1}\Phi=c\Phi,\qquad \overline c c=1,
\]
then
\[
\begin{aligned}
  \langle\Phi,B^2\Phi\rangle
    &= \langle\Phi,T A^2 T^{-1}\Phi\rangle \\
    &= c\,\langle\Phi,T A^2\Phi\rangle \\
    &= c\,\langle T^\dagger\Phi,A^2\Phi\rangle \\
    &= c\overline c\,\langle\Phi,A^2\Phi\rangle
     = \langle\Phi,A^2\Phi\rangle .
\end{aligned}
\]
The new wrapper specializes this to \(A=\hat S_x^{(2)}\),
\(B=\hat S_x^{(1)}\), and then to the concrete z-axis rotation of PR~\#4061.
The remaining Problem~2.5.c input is therefore a ray-invariance theorem for the
AFM ground state under the lifted rotations.  Reference: Tasaki, Springer 2020,
Problem~2.5.c, p.~43, and the Theorem~2.4 symmetry context, pp.~43--44.

\textbf{Problem 2.5.c one-dimensional eigenspace phase bridge (PR~\#4063)}:
the new file \texttt{Quantum/SpinS/Problem25cEigenspacePhaseBridge.lean}
formalizes the abstract linear-algebra source of the ray-invariance hypothesis.
If the Hamiltonian eigenspace at \(\mu\) has \(\dim \le 1\), a non-zero
eigenvector \(\Phi\) spans that eigenspace.  Therefore, whenever a symmetry
\(U\) commutes with \(H\), the vector \(U\Phi\) is again a \(\mu\)-eigenvector
and hence
\[
  U\Phi=c\Phi
\]
for some \(c\in\mathbb C\).  If, in addition, \(U^\dagger U=1\) and
\(\langle\Phi,\Phi\rangle=1\), then
\[
  1=\langle U\Phi,U\Phi\rangle
   =\langle c\Phi,c\Phi\rangle
   =\overline c c,
\]
so the scalar is exactly the unit-modulus phase required by PR~\#4062.  The
remaining concrete input is to show that the lifted rotations commute with the
AFM Hamiltonian.  Reference: Tasaki, Springer 2020, Problem~2.5.c, p.~43, and
the Theorem~2.4 symmetry context, pp.~43--44.

\textbf{Problem 2.5.c lifted z-axis rotation commutation (PR~\#4064)}:
the new file
\texttt{Quantum/SpinS/Problem25cZAxisRotationCommutation.lean} supplies the
concrete Hamiltonian-side commutation input for the phase bridge.  The
single-site rotation of PR~\#4061 is diagonal:
\[
  \exp(-i\theta \hat S^{(3)})_{kk}
    = \exp\!\left(-i\theta\left(\frac N2-k\right)\right).
\]
Therefore the lifted tensor rotation has diagonal entry
\[
  \prod_{x\in V}
  \exp\!\left(-i\theta\left(\frac N2-\sigma_x\right)\right)
  =
  \exp\!\left(-i\theta\left(\frac{|V|N}{2}
    -\sum_{x\in V}\sigma_x\right)\right),
\]
which is exactly the diagonal entry of
\[
  \exp(-i\theta\,\hat S^{(3)}_{\mathrm{tot}}).
\]
This proves
\[
  \bigotimes_{x\in V}\exp(-i\theta \hat S^{(3)})
  =
  \exp(-i\theta\,\hat S^{(3)}_{\mathrm{tot}}).
\]
The existing SU(2)-invariance theorem gives
\[
  [\hat H_J,\hat S^{(3)}_{\mathrm{tot}}]=0,
\]
and the standard exponential commutation lemma then yields
\[
  \left[\hat H_J,
    \bigotimes_{x\in V}\exp(-i\theta \hat S^{(3)})\right]=0.
\]
This is the concrete commutation theorem needed to feed the lifted z-axis
rotation into the one-dimensional eigenspace phase bridge.  Reference:
Tasaki, Springer 2020, Problem~2.5.c, p.~43, and the Theorem~2.4 symmetry
context, pp.~43--44.

\textbf{Problem 2.5.c z-axis ground-state phase input (PR~\#4065)}:
the new file
\texttt{Quantum/SpinS/Problem25cZAxisGroundStatePhase.lean} connects the
abstract one-dimensional eigenspace phase bridge with the concrete z-axis
rotation commutation theorem.  First, the lifted z-axis rotation is unitary:
\[
  \left(\bigotimes_{x\in V}\exp(-i\theta \hat S^{(3)})\right)^\dagger
  \left(\bigotimes_{x\in V}\exp(-i\theta \hat S^{(3)})\right)=1.
\]
The proof uses the tensor adjoint formula, the single-site adjoint identity
\[
  \exp(-i\theta\hat S^{(3)})^\dagger
  = \exp(i\theta\hat S^{(3)}),
\]
and tensor functoriality.  If $\Phi$ is a normalized non-zero eigenvector of
$\hat H_J$ in a one-dimensional eigenspace at energy $\mu$, then the
commutation
\[
  [\hat H_J,\bigotimes_{x\in V}\exp(-i\theta \hat S^{(3)})]=0
\]
and the abstract bridge from PR~\#4063 give a phase $c$ such that
\[
  \left(\bigotimes_{x\in V}\exp(-i\theta \hat S^{(3)})\right)\Phi=c\Phi,
  \qquad \overline c c=1.
\]
At $\theta=\pi/2$, this supplies the exact phase input required by the
Problem~2.5.c squared-expectation wrapper from PR~\#4062.  The resulting
theorem still keeps the lifted axis-swap state invariance as an explicit
hypothesis; the next remaining state-side input is the corresponding
axis-swap ground-state invariance or ray-invariance.  Reference: Tasaki,
Springer 2020, Problem~2.5.c, p.~43, and the Theorem~2.4 symmetry context,
pp.~43--44.

\textbf{Problem 2.5.c axis-swap ground-state phase input (PR~\#4066)}:
the new file
\texttt{Quantum/SpinS/Problem25cAxisSwapGroundStatePhase.lean} removes the
last exact axis-swap state-invariance hypothesis from the previous wrapper.
At the SU(2) point $\lambda=1$, $D=0$, the axis-swapped Hamiltonian
\[
  \sum_{x,y}J_{x,y}\left(\hat S_x^{(1)}\hat S_y^{(1)}
    +\lambda\hat S_x^{(2)}\hat S_y^{(2)}
    +\hat S_x^{(3)}\hat S_y^{(3)}\right)
  +D\sum_x(\hat S_x^{(2)})^2
\]
also reduces to the isotropic Heisenberg Hamiltonian $\hat H_J$.  Combined
with the existing axis-swap gauge equivalence, this gives
\[
  \hat H_J\Theta^{-1}=\Theta^{-1}\hat H_J
\]
for the inverse lifted axis-swap unitary $\Theta^{-1}$.  Since
$(\Theta^{-1})^\dagger\Theta^{-1}=1$, the one-dimensional eigenspace phase
bridge from PR~\#4063 gives a unit-modulus phase $c_{\mathrm{swap}}$ with
\[
  \Theta^{-1}\Phi=c_{\mathrm{swap}}\Phi,\qquad
  \overline{c_{\mathrm{swap}}}c_{\mathrm{swap}}=1.
\]
The phase-invariant expectation bridge then gives the axis-swap equality
$E_3=E_2$.  Together with the z-axis phase equality $E_1=E_2$ from
PR~\#4065 and the single-site Casimir identity, a normalized non-zero vector
in a one-dimensional Heisenberg eigenspace satisfies
\[
  E_1=E_2=E_3=\frac{N(N+2)}{12}.
\]
Reference: Tasaki, Springer 2020, Problem~2.5.c, p.~43, and the
Theorem~2.4 symmetry context, pp.~43--44.

\textbf{Problem 2.5.c MLM ground-state wrapper (PR~\#4067)}:
the new file
\texttt{Quantum/SpinS/Problem25cMLMGroundStateWrapper.lean} connects the
conditional one-dimensional-eigenspace theorem above to the balanced
Marshall--Lieb--Mattis/SU(2) endpoint.  Under the balanced bipartite
Theorem~2.3 hypotheses, the existing balanced uniqueness endpoint identifies
the Heisenberg Hermitian minimum with a common energy $\mu$ and proves
\[
  \dim_{\mathbb C}\ker(\hat H_J-\mu)\le 1.
\]
Therefore any normalized non-zero ground-state eigenvector
\[
  \hat H_J\Phi=E_{\min}\Phi,\qquad \langle\Phi,\Phi\rangle=1
\]
satisfies the hypotheses of the PR~\#4066 phase wrapper.  The z-axis rotation
phase gives $E_1=E_2$, the axis-swap phase gives $E_3=E_2$, and the
single-site Casimir identity yields
\[
  E_1=E_2=E_3=\frac{N(N+2)}{12}.
\]
Reference: Tasaki, Springer 2020, Problem~2.5.c, p.~43, and
Theorems~2.3--2.4, pp.~42--44.

\textbf{Problem 2.5.c balanced structural wrapper (PR~\#4081)}:
the file
\texttt{Quantum/SpinS/Problem25cBalancedStructuralWrapper.lean} removes the
explicit Theorem~2.3 witness from the preceding MLM ground-state wrapper.
In the balanced case \(|A|=|\bar A|\), the orientation
\(|\bar A|\le |A|\) is immediate, and the hypotheses
\(|\bar A|\ge 1\) and \(N\ge 1\) give
\[
  0 < |\bar A|\,N/2 .
\]
These are exactly the extra inputs needed by the structural theorem
\texttt{tasaki\_2\_5\_theorem\_2\_3\_of\_bipartiteCompletePositive}, which
supplies the \texttt{hT23} witness for the general real symmetric
non-negative bipartite coupling \(J\).  Feeding that witness to
\texttt{singleSiteSpinSquareExpectationS\_all\_axes\_eq\_of\_tasaki23\_balanced\_MLM\_groundState}
proves that every normalized non-zero Heisenberg ground-state eigenvector
satisfies
\[
  E_1=E_2=E_3=\frac{N(N+2)}{12}
\]
under the standard balanced bipartite antiferromagnetic hypotheses.
Reference: Tasaki, Springer 2020, Problem~2.5.c, p.~43, and
Theorems~2.3--2.4, pp.~42--44.

\textbf{Problem 2.5.d correlation sign bridge (PR~\#4068)}:
the new file \texttt{Quantum/SpinS/Problem25dCorrelationSignBridge.lean}
starts Tasaki Problem~2.5.d, p.~40, following the sign conversion in the
solution on p.~498, equations (S.22)--(S.23).  It defines
\[
  C_{xy}(\Phi)=
  \langle\Phi,\hat S_x\cdot\hat S_y\,\Phi\rangle
\]
as \texttt{twoSpinCorrelationS}, and the bipartite Marshall-gauge sign
\(\epsilon_A(x)=1\) on \(A\), \(\epsilon_A(x)=-1\) off \(A\), as
\texttt{bipartiteGaugeSign}.  The theorem
\texttt{twoSpinCorrelationS\_sign\_cases\_of\_bipartite\_signed\_re\_pos}
proves the elementary last step:
if \(\operatorname{Re}(\epsilon_A(x)\epsilon_A(y)C_{xy}(\Phi))>0\), then
\(\operatorname{Re} C_{xy}(\Phi)>0\) for same-sublattice pairs and
\(\operatorname{Re} C_{xy}(\Phi)<0\) for cross-sublattice pairs.  The remaining
Problem~2.5.d input is the Marshall-positive coefficient expansion that
supplies the signed positivity hypothesis.

\textbf{Problem 2.5.d ladder positivity expansion (PR~\#4069)}:
the new file \texttt{Quantum/SpinS/Problem25dLadderPositivity.lean}
packages the finite-sum positivity step in Tasaki's solution of
Problem~2.5.d, p.~40, solution p.~498, equations (S.22)--(S.23).  It defines
\[
  C^{+-}_{xy}(\Phi)=
  \langle\Phi,\hat S_x^+\hat S_y^-\,\Phi\rangle
\]
as \texttt{twoSpinPlusMinusCorrelationS}.  The theorem
\texttt{signedExpectation\_re\_eq\_sum} expands a signed expectation
\[
  \operatorname{Re}\bigl(g\langle \Phi,O\Phi\rangle\bigr)
\]
into the configuration-basis double sum of the signed terms, and
\texttt{signedExpectation\_re\_pos\_of\_term\_re\_nonneg\_exists\_pos}
proves the elementary positivity principle: all terms non-negative and one
term strictly positive imply a strictly positive signed expectation.  The
coefficient form
\texttt{signedExpectation\_re\_pos\_of\_positive\_coefficients} makes the
strictly positive factors \(c_\sigma c_\tau\) in (S.23) explicit.  The
Problem~2.5.d wrapper
\texttt{twoSpinPlusMinusCorrelationS\_bipartite\_signed\_re\_pos\_of\_marshall\_coefficients}
specializes this to the bipartite gauge sign and the lifted
\(\hat S_x^+\hat S_y^-\) operator.  The remaining input is the concrete
Marshall-gauge matrix-entry non-negativity and strict witness for the
ground-state vector.

\textbf{Problem 2.5.d ladder entry sign bridge (PR~\#4071)}:
the new file \texttt{Quantum/SpinS/Problem25dLadderEntrySign.lean}
supplies the concrete matrix-entry sign input required by the PR~\#4069
finite-sum wrapper, again following Tasaki, Springer 2020,
Problem~2.5.d, p.~40, and solution p.~498, equations (S.22)--(S.23).
For a cross-sublattice pair \(A(x)\ne A(y)\), the theorem
\texttt{twoSpinPlusMinus\_ladder\_signed\_entry\_re\_nonneg\_of\_bipartite\_ne}
proves
\[
  0\le
  \operatorname{Re}\!\left(
    \epsilon_A(x)\epsilon_A(y)
    m_A(\sigma)
    \langle\sigma,\hat S_x^+\hat S_y^- \tau\rangle
    m_A(\tau)
  \right),
\]
where \(m_A(\sigma)\) is the spin-\(S\) Marshall sign.  On a non-zero
ladder matrix entry the configurations agree away from \(x,y\), the local
indices differ by one at both sites, and the two odd local parity changes
make the Marshall sign product equal to \(-1\).  Since
\(\epsilon_A(x)\epsilon_A(y)=-1\) for \(A(x)\ne A(y)\), the two signs
cancel and the expression reduces to the bare non-negative real ladder
entry; all remaining cases are zero by the local raising/lowering action.

The theorem
\texttt{exists\_twoSpinPlusMinus\_ladder\_signed\_entry\_re\_pos\_of\_bipartite\_ne}
chooses explicit configurations for \(N\ge 1\), with
\(\sigma_x=0,\tau_x=1,\sigma_y=1,\tau_y=0\), and all other sites equal,
giving a strictly positive signed entry.  Combining the non-negativity
and witness with PR~\#4069 gives
\texttt{twoSpinPlusMinusCorrelationS\_bipartite\_signed\_re\_pos\_of\_marshall\_coefficients\_cross}:
strictly positive Marshall coefficients imply a positive
\(\epsilon_A(x)\epsilon_A(y)\)-signed \(\hat S_x^+\hat S_y^-\)
correlation on cross-sublattice pairs.  The remaining Problem~2.5.d
input is the SU(2) / ladder reduction that transfers this signed ladder
positivity to the full two-spin dot-product correlation.

\textbf{Problem 2.5.d ladder-to-dot reduction (PR~\#4072)}.
The file \path{Quantum/SpinS/Problem25dLadderDotReduction.lean}
packages the algebraic reduction from the positive signed ladder correlation
to positive signed dot-product correlation.  It introduces the component
definitions
\[
  \texttt{twoSpinMinusPlusCorrelationS}
  \quad\text{and}\quad
  \texttt{twoSpinZZCorrelationS},
\]
and lifts
\[
  \hat S_x\cdot\hat S_y
  =
  \frac12(\hat S_x^+\hat S_y^-+\hat S_x^-\hat S_y^+)
  + \hat S_x^{(3)}\hat S_y^{(3)}
\]
to the expectation-level theorem
\texttt{twoSpinCorrelationS\_eq\_ladder\_components}.  The signed form
\texttt{signed\_twoSpinCorrelationS\_re\_pos\_of\_ladder\_component\_equalities}
states the SU(2)-symmetric sign transfer used in Tasaki's solution: if
\[
  \operatorname{Re}(g C^{-+}_{xy})=\operatorname{Re}(g C^{+-}_{xy}),
  \qquad
  \operatorname{Re}(g C^{33}_{xy})
    =\frac12\operatorname{Re}(g C^{+-}_{xy}),
\]
and \(\operatorname{Re}(g C^{+-}_{xy})>0\), then
\[
  \operatorname{Re}(g C_{xy})
  =
  \frac32\operatorname{Re}(g C^{+-}_{xy})>0.
\]
The wrapper
\texttt{twoSpinCorrelationS\_bipartite\_signed\_re\_pos\_of\_marshall\_cross\_components}
combines this reduction with the PR~\#4071 Marshall-coefficient signed ladder
positivity on cross-sublattice pairs.  The remaining input is to discharge the
two component equalities from the ground-state SU(2) phase invariance.

\textbf{Problem 2.5.d ladder adjoint equality (PR~\#4073)}.
The file \path{Quantum/SpinS/Problem25dLadderAdjointEquality.lean}
discharges the first component equality left by PR~\#4072.  For distinct
sites, the two-site ladder operator satisfies
\[
  (\hat S_x^+\hat S_y^-)^\dagger=\hat S_x^-\hat S_y^+,
\]
formalized as \texttt{twoSpinPlusMinus\_ladder\_conjTranspose}.  The generic
bridge \path{dotProduct_star_conjTranspose_mulVec_eq_star} identifies
the expectation of \(M^\dagger\) with the complex conjugate of the expectation
of \(M\), so
\[
  C^{-+}_{xy}=\overline{C^{+-}_{xy}},
\]
and, because the cross-sublattice bipartite gauge product is the real scalar
\(-1\), the signed real parts agree:
\[
  \operatorname{Re}(g C^{-+}_{xy})
  =\operatorname{Re}(g C^{+-}_{xy}).
\]
The wrapper
\texttt{twoSpinCorrelationS\_bipartite\_signed\_re\_pos\_of\_marshall\_cross\_ladderAdjoint}
combines this equality with the PR~\#4071 ladder positivity and the PR~\#4072
ladder-to-dot reduction.  The only remaining component-equality input for the
Problem~2.5.d sign theorem is now the longitudinal
\(\operatorname{Re}(gC^{33}_{xy})=\frac12\operatorname{Re}(gC^{+-}_{xy})\)
identity from ground-state SU(2) phase invariance.

\textbf{Problem 2.5.d longitudinal component equality (PR~\#4074)}.
The file \path{Quantum/SpinS/Problem25dLongitudinalComponentEquality.lean}
discharges the second component equality left by PR~\#4072 under the same
phase-invariance interface used in Problem~2.5.c.  The generic bridge
\path{twoSpinProductCorrelationS_eq_of_conj_phaseInvariant} says that if a
unitary conjugates one two-site product to another and the state is fixed by
the inverse unitary up to a unit-modulus phase, then the two expectations are
equal.  The bridge
\path{twoSpinProductCorrelationS_spinSOp3_eq_twoSpinZZCorrelationS} identifies
the axis-3 product expectation with the existing longitudinal correlation
\path{twoSpinZZCorrelationS}.  Applying the phase-invariant bridge to the
axis-swap unitary gives
\[
  C^{33}_{xy}=C^{22}_{xy},
\]
while the lifted \(z\)-axis rotation gives
\[
  C^{11}_{xy}=C^{22}_{xy}.
\]
The two-site transverse ladder identity
\[
  C^{11}_{xy}+C^{22}_{xy}
  =\frac12(C^{+-}_{xy}+C^{-+}_{xy})
\]
together with PR~\#4073's signed real-part equality for
\(C^{-+}_{xy}\) and \(C^{+-}_{xy}\) yields
\[
  \operatorname{Re}(gC^{33}_{xy})
  =\frac12\operatorname{Re}(gC^{+-}_{xy}).
\]
The wrapper
\path{twoSpinCorrelationS_bipartite_signed_re_pos_of_marshall_cross_axis_phases}
therefore removes the last component-equality hypothesis from the PR~\#4072
dot-product positivity transfer, leaving only the ground-state phase-invariance
inputs to be connected to the balanced MLM endpoint.

\textbf{Problem 2.5.d ground-state phase wrapper (PR~\#4075)}.
The file \path{Quantum/SpinS/Problem25dGroundStatePhaseWrapper.lean}
connects the phase-invariance inputs left by PR~\#4074 to the
one-dimensional Heisenberg eigenspace phase bridges already used for
Problem~2.5.c.  If
\[
  \Phi(\sigma)=m_A(\sigma)c_\sigma,\qquad c_\sigma>0,
  \qquad \hat H_J\Phi=\mu\Phi,
\]
is normalized, non-zero, and lies in a Heisenberg eigenspace of complex
dimension at most one, then the existing phase bridges provide
\[
  U_{\mathrm{swap}}^{-1}\Phi=c_{\mathrm{swap}}\Phi,\qquad
  U_3(\pi/2)\Phi=c_{\mathrm{rot}}\Phi,
\]
with \(\overline{c_{\mathrm{swap}}}c_{\mathrm{swap}}=
\overline{c_{\mathrm{rot}}}c_{\mathrm{rot}}=1\).  Substituting these witnesses
into the PR~\#4074 axis-phase wrapper gives
\[
  \operatorname{Re}\{\epsilon_A(x)\epsilon_A(y)
    \langle\Phi,\hat S_x\cdot\hat S_y\Phi\rangle\}>0
\]
for cross-sublattice pairs \(A(x)\ne A(y)\).  The companion balanced MLM
wrapper uses the Theorem~2.3/SU(2) endpoint to supply the one-dimensional
ground eigenspace at the Hermitian minimum.  This wrapper is for a full-space
coefficient function; the sector-supported Perron--Frobenius ground-state
shape is handled by the next wrapper.  Reference: Tasaki, Springer 2020,
Problem~2.5.d, p.~40, solution pp.~498--499, equations (S.22)--(S.23), with
Theorems~2.3--2.4, pp.~42--44.

\textbf{Problem 2.5.d sector-supported wrapper (PR~\#4076)}.
The file \path{Quantum/SpinS/Problem25dSectorSupportedWrapper.lean}
adapts the PR~\#4069--\#4075 positivity chain to the actual
Perron--Frobenius ground-state shape.  Let \(M\) be a magnetization sector and
let
\[
  \Phi=\iota_M(\sigma\mapsto m_A(\sigma)c_\sigma),
  \qquad c_\sigma>0\quad(\sigma\in M),
\]
where \(\iota_M\) is zero-extension from the sector.  In the finite-basis
expansion of
\(\langle\Phi,\hat S_x^+\hat S_y^-\Phi\rangle\), every term with an
off-sector configuration vanishes.  Every on-sector term is
\(c_\sigma c_\tau\) times the signed matrix-entry expression already
controlled by PR~\#4071.  Thus a sector-local strict ladder-entry witness
supplies the same positive signed ladder correlation used by PR~\#4072.
Combining this with the PR~\#4073 adjoint equality and PR~\#4074 longitudinal
axis-phase equality gives the sector-supported signed dot-product positivity
wrapper.  Finally, a normalized non-zero sector-supported eigenvector in a
one-dimensional Heisenberg eigenspace obtains the same axis-swap and
z-rotation phases as in PR~\#4075.  Reference: Tasaki, Springer 2020,
Problem~2.5.d, p.~40, solution pp.~498--499, equations (S.22)--(S.23), with
Theorems~2.3--2.4, pp.~42--44.

\textbf{Problem 2.5.d balanced-sector ladder witness (PR~\#4077)}.
The file \path{Quantum/SpinS/Problem25dBalancedSectorWitness.lean}
supplies the sector-local strict ladder entry required by PR~\#4076 at the
balanced sector \(M_0=|A|N\).  First,
\path{twoSpinPlusMinus_ladder_signed_entry_re_pos_of_off_two_site_raise_lower}
isolates the strict-entry calculation: for \(A(x)\ne A(y)\), agreement away
from \(x,y\) together with
\[
  \sigma_x+1=\tau_x,\qquad \tau_y+1=\sigma_y
\]
makes the bare \(\hat S_x^+\hat S_y^-\) matrix entry strictly positive, while
the Marshall sign and bipartite gauge sign cancel to \(+1\).

For the sector witness, take
\(\eta=\texttt{neelConfigOfS}(\neg A,N)\).  Then
\[
  \sum_z \eta_z = |A|N
\]
by \path{magSumS_neelConfigOfS_complement}.  If \(x\notin A\) and \(y\in A\),
use \(\sigma=\eta\) and update \(\tau\) by setting
\(\tau_x=1,\tau_y=N-1\).  If \(x\in A\) and \(y\notin A\), reverse the
orientation: use \(\tau=\eta\) and update \(\sigma_x=N-1,\sigma_y=1\).
The two-site sum identity \path{magSumS_configUpdateTwo_eq} keeps both
configurations in \(M_0\), and the raise/lower equalities above feed the
strict-entry lemma.  The theorem
\path{exists_twoSpinPlusMinus_ladder_signed_entry_re_pos_of_bipartite_ne_balanced_sector}
packages this witness, and
\path{twoSpinCorrelationS_bipartite_signed_re_pos_of_marshall_balanced_sector_cross_eigenphase}
feeds it into the sector-supported eigenspace-phase wrapper so the balanced
sector no longer needs an explicit strict-entry hypothesis.  Reference:
Tasaki, Springer 2020, Problem~2.5.d, p.~40, solution pp.~498--499, equations
(S.22)--(S.23), with Theorems~2.3--2.4, pp.~42--44.

\textbf{Problem 2.5.d balanced Perron--Frobenius endpoint (PR~\#4078)}.
The file \path{Quantum/SpinS/Problem25dBalancedPFEndpoint.lean}
connects the balanced-sector PF vector produced by the Theorem~2.3 common
energy endpoint to the sector-supported Problem~2.5.d positivity chain.
The Theorem~2.3 data return the embedded vector in the real-sign form
\[
  \Psi_v=\iota_{M_0}\bigl(\sigma\mapsto
    ((m_A(\sigma))_{\mathrm{re}}v_\sigma:\mathbb R):\mathbb C\bigr).
\]
Since \path{marshallSignS_eq_ofReal_re} says that the Marshall sign is the
complex embedding of its real part, extensionality of
\path{magSectorEmbedding} identifies this vector with
\[
  \Phi_v=\iota_{M_0}\bigl(\sigma\mapsto m_A(\sigma)v_\sigma\bigr),
  \qquad v_\sigma>0.
\]
The new non-normalized phase bridge
\path{phase_unit_of_unitary_mulVec_eq_smul_of_ne_zero} avoids introducing a
separate normalized PF vector: if a unitary \(U\) satisfies
\(U\Phi=c\Phi\) and \(\Phi\ne0\), then norm preservation gives
\((\bar c c)\langle\Phi,\Phi\rangle=\langle\Phi,\Phi\rangle\), and the
non-zero norm cancels to \(\bar c c=1\).  Together with the existing
one-dimensional eigenspace scalar-multiple bridge, this supplies the
axis-swap and \(z\)-rotation phases consumed by the PR~\#4076
sector-supported axis-phase wrapper.  The endpoint
\path{twoSpinCorrelationS_bipartite_signed_re_pos_of_tasaki23_balanced_pf_cross}
returns the Hermitian minimum energy, the strictly positive balanced
coefficients, the non-zero embedded eigenvector, full eigenspace
\(\mathrm{finrank}\le1\), and signed cross-sublattice two-spin correlation
positivity for this concrete PF vector.  Reference: Tasaki, Springer 2020,
Problem~2.5.d, p.~40, solution pp.~498--499, equations (S.22)--(S.23), with
Theorems~2.3--2.4, pp.~42--44.

\textbf{Problem 2.5.d balanced Perron--Frobenius sign cases (PR~\#4079)}.
The file \path{Quantum/SpinS/Problem25dBalancedPFSignCases.lean} applies the
final bipartite-gauge sign conversion from PR~\#4068 to the concrete balanced
PF endpoint from PR~\#4078.  The new theorem
\path{twoSpinCorrelationS_sign_cases_of_tasaki23_balanced_pf_cross} preserves
the output package from
\path{twoSpinCorrelationS_bipartite_signed_re_pos_of_tasaki23_balanced_pf_cross}:
the Hermitian minimum, the strictly positive balanced sector coefficients, the
non-zero sector-embedded eigenvector, and the full eigenspace
\(\mathrm{finrank}\le1\).  Its last field replaces the signed positivity
hypothesis by the four Boolean conclusions of
\path{twoSpinCorrelationS_sign_cases_of_bipartite_signed_re_pos}: positive
real two-spin correlation on the same Boolean sublattice and negative real
two-spin correlation on opposite Boolean sublattices, after the corresponding
Boolean assumptions are supplied.  Reference: Tasaki, Springer 2020,
Problem~2.5.d, p.~40, solution pp.~498--499, equations (S.22)--(S.23), with
Theorems~2.3--2.4, pp.~42--44.

\textbf{Problem 2.5.d balanced Perron--Frobenius cross sign (PR~\#4080)}.
The file \path{Quantum/SpinS/Problem25dBalancedPFCrossSign.lean} extracts the
non-vacuous cross-sublattice conclusion from the Boolean sign cases of
PR~\#4079.  The theorem
\path{twoSpinCorrelationS_re_neg_of_tasaki23_balanced_pf_cross} assumes
\(A(x)\ne A(y)\), obtains the balanced PF package from
\path{twoSpinCorrelationS_sign_cases_of_tasaki23_balanced_pf_cross}, and
performs a Boolean case split on \(A(x)\).  If \(A(x)=\mathrm{true}\), then
the cross hypothesis forces \(A(y)=\mathrm{false}\); otherwise
\(A(x)=\mathrm{false}\) and it forces \(A(y)=\mathrm{true}\).  In both cases
the appropriate cross-sublattice field of the PR~\#4079 theorem gives
\[
  \operatorname{Re}\langle \Phi_v,
    \widehat{\mathbf S}_x\cdot\widehat{\mathbf S}_y\Phi_v\rangle < 0 .
\]
The Hermitian minimum, strictly positive balanced coefficients, non-zero
embedded eigenvector, and full eigenspace \(\mathrm{finrank}\le1\) are
preserved unchanged.  Reference: Tasaki, Springer 2020, Problem~2.5.d, p.~40,
solution pp.~498--499, equations (S.22)--(S.23), with Theorems~2.3--2.4,
pp.~42--44.

\textbf{Spin-`S` $\leftrightarrow$ spin-`1/2` bridge at $N = 1$
(PRs~\#922 + \#923)}: the new file
\texttt{Quantum/SpinS/SpinHalfSpecialization.lean} proves
\texttt{spinSOp\{Plus, Minus, 1, 2, 3\} 1 = spinHalfOp\{Plus, Minus, 1, 2, 3\}}
by entrywise \texttt{ext + fin\_cases} on $\textsf{Fin}\,2$. At
$N = 1$ ($S = 1/2$), each generic spin-`S` operator reduces to the
corresponding half-Pauli matrix, providing a direct bridge between
the generic and spin-`1/2` formalisms.

\textbf{Single-site $\hat{S}^{(3)}_x$ expectations on $|c..c\rangle$
(PR~\#925)}: the new file
\texttt{Quantum/SpinS/SingleSiteZExpectation.lean} proves
$\hat{S}^{(3)}_x|c..c\rangle = (N/2 - c)\,|c..c\rangle$ and the
corresponding expectation values $\langle\hat{S}^{(3)}_x\rangle = N/2 - c$,
$\langle(\hat{S}^{(3)}_x)^2\rangle = (N/2-c)^2$. For $c=0$ this gives
$\langle\hat{S}^{(3)}_x\rangle = S$ and $\langle(\hat{S}^{(3)}_x)^2\rangle = S^2$
(zero variance about the mean $S$).

\textbf{Single-site xy-plane Casimir expectation (PR~\#926)}: the
new file \texttt{Quantum/SpinS/SingleSiteXYExpectation.lean}
combines \#920 (universal Casimir $\hat{S}_x\cdot\hat{S}_x = S(S+1)$)
with \#925 ($\langle(\hat{S}^{(3)}_x)^2\rangle = (N/2-c)^2$) to give
$\langle(\hat{S}^{(1)}_x)^2 + (\hat{S}^{(2)}_x)^2\rangle =
N(N+2)/4 - (N/2-c)^2$ on $|c..c\rangle$. For $c=0$ this is $S/2$.

\textbf{Transverse mean vanishes (PR~\#927)}: the new file
\texttt{Quantum/SpinS/SingleSiteTransverseMeanZero.lean} proves
$\langle\text{basisVecS}\,\sigma, \hat{S}^{(\alpha)}_x\,
\text{basisVecS}\,\sigma\rangle = 0$ for $\alpha = 1, 2$ and any
configuration $\sigma$ (transverse operators are purely off-diagonal,
\texttt{spinSOp\{1,2\}\_apply\_diag = 0}). Together with PR \#925,
the saturated-ferromagnet $|c..c\rangle$ has per-axis spin profile
$(0, 0, N/2-c)$ — only the z-component is non-zero.

\textbf{Singlet correlation sum (PR~\#929)}: the new file
\texttt{Quantum/SpinS/SingletCorrelationSum.lean} proves that for
any singlet state $\Phi$ ($\hat{S}_{\rm tot}^2\,\Phi = 0$),
$\sum_{x, y}\langle\Phi, \hat{S}_x\cdot\hat{S}_y\,\Phi\rangle = 0$.
Direct combination of \texttt{totalSpinSSquared\_eq\_sum\_spinSDot}
with \texttt{Matrix.sum\_mulVec} + \texttt{dotProduct\_sum} linearity
and \texttt{dotProduct\_zero}.

\textbf{Universal correlation = Casimir expectation (PR~\#930)}:
the new file \texttt{Quantum/SpinS/CorrelationSumCasimir.lean}
generalises \#929 to ANY state:
$\sum_{x, y}\langle\Phi, \hat{S}_x\cdot\hat{S}_y\,\Phi\rangle =
\langle\Phi, \hat{S}_{\rm tot}^2\,\Phi\rangle$. Specialised to
$|\sigma_\top\rangle$ and $|\sigma_\bot\rangle$ gives the Casimir
eigenvalue $m_{\max}(m_{\max}+1)$.

\textbf{Eigenvector correlation sum (PR~\#931)}: the new file
\texttt{Quantum/SpinS/CorrelationEigenvector.lean} parameterises
the result of \#930 by the eigenvalue: for $\hat{S}_{\rm tot}^2\Phi
= \lambda\,\Phi$, $\sum_{x, y}\langle\Phi, \hat{S}_x\cdot\hat{S}_y\,\Phi\rangle
= \lambda\,\langle\Phi, \Phi\rangle$. For normalized $\Phi$, the
sum equals $\lambda$.

\textbf{Off-diagonal correlation sum (PR~\#933)}: the new file
\texttt{Quantum/SpinS/CorrelationOffDiagonal.lean} introduces the
universal diagonal sum
$\sum_x\langle\Phi, \hat{S}_x\cdot\hat{S}_x\,\Phi\rangle =
|V|\cdot S(S+1)\cdot\langle\Phi,\Phi\rangle$ (from \#920) and
subtracts from \#931 to give: for $\hat{S}_{\rm tot}^2\Phi =
\lambda\,\Phi$, $\sum_{x\neq y}\langle\Phi, \hat{S}_x\cdot\hat{S}_y\,\Phi\rangle
= (\lambda - |V|\cdot S(S+1))\cdot\langle\Phi,\Phi\rangle$.

\textbf{Saturated off-diagonal correlation sum (PR~\#934)}: the
new file
\texttt{Quantum/SpinS/SaturatedOffDiagonalCorrelation.lean}
specialises \#933 to $|\sigma_{\top/\bot}\rangle$ using \#882
(Casimir at $k=0$) + \#913 (normalization), yielding
$\sum_{x\neq y}\langle|\sigma_\top\rangle,
\hat{S}_x\cdot\hat{S}_y\,|\sigma_\top\rangle\rangle =
m_{\max}(m_{\max}+1) - |V|\cdot S(S+1) = N^2\,|V|(|V|-1)/4$.

\textbf{Fermion $\leftrightarrow$ spin-`S` bridge at $N = 1$
(PR~\#936)}: the new file
\texttt{Fermion/SpinSBridge.lean} composes the existing
\texttt{fermionAnnihilation\_eq\_spinHalfOpPlus} (Mode.lean) with
\#922 \texttt{spinSOpPlus\_one\_eq\_spinHalfOpPlus} to give
\texttt{fermionAnnihilation\_eq\_spinSOpPlus\_one} (and the dual
for creation). Direct identification of single-mode fermion
operators with generic spin-`S` operators at $N = 1$.

\textbf{Fermion number as $1/2 - \hat{S}^{(3)}$ (PR~\#937)}: the
new file \texttt{Fermion/NumberSpinSBridge.lean} proves
\texttt{fermionNumber\_eq\_half\_smul\_one\_sub\_spinSOp3\_one}:
$n = (1/2)\,I - \hat{S}^{(3)}$, the standard physics identification
$n = (I - \sigma^z)/2$ lifted to generic spin-`S` at $N = 1$.

\textbf{Per-pair correlation $\hat{S}_x\cdot\hat{S}_y\,|\sigma_\top\rangle
= (N^2/4)\,|\sigma_\top\rangle$ for $x\neq y$ (PR~\#939)}: the new
file \texttt{Quantum/SpinS/SpinSDotAllAlignedZero.lean} computes the
per-pair eigenvalue directly via \texttt{spinSDot\_eq\_plus\_minus}:
the $(\pm,\mp)$ ladder cross-terms vanish on $|\sigma_\top\rangle$
(via \#877 highest-weight annihilation
\texttt{onSiteS\_spinSOpPlus\_mulVec\_allAlignedStateS\_zero})
and an entrywise off-site agreement analysis killing the residual
$(+,-)$ contribution; the $(3)(3)$ term gives $(N/2)^2$.

\textbf{Symmetric raising-side per-pair eigenvalue (PR~\#940)}: the
new file \texttt{Quantum/SpinS/SpinSDotAllAlignedLast.lean} mirrors
\#939 for $|\sigma_\bot\rangle$ via $\hat{S}^+ \leftrightarrow
\hat{S}^-$, $0 \leftrightarrow N$ swap. The $(3)(3)$ coefficient
$(-N/2)^2$ also gives $N^2/4$.

\textbf{Per-pair correlation expectation (PR~\#941)}: the new file
\texttt{Quantum/SpinS/PerPairCorrelationExpectation.lean} packages
\#939/\#940 with the normalization (\#913):
$\langle|\sigma_{\top/\bot}\rangle, \hat{S}_x\cdot\hat{S}_y\,
|\sigma_{\top/\bot}\rangle\rangle = N^2/4 = S^2$ for $x \neq y$.
The total off-diagonal sum then equals $|V|\cdot(|V|-1)\cdot S^2$,
matching \#934.

\textbf{Heisenberg expectation on saturated states (PR~\#943)}: the
new file \texttt{Quantum/SpinS/SaturatedHeisenbergExpectation.lean}
gives the H-eigenvalue expectation:
$\langle|\sigma_\top\rangle, H\,|\sigma_\top\rangle\rangle =
\texttt{saturatedFerromagnetEigenvalueS}\,J\,N$ (defined as
$H(\sigma_\top, \sigma_\top)$ in \#899). Symmetric for
$|\sigma_\bot\rangle$, with eigenvalue $H(\sigma_\bot, \sigma_\bot)$.

\textbf{Uniform-J Heisenberg equals Casimir (PR~\#945)}: the new
file \texttt{Quantum/SpinS/HeisenbergUniformCasimir.lean} proves
\texttt{heisenbergHamiltonianS (fun \_ \_ => 1) N = totalSpinSSquared V N}.
Direct from the definition $H = \sum_{x, y} J\,\hat{S}_x\cdot\hat{S}_y$
+ \texttt{totalSpinSSquared\_eq\_sum\_spinSDot}.

\textbf{H-diagonal symmetry (PR~\#946)}: the new file
\texttt{Quantum/SpinS/SaturatedHeisenbergSymmetric.lean} shows
$H(\sigma_\bot, \sigma_\bot) =
\texttt{saturatedFerromagnetEigenvalueS}\,J\,N$. Both \#875 and
\#876 produce the same explicit eigenvalue $E$, and uniqueness of
the eigenvalue on a non-zero eigenvector (via
\texttt{allAlignedStateS\_ne\_zero}) extracts
$E = H(\sigma_\top, \sigma_\top) = H(\sigma_\bot, \sigma_\bot)$.

\textbf{Symmetric H expectation cleanup (PR~\#948)}: the new file
\texttt{Quantum/SpinS/SaturatedHeisenbergExpectationClean.lean}
chains \#943 and \#946 to give the clean form
$\langle|\sigma_\bot\rangle, H\,|\sigma_\bot\rangle\rangle =
\texttt{saturatedFerromagnetEigenvalueS}\,J\,N$, matching the
$|\sigma_\top\rangle$ statement.

\textbf{Uniform-J saturated eigenvalue equals Casimir (PR~\#949)}:
the new file
\texttt{Quantum/SpinS/SaturatedHeisenbergUniformEigenvalue.lean}
proves
$\texttt{saturatedFerromagnetEigenvalueS}\,(\lambda\_\_, 1)\,N =
\texttt{saturatedFerromagnetCasimirEigenvalueS}\,V\,N =
m_{\max}(m_{\max}+1)$. Direct from \#945 ($H_{\rm uniform} =
\hat{S}_{\rm tot}^2$) plus the diagonal entry of $\hat{S}_{\rm tot}^2$
at $\sigma_\top$.

\textbf{Explicit form of saturated FM eigenvalue (PR~\#951)}: the
new file
\texttt{Quantum/SpinS/SaturatedEigenvalueExplicit.lean} extracts
the explicit combinatorial form:
$\texttt{saturatedFerromagnetEigenvalueS}\,J\,N = \sum_x \sum_y
J(x,y) \cdot
(\textsf{if}\ x=y\ \textsf{then}\ N(N+2)/4\ \textsf{else}\ (N/2)^2)$.
Direct application of \#875 + \#899 uniqueness on $|\sigma_\top\rangle \neq 0$.
In the Tasaki §2.5 Theorem~2.3 chain, the same explicit formula gives
\(\texttt{saturatedFerromagnetEigenvalueS\_im\_zero}\) and
\(\texttt{saturatedFerromagnetEigenvalueS\_exists\_real}\): if every coupling
\(J(x,y)\) has zero imaginary part, each summand has zero imaginary part, so
the saturated ferromagnetic energy is represented by its real part as a real
scalar.

\textbf{Explicit uniform-J eigenvalue simplification (PR~\#953)}:
the new file
\texttt{Quantum/SpinS/SaturatedExplicitUniformSimp.lean} combines
\#951 (explicit form) with \#949 (uniform = Casimir) to give
the closed-form simplification
$\sum_x \sum_y (\textsf{if}\ x=y\ \textsf{then}\ N(N+2)/4\
\textsf{else}\ (N/2)^2) = m_{\max}(m_{\max}+1)$.

\textbf{Explicit H expectation on saturated states (PR~\#954)}: the
new file \texttt{Quantum/SpinS/HExpectationExplicit.lean} chains
\#943/\#948 with \#951 to give the explicit combinatorial form of
$\langle|\sigma_{\top/\bot}\rangle, H\,|\sigma_{\top/\bot}\rangle\rangle$
in terms of the $J(x,y)$ couplings.

\textbf{Distinct constant-spin states are different (PR~\#956)}:
the new file \texttt{Quantum/SpinS/AllAlignedStateDistinct.lean}
proves \texttt{allAlignedConfigS\_injective} (when $V$ is non-empty)
and the corresponding \texttt{allAlignedStateS\_ne\_of\_ne}.
Used as a building block for the LI of the all-aligned family.

\textbf{All-aligned states form a LI family (PR~\#957)}: the new
file \texttt{Quantum/SpinS/AllAlignedStateLI.lean} proves
\texttt{allAlignedStateS\_linearIndependent}: the
$(N+1)$-element family $\{|c..c\rangle : c \in \textsf{Fin}\,(N+1)\}$
is linearly independent over $\mathbb{C}$. Direct application of
\texttt{Module.End.eigenvectors\_linearIndependent'} with the
injective $\hat{S}^z_{\rm tot}$ eigenvalue function
$c \mapsto |V|\cdot N/2 - |V|\cdot c$.

\textbf{All-aligned span finrank (PR~\#959)}: the new file
\texttt{Quantum/SpinS/AllAlignedStateSpan.lean} applies
\texttt{finrank\_span\_eq\_card} to PR \#957's LI family,
yielding $\dim_\mathbb{C}\,\texttt{Submodule.span}\,\mathbb{C}\,
(\textsf{Set.range}\,\texttt{allAlignedStateS}) = N + 1$.

\textbf{All-aligned states orthogonal (PR~\#960)}: the new file
\texttt{Quantum/SpinS/AllAlignedStateOrthogonal.lean} proves
$\langle|c_1..c_1\rangle, |c_2..c_2\rangle\rangle = 0$ for
$c_1 \neq c_2$ via \texttt{basisVecS\_inner\_of\_ne} (PR \#914) +
\texttt{allAlignedConfigS\_injective} (PR \#956).

\textbf{All-aligned magnetization subspace membership (PR~\#962)}:
the new file
\texttt{Quantum/SpinS/AllAlignedStateMagSubspace.lean} proves
\texttt{allAlignedStateS\_mem\_magSubspaceS}: for any constant
$c$, $|c..c\rangle \in \texttt{magSubspaceS}\,V\,N\,(|V|\cdot N/2
- |V|\cdot c)$. Generalises \#891's extremal cases to general $c$.

\textbf{Fermion identity $c\cdot c^\dagger = 1 - n$ (PR~\#963)}:
the new file \texttt{Fermion/AnnihilationCreationIdentity.lean}
proves the standard hole-occupation identity from
\texttt{fermionAnticomm\_self} ($c\,c^\dagger + c^\dagger\,c = 1$)
and \texttt{fermionNumber\_eq\_creation\_mul\_annihilation}
($n = c^\dagger\,c$).

\textbf{$c\,c^\dagger$ in spin-`S` form (PR~\#965)}: the new file
\texttt{Fermion/CCDaggerSpinSBridge.lean} chains \#937 and \#963
to give $c\cdot c^\dagger = (1/2)\,I + \hat{S}^{(3)}_1$.

\textbf{$c\,c^\dagger$ action on basis states (PR~\#966)}: the new
file \texttt{Fermion/CCDaggerAction.lean} proves
$(c\cdot c^\dagger) |0\rangle = |0\rangle$ and
$(c\cdot c^\dagger) |1\rangle = 0$, the vacuum/occupied eigenstate
identifications.

\textbf{Vacuum/occupied orthonormality (PR~\#968)}: the new file
\texttt{Fermion/StatesOrthonormal.lean} computes the standard
orthonormality of $|0\rangle, |1\rangle$ entrywise.

\textbf{Number expectations (PR~\#969)}: the new file
\texttt{Fermion/NumberExpectations.lean} gives $\langle 0| n |0\rangle = 0$
and $\langle 1| n |1\rangle = 1$.

\textbf{$c\,c^\dagger$ expectations (PR~\#971)}: the new file
\texttt{Fermion/CCDaggerExpectations.lean} gives
$\langle 0| c c^\dagger |0\rangle = 1$ and
$\langle 1| c c^\dagger |1\rangle = 0$.

\textbf{Projection sum (PR~\#972)}: the new file
\texttt{Fermion/ProjectionSum.lean} proves
$n + c\,c^\dagger = 1$, the resolution of identity in the
occupation basis.

\textbf{$c\,c^\dagger$ idempotent (PR~\#974)}: the new file
\texttt{Fermion/CCDaggerIdempotent.lean} shows
$(c\,c^\dagger)^2 = c\,c^\dagger$. From $c\,c^\dagger = 1 - n$
and $n^2 = n$: $(1-n)^2 = 1 - 2n + n^2 = 1 - n$.

\textbf{Projection orthogonality (PR~\#976)}: the new file
\texttt{Fermion/ProjectionsOrthogonal.lean} shows
$n \cdot (c\,c^\dagger) = 0$ and $(c\,c^\dagger) \cdot n = 0$ via
$c^2 = 0$ and $(c^\dagger)^2 = 0$.

\textbf{Projections commute (PR~\#978)}: trivially from \#976
both products vanish, hence \texttt{Commute n (c $\cdot$ c$^\dagger$)}.

\textbf{$c\,c^\dagger$ Hermitian (PR~\#980)}: $(c\,c^\dagger)^H =
(c^\dagger)^H \cdot c^H = c \cdot c^\dagger$. Direct from
\texttt{Matrix.conjTranspose\_mul} +
\texttt{fermionAnnihilation\_conjTranspose} +
\texttt{fermionCreation\_conjTranspose}.

\textbf{$n\cdot c = 0$ and $c\cdot n = c$ (PR~\#982)}: in the file
\texttt{Fermion/AnnihilationNumberIdentities.lean}.
$n\cdot c = (c^\dagger c)\,c = c^\dagger(c\,c) = c^\dagger\cdot 0 = 0$
via \texttt{fermionAnnihilation\_sq}.
$c\cdot n = c(c^\dagger c) = (c\,c^\dagger)\,c = (1-n)\,c
 = c - n\cdot c = c$
via \texttt{fermionAnticomm\_self} (so $c\,c^\dagger = 1 - n$) and
the previous identity.

\textbf{$c^\dagger\cdot n = 0$ and $n\cdot c^\dagger = c^\dagger$
(PR~\#984)}: companion to PR~\#982, in the file
\texttt{Fermion/CreationNumberIdentities.lean}.
$c^\dagger\cdot n = c^\dagger(c^\dagger c) = (c^\dagger)^2\,c
 = 0\cdot c = 0$
via \texttt{fermionCreation\_sq}.
$n\cdot c^\dagger = (c^\dagger c)\,c^\dagger = c^\dagger(c\,c^\dagger)
 = c^\dagger(1-n) = c^\dagger - c^\dagger\cdot n = c^\dagger$
via \texttt{fermionAnticomm\_self} and the previous identity.

\textbf{Partial-isometry relations $c\,c^\dagger c = c$ and
$c^\dagger c\,c^\dagger = c^\dagger$ (PR~\#986)}: in the file
\texttt{Fermion/PartialIsometry.lean}.
$c\,c^\dagger c = c\,(c^\dagger c) = c\,n = c$ by PR~\#982.
$c^\dagger c\,c^\dagger = (c^\dagger c)\,c^\dagger = n\,c^\dagger = c^\dagger$
by PR~\#984. These are the standard $S\,S^*S = S$ /
$S^*S\,S^* = S^*$ partial-isometry identities of operator algebra.

\textbf{Ladder commutators $[n,c] = -c$ and $[n,c^\dagger] = c^\dagger$
(PR~\#988)}: in the file
\texttt{Fermion/NumberLadderCommutators.lean}.
$[n,c] = n\cdot c - c\cdot n = 0 - c = -c$ by PR~\#982.
$[n,c^\dagger] = n\cdot c^\dagger - c^\dagger\cdot n = c^\dagger - 0
= c^\dagger$ by PR~\#984. These are the fermion analogues of the
bosonic number-ladder commutators $[N,a] = -a$, $[N,a^\dagger] = a^\dagger$,
expressing that the adjoint action of $n$ (commutation with $n$)
shifts the fermion number of $c$ down by one and $c^\dagger$ up by one.

\textbf{$c$--$c^\dagger$ commutator $[c, c^\dagger] = 1 - 2\,n$
(PR~\#989)}: in the file
\texttt{Fermion/CCDaggerCommutator.lean}.
$[c, c^\dagger] = c\,c^\dagger - c^\dagger c = (1 - n) - n = 1 - 2n$
via PR~\#963 ($c\,c^\dagger = 1 - n$) and
\texttt{fermionNumber\_eq\_creation\_mul\_annihilation}
($c^\dagger c = n$). Together with the anticommutator
$\{c, c^\dagger\} = 1$ (\texttt{fermionAnticomm\_self}) this encodes
the algebraic difference between fermions and bosons: the bosonic
commutator $[a, a^\dagger] = 1$ is number-independent, while the
fermion commutator $1 - 2n$ takes value $+1$ on the vacuum
($n = 0$) and $-1$ on the occupied state ($n = 1$).

\textbf{Ladder anticommutators $\{n,c\} = c$ and $\{n,c^\dagger\} =
c^\dagger$ (PR~\#990)}: in the file
\texttt{Fermion/NumberLadderAnticommutators.lean}. Duals of the
ladder commutators of PR~\#988.
$\{n,c\} = n\cdot c + c\cdot n = 0 + c = c$ by PR~\#982.
$\{n,c^\dagger\} = n\cdot c^\dagger + c^\dagger\cdot n = c^\dagger
+ 0 = c^\dagger$ by PR~\#984. In each case exactly one of the
two products vanishes and the other reproduces the ladder
operator, so the anticommutator equals the ladder operator
itself. Combined with the commutators of PR~\#988 this is the
full per-product breakdown of $n\cdot c$, $c\cdot n$,
$n\cdot c^\dagger$, $c^\dagger\cdot n$.

\textbf{Single-mode trace identities (PR~\#991)}: in the file
\texttt{Fermion/Traces.lean}. From the explicit matrices
$c = \texttt{!![0,1;0,0]}$, $c^\dagger = \texttt{!![0,0;1,0]}$,
$n = \texttt{!![0,0;0,1]}$, and $c\,c^\dagger = 1 - n =
\texttt{!![1,0;0,0]}$ (PR~\#963), one reads off
$\operatorname{tr}(c) = 0$, $\operatorname{tr}(c^\dagger) = 0$
(off-diagonal), $\operatorname{tr}(n) = 1$, and
$\operatorname{tr}(c\,c^\dagger) = 1$. The trace sum
$\operatorname{tr}(n) + \operatorname{tr}(c\,c^\dagger) = 2 =
\operatorname{tr}(I)$ matches the projection-sum identity
$n + c\,c^\dagger = 1$ (PR~\#972). These are the building blocks
for the single-mode partition function $Z = 1 + e^{-\beta}$ and
reduced density matrix entries.

\textbf{Idempotent projection powers $n^{k+1} = n$ and
$(c\,c^\dagger)^{k+1} = c\,c^\dagger$ (PR~\#992)}: in the file
\texttt{Fermion/ProjectionPow.lean}. Both identities are
inductions in $k$ from the squared cases
$n^2 = n$ (\texttt{fermionNumber\_sq}) and
$(c\,c^\dagger)^2 = c\,c^\dagger$ (PR~\#974), with the step
$n^{k+2} = n^{k+1}\cdot n = n\cdot n = n$ (and similarly for
$c\,c^\dagger$). These power identities are the building blocks
for collapsing arbitrary positive-degree polynomials in $n$ (or
$c\,c^\dagger$) into the basis $\{1, n\}$ (resp.\ $\{1,
c\,c^\dagger\}$); that linear collapse itself is not proved here.

\textbf{Multi-mode (JW) ladder anticommutators $\{n_i,c_i\} = c_i$
and $\{n_i,c_i^\dagger\} = c_i^\dagger$ (PR~\#993)}: multi-mode
mirror of PR~\#990, in the file
\texttt{Fermion/JordanWigner/NumberAnticommutators.lean}. The
proof reuses the per-product computations $n_i\cdot c_i = 0$
(from $\sigma^+_i\,\sigma^+_i = 0$) and $c_i\cdot n_i = c_i$
(from $\sigma^+_i\,\sigma^-_i\,\sigma^+_i = \sigma^+_i$) already
encapsulated in the commutator proofs of
\texttt{Fermion/JordanWigner/Number.lean}; the dual pair via
\texttt{Matrix.conjTranspose}.

\textbf{Multi-mode (JW) $c_i$--$c_i^\dagger$ commutator
$[c_i,c_i^\dagger] = 1 - 2\,n_i$ (PR~\#994)}: multi-mode mirror
of PR~\#989, in the file
\texttt{Fermion/JordanWigner/CDaggerCCommutator.lean}. From the
same-site CAR $c_i\,c_i^\dagger + c_i^\dagger c_i = 1$
(\texttt{fermionMultiAnticomm\_self}) one gets
$c_i\,c_i^\dagger = 1 - c_i^\dagger c_i = 1 - n_i$, where the
last equality uses the definitional
$n_i = c_i^\dagger c_i$. Subtracting $c_i^\dagger c_i = n_i$
yields $[c_i,c_i^\dagger] = (1 - n_i) - n_i = 1 - 2 n_i$.

\textbf{Multi-mode (JW) number-operator power $n_i^{k+1} = n_i$
(PR~\#995)}: multi-mode mirror of (the first half of) PR~\#992,
in the file \texttt{Fermion/JordanWigner/NumberPow.lean}. Direct
induction on $k$ from $n_i^2 = n_i$
(\texttt{fermionMultiNumber\_sq}, in
\texttt{Fermion/JordanWigner/Operators.lean}).

\textbf{Multi-mode (JW) hole occupation $c_i\,c_i^\dagger = 1 -
n_i$ and resolution of identity $n_i + c_i\,c_i^\dagger = 1$
(PR~\#996)}: multi-mode mirror of PRs~\#963 and \#972, in the
file \texttt{Fermion/JordanWigner/CDaggerCIdentity.lean}. From
the same-site CAR $c_i\,c_i^\dagger + c_i^\dagger c_i = 1$
(\texttt{fermionMultiAnticomm\_self}) and the definitional
$n_i = c_i^\dagger c_i$, one extracts $c_i\,c_i^\dagger = 1 -
n_i$ via \texttt{eq\_sub\_of\_add\_eq}; rearranging gives the
projection-sum form $n_i + c_i\,c_i^\dagger = 1$.

\textbf{Multi-mode (JW) hole-projection idempotency
$(c_i\,c_i^\dagger)^2 = c_i\,c_i^\dagger$ and power
$(c_i\,c_i^\dagger)^{k+1} = c_i\,c_i^\dagger$ (PR~\#997)}:
multi-mode mirror of PRs~\#974 and the second half of \#992,
in the file
\texttt{Fermion/JordanWigner/CDaggerCProjection.lean}. The
squared identity uses $c_i\,c_i^\dagger = 1 - n_i$ (PR~\#996)
and $n_i^2 = n_i$ (\texttt{fermionMultiNumber\_sq}) to give
$(1 - n_i)^2 = 1 - 2n_i + n_i^2 = 1 - n_i$. The power identity
is then a direct induction in $k$ from the squared case.

\textbf{Multi-mode (JW) hole projection Hermitian
$(c_i\,c_i^\dagger)^H = c_i\,c_i^\dagger$ (PR~\#998)}:
multi-mode mirror of PR~\#980, in the file
\texttt{Fermion/JordanWigner/CDaggerCHermitian.lean}. From
$(c_i)^H = c_i^\dagger$
(\texttt{fermionMultiAnnihilation\_conjTranspose}) and
$(c_i^\dagger)^H = c_i$
(\texttt{fermionMultiCreation\_conjTranspose}), one has
$(c_i\,c_i^\dagger)^H = (c_i^\dagger)^H \cdot (c_i)^H =
c_i\,c_i^\dagger$ via \texttt{Matrix.conjTranspose\_mul}.

\textbf{Multi-mode (JW) orthogonal projections
$n_i\cdot(c_i\,c_i^\dagger) = 0$ and
$(c_i\,c_i^\dagger)\cdot n_i = 0$ (PR~\#999)}: multi-mode
mirror of PR~\#976, in the file
\texttt{Fermion/JordanWigner/ProjectionsOrthogonal.lean}. Direct
from $c_i^2 = 0$ (\texttt{fermionMultiAnnihilation\_sq}) and
$(c_i^\dagger)^2 = 0$ (\texttt{fermionMultiCreation\_sq}) using
the definitional $n_i = c_i^\dagger c_i$:
$n_i(c_i\,c_i^\dagger) = c_i^\dagger\,c_i^2\,c_i^\dagger = 0$,
$(c_i\,c_i^\dagger)n_i = c_i\,(c_i^\dagger)^2\,c_i = 0$.

\textbf{Multi-mode (JW) projections commute
$\texttt{Commute}\,n_i\,(c_i\,c_i^\dagger)$ (PR~\#1000)}:
multi-mode mirror of PR~\#978, in the file
\texttt{Fermion/JordanWigner/ProjectionsCommute.lean}. Trivially
from PR~\#999: both products $n_i\cdot(c_i\,c_i^\dagger)$ and
$(c_i\,c_i^\dagger)\cdot n_i$ are zero, hence equal.

\textbf{Multi-mode (JW) number-annihilation identities
$n_i\cdot c_i = 0$ and $c_i\cdot n_i = c_i$ (PR~\#1001)}:
multi-mode mirror of PR~\#982, in the file
\texttt{Fermion/JordanWigner/AnnihilationNumberIdentities.lean}.
$n_i\cdot c_i = (c_i^\dagger c_i)c_i = c_i^\dagger\,c_i^2 = 0$
via \texttt{fermionMultiAnnihilation\_sq}. $c_i\cdot n_i =
c_i\,(c_i^\dagger c_i) = (c_i\,c_i^\dagger)c_i = (1-n_i)c_i =
c_i - n_i\,c_i = c_i$ via PR~\#996 and the previous identity.

\textbf{Multi-mode (JW) number-creation identities
$c_i^\dagger\cdot n_i = 0$ and $n_i\cdot c_i^\dagger =
c_i^\dagger$ (PR~\#1002)}: multi-mode mirror of PR~\#984, in
the file
\texttt{Fermion/JordanWigner/CreationNumberIdentities.lean}.
$c_i^\dagger\cdot n_i = c_i^\dagger(c_i^\dagger c_i) =
(c_i^\dagger)^2\,c_i = 0$ via
\texttt{fermionMultiCreation\_sq}.
$n_i\cdot c_i^\dagger = (c_i^\dagger c_i)\,c_i^\dagger =
c_i^\dagger(c_i\,c_i^\dagger) = c_i^\dagger(1-n_i) =
c_i^\dagger - c_i^\dagger\,n_i = c_i^\dagger$ via PR~\#996 and
the previous.

\textbf{Multi-mode (JW) partial-isometry identities
$c_i\,c_i^\dagger c_i = c_i$ and $c_i^\dagger c_i\,c_i^\dagger
= c_i^\dagger$ (PR~\#1003)}: multi-mode mirror of PR~\#986, in
the file \texttt{Fermion/JordanWigner/PartialIsometry.lean}.
Direct one-line corollaries of PRs~\#1001 and \#1002 via the
definitional $n_i = c_i^\dagger c_i$:
$c_i\,c_i^\dagger c_i = c_i\,(c_i^\dagger c_i) = c_i\,n_i = c_i$
by PR~\#1001;
$c_i^\dagger c_i\,c_i^\dagger = (c_i^\dagger c_i)\,c_i^\dagger
= n_i\,c_i^\dagger = c_i^\dagger$ by PR~\#1002. Standard
$S\,S^*S = S$ / $S^*S\,S^* = S^*$ partial-isometry identities
of operator algebra, lifted to JW multi-mode at the same site.

\textbf{Multi-mode (JW) hole projections commute
$\texttt{Commute}\,(c_i\,c_i^\dagger)\,(c_j\,c_j^\dagger)$
(PR~\#1004)}: in the file
\texttt{Fermion/JordanWigner/HoleProjectionsCommute.lean}. From
$c_i\,c_i^\dagger = 1 - n_i$ (PR~\#996) and
$\texttt{Commute}\,n_i\,n_j$
(\texttt{fermionMultiNumber\_commute},
\texttt{Fermion/JordanWigner/CAR/SameSite.lean}), expanding
$(1 - n_i)(1 - n_j) = 1 - n_j - n_i + n_i\,n_j$ and the
opposite ordering gives $1 - n_i - n_j + n_j\,n_i$; the two
agree by $n_i\,n_j = n_j\,n_i$.

\textbf{Hubbard same-site double-occupancy operator
$\texttt{Commute}\,n_\uparrow(i)\,n_\downarrow(i)$ and
$(n_\uparrow(i)\,n_\downarrow(i))^2 = n_\uparrow(i)\,
n_\downarrow(i)$ (PR~\#1005)}: in the file
\texttt{Fermion/JordanWigner/Hubbard/DoubleOccupancyProjection.lean}.
The two spin-species number operators at site $i$ live at
distinct underlying JW positions $2i$ and $2i+1$, so they
commute by \texttt{fermionMultiNumber\_commute}. Their product
$n_\uparrow n_\downarrow$ is then the on-site
double-occupancy projection: $(n_\uparrow n_\downarrow)^2 =
n_\uparrow (n_\downarrow n_\uparrow) n_\downarrow =
n_\uparrow^2 n_\downarrow^2 = n_\uparrow n_\downarrow$ via the
commute and \texttt{fermionMultiNumber\_sq} for both factors.
The eigenvalue-$1$ subspace is the doubly-occupied site
state, so the on-site Hubbard interaction
$U \sum_i n_\uparrow(i) n_\downarrow(i)$ is $U$ times a sum
of projections (each summand is a projection by this PR;
cross-site commutation of summands is not proved here).

\textbf{Cross-site Hubbard double-occupancy commute
$\texttt{Commute}\,(n_\uparrow(i) n_\downarrow(i))\,
(n_\uparrow(j) n_\downarrow(j))$ (PR~\#1006)}: in the file
\texttt{Fermion/JordanWigner/Hubbard/DoubleOccupancyCommute.lean}.
The four underlying single-species number operators
$n_\uparrow(i), n_\downarrow(i), n_\uparrow(j), n_\downarrow(j)$
live at JW positions $2i$, $2i+1$, $2j$, $2j+1$ and pairwise
commute by \texttt{fermionMultiNumber\_commute}; the products
commute by \texttt{Commute.mul\_left} and
\texttt{Commute.mul\_right}. Combined with PR~\#1005, the
on-site Hubbard interaction $U\sum_i n_\uparrow(i)
n_\downarrow(i)$ is $U$ times a sum of pairwise commuting
projections, hence simultaneously diagonalisable.

\textbf{Spinful Hubbard number-operator Hermiticity (PR~\#1007)}:
in the file
\texttt{Fermion/JordanWigner/Hubbard/SpinfulNumberHermitian.lean}.
Standalone named Hermiticity lemmas
$(n_\uparrow(i))^H = n_\uparrow(i)$,
$(n_\downarrow(i))^H = n_\downarrow(i)$, and
$(n_\uparrow(i) n_\downarrow(i))^H = n_\uparrow(i)
n_\downarrow(i)$. Each is a one-line corollary of
\texttt{fermionMultiNumber\_isHermitian} and
\texttt{fermionMultiNumber\_mul\_isHermitian}.

\textbf{Cross-site multi-mode (JW) hole projection commutes
with ladder operators
$\texttt{Commute}\,(c_i\,c_i^\dagger)\,c_j$ and
$\texttt{Commute}\,(c_i\,c_i^\dagger)\,c_j^\dagger$ for $i\neq j$
(PR~\#1008)}: in the file
\texttt{Fermion/JordanWigner/HoleProjectionCommuteLadder.lean}.
From $c_i\,c_i^\dagger = 1 - n_i$ (PR~\#996) and the cross-site
$\texttt{Commute}\,n_i\,c_j$
(\texttt{fermionMultiNumber\_commute\_fermionMultiAnnihilation\_of\_ne})
/ $\texttt{Commute}\,n_i\,c_j^\dagger$
(\texttt{fermionMultiNumber\_commute\_fermionMultiCreation\_of\_ne}),
expanding $(1-n_i)c_j$ vs $c_j(1-n_i)$ and similarly for
$c_j^\dagger$ closes both via $n_i\,c_j = c_j\,n_i$ (resp.\ for
$c_j^\dagger$).

\textbf{Multi-mode (JW) hole projection commutes with number
$\texttt{Commute}\,(c_i\,c_i^\dagger)\,n_j$ and
$\texttt{Commute}\,n_i\,(c_j\,c_j^\dagger)$ (PR~\#1013)}:
mixed-form siblings of PR~\#1004, in the file
\texttt{Fermion/JordanWigner/HoleProjectionCommuteNumber.lean}.
Direct from $c_i\,c_i^\dagger = 1 - n_i$ (PR~\#996) and
$\texttt{Commute}\,n_i\,n_j$
(\texttt{fermionMultiNumber\_commute}): expand
$(1 - n_i)\,n_j = n_j - n_i\,n_j = n_j - n_j\,n_i = n_j(1 -
n_i)$.

\textbf{Single-mode ladder-on-hole-projection vanishing
$c\,(c\,c^\dagger) = 0$ and $(c\,c^\dagger)\,c^\dagger = 0$
(PR~\#1009)}: in the file
\texttt{Fermion/CCDaggerLadderZero.lean}. Direct from
$c^2 = 0$ (\texttt{fermionAnnihilation\_sq}) and
$(c^\dagger)^2 = 0$ (\texttt{fermionCreation\_sq}):
$c\,(c\,c^\dagger) = c^2\,c^\dagger = 0$ and
$(c\,c^\dagger)\,c^\dagger = c\,(c^\dagger)^2 = 0$. Together
with $(c\,c^\dagger)\,c = c$ (PR~\#982) and $c^\dagger\,
(c\,c^\dagger) = c^\dagger$ (PR~\#984) this is the full
four-way breakdown of left/right multiplication of
$c\,c^\dagger$ by the ladder operators.

\textbf{Single-mode $(c + c^\dagger)^2 = 1$ (X-Pauli analog,
PR~\#1021)}: in the file
\texttt{Fermion/CPlusCDaggerSq.lean}. The symmetric Hermitian
combination is an involution. Direct expansion
$(c + c^\dagger)^2 = c^2 + c\,c^\dagger + c^\dagger c +
(c^\dagger)^2 = 0 + 1 + 0 = 1$ via $c^2 = (c^\dagger)^2 = 0$
and $\{c, c^\dagger\} = 1$
(\texttt{fermionAnticomm\_self}). In the spin-$1/2$ identification
$c = \sigma^+, c^\dagger = \sigma^-$ this is the algebraic
shadow of $\sigma_x^2 = I$. Useful for quench dynamics and
Majorana decompositions.

\textbf{Multi-mode (JW) $(c_i + c_i^\dagger)^2 = 1$
(PR~\#1022)}: multi-mode mirror of PR~\#1021, in the file
\texttt{Fermion/JordanWigner/CPlusCDaggerSq.lean}. Same
expansion using \texttt{fermionMultiAnnihilation\_sq},
\texttt{fermionMultiCreation\_sq}, and
\texttt{fermionMultiAnticomm\_self}. In the JW representation
$c_i + c_i^\dagger = \mathrm{jwString}_i \cdot \sigma^x_i$, and
$(\mathrm{jwString}_i \cdot \sigma^x_i)^2 = 1$ since
$\mathrm{jwString}_i^2 = 1$ (\texttt{jwString\_sq}).

\textbf{Single-mode $(c - c^\dagger)^2 = -1$ (iY-Pauli analog,
PR~\#1023)}: in the file
\texttt{Fermion/CMinusCDaggerSq.lean}. Companion to PR~\#1021.
The antisymmetric anti-Hermitian combination squares to $-1$.
$(c - c^\dagger)^2 = c^2 - c\,c^\dagger - c^\dagger c + (c^\dagger)^2
= -(c\,c^\dagger + c^\dagger c) = -1$ via $c^2 = (c^\dagger)^2 = 0$
and \texttt{fermionAnticomm\_self}. In spin-$1/2$ identification
$c - c^\dagger = \sigma^+ - \sigma^- = i\,\sigma_y$ and
$(i\,\sigma_y)^2 = -\sigma_y^2 = -I$.

\textbf{Multi-mode (JW) $(c_i - c_i^\dagger)^2 = -1$
(PR~\#1024)}: multi-mode mirror of PR~\#1023, in the file
\texttt{Fermion/JordanWigner/CMinusCDaggerSq.lean}. Same
expansion using \texttt{fermionMultiAnnihilation\_sq},
\texttt{fermionMultiCreation\_sq}, and
\texttt{fermionMultiAnticomm\_self}. Companion to PR~\#1022.

\textbf{Single-mode mixed Pauli-X\,$\cdot$\,iY products
$(c + c^\dagger)(c - c^\dagger) = 2n - 1$ and
$(c - c^\dagger)(c + c^\dagger) = 1 - 2n$ (PR~\#1025)}: in the
file \texttt{Fermion/CPlusCDaggerMulCMinusCDagger.lean}.
Direct expansion with $c^2 = (c^\dagger)^2 = 0$,
$c\,c^\dagger = 1 - n$ (PR~\#963), and $c^\dagger c = n$:
$(c + c^\dagger)(c - c^\dagger) = 0 - (1 - n) + n - 0 =
2n - 1$. Physics identification:
$(c + c^\dagger)(c - c^\dagger) = \sigma_x \cdot i\sigma_y =
i^2 \sigma_z = -\sigma_z$; in our convention
$n = (I - \sigma_z)/2$, so $-\sigma_z = 2n - 1$.

\textbf{Single-mode Pauli-X / iY analog Hermiticity +
anticommute structure (PR~\#1026)}: in the file
\texttt{Fermion/CPlusMinusCDaggerHermitian.lean}. Three lemmas:
$(c + c^\dagger)$ Hermitian, $(c - c^\dagger)^H = -(c -
c^\dagger)$ (anti-Hermitian), and $\{c + c^\dagger, c -
c^\dagger\} = 0$ (anticommute). The first two follow from
$c^H = c^\dagger$ and $(c^\dagger)^H = c$ via
\texttt{Matrix.conjTranspose\_add/\_sub}; the anticommutator
follows from PR~\#1025: $(2n - 1) + (1 - 2n) = 0$. Physics
identification: $\{\sigma_x, \sigma_y\} = 0$ in the spin-$1/2$
representation.

\textbf{Multi-mode (JW) $(c_i \pm c_i^\dagger)$ Pauli-X/iY
analog full structure (PR~\#1027)}: multi-mode mirror of
PRs~\#1025 and \#1026, in the file
\texttt{Fermion/JordanWigner/CPlusMinusCDaggerPauli.lean}.
Bundles all five identities: mixed products $(c_i + c_i^\dagger)
(c_i - c_i^\dagger) = 2 n_i - 1$ and the reverse, plus
Hermiticity / anti-Hermiticity / anticommute. Same proof
patterns lifted to multi-mode using the per-site identities of
PRs~\#996 and the existing
\texttt{fermionMultiAnnihilation\_conjTranspose} /
\texttt{fermionMultiCreation\_conjTranspose}.

\textbf{Single-mode $(1 - 2\,n)^2 = 1$ ($\sigma_z$-analog
involution, PR~\#1028)}: in the file
\texttt{Fermion/OneSubTwoNumberSq.lean}. Direct from $n^2 = n$
(\texttt{fermionNumber\_sq}): $(1 - 2n)^2 = 1 - 4n + 4n^2 =
1 - 4n + 4n = 1$. Physics identification: $1 - 2n = \sigma_z$
in our convention $n = (I - \sigma_z)/2$, and $\sigma_z^2 = I$.
Together with PRs~\#1021 ($\sigma_x^2 = 1$) and \#1023
($(i\sigma_y)^2 = -1$) this completes the single-mode Pauli-trio
$\sigma_\alpha^2 = \pm 1$ involution identities.

\textbf{Multi-mode (JW) $(1 - 2\,n_i)^2 = 1$ (PR~\#1029)}:
multi-mode mirror of PR~\#1028, in the file
\texttt{Fermion/JordanWigner/OneSubTwoNumberSq.lean}. Same
algebra using \texttt{fermionMultiNumber\_sq}. Together with
PRs~\#1022 and \#1024 this completes the multi-mode Pauli-trio
$\sigma_\alpha^2 = \pm 1$ involution identities.

\textbf{Symmetric `\_of\_ne` cross-site CAR (PR~\#1030)}: in
the file
\texttt{Fermion/JordanWigner/CAR/CrossSiteOfNe.lean}. Lifts the
four `\_lt` cross-site CAR identities to the symmetric form
accepting any $i \neq j$ by trichotomy on $i.\mathrm{val}$ vs
$j.\mathrm{val}$: the $\_lt$ form covers $i < j$ directly, and
$j < i$ follows by swapping $i, j$ and applying \texttt{add\_comm}.
Foundational for upstream Pauli-analog cross-site lemmas.

\textbf{Cross-site Pauli-X-analog anticommute
$\{c_i + c_i^\dagger, c_j + c_j^\dagger\} = 0$ for $i \neq j$
(PR~\#1031)}: in the file
\texttt{Fermion/JordanWigner/CPlusCDaggerAnticomm.lean}. Direct
expansion using the four `\_of\_ne` cross-site CAR identities
(PR~\#1030):
$\{c_i + c_i^\dagger, c_j + c_j^\dagger\} = \{c_i, c_j\} +
\{c_i, c_j^\dagger\} + \{c_i^\dagger, c_j\} +
\{c_i^\dagger, c_j^\dagger\} = 0$. JW-string fingerprint: while
on-site $\sigma_x$ operators commute, the JW-dressed
combinations $c_i + c_i^\dagger = \mathrm{jwString}_i \cdot
\sigma^x_i$ pick up the JW string sign and anticommute across
sites.

\textbf{Cross-site mixed Pauli anticommutators (PR~\#1032)}: in
the file
\texttt{Fermion/JordanWigner/CMinusCDaggerAnticomm.lean}.
Companion to PR~\#1031: $\{c_i - c_i^\dagger, c_j - c_j^\dagger\}
= 0$ and $\{c_i + c_i^\dagger, c_j - c_j^\dagger\} = 0$ for
$i \neq j$. Same expansion technique with appropriate sign
combinations: $\{X, X\} - \{X, X^\dagger\} - \{X^\dagger, X\}
+ \{X^\dagger, X^\dagger\}$ (resp.\ similar for the mixed
case). Together with PR~\#1031 covers all four sign
combinations of cross-site Pauli-X/iY-analog anticommutators.

\textbf{Cross-site $\texttt{Commute}\,n_i\,(c_j \pm
c_j^\dagger)$ for $i \neq j$ (PR~\#1033)}: in the file
\texttt{Fermion/JordanWigner/NumberCommutePauliOfNe.lean}.
Direct from the existing
\texttt{fermionMultiNumber\_commute\_fermionMultiAnnihilation\_of\_ne}
and \texttt{\_Creation\_of\_ne}, combined via
\texttt{Commute.add\_right} / \texttt{.sub\_right}.

\textbf{Multi-mode (JW) ladder-on-hole-projection vanishing
$c_i\,(c_i\,c_i^\dagger) = 0$ and $(c_i\,c_i^\dagger)\,
c_i^\dagger = 0$ (PR~\#1010)}: multi-mode mirror of PR~\#1009,
in the file
\texttt{Fermion/JordanWigner/CDaggerCLadderZero.lean}. Direct
from $c_i^2 = 0$ (\texttt{fermionMultiAnnihilation\_sq}) and
$(c_i^\dagger)^2 = 0$ (\texttt{fermionMultiCreation\_sq}).
Together with $(c_i\,c_i^\dagger)\,c_i = c_i$ (PR~\#1001) and
$c_i^\dagger\,(c_i\,c_i^\dagger) = c_i^\dagger$ (PR~\#1002)
this is the full four-way breakdown of left/right
multiplication of the multi-mode hole projection by the
same-site ladder operators.

\textbf{Hubbard per-site 4-state partition of identity
$p_\emptyset(i) + p_\uparrow(i) + p_\downarrow(i) +
p_\Uparrow(i) = 1$ (PR~\#1011)}: in the file
\texttt{Fermion/JordanWigner/Hubbard/SitePartitionIdentity.lean}.
The four occupation projections
$p_\emptyset(i) := (1 - n_\uparrow(i))(1 - n_\downarrow(i))$,
$p_\uparrow(i) := n_\uparrow(i)(1 - n_\downarrow(i))$,
$p_\downarrow(i) := (1 - n_\uparrow(i)) n_\downarrow(i)$,
$p_\Uparrow(i) := n_\uparrow(i) n_\downarrow(i)$
sum to the identity by distributing
$((1 - n_\uparrow) + n_\uparrow)((1 - n_\downarrow) +
n_\downarrow) = 1\cdot 1 = 1$. Standard 4-state per-site
decomposition underlying the Hubbard Hilbert-space tensor
structure $\mathbb{C}^4 = |0,0\rangle \oplus |1,0\rangle
\oplus |0,1\rangle \oplus |1,1\rangle$.

\textbf{Hubbard empty/up/down per-site projections idempotent
$(p_\emptyset)^2 = p_\emptyset$, $(p_\uparrow)^2 = p_\uparrow$,
$(p_\downarrow)^2 = p_\downarrow$ (PR~\#1012)}: in the file
\texttt{Fermion/JordanWigner/Hubbard/SiteProjectionsIdempotent.lean}.
Three companion idempotency lemmas to PR~\#1005 ($(p_\Uparrow)^2
= p_\Uparrow$). Each proof uses the same-site $\texttt{Commute}\,
n_\uparrow(i)\,n_\downarrow(i)$ (PR~\#1005) to swap the middle
factors, then applies the squared identities
\texttt{fermionMultiNumber\_sq} ($n_\sigma^2 = n_\sigma$) and
\texttt{fermionMultiAnnihilation\_mul\_fermionMultiCreation\_sq}
together with PR~\#996 ($c_\sigma c_\sigma^\dagger = 1 -
n_\sigma$) for the complementary $(1 - n_\sigma)^2 = 1 -
n_\sigma$. Together with PR~\#1005 these establish that all
four per-site occupation operators are projections.

\textbf{Hubbard per-site doubly-occupied and empty projections
orthogonal $p_\Uparrow\cdot p_\emptyset = 0$ and
$p_\emptyset\cdot p_\Uparrow = 0$ (PR~\#1014)}: in the file
\texttt{Fermion/JordanWigner/Hubbard/SiteProjectionsDoublyEmpty.lean}.
Direct from the same-site $\texttt{Commute}\,n_\uparrow\,
n_\downarrow$ (PR~\#1005) and $n_\sigma(1 - n_\sigma) =
(1 - n_\sigma)n_\sigma = 0$ (specialised from PR~\#999 via
$c_\sigma c_\sigma^\dagger = 1 - n_\sigma$, PR~\#996). Most
extreme orthogonal pair of the per-site 4-state decomposition
(PR~\#1011): "doubly occupied" cannot share with "empty".

\textbf{Hubbard per-site projections $p_\emptyset$,
$p_\uparrow$, $p_\downarrow$ Hermitian (PR~\#1015)}: in the
file
\texttt{Fermion/JordanWigner/Hubbard/SiteProjectionsHermitian.lean}.
Three companions to PR~\#1007 ($p_\Uparrow$ Hermitian). Each
per-site projection is a product of two commuting Hermitian
factors --- $(1 - n_\sigma)$ Hermitian via
\texttt{Matrix.IsHermitian.sub} with
\texttt{Matrix.isHermitian\_one}, and pairs commute by ring
algebra from PR~\#1005's $\texttt{Commute}\,n_\uparrow\,
n_\downarrow$. Closed by
\texttt{Matrix.IsHermitian.mul\_of\_commute}. Together with
PR~\#1007 all four per-site occupation projections are
Hermitian.

\textbf{Hubbard per-site only-up and only-down projections
orthogonal $p_\uparrow\cdot p_\downarrow = 0$ and
$p_\downarrow\cdot p_\uparrow = 0$ (PR~\#1016)}: in the file
\texttt{Fermion/JordanWigner/Hubbard/SiteProjectionsUpDown.lean}.
Swap the middle factors via the same-site
$\texttt{Commute}\,(1 - n_\downarrow)\,(1 - n_\uparrow)$
(propagated from PR~\#1005), then collapse via
$n_\sigma(1 - n_\sigma) = 0$ (resp.\ $(1 - n_\sigma)
n_\sigma = 0$) specialised from PR~\#999. Together with
PR~\#1014 ($p_\Uparrow\cdot p_\emptyset = 0$) this covers two
of the six pairs of cross-projection orthogonality for the
per-site 4-state decomposition (PR~\#1011).

\textbf{Hubbard per-site empty projection orthogonal to both
single-occupancy projections (PR~\#1017)}: in the file
\texttt{Fermion/JordanWigner/Hubbard/SiteProjectionsEmptySingle.lean}.
Four orthogonality lemmas (both orderings):
$p_\emptyset\cdot p_\uparrow = 0$,
$p_\uparrow\cdot p_\emptyset = 0$,
$p_\emptyset\cdot p_\downarrow = 0$,
$p_\downarrow\cdot p_\emptyset = 0$. Same proof template as
PRs~\#1014 and \#1016. Together with those PRs this covers
four of the six pairs of cross-projection orthogonality for
the per-site 4-state decomposition (PR~\#1011); the remaining
two ($p_\uparrow\cdot p_\Uparrow$, $p_\downarrow\cdot
p_\Uparrow$) follow the same template.

\textbf{Hubbard single-occupancy $\perp$ doubly-occupied
per-site projections (PR~\#1018, completes 6/6 orthogonality)}:
in the file
\texttt{Fermion/JordanWigner/Hubbard/SiteProjectionsSingleDoubly.lean}.
Final two pairs:
$p_\uparrow\cdot p_\Uparrow = p_\Uparrow\cdot p_\uparrow = 0$
and $p_\downarrow\cdot p_\Uparrow = p_\Uparrow\cdot
p_\downarrow = 0$. Same proof template; collapse uses
\texttt{fermionMultiNumber\_sq} ($n_\sigma^2 = n_\sigma$) and
$n_\sigma(1 - n_\sigma) = 0$. Together with PRs~\#1014,
\#1016, \#1017 \textbf{all six pairs of cross-projection
orthogonality are now formalised}, completing the proof that
$\{p_\emptyset, p_\uparrow, p_\downarrow, p_\Uparrow\}$ is a
complete set of mutually orthogonal Hermitian projections
summing to the identity (combined with PRs~\#1007, \#1011,
\#1015).

\textbf{Hubbard per-site projection-aggregation identities
(PR~\#1019)}: in the file
\texttt{Fermion/JordanWigner/Hubbard/SiteProjectionsSpinResolved.lean}.
Four spin-resolved aggregation identities of the per-site
projections:
$p_\uparrow + p_\Uparrow = n_\uparrow$,
$p_\downarrow + p_\Uparrow = n_\downarrow$,
$p_\emptyset + p_\uparrow = 1 - n_\downarrow$,
$p_\emptyset + p_\downarrow = 1 - n_\uparrow$.
Direct factoring out the common left/right factor and using
$(1 - n_\sigma) + n_\sigma = 1$. Shows how the per-site 4-state
projections aggregate to the spin-resolved number operators
(and their complements), complementing the partition-of-identity
(PR~\#1011) $\sum_\alpha p_\alpha = 1$.

\textbf{Hubbard same-site per-site projections pairwise commute
(PR~\#1020)}: in the file
\texttt{Fermion/JordanWigner/Hubbard/SiteProjectionsCommute.lean}.
The four per-site occupation projections at the same site $i$
pairwise commute: $\texttt{Commute}\,p_\alpha(i)\,p_\beta(i)$
for any $\alpha, \beta \in \{\emptyset, \uparrow, \downarrow,
\Uparrow\}$. For $\alpha \neq \beta$ trivial from
PRs~\#1014, \#1016, \#1017, \#1018: both products vanish, so
the commutator is $0$. Together with the partition / idempotency
/ Hermiticity / orthogonality results, the per-site projection
algebra is fully commutative on top of being a complete
orthogonal partition.

\textbf{Hubbard per-site projection power identities
(PR~\#1034)}: in the file
\texttt{Fermion/JordanWigner/Hubbard/SiteProjectionsPow.lean}.
For every $k : \mathbb{N}$ each of the four per-site occupation
projections satisfies $(p_\alpha(i))^{k+1} = p_\alpha(i)$.
Each by direct induction on $k$ from the squared identities of
PRs~\#1005 and \#1012. Mirrors PR~\#992 (single-mode
$n^{k+1} = n$ and $(c\,c^\dagger)^{k+1} = c\,c^\dagger$).

\textbf{Cross-site Hubbard empty projection commute (PR~\#1035)}:
in the file
\texttt{Fermion/JordanWigner/Hubbard/EmptyProjectionCommute.lean}.
For any $i, j$,
$\texttt{Commute}\,(p_\emptyset(i))\,(p_\emptyset(j))$. All four
underlying $n_\sigma$ pairwise commute (and so do their
$(1 - n_\sigma)$-complements by ring algebra), and the products
commute via \texttt{Commute.mul\_left} /
\texttt{.mul\_right}. Companion to PR~\#1006 for $p_\Uparrow$.

\textbf{Cross-site Hubbard single-occupancy projections commute
(PR~\#1036)}: in the file
\texttt{Fermion/JordanWigner/Hubbard/SingleProjectionsCommute.lean}.
For any $i, j$, $\texttt{Commute}\,(p_\uparrow(i))\,
(p_\uparrow(j))$ and the analogous $p_\downarrow$ statement.
Same proof technique as PRs~\#1006 and \#1035, with three
auxiliary lemmas covering the relevant
$\texttt{Commute}$ patterns between $n_\sigma$ and
$(1 - n_\tau)$ at distinct JW positions. Together with
PRs~\#1006 and \#1035 this completes the cross-site diagonal
Commute relations for all four per-site occupation projections.

\textbf{Cross-projection $\texttt{Commute}\,(p_\uparrow(i))\,
(p_\downarrow(j))$ for any $i, j$ (PR~\#1037)}: in the file
\texttt{Fermion/JordanWigner/Hubbard/UpDownProjectionCommute.lean}.
Non-diagonal extension of the diagonal cross-site Commute
relations. Same proof template as PRs~\#1006, \#1035, \#1036:
the four underlying $n_\sigma$ (and complements) pairwise
commute, so the products commute via \texttt{Commute.mul\_left}
/ \texttt{.mul\_right}.

\textbf{Remaining 5 cross-projection commute relations (PR~\#1038,
completes 16/16 commute matrix)}: in the file
\texttt{Fermion/JordanWigner/Hubbard/RemainingProjectionCommutes.lean}.
Five remaining pairs:
$\texttt{Commute}\,(p_\emptyset(i))\,(p_\uparrow(j))$,
$\texttt{Commute}\,(p_\emptyset(i))\,(p_\downarrow(j))$,
$\texttt{Commute}\,(p_\emptyset(i))\,(p_\Uparrow(j))$,
$\texttt{Commute}\,(p_\uparrow(i))\,(p_\Uparrow(j))$,
$\texttt{Commute}\,(p_\downarrow(i))\,(p_\Uparrow(j))$.
Same proof template throughout. Combined with PRs~\#1006,
\#1020, \#1035, \#1036, \#1037, all 16 combinations of
$\texttt{Commute}\,(p_\alpha(i))\,(p_\beta(j))$ are formalised,
making the 4-state per-site occupation algebra fully commutative
across arbitrary site pairs.

\textbf{All-down spin state for the Hubbard model (PR~\#1039)}:
in the file
\texttt{Fermion/JordanWigner/Hubbard/AllDownState.lean}. Mirror
of \texttt{AllUpState.lean}: defines the fully spin-polarised
all-down state $|\!\downarrow\cdots\downarrow\rangle$ (basis
vector with occupation $1$ on odd JW indices), and proves
$n_{i,\downarrow}\,|\!\downarrow\cdots\downarrow\rangle =
|\!\downarrow\cdots\downarrow\rangle$, $n_{i,\uparrow}\,
|\!\downarrow\cdots\downarrow\rangle = 0$, plus the no-double-
occupancy result $H_{\rm int}\,|\!\downarrow\cdots\downarrow
\rangle = 0$ (and the relevant ladder annihilations). The
all-down state belongs to the kernel of $H_{\rm int}$ by the
same argument as for the all-up state, manifesting up/down
symmetry.

\textbf{All-down state total-number / spin operators (PR~\#1040)}:
in the file
\texttt{Fermion/JordanWigner/Hubbard/AllDownStateTotalNumber.lean}.
Mirror of \texttt{SaturatedFerromagnetism.lean}'s all-up
versions: $\hat N_\downarrow\,|\!\downarrow\cdots\downarrow
\rangle = (N+1)\,|\!\downarrow\cdots\downarrow\rangle$,
$\hat N_\uparrow\,|\!\downarrow\cdots\downarrow\rangle = 0$,
$\hat S^z_{\rm tot}\,|\!\downarrow\cdots\downarrow\rangle =
-(N+1)/2\,|\!\downarrow\cdots\downarrow\rangle$, and
$\hat S^-_{\rm tot}\,|\!\downarrow\cdots\downarrow\rangle = 0$.
The all-down state is the lowest-weight state of the global
SU(2) spin algebra (lowest $S^z$, annihilated by $S^-$).

\textbf{Total-spin operator Hermiticity (PR~\#1041)}: in the
file
\texttt{Fermion/JordanWigner/Hubbard/SpinTotHermitian.lean}.
Three Hermiticity-related lemmas: $(\hat S^-_{\rm tot})^H =
\hat S^+_{\rm tot}$ (dual of \texttt{fermionTotalSpinPlus\_conjTranspose}
via $((\cdot)^H)^H = \mathrm{id}$), $\hat S^z_{\rm tot}$
Hermitian (sum of Hermitian numbers scaled by real $1/2$),
and $\hat S^2_{\rm tot}$ Hermitian (using $(\hat S^- \hat S^+)^H =
\hat S^- \hat S^+$ and that $\hat S_z(\hat S_z + 1)$ is a
polynomial in the Hermitian $\hat S_z$).

\textbf{Linear independence of the first ladder pair (PR~\#891)}:
combines PRs~\#875, \#886, \#889, \#890 into the smallest
non-trivial linear-independence statement on the saturated-
ferromagnet ladder. In the new file
\texttt{Quantum/SpinS/SaturatedPairLinearIndependent.lean} the
all-aligned state non-vanishing
\texttt{allAlignedStateS\_ne\_zero} (its value at the all-aligned
configuration is $1$) is paired with the eigenspace bridges
\texttt{allAlignedStateS\_zero\_mem\_eigenspace\_mMax} and
\texttt{totalSpinSOpMinus\_mulVec\_allAlignedStateS\_zero\_mem\_eigenspace\_mMaxSubOne}
to satisfy the \texttt{Module.End.HasEigenvector} hypotheses of
mathlib's \texttt{Module.End.eigenvectors\_linearIndependent'}.
The eigenvalue pair $(m_{\max},m_{\max}-1)$ is injective, hence
the family $![|\sigma_\top\rangle,\hat S^-_{\rm tot}|\sigma_\top\rangle]$
is \texttt{LinearIndependent}~$\mathbb{C}$ for $0<N$ and
$[\texttt{Nonempty V}]$
(\texttt{allAlignedStateS\_zero\_totalSpinSOpMinus\_mulVec\_linearIndependent}).
The symmetric raising-side pair is
\texttt{allAlignedStateS\_last\_totalSpinSOpPlus\_mulVec\_linearIndependent}.
Generalises by induction (with the higher-$k$ extension of
PR~\#890) to the full $(2m_{\max}+1)$-vector family
$\{(\hat S^\mp_{\rm tot})^k\,|\sigma_{\top/\bot}\rangle:
 k=0,\ldots,2m_{\max}\}$.

\subsection*{DONE: Tasaki Problem 2.2.c}

\textbf{Statement} (Tasaki §2.2, eq.~(2.2.15), p.~24): The
SU(2)-averaged two-site state
\[
  \frac{1}{4\pi} \int_0^{2\pi}\!\!d\varphi \int_0^\pi d\theta \sin\theta\;
    \hat U^{(3)}_{\varphi,{\rm tot}}\,\hat U^{(2)}_{\theta,{\rm tot}}\,
    |\!\uparrow\downarrow\rangle
  \;=\;
  \tfrac{1}{2}\bigl(|\!\uparrow\downarrow\rangle
                   - |\!\downarrow\uparrow\rangle\bigr).
\]

\textbf{Status}: Formalized in
\texttt{Quantum/SU2Integral.lean} as \texttt{problem\_2\_2\_c}.
The proof decomposes into four trigonometric integral lemmas
(\texttt{integral\_cexp\_\{I,neg\_I\}\_mul\_zero\_two\_pi},
\texttt{integral\_sin\_mul\_\{cos,sin\}\_sq\_half\_zero\_pi}) and the
factorization of the global rotation via
\texttt{noncommProd\_mul\_distrib}; see the detailed outline in the
Tasaki §2.2 supplementary subsection above.

\subsection*{Roadmap status note: P1f''' (eqs.~(2.2.12)--(2.2.13))}

The fully general operator-level statement is now formalised in
three pieces:
\texttt{totalSpinHalfRot\{1,2,3\}\_commute\_of\_commute}
($\mathrm{Commute}\,A\,\hat S^{(\alpha)}_{\rm tot}
  \Rightarrow \mathrm{Commute}\,A\,\hat U^{(\alpha)}_{\theta,{\rm tot}}$),
\texttt{totalSpinHalfRot\{1,2,3\}\_conjTranspose\_mul\_self}
(unitarity $\hat U^\dagger\hat U = \hat 1$, via
\texttt{exp\_mem\_unitary\_of\_mem\_skewAdjoint}), and
\texttt{totalSpinHalfRot\{1,2,3\}\_conj\_eq\_self\_of\_commute}
(the Tasaki~(2.2.13) finite form, combination of the previous two).
P1f''' is now fully closed without axioms.

\section{Progress table}
\label{sec:progress}
%======================================================================

The table below tracks the status of each planned area. Entries are
updated when a formalization completes, via the project workflow.

\begin{center}
\begin{tabular}{lll}
\hline
Area & Stage & Status \\
\hline
Pauli matrices and their algebra & finite-dim & \emph{formalized} \\
Spin-1/2 operators and $[\hat S^{(\alpha)},\hat S^{(\beta)}]$ (Tasaki §2.1) & finite-dim & \emph{formalized} \\
Basis states, $\hat S^{(3)}$ eigen-eqn, $\hat S^\pm$ (Tasaki §2.1) & finite-dim & \emph{formalized} \\
$S = 1$ matrix representations (Tasaki §2.1 eq.~2.1.9) & finite-dim & \emph{formalized} \\
Pauli basis of $M_2(\mathbb{C})$ (Tasaki Problem 2.1.a, $S=1/2$) & finite-dim & \emph{formalized} \\
$S = 1/2$ rotation (closed form, $\theta=0$, adjoint, $2\pi$) & finite-dim & \emph{formalized} \\
$S = 1/2$ rotation (group law, unitarity) & finite-dim & \emph{formalized} \\
$S = 1/2$ $\pi$-rotations (eq.~2.1.25) & finite-dim & \emph{formalized} \\
$S = 1/2$ $\pi$-rotation products and conjugations (Tasaki §2.1) & finite-dim & \emph{formalized} \\
Abstract lattice + total spin (Tasaki §2.2, partial) & finite-dim & \emph{formalized} \\
SU(2)-averaged singlet, Problem 2.2.c (Tasaki eq.~(2.2.15)) & finite-dim & \emph{formalized} \\
Multi-body operator space / \texttt{onSite} & finite-dim & \emph{formalized} \\
Site commutativity at distinct sites & finite-dim & \emph{formalized} \\
1D open-chain quantum Ising, Hermiticity & finite-dim & \emph{formalized} \\
Gibbs state $\rho = e^{-\beta H}/Z$, $\operatorname{Tr}(\rho)=1$, $\langle\mathbf{1}\rangle=1$, $Z(0)=\dim$, $Z(0)\neq 0$, linearity $\langle O_1+O_2\rangle$, $\langle cO\rangle$ & finite-dim & \emph{formalized} \\
Periodic boundary conditions & finite-dim & \emph{formalized} \\
Quantum Heisenberg chain (open and periodic BC) & finite-dim & \emph{formalized} \\
Quantum Heisenberg lattice (2D / 3D, full Gibbs companion family) & finite-dim & \emph{formalized} \\
Néel state on bipartite chain / 2D / 3D + per-bond expectation $-1/4$ (Tasaki §2.5) & finite-dim & \emph{formalized} \\
Generic graph-centric Néel state \texttt{neelStateOf} (Tasaki §2.5 (2.5.2)) & finite-dim & \emph{formalized} \\
Marshall sign, generic + chain / 2D / 3D, flip + time-reversal bridge (Tasaki §2.5) & finite-dim & \emph{formalized} \\
Spin-$S$ Marshall--Lieb--Mattis Theorem 2.2 (Tasaki §2.5), magnetization-sector form (real and complex) and full-Hilbert-space form & finite-dim & \emph{formalized} \\
Spin-$S$ saturated ferromagnetic state (Tasaki §2.4 generalised): all-up/all-down as Heisenberg eigenvector for any coupling, ladder annihilation $\hat S^\pm_{\rm tot}$, Casimir eigenvalue $(|V|\cdot N/2)(|V|\cdot N/2 + 1)$ & finite-dim & \emph{formalized} \\
Time-reversal map for many-body $S=1/2$ + multi-body Kramers degeneracy (Tasaki §2.3) & finite-dim & \emph{formalized} \\
Multi-mode fermion via Jordan--Wigner (CAR + number op + Hubbard chain) & finite-dim & \emph{formalized} \\
Hubbard chain, finite volume (open and periodic BC) & finite-dim & \emph{formalized} \\
Hubbard SU(2) symmetry: $[\hat{S}^{\pm}_{\mathrm{tot}}, H] = 0$ (Tasaki §9.3.3) & finite-dim & \emph{formalized} \\
Hubbard all-up-spin state: $H_{\mathrm{int}}|\!\!\uparrow\cdots\uparrow\rangle = 0$ (Tasaki §11.1.1) & finite-dim & \emph{formalized} \\
Hubbard hard-core subspace: $H_{\mathrm{int}}\psi=0$ for no-double-occupancy states (unnumbered infrastructure for Tasaki Theorems 11.5 and 11.7; 1st ed., §11.2, pp.~381--388) & finite-dim & \emph{formalized} \\
Partial trace & finite-dim & not started \\
Quantum canonical ensemble wrapper & finite-dim & not started \\
BCS model (mean field) & finite-dim & not started \\
CAR algebra & algebraic & not started \\
Quasi-local C*-algebras, KMS states & algebraic & not started \\
Thermodynamic limit, phase transitions & infinite & not started \\
Lattice QCD & survey & not started \\
\hline
\end{tabular}
\end{center}

%======================================================================
\section{References}
\label{sec:refs}
%======================================================================

\begin{thebibliography}{10}

\bibitem{Tasaki:QMB}
H.~Tasaki.
\emph{Physics and Mathematics of Quantum Many-Body Systems}.
Graduate Texts in Physics, Springer, 2020. Primary reference for
quantum spin systems and for the one-dimensional transverse-field
quantum Ising model (§3.3).

\bibitem{NielsenChuang}
M.~A.~Nielsen and I.~L.~Chuang.
\emph{Quantum Computation and Quantum Information}, 10th Anniversary
Edition. Cambridge University Press, 2010. Used as a cross-check for
the standard Pauli relations: §2.1.3, Figure~2.2 (pp.~65--66) for the
definitions, Ex.~2.19 (p.~70) for Hermiticity, Ex.~2.40 (p.~77) for
the commutator, Ex.~2.41 (p.~78) for $\sigma^2 = I$ and
anticommutation.

\bibitem{FriedliVelenik}
S.~Friedli and Y.~Velenik.
\emph{Statistical Mechanics of Lattice Systems: A Concrete
Mathematical Introduction}. Cambridge University Press, 2017.
Reference for classical lattice models; not yet used at the current
stage.

\bibitem{GlimmJaffe}
J.~Glimm and A.~Jaffe.
\emph{Quantum Physics: A Functional Integral Point of View}. Springer,
1987. Reference for constructive field theory; not yet used at the
current stage.

\bibitem{ArakiMoriya}
H.~Araki and H.~Moriya.
\emph{Equilibrium Statistical Mechanics of Fermion Lattice Systems}.
Reviews in Mathematical Physics, 2003. Planned reference for the
CAR-algebra stage; not yet used.

\bibitem{BruPedra}
J.-B.~Bru and W.~de~Siqueira~Pedra.
\emph{C*-Algebras and Mathematical Foundations of Quantum Statistical
Mechanics: An Introduction}. Planned reference for the algebraic
stage; not yet used.

\bibitem{Simon:LGI}
B.~Simon.
\emph{The Statistical Mechanics of Lattice Gases, Volume~I}. Princeton
University Press, 1993. Reference for classical lattice gases; not yet
used at the current stage.

\end{thebibliography}

\begin{remark}
The bibliography lists every reference named in either
\texttt{CLAUDE.local.md} (local planning notes) or this guide. Entries
labelled \emph{not yet used} are retained so that future work can cite
them without re-establishing the bibliography; the current code uses
only the entries explicitly referenced above.
\end{remark}

\end{document}
